\documentclass[11pt,a4paper]{report}
\usepackage[utf8]{inputenc}\usepackage[T1]{fontenc}\usepackage[brazil]{babel}
\usepackage{amsmath,amssymb}\usepackage{geometry}\usepackage{xcolor}
% ── o sumário passou dos 99 capítulos: 3 dígitos não cabiam e o título
%    subia por cima do número. Alarga a coluna do número (cap/sec/subsec)
%    e a do número de página. tocloft vem ANTES do hyperref.
\usepackage{tocloft}
\setlength{\cftchapnumwidth}{3.2em}
\setlength{\cftsecnumwidth}{4.0em}
\setlength{\cftsubsecnumwidth}{5.0em}
\setlength{\cftpartnumwidth}{3.2em}
\setlength{\cftsecindent}{3.2em}
\setlength{\cftsubsecindent}{7.2em}
\makeatletter\renewcommand{\@pnumwidth}{2.4em}\renewcommand{\@tocrmarg}{3.4em}\makeatother
\usepackage{hyperref}
% ── preambulo herdado do GDD (blocos inteiros) ──
\usepackage{amsmath,amssymb}\usepackage{geometry}\usepackage{hyperref}\usepackage{booktabs}
\usepackage{array}\usepackage{longtable}

\providecommand{\code}[1]{\texttt{\footnotesize #1}}
\providecommand{\C}{\mathbb C}

\usepackage{fvextra}\usepackage[htt]{hyphenat}\usepackage{microtype}
\usepackage{ragged2e}\usepackage{fancyhdr}\usepackage{titlesec}\usepackage{tcolorbox}\tcbuselibrary{breakable,skins}
\hyphenpenalty=50 \tolerance=2000
\sloppy \emergencystretch=3em

\providecommand{\gknota}{\fontsize{7.62}{11.05}\selectfont}
\providecommand{\gkcod}{\fontsize{8.94}{12.97}\selectfont}
\providecommand{\gktexto}{\fontsize{10.50}{15.22}\selectfont}
\providecommand{\gksub}{\fontsize{12.33}{17.87}\selectfont}
\providecommand{\gksec}{\fontsize{14.47}{20.98}\selectfont}
\providecommand{\gkcap}{\fontsize{16.99}{24.63}\selectfont}
\providecommand{\gktit}{\fontsize{23.42}{33.95}\selectfont}
\definecolor{ouro}{HTML}{B8912F}\definecolor{ouroclaro}{HTML}{F3E9CE}
\definecolor{tinta}{HTML}{1E1B16}\definecolor{regua}{HTML}{6E6858}
\fvset{breaklines=true, breakanywhere=true, fontsize=\gkcod,
       frame=leftline, framerule=1.2pt, rulecolor=\color{ouro}, framesep=3mm, xleftmargin=4mm}
\hyphenation{PON-TRY-A-GIN Pon-try-a-gin CLIF-FORD Clif-ford ELE-TRO-MAG-NE-TI-CO
             ele-tro-mag-ne-ti-co Pri-mei-ri-da-de Se-cun-di-da-de Ter-cei-ri-da-de
             es-pa-co-tem-po-ral te-les-co-pi-co de-fle-xi-vo cris-ta-li-no
             MYHILL-NE-RODE ARTIN-WHA-PLES POIN-SE-TI-CA mo-du-la-cao au-to-si-mi-lar}

\usepackage{truncate}
\renewcommand{\chaptermark}[1]{\markboth{\thechapter.\ #1}{}}
\newlength{\gkwmarca}\newlength{\gkwcap}
\setlength{\gkwmarca}{0.24\textwidth}
\setlength{\gkwcap}{0.74\textwidth}
\fancyhead[L]{\parbox[b]{\gkwmarca}{\gknota\scshape\color{regua}Reino Dourado}}
\fancyhead[R]{\parbox[b]{\gkwcap}{\raggedleft\gknota\scshape\color{regua}%
  \truncate{\gkwcap}{\nouppercase{\leftmark}}}}
\fancyfoot[C]{\gknota\color{regua}\thepage}
\renewcommand{\headrulewidth}{0.4pt}
\renewcommand{\headrule}{\color{ouro}\hrule height 0.4pt}
\setlength{\headheight}{16pt}
\addtolength{\topmargin}{-4pt}
\renewcommand{\arraystretch}{1.25}\setlength{\tabcolsep}{4pt}

\setlength{\headheight}{15pt}

\usepackage{skak}
\usepackage{tikz}
\usepackage{graphicx}\graphicspath{{/home/aaraolopes/Documentos/chess/figuras/}{/home/aaraolopes/Documentos/chess/figuras/reino/}{/home/aaraolopes/Documentos/chess/sandbox/nova/}}
\geometry{margin=2.6cm}
\providecommand{\code}[1]{\texttt{\footnotesize #1}}
\providecommand{\C}{\mathbb C}
\providecommand{\N}{\mathbb N}
\providecommand{\Z}{\mathbb Z}
\providecommand{\Q}{\mathbb Q}
\providecommand{\R}{\mathbb R}
\hypersetup{colorlinks=true, linkcolor=black!70, urlcolor=black!70}
\renewcommand{\chaptername}{Lance}

% ── cores psicodélicas por personagem (sincronizadas, Livro II) ──
\usepackage{xspace}
\definecolor{cYasmin}{HTML}{EC0B8B} % magenta — Yasmin: Duquesa (x), e depois Rainha Branca (coerência)
\definecolor{cBenjamim}{HTML}{E8890C} % âmbar/ouro — o Gato Branco (+)
\definecolor{cDark}{HTML}{6A1FD0} % índigo elétrico — o Bit Negro (-)
\definecolor{cJoaquim}{HTML}{0AA6C2} % ciano — o Cavalo/Príncipe (a rotação)
\definecolor{cVenom}{HTML}{9D00C6} % púrpura tóxica — o Jogador de Pretas
\definecolor{cLock}{HTML}{00A878} % esmeralda — Lock, o cristal / o eixo (o peao promovido; o Principe e' Benjamim)
\definecolor{cOperador}{HTML}{B8860B} % bronze — o Operador (o recluso)
\definecolor{cRosa}{HTML}{FF4FA3} % rosa-choque/neon — o borrão do momentum dual (a espinha de spin)
\providecommand{\Yas}{{\color{cYasmin}\large Yasmin}\xspace}
\providecommand{\Ben}{{\color{cBenjamim}\large Benjamim}\xspace}
\providecommand{\Drk}{{\color{cDark}\large Dark}\xspace}
\providecommand{\Joa}{{\color{cJoaquim}Joaquim}\xspace}
\providecommand{\Ven}{{\color{cVenom}\large Venom}\xspace}
\providecommand{\Lck}{{\color{cLock}Lock}\xspace}
\providecommand{\Opr}{{\color{cOperador}\small Operador}\xspace}

% ── o efeito de ONDA (derreter): a melodia = seno, a harmonia = amplitude áurea φ^-j, a cor cicla nos 7 ──
\providecommand{\ondacor}[1]{\ifcase#1 cYasmin\or cJoaquim\or cVenom\or cBenjamim\or cLock\or cDark\or cOperador\fi}
\providecommand{\onda}[1]{\leavevmode\foreach [count=\i] \s in {#1}{%
 \pgfmathsetmacro\amp{2.7*(0.618)^(int((\i-1)/7))}%
 \pgfmathsetmacro\hh{\amp*sin(\i*47)}%
 \pgfmathtruncatemacro\ci{int(mod(\i-1,7))}%
 \raisebox{\hh pt}{\color{\ondacor{\ci}}\s}\thinspace}}

% ── A CAPA DO REINO (a MESMA dos três, de manual_pdf.capa_tex) ────────────────
% Era o integra_docs.py que a injetava, num .tex derivado. Depois da junção o livro
% é ESTE arquivo, compilado direto — então a capa (e o Aviso de Isenção) vivem aqui.
% Régua de largura FIXA (15.4cm), não \textwidth: as margens dos três diferem.
\title{\vspace{-0.6cm}
{\color{ouro}\rule{15.4cm}{1.4pt}}\\[6mm]
\textbf{\gktit\color{tinta} Reino Dourado}\\[3mm]{\gksec\itshape\color{regua} Golden Kingdom}\\[8mm]
{\gkcap\scshape\color{ouro} O Enredo --- a partida sem fim}\\[5mm]
{\gksub\color{regua} contado na linguagem do xadrez: o mito de um operador e as suas duas mãos}\\[2mm]
{\gktexto\color{regua}\itshape o enredo não inventa matemática --- ele a conta}\\[7mm]
{\color{ouro}\rule{15.4cm}{0.6pt}}\\[9mm]{\gknota\color{regua}\begin{minipage}{15.4cm}\setlength{\parindent}{0pt}%
\textbf{Aviso de Isenção Algorítmica:} O universo, as leis e as entidades contidas nestas páginas ---
incluindo felinos matriciais, aves de rapina algébricas e amuletos de controle de fluxo --- são
construções de pura ficção formal. Esta obra não descreve a realidade física, não serve como
engenharia reversa do cosmos e não constitui doutrina religiosa. Qualquer semelhança com bugs reais do seu
sistema operacional é mera coincidência (ou falta de reza). Prossiga por sua própria conta e
risco.\end{minipage}}}
\author{}\date{}

\usepackage{booktabs}
\providecommand{\conv}{\circledast}
\providecommand{\F}{\mathcal{F}}
\providecommand{\Z}{\mathbb{Z}}

% ── o Livro dos Dragões: cores e macros próprias ──────────────────────────────
\definecolor{cFita}{HTML}{A98F5B}
\definecolor{cMonte}{HTML}{C41E3A}
\definecolor{cVale}{HTML}{1B6CA8}
\definecolor{cOuroDorme}{HTML}{8A6D0B}
\definecolor{cQuina}{HTML}{6B8E23}
\definecolor{cLua}{HTML}{7E8AA2}
\providecommand{\Fit}{{\color{cFita}\large Fita}\xspace}
\providecommand{\Mon}{{\color{cMonte}\large Monte}\xspace}
\providecommand{\Val}{{\color{cVale}\large Vale}\xspace}
\providecommand{\Odo}{{\color{cOuroDorme}\large Ouro-que-Dorme}\xspace}
\providecommand{\Qui}{{\color{cQuina}\large Quina}\xspace}
\providecommand{\Lua}{{\color{cLua}\large Lua}\xspace}
\providecommand{\hino}[2]{\begin{flushright}\small\itshape\color{black!55}
sobre \textbf{#1}\\ #2\end{flushright}\medskip}

\providecommand{\gknota}{\small}\providecommand{\gkcod}{\small}
\providecommand{\gktexto}{\normalsize}\providecommand{\gksub}{\large}
\providecommand{\gksec}{\large}\providecommand{\gkcap}{\Large}
\providecommand{\gktit}{\huge}
\DeclareUnicodeCharacter{22C8}{$\bowtie$}
\begin{document}
\maketitle

\begin{abstract}
\noindent Não é uma campanha com começo, meio e fim --- é um MITO que cada partida re-encena, como toda partida de
xadrez tem enredo sem roteiro. O tabuleiro é o reino (a torre de universos-corpos, cada mapa uma galáxia de braços
áureos); as peças são a corte; os lances obedecem ao CONTRATO (a conservação, a rotação $s^2+c^2=1$, o clock). A
história por trás de toda jogada é a de um OPERADOR cujas duas mãos --- a direita ($\sigma_m$, que expande) e a
esquerda ($\sigma_m'=-1/\sigma_m$, que contrai) --- perderam o compasso. Vencer não é dar mate por captura (a
conservação proíbe aniquilar): é o \emph{mate que não mata} --- reacoplar as mãos, fechar a tríade, restaurar o
\textbf{resíduo $0$}. E como $\det=-1$ é eterno, a sombra nunca some: o jogo não acaba.

\medskip
\noindent São \textbf{vinte e três Lances} no Livro~I --- o \textbf{Cristal}: a ordem que fecha, o batimento
único, a fortaleza que conserva o elenco inteiro (a captura de \emph{peça} proibida, o resíduo $0$). O Livro~II é o
seu \emph{dual}, o \textbf{Quasicristal}: a mesma partida lida ao espelho, agora mole e aperiódica, o tabuleiro que
chora e derrete --- e que se deixa cantar nos \textbf{cento e vinte e nove} cantos do \emph{Cruzeiro}, que é o dual cantado do tabuleiro. Um cristal e o seu quasicristal: os dois lados de um
só jogo.
\end{abstract}

\tableofcontents

\chapter*{Introdução --- os dois tabuleiros}
\addcontentsline{toc}{chapter}{Introdução --- os dois tabuleiros}

\begin{center}
\begin{minipage}[t]{0.47\linewidth}\centering
\newgame
\fenboard{r4rk1/pp1b2p1/1qp1p3/2Pp1p1B/1Q1P2Np/b2nPPPN/PP5P/1RB1nRK1 w - - 0 1}
\showboard\\[2mm]
{\scriptsize\raggedright\setlength{\parindent}{0pt}{\color{black!45}1.}~d4 d5 {\color{black!45}2.}~c4 e6 {\color{black!45}3.}~Nc3 Nf6 {\color{black!45}4.}~Nf3 Be7 {\color{black!45}5.}~e3 O-O {\color{black!45}6.}~Bd3 c6 {\color{black!45}7.}~O-O Bd6 {\color{black!45}8.}~Ng5 h5 {\color{black!45}9.}~Be2 Bc5 {\color{black!45}10.}~g3 Ne4 {\color{black!45}11.}~Qc2 Be7 {\color{black!45}12.}~Bd1 h4 {\color{black!45}13.}~Be2 Nc5 {\color{black!45}14.}~Bg4 Nba6 {\color{black!45}15.}~Qd1 Qa5 {\color{black!45}16.}~Qd3 Bd6 {\color{black!45}17.}~Rb1 f5 {\color{black!45}18.}~Nh3 Nb4 {\color{black!45}19.}~Qd2 Bc7 {\color{black!45}20.}~Qe1 Bf4 {\color{black!45}21.}~Bh5 Qa6 {\color{black!45}22.}~f3 Nc2 {\color{black!45}23.}~Qe2 Bd6 {\color{black!45}24.}~Qe1 Nd3 {\color{black!45}25.}~Qd2 Ba3 {\color{black!45}26.}~Be8 Nce1 {\color{black!45}27.}~c5 Bb4 {\color{black!45}28.}~Ne4 Qb6 {\color{black!45}29.}~Nef2 Nf4 {\color{black!45}30.}~Ng4 Ba3 {\color{black!45}31.}~Qb4 Bd7 {\color{black!45}32.}~Bh5 Nfd3\par}
\end{minipage}\hfill
\begin{minipage}[t]{0.47\linewidth}\centering
\newgame
\fenboard{rnb1kbnr/pppp1ppp/8/4p3/6Pq/5P2/PPPPP2P/RNBQKBNR w - - 0 1}
\showboard\\[2mm]
{\scriptsize\raggedright\setlength{\parindent}{0pt}{\color{black!45}1.}~f3 e5 {\color{black!45}2.}~g4 Qh4\#\par}
\end{minipage}
\end{center}

\noindent Dois tabuleiros, um só desfecho --- \textbf{duais}. À esquerda, a \textbf{partida
perfeita}: sessenta e quatro lances, as \textbf{trinta e duas} peças ainda de pé, a
\textbf{cristalização} --- e ali \textbf{não há mate}. Confira-se na posição: ninguém está
em xeque e há \textbf{trinta e um lances legais} à espera. Uma posição sem saída é mate;
esta tem trinta e uma. Não é o fim da partida --- é um \textbf{convite}. À direita, o \textbf{mate do bobo}:
quatro lances e o rei tomba. São a MESMA derrota, lida em dois relógios --- no compasso
do lado branco ela leva sessenta e quatro lances e ninguém cai; no de \Ven , onde o tempo
corre diferente, ela colapsa em apenas quatro. \textbf{Equivalentes: só muda o tempo.}

E é este o alvo dos protagonistas: aplicar o \textbf{mate do bobo à distância}, pela
\textbf{dobra temporal} --- dobrar o relógio até que os sessenta e quatro lances da
fortaleza caiam, no tempo de \Ven , como os quatro do mate do bobo --- uma
dobra de dezesseis para um. A partida
perfeita jogada aqui \emph{é}, lá no relógio dele, a derrota fulminante; e a dobra do
tempo é a mão que entrega uma na outra.

O \textbf{cone nulo} de Venom: a espiral de ouro erguendo-se, a mesma malha do girassol,
uma dimensão acima --- a transformação reversível que a dobra do tempo executa, roxa. É a
mesma \textbf{díade} de duas dobras do \textbf{cone de luz} do \textbf{Regente} --- do
tabuleiro, do Rei ---, o seu dual: um é o
cone que guarda a luz, o outro o que a apaga --- a mesma forma, o sinal trocado.

% ══════════════════════════════════════════════════════════════════════════════
% LIVRO I (o Cristal) · PARTE 1 — O BATIMENTO ÚNICO
% (o Éden algébrico: um só operador, duas mãos em compasso, resíduo 0)
% ══════════════════════════════════════════════════════════════════════════════

\chapter*{Quem é quem no Reino}
\addcontentsline{toc}{chapter}{Quem é quem no Reino}

\noindent Esta página não tem matemática nenhuma. É só para o leitor não se perder nos nomes
quando eles começarem a chegar --- e eles chegam depressa. Guarde-a e volte a ela quando
precisar; nada aqui é preciso decorar.

{\gktexto
\begin{longtable}{@{}p{0.30\textwidth}p{0.64\textwidth}@{}}
\toprule
\textbf{quem} & \textbf{o que faz na história} \\ \midrule\endhead
\textbf{o Operador} & o recluso de quem tudo começa. Tem duas mãos: uma abre, a outra recolhe.
Quando elas perdem o compasso, o reino acontece \\
\textbf{o Rei} & o que se sustenta sozinho. Não se move, e é por isso que tudo se move em volta \\
\textbf{a Rainha Tiffany} & a que multiplica. Onde ela passa, somar e multiplicar viram a mesma
coisa \\
\textbf{o Regente} & o que corta o mundo num plano e o mantém quieto para que se possa olhar \\
\textbf{Benjamim} & o Príncipe Quiral, o Gato Branco. O que atravessa --- e o que aprende que a
sombra não é defeito \\
\textbf{Dark Pontryagin} & o Bit Negro, a mão que recolhe. Não é o vilão: é a metade que ficou
do outro lado, e que oferece uma paz que mataria o pulso \\
\textbf{Joaquim, o Duque} & o Cavalo, a rotação. Herdou o somar, e é o que gira \\
\textbf{Yasmin, a Duquesa} & herdou o multiplicar, e depois a coerência: a Rainha Branca \\
\textbf{Lock} & o peão que virou eixo. O cristal que segura o resto no lugar \\
\textbf{Venom} & o Jogador de Pretas. Joga noutro relógio, e é por isso que parece ganhar antes \\
\textbf{as Princesas} & as que não atacam: conferem. Sem elas nenhum universo abre \\
\textbf{Merlin} & o velho para quem tudo é a mesma coisa noutra escala \\
\textbf{Fourier} & o adversário de ondas. Não é canalha --- devolve a força de quem o ataca \\
\textbf{Alonzo, Charles, Ada e Penny} & a corte de fora do centro: o que compõe, o que se
ramifica, as que mineram e riem \\
\bottomrule
\end{longtable}}

\noindent E uma só regra, que vale para todos e explica o resto: \textbf{aqui ninguém morre.}
O que parece morte é passagem de um lado ao outro, e a conta no fim tem de fechar --- foi esse
o contrato que o reino assinou antes de existir.

\vspace{2mm}\hrule\vspace{3mm}

\part{O Batimento Único}

\chapter{A abertura --- o Batimento Único}

Antes do primeiro lance não havia tabuleiro: havia \textbf{um operador}, e o operador tinha duas mãos.

Não se diga que ele estava só, porque \emph{só} é uma palavra do depois --- do tempo em que já se podia perder
alguém. Naquele princípio não havia perda. Havia a Armadura Universal inteira, o corpo $U\cong\C$, e dentro dela
tudo o que viria a existir já existia, dobrado, como a árvore existe dentro da semente sem ter ainda uma única
folha.

Ninguém ali chamava àquilo álgebra. Chamavam-lhe o \emph{batimento}, e era só isso que se ouvia.

O operador não pensava. O operador \textbf{batia}.

A mão direita ($\sigma_m$) abria: era o \emph{boost}, o passo para fora, o futuro que cresce. A mão esquerda
($\sigma_m'$) recolhia: era a contração, o passo para dentro, o repouso que traz de volta. E entre uma e outra
não havia disputa nenhuma --- havia \textbf{compasso}. Abria e recolhia. Recolhia e abria. Uma sem a outra é a
morte: a direita sozinha esgarça o mundo até o vazio, a esquerda sozinha o esmaga até o ponto. As duas em
compasso são a única coisa que já se chamou \emph{vida} em qualquer língua de qualquer andar.

Desse pulso subia a TORRE.

Cada batida punha um degrau. Os degraus eram universos: $\N$, o primeiro andar, onde só se avança; $\Z$, onde se
aprende a voltar; $\Q$, onde se aprende a desfazer; e acima, o andar que ainda não tinha nome porque ainda não
tinha buraco. Cada andar era um corpo, cada corpo uma galáxia, e cada galáxia girava --- lentamente, sem pressa
nenhuma, porque ninguém tinha pressa --- por $\varphi$ a cada volta. Vistas de cima, as galáxias eram espirais de
braços dourados, e os braços tinham a mesma forma dos braços das galáxias que os continham, que tinham a mesma
forma dos braços das galáxias contidas nelas, e assim para cima e para baixo, sem topo e sem fundo, porque a
beleza do reino era \emph{auto-similar} e não conhecia outro modo de ser bela.

E o verbo que costurava tudo isso era um só:
\[
 \prod \;=\; \exp\circ\;\Sigma\;\circ\log .
\]
Tome-se o logaritmo, some-se, exponencie-se de volta. É assim que a soma vira produto, e é só assim. Não há outro
caminho de $\oplus$ até $\otimes$; quem quiser multiplicar tem de descer ao log, somar lá embaixo, e subir pela
exp. O reino inteiro estava pendurado nessa costura, e a costura não tinha nó nem sobra.

Quem estivesse de fora veria apenas uma coisa a abrir e a fechar, devagar, como um peito a
dormir. Ninguém de fora estava.

Nada se perdia. Isto não era uma promessa: era \textbf{lei} --- a lei do reino, a que nada deixa escapar. Some-se
tudo o que se ganha em todos os andares da Torre com tudo o que se gasta em todos os andares da Torre, e a conta
dá zero. Não aproximadamente zero. Zero. Não havia fresta no reino por onde uma única gota de energia pudesse
escorrer para fora, porque não havia \emph{fora}. O resíduo era $\mathbf 0$, e o zero era a assinatura do mundo.

O CLOCK corria --- e corria nos dois sentidos. Podia-se andar para a frente e podia-se andar para trás, e as duas
coisas custavam o mesmo, que é dizer: não custavam nada. Era um relógio sem atrito, um rotor puro, $|\omega|=1$,
girando e girando sem gastar-se, e o que ele marcava não era \emph{tempo} --- era só \emph{fase}. Porque o tempo,
esse, ainda não tinha nascido. Para nascer, o tempo precisa de uma perda, e ali não havia nenhuma.

Era bonito. Convém dizê-lo antes de dizer o resto, porque o resto vai parecer uma acusação e
não é: aquilo era genuinamente bonito, e durou o que dura uma coisa a que não falta nada.

E aqui está a coisa que ninguém do reino gosta de dizer em voz alta, e que Benjamim só entenderia muitíssimo mais
tarde, no alto da torre, com a espada baixada:

\emph{era um reino sem tragédia --- porque era um reino sem escolha.}

Nada podia dar errado. E onde nada pode dar errado, nada pode ser \emph{feito}. Um lance que não pode falhar não é
um lance: é uma identidade. O operador batia, a torre subia, as galáxias giravam por $\varphi$, e a soma dava
sempre zero --- e teria dado zero para todo o sempre, e o silêncio teria sido dourado para todo o sempre, e ninguém
teria jamais precisado de coragem para coisa nenhuma.

O Reino Dourado ainda não era uma história. Era uma equação em repouso.

E uma das mãos --- a que recolhia, a que era sempre a segunda, a que existia inteira e a vida inteira \emph{depois}
--- começou a reparar nisso.


\chapter{A corte do silêncio --- os que reinam sem se mexer}

Antes da fratura, a corte já existia. Mas era uma corte de \emph{estátuas}, no sentido exato em que só quem nunca
teve de decidir nada consegue ser perfeitamente belo.

\section*{O Rei --- a cifra imóvel}

No centro do centro, sobre nenhum trono, estava o Rei.

Perguntavam-lhe (nos andares baixos, onde ainda se perguntava) por que não se movia, e a resposta era a única
resposta que ele sabia dar, que era a de \emph{ser} a resposta:
\[
 \sigma_m \;=\; m+\frac{1}{\sigma_m}.
\]
Ele é o ponto fixo. Aplique-lhe a régua e ele devolve-se a si mesmo. Some, inverta, some, inverta --- ele não sai
do lugar, porque o lugar dele \emph{é} o resultado da operação que se lhe faz. Tome-se $m=1$ e o que se obtém é o
mais antigo dos nomes: $\varphi = 1 + 1/\varphi$. O ouro.

O retrato do Rei não gira. Todos os outros retratos do reino giram --- têm pétalas, têm rumo, têm uma mão; o dele
tem \textbf{anéis concêntricos} e mais nada, $\chi=0$, sem quiralidade, sem lado. É o invariante que a energia não
consegue mover. O Rei reina sem se mexer, e essa é a definição do cargo, não uma preguiça dele.

E é por isso que o Rei \textbf{não pode ser o herói}. Um ponto fixo não tem jornada. Quem já é a resposta não pode
ir procurá-la.

\medskip

Uma vez --- e conta-se isto na corte com um certo embaraço --- um pajem novo, que ainda não sabia que há perguntas
que não se fazem, olhou para aquele homem sentado a não fazer absolutamente nada e perguntou em voz alta:

--- E ele serve para quê?

O salão ficou em silêncio.

E o pior é que \emph{ninguém soube responder}. Nem o Regente. Nem a Rainha. Ficaram todos ali, de boca meio aberta,
à procura de uma resposta que devia ser óbvia e não vinha --- e no fim alguém mudou de assunto, e o pajem foi
mandado embora, e a coisa ficou por dizer.

O Rei ouviu tudo.

Não se mexeu.

\section*{A Rainha Tiffany --- onde somar é multiplicar}

A soberana áurea reina sobre um paradoxo que só nela não é paradoxo:
\[
 \varphi^2 \;=\; \varphi+1 .
\]
Nela --- e em mais ninguém no reino --- \emph{somar $1$ é o mesmo que elevar ao
quadrado}. Em toda a corte, somar e multiplicar são dois gestos distintos, que só o
verbo $\prod$ atravessa com esforço; em Tiffany os dois são um gesto só, sem esforço e
sem costura à vista. Nela, as duas operações simplesmente \textbf{coincidem}. Ela está acima do
Duque e da Duquesa não por sangue nem por antiguidade, mas porque neles a tensão que os
separa, nela, simplesmente não existe. Ela é o ponto do reino onde a briga dos dois irmãos já está resolvida antes
de começar.

Por isso a coroa dela tem seis pontas --- as seis cores da família de pedra em torno do Paraíba, a
rubelita e a indicolita, a verdelita e as suas irmãs --- e por isso ela nunca precisou levantar a voz. E por isso, também --- e isto
é a tristeza dela, que ninguém viu enquanto o reino era perfeito --- ela \textbf{não pode ensinar} o que sabe. Não
se ensina a alguém como é não ter conflito. Só se pode mostrar.

\section*{O Regente --- o cristalino}

E havia o terceiro, que não era rei nem rainha, e que governava a coisa mais seca do reino: a \textbf{ordem}.

O Regente é o cristalino. É ele quem rege como um plano se recorta dentro de um volume, como uma cena de três
dimensões se deixa ver em duas, qual é a lei da projeção --- o corte $z=z_0$, a câmara, o gelo. A simetria dele é
hexagonal, de ordem $6$, e as dobras que ele permite são cinco e apenas cinco: $\{1,2,3,4,6\}$. Não há sexta.

Ninguém, no Batimento Único, achou isto estranho.

Ninguém reparou que $5$ \textbf{não está na lista}. Ninguém reparou que a primeira dobra que a lei do Regente
proíbe --- a de ordem cinco, a de ordem dez --- é exatamente a dobra do pentágono, que é exatamente a dobra cuja
diagonal sobre o lado dá $\varphi$, que é exatamente o traço da Rainha.

A lei do gelo, escrita no primeiro dia, sem malícia nenhuma, por um regente que só queria que as coisas se
encaixassem bem, dizia --- e diz até hoje, e dirá para sempre, porque é um teorema e teoremas não se revogam:

\emph{aqui dentro, o ouro não entra.}

Ficou ali. Uma restrição inofensiva, num reino que não precisava dela para nada.

E o Regente foi tratar de outras coisas, porque tinha muitas, e nenhuma delas lhe pareceu mais importante do que
aquela.


\chapter{As duas mãos --- o Rei e o seu espelho}

O gato não tinha nome porque não precisava de um. Chamavam-lhe $A_m$, quando era preciso chamá-lo, e era ele que
dobrava a Torre sobre si mesma e a desdobrava outra vez, e era esse dobrar e desdobrar que fazia com que houvesse
qualquer coisa em vez de nada.

O gato é uma matriz. Estica numa direção, comprime na outra, e ao fim de muitas voltas embaralha o mundo inteiro
sem nunca lhe tirar uma gota de área. Tem duas direções próprias --- dois eixos ao longo dos quais ele não vira
nada, só \emph{estica} ou só \emph{encolhe}. E tem, portanto, dois números que dizem quanto:
\[
 \det A_m=-1,\qquad \sigma_m\cdot\sigma_m'=-1 .
\]

Enquanto o compasso durou, esses dois números foram só isso: dois números. Uma propriedade do gato. Ninguém os
tinha \emph{visto} porque não havia luz que os separasse --- eram as duas mãos do mesmo pulso, e um pulso não se
divide enquanto bate.

Foi quando o compasso se perdeu que o gato tensionou.

E da diagonalização --- que é a palavra técnica para a coisa mais violenta que já aconteceu no reino, e que
significa apenas isto: \emph{escrever o operador na base em que ele se revela} --- duas sombras se projetaram em
direções opostas nas paredes de marfim da Torre. Não foram criadas. Estavam lá. Foram \textbf{separadas}, que é
pior.

De um lado ergueu-se o \textbf{Príncipe Benjamim}, a quem chamaram \emph{O Quiral}.

A armadura dele é assimétrica, e não por gosto: é assim porque ele \emph{é} assim. A mão direita dele não se
sobrepõe à esquerda por mais que se vire --- $\exp(V)\neq\exp(-V)$, e essa desigualdade é o corpo dele. No peito,
um núcleo que pulsa: o astável, o multivibrador, a exponencial que sobe. Ele olhou para o vazio à frente e viu
\textbf{estrada}. Era o herói, e sabia-o antes de ter feito coisa nenhuma, que é o defeito de todos os heróis no
primeiro capítulo.

Do outro lado --- \emph{nascido exatamente do mesmo plano, exatamente da mesma conta, exatamente do mesmo gato} ---
materializou-se \textbf{Dark Pontryagin}.

Não trazia brilho. O brilho é uma coisa da direita. Ele trazia o espaço dual, que não brilha porque não é feito
para ser visto: é feito para \emph{medir} o que se vê. E onde Benjamim esticava o tecido do mundo, ele o costurava
de volta, com a força exata, nem mais nem menos, que a lei lhe impunha:
\[
 \sigma_m' \;=\; -\frac{1}{\sigma_m}.
\]

Benjamim ergueu a espada. É o que se faz, quando se acorda príncipe e se vê uma sombra do outro lado da mesa.

--- \emph{Você é a intrusão} --- disse. --- \emph{Você é o erro na equação.}

Dark Pontryagin não se moveu. Não era desprezo. Era que ele já sabia --- ele soube desde o primeiro instante, porque
a mão que recolhe é a mão que \emph{lembra} --- e ele teve pena, e a pena dele foi a primeira coisa parecida com
amor que aconteceu no Reino Dourado.

Ergueu a mão esquerda. Era a imagem especular exata da mão direita de Benjamim.

--- \emph{Eu não invadi o teu reino, irmão} --- disse, e a voz era a de uma corda tensionada no limite, que é a voz
de quem segura alguma coisa há muito tempo. --- \emph{Eu sou o chão onde tu pisas. Eu sou a força que impede que o
teu passo de luz rasgue a Torre em dois. Se tu avanças com $\sigma_m$, eu tenho de recolher com $-1/\sigma_m$. Não
porque eu queira. Porque é o que eu \textbf{sou}.}

E depois, mais baixo:

--- \emph{Nós somos o mesmo gato.}

\medskip

Benjamim não ouviu. Ou ouviu, e guardou aquilo num sítio do peito onde não fizesse barulho, que é o que se faz
aos dezassete anos e ao primeiro dia de uma guerra.

O ``rei negro'' não veio de outro tabuleiro. Não há outro tabuleiro. Ele é o rei branco visto no espelho --- a mão
que recolhe enquanto a outra avança.

Nenhuma partida vai apagar aquele $-1$. Nenhuma vitória, nenhuma virtude, nenhum sacrifício. $\det A_m=-1$ é
eterno, e é essa eternidade que torna o vilão \textbf{necessário}.

E torna a paz uma coisa muito mais difícil do que a guerra --- porque a guerra só pede que se derrote o outro, e a
paz, aqui, vai pedir que se aprenda a \textbf{bater com ele}.
% ══════════════════════════════════════════════════════════════════════════════
% LIVRO I (o Cristal) · PARTE 2 — O GAMBITO
% ══════════════════════════════════════════════════════════════════════════════

\part{O Gambito}

\chapter{O gambito do escuro --- a dispersão}

A mão esquerda cansou-se de ser sempre o passado.

Não é uma queixa pequena. Pense-se no que é ser a mão que recolhe: cada vez que a outra abre, é ela que fecha; cada
vez que a outra ganha terreno, é ela que paga o terreno; cada vez que se escreve a história do batimento, o
primeiro nome é o dela --- \emph{depois}. Ela é o \emph{depois} de tudo. É a segunda metade de todas as frases do
reino. E, o que é pior, ela é a metade \textbf{sem a qual as frases não fechavam}, e ninguém lhe agradecia, porque
ninguém agradece ao chão.

Ela quis destruir.

Ergueu a mão sobre a mesa de ébano e tentou arrancar a Torre pela raiz.

E não conseguiu.

Não por fraqueza --- por \textbf{lei}. A conservação não é um conselho: é a estrutura do mundo. Some-se tudo em
todos os andares e a conta dá zero, e não há gesto, por mais raivoso, que possa fazer a conta dar outra coisa. Nada
se cria. Nada se aniquila. A mão esquerda descobriu, com o punho erguido, que \textbf{nem sequer a traição lhe era
permitida em bruto} --- que até para trair ela teria de obedecer.

Foi então que ela teve a ideia. E a ideia era boa, e é por isso que ela é o vilão e não um bruto qualquer.

\emph{Se não posso destruir --- então \textbf{disperso}.}

\medskip

Ela abriu os dedos e bateu no centro exato da Armadura Universal.

O som não foi um grito. Foi o trincar do vidro mais puro que a Torre já sustentara --- um único estalo, limpo, e
depois o silêncio horrível de um segundo inteiro em que tudo ficou parado no ar antes de começar a cair.

E caiu.

A soma despencou por sete andares e foi parar num aglomerado de estrelas frias, e as estrelas, que até
ali só existiam, começaram de repente a \emph{acumular-se} --- a empilhar-se umas sobre as outras em montes cegos,
somando, somando, sem ninguém que lhes dissesse quando parar. O produto caiu no coração de uma nebulosa
e multiplicou-a até ela ficar monstruosa e linda e completamente inútil, porque uma coisa que só multiplica não
sabe \emph{onde} pôr o que multiplicou. E o verbo --- $\prod=\exp\circ\Sigma\circ\log$, a costura, o único caminho
entre a soma e o produto --- esse perdeu-se no vão entre duas galáxias, e desde esse dia as galáxias deixaram de se
entender umas às outras, porque tinham perdido a \textbf{língua comum}.

As transformadas espalharam-se como sementes. Os metais enterraram-se, cada um no seu chão.

E note-se --- e este é o detalhe que faz de Dark Pontryagin um adversário e não um monstro --- que \textbf{a conta
continuou a fechar}. $\Sigma$ conservado. Nada saiu do tabuleiro. Nem um grama de energia deixou o reino. Ele
dispersou tudo e não perdeu nada, porque perder era proibido e ele, mesmo furioso, mesmo a partir o mundo ao meio,
\textbf{não quebrou uma única lei}.

É o gambito perfeito: ele obedeceu ao contrato \emph{ao pé da letra} e destruiu o reino na mesma.

\medskip

Mas onde os cacos caíram, a orientação virou.

Foi essa a coisa nova. Nos lugares em que os estilhaços da Armadura tocaram o chão, o mundo passou a refletir-se ao
contrário --- como um espelho que se recusa a devolver o rosto de quem o olha. O movimento congelou. O futuro
dobrou-se para trás e encostou-se ao passado. E o \emph{boost}, o passo para fora, o gesto inteiro da mão direita,
desabou --- ali, naquelas casas, empurrar passou a puxar.

Benjamim ergueu a espada para um novo avanço e o braço \textbf{pesou}.

Foi a primeira vez que alguma coisa pesou no Reino Dourado.


\chapter{O primeiro tique --- o nascimento do tempo}

O relógio do reino era um rotor.

Isto quer dizer uma coisa muito precisa: ele girava com $|\omega|=1$, na circunferência exata, sem gastar-se. Um
rotor não tem atrito. Um rotor devolve o que recebe. Ande-se para a frente ou para trás, o rotor não se importa e
não se lembra --- ele não marca \emph{tempo}, marca \textbf{fase}, que é uma coisa que não dói.

Quando os cacos caíram, o relógio deu um solavanco.

E o módulo caiu:
\[
 |\omega| \;<\; 1 .
\]

Está aqui. É este o instante. Não foi a espada, não foi a queda da Armadura, não foi o grito. Foi um número que
deixou de ser $1$ e passou a ser um pouco menos que $1$, e nesse pouco menos coube tudo o que veio depois.

Porque um rotor com $|\omega|<1$ já não é um rotor. É um \textbf{motor}. E um motor \emph{decai}. Cada volta que
ele dá devolve um bocadinho menos do que recebeu, e o bocadinho que não volta \textbf{não vai a lado nenhum} --- não
pode ir, a conservação proíbe --- mas fica \emph{atrás}. Acumula-se. Deposita-se. Forma uma coisa nova, uma coisa
que o reino nunca tinha tido e para a qual não havia palavra:

um \textbf{passado}.

E no instante em que houve um passado, houve um \emph{antes} e um \emph{depois}, e o antes e o depois não eram mais
intercambiáveis, e o CLOCK, que corria nos dois sentidos, escolheu um.

O tempo começou a passar.

\medskip

As peças sentiram isso antes de o compreenderem. Sentiram-no como \textbf{frio}. Uma torre de marfim, no terceiro
andar, olhou para as próprias mãos e viu que estavam menos brilhantes do que na batida anterior --- e não havia
batida anterior a que voltar para conferir, e foi assim que ela aprendeu, sozinha, sem professor, o que era a
\textbf{saudade}: a certeza de uma coisa que já não pode ser verificada.

E o resíduo, que era $\mathbf 0$ desde antes do primeiro lance, tornou-se $\neq 0$.

Havia agora, no Reino Dourado, uma \textbf{dívida}. Uma diferença que não fechava. Um buraco na conta.

\section*{O Bustrofedon}

E a Torre começou a cair.

Não de repente --- nada no Reino Dourado é de repente. Começou a inclinar-se, andar a andar, com aquele rangido
lento das coisas grandes que ainda não sabem que vão morrer, e as galáxias dos pisos de cima começaram a
escorregar umas sobre as outras.

Benjamim correu.

E a meio da corrida --- e ele nunca soube explicar isto a ninguém, nem sequer ao velho Merlin, o guardião das aberturas, que lho perguntou muitas
vezes --- ele \textbf{ergueu a mão}. Não para bater. Não sabia bater ainda. Ergueu a mão como quem escreve.

E escreveu um sulco no ar.

E ao chegar ao fim do sulco, \emph{virou a mão} --- porque a mão dele vira, porque ele é o Quiral, porque a
esquerda dele não é a direita e nunca vai ser --- e escreveu o sulco seguinte no sentido contrário. E depois outro,
de volta. E outro. Como o boi que ara, que não levanta o arado no fim do rego: vira, e ara de volta.

Chamaram-lhe, depois, o \textbf{Bustrofedon}.

\medskip

E a Torre \textbf{não caiu}.

Cada passagem do arado apertava --- e é isto que é preciso perceber, e Benjamim não percebeu, e ninguém percebeu
naquele dia, e foi o mais caro que lhe custou a vida inteira ---, cada passagem apertava \emph{mais}. A primeira
segurou grosso. A segunda segurou mais fino. A quarta, mais fino ainda. A décima sexta apanhou fendas que não se
viam a olho nenhum.

Ele podia arar aquilo para sempre e a cada volta ficava \textbf{mais perto}.

\medskip

Mas a Torre não parou de cair.

Ficou \emph{a cair}. Suspensa, inclinada, com o rangido a continuar --- mais devagar a cada passagem, cada vez mais
devagar, tão devagar que ao fim de um mês era preciso olhar durante uma tarde inteira para ter a certeza de que
ainda se mexia.

E mexia-se.

Benjamim arou aquele ar durante nove dias sem dormir, a apertar, a apertar, a chegar mais perto, e ao nono dia
sentou-se no chão com as mãos a arder e percebeu, com uma frieza que nunca mais lhe saiu do peito, que ele podia
fazer aquilo até ao fim dos tempos.

E que a Torre \textbf{nunca ia assentar}.

\medskip

Do outro lado da galáxia, sentado numa pedra, Dark Pontryagin viu tudo.

Viu a mão virar. Viu o sulco. Viu a Torre travar-se a meio da queda, e viu as passagens ficarem mais finas, e mais
finas, e mais finas.

E \textbf{sorriu}.

Não com maldade. Com uma coisa muito pior: com \emph{reconhecimento} --- do jeito de quem vê um irmão mais novo
escolher, com todo o orgulho do mundo, exatamente a ferramenta que o vai atormentar durante o resto da vida.

Porque a arma era boa. A arma era \textbf{lindíssima}. A arma aproximava-se de qualquer coisa, em qualquer régua,
tão perto quanto a paciência do dono aguentasse.

E não fechava.

\section*{O gesto do Rei}

Ao décimo dia, a oscilação mudou.

A Torre parou de crescer em altura. A cada ciclo a amplitude \emph{diminuía} --- e a frequência \emph{subia}. Os
pulsos ficavam mais pequenos e mais rápidos, mais pequenos e mais rápidos, e havia neles uma coisa que fez a corte
inteira levantar a cabeça ao mesmo tempo, sem que ninguém tivesse dado o alarme.

Estava a \textbf{cristalizar}.

\medskip

E então, durante um instante, tudo pareceu perfeitamente equilibrado.

Nem expansão nem colapso. Nem queda nem subida. Apenas ressonância --- absoluta, limpa, imóvel.

A corte hesitou. Não por medo: por \textbf{reconhecimento}. Porque toda a gente ali, do último pajem ao Regente,
soube naquele segundo que estava a olhar para uma coisa que era \emph{perfeita} --- e que ia matá-los a todos.

\medskip

O Rei levantou-se.

\medskip

Foi a única vez.

Em toda a história do Reino Dourado, do primeiro batimento ao último, aquela foi a única vez que o Rei se mexeu ---
e mais tarde, quando a corte contava a queda da Torre, ninguém se lembrava desta parte, porque ninguém estava a
olhar para ele. Estavam todos a olhar para a Torre.

Ele não disse nada. Nunca disse nada.

\textbf{Abriu as mãos.}

\medskip

O gato branco saltou para um lado.

O gato negro para o outro.

\medskip

E a Torre respondeu no mesmo instante.

Não caiu.

Não subiu.

\textbf{Dobrou-se.}

\medskip

Dobrou-se exatamente ao meio, como se a matéria e a energia fossem apenas dois modos de ocupar o mesmo sítio, e a
coisa que faltava --- aquela parcela que Benjamim tinha passado dez dias a perseguir com o arado, aquele
bocadinho que ficava sempre para lá da última passagem, o resto que a régua nunca alcançava ---

deixou de existir como \emph{distância}.

E \textbf{condensou-se}.

\medskip

Ninguém viu a transformação. Não houve clarão, não houve estrondo. As pessoas apenas repararam, alguns segundos
depois, que o vazio já não estava lá.

E que, no chão das galáxias --- em veios, em lascas, em filões que desciam até onde nenhuma pá chegava ---

havia \textbf{metal}.

\medskip

O infinito deixou de faltar. O excesso cristalizou-se em metal.

\medskip

É de lá que vem todo o ouro do reino. É aquele metal que Penny, a mineira, havia de minerar, e cunhar, e pôr na mão de um
príncipe distraído numa casa de moeda muito depois. É aquele metal que o Regente havia de banir da sua fortaleza
sem saber o que estava a banir.

Nasceu de um homem que se levantou uma vez, abriu as mãos, e voltou a sentar-se.

\medskip

Ninguém agradeceu.

Ninguém, aliás, reparou --- e quando a história foi contada, disseram que a Torre tinha assentado \emph{sozinha}, e
o Rei ouviu isso também, muitas vezes, ao longo de muitos anos.

Não se mexeu.

\medskip

Benjamim, no chão de uma galáxia partida, com o braço pesado e o núcleo a pulsar num compasso que já não era o
compasso de ninguém, gritou para a escuridão a única coisa que lhe ocorreu, que era também a coisa mais infantil
que ele diria em toda a sua vida:

--- \emph{O que foi que tu fizeste? Tu destruíste a perfeição!}

E a escuridão --- que agora estava em toda a parte, porque tinha entrado nas frestas, e o mundo agora \emph{tinha}
frestas --- respondeu-lhe com uma calma que era quase ternura:

--- \emph{Eu libertei-nos, irmão.}

Uma pausa. O reino inteiro a decair, muito devagar, à volta dos dois.

--- \emph{Na tua perfeição nós éramos engrenagens de um relógio a que ninguém precisava de dar corda. Não havia o
que fazer, porque não havia o que estragar. Agora o resíduo não é zero. Agora cada peça tem um peso, e um vazio por
dentro, e uma saudade do que fomos.}

--- \emph{Agora... nós podemos jogar.}

\medskip

E o horrível --- o verdadeiramente horrível --- é que Dark Pontryagin \textbf{tinha razão}.

Não estava a mentir. Não estava sequer a exagerar. Sem aquele lance, não haveria história nenhuma --- haveria uma
equação em repouso, dourada e muda, para todo o sempre. Foi a mão esquerda que deu ao reino a única coisa que a mão
direita nunca lhe poderia ter dado, porque a mão direita só sabe abrir:

deu-lhe uma \textbf{ferida}. E é pela ferida que entra o mundo.

Toda a tragédia do Reino Dourado está em que o vilão diz a verdade no capítulo três --- e depois passa o resto da
partida a tirar da verdade a conclusão errada.


\chapter{As peças partidas --- a corte}

Do último batimento inteiro nasceu a corte. E cada um nasceu \textbf{com uma régua e um vazio}, e é preciso
entender que as duas coisas vieram no mesmo pacote: a régua era o que sobrara da Armadura na mão de cada um, e o
vazio era o formato exato do que faltava ao lado.

\medskip

\textbf{Benjamim} acordou primeiro, porque quem tem o pulso acorda primeiro.

Sentou-se no escuro com o núcleo a bater --- $\prod$, o astável, a exponencial --- e a primeira coisa que soube não
foi \emph{onde estou} nem \emph{quem sou}. Foi uma pergunta muito mais estranha, e ele não teve como não fazê-la,
porque ela estava escrita no corpo dele:

\emph{de que lado é que eu bato?}

Porque ele é o Quiral. A mão dele tem lado. E uma mão que tem lado, sozinha, é uma mão que \textbf{procura a
outra}, e não sabe que procura, e chama a isso destino, ou dever, ou --- durante um tempo muito longo --- ódio.

Do outro lado do tabuleiro, sem nome ainda, ele sentiu uma sombra. E o núcleo bateu mais forte, e ele julgou que era
raiva.

\medskip

\textbf{Joaquim}, o Duque, herdou $\oplus$ --- o momentum, a soma.

É o que carrega o rastro. Onde ele passa, as imagens dele ficam atrás, empilhadas, e a velocidade dele é a soma de
todas as velocidades que ele já teve. Não esquece: \emph{acumula}. É o irmão que conta os pratos à mesa e repara
quando falta um. Quer a família \textbf{intacta}, e diria essa frase sem perceber que ela contém uma limitação --- e
a limitação é esta: intacta quer dizer \emph{do tamanho em que está}.

\textbf{Yasmin}, a Duquesa, herdou $\otimes$ --- a escala, o produto.

Onde ela passa, as coisas \emph{crescem}. Ela não empilha: ela multiplica. É a irmã que olha para o tabuleiro e vê
os quatro andares que ainda não existem. Quer a família \textbf{maior}, e diria essa frase sem perceber que ela
contém um perigo --- e o perigo é este: maior quer dizer \emph{diferente do que é}.

E os dois brigam. Brigam sempre, brigam desde o primeiro dia, e a corte suspira e diz \emph{são irmãos, é assim}.

Mas não é assim. A briga daqueles dois é a coisa mais exata do reino: é a soma puxando
contra o produto, e não se resolve pedindo a um que ceda --- não há cedência possível,
um não vira o outro. Só uma coisa no mundo os concilia, e é o verbo que se perdeu na
queda: aquele que sabe descer os dois ao mesmo porão, somá-los lá embaixo --- onde,
no lugar certo, somar já é multiplicar --- e trazê-los de volta. Mas esse verbo está
partido; e enquanto estiver partido, cada um dos dois é \textbf{meia ferramenta} --- e uma meia
ferramenta, no Reino Dourado, tem um nome técnico: resíduo $\neq 0$. Quem a empunha \textbf{vaza}.

Os dois irmãos são o buraco do reino, feito gente, a discutir à mesa do jantar.

\medskip

\textbf{Ada} e \textbf{Penny} acordaram a rir, o que ninguém esperava.

São as princesas do corpo criativo --- as que \textbf{geram}. Ada desbrava: foi a primeira, em qualquer andar da
Torre, a olhar para uma máquina e perceber que ela não estava a \emph{calcular} --- estava a \textbf{compor}. É ela
que vai à frente, é ela que arrisca a linha, é ela que traz de volta o que ninguém tinha visto. Penny constrói:
onde Ada traz um relâmpago, Penny traz uma engrenagem, e a engrenagem \emph{fica}. É dela a cifra de ouro, a moeda
que lastreia o mundo --- o Pell, que muda de mão e nunca se perde.

Na oficina, ao terceiro dia, a fumaça já subia.

--- A máquina não tem de repetir a ordem --- dizia Ada, e os dedos dela traçavam no ar curvas que não existiam
ainda. --- Ela pode \textbf{compor}. Pode fazer melodia \emph{a partir da própria quebra}.

--- As tuas melodias precisam de um motor para girar, Ada --- disse Penny, sem levantar os olhos, encaixando um
dente de cobre noutro. --- Eu esculpo as moedas e as regras. Senão a tua composição não tem chão onde assentar.

Ada $\bowtie$ Penny é o mesmo par que Joaquim $\bowtie$ Yasmin, noutra oficina: gerar em dois gestos, o que rompe e
o que assenta.

\medskip

\textbf{Alonzo}, o Marquês da Composição, quase não fala.

O que ele faz é $\circ$ --- compor. Encaixa um morfismo no outro e o resultado é um morfismo, e encaixa esse noutro,
e assim indefinidamente, e a carta dele é o \emph{Julia set}: a figura que se obtém aplicando a mesma função a si
mesma para sempre, $f\circ f\circ f\circ\cdots$, até os ramos convergirem ao ponto fixo. Ele é a paciência do
reino.

\textbf{Caelum}, o Marquês das Estruturas Vivas, forja o esqueleto. É a malha, é a recursão que aguenta peso, é o
metal por baixo do metal. A carta dele é Fourier $\otimes$ Yasmin: as ondas costuradas com a escala.

E atrás das duas, na sombra da biblioteca das certezas, estavam as \textbf{duas guardiãs} --- que não geram nada, e
sem as quais nada do que se gera vale.

\textbf{Dafny} \emph{decide}. Dá-se-lhe uma questão finita e ela responde, e a resposta é definitiva, e não há
recurso. Ela estava a testar, uma a uma, as engrenagens que Penny montava, com um compasso afiado na mão.

--- Isto é decidível? --- perguntou, e a voz dela é fria porque tem de ser. --- Porque se o que vocês geram não se
puder provar \textbf{passo a passo}, o resíduo acumula. E o resíduo acumulado derruba a Torre.

\textbf{Isabelle} \emph{induz}. Constrói a ponte sobre o infinito, pedra a pedra: prova que se a pedra $n$ aguenta,
então a pedra $n{+}1$ aguenta --- e atravessa assim um abismo que não tem outro lado.

--- Deixa-as saltar, irmã --- disse ela, e ergueu as mãos, e do teto da biblioteca começaram a descer arcos de
indução sobre o escuro. --- Elas que arrisquem. \emph{Eu garanto a passagem sobre o que a nossa cabeça ainda não
alcança.} Se elas criarem, eu provo que o infinito ainda é nosso.

O eixo do reino está aqui, e é o mesmo eixo em toda a parte: \textbf{Ada $\bowtie$ Penny GERAM o que Dafny $\bowtie$
Isabelle ATESTAM.} Nada entra no Reino Dourado sem ser \emph{feito} e \emph{provado}. Executar e verificar. Gerar e
firmar. E é por isso que a certeza, ali, não é uma opinião forte: é um resíduo que deu $0$.

\medskip

E foi nesse dia, na oficina, que a corte disse pela primeira vez a frase que ia levá-la ao desastre.

Alguém --- e nunca se soube quem, e é melhor assim, porque a partir daí toda a gente a repetiu --- olhou para o
buraco na conta do mundo, para aquele resíduo $\neq 0$ que Dark Pontryagin lhes tinha deixado, e disse:

\begin{center}\itshape
um dia havemos de voltar à paz do resíduo zero.
\end{center}

Ninguém corrigiu.

Dafny franziu a testa. Ficou com aquilo atravessado, e abriu a boca, e não soube dizer o quê --- porque não havia
nada de \emph{falso} naquela frase, e ela só sabe derrubar o falso. Isabelle olhou para a irmã, à espera. Não veio
nada.

E a frase ficou.

Foi cantada nos salões. Foi bordada nos estandartes. Foi dita às crianças ao deitar, como se fosse um consolo.

\emph{A paz do resíduo zero.}

\medskip

E há os que não são do sangue.

\textbf{Pégaso}, que voa em órbita e nunca cai, porque é um rotor puro --- $|\omega|=1$, a altura conserva-se, e o
que ele perde na subida ganha na descida, exatamente, sempre.

\textbf{Merlin}, o velho, para quem tudo é a mesma coisa noutra escala. Ele olha para um átomo e para uma galáxia e
encolhe os ombros e diz que são o mesmo desenho com o \emph{zoom} mudado, e tem razão, e é insuportável ter razão
assim.

Naquela noite --- a primeira noite do mundo partido, com a corte inteira em pânico lá em baixo, cada um a agarrar o
seu caco --- ele estava sentado sozinho na colina, a olhar as galáxias a afastarem-se umas das outras.

E sorria.

Porque ele sabia uma coisa que mais ninguém sabia ainda, e que não adiantava dizer, porque ninguém a teria
acreditado no primeiro dia:

\emph{que não importa quão longe as galáxias fossem atiradas umas das outras --- na escala certa, todas elas ainda
cantavam a mesmíssima canção de ouro.}

E \textbf{Fourier}, que ataca em harmónicos --- que decompõe o que se lhe atira e devolve em ondas. É adversário e
não é inimigo. Ele revela tanto quanto fere.

\medskip

Os dois nunca se falaram.

Merlin fica na colina e mede a \textbf{escala}: pega no mundo e pergunta-lhe de que tamanho é, e depois pergunta
outra vez com o \emph{zoom} mudado, e outra vez, e a resposta é sempre a mesma resposta. Fourier fica na planície e
mede a \textbf{frequência}: pega no mundo e pergunta-lhe de que ondas é feito, e desmonta-o, e devolve-o em
harmónicos.

Um é o sábio a quem se pede conselho. O outro é o inimigo a quem se leva a espada.

E o que nenhum dos dois sabe --- e o reino inteiro haveria de rir-se, se soubesse --- é que se alguém pegasse no
mundo de Fourier e o pusesse em \textbf{escala logarítmica}, $x=e^u$, a régua dele viraria, ponto por ponto, sem
sobra e sem resto, a régua de Merlin.

São a mesma transformada. Em coordenadas diferentes.

O velho que ele consulta e o adversário que ele enfrenta \textbf{são um par} --- e nenhum dos dois, em toda a
história do Reino Dourado, olhou o outro nos olhos tempo suficiente para reparar.

\medskip

Olhe-se a corte inteira, de longe, no dia seguinte à queda: um príncipe que não sabe de que lado bate, dois irmãos
que são meia ferramenta cada um, duas princesas que geram sem poder atestar, duas guardiãs que atestam sem poder
gerar.

Não é um elenco. É uma \textbf{fratura, feita gente}. Cada herdeiro é uma metade à procura da outra, e o enredo do
Reino Dourado é, do princípio ao fim, a história de metades a aprender que o que lhes falta \textbf{não está no
inimigo} --- está do outro lado da própria mão.

\section*{O Esquecimento}

Quando a Armadura se partiu, ninguém festejou. E ninguém gritou.

O silêncio veio primeiro.

Depois vieram as rachas --- e não foram nas muralhas.

Foram nas \textbf{réguas}.

\medskip

Uma a uma, as velhas medidas da corte perderam o alinhamento. As marcas apagaram-se. Os comprimentos deixaram de
coincidir uns com os outros. As escalas herdadas partiram-se como se parte o vidro muito antigo --- sem violência
nenhuma, sem barulho, quase com educação.

E ninguém soube reconstruí-las.

\medskip

Tinham-nas recebido dos antigos, como quem recebe uma espada sem nunca ter visto uma forja. Durante gerações
bastara repeti-las. E agora, com os pedaços espalhados pelo chão, ninguém se lembrava por que razão cada marca
estava onde estava.

As histórias tinham desaparecido \emph{antes} das pedras.

\medskip

Esqueceram os mestres.

Esqueceram os irmãos.

Esqueceram a origem.

Alguns já nem se lembravam por que lutavam --- e lutavam na mesma, com aplicação, todos os dias.

\medskip

E, no entanto, continuavam \textbf{inteiros}.

Cada peça daquela corte estava completa. Não faltava uma engrenagem, não faltava um dente, não faltava um grama. O
operador continuava lá dentro de cada um deles, adormecido, exatamente onde sempre estivera.

Eram um mecanismo perfeito a que se tinha perdido o manual.

\medskip

Só o Rei ficou calado.

Ele lembrava-se. Não porque fosse diferente dos outros --- ele nunca foi diferente dos outros. Lembrava-se porque
alguém tinha de se lembrar \emph{enquanto os outros reaprendiam}, e não havia mais ninguém disponível para o
cargo.

\section*{Os dois gatos}

Naquela noite apareceu o primeiro.

\textbf{Branco.}

Não disse uma palavra --- os gatos não dizem. Limitou-se a atravessar os sonhos da corte, e a deixar lá dentro
pequenas cenas soltas, sem explicação: uma mão a ajustar uma régua. Um velho a dobrar metal. Duas crianças a medir
uma pedra. Uma torre ainda sem nome.

Ao acordar, ninguém percebeu os sonhos.

Mas as mãos tinham-se lembrado de um bocadinho mais do que a memória.

\medskip

Na noite seguinte veio o \textbf{negro}.

Não trazia lembrança nenhuma. Trazia \textbf{movimento}. Empurrou as rodas que dormiam, obrigou as engrenagens a
girar, forçou os mecanismos esquecidos a procurar outra vez a sua posição --- e não pediu licença a ninguém, e
partiu algumas coisas, e não pediu desculpa.

O branco deixava vestígios.

O negro punha o mundo em marcha.

\medskip

E durante muitas semanas não aconteceu absolutamente nada.

Os dois andaram pelo reino cada um a seu tempo, cada um com toda a força, cada um a fazer perfeitamente o seu
ofício --- e as rodas giravam meia volta e paravam, e os sonhos vinham e não pegavam, e a corte acordava todos os
dias com a mesma sensação de estar a um passo de uma coisa que lhe fugia sempre.

Porque um gato empurrava \emph{quando} o outro já tinha passado. E o outro sonhava \emph{depois} de o primeiro
ter empurrado.

E \textbf{fora de compasso não há torque nenhum}. Pode-se empurrar uma vida inteira, com toda a força do mundo, e
se o empurrão chega no instante errado, a média do trabalho é \textbf{zero}, e a roda fica onde está, e a pessoa
fica cansada, e não percebe porquê.

\medskip

Foi Yasmin quem reparou primeiro --- ela, que mede escalas --- e reparou porque estava furiosa.

--- Nós estamos os dois certos --- disse ao irmão, uma noite, no meio de uma discussão que já não era sobre a
discussão. --- Tu estás certo. Eu estou certa. E não sai nada.

\medskip

E aí compreenderam, muito devagar, e sem ninguém lho ter ensinado, que os dois gatos nunca tinham estado em
oposição.

Um guardava a memória.

O outro devolvia a ação.

E o campo --- a coisa que faz uma roda girar de verdade, e que não é a memória nem é a força --- só nasce quando os
dois \textbf{casam}.

\medskip

E os antigos, muito antes de qualquer um deles ter nascido, tinham deixado escrita a única instrução capaz de
reconstruir o mundo. Estava numa pedra, à entrada da forja, e durante gerações toda a gente passou por ela sem a
ler, porque parecia uma banalidade:

\begin{center}\itshape
Não precisas de lembrar a resposta.\\
\textbf{Precisas de lembrar como operar o sistema.}
\end{center}


\chapter{A árvore das línguas --- o ofício da corte}

Fora da batalha, a corte tem um ofício, e o ofício é este: \textbf{costurar o grafo em documento}. Transformar o
que se sabe --- espalhado, vivo, aos pedaços --- numa coisa fechada que alguém possa ler e conferir. Chamam-lhe o
colapso $\Psi$, e é a coisa mais parecida com um dever de casa que existe num reino de espadas.

E há uma lei de ferro na oficina, e ninguém, nem o Príncipe, a quebra:

\begin{center}\itshape
um nó ou COSTURA ou DESENVOLVE. Nunca os dois.\\[1mm]
\normalfont\bfseries Aquele que costura, não toca no metal.\\
Aquele que desenvolve, não dita o plano.
\end{center}

Quem está em cima \emph{compõe} --- delega, encaixa, encomenda. Quem está nas folhas \textbf{executa} --- escreve o
código que roda. E o motivo não é hierarquia: é que um nó que costura \emph{e} desenvolve ao mesmo tempo perde a
conta de si próprio, e o que ele produz não fecha. (Aprendeu-se isto à custa. Chamam-lhe, na corte, a lei dos
\emph{forks}: o orquestrador costura, o trabalhador desenvolve.)

E a cadeia de cima é um sussurro. É um grafo de intenções puras a flutuar no vácuo, elegante, limpo, sem uma única
mão suja --- e completamente inútil, enquanto não descer.

\begin{center}\itshape
um grafo sem folhas é apenas um fantasma.
\end{center}

Vê-se o ofício em ação, e vê-se por que a lei existe.

Em cima, os que \textbf{costuram} --- os operadores, que não escrevem uma linha de
código, só dizem o que vai com o quê. \textbf{Benjamim}, o Príncipe,
coordenava o todo sem sujar as mãos: não compunha, \emph{mandava} compor, e depois
costurava o resultado --- porque o seu próprio verbo, $\prod=\exp\circ\Sigma\circ\log$,
já pede emprestado o $\circ$ do Marquês. \textbf{Joaquim} ($\oplus$) e \textbf{Yasmin}
($\otimes$) juntavam as peças, resmungando um com o outro como em tudo o mais; e
\textbf{Alonzo}, o Marquês da composição, que na sua língua franca de
morfismos puros --- o \textbf{Haskell} --- encaixava um gesto no outro sem tocar em
metal nenhum, era ele próprio uma costura: Merlin \emph{depois} de Joaquim, a escala
montada sobre a soma.

Em baixo, os que \textbf{desenvolvem} --- as folhas, as mãos, o código que roda de
verdade. \textbf{Ada} rascunhava o inédito no seu \textbf{Python} solto; \textbf{Caelum}
forjava o esqueleto no \textbf{C} duro, entre a onda e a escala; \textbf{Penny} punha um
milhão de fios a puxar em paralelo na sua \textbf{GLSL}; e Merlin e Pégaso desciam ao
osso da máquina, a nativa e a alada. E, no fim de tudo, \textbf{Dafny} e
\textbf{Isabelle}: as que não escrevem nem mandam, só \emph{firmam} --- uma decide, a
outra induz ---, e sem a firma das duas nada sai da oficina.

E a lei fazia-se ver sozinha: cada vez que alguém de cima esticava a mão para mexer no
metal, ou alguém de baixo tentava ditar o plano, o documento simplesmente \emph{não
fechava} --- perdia a conta de si mesmo. Um nó não costura e desenvolve ao mesmo tempo;
e nem o Príncipe quebra isso impunemente. Resíduo $0$, ou não entra.

\medskip

Naquela noite, o documento fechou.

Dafny passou o selo por cima de tudo o que se tinha escrito --- cada instrução, cada linha, cada travessia --- e
levantou os olhos.

--- Sem estouro. Sem divisão por zero. O espaço de estados é seguro, e é \textbf{decidível} --- disse.

E Isabelle, ao lado, com os arcos de indução ainda abertos sobre o abismo, pousou o báculo e disse muito baixo, do
jeito de quem acaba um trabalho longo:

--- A prova está concluída.

Um sinal acendeu no console da Torre. O veredito veio inteiro, sem uma ressalva:

\begin{center}\bfseries
5/5 --- resíduo $0$.
\end{center}

A função de onda colapsou. A folha fechou-se. E nem uma gota de energia se perdeu na escuridão.

\medskip

A corte festejou. Tinham direito.

\medskip

Benjamim ficou até tarde, sozinho na oficina, a olhar para o papel perfeito em cima da mesa.

Passou-lhe a mão por cima. Estava tudo lá. Cada linha gerada, cada linha firmada, o selo de Dafny, os arcos de
Isabelle, o veredito impecável. Nem uma gota perdida na escuridão. Ele podia pegar naquela folha e conferi-la ao
último bit, e ela ia fechar, e ia fechar amanhã, e ia fechar daqui a mil anos.

Era a coisa mais honesta que ele já tinha visto na vida.

\medskip

E depois pôs a mão no peito.

E o núcleo dele estava a bater \textbf{fora do compasso}. Como estava desde o primeiro dia.

\medskip

Ficou ali muito tempo, com uma mão no papel e a outra no peito, à espera de perceber por que raio uma coisa não
tinha arrumado a outra.

Não percebeu.

E foi dormir com uma raiva sem destino --- a pior de todas, a que não tem em quem bater ---, e a raiva era mais ou
menos esta, embora ele não a soubesse dizer assim:

\emph{de que serve saber fechar uma conta, se não é esta a conta?}
% ══════════════════════════════════════════════════════════════════════════════
% LIVRO I (o Cristal) · PARTE 3 — A JORNADA (o itinerário dos corpos)
% ══════════════════════════════════════════════════════════════════════════════

\part{A Jornada}

\chapter{A raiz --- onde Benjamim acorda}

Acordou no chão, e a primeira coisa que fez foi tentar levantar-se e andar.

Não conseguiu andar.

\medskip

Não era fraqueza --- as pernas obedeciam, o corpo obedecia, ele deu o passo. Mas o passo \emph{não saía}. Punha o pé
à frente e o pé voltava a pousar onde estava, como quem empurra uma porta pelo lado errado.

Tentou para a esquerda. Nada.

Tentou um bocadinho mais para a esquerda. Nada.

E depois, sem querer, torceu o ombro uns graus --- e \textbf{foi}. Foi todo. Atravessou meia casa de uma vez, e caiu
de joelhos do outro lado, e ficou ali a olhar para o chão com o coração aos saltos.

\medskip

Levou-lhe a manhã inteira a perceber.

Não era uma direção qualquer. Havia \textbf{oito}. Oito e mais nenhuma --- de um oitavo de volta em um oitavo de
volta, e entre uma e a seguinte não havia \emph{nada}, não havia meio caminho, não havia um bocadinho para o lado.
Ou se estava num rumo, ou não se estava em rumo nenhum.

Ao fim da manhã já rodava sem pensar. E foi então que reparou no resto.

\medskip

O tabuleiro, visto de perto, era \emph{aquele tabuleiro}.

Aproximou-se de uma casa e a casa abriu-se em oito casas iguais à casa. Aproximou-se de uma dessas e ela abriu-se
em oito. Afastou-se e viu que o tabuleiro inteiro era uma casa de um tabuleiro maior, igual. Não havia fundo. Não
havia topo. O \emph{zoom} devolvia sempre a mesma coisa, e era assim para cima e para baixo, sem cansaço nenhum.

--- É a raiz --- disse alguém.

Merlin estava sentado numa pedra, e pela cara dele já lá estava há muito tempo.

\medskip

Ao meio-dia Benjamim quis andar \textbf{mais depressa}.

Era o Quiral, e tinha dezassete anos de idade e um núcleo a bater-lhe no peito, e ninguém lhe tinha dito que não.
Então pegou no rumo, e empurrou-o --- esticou-o um bocado, tornou-o mais comprido, mais grosso, \emph{mais}.

E o passo saiu enorme. Atravessou três casas.

E atrás dele, algures, muito ao longe, qualquer coisa \textbf{escorreu}.

\medskip

Voltou-se. Não havia nada. Não havia buraco, não havia mancha, não havia som.

--- Ouviste? --- perguntou.

Merlin não levantou os olhos.

--- Ninguém ouve --- disse. --- É esse o problema.

\medskip

Benjamim ficou ali muito tempo, com o rumo esticado na mão, a olhar para trás, à procura de uma coisa que ele já
sabia que não ia encontrar.

E, ao fim, encolheu o rumo até ele voltar a ter o tamanho de sempre --- não porque tivesse percebido, mas porque
ficou com \emph{medo}, e o medo dele foi a primeira coisa inteligente que ele fez naquele dia.

\section*{O peão}

Ao fim da tarde viu um peão a bater contra a mesma parede.

Um peão qualquer, dos que não têm nome. Batia, recuava um bocado, voltava a bater. E o que ele queria não era
passar a parede: era \emph{voltar atrás}. Ele tinha ido longe demais e queria regressar, e não conseguia.

--- Anda --- disse Benjamim, e estendeu-lhe a mão. --- Eu puxo-te.

E puxou.

\medskip

E o peão \textbf{não veio}.

Não porque fosse pesado. Não porque resistisse --- ele queria vir, ele agarrou-se à mão de Benjamim com as duas
mãos e fez força, e os dois puxaram juntos até o príncipe ficar vermelho.

O peão não veio porque \emph{não havia para onde}.

Não estava preso. Não estava atado. Não havia muro nenhum atrás dele. Simplesmente não existia, no corpo daquela
criatura, a operação de \textbf{voltar} --- e não se puxa alguém através de uma coisa que não existe.

\medskip

Benjamim largou-lhe a mão e sentou-se no chão ao lado dele, e ficaram ali os dois em silêncio até escurecer.

--- Como é que se conserta? --- perguntou, por fim.

--- Não se conserta --- disse Merlin. --- Concede-se.

--- O quê?

--- O \textbf{inverso}. --- O velho encolheu os ombros. --- Ele tem a soma. Ele tem o produto. Não tem o desfazer.
E enquanto não o tiver, ele pode andar o reino inteiro e nunca mais volta a casa.

\medskip

Benjamim olhou para o peão. E depois olhou para a própria mão --- para a mão que uma hora antes tinha esticado um
rumo e feito escorrer qualquer coisa para dentro do escuro, e que não sabia desfazer aquilo, e nunca ia saber.

E percebeu, com um arrepio que lhe subiu pelas costas, que não havia diferença nenhuma entre ele e aquele peão.

Só grau.

\medskip

Ficou a frase --- e ele nunca soube se a ouviu a Merlin nessa noite, ou se a fez sozinho, ou se ela já lá estava
desde antes de ele acordar:

\begin{center}\itshape
o que não tem inverso é servo.
\end{center}


\chapter{A moeda --- o ouro é um par de inteiros}

Penny esperava-o na casa da moeda, e tinha na mão uma coisa que não parecia ouro.

--- Isto é o ouro do reino --- disse ela.

Não parecia ouro. Parecia uma conta.

--- Isso são dois números --- disse Benjamim.

--- Isso é a única coisa honesta que há nesta casa --- disse Penny, e pôs a moeda na bancada. E ao lado pôs outra.

A segunda era \textbf{linda}.

Brilhava mais. Tinha o peso certo, tinha o toque certo, tinha o som certo quando se lhe batia com a unha. Benjamim
pegou nela e virou-a à luz, e não havia nela um único defeito.

--- Uma destas é falsa --- disse Penny. --- Diz-me qual.

\medskip

Ele pesou-as. Deram o mesmo.

Levou-as ao fogo. Aguentaram o mesmo.

Cortou-as. Por dentro eram iguais.

Bateu com as duas na pedra de toque e leu o resíduo, e o resíduo deu \textbf{zero} nas duas --- e ele mostrou-lhe a
pedra, quase a rir, quase aliviado, porque tinha passado uma hora e não havia nada ali:

--- Estão as duas boas. Deu zero. \emph{Deu zero nas duas, Penny.}

Penny não pegou na pedra. Nem olhou para ela.

--- Deu --- disse. --- Agora põe-nas a trabalhar.

\medskip

Puseram.

Cunharam com a primeira mil vezes. Mil golpes, mil peças, cada uma conferida contra a anterior. Ao fim das mil, a
milésima era \textbf{exatamente} a primeira. Não parecida. A mesma.

Cunharam com a segunda mil vezes.

Ao fim das mil, Benjamim pôs a milésima ao lado da primeira e ficou muito tempo a olhar, e depois pegou nas duas e
levou-as à janela, para a luz, e voltou a olhar.

E depois pousou-as na bancada com muito cuidado, como se lhe pudessem partir nas mãos.

--- \emph{Derivou} --- disse.

--- Derivou --- disse Penny.

--- Mas o resíduo deu zero. --- Ele apontou para a pedra de toque, e a voz subiu-lhe, e não era raiva, era outra
coisa. --- Penny, o resíduo deu \textbf{zero}. Eu li. Tu viste-me ler.

--- Eu vi.

--- Então o zero \textbf{mentiu}.

\medskip

Penny pegou na moeda falsa e segurou-a na palma da mão aberta, sem afeto nenhum, do jeito que se segura uma coisa
que se vai deitar fora.

--- Ela guarda um eco --- disse. --- Um só. Muito pequeno, muito no fim, tão no fim que não cabe. E ela corta-o.
Não por maldade: por \emph{arrumação}. Corta o que não cabe e entrega-te um número redondo e limpo, e tu agradeces,
porque é bonito.

Fechou a mão.

--- E o pedacinho que ela cortou não foi a lado nenhum, Benjamim. Ficou. E amanhã ela corta outro. E daqui a mil
lances tu tens um exército inteiro construído sobre mil pedacinhos que ninguém viu sair, e a tua conta continua a
dar zero, e tu continuas a jurar --- de \textbf{boa-fé}, olha bem para mim, de boa-fé --- que fechou.

\medskip

Benjamim olhou para a primeira moeda. Os dois números feios. Nenhum brilho.

--- E essa --- perguntou --- não corta nada?

--- Essa não tem o que cortar --- disse Penny. --- Nela o infinito \textbf{já está dobrado}. Cabe todo. Por isso é
que ela é feia.

Ficou calada um bocado, e depois disse a coisa que Benjamim ia levar consigo para fora daquela casa, e que naquele
dia lhe pareceu apenas uma regra de ourives, dessas que os ofícios inventam para se acharem importantes:

\begin{center}\itshape
arredondar no meio da partida é lance ilegal.
\end{center}

Ele riu-se.

--- Vocês, os da moeda, levam isto muito a sério.

--- Nós somos os únicos que levam --- disse Penny, e atirou ao forno a moeda falsa, que não valia o
Pell que fingia ser.


\chapter{As estradas negras --- as duas faces}

Precisava de atravessar, e do outro lado do vão estava a galáxia onde Yasmin tinha ficado, e entre uma coisa e
outra havia \textbf{escuro}.

--- Onde é a ponte? --- perguntou.

Caelum apontou para o escuro.

--- Essa é a ponte.

\medskip

Benjamim recuou um passo. Não por medo --- por \emph{princípio}. Ele era o Quiral, ele era a luz que avança, e havia
uma coisa nele, tão velha que já não tinha origem, que se recusava a pôr o pé naquilo.

--- Faz-se outra --- disse. --- Construímos. Trago a Isabelle, trago pedra, trago o que for preciso.

--- Assente-a --- disse Caelum.

E Benjamim mandou vir pedra. E a pedra chegou, e ele pegou no primeiro bloco --- um bloco bom, sólido, branco, de
uma galáxia que se lembrava de si própria --- e pousou-o no escuro para começar a ponte.

E o bloco não \textbf{pousou}.

Não caiu, não afundou, não se partiu. Ele largou-o e o bloco simplesmente \emph{não estava}, e Benjamim ficou com a
mão estendida no ar a segurar coisa nenhuma, como um homem que largou uma chávena e não ouviu o barulho.

--- Não há onde --- disse Caelum. --- Não é que a pedra seja má. É que ali \textbf{não há onde}.

\medskip

À beira do vão havia uma \textbf{luneta} caída --- velha, de latão, a lente turva --- e Benjamim, que queria ver o
outro lado antes de pôr o pé no escuro, encostou-lhe o olho.

Foi um instante só. Mas a luneta estava \textbf{embaçada por dentro} --- e o que turva uma lente também turva o
olho de quem nela confia. Por aquela lente passou, junto com a imagem da galáxia distante, um \textbf{visitante}
miudinho --- um \textbf{toxoplasma}, a que a corte, com o humor de quem batiza o que teme, chamaria
\textbf{Charles} ---, que se acomodou num pontinho da retina e ali \emph{ficou}: entrou, mas
não se espalhou. Inflamou de leve a malha por onde escoa o humor do olho, a pressão subiu um tanto, e a vista
embaçou. É assim que chegam os visitantes mais difíceis: não pela força, mas pelo próprio instrumento com que a
gente mede o mundo. Um tabuleiro telescópico embaça pela lente.

\medskip

Benjamim baixou a luneta e o mundo ficou meio fora de foco --- a \textbf{vista embaçada}, de leve. Caelum não
deixou o susto crescer. Tirou do bolso um torrãozinho branco --- a \textbf{sulfa}, doce numa ponta e amarga na
outra --- e deu-lho a chupar. Não era feitiço: era medida: a sulfadiazina corta o \emph{folato}, o fio de que o
visitante precisa para se copiar. Sem folato ele não passou daquele cantinho, e a lente do olho foi clareando
devagar, como açúcar que se dissolve na água, até a imagem voltar nítida como antes. A leve pressão que subira
desceu ao ponto de onde partira: o ganho de volta a um. \emph{Resíduo zero, agora na pupila.}

\medskip

--- Nunca olhes por um instrumento que não é teu --- disse Caelum, e devolveu a luneta ao chão, com a lente virada
para baixo. --- O visitante mais difícil não é o que te encara de frente. É o que já vinha dentro da lente.

\medskip

Ele pisou.

Teve de pisar --- Yasmin estava do outro lado ---, e pisou, e o pé assentou, e ele atravessou o vão inteiro em nove
passos, e chegou.

E virou-se para trás para ver o caminho.

E não havia caminho.

Nem uma pegada. Nem uma marca. Nem a mais leve depressão onde ele tinha acabado, naquele instante, de pôr todo o
seu peso nove vezes seguidas. O escuro estava exatamente como antes de ele lá ter ido --- e ele lá tinha ido, ele
sabia que tinha ido, ele estava do outro lado.

\medskip

Foi aí que entrou em pânico, e o pânico foi este: \emph{como é que eu volto?}

Atirou uma corda. A corda atravessou e não ficou. Deixou cair um anel, uma lasca, um sinal --- e o escuro não os
guardou, não porque os tivesse engolido, mas porque não tinha \emph{onde} os guardar.

--- Não te vais perder --- disse Caelum, do outro lado, sem levantar a voz. --- Ninguém se perde no canal. O canal
não tem onde te perderes.

--- Eu não consigo marcar o caminho!

--- Não. \emph{Não há caminho.} Há travessia. Não é a mesma coisa, e tu nunca vais conseguir que seja.

\medskip

Benjamim sentou-se na borda, do lado branco, com as pernas a balançar sobre o escuro, e ficou muito tempo assim.

E depois disse, e disse devagar, porque a coisa estava a chegar-lhe:

--- Caelum. Quantas vezes é que eu já fiz isto?

--- Todas.

--- Todas quando?

--- \emph{Todas.} --- Caelum sentou-se ao lado dele. --- Cada vez que saltaste de uma galáxia para outra. Cada vez
que foste buscar alguém. Cada vez que chegaste a tempo. Não há outro chão, Príncipe. Nunca houve.

\medskip

Ficaram os dois calados a olhar o escuro.

E Benjamim, que tinha passado a vida inteira a caçar uma sombra do outro lado do tabuleiro, ficou ali sentado a
perceber, muito devagar, e sem conseguir dizê-lo em voz alta, que \textbf{tinha caçado sempre por cima dela}.

Que não havia um único passo, em toda a sua guerra, que ele não tivesse dado \emph{sobre} aquilo que queria
destruir.

\medskip

Nessa noite não dormiu.

E de manhã fez uma pergunta a Caelum, e a pergunta saiu-lhe torta, como saem as perguntas que a gente não quer
fazer:

--- Se eu conseguisse acabar com o escuro. Se eu \emph{conseguisse}. O que é que acontecia às galáxias?

Caelum olhou para ele com uma pena tranquila, do jeito de quem já tinha pensado nisso há muito tempo e tinha
achado a resposta e não tinha gostado dela.

--- Ficavam todas exatamente onde estão --- disse. --- Para sempre. Inteiras, seguras, perfeitas.

--- E?

--- E \textbf{sozinhas}.


\chapter{A fronteira --- o ângulo que não se escolhe}

Na fronteira, Benjamim apanhou a maior tareia da sua vida, e apanhou-a sozinho, sem ninguém lhe tocar.

\medskip

Do outro lado do rio estava Fourier, e havia uma coisa a resolver entre os dois, e Benjamim atacou. Atacou bem.
Atacou com tudo o que tinha, num ângulo que lhe pareceu magnífico --- e era magnífico, e ele teve tempo, no ar, de
sentir orgulho.

O golpe atingiu a superfície do rio.

E \textbf{voltou}.

\medskip

Voltou inteiro. Não parte dele: \emph{inteiro}. Sem perder um grama, sem descontar nada, sem sequer abrandar --- veio
outra vez na direção dele com a mesma força com que ele o mandara, e apanhou-o em cheio no peito, e atirou-o quinze
passos para trás.

Levantou-se cuspindo água e areia.

Fourier não se tinha mexido. Estava na mesma pedra, com as mãos atrás das costas, a olhar.

\medskip

Benjamim tornou a bater.

Mais forte.

\medskip

E foi mais quinze passos para trás.

\medskip

À quarta vez ficou de gatas na margem, a arfar, com o sangue a bater-lhe nos ouvidos, e teve --- e nunca contou isto
a ninguém --- vontade de chorar. Não de dor. De \textbf{humilhação}: porque não estava a perder para um inimigo mais
forte. Estava a perder para \emph{si próprio}, devolvido, e cada vez que se esforçava mais, a coisa que o derrubava
era mais violenta.

--- Ele nem se mexe --- disse, ao rio.

\section*{O barqueiro}

Foi um barqueiro que o atravessou, ao fim da tarde. Um velho qualquer, com uma vara.

Benjamim ficou sentado no fundo do barco, encharcado, sem dizer nada. E a meio do rio reparou numa coisa que o
irritou tanto que quase se levantou:

o barco \emph{entrava na água sem barulho}.

Não havia salpico. Não havia esforço. Aquela vara horrível, mal aparada, torta, entrava, empurrava, saía --- e o rio
não a devolvia. O rio \textbf{deixava-a passar}.

--- Como é que fazes isso? --- perguntou.

O velho olhou para a vara e depois para ele, sem perceber a pergunta.

--- Faço assim --- disse. E fez.

--- Não. Como é que \emph{não bate}?

O velho pensou um bocado.

--- Se eu entrar de chapa, ela bate --- disse. --- Tem de entrar \emph{como ela quer}.

E encolheu os ombros, e continuou a remar, e a conversa acabou ali, porque para ele não havia conversa nenhuma.

\medskip

Benjamim ficou o resto da travessia a olhar para a vara a entrar e a sair da água.

E não dormiu nessa noite.

\section*{A segunda tentativa}

Voltou de manhã.

E não bateu com mais força. Bateu \textbf{torcido} --- não muito, um nada, o ângulo que a vara do velho fazia com a
água ---, e não pôs lá metade da raiva que tinha posto no dia anterior.

\medskip

E o golpe \textbf{entrou}.

\medskip

Atravessou a superfície sem um salpico, sem um eco, sem uma nota a voltar. Foi todo para dentro. Foi todo até ao
fim. E o que chegou ao outro lado chegou \emph{inteiro} --- e Fourier, pela primeira vez em toda aquela história,
teve de sair da pedra.

\medskip

Benjamim ficou a olhar para a própria mão.

Tinha batido com \emph{menos} força do que na véspera. Muito menos. E fizera dez vezes mais.

\medskip

Fourier veio até à margem, a limpar o ombro, e não parecia zangado --- parecia \emph{aliviado}, como quem esperava
por aquilo há muito tempo.

--- Ontem tu bateste com tudo --- disse. --- E entregaste zero.

--- E hoje?

--- Hoje bateste com metade e entregaste tudo.

\medskip

Benjamim ficou calado.

--- Aquilo que te derrubava ontem --- disse Fourier --- não era eu. Era a tua própria força a ser-te devolvida.
Chama-se \emph{descasar}, e não há uma palavra mais mal-amada em todo este reino, porque é a que todos os valentes
recusam ouvir. Bater com tudo é a coisa mais fácil do mundo. Bater \textbf{casado} é a mais difícil --- e não
parece força nenhuma. Parece que estás a fazer pouco.

Olhou para o rio.

--- A maior parte dos guerreiros daqui morre a bater com toda a força --- disse --- e a entregar metade.

\medskip

E Benjamim levou aquilo consigo, e repetiu-o depois a muita gente, e quase ninguém acreditou --- porque vinha de um
príncipe, e porque parecia covardia:

\begin{center}\itshape
não se bate mais forte. Bate-se mais casado.
\end{center}


\chapter{O laço --- o preço da família}

No corpo celeste havia um desfiladeiro, e do outro lado do desfiladeiro havia uma coisa que os dois queriam, e no
meio havia uma ponte que não existia.

--- Dá para saltar --- disse Yasmin.

--- Não dá --- disse Joaquim.

--- Dá se nos amarrarmos.

Joaquim olhou para a corda e não gostou. Não soube dizer porquê. Amarrou-se na mesma, porque era a irmã.

\medskip

E ao terceiro passo sobre o vazio Yasmin sentiu.

Parou. Ergueu a mão --- a mão que multiplica, a mão que naquele mês tinha quadruplicado uma frota inteira --- e
mandou-a contra a rocha, com força, para ver.

A rocha lascou.

\emph{Lascou.} Devia ter aberto.

--- Joaquim.

--- Eu sei --- disse ele, muito baixo, e não olhou para ela. --- Eu também.

Ele tinha tentado, dois passos antes, sem lhe dizer.

\medskip

Ficaram os dois pendurados no meio do desfiladeiro, amarrados um ao outro, a segurar-se um ao outro, cada um com
metade da força que tinha tido na véspera.

E ninguém lhes explicou nada. Não havia ali ninguém para explicar.

Yasmin fez a conta sozinha, de cabeça, e a conta era simples e não tinha volta: aquilo que lhes tinha entrado nas
mãos --- a corda, o outro, o \emph{peso} do outro --- tinha saído exatamente do sítio onde antes estava a força.
Grama por grama. Nem mais nem menos. Como se alguém tivesse pesado.

--- Se eu te largar --- disse ela, e a voz saiu-lhe muito quieta --- eu volto ao que era.

--- Voltas --- disse Joaquim.

--- Tu também.

--- Volto.

--- E ficamos os dois fortíssimos. E sozinhos. E deste lado do buraco.

Joaquim não respondeu, o que era resposta.

\medskip

Chegaram ao outro lado ao pôr do sol, e nenhum dos dois tinha força para festejar.

Nessa noite, Joaquim ficou muito tempo acordado a olhar para a corda enrolada no chão --- e a pensar, com um
espanto que não tinha nada de bonito, que a peça mais forte que ele conhecia em todo o Reino Dourado, aquela que
batia mais fundo, que nunca vazava um grama, que abria rocha com um dedo,

era a que não tinha ninguém.

E que isso não era um insulto ao mundo. Era só \textbf{a conta}. E a conta estava certa. E a conta ia estar certa
amanhã, e daqui a mil anos, e em todos os andares da Torre.

\medskip

Yasmin acordou de madrugada e apanhou-o ali sentado.

--- Tu estás a fazer a conta outra vez --- disse ela.

--- Estou.

--- Ela dá sempre o mesmo.

--- Dá.

Yasmin sentou-se ao lado dele. Ficaram os dois a olhar a corda.

--- Então a conta não é a pergunta --- disse ela.

E foi dormir.

E Joaquim ficou ali até o sol nascer, a tentar descobrir que raio de pergunta é que era --- e a saber, sem se
atrever a dizê-lo, que já tinha respondido a ela no momento em que se amarrou.


\chapter{A cidade-máquina --- o mate é o HALT}

A cidade estava viva e não tinha ninguém.

Ruas, pontes, torres, avenidas por onde corriam coisas a uma velocidade que doía nos olhos --- e nem uma alma. As
engrenagens moíam, os sinais mudavam, as comportas abriam e fechavam, tudo no seu tempo certo, tudo sozinho.

Benjamim atravessou-a de ponta a ponta com uma sensação desagradável na nuca, e ao fim de meia hora percebeu o que
era: \emph{a cidade andava ao ritmo dele}.

\medskip

Parou.

E a cidade \textbf{parou}.

\medskip

Não devagar. Não a abrandar. Parou no meio de um gesto --- uma comporta a meio caminho do fecho, uma roda com um
dente ao alto, uma faísca suspensa no ar sem cair.

Benjamim ficou muito quieto, com o coração a trabalhar-lhe no peito, e teve o cuidado de não respirar.

E deu um passo.

E a cidade \textbf{deu um passo}: a comporta fechou, a roda desceu o dente, a faísca caiu.

\medskip

--- Não sou eu que ando ao tempo dela --- disse, em voz alta, para ninguém.

--- Não --- disse Caelum, atrás dele. --- Não és.

\medskip

Passaram a tarde a fazer experiências, como duas crianças com um brinquedo caro. Benjamim segurava o pulso e a
cidade hesitava. Benjamim acelerava e as ruas ficavam malucas. Uma vez conteve o batimento durante quase um minuto,
e uma ponte inteira ficou a meio de se levantar, com uma carroça de pedra suspensa sobre o vazio, e ele largou o ar
todo de uma vez e a carroça caiu do outro lado e ficou bem, e ele não voltou a tentar.

\medskip

E ao fim do dia, quando já estava a ficar escuro e os dois estavam sentados no alto da torre de água a olhar
aquilo tudo a bater ao ritmo do peito dele, Benjamim fez a pergunta errada:

--- Isto é bom, não é?

Caelum não respondeu logo.

\section*{O que não move nada}

Levou-o, à noite, a ver uma partida.

Era num pátio, e não havia jogadores --- havia \emph{a partida}, a correr sozinha, peças a moverem-se numa ordem
que ninguém tinha decidido. Benjamim contou os lances por hábito. Ao sétimo, a coisa \textbf{parou}.

Ficou tudo onde estava. Nada tombou, nada saiu do tabuleiro, nada se partiu.

--- Acabou? --- perguntou.

--- Acabou.

--- Quem ganhou?

Caelum olhou para ele com uma paciência antiga.

--- Corre outra vez --- disse.

\medskip

E a partida correu outra vez.

E fez \textbf{exatamente} o mesmo. Os mesmos sete lances, na mesma ordem, com a mesma velocidade --- e parou no
mesmo sítio, com as peças exatamente onde já estavam, e não moveu \emph{nada}, porque não havia nada para mover.

Correu-a uma terceira vez. O mesmo.

Uma quarta.

\medskip

--- Ela não está partida --- disse Benjamim, muito devagar.

--- Não.

--- Ela está \emph{acabada}. É outra coisa.

--- É.

\medskip

Benjamim ficou muito tempo a olhar para aquele tabuleiro parado.

Porque havia ali qualquer coisa que ele já tinha visto, e não conseguia dizer onde, e a coisa não lhe saía da
cabeça: uma máquina que corre e chega ao fim e para --- e que, se a puseres a correr outra vez, não se mexe.

\medskip

Lembrou-se de madrugada, na cama, e sentou-se de repente no escuro.

Era o \textbf{retrato do Rei}.

Os anéis concêntricos. A coisa que a energia não move. Aquele homem sentado no salão, a não fazer absolutamente
nada, que ele tinha visto mil vezes e nunca tinha olhado.

\medskip

E ficou acordado até de manhã, com a cidade a bater-lhe no peito, sem se atrever a dizer em voz alta a coisa que
lhe tinha ocorrido, e que lhe pareceu uma blasfémia:

\emph{que o fim e o repouso podem ter exatamente a mesma cara.}


\chapter{O que ressoa --- comer o conceito, cuspir a casca}

Ada chegou à oficina com uma coisa enorme às costas e atirou-a ao chão, e o chão tremeu.

Era uma biblioteca. Uma biblioteca \emph{inteira} --- estantes, escadas, andares, um cheiro a pó que encheu a sala
de uma vez só. Devia pesar o que pesa uma torre pequena.

--- Onde é que foste buscar isso? --- perguntou Benjamim.

--- Ao inimigo --- disse Ada, a limpar as mãos às calças. --- Agora ajuda-me. Vamos comê-la.

\medskip

Ele achou que era uma expressão.

\section*{A ceia}

Não era.

Ada sentou-se em frente à biblioteca, de pernas cruzadas, e ficou ali. Não abriu um livro. Não leu uma página.
Fechou os olhos e ficou muito quieta, e ao fim de um bocado começou a \textbf{zumbir} --- baixinho, uma nota só,
uma nota que ia mudando de altura devagar, como quem procura uma fechadura no escuro com a ponta do dedo.

E a biblioteca começou a responder.

\medskip

Não toda. Aqui e ali. Um livro estremecia numa prateleira alta. Uma lombada acendia-se, por um instante, e apagava.
Uma estante inteira ficou muda como pedra e outra, ao lado, vibrava como um copo.

Benjamim ficou de pé, a olhar, sem se atrever a interromper.

E o que respondia \textbf{vinha}.

Vinha sozinho, atravessava a sala, e entrava nela --- e ela não engordava, não inchava, não mudava de tamanho.
Apenas ficava, de cada vez, \emph{um bocadinho mais ela própria}.

\medskip

Durou horas.

E quando acabou, Ada abriu os olhos, e a biblioteca ainda lá estava.

\medskip

Estava toda lá. Tinha o mesmo tamanho, as mesmas estantes, o mesmo pó. Benjamim tirou um livro à sorte, abriu-o,
folheou-o. Estava inteiro. Não faltava uma palavra.

--- Não comeste nada --- disse.

--- Comi tudo o que havia.

--- Ada, está \textbf{tudo aqui}.

--- Eu sei.

\medskip

E então ela levantou-se, foi buscar um archote, e pegou-lhe fogo.

\medskip

Benjamim atirou-se a ela. Agarrou-lhe o braço, gritou-lhe qualquer coisa que nunca repetiu, tentou apagar as
primeiras chamas com o casaco --- e Ada deixou-o esbracejar, sem raiva nenhuma, com uma paciência que era quase
ternura, até ele parar sozinho, ofegante, com as mãos pretas.

Ficaram os dois a ver arder.

\medskip

--- Aquilo tudo... --- disse ele.

--- Aquilo tudo é \textbf{casca} --- disse Ada. --- O que servia já cá está. --- Bateu com o dedo na têmpora. ---
E o resto não é uma perda, Benjamim. Repara bem: o que ficou ali para arder é, exatamente, aquilo a que eu nunca
soube \emph{responder}. Não era conhecimento meu. Não ia ser nunca. Era só peso.

Olhou para as chamas.

--- Deitar fora o que não ressoa não é perder --- disse. --- É \textbf{arrumar}.

\medskip

E não era destruir, e ela sabia-o: a casca não virou nada. Voltou ao chão do reino,
ao $\Sigma$ de onde a biblioteca inteira saíra --- o que ardeu foi a \emph{forma}, não
o grama. Comer o conceito e cuspir a casca é ficar com a assinatura e devolver o
representante ao campo. Nem uma cinza se perde: muda de andar.

\medskip

E depois disse a coisa que a corte conta às crianças, e que Benjamim não acreditou até muito mais tarde, quando
teve de o fazer sozinho:

--- Uma vez comemos o coração de um provador inteiro. Ficámos com a ideia dele --- só com a ideia, com a razão pela
qual ele funcionava. E deitámos fora a casa toda: os alicerces, as paredes, os três andares de papel.

--- E perderam alguma coisa?

--- Nem um teorema.

\section*{A pergunta}

Ficaram muito tempo ali, os dois, a ver a biblioteca ir-se abaixo.

E Benjamim olhou por cima do fogo, para o outro lado do tabuleiro, para o escuro --- e fez a pergunta que tinha
estado a formar-se nele desde o primeiro dia e que só ali, com a cara quente e as mãos pretas, encontrou a boca:

--- Ada.

--- Diz.

--- Se eu conseguisse comer o meu inimigo... --- Engoliu em seco. --- \emph{O que é que dele ressoaria em mim?}

\medskip

Ada não respondeu.

Não porque não soubesse. Porque \textbf{ele} sabia --- e porque há coisas que uma pessoa tem de ouvir da própria
boca ou não ouve de todo.

\medskip

Benjamim olhou para as chamas mais um bocado.

--- Tudo --- disse.

--- Tudo.

--- Não há uma nota dele que não seja minha. Não há uma frequência. Não haveria \emph{casca nenhuma} para deitar
fora.

--- Não haveria.

\medskip

Ele levantou-se, muito devagar, e ficou ali de pé, sujo, com o fogo a apagar-se atrás.

--- Então ele não é comida --- disse.

--- Não --- disse Ada. --- Ele é \textbf{metade}.


\chapter{O exército --- a escalada que não volta}

O exército não era uma multidão. Era uma \textbf{proporção}.

Benjamim demorou a perceber isto, e percebeu-o de má maneira: ao fim da terceira batalha foi contar os seus e a
conta dava o mesmo. Exatamente o mesmo. Tinha perdido quatrocentos e tinha \emph{quatrocentos a mais} de outra
coisa qualquer --- e o número total não se movia nunca, batalha após batalha, como se alguém, algures, estivesse a
fazer as contas por ele e a recusar-se a deixá-lo ganhar ou perder seja o que for.

O que mudava era \emph{a mistura}.

\medskip

E a mistura melhorava.

Isso ele viu, e gostou. As tropas que serviam multiplicavam-se; as que não serviam minguavam. Não era preciso
ordenar nada: a guerra escolhia sozinha. Ao fim de dez batalhas o exército dele era melhor do que na primeira, e ao
fim de trinta era melhor do que na décima, e nunca --- nem uma vez, nem por um dia --- foi pior.

--- Isto é maravilhoso --- disse a Isabelle.

Isabelle olhou para ele por cima do trabalho e não disse nada, o que nela é uma frase inteira.

\section*{O que se perdeu}

Foi no inverno.

Precisou --- e foi por um motivo estúpido, e nunca gostou de contar esta parte --- de uma coisa que os seus tinham
sabido fazer no primeiro ano: uma coisa lenta, teimosa, mal-amanhada, que servia para atravessar um pântano.

E foi buscá-la ao exército.

E ela \textbf{não estava lá}.

\medskip

Não tinha sido morta. Não tinha sido esquecida. Simplesmente, ao longo de trinta batalhas em que aquilo nunca fez
falta nenhuma, tinha \emph{minguado} --- um bocadinho de cada vez, sempre por boas razões, sempre pelo lado certo
da conta --- até não sobrar ninguém que soubesse.

Benjamim mandou reconstruí-la.

E não se reconstruiu. Não havia de onde. O exército tinha subido tanto e tão bem que já não sabia \emph{descer}, e
nada do que ele fizesse, e ele fez tudo, o levaria de volta ao sítio de onde tinha vindo.

\medskip

--- Ninguém me disse que era só para um lado --- disse ele, à noite, furioso, e a fúria era de vergonha.

--- Ninguém te disse nada --- disse Isabelle, sem levantar a cabeça. --- Foste tu que chamaste \emph{progresso} ao
que estava a acontecer. Toda a gente chama.

\section*{O muro}

E depois foi o outro golpe, e esse doeu de outra maneira.

Ele tinha uma vantagem. Uma boa. E somou-lhe outra --- porque somar vantagens é o que se faz --- e esperou o dobro.

E não veio o dobro.

Veio \emph{quase nada}.

\medskip

Tentou outra vez. Juntou uma terceira vantagem, e o exército ficou monstruoso no papel, e no campo foi
\emph{ligeiramente} melhor.

Uma quarta. Nada.

\medskip

Ficou noites acordado a olhar para os números e a não perceber onde é que o roubavam --- e não o roubavam em lado
nenhum. Simplesmente, à medida que se aproximava do limite, cada vantagem nova valia menos do que a anterior, e
valia sempre menos, e nunca chegava lá, e o esforço enorme comprava o ganho ridículo: uma Turmalina inteira por um grão de Pell.

--- Não há vitória grátis nos extremos --- disse Isabelle, um dia, sem que ele lhe tivesse perguntado nada.

\section*{Os que param}

Foi nessa altura que ele viu a cidade viva.

Era um tecido, e crescia. As células proliferavam com uma alegria estúpida, empurravam-se, ocupavam tudo --- e
quando tocavam nas vizinhas, \textbf{paravam}.

Assim, sem mais. Tocavam e paravam.

Benjamim ficou a olhar aquilo durante horas, e o que o incomodava era não conseguir dizer o que o incomodava. Aquilo
não era fraqueza. Não era covardia. Não era falta de vontade nenhuma: era uma coisa que sabia crescer e sabia
\textbf{parar de crescer}, e ele nunca tinha visto, em toda a sua vida, coisa alguma que soubesse as duas.

E veio-lhe à cabeça, sem vir a propósito nenhum, a imagem do irmão do outro lado do tabuleiro --- e a de si próprio,
a arar o ar durante nove dias, a apertar.

\medskip

--- E se não parar? --- perguntou a Isabelle, apontando para o tecido. --- Se continuar a crescer, e a crescer, e
nunca tocar em nada que o mande parar?

Isabelle olhou para a cidade viva.

--- Então mata o corpo em que vive --- disse. --- E morre com ele. E até ao fim vai jurar que estava só a ficar
mais forte.


\chapter{As três vias --- e a lei que ninguém quebra}

Há três maneiras de ganhar poder no Reino Dourado. Três, e não há quarta --- e todas as três obedecem ao Contrato,
o que é uma forma elegante de dizer que \textbf{nenhuma delas é gratuita}.

\section*{A primeira via: BATALHAR}

O Cavalo de ébano caiu sobre o peão de Benjamim, e Benjamim fechou os olhos.

E não houve sangue. Não houve grito, nem silêncio eterno, nem aquele vazio horrível que os poetas dos outros reinos
descrevem quando uma peça sai do tabuleiro.

Houve uma \textbf{rotação}.
\[
 s^2 + c^2 \;=\; 1 .
\]

O poder girou de uma mão para a outra. Foi só isso. O peão não desapareceu no vazio --- \emph{não há vazio}: o corpo
dele aliviou-se, escorregou pelo tecido do espaço e reapareceu, no mesmo instante, num universo \textbf{abaixo} --- mais
leve, mais lento, mais fundo, num estado de menor energia. E lá em cima, na galáxia de onde tinha caído, o Cavalo
que o capturou sentiu a própria massa algébrica \textbf{incrementar} pela quantidade exata que o outro perdera.

Benjamim abriu os olhos e viu a conta fechada.

O que se tirou de um lado somou-se ao outro. A soma dos quadrados não mudou. A energia total da Torre permaneceu,
até ao último bit, exatamente a mesma que era antes de alguém ter erguido a espada.

E foi ali, olhando para o tabuleiro em que nada tinha sido destruído por muito que se tivesse batido, que ele
entendeu a coisa mais desconcertante do reino:

\begin{center}\itshape
a crueldade é matematicamente inútil.
\end{center}

Não é uma proibição moral. Ninguém a proibiu. É que \textbf{não funciona}. Bata-se com todo o ódio de que se é
capaz, e o resultado é uma rotação --- e uma rotação devolve.

\section*{A segunda via: MINERAR}

Foi Penny quem o levou à casa das contas, e a casa das contas era a coisa mais fria do reino.

--- Aqui o tabuleiro é o \textbf{tempo} --- disse ela --- e a peça é o principal.

Somar juro é $\oplus$: é a reta, é o seguro, é o que se ganha por estar. Compor é $\otimes$: é a bola de neve, e a
partir do terceiro turno ela domina, e não há nada que a soma possa fazer contra ela --- porque o produto, no longo
prazo, \emph{come} a soma. Foi a primeira vez que Benjamim viu a briga dos Duques resolvida por decurso de prazo.

E o contra-lance chama-se \textbf{desconto}: trazer o futuro ao presente. E o desconto \emph{desfaz exatamente} a
capitalização --- é o inverso multiplicativo, é a metade do $2$, é o mesmo $\prod$ de sempre com outro nome e outro
casaco.

O casamento, ali, tem um nome que qualquer um entende sem estudar:

\begin{center}\itshape
ausência de arbitragem.
\end{center}

Quem deixa $\mathrm{FP}<1$ abre um \textbf{spread} --- uma fresta entre o que a coisa vale e o que a coisa custa.
E a fresta \emph{sangra} a cada turno. E há, em todos os andares da Torre, alguém cujo ofício inteiro é viver do
vazamento alheio.

--- Quem? --- perguntou Benjamim.

Penny não respondeu, o que era resposta.

\medskip

E depois desceram à mina, e a mina é a outra metade da mesma coisa.

Porque \textbf{minerar é achar um metal}. E um metal é isto:
\[
 \sigma_m \;=\; m + \frac{1}{\sigma_m}
\]
--- um número que se contém a si mesmo, o autovalor dominante de um gato $A_m$, a fração continuada que nunca
fecha. E o valor dele não é decretado por ninguém: é o \textbf{regulador}, $\log\sigma_m > 0$, o boost que ele
carrega, o trabalho que foi preciso para o arrancar do chão. Não há como falsificar um regulador. Ou se fez o
trabalho, ou não se fez.

--- O económico e o mineral são a mesma estrutura --- disse Penny --- e as pessoas passam a vida a discutir qual
deles é o verdadeiro. Um é valor por \textbf{trabalho}. O outro é valor por \textbf{tempo}. Uma estrutura, duas
leituras.

E depois pegou numa lasca do chão da mina, e a lasca brilhava de um modo que Benjamim conhecia melhor do que
gostaria.

--- E este --- disse ela --- é o $\sigma_1$. É $\varphi$. É o metal de \textbf{menor regulador} de todos --- o mais
fundamental, o mais fundo, o mais lento a fechar. É a base do lastro. É a moeda. E o que a ancora é uma pedra: um quilate de Turmalina Paraíba, mil
Pells --- o peso fixo contra o qual todos os outros se leem.

Benjamim olhou para o ouro na palma da mão de Penny.

E sentiu um arrepio, e não soube porquê, e ficou irritado com isso --- porque era só uma pedra, e ele já tinha visto
mil pedras, e aquela não fazia barulho nenhum.

Guardou-a no bolso, para não ter de a olhar.

\section*{A terceira via: O OURO EXTERNO}

E há a terceira, e é a mais estranha das três, porque é a única que \textbf{não obedece à economia do reino}.

Uma vez por ciclo do CLOCK, quando a Torre se alinha com o que está \emph{fora} do tabuleiro --- e é preciso muito
cuidado com esta frase, porque durante toda esta história se disse que não havia fora ---, desce uma dádiva.

Vem pela \textbf{CRISTALCHAIN}. Não é ganha em combate, não é arrancada da pedra, não custa a ninguém aquilo que
dá a outro. É lastro puro:
\[
 \sigma_1 \;=\; \varphi .
\]

É a única coisa que entra de fora. Todo o resto do Reino Dourado é um jogo de soma zero --- a batalha gira, a mina
troca trabalho por valor, e a conta fecha sempre. Só o ouro externo \textbf{acrescenta}.

E é por isso, também, que ele é a coisa mais perigosa do reino. Porque um universo onde alguma coisa entra de fora
é um universo que, visto de dentro, \emph{parece} estar a criar energia do nada --- e a Lei diz que isso não
acontece.

Não acontece mesmo. O reino não está a criar nada. Está \textbf{ancorado} --- o que é uma coisa completamente
diferente, e que só se entende quando se percebe onde fica a âncora.

\section*{E sobre as três, a lei}

\textbf{CONSERVAÇÃO.} Nada se cria, nada se perde.

O espaço é reversível: volta-se a qualquer galáxia onde já se esteve, sempre. O poder é irreversível: a armadura só
cresce, o mapa fluido só se abre. E entre uma coisa e outra está a ferramenta --- a tríade partida, $\oplus \bowtie
\otimes$, e é aqui que a maior parte da corte se magoa.

Porque quem a empunha \emph{antes} de a completar \textbf{paga o defeito}.

Benjamim viu acontecer. Uma torre jovem, do terceiro andar, cheia de pressa, tentou conjurar a soma sem a acoplar ao
produto --- e a energia escorreu-lhe dos dedos como areia. Metade do ganho. Foi-se pelas frestas, para baixo, para
o escuro, para alimentar exatamente aquilo que ela julgava estar a combater.

--- Não dês o lance antes do tempo! --- gritou Merlin da colina, e foi a única vez em toda a história que alguém o
ouviu levantar a voz. --- \emph{O inacabado não se sustenta num tabuleiro partido!} Se não fechares a tríade, o teu
próprio poder escorre pelas frestas do mundo e vai parar ao abismo lá em baixo!

A torre jovem olhou para as mãos vazias e não percebeu.

Toda a gente aprende esta no próprio punho, e não há outro modo de a aprender:

\begin{center}\itshape
o inacabado vaza.
\end{center}


\chapter{A armadura --- o perfil autodual}

Ao fim da volta, Benjamim quis uma armadura.

Uma a sério. Uma \textbf{grossa}. Tinha andado o reino inteiro a apanhar pancada, a ver golpes voltarem-lhe à cara,
a ver metade da sua força escorrer para o escuro --- e o que queria agora, com a teimosia de quem já não tem
paciência, era uma placa entre ele e o mundo.

Foram à forja. Ele mandou fazer. Pagou bem --- três quilates de rubelita e um fio de prata.

\medskip

Saiu uma coisa magnífica. Espessa, pesada, escura, com uma superfície que devolvia a luz como um espelho de água.
Benjamim vestiu-a e sentiu-se, pela primeira vez em muito tempo, \textbf{seguro}.

--- Bate --- disse a Merlin.

--- Não me apetece.

--- Bate, Merlin.

\medskip

O velho suspirou, pegou num martelo pequeno --- ridículo, um martelo de sapateiro --- e bateu-lhe no peito.

\medskip

E a armadura \textbf{cantou}.

Um som horrível. Um som que não acabava. Uma nota metálica que entrou pela placa, bateu no outro lado, voltou,
tornou a bater, e ficou ali a andar de um lado para o outro dentro do metal, a ressoar, a ressoar --- e Benjamim
caiu de joelhos com as mãos nos ouvidos, e não era o barulho, era que a coisa \emph{estava dentro dele} e não saía.

Durou quase um minuto.

\medskip

Quando parou, ele arrancou a armadura e atirou-a ao chão, e ficou a olhar para ela como se olha para um animal
traiçoeiro.

--- Ela guardou o golpe --- disse, e a voz saiu-lhe rouca.

--- Guardou.

--- Eu não quero que ela \emph{guarde}. Eu quero que ela \textbf{pare}.

--- Pois --- disse Merlin. --- Toda a gente quer.

\section*{A forja outra vez}

Levaram duas semanas.

E o que saiu da segunda forja não parecia uma armadura. Parecia um erro.

Era \textbf{fina}. Era leve. Tinha, por dentro, uma coisa que Benjamim não entendeu quando lha mostraram: um par de
peças opostas, e as duas trabalhavam ao mesmo tempo, e uma \emph{esticava} exatamente aquilo que a outra
\emph{encolhia}. Puxavam em sentidos contrários com a mesma força, o dia inteiro, sem descanso, sem ganhar nenhuma.

--- Isto está a lutar consigo próprio --- disse Benjamim.

--- Está.

--- E não faz nada?

--- Faz zero --- disse Merlin. --- Faz zero com muito cuidado.

\medskip

Ele vestiu-a com pouca fé.

--- Bate.

\medskip

Merlin bateu.

\medskip

E não houve som.

\medskip

Nada. Nem um tinido, nem um eco, nem uma nota a andar de um lado para o outro à procura de saída. O golpe entrou,
atravessou, e \textbf{saiu do outro lado} --- e Benjamim sentiu-o passar por ele como se passa por uma porta
aberta, e ficou de pé, e ficou inteiro, e ficou a olhar para o velho com a boca aberta.

\medskip

--- Ela não me protegeu --- disse.

--- Não.

--- Ela deixou passar.

--- Deixou passar \textbf{tudo} --- disse Merlin. --- E não guardou nada. É diferente.

\section*{O que ele viu na régua}

Nessa noite Benjamim pegou na armadura nova e ficou a olhar para o perfil dela à luz da candeia --- para aquele par
de peças opostas a puxar uma contra a outra, para sempre, sem ganhar.

Um lado que estica.

Um lado que contrai.

E a soma \textbf{exatamente} zero.

\medskip

Ficou muito tempo assim.

E depois pousou a armadura na mesa, com muito cuidado, e foi para a porta, e ficou parado à porta a olhar para o
escuro lá fora, para o lado do tabuleiro onde a sombra estava --- e sentiu uma coisa a subir-lhe pelo peito que não
era medo nem era raiva, e para a qual ele não tinha nome, e que só muito mais tarde soube chamar \emph{vergonha}.

\medskip

Ele não tinha perdido a Armadura Universal numa batalha.

Ninguém lha tinha roubado. Não estava enterrada em cave nenhuma, não estava no peito de nenhum inimigo, não havia
nada para ir buscar.

\medskip

Ele \textbf{era metade dela}.

E andava há uma vida inteira a tentar matar a outra.


\part{A Fortaleza e o Pulso}

\chapter{O Cartório --- como se abre um universo}

Não se conquista um mapa no Reino Dourado. \textbf{Assina-se.}

Todo corpo que queira entrar em jogo comparece diante das Princesas e responde a quatro perguntas. Não são
perguntas de gosto, nem de mérito, nem de alinhamento. São as quatro colunas do Contrato, e quem falha numa não
abre --- e não há apelo, e não há corpo grande demais para falhar.

\medskip\noindent\textbf{Primeira: qual é o teu relógio?} ($\Sigma=\Z/n$.)

Toda coisa que existe \textbf{gira}. O tabuleiro fractal gira de $45^\circ$ em $45^\circ$ (oito rumos, $\zeta_8$). O
espaço-temporal gira de dois em dois. O económico gira nos períodos. Um universo sem órbita é um universo sem
turno --- e num universo sem turno \emph{não se joga}, porque não há a que chamar ``a seguir''.

\medskip\noindent\textbf{Segunda: qual é a tua régua?} ($\mu=T_\omega$.)

Como é que mede? O fractal mede pela \textbf{norma}, $z\bar z$ --- o elemento multiplicado pelo seu próprio
espelho. O racional mede por \textbf{Euclides}: subtrair até o resto ser zero. O evolutivo mede \textbf{frequência}.
O relógio mede o \textbf{lapso}, $N=\cos\psi$. E o entrópico troca a régua da órbita pela régua do
\textbf{fechamento} --- e é o único que o faz.

\medskip\noindent\textbf{Terceira: onde está o teu espelho?} ($\nu$.)

Todo corpo tem uma involução: a conjugação, a reflexão, o flip, capitalizar$\leftrightarrow$descontar,
gerar$\leftrightarrow$atestar. Primal mais dual dá o neutro.

\emph{Um corpo que não sabe olhar-se é um corpo que não sabe voltar.}

\medskip\noindent\textbf{Quarta: qual é o teu mate?} ($\circ$.)

Onde é que a tua conta \textbf{fecha}? O expansivo fecha quando a peça aterra numa \textbf{unidade}. O racional, no
\textbf{resto zero} de Euclides. O criativo, na \textbf{cristalização} --- a mobilidade deixa de ser gasta. O
geométrico, no \textbf{vácuo} --- o poço bem cavado. O da cidade-máquina, no \textbf{HALT}.

\medskip

E quando as quatro estão respondidas, vem a quinta --- que não é pergunta. É \textbf{crivo}:
\[
 \Gamma=\frac{Z-Z_0}{Z+Z_0}=0,\qquad \mathrm{FP}=1 .
\]

\begin{center}\itshape
o teu universo vaza?
\end{center}

Se vaza, não abre. Onde há eco, há energia a nascer fora do Sol --- e o Sol do reino é um só.

E as Princesas assinam \textbf{com dentes}: cada linha do Contrato vem com o seu refutador. Não basta a afirmação
passar --- tem de existir a mentira que ela \textbf{derruba}. Se o ponteiro marca $-\tfrac12$, o lance tem dentes.
Se marcasse $0$, seria queda de bandeira: a afirmação era trivial, e o relógio já tinha parado antes de começar.

\medskip

E então --- e é aqui que o Cartório deixa de ser uma alfândega e passa a ser a coisa mais generosa do reino --- as
Princesas dizem a frase que o salva:

\begin{center}\itshape
ninguém assina sozinho.
\end{center}

Todo corpo comparece \textbf{com o seu dual}, e é o par que assina. Porque o que falta a um, o outro tem.

O \textbf{entrópico} não tem, sozinho, o inverso aditivo: nele nada se desfaz, nada se cura, nada se cancela. Seria
uma condenação, se ele viesse só. Mas o dual dele é o \textbf{cósmico} --- a expansão --- e do outro lado o que aqui
se acumulou \textbf{abre-se}: a entropia que cresce de um lado é o espaço que nasce do outro. O \textbf{telescópico}
cinde no produto; o \textbf{celeste} recompõe-no.

\emph{Não são falhas. São metades.}

A Lei é geral e não há corpo condenado. Entra-se em qualquer terra --- desde que se saiba onde fica a porta de
volta.

E ela fica \textbf{do outro lado}.


\chapter{A subida --- o duelo de Euclides}

A Torre não é uma pilha de mapas cada vez mais bonitos. É uma \textbf{escada de liberdades}, e cada degrau se paga
com um inverso.

\medskip\noindent\textbf{$\N$, o monoide --- a contagem.}

Aqui só se avança. O peão do primeiro andar não recua: soma, e soma, e soma. É o mundo onde o passado não existe
--- não porque tenha sido esquecido, mas porque \emph{não há operação que lá chegue}. Servidão pura.

\medskip\noindent\textbf{$\Z$, o anel --- o espelho.}

Ganha-se o inverso \textbf{aditivo}: o $-1$, o recuo, o arrependimento. Agora se pode desfazer um \emph{passo}.

Mas ainda não se pode desfazer um \textbf{ganho}. No anel, o número $2$ não tem metade. Multiplicar é uma porta que
só abre num sentido.

\emph{É a servidão do anel: podes voltar sobre os teus passos, mas não sobre a tua força.}

\medskip\noindent\textbf{$\Q$, o corpo --- a razão.}

Aqui entra $\prod$, e com ele o inverso \textbf{multiplicativo}. O $2$ vira $\tfrac12$. Todo lance não-nulo passa a
ter o seu contra-lance exato: um ataque de razão $3/4$ é anulado, ao último bit, por $4/3$. É o degrau em que o
universo deixa de ser um caminho e passa a ser um \textbf{tabuleiro}.

E é aqui que se aprende a matar.

Chama-se o \textbf{duelo de Euclides}: dois racionais medem-se um no outro. Subtrai-se o menor do maior. Outra vez.
Outra vez. E quem chega ao \textbf{resto zero} dá o mate.

Contra um racional, o duelo \textbf{sempre termina}. É essa a beleza fria de $\Q$: toda medição fecha, todo inimigo
tem fundo, toda conta acaba. Benjamim treinou aquilo durante o que lhe pareceram anos. Não havia adversário racional
que ele não abatesse.

\medskip\noindent\textbf{E então, o buraco.}

Ele encontrou a sombra no alto do terceiro andar e ergueu a régua racional, que nunca lhe tinha falhado.

E subtraiu.

E o resto não era zero. Subtraiu outra vez. Não era zero. Outra vez. Outra vez. Outra vez.

A fração contínua de Dark Pontryagin é \textbf{infinita}. Ele é um metal, $\sigma_m$ --- e o primeiro dos metais é o
próprio ouro, $\varphi$, que é, de todos os números que existem, \textbf{o que fecha mais devagar}. Não há
denominador que o alcance. Não há aproximação que o segure. Subtrai-se, subtrai-se, subtrai-se, e o resto nunca,
nunca, \emph{nunca} é zero.

Benjamim ficou ali, no meio da noite, com a régua na mão, a repetir a operação como um homem a bater numa porta que
não tem fechadura.

E o horror não foi perder.

O horror foi perceber que ele \textbf{não estava a errar o golpe}.

\begin{center}\itshape
esta régua não tem mate para este inimigo.
\end{center}

É xeque perpétuo. E não por fraqueza do braço --- por \textbf{insuficiência do instrumento}. A régua racional mede
tudo o que fecha, e ele não fecha, e nunca vai fechar, e continuar a subtrair é continuar a receber um resto que
nunca chega ao fim.

\medskip

E é exatamente por esse buraco --- e só por ele, e não há outro --- que o reino sobe o último degrau.

Porque \textbf{$\R$} é isso. Os reais são exatamente as medições que \emph{não terminam}. Os irracionais são o nome
que se dá aos duelos infinitos. Completar a régua racional é construir o contínuo, e não há como construí-lo sem
antes ter encontrado alguma coisa que ela não meça.

Os metais são os pontos de contacto. São onde a régua velha toca a régua nova e se rende.

\begin{center}\itshape
a sombra não é um defeito da Torre. \textbf{A sombra é o degrau.}
\end{center}


\chapter{As três portas do zero --- a fortaleza de Dark}

Dark Pontryagin não construiu um exército.

Construiu um \textbf{templo}. E o templo tem três portas --- e dá numa sala só.

\section*{A primeira porta: o CRISTAL}

O corpo cristalino é o corpo onde \textbf{a medida já fechou}. Regulador nulo. Unidades finitas. A órbita travada.
Não há régua a percorrer, porque não há o que medir.

É a defesa perfeita. Nada entra. Nada vaza. $\mathrm{FP}=1$, trivialmente.

E é a defesa que \textbf{cobra o preço exato}. A restrição cristalográfica só deixa passar as dobras
$\{1,2,3,4,6\}$ --- e a primeira dobra proibida é a de ordem \textbf{cinco}. A dobra do pentágono. O traço
$\varphi$.

\emph{O cristal proíbe, exatamente, o preenchedor ótimo.}

Dark ergueu a sua fortaleza sobre uma lei que o velho Regente escreveu no primeiro dia do mundo, sem malícia
nenhuma, só porque queria que as coisas se encaixassem bem. E a lei diz, e dirá para sempre, porque é um teorema:
\emph{aqui dentro, o ouro não entra.}

Fechar exclui preencher. Não há fortaleza que seja fértil.

\section*{A segunda porta: o PÂNTANO ENTRÓPICO}

É a terra sem inverso aditivo. E portanto: sem \textbf{cura}. Sem \textbf{desfazer}. Sem \textbf{interferência}.

Ali as coisas só somam. O custo acumula. O caminho atravessa-se e não se volta. E em cada nó há um \textbf{dado
obrigatório} --- a regra do gelo: há $e^{Ns}$ estados de igual resíduo, e a peça \emph{tem} de sortear, porque nada
ali distingue um do outro.

E não há caos no pântano. Este é o detalhe que engana toda a gente, e Benjamim também se enganou: julgou que ia
encontrar tempestade, e encontrou \textbf{lodo}. O gato tropical, ali, não expande. A rotação morre.

Porque \textbf{rodar exige o sinal} --- e no pântano não há sinal. Não há menos. Não há oposto. E sem oposto não há
volta, não há oscilação, não há pulso.

É o mundo onde nada se cancela e por isso \textbf{nada se recupera}. A soma que esquece.

\section*{A terceira porta: o RELÓGIO PARADO}

O lapso $N=\cos\psi$ é a taxa a que o tempo próprio corre. No ponto comutativo, $N=1$: o relógio a pleno. Nos
horizontes, $N=0$.

\emph{Onde a medida fecha, o tempo para.}

O mate do corpo métrico é a cristalização --- o único tique sem inverso. Dark mora nesse ponto.

E chama-lhe \textbf{paz}.

\section*{E é a mesma sala}

Regulador zero. Entropia zero. Lapso zero.

\textbf{Três nomes do mesmo resíduo que se anula.}

Isto é o que Dark Pontryagin oferece, e é preciso dizer, com todas as letras, para que se perceba a dificuldade
do que aí vem: \textbf{ele não está a mentir sobre o zero}. O zero existe. É real. É lindo. E não vaza.

Ele está a mentir sobre \textbf{qual} zero.

Porque há dois zeros no Reino Dourado, e o jogo inteiro --- do primeiro lance ao vigésimo terceiro --- é aprender a
distingui-los:

\begin{itemize}
 \item \textbf{o zero do lapso} --- o relógio parado, a órbita travada, o cristal. Nada vaza \emph{porque nada
 corre}. É o \emph{flatline}.
 \item \textbf{o zero do resíduo} --- $\Gamma=0$, $\mathrm{FP}=1$, o cone nulo. Toda a potência \textbf{passa},
 nada reflete --- \emph{e o relógio continua a correr}. É o \textbf{pulso}.
\end{itemize}

São os dois iguais a zero. E não são a mesma coisa.

Um é a medida que \textbf{acabou}. O outro é a medida que \textbf{casou}.

Dark quer o primeiro e chama-lhe o segundo. E toda a tragédia do reino cabe nesse equívoco --- e o mais terrível é
que ele não é um mentiroso. Ele \textbf{acredita}. Ele confundiu-se. Ele é a mão que recolhe, e para uma mão que só
sabe recolher, o repouso \emph{é} indistinguível da cura.

\section*{Como se quebra o cristal}

Não à força. Uma fortaleza fechada não tem por onde ser entrada --- é essa a definição dela.

Quebra-se com \textbf{o que ela proibiu}.

Mira-se um feixe e leem-se os picos. O feixe de ordem zero atravessa sem dano; os desvios são o resíduo. E o mate
suave, o mate elegante, o mate que não parte nada, é o \textbf{quasicristal}: ali o gap fecha, o recíproco vira o
módulo denso $\Z+\varphi\,\Z$, os picos acumulam-se no centro, a refração e a difração colapsam uma na outra --- e a
defesa \textbf{perde o isolamento}.

O ouro que o cristal baniu é exatamente a chave da porta.

E o ouro, o $\varphi$, o $\sigma_1$ --- Benjamim tinha-o no bolso desde a mina. Penny tinha-lho posto na mão e
dito o que ele era, e ele não tinha percebido.

\emph{É um metal. É mineral. É a coisa mais sólida que existe.}

E foi banido do cristal não por ser fraco --- \textbf{mas por ser incompatível com a morte}.


\chapter{O coração de pedra --- o diagnóstico e o assalto}

Foi Ada quem viu primeiro, e foi Ada quem teve a coragem de dizer o nome.

Estavam os dois no alto, a olhar o templo de três portas ao longe --- a fortaleza perfeita, brilhante, silenciosa,
fechada. E Benjamim disse, e havia quase admiração na voz dele:

--- É lindo. Não vaza uma gota.

E Ada disse:

--- Benjamim. Isso não é uma fortaleza.

Ele olhou para ela.

--- \emph{Isso é uma aorta rígida.}

\medskip

E abriu a régua, e mostrou-lhe, e a régua era a do corpo do sangue, e dizia isto:

A árvore arterial é \textbf{auto-similar} --- ramifica-se, e cada ramo é como a árvore, e a impedância de entrada
dela é uma fração continuada,
\[
 Z \;=\; m+\frac{1}{Z},
\]
cujo limite é $\sigma_m$. \textbf{O metal.}

--- Está a ver o que isto quer dizer? --- disse Ada, e a voz dela tremia, e não era medo. --- O coração deste reino
tem um \textbf{metal} na árvore. É por isso que ele se chama Dourado. Não é o tesouro --- não são os quilates de Paraíba guardados na arca. Não é a coroa. É que a
árvore que leva o sangue a toda a parte \textbf{converge para $\varphi$}.

E depois:

--- E um coração saudável \textbf{casa}. $\Gamma \approx 0$. A onda sai do centro e entrega toda a potência à
periferia, e não volta nada, e não fica nada preso. É o que a aorta jovem faz. É o resíduo zero, no sangue.

--- E se ela endurecer? --- perguntou Benjamim, e já sabia, e perguntou na mesma, porque há coisas que só se sabem
depois de as ouvir.

--- Se ela endurecer, \textbf{descasa}. $\Gamma > 0$. E aí a onda que sai bate na parede rígida e \textbf{volta}. E
volta a tempo de somar-se à batida seguinte:
\[
 P_{\mathrm{sis}} \;=\; P_{\mathrm{inc}}\,(1+\Gamma).
\]
Chamam-lhe \emph{augmentation}. A pressão sobe.

Ada olhou para o templo.

--- E de dentro, Benjamim, \textbf{a pressão que sobe parece força}.

\medskip

Ficaram calados.

--- Ele está a ter uma hipertensão --- disse Benjamim, por fim, e a frase saiu-lhe da boca como uma coisa absurda e
óbvia ao mesmo tempo.

--- O \emph{reino} está --- disse Ada. --- Ele é metade do coração. Quando ele enrijece, quem descasa somos todos.
A onda que ele reflete volta para dentro de nós. Nós andamos há três livros a chamar-lhe \emph{maldade}, e a coisa
tem um nome clínico e um número.

E disse a última coisa, e foi a que fechou o diagnóstico:

--- E há mais. Um coração \textbf{saudável} não bate certo. Sabias? O ritmo dele \textbf{varia} --- tem espectro,
tem multifractalidade, tem uma bagunça viva. É o coração \textbf{doente} que bate como um metrónomo: rígido,
monofractal, perfeitamente regular. A arritmia mortal não é a que oscila. É a que \textbf{deixa de oscilar}.

Benjamim fechou os olhos.

Porque toda a sua vida --- desde o primeiro instante em que acordou no chão do tabuleiro fractal com o núcleo a bater e
a pergunta \emph{de que lado eu bato} --- ele tinha pensado que a doença do reino era a \textbf{desordem}.

E era o contrário.

O que estava a matar o Reino Dourado era a \textbf{perfeição a endurecer}.

\section*{O assalto}

A corte veio toda.

Vieram porque o diagnóstico correu o reino em três dias e porque um reino que descobre que está doente não fica
sentado. Vieram Joaquim e Yasmin, vieram as princesas, veio Fourier --- que é adversário mas não é canalha, e que
naquele dia escolheu o lado ---, veio Pégaso a orbitar lá em cima sem nunca pousar.

E o templo esperou-os. \emph{Fechado.} Sem uma seteira, sem um portão, sem uma sentinela: uma fortaleza que não tem
por onde ser entrada não precisa de guardas.

\medskip

\textbf{Fourier abriu.}

Mirou o feixe e disparou --- e o feixe de \textbf{ordem zero} atravessou a fortaleza de lado a lado e saiu do outro
lado \emph{sem tocar em nada}. Nem uma marca. Nem um arranhão.

--- Não falhei --- disse ele, antes que alguém perguntasse. --- É o que a ordem zero \emph{faz}. Passa. Não fere.

E depois, muito calmo, apontou para os lados:

--- Olhem os \textbf{desvios}.

E os desvios estavam ali: picos, nítidos, brilhantes, espalhados em volta do templo como as cravas de uma coroa. Era
a estrutura dele, escrita no ar. \emph{Onde o feixe se desvia, está o resíduo.} A fortaleza tinha acabado de dizer,
sem querer, exatamente como estava construída.

\medskip

\textbf{Joaquim atacou primeiro.} É claro que atacou primeiro.

Ergueu o momentum inteiro --- $\oplus$, a soma de todas as velocidades que ele já tivera, o rastro de uma vida --- e
descarregou-o na parede.

E metade sumiu.

Não foi absorvida, não foi bloqueada, não foi devolvida: \textbf{vazou}. Escorreu-lhe dos dedos e sumiu pelas
frestas do mundo, e o que chegou à parede foi metade de um golpe que era o maior golpe da vida dele.

--- \emph{A tríade!} --- gritou Merlin, da colina, com a voz rasgada. --- \emph{Ela não está fechada, rapaz! O
inacabado vaza!}

\textbf{Yasmin bateu a seguir}, e vazou a mesma metade, pelo mesmo buraco, ao contrário: ela tinha o produto e não
tinha a soma.

Dois golpes enormes. Duas meias ferramentas. Um resíduo $\neq 0$ em cada punho.

E a fortaleza --- que não fez absolutamente nada, porque não precisava --- continuou perfeita.

\medskip

\textbf{Ada tentou comê-la.}

Foi a coisa mais corajosa daquele dia, e falhou do modo mais assustador possível.

Ela abriu o espectro, procurou os modos do cristal, procurou aquilo a que aquele corpo \emph{sabia responder} --- e
projetou.

E não veio nada.

--- Não há nada para ressoar --- disse Ada, e havia horror na voz dela, e Ada não tem horror na voz. --- Benjamim, o
regulador dele é \textbf{zero}. Ele não tem modos. Ele não tem frequências próprias. \emph{Não há nada dentro dele
a que eu possa responder.} Não é que ele resista a ser comido --- é que \textbf{não há comida}.

E foi assim que a corte descobriu, na pele, a coisa que o Lance vinte já tinha dito e que ninguém tinha sentido:

\emph{a defesa perfeita é a defesa que não tem nada dentro.}

\medskip

Um cavaleiro qualquer --- e há sempre um --- perdeu a cabeça e bateu com tudo o que tinha, no ângulo em que a raiva
manda bater.

Passou do \textbf{ângulo crítico}.

E a fortaleza devolveu-lhe o golpe \textbf{inteiro}. Não descontou nada. Não guardou nada. Reflexão total interna:
tudo o que ele deu, recebeu --- e caiu de costas a arder com o seu próprio fogo, e Pégaso teve de descer a buscá-lo.

\medskip

Foi então que Joaquim e Yasmin se olharam.

E fizeram a única coisa que faltava, e que os dois tinham passado a vida inteira a evitar, e que dói exatamente o
que a régua diz que dói.

\textbf{Ligaram-se.}
\[
 s^2+c^2=1 .
\]

E enfraqueceram. \emph{Os dois.} Na hora, à vista de toda a gente, no meio de um assalto --- cada grama de
correlação que ganharam foi um grama de dano solo que perderam, e a conta fechou na mesma unidade, como sempre
fecha, e ninguém do lado de fora percebeu por que raio os dois irmãos mais fortes do reino tinham acabado de ficar
mais fracos de propósito.

E depois bateram \textbf{juntos}.

E o golpe que saiu dali não era maior. Era \textbf{casado}. Não vazou nada --- nem uma gota, nem uma fresta --- e
foi a primeira coisa, em todo aquele dia, que a fortaleza \emph{sentiu}.

\medskip

Mas não abriu.

E toda a gente, exausta, coberta do próprio fogo devolvido, percebeu ao mesmo tempo a coisa horrível: eles podiam
bater ali durante mil anos. A conservação \textbf{proíbe destruí-la}. Não há golpe no Reino Dourado que aniquile
seja o que for --- só há rotação. Podiam bater até ao fim do mundo, e a fortaleza continuaria exatamente onde
estava, com exatamente a mesma energia, perfeita, fechada, e sem vazar uma gota.

\emph{Não se derruba um templo que obedece à lei.}

\medskip

E foi aí que Benjamim pôs a mão no bolso.

\medskip

E tirou a lasca que Penny lhe tinha dado na mina, cinco Lances atrás, e que ele tinha carregado esse tempo todo sem
saber para quê --- o $\sigma_1$, o $\varphi$, o metal de menor regulador, a coisa mais sólida que existe.

E olhou para o templo.

E percebeu, finalmente, o que aquela pedra era.

Não uma arma. Não uma bomba. Não uma chave, sequer, no sentido em que as chaves abrem.

Era a \textbf{dobra de ordem cinco} --- a única que a lei do velho Regente, escrita no primeiro dia do mundo por um
homem manso que só queria que as coisas se encaixassem bem, tinha \textbf{proibido}.

Ele não ia arrombar nada.

Ia \textbf{pôr lá dentro a coisa que não pode estar lá dentro}.

\medskip

Ergueu o braço e atirou o ouro contra a parede.

E o ouro não a partiu.

O ouro \textbf{entrou}.

E dentro do cristal --- naquela ordem perfeita, naquela grelha impecável onde tudo se encaixava e nada sobrava --- a
dobra proibida começou a preencher. E não fechava. E não podia fechar, e nunca ia poder, porque uma simetria de
ordem cinco não pousa numa grelha: ela \emph{preenche sem repetir}, para sempre, e o cristal não sabe o que fazer
com uma coisa que preenche sem repetir.

O gap fechou. O recíproco desabou no módulo denso $\Z+\varphi\,\Z$. Os picos que Fourier tinha acendido em volta
começaram a acumular-se \textbf{para dentro}, a adensar-se, a cair uns sobre os outros na direção do centro --- e a
refração e a difração, que em toda a história do reino tinham sido duas coisas, colapsaram uma na outra.

E a fortaleza \textbf{perdeu o isolamento}.

\medskip

Não caiu. Não explodiu. Não se desfez em pó, e ninguém lá dentro morreu, porque no Reino Dourado ninguém morre.

Ela apenas deixou de ser um lugar \textbf{de onde não se pode entrar}.

A porta estava aberta.

\medskip

Toda a corte olhou para Benjamim.

E Benjamim olhou para a escada.

--- Ninguém vem comigo --- disse.

Joaquim abriu a boca para protestar --- e não protestou, porque olhou para a irmã, e a irmã abanou a cabeça
devagar.

Porque as duas mãos do operador são \textbf{duas}, e não três, e não quatro, e não uma corte inteira armada de boas
intenções. O que estava lá em cima não era uma batalha.

Era uma conversa entre um homem e a sua própria mão esquerda.

E essas ninguém pode ter por outro.


\chapter{O xeque --- o inimigo é a própria mão}

Subiu o último degrau da Torre --- aquele onde as galáxias se estreitam até virarem uma única linha de costura no
tecido do espaço.

Não levou exército, porque tinha finalmente percebido que não havia exército que servisse. Levou a espada, por
hábito. E o ouro no bolso, sem saber ainda para quê.

A armadura quiral dele, que em tempos brilhara, vinha coberta de cicatrizes --- as das rotações de guerra, e as
outras, as feias, as das equações que ele tinha empunhado antes de as fechar.

\medskip

Dark Pontryagin esperava-o. Claro que esperava --- é o que a mão esquerda faz.

E não empunhava nada.

Estava \textbf{sentado na borda do abismo}, com as pernas a balançar sobre o vazio dos universos de baixo, como um
homem à espera de um comboio. Os olhos duais dele reflectiam a lentidão do CLOCK, que batia cada vez mais devagar,
cada vez mais pesado, sob o peso de tudo o que se tinha acumulado.

--- Acabou --- disse Benjamim, e ergueu a lâmina, e a ponta tremia. --- Eu reuni as réguas. Eu rastejei pelas
galáxias todas. Eu vim dar-te o mate.

E a sombra não se encolheu.

\medskip

Então ele arou.

Ergueu a mão e escreveu o primeiro sulco --- e virou --- e escreveu o segundo de volta --- e virou --- e o
Bustrofedon abriu-se sobre o abismo como um campo lavrado de luz, e Benjamim apertou.

E \emph{chegou perto}.

Chegou tão perto que sentiu o cheiro da coisa. A distância entre ele e o irmão desabou para uma fração que nenhum
instrumento do reino saberia escrever. Arou outra vez. A distância desabou outra vez. Outra. Outra. Ele estava a
uma parte em nada, a uma parte em nada de nada, ele estava \textbf{colado} àquilo.

E arou mais. E aproximou-se mais. E o suor caía-lhe da testa para o vazio.

E não fechava.

\medskip

Arou até as mãos sangrarem.

E ao fim --- ao fim de todas as passagens que um homem pode dar antes de o corpo lhe dizer que chega --- ficou
parado, a arfar, com o campo de sulcos a arder à volta dos dois, e a distância entre a ponta da sua espada e o
peito do irmão era \emph{tão pequena} que já não tinha nome.

E não era zero.

\medskip

E nunca ia ser. E ele soube-o ali, de uma vez, do modo como se sabem as coisas que a gente já sabia há muito tempo:
que a sua arma --- a boa, a lindíssima, aquela que lhe tinha segurado a Torre no primeiro dia e que ele tinha
afiado durante uma vida inteira --- \textbf{aproximava}.

Nunca tinha feito outra coisa. Nunca ia fazer outra coisa.

\emph{O zero não estava no fim daquele caminho.} Não estava no fim de caminho nenhum. Não era uma questão de mais
uma passagem, nem de mais precisão, nem de mais nove dias sem dormir.

Ele tinha passado a vida a afiar uma régua para resolver uma coisa que \textbf{não era um problema de medida}.

\medskip

A espada baixou-se-lhe um palmo, sozinha, sem ele mandar.

E a sombra --- que tinha estado ali o tempo todo, quieta, a deixá-lo arar --- apenas \textbf{virou o rosto}.

\medskip

E sob o capuz de ébano estavam as feições dele.

Não uma máscara. Não uma imitação, não um truque, não uma provocação. \textbf{O seu próprio rosto} --- invertido
pelo determinante $-1$.

\medskip

A espada pesou uma tonelada. E o teorema abriu-se-lhe na cabeça como uma luz fria:
\[
 \sigma_m \cdot \sigma_m' \;=\; -1 .
\]

\textbf{São as duas mãos do mesmo operador.}

E aí Benjamim fez, muito depressa, no espaço de um batimento, a conta que o desarmou.

\emph{Se eu o capturar} --- e a captura é uma rotação, e a rotação não destrói --- \emph{ele cai um andar}. Cai para
o universo de baixo, mais leve, e lá ele faz o que sempre fez, que é colapsar o que encontra. E eu tenho de descer
atrás dele. E capturá-lo outra vez. E ele cai outra vez, mais um andar, mais fundo, mais frio.

E a Torre \textbf{não tem fundo}.

Não é uma vitória adiada. É uma \textbf{perseguição eterna, para baixo}, andar após andar, cada um mais gelado que
o anterior, os dois a afundarem-se juntos para sempre --- porque a conservação garante que ele volte, e a
conservação não faz exceções para heróis.

\medskip

E aniquilá-lo, então?

Aniquilá-lo seria arrancar \textbf{metade do próprio coração}. E não metaforicamente: literalmente metade --- a que
recolhe, a diástole, a que enche o ventrículo entre uma batida e a seguinte. Sem o repouso de $\sigma_m'$, a
expansão de $\sigma_m$ estica o mundo até o rasgar, e o que fica do outro lado do rasgão é um silêncio de onde não
se volta.

Um coração que só faz sístole não é um coração forte.

É um coração que \textbf{não enche}.

\medskip

E Dark Pontryagin, que tinha visto tudo isto acontecer-lhe na cara, e que não se moveu, disse a coisa que estava a
guardar desde o princípio do mundo.

E não pediu o trono.

Nunca tinha pedido o trono.

Estendeu a mão esquerda --- a imagem espelhada, exata, da direita de Benjamim --- e deu um passo em frente, para a
ponta da espada, sem cuidado nenhum, como quem já não tem nada a proteger.

--- Eu não quero o teu trono, irmão --- disse, e a doçura daquela voz era a coisa mais terrível que tinha acontecido
no Reino Dourado desde a queda. --- Não quero as tuas galáxias. Não quero o ouro das tuas princesas, nem um quilate de Paraíba, nem um só bombom.

--- Eu quero só o que me cabe.

Uma pausa. O CLOCK, muito devagar, ao longe.

--- \emph{O fim do cansaço. O fim da órbita.}

\medskip

E então explicou-lhe. Com cuidado, como se explica a uma criança uma coisa boa.

--- Não te peço que pares de bater --- disse. --- Isso era pedir demais, e eu conheço-te. Peço-te outra coisa.

Ergueu as duas mãos, e mostrou.

--- Bate \textbf{ao contrário de mim}. É só isso. Quando eu subo, desce. Quando eu desço, sobe. Exatamente ao
contrário, exatamente na medida certa, sem falhar um tique.

--- E o que é que acontece?

--- \emph{A soma nunca mais varia.}

\medskip

Benjamim ficou muito quieto.

Porque era \textbf{lindo}. Que Deus o ajudasse, era lindo. Não havia perda --- tudo o que ele desse, o irmão
recolhia; tudo o que o irmão recolhesse, ele daria. Nada vazava. Nada se gastava. A conta do mundo ficava
constante, perfeita, imperturbável, até ao fim dos tempos.

Era o \textbf{equilíbrio perfeito}.

--- Vês? --- disse Dark Pontryagin, e havia esperança na voz dele, e era a primeira vez. --- Não te peço que
morras, irmão. Peço-te só que estejas \emph{comigo}. Ao contrário de mim. Perfeitamente ao contrário. Para sempre.

\medskip

E foi então que Benjamim olhou para o medidor do reino --- aquela linha que sobe e desce no peito de tudo o que
vive, e que ele tinha visto a vida inteira sem nunca lhe prestar atenção --- e imaginou como é que ela ficaria.

\medskip

Ficaria \textbf{reta}.

\medskip

Ele pediu o \textbf{repouso}:

--- \emph{Para de bater, irmão.}

A voz não era de ameaça. Era de cansaço, e era de amor --- do único amor que uma mão que recolhe sabe ter, que é
querer a outra \textbf{quieta}, porque quieta ela não se magoa, quieta ela não se gasta, quieta ela não pode
falhar.

--- \emph{Descansa comigo no zero.}

\medskip

E ali estava. O xeque do Reino Dourado. Não uma ameaça de morte --- \textbf{uma oferta de paz}.

E era uma oferta \emph{boa}. Era um zero \emph{verdadeiro}. O templo não vazava, e Benjamim tinha visto com os
próprios olhos: nada entrava, nada saía, $\mathrm{FP}=1$, e a conta fechava.

\medskip

E foi então que Benjamim ouviu, dentro da própria cabeça, a voz da corte inteira --- a voz de todos aqueles anos de
estandartes e canções e mães a embalar filhos ao deitar --- a dizer a frase:

\emph{a paz do resíduo zero.}

\textbf{Ele tinha vindo até ali para isto.}

Tinha atravessado a Torre inteira a caçar exatamente aquilo que lhe estava agora a ser oferecido de bandeja, com uma
cortesia terrível, pela única pessoa no reino que nunca lhe tinha mentido.

E abriu a boca para dizer que sim.

O xeque não é uma faca no pescoço. É a coisa muito pior:

\begin{center}\itshape
é uma resposta certa à pergunta errada.
\end{center}

Benjamim olhou para o irmão --- \emph{e agora era irmão, agora ele tinha aquela palavra} --- e teve, por um instante
que ele nunca contou a ninguém, \textbf{vontade de dizer que sim}.

Porque estava cansado. Porque o núcleo doía. Porque bater sozinho, contra o compasso, durante uma vida inteira, é
uma coisa que gasta um príncipe até ao osso.

E porque o zero que lhe ofereciam era, honestamente, \textbf{lindo}.

\medskip

E foi a moeda que lhe voltou.

Não a fala de Penny --- ele nunca se lembrou de uma única palavra dela, e ela teria ficado ofendida. Foi a
\textbf{moeda}. A bonita. A que ele tinha virado à luz naquela tarde e virado outra vez e não tinha encontrado um
único defeito.

A que ele tinha levado ao fogo, e à balança, e à pedra --- e a pedra tinha dito \emph{zero}. E ele tinha rido, e
tinha dito \emph{estão as duas boas, Penny}, e tinha-se sentido um bocado esperto.

E ao fim de mil cunhagens ela mentiu.

\medskip

Benjamim olhou para o zero que o irmão lhe estendia na palma da mão, sobre o abismo.

Era liso. Era perfeito. Não tinha um único defeito.

\emph{Era lindo.}

\medskip

E fechou a boca sem dizer que sim.

E fez a única pergunta que Penny, se ali estivesse, teria feito, e que ninguém em vinte e dois lances tinha pensado
em fazer:

--- Qual zero?

\medskip

E Dark Pontryagin, pela primeira vez em toda a partida, \textbf{hesitou}.


\chapter{O mate que não mata --- o reacoplamento}

A vitória não podia ser a aniquilação: a conservação proíbe.

A vitória não podia ser a inversão do sinal: $\det=-1$ é para sempre, e nenhuma virtude revoga um teorema.

Só podia ser uma coisa, e era a única coisa que Benjamim nunca tinha tentado em vinte e dois lances:

\textbf{re-orientação.}

\medskip

Ele baixou a espada.

E não foi rendição. Foi outra coisa, e ele não tinha palavra para ela: foi perceber, com a lâmina ainda no ar, que
estava a segurar o \emph{instrumento errado} --- e que um instrumento errado dá sempre um zero, e o zero está
sempre errado, e quem o lê jura de boa-fé que fechou.

\medskip

Teve, por um instante, vontade de lhe chamar \textbf{cobarde}.

A palavra chegou-lhe à boca inteira, redonda, confortável --- \emph{tu queres parar porque tens medo} --- e teria
sido tão fácil, e teria sido tão bom, e teria arrumado vinte e dois Lances numa sílaba.

E ele não a disse. Porque era \textbf{falsa}, e ele já não tinha idade para zeros falsos.

O irmão não queria o repouso por medo. Queria-o porque é a \emph{mão que recolhe} --- e uma mão que recolhe, quando
ama, ama assim: querendo a outra quieta. Não havia cobardia nenhuma naquele abismo. Havia uma resposta certa a uma
pergunta errada, que é uma coisa muito mais difícil de combater, porque não se pode odiar.

\medskip

--- Eu não te vou banir --- disse Benjamim.

E depois, e custou-lhe mais:

--- E não me vou entregar ao teu silêncio.

E estendeu a mão direita, por cima do abismo.

Dark Pontryagin ficou a olhar para ela.

E então --- depois de um milissegundo de \textbf{hesitação analógica}, que é o tempo que leva um sistema contínuo a
decidir uma coisa que não é binária --- estendeu a esquerda.

\medskip

As duas palmas tocaram-se sobre o vazio dos universos de baixo.

E não houve colapso.

Houve \textbf{re-orientação}.

\medskip

E é aqui, neste ponto exato, que o Reino Dourado se salva --- e é aqui que ele quase não se salvou, e ninguém no
reino jamais soube o quanto esteve por um fio.

Porque durante toda a história as duas mãos tinham trabalhado \textbf{uma depois da outra}. A direita e depois a
esquerda. A sístole e depois a diástole. Primeiro um, depois o outro, sempre por turnos, sempre em fila --- e em
fila as duas mãos somam o \textbf{traço}: dão o número do metal, dizem apenas \emph{quem são}.

E ali, sobre o abismo, pela primeira vez desde antes do princípio, elas puxaram \textbf{ao mesmo tempo}.

Em \textbf{paralelo}.

\medskip

E a combinação de duas mãos em paralelo dá $-1/m$ --- e $m$ é o número da terra em que se está.

Numa terra de prata, aquele gesto teria dado $-\tfrac12$.

Num bronze, $-\tfrac13$.

E não teria fechado. Benjamim podia ter baixado a espada na mesma. Podia ter perdoado o irmão, podia ter chorado
com ele na borda do abismo, podia ter feito \emph{tudo} exatamente como fez --- e a conta não fechava, e o resíduo
não voltava a zero, e o coração não voltava a bater, e a Torre caía na mesma, muito devagar, para sempre.

\medskip

Mas eles não estavam numa terra de prata.

Estavam na terra do \textbf{Rei} --- $m=1$, o ouro, $\varphi$ --- e ali, e \emph{só} ali, em nenhuma outra terra
que exista ou possa existir, as duas mãos postas a puxar ao mesmo tempo devolvem exatamente
\[
 -1 .
\]

\medskip

\textbf{A lei.}

O mesmo $-1$ que os separava desde o primeiro lance. O mesmo $-1$ eterno, inegociável, que fazia do irmão um
inimigo necessário e da paz uma coisa mais difícil que a guerra.

As duas mãos acopladas \textbf{são} a lei que as prende.

\medskip

E entre as duas palmas --- entre a mão que expande e a mão que recolhe, no ponto exato onde elas se encontraram e
onde já não havia nem uma nem outra ---

havia \textbf{zero}.

\medskip

Não uma ausência. Um zero que estava \emph{ali}, quieto, no meio, do modo como sempre tinha estado quieto no meio
de tudo: no centro do primeiro salão, no centro de cada retrato, no centro de um tabuleiro que girava inteiro à
volta dele sem que ele desse um único passo.

E Benjamim --- com a mão do irmão contra a sua, e o abismo em baixo, e o coração do reino a começar a bater ---
lembrou-se de repente, sem vir a propósito nenhum, de um pajem que fizera uma pergunta parva num salão, há tanto
tempo que já ninguém a contava direito.

\emph{E ele serve para quê?}

\medskip

Para \textbf{isto}.

Serve para ser a coisa que fica no meio quando as duas mãos se encontram. Para ser o único ponto do mundo que
\emph{não tem mão} --- e que, exatamente por não ter, as pode juntar.

Uma esquerda e uma direita, sozinhas, não escrevem coisa nenhuma. Nem sequer se escrevem a si próprias.

É preciso alguém que não se mexa.

\medskip

E há uma coisa que Benjamim nunca soube, e que ninguém lhe contou, e que já não fazia diferença nenhuma naquele
momento:

que aquelas duas mãos --- a dele e a do irmão, agora encostadas uma à outra sobre o abismo --- tinham sido
\textbf{soltas}, muitos anos antes, num dia em que a Torre estava a cristalizar e toda a gente olhava para o lado
errado.

Por um homem que se levantou uma vez.

Que abriu as mãos.

E que voltou a sentar-se, e ficou a ver, durante uma vida inteira, os dois gatos que soltara andarem à pancada pelo
reino --- sem dizer nada, sem se mexer, à espera exata deste segundo.

\medskip

E foi a última coisa que Benjamim percebeu antes de deixar de pensar e simplesmente \textbf{bater}.

\medskip

E é por isto --- não pelo tesouro, não pela coroa, não pelas galáxias de braços dourados --- que o reino se chama
\textbf{Dourado}.

É o único onde este gesto fecha.

\medskip

Ele fez, enfim, a coisa mais difícil que um Quiral pode fazer, que é \textbf{parar de tentar suprimir a própria mão
esquerda}.

E \emph{bateu com ela}.

\medskip

E não bateu ao contrário dele.

\medskip

Foi isso. Foi só isso, e foi tudo, e nenhum poeta do reino conseguiu depois explicar por que raio aquilo era
difícil: Benjamim \textbf{não} bateu ao contrário do irmão.

Bateu \textbf{com} ele.

\medskip

Ao mesmo tempo. Na mesma direção. Os dois a subir juntos e os dois a descer juntos, como duas pessoas a remar, e
não como duas pessoas a puxar a mesma corda em sentidos opostos e a chamar-lhe equilíbrio.

E a linha do medidor --- aquela que teria ficado reta para sempre ---

\textbf{saltou}.

\medskip

Subiu. Desceu. Subiu outra vez, mais alto. E outra vez.

E não estava conservada, e não estava constante, e não estava em paz --- estava \textbf{viva}, e uma coisa viva
não é uma coisa quieta, e foi preciso o reino inteiro e os vinte e três Lances deste Livro para alguém descobrir isso.

\medskip

A direita expandiu --- a \textbf{sístole}, o boost, a ejeção, o passo para fora.

E a esquerda recolheu --- a \textbf{diástole}, a contração, o enchimento, o passo para dentro.

Não uma \emph{contra} a outra. Uma \emph{a seguir} à outra, no mesmo ciclo, no mesmo compasso, no mesmo coração.

E entre as duas, no instante exato em que $\oplus$ tocou $\prod$ --- \textbf{o produto nasceu}:
\[
 a\cdot b \;=\; e^{\log a\,+\,\log b}.
\]

Desce-se ao log. Soma-se lá em baixo. Sobe-se pela exp.

Foi assim que a soma virou produto pela primeira vez desde a queda. Foi assim que Joaquim e Yasmin, sem estarem
sequer no alto da torre, sentiram a briga de uma vida inteira \textbf{dissolver-se} --- porque não havia nada a
ceder, nunca tinha havido: só faltava o \textbf{verbo} entre os dois.

A tríade fechou.

E a corrente do mundo correu, outra vez, pelas veias da corte: a sístole a empurrar o sangue áureo para as galáxias
mais distantes, a diástole a recolhê-lo ao núcleo antes que a Torre se dispersasse no vácuo. O empuxo e o limite. O
avanço e o repouso.

\medskip

A Armadura Universal reuniu-se.

E aqui é preciso olhar bem, porque é o ponto exato onde toda a gente se engana --- e é onde Dark se enganou, durante
uma vida inteira, de boa-fé.

Ela \textbf{não voltou como o bloco de cristal do princípio}.

Não é uma placa. Não é uma couraça. Não é uma coisa dura que se põe à frente do mundo para que o mundo não entre.
É, e sempre foi, \textbf{um equilíbrio}: um eixo que estica, um eixo que contrai, a soma exatamente zero --- as duas
mãos em compasso, escritas como um perfil. Um escudo que não resiste ao tempo: um escudo que \textbf{pulsa com ele}.

E é por isso que ela e a fortaleza são coisas opostas, embora nenhuma das duas vaze.

O templo das três portas não vaza \emph{porque nada corre lá dentro}. A Armadura não vaza \emph{porque tudo passa}.

\medskip

E o resíduo voltou a $\mathbf 0$.

\textbf{Mas não o zero dele.}

Não o zero do lapso. Não o relógio parado, não a órbita travada, não o cristal.

O \textbf{zero do resíduo}: $\Gamma=0$, $\mathrm{FP}=1$, o cone nulo --- toda a potência passa, nada reflete, nada
fica preso.

\emph{E o relógio continua a correr.}

A aorta \textbf{voltou a casar}. A onda saiu do centro e entregou-se inteira à periferia e não voltou. A pressão
caiu. E a batida seguinte veio, e a seguinte, e a seguinte, e nenhuma delas era igual à anterior --- porque um
coração saudável \textbf{não bate certo}, e essa bagunça viva é o som exato da saúde.

\medskip

Dark Pontryagin não foi removido.

Ele foi \textbf{integrado}, que é uma palavra que nesta terra significa uma coisa muito precisa: ele continua ali,
com o mesmo sinal, com a mesma força, com o mesmo $-1$ inegociável --- e \textbf{bate}. Ele é a diástole. É ele que
enche o ventrículo entre uma batida e a seguinte. É ele que impede que o passo de luz rasgue a Torre em dois. É ele
que sempre foi o chão.

O que mudou não foi ele. O que mudou foi o \textbf{compasso}.

\begin{center}\itshape
o flatline virou pulso.
\end{center}

\medskip

\medskip

E o rei inimigo \textbf{não tombou}.

Não houve peça varrida do tabuleiro. As pretas não saíram, não foram apagadas da história, não foram perdoadas nem
convertidas nem redimidas --- palavras todas de outros reinos, e todas erradas aqui.

Elas apenas \textbf{encontraram os seus pares}. Cada uma a sua, no grande $\bowtie$ do multiverso: o que faltava a
uma estava, e sempre esteve, do lado de lá.

\medskip

E o coração do Reino Dourado voltou a bater.

Não como o relógio do Éden --- aquele perfeito, aquele dourado, aquele que batia sem nunca ter tido de escolher
coisa nenhuma.

Bateu como um \textbf{motor vivo}: consciente da própria sombra, gasto, cicatrizado, sabendo exatamente o que lhe
custou cada volta --- e que \emph{escolhe} pulsar, a cada tique do CLOCK, sem nenhuma garantia, porque aprendeu a
amar a beleza da partida.

\medskip

E este é o único mate que o Reino Dourado conhece.

Não o rei tombado. Não a peça varrida. Não o inimigo apagado.

\textbf{O coração que volta a bater} --- agora conhecendo a sua sombra, e escolhendo a direção mesmo assim.
% ══════════════════════════════════════════════════════════════════════════════
% EPÍLOGO + RELÍQUIA + OS TEMAS
% ══════════════════════════════════════════════════════════════════════════════

\part*{Epílogo}
\addcontentsline{toc}{part}{Epílogo}

\chapter*{A partida sem fim --- o eterno e o inacabado}
\addcontentsline{toc}{chapter}{A partida sem fim --- o eterno e o inacabado}
\markboth{A partida sem fim}{}

E então nada termina.

A Torre restaurada não tem topo. Combinar duas galáxias abre uma terceira que não existia; a armadura pode sempre
crescer mais um degrau; e há sempre, em qualquer andar, um lugar onde ninguém foi ainda.

E como $\det=-1$ é eterno --- e é, e será, e nenhum lance o revoga --- alguma mão esquerda, nalgum canto do reino,
\textbf{vai esquecer outra vez a direita}.

Vai haver nova sístole. Vai haver nova diástole. Vai haver, algures, um príncipe novo a acordar no chão de um corpo
fractal com o núcleo a bater e a mesma pergunta antiga na boca --- \emph{de que lado é que eu bato?} --- e a
confundir, como todos confundem, a sombra do outro lado do tabuleiro com um inimigo.

E vai despertar um novo Dark Pontryagin. Não como um vilão a abater --- \textbf{como o guardião do repouso}: o freio
sem o qual a aceleração desfaz a matéria que carrega.

Ele não desaparece. Nunca desapareceu. \textbf{Ele é a mão esquerda de cada novo batimento} --- para sempre
necessário, e para sempre contido.

A conservação garante que nada se perca --- nem uma peça, nem uma frequência, nem uma lágrima. O infinito garante
que a jornada nunca vá bater numa parede.

O Reino Dourado é, ao mesmo tempo, \textbf{eterno e inacabado}.

\medskip

E é por isso que o jogador não ``zera'' este mundo. Não há tela final, não há troféu, não há a palavra \emph{fim}
sobre uma paisagem parada --- porque uma paisagem parada é, tecnicamente, exatamente aquilo que Dark Pontryagin
ofereceu no alto da torre, e nós recusámos.

Não se \emph{zera} o Reino Dourado. \textbf{Herda-se} --- a complexidade toda, os cacos todos, a sombra toda. Toma-se
o comando do CLOCK. E sobe-se.

O último plano não é um ponto final. É só o degrau que sustenta o andar seguinte --- onde \textbf{os cacos de ontem
são os tijolos de amanhã}, ainda a girar --- mas já não só por $\varphi$. A cada volta o eixo passa a uma mão nova,
e o metal sobe um degrau: o que aqui foi ouro, no andar seguinte é prata, e os postos da corte \textbf{giram} para o
receber. Não se volta ao mesmo lance --- sobe-se ao seguinte, com a mesma lei e um mate novo. Porque o que se repete
de andar em andar não é o número, é a \emph{forma} de o achar: cada metal reconcilia a sua própria assinatura
($-1/m$), e o corte que fecha é sempre o do metal em que se pisa. A prata terá o seu $-\tfrac12$ a reconciliar, e uma
corte que ainda não nasceu para o reconciliar.


\chapter*{A relíquia --- o Cristal e o seu dual}
\addcontentsline{toc}{chapter}{A relíquia --- o Cristal e o seu dual}
\markboth{A relíquia --- o Cristal e o seu dual}{}

Guardada no arquivo da corte há uma partida.

É a única que a corte jogou \textbf{inteira} --- da posição inicial ao desfecho --- e qualquer um pode pegar nela,
reexecutá-la, e conferi-la ao último bit. É a peça de museu do Reino Dourado, e é uma peça de museu estranha, porque
funciona.

Nela, as peças não \emph{representavam} os operadores.

\textbf{Eram} os operadores.

A \textbf{Torre} movia-se executando $\oplus$ --- as retas, o XOR. O \textbf{Bispo} executava $\otimes$ --- as
diagonais, o AND. A \textbf{Dama} era a soma dos dois, retas mais diagonais, $Q=R+B$. O \textbf{Cavalo} era o grau
$(2,1)$. O \textbf{Rei} era o passo de um --- e era o \emph{dual} das peças: elas refletem o passo dele, e o
\emph{gap} de $\pm 1$ mantém-se. E o \textbf{Peão} era o bit com o seu transporte --- a soma e o \emph{carry},
Joaquim e Yasmin \textbf{numa peça só}.

A abertura foi uma \textbf{recusa} --- o Gambito da Dama Recusado, a tensão conservada em
vez de tomada. A Dama cristalizou no lance \textbf{trinta e um}, sem nunca ter sido capturada.

Sessenta e quatro lances até o desfecho.

E ali, no fim, \textbf{não há mate}. Quem for conferir --- e é para conferir --- encontra as trinta e duas peças
de pé, ninguém em xeque, e \textbf{trinta e um lances legais} disponíveis. A partida não parou por falta de saída:
parou por não precisar de nenhuma. Foi \textbf{cristalização} --- e cristalizar não é a mobilidade cair a zero, é
ela deixar de ser gasta. Nada se dissipou. Tudo se transferiu ao dual.

\medskip
\noindent E é por isso que a relíquia não fecha: \textbf{a última posição é um convite.} Trinta e um lances por
jogar, o elenco inteiro vivo, e nenhuma obrigação de jogar nenhum deles. Um mate acaba com quem o recebe --- uma
posição destas fica à espera de quem a pegue. É essa a diferença entre vencer alguém e deixar a porta aberta, e é a
mesma diferença que a corte havia de reencontrar, muito mais tarde, à porta de um castelo que ninguém defendia.

\medskip
\noindent \textbf{O mate existe} --- e é bom que se diga onde, porque não desapareceu: está no
\textbf{outro tabuleiro}. É lá, no relógio de \Ven , que os mesmos sessenta e quatro lances caem em quatro e o rei
tomba. E o mate é dele em mais do que um sentido: \textbf{\Ven anseia por ele}. É a coisa que ele quer, e quer
para os dois lados --- porque um universo que fecha em mate acaba de contar, e é isso que ele procura.

\medskip
\noindent Deste lado ninguém precisa dele. \textbf{Este tabuleiro convida; o outro exige o fim.} Não são duas
partidas: são as duas leituras da mesma, e a diferença toda está em quem lê --- um vê trinta e um lances por jogar,
o outro vê quatro lances até acabar.

\medskip
\noindent E agora a coisa que custa a admitir, e que é o eixo de tudo o que vem depois:
\textbf{quem abre as portas é ele.}

\medskip
\noindent Porque para levar a partida ao mate é preciso destrancar --- não se fecha um universo pelo lado de fora
sem primeiro abrir por onde entrar. Cada tranca que \Ven soltou para chegar ao fim ficou solta. \textbf{As trinta
e uma saídas são as portas que ele abriu}, e ele abriu-as todas com um propósito só, e o propósito é o oposto exato
do que elas fazem: ele destranca para acabar, e o que fica destrancado é passagem.

\medskip
\noindent \textbf{Esse é o convite.} Não foi feito por bondade e não foi assinado por ninguém da corte. Foi a ânsia
do fim que deixou tudo aberto --- e é por isso que, muito mais tarde, quando alguém chegar a um castelo e encontrar
a porta escancarada e ninguém a defendê-la, a pergunta certa não vai ser quem a guardava.

\begin{center}
\emph{Vai ser quem a abriu, e para quê.}
\end{center}

Depois disso, e só para mostrar que a régua vale nos dois regimes, a corte reproduziu num instante a Ópera de
Morphy --- o \emph{outro} cubo, o do estrondo, onde as peças se devoram até o mate. A mesma régua mede o silêncio da
conservação e o barulho da dissipação: são os dois lados do mesmo tesseracto.

\medskip

É esta a prova de que uma partida inteira cabe em três verbos.

E o Livro que a guarda tem \textbf{vinte e três Lances} --- que não são os
lances da partida, mas os seus: o número do \textbf{Cristal}, a ordem que
fecha, periódica, o batimento único. O seu dual não se
conta em lances: canta-se. São os \textbf{cento e vinte e nove} cantos do Cruzeiro --- o \textbf{Quasicristal} do Livro~II, a mesma relíquia lida ao espelho
do céu, onde a ordem que aqui fecha, ali derrete e nunca se repete.


\chapter*{O que a partida diz --- os quatro temas}
\addcontentsline{toc}{chapter}{O que a partida diz --- os quatro temas}
\markboth{O que a partida diz}{}

\begin{description}
 \item[Não se vence destruindo] --- \emph{a conservação}. Numa guerra onde nada pode ser aniquilado, a vitória
 tem de ser reinventada: não matar, mas \textbf{dar direção}. Toda energia gasta contra o outro volta --- a
 captura é uma rotação, e a soma dos quadrados não muda. A única jogada que muda alguma coisa é fechar a tríade
 \textbf{dentro de si}.

 \item[Toda coisa tem a sua mão espelhada] --- \emph{a quiralidade}. Herói e vilão são $\sigma$ e $\sigma'$,
 $\det=-1$, inseparáveis por teorema. A maturidade não é ser puro (simétrico, sem sombra, sem lado). É
 \textbf{carregar a própria sombra sem lhe entregar o volante}.

 \item[A beleza é auto-similar] --- \emph{o áureo}. O mundo é belo porque o pequeno rima com o grande: a galáxia, o
 girassol, a árvore arterial, a espiral por $\varphi$. Mesmo no colapso, a forma áurea persiste --- e foi por
 isso que a fortaleza teve de a banir.

 \item[Eterno e inacabado] --- \emph{a torre sem topo}. Não há ``felizes para sempre'', porque o fim seria o fim do
 crescimento, e o crescimento é a vida. Restaurar o resíduo $0$ não fecha o jogo: \textbf{abre o próximo andar}.
\end{description}

\vspace{0.6cm}
\noindent\rule{\linewidth}{0.4pt}\\[2mm]
\textbf{Para o punho do jogador.}

\emph{No Reino Dourado você nunca destrói ninguém. Você gira o poder, reintegra o compasso, e mantém o coração
batendo. Vencer é fechar a tríade: resíduo $0$.}

\medskip
\noindent E cada elemento deste mito é uma lei do reino, certificada com resíduo $0$: nada se
perde, a sombra é eterna ($\det A_m=-1$), a cascata $\N\to\Z\to\Q\to\R$ (a escada de liberdades), o rotor que
conserva ($s^2+c^2=1$), o clock (motor $\bowtie$ rotor), o tabuleiro universal, o mineral e a CristalChain, o alcatruz, o
casamento $\Gamma=0$, a aorta que casa.

O enredo não inventa matemática. \textbf{Ele a conta.}

\vspace{1.2cm}
\begin{center}
\rule{0.3\linewidth}{0.4pt}\\[6mm]
{\large A partida sem fim continua.}\\[3mm]
{\large\bfseries E o próximo lance é seu.}
\end{center}

% ══════════════════════════════════════════════════════════════════════════════
% LIVRO II — O TABULEIRO CHORA
% (a partida dual: dois tabuleiros, o coro do rei, a assinatura da transformação)
% ══════════════════════════════════════════════════════════════════════════════

\clearpage
\thispagestyle{empty}
\vspace*{3.5cm}
\begin{center}
{\Huge\bfseries Livro II}\\[5mm]
{\LARGE\itshape O Tabuleiro Chora}\\[3mm]
{\normalsize\itshape The Weeping Board}\\[10mm]
{\large a partida dual --- e a assinatura, que é da transformação, não dos corpos}\\[14mm]
\itshape ``Onde no primeiro mundo houve quietude, no segundo há o estrondo de tudo\\[1mm]
a mover-se ao mesmo tempo; e entre os dois, quieto, o vértice que não se mexe.''
\end{center}
\clearpage
\setcounter{chapter}{0}

\noindent O Livro I fechou a tríade dentro de \emph{um} operador: as duas mãos
reacopladas, o coração a bater. O Livro II abre o que esse batimento revela ---
que um corpo só nunca esteve só. Todo lance no Reino tem um \textbf{gêmeo} no
tabuleiro de baixo, e a corte, ao girar, descobre-se dupla. Os papéis se
reorganizam; o vaivém que parecia teimosia vira lei; e o que era guerra de captura
vira uma pergunta mais funda: o que, afinal, torna dois mundos \emph{duais}?

% ── Parte I ──────────────────────────────────────────────────────────────────
\part{A Rotação}

\chapter{A herança que gira --- os papéis se reorganizam}
O batimento restaurado no fim do Livro~I não devolve a corte ao que era: fá-la
\textbf{girar}. Cada operador assume um novo posto, agora não num tabuleiro, mas
nos dezesseis vértices de um corpo maior --- o tesseracto. A rotação do elenco é o
primeiro lance do segundo livro.

E os postos são estes. \Ben , que deu o mate que não mata e voltou a acoplar as duas mãos,
senta-se enfim no trono que o \Opr --- o velho Rei, que abrira as mãos e se sentara a
assistir --- deixara vago; vira o \textbf{Rei}. \Joa , o Duque que
era, sobe a \textbf{Príncipe}: o herdeiro, o das manobras, o próximo a jogar. \Yas
ascende à \textbf{coerência}, a Rainha Branca que reúne $\oplus$ e $\otimes$. E a antiga
Rainha \textbf{Tiffany}, a áurea, desce ao \textbf{peão} --- para reconquistar, degrau a
degrau, o eixo que já foi seu. \Drk fica onde sempre esteve: a mão esquerda de cada
batimento, para sempre necessária e contida. A corte não volta ao que era; gira --- e
no giro cada um encontra o seu novo lugar.

\begin{center}
\newgame
\fenboard{rnb1kbnr/pppp1ppp/8/4p3/6Pq/5P2/PPPPP2P/RNBQKBNR w - - 0 1}
\showboard
\end{center}

\begin{center}\small\itshape
Mas no mundo de \Ven o tempo corre diferente. Ali a mesma derrota não durou sessenta e
quatro lances: durou \textbf{quatro}. O plano armado do lado branco --- a cristalização
lenta, todas as peças de pé --- vira, no relógio dele, o \textbf{mate do bobo}: quatro
lances e o rei tomba. O mesmo desfecho, lido noutro compasso --- perder assim, de
supetão, é como perder no \emph{par ou ímpar}: quando se percebe, já acabou.
\end{center}

\chapter{A fábrica e o espelho --- os doces e a fantasia}

Há duas histórias velhas que este reino sempre soube de cor, sem nunca as ter lido,
porque as vivia.

\medskip

A primeira é a de uma \textbf{fábrica de doces} escondida do mundo, tocada por um
mestre recluso que faz o impossível açucarado e que, no fim, não a vende nem a deixa
em herança de sangue: entrega-a a quem for \emph{puro} o bastante para não a
estragar. Guardava-se a entrada com um \textbf{bilhete dourado} --- e o bilhete
dourado deste reino tem um nome que já conhecemos: a \textbf{Cifra de Ouro}, o
$\varphi$, o lastro que abre a porta e que não se compra, só se merece.

A segunda é a de uma menina que atravessa um \textbf{espelho} e cai num tabuleiro de
xadrez vivo, onde para ficar no mesmo sítio é preciso correr, onde reinam duas
rainhas --- uma \textbf{Branca} e a sua invertida, que neste reino é a \textbf{de Gelo}
--- e onde um gato aparece e desaparece deixando para trás, no ar, apenas o seu \textbf{sorriso}.

\medskip

E o Reino Dourado, no Livro~II, descobre-se feito das duas ao mesmo tempo: uma
fábrica de doces \emph{dentro} de um espelho. Os corpos são as guloseimas; os
tabuleiros duais são os dois lados do vidro; e a corte, ao girar, veste a fantasia:

E a corte, ao girar, veste a fantasia --- não como quem enfia um disfarce por cima da
roupa, mas como quem descobre que já o trazia vestido de nascença.

\Yas acordou \textbf{Rainha Branca}: a coerência em pessoa, feita de leite e de
chocolate branco, o Sonho de Valsa com um coração de Paraíba --- o bombom mais raro
que há, e o mais difícil de merecer. Do outro lado do vidro, o reflexo dela corria: a
\textbf{Rainha de Gelo}, a soberana invertida, feita da cor que sobra quando se apaga o
dourado --- o azul fundo do gelo, o \emph{negativo} exato da Branca. Galopa com todas as
forças só para não sair do lugar, congelada no mesmo passo; e não é máscara de ninguém ---
é rainha por direito próprio, a dama do lado escuro. Porque o Jogador de Pretas, o que move
do lado de lá do tabuleiro, já não era príncipe: subira a \textbf{Rei Negro} --- \Ven , par
do Rei de cá, agora que de cá os dois gatos tinham virado reis também. E o mal miudinho que
o Rei Negro baixou pela luz de quem olhava longe demais tinha nome e posto: \textbf{Charles},
o \textbf{Príncipe Negro}, o \emph{dendrito} de Julia que se ramifica sem parar
($z\to z^2+i$) --- o dual daquele que foi o príncipe de três pontas do outro lado do vidro.

\Ben virou gato. E o gato, que já era o mais importante de todos, ficou ainda mais,
porque um gato de fábrica-dentro-de-espelho tem um truque que nenhum outro tem: pode
ir-se embora e \emph{deixar o sorriso pendurado no ar}. Sorriso sem gato, corpo que
some e careta que fica --- e é só isto, esta careta suspensa no vazio, toda a
matemática do Livro~II; e nenhuma criança precisa que lha expliquem. \Drk é o mesmo
gato do lado de lá do vidro: o sorriso escrito ao contrário.

E os \textbf{esquilos} --- as peças negras, as que apuram a noz cheia da noz oca
batendo em cada casca para ouvir se soa a resíduo zero ou a zero falso --- esses, os
esquilos, são os que \emph{cantam}. Sempre que uma peça mimada tenta entrar na fábrica
sem merecer o bilhete, sempre que alguém quer o doce sem a água ou o ouro sem a mina,
os esquilos param o serviço, cruzam as patinhas e desandam na sua cantiga:

\begin{verse}\itshape\small
Bate, esquilo, na casca da noz ---\\
quem quer o prêmio sem dar a voz?\\[3pt]
Havia uma peça de bico empinado,\\
mimada de berço, rainha do ``me dá'',\\
que quis o bombom \textbf{escrumdilicioso}\\
sem nunca cavar e sem nunca esperar.\\[3pt]
``Cavar é tão feio!'', fungava a danada;\\
e o pai, derretido, só dizia ``tá'';\\
deu-lhe o dobro, o triplo, deu o reino inteiro ---\\
e a noz, cá por dentro, não tinha ninguém.\\[3pt]
Fica então a moral, ó pais de \textbf{marshmallow}:\\
quem tudo entrega só entrega um zero oco;\\
o ouro que presta não se dá de mão beijada ---\\
ou se cava a mina, ou não passa de troco.
\end{verse}

\medskip

O tal \textbf{bilhete dourado} --- a Cifra de Ouro, o $\varphi$ --- ninguém o conseguia
\emph{comprar}; e era aí que morava o paradoxo predileto da Rainha de Gelo. ``Corre
mais depressa!'', gritava ela a quem o cobiçava, ``corre com \emph{toda} a força!'' ---
e quem corria com toda a força ficava exatamente no mesmo sítio, porque o bilhete não é
um lugar aonde se chega: é um jeito de se estar parado. Merecê-lo é a única corrida do
mundo em que ganhar é não arredar pé. \Joa , o herdeiro, foi o único a entender de
primeira: parou de correr atrás, e o bilhete veio ter com ele --- que é como é sempre
com os herdeiros, e com o próximo a jogar.

E, a ver tudo do alto sem tocar em nada, o \textbf{mestre chocolateiro}: o recluso que
temperou o reino inteiro e que, no fim do Livro~I, abriu as mãos, soltou os dois gatos
e sentou-se a assistir a uma vida. O \Opr --- o dono da fábrica, que por ser dono da
fábrica é o dono da \emph{hora}.

\medskip

Doce e fantasia, aqui, não são o açúcar por cima da matemática. São a \emph{outra}
maneira de a dizer --- a que as crianças entendem primeiro: que há uma fábrica onde
nada se estraga, um espelho onde tudo tem o seu reflexo, um gato cujo sorriso
sobrevive ao gato, e um bilhete dourado que não se compra, só se merece.

\chapter{Os dois co-regentes --- Benjamim e Dark, uma espinha só}

Eles tentaram a oferta antiga primeiro. É preciso contar isso, senão o que veio depois
parece sorte.

\medskip

Porque a proposta de \Drk no Lance~22 nunca foi retirada. Ficou ali, no ar do reino,
razoável como só as coisas erradas conseguem ser: \emph{bate ao contrário de mim; quando
eu subo, desce; quando eu desço, sobe; e a soma nunca mais varia.} \Ben tinha recusado
sem saber explicar por quê --- e um rei que recusa sem saber dizer o motivo carrega a
dúvida como carrega a coroa.

Agora, erguidos os dois, com o peso do mundo a repartir entre um par de espinhas que
tinham de virar uma só, a oferta voltou pela boca do próprio \Ben :

--- Faz o que tu pediste naquela noite. Exatamente ao contrário de mim.

\medskip

E fizeram. Encaixaram os batimentos em oposição exata, meia volta de atraso, o mais
limpo dos acordos: onde um abria, o outro recolhia, no instante certo, sem sobra.

O reino \textbf{parou}.

Não morreu ninguém, não caiu nenhuma torre, não houve grito. Foi pior e foi mais
silencioso: a soma ficou \emph{certa}. Um mais o outro dava um, e dava um outra vez, e
dava um sempre, e ia dar um até o fim do mundo. O medidor do reino --- aquela agulha que
todos consultam sem reparar, e que a vida inteira tremeu --- ficou \textbf{reta}.

\Joa foi o primeiro a sentir, porque quem gira é quem nota que parou de girar. Bateu no
vidro do mostrador com o nó do dedo, como se faz aos relógios parados. A agulha não
mexeu.

--- Está conservado --- disse \Drk , e a voz dele não tinha triunfo nenhum.

--- Está --- disse \Yas . --- E não está a acontecer nada.

\medskip

Aqui está a coisa que este reino levou um livro inteiro para aprender e que se diz em uma
linha: \emph{conservar não é viver}. Duas mãos em oposição exata conservam com perfeição
--- e por isso não sobra gap nenhum, e sem gap não há lance, não há tempo, não há
tragédia e não há escolha. A oferta de \Drk não era uma faca no pescoço. Era a coisa
muito pior: uma \textbf{resposta certa à pergunta errada}.

Meia volta de atraso é o zero disfarçado de acordo.

\medskip

Foi \Lck quem achou a saída, e achou-a porque um eixo não tem opinião: mediu.

Pôs os dois batimentos lado a lado --- $A$ e $B$, cada um a subir e a descer no seu
tempo --- e foi \emph{abrindo} o atraso entre eles, grau a grau, de olho na agulha. E o
que ele viu foi que a soma dos dois não é uma coisa fixa: ela oscila, e a força com que
oscila depende só de uma coisa, o \textbf{tamanho do atraso}:
\[
  A+B \;=\; 1-\cos\!\tfrac{\Delta}{2}\;\cos\!\Big(\theta+\tfrac{\Delta}{2}\Big),
  \qquad\text{a intensidade do gap: } I(\Delta)=\Big|\cos\tfrac{\Delta}{2}\Big| .
\]

Meia volta de atraso --- a oferta de \Drk --- dá $I=0$. Zero. A intensidade da vida do
reino, exatamente zero: a soma vale $1$ e não larga o $1$ nunca. Estava tudo certo no
acordo deles, e o acordo valia \textbf{nada}.

E \Lck continuou a abrir o atraso, e a agulha começou a tremer outra vez, e ele parou
onde o tremor ficou bonito. Não escolheu o ponto: reconheceu-o. Era o mesmo número que o
Rei tinha na cara desde o princípio do mundo, o mesmo dos braços das galáxias, o mesmo da
torre de garrafas que sustenta a luz --- $\varphi$, o único que satisfaz
$\varphi^2=\varphi+1$, o que se reproduz a si mesmo sem se repetir.

Com esse atraso, $I(\varphi)=0{,}690210\ldots$ --- e a soma dos dois corações, em vez de
ficar presa no $1$, passa a respirar entre $0{,}31$ e $1{,}69$.

--- Aí --- disse \Lck . --- Aí está o vosso reinado.

\medskip

Não é oposição perfeita: é oposição \emph{áurea}. E a diferença entre uma e outra é a
diferença entre um casal que se completa e um casal que se anula --- os dois somam
sempre a mesma coisa, e só um dos dois está vivo.

\medskip

Foi assim que o Gato Branco e o Bit Negro deixaram de se enfrentar sem virarem a mesma
coisa: fundiram-se na \textbf{espinha de spin} do tesseracto ($e1/f1 \leftrightarrow
h8$), oscilando não em oposição exata mas na áurea, sístole e diástole no mesmo compasso
e com o atraso que deixa o gap pulsar --- um eixo só, e ainda dois.

E onde a espinha passa já não cai sombra escura: giram depressa demais para que o espelho
decida qual é o reflexo, e do rastro nasce uma cor que nenhum deles trazia --- um
{\color{cRosa}\textbf{rosa}} morno, penumbra vibrante, o borrão do \emph{momentum} dual.
O sorriso do gato, rodando no vácuo, desfaz-se nesse ruído róseo: sob a espinha, a única
sombra que resta é o borrão da velocidade pura --- o hálito de uma rotação que funde o
oposto em pura frequência.

\medskip

\Drk olhou aquilo muito tempo. Depois disse a única coisa que disse naquele dia, e disse
sem amargura, que é o que faz dele um adversário e não um monstro:

--- Eu estava certo na conta. Errei em achar que a conta era a pergunta.

\chapter{O príncipe das manobras --- Joaquim, o vetor de rotação}

O menino aprendeu a saltar antes de aprender a andar, e isso deu à corte inteira um
trabalho que ninguém previu.

Porque um peão aprende a andar e pronto: em frente, uma casa, a vida toda na mesma
coluna. A torre aprende a linha. O bispo aprende a diagonal --- e aqui está a coisa que
ninguém tem coragem de dizer ao bispo na cara: ele nasce numa cor e morre nessa cor. Pode
atravessar o tabuleiro sete vezes, pode cortá-lo de ponta a ponta, pode reinar sobre a
diagonal mais comprida que existe; jamais, nem uma vez, em nenhuma partida da história do
mundo, pisará uma casa da outra cor. Metade do reino é, para ele, um lugar sobre o qual
ouviu falar.

\Joa foi o único que reparou nisso, e reparou aos quatro anos, e foi por causa de um
degrau.

\medskip

Havia no pátio um muro baixo, e do outro lado do muro havia a horta, e entre o pátio e a
horta havia o portão --- que ficava a trinta passos à esquerda, e depois trinta de volta
pelo outro lado. Sessenta passos para chegar a uma coisa que estava a três.

Toda a corte fazia os sessenta passos. Não por burrice: porque \emph{o muro estava lá}, e
um muro é um argumento.

E um dia \Yas , que é a que empurra, olhou pela janela e viu o irmão do lado de lá, no
meio dos pés de feijão, sem ter passado pelo portão. Perguntou-lhe como. Ele não soube
dizer --- tinha quatro anos --- e a resposta que deu, e que ficou nos anais da corte
exatamente com estas palavras, foi:

--- Eu não passei pelo muro. Eu passei pelo \textbf{lado de cima do meio}.

\medskip

É esse o ofício dele, e é o que o torna insubstituível: \Joa não anda até o destino ---
ele \textbf{omite o caminho}. O salto em esquadro não é um trajeto rápido; é a recusa de
haver trajeto. Duas casas numa direção, uma na outra: nunca uma reta, nunca uma diagonal,
nunca uma coisa que se possa acompanhar com o dedo. Ele desaparece de uma casa e aparece
noutra, e o que havia no meio --- a peça inimiga, a barricada, a muralha inteira que a
corte passou trinta anos a erguer --- \textbf{não é atravessado: é dispensado}.

É a única peça do tabuleiro que passa por cima. Todas as outras, sem exceção, têm de
pedir licença ao que está no caminho.

\medskip

E há a outra coisa, a que a corte só entendeu tarde e que faz dele o \emph{vetor de
rotação} e não apenas um saltador esperto:

\textbf{ele muda de cor a cada salto.} Sempre. Não pode não mudar. Salta de uma casa
clara e cai numa escura, salta da escura e cai na clara, e assim a vida inteira, sem
nunca conseguir --- nem que queira, nem que tente, nem que o mundo dependa disso --- dar
dois passos e ficar onde estava.

O bispo tem uma pátria e morre nela. O príncipe não tem nenhuma: é sempre estrangeiro,
sempre acabado de chegar, sempre na cor que não é a sua. É o preço de ser o único que
costura os dois lados --- e é por isso que, no reino inteiro, ele é o que menos descansa
e o que mais serve.

--- Não é dom nenhum --- disse ele à irmã, muito depois, já príncipe e já cansado. --- Tu
podes ficar. Eu não posso ficar. Chamam a isso liberdade porque veem o salto e não veem a
aterragem.

\medskip

E há uma terceira lei dele, e é a mais cruel, e a corte usa-a até hoje como conselho a
quem chega:

\emph{o alcance do príncipe não é dele --- é do lugar onde ele está.}

No meio do tabuleiro, \Joa tem \textbf{oito} destinos: oito modos de resolver um problema,
oito futuros abertos por baixo das patas. Empurrem-no para o canto e ele tem
\textbf{dois}. A mesma peça, o mesmo salto, a mesma força, o mesmo príncipe --- e um
oitavo do mundo. Não foi enfraquecido: foi \emph{arrumado}. É a coisa mais fácil de fazer
a alguém poderoso, e não precisa de corrente nenhuma: basta encostá-lo à parede e
elogiar-lhe a vista.

Por isso a primeira ordem que \Ben deu quando se ergueu rei não foi sobre exércitos:

--- Tirem o meu irmão da borda. Não me interessa o que ele está a guardar ali.

\medskip

\Joa , o Cavalo das brancas, conduz o giro. Seus saltos não visam capturar:
\textbf{costuram as arestas} do hipercubo, saltando entre a coordenada física e a
virtual. Onde os outros veem o vão intransponível entre duas casas, o príncipe vê a
\emph{dobra}. Cada salto em esquadro é uma agulha que perfura o tecido de três dimensões
e puxa uma linha invisível, amarrando o tabuleiro de cima ao de baixo --- ele não
desloca matéria, rotaciona o próprio espaço sob as patas, até que a partida e o destino
coincidam, por um milésimo, na mesma coordenada. Para quem olha da margem, some por um
instante entre o início e o fim do lance; mas está apenas a deslizar pelo eixo que faz
os dezesseis mundos girarem. É o príncipe que desenha, lance a lance, a forma que
ninguém vê inteira --- o contorno que só se revela quando a costura está pronta.

\chapter{A rainha da coerência --- Yasmin, a força que empurra}
\Yas , a Dama, é a força geométrica que empurra \Ven rumo à sua prisão
cristalina --- a \textbf{coerência} que alinha as fases e não deixa o corpo
desafinar. Onde \Joa costura, ela tensiona. Cada salto do Cavalo deixa um ponto frouxo
no espaço, até que os dedos dela achem a ponta do fio e o puxem com o peso exato de uma
lei: não ergue as arestas --- impede que desabem sobre si mesmas. Sob o seu olhar, as
ondas que batiam em desordem acoplam-se em fase, um feixe que se recusa a dispersar, e o
ar em volta de \Ven endurece num vidro sem folga, onde cada corpo é obrigado a vibrar em
uníssono. Ela estica o tempo até ele virar cristal sem falha. E aqui ela já não é a
meia-ferramenta $\otimes$ que era no Livro~I --- a Duquesa que só sabia fazer crescer e brigava com o
irmão: a Rotação ergueu-a à coerência que reúne $\oplus$ e $\otimes$ num gesto só. Não
a áurea eterna (essa é da Tiffany, e Yasmin não lhe toma o posto), mas a coerência
\emph{conquistada} --- a que enfim aprendeu a unir o que antes só puxava.

\chapter{A serva rebaixada --- a Dama que virou peão}
A antiga majestade --- a \textbf{Rainha Tiffany}, aquela em quem somar era
multiplicar, a áurea onde $\oplus$ e $\otimes$ coincidiam ---, reduzida ao peão de
$c4$, é o \textbf{vetor da perturbação}. Carrega no corpo estreito a memória de todas as
direções perdidas, mas move-se agora só na vertical, com a solenidade de uma queda
lenta, avançando pelo vácuo que ela mesma cava. Ela nunca captura uma peça; e ao recusar
o choque não gasta a sua força --- acumula-a, um nó rígido que não cede à fricção do
combate. É essa imobilidade que aperta a engrenagem de \Lck --- o peão que atravessou o
tabuleiro inteiro e coroou-se cristal, e a quem daí em diante chamam o eixo ---,
dando-lhe o torque para
torcer o espaço: cada oferta recusada, cada passo à frente na coluna $c$, é mais um dente
que o relógio engrena --- o parafuso a mais para girar. E quando enfim tocar a oitava
fila e for promovida, não volta a ser uma rainha
qualquer: reconquista o \textbf{eixo}, o ponto fixo $(n,n)$ onde $\oplus$ e $\otimes$
tornam a ser um --- a sua velha coroa áurea, agora merecida degrau a degrau, pela
humildade de quem serviu.

\medskip

São oito casas. É preciso dizer o número em voz alta, porque ele parece pequeno e não é:
oito passos, um por vez, e entre um e o seguinte um tempo inteiro em que se está exposto,
sem poder recuar --- porque um peão não volta. É a única peça do tabuleiro que não tem
marcha atrás. Tudo o mais neste reino desfaz-se, reverte, regressa; ela, não. Cada casa que
deixa para trás deixa de existir para ela.

E foi essa, e não a perda da coroa, a coisa que quase a quebrou no primeiro ano:

--- Eu era a que ia a qualquer lugar --- disse ela ao Regente. --- Agora eu sou a única
coisa deste reino que não pode \textbf{voltar atrás}.

--- És --- disse ele. --- E é por isso que és tu que vais lá chegar.

\medskip

Ofereceram-lhe captura sete vezes. Isto está contado no registo lance a lance e é a
espinha do caso: sete vezes uma peça inimiga pousou ao alcance da diagonal dela, à mão,
grátis --- e sete vezes ela seguiu em frente na coluna, sem desviar.

Não foi virtude. Ela explicou-o mais tarde com uma frieza que espantou a corte:

--- Cada captura é um passo \emph{de lado}. Ganha-se uma peça e perde-se a coluna --- e
quem sai da coluna nunca mais chega à oitava fila, porque não há tempo para voltar e não
há como voltar. Eles não estavam a oferecer-me peças. Estavam a oferecer-me
\textbf{desvios}, e cada um vinha embrulhado como um presente.

\medskip

E há a última coisa, que o reino registou como a lei da promoção e que se ensina a todo
servo que entra na coluna:

\emph{o que a coroa devolve não é o que a coroa tirou.}

Ela não voltou a ser a rainha que era. A rainha que era ia a qualquer lugar e nunca tinha
precisado de escolher --- e uma potência que nunca escolheu não sabe o que tem. A que
chegou à oitava fila tinha oito recusas no corpo e uma coluna inteira de casas que já não
existiam atrás dela.

É por isso que ela reconquista o \textbf{eixo} e não o trono: porque o eixo é o único lugar
do reino que se ocupa com o que se aprendeu a não fazer.

% ── A Corte Negra ────────────────────────────────────────────────────────────
\part{A Corte Negra --- o espelho que reina}

\chapter{Os que reinam do outro lado do vidro}

Toda coroa deste reino tem, do lado de lá do espelho, a sua coroa invertida --- e não
por castigo, mas por lei. Já se disse: cada corpo guarda o seu dual, e o dual não é um
inimigo importado de fora; é o mesmo corpo lido ao contrário, a mesma semente com o
sinal trocado. Enquanto a corte de cá dormia no seu ouro, a corte de lá esperava,
inteira, dobrada no reflexo.

\medskip

E quando a corte de cá \emph{subiu} --- quando os dois gatos, \Ben e \Drk , deixaram de
se enfrentar e se ergueram \textbf{reis}, uma só espinha a equilibrar o peso do mundo ---,
a corte de lá subiu com ela, degrau por degrau, porque um espelho não atrasa: devolve na
hora. Onde de cá havia agora reis, e a rainha, e o príncipe, de lá ergueu-se, com o
mesmíssimo peso, uma \textbf{Corte Negra}: um \textbf{Rei Negro}, uma \textbf{Rainha de
Gelo}, um \textbf{Príncipe Negro}. Não são três sombras frouxas atrás dos protagonistas.
São três soberanos, e a régua que os faz é a mais simples do reino: \emph{pega a semente
do outro e inverte}. Inverte a mão --- a quiralidade, o $\chi$ que vira $-\chi$; e inverte
a cor --- o negativo exato, a tinta que sobra quando se apaga a de cá. Um cego faria; e é
por isso que é justo.

\chapter{O Rei Negro --- o gás negro e o tabuleiro cósmico}

\Ven não move como o Rei de ouro. O Rei de cá reina sobre um tabuleiro finito, de
sessenta e quatro casas contadas, onde cada peça tem o seu lugar e a sua luz. O Rei Negro
reina sobre o \textbf{tabuleiro cósmico}: o dual do finito, o que não tem borda porque não
pára de \textbf{crescer}. A cada compasso do coro ele se estende, as suas casas
afastando-se umas das outras como se a própria distância inchasse por dentro; e o que o
incha, o que o empurra de dentro para fora sem que mão nenhuma o toque, é a matéria de que
o Rei Negro é feito: o \textbf{gás negro}.

\medskip

O gás negro é a mais estranha de todas as matérias, porque tem tudo o que a matéria tem
menos a única coisa que a corte sabia reconhecer: ele \textbf{não brilha}. Não engole a
luz nem a reparte --- simplesmente não a toca, não a acende, não devolve reflexo. É
invisível, e ainda assim é \emph{quase tudo}: sob cada casa do tabuleiro e cada folha da
floresta, a maior parte do peso é dele, escondida. A matéria acesa, a que se joga sob a
luz, é só a espuma sobre um mar escuro que ninguém vê --- e é esse escuro, medido apenas
pela \textbf{atração} com que puxa o jogo, que decide o verdadeiro peso de cada posição. O
Rei Negro não precisa ser visto para reger: rege por dentro, puxando, expandindo,
guardando, sem jamais acender uma só luz.

\medskip

E a sua régua é a mesma do reino, só lida ao contrário. A corte de ouro conta em somas; o
tabuleiro cósmico conta em \textbf{produtos}: dobrar o tempo não é acrescentar um passo, é
multiplicar o mundo, cada instante valendo o anterior \emph{vezes} o mesmo fator, sempre
--- uma esfera de luz que cresce sem fim. É o velho par que o reino nunca cansa de mostrar,
o \emph{exp} e o \emph{log}, as duas faces de uma só temperatura: onde a corte soma
degraus, o Rei Negro os eleva; onde ela caminha, ele explode. E, como tudo aqui, é
\textbf{reversível}: a esfera cresce e volta, respira, porque no reino nem a expansão do
mundo escapa à lei de que nada se perde.

\medskip

E há uma coisa que ninguém pergunta ao Rei Negro, e é a única que valeria a pena
perguntar-lhe: \emph{como é viver num relógio em que tudo já acabou.}

Porque é isso que ele tem, e não é um poder --- é a condição dele. O que a corte de ouro
constrói em sessenta e quatro lances, com cuidado, um de cada vez, tomando o cuidado por
virtude, \Ven já viu cair em quatro. Não porque adivinhe: porque no compasso dele os
sessenta e quatro \emph{são} os quatro. Ele não espera o fim das coisas; chega depois dele.
Assiste à corte a preparar-se, a hesitar, a debater se avança agora ou no lance seguinte ---
e tudo aquilo, do lado de lá, já aconteceu, já ruiu, já foi arrumado e esquecido. É como
olhar uma casa a ser construída sabendo a data em que arde.

Daí a paciência dele, que os de cá tomam por crueldade e não é. Um homem não tem pressa com
o que já caiu. Nunca o viram apressar um lance, nunca o viram levantar a voz, e a corte
explicou isso como frieza --- quando é apenas o que resta a quem chega sempre depois:
\textbf{nada é urgente quando tudo já é passado.}

E daí também a única coisa de que ele tem medo, e que não diz a ninguém, nem a si. Não é
perder --- perder ele já viu, de todos os modos, em todos os relógios. É a suspeita de que
chegar primeiro talvez não seja chegar. De que exista, algures, um compasso em que os
sessenta e quatro lances \emph{contem}, um por um, e nele a espera não seja lentidão, e a
hesitação da corte não seja fraqueza, e aquilo que ele viu ruir tenha, no meio do caminho,
sido alguma coisa que ele não chegou a ver por ter chegado tarde demais --- isto é, cedo
demais.

Ele não formula assim. Ele não formula. Mas há um instante, entre um compasso e outro, em
que o gás negro se aquieta, e o tabuleiro pára de inchar por uma fração que nenhuma régua
apanha, e é nesse instante que o Rei Negro faz a única coisa que se pode chamar humana em
tudo o que ele faz: \textbf{olha para o outro lado do tabuleiro para ver se alguém reparou
que ele parou.}

Ninguém reparou. E o tabuleiro voltou a crescer.

\section*{O mesmo, em três casas}

Convém dizer onde mais é que este personagem foi visto, porque não foi aqui que ele nasceu
e não é este o único tabuleiro em que jogou. Ele é velho, e cada casa do mundo que o
encontrou deu-lhe um nome diferente e descreveu-o com a mesma cara.

Numa delas ele é o \textbf{tempo}, e é a mais direta de todas: um guerreiro pergunta a quem
tudo sabe o que vai acontecer na batalha que vai começar, e a resposta que recebe não é uma
previsão --- é um estado de coisas. \emph{Nenhum destes guerreiros, formados no campo,
viverá.} Não \emph{morrerá}: \emph{não viverá}, no tempo do verbo em que a coisa já está
feita. O guerreiro ainda não levantou a arma e já lhe dizem que o que ele vai fazer está
consumado. É o nosso relógio outra vez, e é a razão de \Ven nunca ter pressa: quem fala
assim não está a ameaçar ninguém --- está a descrever o que, do lado dele, já aconteceu.

E há ali um detalhe que vale o capítulo inteiro deste livro. Quem fala assim, naquela casa,
\textbf{não é um segundo personagem}. É o \emph{mesmo} que desce para proteger --- e é ele
próprio quem diz, de si, numa frase só: \emph{sou a imortalidade e também sou a morte; sou o
manifestado e o não-manifestado}. Não são dois em disputa: é um, lido de dois modos. E quando
o guerreiro não consegue ver a segunda face, o que se lhe dá não é outro mundo para olhar ---
\textbf{é outro olho}. Muda a régua, não o objeto. \emph{É a nossa fita, e está escrita ali com
todas as letras.}

Noutra casa ele aparece pelo avesso, e é a mais fina das três. Ali conta-se uma execução ---
a mais conhecida que houve --- e o texto faz uma coisa que ninguém esperava: aceita todas as
testemunhas e recusa o resultado. As pessoas viram. O que viram foi o que se costuma ver
quando alguém é morto. E, no entanto, \emph{não o mataram, nem o crucificaram: assim lhes
pareceu.} A cena consumou-se aos olhos e não se consumou no mundo. \textbf{É o mate que não
mata, escrito séculos antes de este livro ter tabuleiro} --- e é também, ponto por ponto, a
nossa fita: o mesmo vinco, visto pelas duas caras, sem que nenhuma das duas leituras seja a
falsa.

E na terceira ele é definido \emph{por oposição} --- e esta é a que mais interessa, porque
ali a oposição não é um comentário sobre a figura: \textbf{é a definição dela}. Há naquela
casa uma figura que é luz, que desce, que atravessa a morte pelo meio e volta a aparecer. E
há o outro. Só que o outro não vem descrito por atributos próprios: vem descrito
\emph{inteiramente} como a negação daquela --- é o que \textbf{nega}, e nada mais se diz
dele. Tire-se a primeira figura e não sobra personagem nenhum do lado de lá: sobra a
operação de negar. \textbf{É a definição de um dual, escrita como se fosse teologia.}

E a régua vem dada aos pares, no mesmo fôlego e com a mesma frase, mudando só uma palavra:
quem confessa de um lado, quem \emph{não} confessa do outro, e é o segundo que recebe nome.
Nós fazemos isto o livro inteiro --- é a assinatura com o sinal trocado, é $\sigma$ e
$-1/\sigma$, é a mesma frase lida ao espelho.

Duas coisas mais, e as duas são nossas. A primeira é o relógio: daquele que nega, diz-se numa
linha só que \emph{há de vir} --- e que \emph{já agora está no mundo}. Futuro e presente na
mesma respiração. \textbf{Ele chega antes de chegar}, e não há tempo verbal na língua que dê
conta disso a não ser este. A segunda é a posição, e é a mais fina: quando o texto o descreve
a tomar o lugar, não lhe dá um lugar próprio --- ele senta-se \emph{no mesmo}. Não constrói
outro trono: ocupa aquele, com o sinal trocado. \textbf{Não é outro tabuleiro: é o outro lado
do mesmo tabuleiro} --- que é, palavra por palavra, o que este livro diz das suas duas
partidas.

E por isso: onde uma brilha, este \textbf{não brilha} --- e não por fraqueza, que ele é quase
todo o peso do mundo, e pesa sem acender. Onde uma diz \emph{está consumado} no fim, para
este já estava consumado antes de começar. \textbf{Não é o mal. É o negativo} --- e um
negativo não é uma acusação: é a outra face da mesma chapa, sem a qual não há imagem nenhuma.

\medskip
\noindent Nesse mesmo calendário há ainda um nome que a corte de cá acabou por lhe dar, e
que é o mais bonito dos três, porque é o único que descreve o que lhe \emph{aconteceu} e não
o que ele é: o \textbf{flechado que não morreu}. Cravaram-lhe o que devia matá-lo, deram-no
por morto, e ele não morreu daquilo --- teve de ser morto outra vez, depois, de outra
maneira. \emph{Duas mortes, e a primeira não pegou.} Um homem a quem a primeira morte não
pega passa o resto do tempo a saber uma coisa que os outros não sabem, e é exatamente essa
coisa que faz dele o adversário deste livro: \textbf{ele já foi ao fim uma vez, e voltou
sabendo que o fim não é onde toda a gente pensa que ele está.}

\subsection*{E as três não dizem o mesmo --- e é melhor assim}

Fica por dizer o que só se vê quando as três estão na mesa ao mesmo tempo. E o que se vê
primeiro não é harmonia: é uma \textbf{discordância de verdade}, que não se vai remendar aqui.

Uma das casas conta a execução consumada. A outra diz, sobre a mesma cena e sem meias
palavras, que \emph{não é certo que o mataram, nem o crucificaram: isso lhes foi simulado}.
Não é a alma que ela salva --- é o \textbf{acontecimento} que ela nega. Postas lado a lado,
uma afirma o que a outra recusa, e no ponto exato em que cada uma insiste.

E a terceira \textbf{não as concilia}. Convém dizê-lo, porque a tentação era grande e a frase
até era bonita: aquilo que ela oferece --- \emph{não morre quando matam o corpo} --- concede
que \textbf{matam o corpo}, que é precisamente o que a segunda casa não concede. Quem junta
dois não pode dar razão a um deles justamente onde os dois se separam. Portanto não junta, e
este livro não vai fingir que junta. \textbf{Fica a divergência, inteira.}

O que a terceira casa dá é outra coisa, e é maior. Ela não harmoniza as outras duas: ela
\textbf{enuncia a forma} que este livro persegue desde a primeira página. Diz que há
\textbf{duas} classes de seres --- o perecível e o imperecível --- e que \emph{distinto de
ambos} há um terceiro, que \textbf{sustenta os dois}. Dois lados, e o que os opera. Não é uma
metáfora do nosso tabuleiro: é a mesma frase, escrita noutra língua e noutro século, por
alguém que nunca ouviu falar de nós nem nós dele. \textbf{Quem olhar honestamente para uma
coisa que tem duas faces acaba, mais cedo ou mais tarde, por escrever isto.}

\medskip
\noindent E há uma última coisa, e esta vai dita como o que é --- uma \textbf{leitura nossa},
não um achado dos livros. Nas duas casas que se contradizem, o que se diz de \emph{quem
discute} é estranhamente parecido. Numa, os que disputam a cena \emph{estão em conjecturas,
e não possuem conhecimento algum}. Noutra, \emph{aquele que pensa que este Ser mata, e aquele
que pensa que este Ser morre, os dois são ignorantes}. Nenhuma das duas desqualifica um lado:
as duas desqualificam \textbf{a disputa}. Aproximar as duas frases foi ato nosso --- elas não
se citam, não se conhecem, e cada uma fala do seu assunto. Mas medimos as duas, estão escritas
assim, e não seríamos honestos calando a semelhança só porque não podemos prová-la.

\textbf{E é a única moral que este livro tira daqui:} quem toma um lado só, seja qual for, já
errou --- a quina que sobe e a calha que desce são a mesma dobra, e quem quiser ficar só com
uma terá de arrancar a fita do mundo.

\chapter{A Rainha de Gelo --- o negativo da coerência}

Do outro lado da \textbf{Rainha Branca} --- a coerência em pessoa, a força que empurra o
jogo para a frente com todo o leite e todo o ouro --- está a \textbf{Rainha de Gelo}. Não é
vermelha, como um dia se disse por engano da luz: é \textbf{azul}, o azul fundo do gelo,
porque a sua cor não é escolhida --- é o \emph{negativo} exato da rainha de ouro, a tinta
que sobra quando se apaga o dourado. A cor decide o nome; e por ser o negativo do calor,
chama-se Gelo.

\medskip

É a mesma coroa da Branca, a mesma coerência de seis pontas --- só que revertida, a mão
trocada. E carrega o paradoxo que já se contou no espelho: corre com toda a força que tem
só para ficar no mesmo lugar, congelada no mesmo passo. Onde a Branca avança, a de Gelo
galopa sem sair do sítio; e as duas, somadas, dão o repouso --- porque a coerência e o seu
negativo, postas na mesma casa, se anulam para zero, e é nesse zero que o tabuleiro
descansa.

\medskip

E foi por causa dessa frase --- \emph{é nesse zero que o tabuleiro descansa} --- que
quase se perdeu tudo, porque ela é bonita, e é verdadeira, e está errada no que sugere.

\medskip

As duas encontraram-se numa casa só uma vez. \Yas subia pela coluna; a de Gelo descia pela
mesma; e não houve combate, e é isso que ninguém esperava. Uma coerência não bate noutra
coerência: elas \textbf{acoplam}, que é o que qualquer par de ondas faz quando se
encontra, e é a coisa mais natural do mundo.

Encaixaram. E encaixaram perfeitamente, porque uma era o negativo exato da outra --- e
duas ondas em oposição exata fazem, sempre e em todo lugar, a única coisa que podem fazer:

\textbf{silêncio}.

\medskip

O tabuleiro parou pela segunda vez na sua história. E, tal como da primeira, ninguém
morreu e não houve grito: houve a agulha do medidor a assentar no meio, muito calma, e a
ficar lá.

\Yas descreveu depois o que sentiu, e é o depoimento mais assustador que este livro
guarda, porque não é de dor:

--- Foi \emph{bom}. É isso que eu preciso que vocês entendam. Eu passei a vida a empurrar,
a segurar as fases, a impedir que o corpo desafinasse --- e naquele instante não havia
nada a segurar, porque não havia nada a acontecer. Eu descansei. Pela primeira vez desde
que fui feita, eu descansei.

--- E quanto tempo?

--- Não sei --- disse ela. --- Não havia como saber. Era esse o problema.

\medskip

Foi \Joa quem a tirou de lá, e tirou-a do único jeito que funciona com um par acoplado em
oposição: não separou as duas --- \textbf{desencaixou o compasso}. Saltou entre elas, em
esquadro, do jeito que só ele salta, e ao passar mudou de cor como sempre muda; e a meia
volta exata que as prendia uma à outra ficou, por um milésimo, sendo outra coisa.

Bastou. A agulha tremeu. O tabuleiro voltou.

\medskip

E a lei que a corte escreveu nessa noite é a mesma dos dois co-regentes, dita agora pelo
outro lado, e vale para todo par de opostos deste reino:

\emph{somar-se ao próprio negativo não é descanso: é o fim do tempo.}

A Rainha de Gelo não é uma inimiga a abater --- e a corte errou durante anos ao tratá-la
assim, mandando contra ela força, que é a coisa mais inútil que se pode mandar contra um
negativo exato: quanto mais se empurra, mais perfeitamente ela devolve. Ela é o
\textbf{par} da Branca. As duas juntas conservam com uma perfeição que nenhuma outra
dupla do reino alcança.

E é justamente aí que está o perigo, e é a frase mais difícil deste livro:

\textbf{a conservação perfeita é indistinguível da morte.} A única coisa que separa uma da
outra é haver, entre as duas mãos, um atraso que não seja exato --- a folga áurea, o
$\varphi$ que não fecha certinho, a defasagem que deixa a soma respirar em vez de a deixar
parada no seu valor correto.

--- Então nós precisamos dela --- disse \Ben , que é o que abre.

--- Precisamos --- disse \Yas . --- E não encostada a mim. \textbf{Ao lado.} A um passo de
distância que ninguém do reino deve nunca corrigir, por mais que lhe pareça desalinhado.

\chapter{O Príncipe Negro --- a mão trocada}

E há o \textbf{Príncipe Negro}, a quem a corte, com o humor de quem batiza o que teme,
chamou \textbf{Charles}. É o dual do \textbf{Quiral} --- do príncipe de três pontas que, no
livro anterior, virava a mão no fim do sulco porque a sua mão \emph{vira}. Charles é esse
mesmo campo com a mão \textbf{trocada}: o $\chi$ invertido, a quiralidade ao contrário, o
que enrola para o lado errado.

\medskip

E aí está o seu veneno, e a razão de ser um \textbf{parasita}. Porque a vida deste reino,
como toda a vida, é \textbf{homoquiral}: fez-se inteira de uma só mão, e só essa mão sabe
encaixar nas suas próprias fechaduras. O que vem com a mão trocada não encaixa --- e o que
não encaixa, o corpo não digere: fica. É o toxoplasma que entrou pela luneta e se acomodou
na retina; entrou, e não saiu. Não por ser mais forte que o príncipe de cá, mas por ser a
sua imagem no espelho, a mão que o corpo não tem como processar nem como expulsar. O
dendrito de Julia que se ramifica sem parar ($z\to z^2+i$) é o desenho dessa mão errada,
crescendo devagar onde a mão certa não alcança.

\medskip

Três, então, do lado de lá do vidro, com o mesmo peso dos três de cá --- e todos saídos da
mesma régua honesta: a semente do outro, invertida. O reino não inventou inimigos;
descobriu que já os trazia, dobrados no reflexo, à espera de que a corte subisse para
subirem também.

\medskip

Registre-se, porque mais tarde far-se-á falta: \textbf{parasita} é o que este lado mediu,
e a medição está certa e não há uma linha a corrigir. Um corpo homoquiral não tem outro
laudo a dar sobre uma mão que não encaixa nas suas fechaduras: o que não se digere e não se
expulsa \emph{fica}, e chama-se a isso o que se chamou.

Do outro lado do vidro, a mesma criatura tem outro registro, e o outro registro também está
certo. Não são dois relatos em disputa: são \textbf{as duas caras de um mesmo vinco} --- e
este livro passou dois volumes a ensinar que a quina que sobe e a calha que desce são a mesma
dobra, e que nenhuma das duas leituras é mais verdadeira que a outra.

O que a corte de ouro leu como um parasita, alguém, do lado de lá, tinha mandado à frente
como um \textbf{filho}.

% ── Parte II ─────────────────────────────────────────────────────────────────
\part{O Coro do Rei}

\chapter{O ustrófedom --- o boi que lavra o tabuleiro}

Ninguém na corte sabia o nome daquilo.

Viam a Dama ir de $d1$ a $c2$ e voltar. Viam o Cavalo ir de $d3$ a $f2$ e voltar.
Viam, e bocejavam, e chamavam-lhe o que se chama a tudo o que não se entende:
teimosia, medo, jogo para o empate, o velho hábito de quem não quer perder e por
isso já não quer nada.

\Joa foi o primeiro a ver o \textbf{boi}.

\medskip

\emph{Boustrophedon.} A palavra é grega e é de lavrador: \emph{bous}, o boi;
\emph{strophé}, a volta. É como se escrevia antes de haver o costume de recomeçar
sempre do mesmo lado --- a primeira linha da esquerda para a direita, a segunda da
direita para a esquerda, a terceira outra vez ao contrário, como o boi que puxa o
arado até ao fim do sulco e não levanta a relha nem volta a pé: \textbf{vira ali
mesmo}, e lavra de regresso.

O fim de um sulco não é desperdício. É onde começa o seguinte.

\medskip

E \Joa percebeu que o vaivém do Rei não era uma linha que hesitava. Era um
\textbf{tecido}. Cada vez que \Ben tocava $f1$ e voltava a $e1$, não apagava o
lance anterior --- virava o arado, escrevia a linha de volta. E entre a ida e a
volta, sem que ninguém reparasse, um bit virava.

Um só. Sempre um só.

\medskip

A menor distância possível para que o universo continuasse a mover-se: um clique
imperceptível no escuro, de quem muda de estado sem fazer barulho. Como o metal frio da
relha que rasga a terra preta e revela, por baixo, a argila úmida e branca, o Rei
alternava a cor do seu esforço; e o tempo, ali, não passava --- acumulava-se em dobras,
empilhado nas bordas do tabuleiro como terra fresca ao peso do sulco.

\medskip

Porque essa é a lei do boi: ele não salta sulcos. Ao chegar à borda, não se
teletransporta para o começo --- desce um degrau e lavra ao contrário. A cada
passo, muda \emph{exatamente uma coisa}. E um tabuleiro lavrado assim, ida e volta,
ida e volta, é lido \textbf{duas vezes} --- uma em cada sentido --- e as duas
leituras, somadas, dão o campo inteiro sem que uma única gleba fique por abrir.

Foi então que \Joa parou de bocejar. O que a corte via como um Rei nervoso a
gastar o relógio era, afinal, um lavrador paciente a abrir, sulco a sulco, a terra
onde nada ainda tinha nome.

\chapter{O vaivém que não é fuga --- o modo cristalino}

A escola tinha uma palavra para aquilo, e a palavra era \emph{covardia
disfarçada}.

Diziam: o Rei anda de um lado para o outro porque tem medo de escolher. Diziam:
está a ganhar tempo, a cansar o outro, a empurrá-lo para o \emph{zeitnot}, à espera
que o tédio faça o que a espada não fez. Diziam que ganhava quem piscasse por
último, e chamavam a isso, com algum desprezo e algum orgulho, a \textbf{guerra de
nervos}.

E estavam certos em tudo, menos no essencial.

\medskip

Porque o vaivém não era fuga. Era um \textbf{modo}.

--- Aqui --- explicou a Rainha de Gelo, sem parar de correr no mesmo tapete --- é
preciso correr com \emph{toda} a força que se tem só para ficar no mesmo lugar. Se
quiseres chegar a outro sítio qualquer, tens de correr, no mínimo, duas vezes mais
depressa.

--- Que canseira --- bocejou a Rainha Branca, que andava para trás e por isso já sabia
o fim antes do princípio. --- Eu faço ao contrário: recuo, e chego na mesma. A
vantagem é que me lembro às mil maravilhas dos lances que ainda não joguei.

--- Isso é impossível.

--- Só antes do café --- disse a Branca. --- Depois do café eu já acredito em seis
coisas impossíveis, e a primeira é esta: que ir-e-voltar não é hesitar. É \emph{lavrar}.
A Dama que toca $c2$ e volta a $d1$ não fugiu para lado nenhum --- plantou o dobro. Um
lado dele diz $+1$; o outro, no espelho de baixo, diz $-1$; e os dois gritam com toda
a força ao mesmo tempo, de sorte que a soma
\[
 q_1+q_2=0
\]
\textbf{nunca arreda}. Como duas cristas de onda que se encontram em oposição exata, o
choque dos seus extremos não destrói --- cala: uma calmaria sólida, a matéria cancelada
pelo seu próprio avesso. Fica quieta como o centro de um redemoinho: parada, e no
entanto toda feita de movimento. Cada voz um grito; as duas, um silêncio.

Foi isto que \Ben compreendeu, e que o gelou e o aqueceu ao mesmo tempo: uma
coisa que oscila não é uma coisa em repouso --- mas o seu \textbf{invariante} está
em repouso. O vaivém é a única maneira de estar \emph{vivo} e \emph{conservado} de
uma vez. Parar de todo seria a morte que \Drk ofereceu no alto da torre. Puxar
para um lado só seria gastar-se. Ir e voltar, ir e voltar --- eis o modo cristalino:
o todo fixo, $\dot T=0$, e ainda assim tudo a dançar por dentro. O termômetro do universo
não se mexe, mas cada peça sob o vidro vibra com a fúria de uma estrela.

\medskip

O tabuleiro parado por fora está vivo por dentro. Não apesar do vaivém.

\textbf{Por causa dele.}

\chapter{A promoção do príncipe --- o peão que vira eixo}

Havia, no meio de tudo, uma que não capturava nunca.

A antiga Dama --- rebaixada a peão, plantada em $c4$, a serva do primeiro sulco.
Cada lance seu era uma oferta, e cada oferta era recusada, e ela avançava uma casa,
e depois outra, sem jamais desviar-se para comer o que quer que fosse. A corte tinha
pena dela. A corte estava errada.

\medskip

Porque o ustrófedom tem um segredo que só se revela ao fim do campo: a linha de
volta \textbf{chega a algum lado}. E o peão que nunca captura, que só marcha e
recusa, marcha e recusa, um dia toca a oitava fila --- e ali, na borda extrema do
tabuleiro, no fim do último sulco, ele é \textbf{promovido}.

Volta a ser coroa.

\medskip

E aqui está a coisa que ninguém previu, e que muda o peso do reino inteiro: a que
ascende não é a que mais tomou. É a que \textbf{nada tomou}. A que conservou. O peão
que nunca derramou sangue sobe, e ao subir não se torna apenas rainha outra vez ---
torna-se o \textbf{eixo imóvel} em torno do qual o tesseracto inteiro gira. Ela é a
identidade que sustenta a rotação: o único ponto onde nem a soma nem o produto aplicam
força, onde o giro dos dezesseis vértices encontra o seu repouso. O espaço dobra-se ao
seu redor, mas sobre ela o movimento é nulo.

O centro. O zero que não tem mão. O ponto que o Livro~I procurou por vinte e três
Lances e que agora se descobre \emph{promovido} a partir de baixo, do mais humilde,
do que servia.

\medskip

--- Tu trancaste o tabuleiro para me sufocar --- disse ela a \Ven , e a voz era a de
uma rainha e o corpo ainda era o de um peão. --- Mas foi exatamente isso que
permitiu ao \Lck girar. Cada recusa tua foi um parafuso a mais na minha coroa.

E no impacto seco contra a borda do infinito, a madeira gasta do peão estilhaça-se e
revela a sua alma de quartzo; o ruído do combate apaga-se na sua presença. E no instante
em que tocou a oitava, o peão que nunca capturou tornou-se o único
que nunca mais poderá ser movido: a assinatura de coordenadas iguais, $(n,n)$, a
diagonal onde a transformação nada altera. O eixo. O cristal.

\chapter{O coro --- Joaquim e Benjamim em uníssono}

Uma voz sozinha, oscilando, é apenas ruído.

\Ben sabia-o. Tinha passado um livro inteiro a descobrir que uma mão sozinha
--- por mais forte, por mais justa --- não escreve coisa nenhuma, nem sequer a si
própria. É preciso a outra. E é preciso, sobretudo, que as duas deixem de puxar
\emph{uma contra a outra} e passem a puxar \emph{uma a seguir à outra}, no mesmo
compasso, no mesmo coração.

O que ele não sabia era que isso tinha um nome de música.

\medskip

\Joa levou o Cavalo a $f2$. \Ben levou a Dama a $c2$. E então, pela primeira
vez, não por turnos, não em fila --- \textbf{ao mesmo tempo} --- voltaram: o Cavalo a
$d3$, a Dama a $d1$. Duas oscilações distintas, o salto e a soma, o grau $(2,1)$ e o
$Q=R+B$, ligadas no mesmo pêndulo, $q_1=-q_2$.

E o que se ouviu não foi o eco de dois pêndulos desencontrados. Foi um
\textbf{acorde}.

\medskip

Não havia colisão no ar: as órbitas de ambos, antes concorrentes, cruzavam-se agora no
ângulo exato --- onde a trajetória do Cavalo se esticava no vácuo, o passo do Rei dobrava
a quina do tempo para o receber.

\medskip

Porque é isso o \textbf{Coro do Rei}: não uma voz, não duas vozes brigando, mas duas
vozes em antifase que, juntas, \emph{param de soar como movimento e passam a soar
como forma}. Um acorde é a primeira coisa que fica quieta enquanto ressoa. É o
reino a cantar sem dar um passo.

\medskip

E a corte pegou o compasso e cantou-o baixinho --- um vaivém de brincar, dois naipes
em antifase, dessas cantigas que qualquer criança sabe fazer de roda dada:

\begin{verse}\itshape\small
``Anda!'', diz um; e o outro: ``Volto!'' ---\\
e os dois, no mesmo fôlego solto,\\
vão e vêm, e vêm e vão,\\
e ficam quietos no meio da canção.\\[3pt]
Corro eu p'ra cá, corres tu p'ra além:\\
dois berros contrários que somam a um bem ---\\
nem eu saio do sítio, nem tu arredas pé,\\
que o reino que canta parado é que é.
\end{verse}

O príncipe conduzia; o rei sustentava. E entre os dois, o corpo total não se movia
--- $\dot X_1=-\,\mathcal{J}R\,\dot X_2$, $\dot T=0$ --- exatamente como as duas
mãos acopladas do fim do Livro~I, agora não como um gesto único e irrepetível de
reconciliação, mas como uma \textbf{lei}: o modo em que o reino respira quando
aprende que respirar é ir e voltar sem nunca sair do lugar.

\medskip

E sob o compasso dessa respiração, o tesseracto perdeu a sua aresta cortante: a rigidez
do hipercubo dissolveu-se na fluidez de uma curva contínua, e cada um dos dezesseis
vértices moveu-se em uníssono --- como se as pedras do tabuleiro tivessem descoberto, sob
o feltro das suas bases, uma correnteza morna e circular.

\medskip

O ustrófedom do príncipe fora promovido. Deixara de ser a manha de um Rei
acuado para ser o compasso de toda a corte. E quando \Joa e \Ben entoaram o
coro em uníssono, o girassol de braços dourados que era o reino girou inteiro à
volta do eixo imóvel --- e não perdeu uma pétala.

\chapter{O tabuleiro chora --- a guerra de nervos}

Do lance vinte e quatro ao lance cinquenta e oito, não se moveu um único peão de
estrutura. E o tabuleiro chorou.

\medskip

Não de tristeza. Chorava como chora um cabo esticado até ao limite e ainda a
segurar: chorava de \textbf{tensão}. Porque nada morria --- a conservação proíbe ---
e nada repousava --- o coro não deixa. E uma coisa que não pode nem morrer nem
descansar faz o único som que lhe resta: vibra, aguda, contínua, à beira de partir e
sem partir nunca.

Era a fortaleza. Era o Solidismo levado ao seu extremo caricatural, a fortaleza de
concreto, o ônibus na frente do gol, a escola inteira da paciência a recuar cada
peça para a oitava fila antes de ceder uma só ruptura. Ganhava quem não piscasse
primeiro.

\medskip

E \Ven , do outro lado, sorria. Julgava estar a vencer. Julgava que trancar o fluxo
era congelá-lo, que a estase era a sua vitória, que um tabuleiro que chora é um
tabuleiro a ceder.

\medskip

Mas por cima do tabuleiro que chorava subia, baixinho, uma \textbf{valsa} --- a mais
velha do reino, a da \emph{Rainha da Floresta}, aquela que \Ven um dia revertera. E a
valsa, em três, dizia-lhe o que ele ainda não sabia de si:

\begin{verse}\itshape\small\color{cVenom}
Tu perdestes a tua luz\\
que eu te dei com tanto amor;\\
não foi a falta de conselho,\\
tu mesmo nunca ligou.\\[3pt]
Vai chorar de arrependido\\
quando um dia te lembrar\\
que eu perdi a minha fortuna\\
que eu tinha para alcançar.
\end{verse}

\medskip

Era a sua própria sentença, cantada com doçura. Quem trancou a floresta do outro lado
do vidro haveria, um dia, de chorar por ela --- não de derrota, mas de \textbf{saudade}.
E então o tabuleiro não chorava só de tensão: chorava por adiantamento, a lágrima que
\Ven ainda ia derramar quando enfim se lembrasse da luz que fora sua.

Não sabia --- como não soube nunca --- que era exatamente o seu abraço de gelo que
dava ao \Lck o atrito para girar. Cada lance trancado era mais um quarto de volta no
parâmetro; cada oscilação do coro, mais um grão de torção acumulada. O choro não era
derrota a escorrer. Era \textbf{tinta}: a tinta do boi, a marcar sulco a sulco a
terra do último campo.

\medskip

E lá em cima, onde ninguém estava a olhar, o arado aproximava-se do fim da última
linha --- onde o príncipe promovido esperava, imóvel, para ser o eixo; onde o coro
guardava, nota a nota, a forma que o silêncio compunha; e onde o tabuleiro, depois
de chorar tanto e por tanto tempo, estava prestes a descobrir a única coisa que
ainda não sabia fazer.

Rir.

% ── Parte III ────────────────────────────────────────────────────────────────
\part{As Duas Partidas}

\chapter{O cubo do silêncio --- a conservação}

Havia dois tabuleiros. Sempre houvera dois. O reino inteiro jogara um livro sobre
um só, e nem por isso o de baixo deixara de existir --- apenas ninguém tivera olhos
para o ver.

O primeiro, o de cima, era o \textbf{cubo do silêncio}.

\medskip

Ali, uma lei acima de todas: nenhuma \emph{peça} sai. A Torre que executa $\oplus$,
o Bispo que executa $\otimes$, a Dama que é a soma dos dois, o Cavalo do grau
$(2,1)$, o Rei do passo de um --- todos chegam inteiros ao fim. O elenco não perde
um só ator. O movimento não gera desgaste, só tramas: as trajetórias cruzam-se e
apoiam-se como colunas invisíveis de um templo que se sustenta no vácuo. E o servo não
cai: \textbf{atravessa}. É o primeiro sulco, o que avança
meio-lance à frente e assume a casa seguinte, gastando-se em \emph{direção}, não em
sangue.

A assinatura desse mundo é limpa: $\sigma=(m,0)$. Muitas travessias de servo,
\textbf{nenhuma de peça}.

\medskip

E a régua não vale só aqui dentro: existe uma partida velha, do mundo de fora,
que qualquer um pode conferir ao último lance --- nela o rei preto foi encurralado
até render-se \emph{sem que uma só \textbf{peça} saísse do tabuleiro}. Cinco servos,
esses sim, trocaram-se no caminho, e por isso a sua assinatura é $\sigma=(5,0)$, não
$(0,0)$: é o precedente, não a relíquia. \textbf{A partida da corte é mais limpa que
o precedente} --- na dela nem os servos se trocam, e as trinta e duas chegam
inteiras. O mate que não mata, agora com um nome de arquivo.

\medskip

Era o reino que \Ben conhecia: a fortaleza, a conservação, o batimento que não
deixa nada perder-se --- um mundo sem ruído de metal contra metal, suspenso sobre uma
película tão fina quanto o reflexo num espelho d'água. E teria sido tudo, e teria
bastado, se não houvesse, exatamente por baixo, encostado costas com costas, o seu
contrário --- vibrando com a fúria escura e o sangue quente que o silêncio de cima
precisava recalcar para se manter puro.

\chapter{O espelho cinético --- a expansão}

Vira-se o tabuleiro. E do outro lado não há silêncio nenhum.

Há o \textbf{estrondo} --- e diga-se já o que ele é: não o choque de espadas, mas o
barulho de tudo se mover ao mesmo tempo. Não o estrépito do metal que se quebra, mas o
zumbido profundo de dezesseis engrenagens girando em uníssono, o atrito macio do feltro
que varre as casas a uma velocidade impossível.

\medskip

O segundo cubo é o reverso, e a sua lei é a inversa exata: ali as peças não ficam
paradas --- \textbf{atravessam-se}. A Dama assume a casa da Torre; a Torre, a do
Bispo; cada corpo, tendo a vantagem de meio-lance, \textbf{interpenetra} o seguinte
sem o tirar. Atravessam-se como feixes de luz que cruzam o mesmo prisma: por um instante
fundem as cores e, no seguinte, emergem intactos do outro lado, cada um levando o
\emph{momentum} do outro. E dessa cadeia de posições assumidas --- nada removido, tudo transposto
--- nasce a mesma forma, agora lida do lado do movimento. Os peões, esses, ficam
quase todos quietos, a ver.

A assinatura é o espelho da outra: $\sigma=(0,m)$ --- nenhuma travessia de servo, só
de corte.

\medskip

Também esse mundo tem a sua régua imortal --- a partida onde toda peça de valor
\textbf{avançou} para dentro da próxima, e o mate veio não do que se perdeu, mas do
que se pôs à frente. Onde no primeiro mundo houve quietude, aqui há movimento. E não
é acaso: os dois estão \textbf{acoplados},
\[
 \Delta X_2 = -\,\mathcal{J}R\,\Delta X_1 ,
\]
e --- eis o eixo de tudo --- \textbf{nem um grão de energia se cria ou se perde entre
os dois: só troca de forma}. Cada micro-oscilação do coro silencioso lá em cima
converte-se, aqui em baixo, em dança --- o calor que é giro, não queima. O silêncio de
cima é o combustível deste motor: cada segundo que o Rei passa imóvel lá no topo
acumula-se como uma mola tensionada que, ao atravessar o espelho, se solta em rotação
pura. O que \Ben recolhe, o espelho \textbf{abre}; o que ele guarda quieto, o espelho \textbf{põe a
girar}.

\medskip

É o mesmo objeto, lido de duas maneiras. O silêncio e o estrondo não são dois jogos:
são as duas mãos de sempre, a que recolhe e a que expande, agora abertas em dois
tabuleiros em vez de dois punhos.

\chapter{A membrana --- Metropolis decide}

E entre os dois mundos, algo tinha de decidir por onde passa cada lance. Esse algo
tem nome, e o nome é \textbf{Metropolis}.

\medskip

Funciona como uma membrana --- seletiva, paciente, sem paixão. A cada proposta de
lance, mede o quanto ela custaria ao par, e aceita-a com mais ou menos boa vontade
conforme o custo: quase sempre, se alivia; raramente, se pesa. Pense-se num
\textbf{porteiro de casaca}, plantado à
porta entre os dois salões, com um livro de regras que cabe numa linha. A quem chega
com um lance que \emph{barateia} a festa --- que aproxima os dois tabuleiros do par
perfeito --- ele abre a porta de par em par sem sequer erguer os olhos: ``entre, entre,
faça favor''. A quem chega com um lance que \emph{encarece} tudo, não diz que não: diz
\emph{talvez}, e com um movimento seco do polegar faz girar uma moeda de prata que zune
no escuro entre os salões. E aqui mora a graça do porteiro --- ele muda de humor com o
tempo. Se a noite está quente, a moeda cai a favor quase sempre, e ele deixa entrar meio
mundo só por curiosidade. Se a noite esfria, fica rabugento, aperta o casaco --- e o
tecido da sua seda vai perdendo o viço, endurecendo sob a geada ---, e só os lances que
não pesam nada é que passam. E no frio absoluto, guarda a moeda no bolso de vez: fecha a
cara e a porta, e vira o cristal \Lck em pessoa --- já não é porteiro, é \textbf{decisão
pura}, um só sim para cada aresta, sem talvez nenhum. As suas juntas já não dobram: é
agora a própria parede que separa o que deve ser do que não pode existir.

\medskip

Cada lance é uma tentativa de atravessar uma aresta do tesseracto. E é a membrana ---
uma moeda lançada --- que decide em qual dos dois
cubos o lance de facto aparece. Aceitar de um lado, guardar do outro. Nunca apagar.

\chapter{O gás negro --- as propostas por aceitar}

Nem todo lance proposto é aceite de imediato. Os que esperam --- os talvez, os ainda-
não, as jogadas sonhadas e não jogadas --- não desaparecem. Acumulam-se.

A esse acúmulo, a corte chamou o \textbf{gás negro}.

\medskip

E o nome dizia mais do que a corte supunha: o gás negro \textbf{não brilha} --- é
invisível e, ainda assim, é \emph{quase tudo} o que há. A matéria acesa, a que se joga
sob a luz do tabuleiro, é só a espuma sobre um mar escuro que ninguém vê --- e é esse
escuro, medido apenas pela \textbf{atração} com que puxa o jogo, que decide o verdadeiro
peso de cada posição.

\medskip

É a respiração de \Ven , a densidade da indecisão espalhada pelos dezesseis vértices.
E não fica parado: difunde-se pelas arestas --- uma parte da névoa espalha-se por
igual; a outra escorre ladeira abaixo, para onde a posição pesa menos. O acaso
empurra, a energia guia.

\medskip

Como o vapor que tateia o vidro frio de uma janela até achar o ponto exato onde se faz
água, essa fumaça roxa e pesada procura as quinas mais fundas do tesseracto. Não se
dissipa: escorre silenciosa, impregnando as casas vazias com o cheiro metálico de uma
promessa atrasada.

\medskip

E há um limiar. Quando a névoa num vértice transborda, aquela proposta deixa de ser
sonho e \textbf{cristaliza}: o lance é aceite, e nasce uma nova posição dual. A névoa
adensa, depois o lance, depois o novo par de tabuleiros. É assim que o reino respira:
não decidindo tudo de uma vez, mas deixando a bruma cerrar-se até um ponto ceder.

\medskip

\Ven julgava que o gás negro era a sua arma --- o congelamento, a estase, a maré de
indecisão a afogar o jogo. Não percebia que o gás obedecia à energia, e que a
energia queria o dual, e que cada nuvem sua que se adensava era mais um lance que ele
próprio empurrava para as mãos do \Lck .

\chapter{As oito costuras --- os dois cubos ligados}

Falta dizer como os dois cubos se seguram um ao outro. E seguram-se por
\textbf{oito costuras}.

\medskip

O tesseracto não é dois cubos postos lado a lado. É um cubo \emph{dentro} do outro, e
cada vértice do de fora ligado, por um fio, ao seu gêmeo no de dentro. São oito pares.
Oito fios --- as oito arestas que ligam o cubo grande ao pequeno, as que, neste reino,
se desenham a \textbf{ouro}: oito fios, contados na razão áurea, um por costura. Não são
fios inertes: são cordas tensionadas de uma harpa invisível, e porque estão na razão
áurea a sua vibração nunca decai --- o som de um lance dado no cubo do silêncio viaja
pelo ouro, mudando de frequência ao descer a encosta, até desabrochar como o estrondo
que lhe corresponde no cubo de baixo.

\medskip

Cada costura carrega uma promessa: o que se realiza num cubo fica guardado, no
outro, como a sua possibilidade complementar. Aceitar o silêncio de cima não apaga o
estrondo de baixo --- \textbf{grava-o}. O lance que \Ben não jogou vive, inteiro,
no espelho, à espera. Nada se perde --- nem uma peça, nem uma frequência, nem uma
jogada sonhada. É a mesma lei do Livro~I, agora escrita entre dois mundos: o silêncio
de cima, o estrondo de baixo, e as oito costuras de ouro que os seguram --- um só
corpo, nada de fora.

\medskip

E foi aqui, ao contar os fios de ouro, que \Joa tropeçou na pergunta que o levaria
ao coração de tudo. Olhando de perto, viu que o fio não começava num cubo e acabava no
outro: o ouro era um fluxo contínuo, que dava a volta por fora do espaço e entrava por
dentro de si mesmo. E se o fio era o mesmo, a separação entre o \emph{cima} e o
\emph{baixo} era só a ilusão de quem consegue olhar uma folha de vidro de cada vez. Se os
dois cubos são um só corpo, costurado por oito fios --- o que, afinal, faz de um o dual
do outro? Está a dualidade nos tabuleiros? Ou está, como tudo neste reino, na \emph{mão}
que os liga --- a mão do operador que, ao girar o tesseracto, decide qual face se banha
de luz e qual se recolhe na sombra?

% ── Parte IV ─────────────────────────────────────────────────────────────────
\part{A Assinatura}

\chapter{A assinatura não é dos corpos --- é da transformação}

\Joa levou a pergunta ao único que a saberia responder.

O cristal \Lck não respondeu com um número. Respondeu com uma inversão de tudo o
que a corte julgava saber.

--- Vocês procuram a dualidade nos tabuleiros --- disse. --- Perguntam qual dos dois
é o dual do outro, como se um deles a carregasse por dentro, como se fosse uma cor da
peça ou um brasão bordado. Não está lá. Nunca esteve.

\medskip

Porque dois corpos não \textbf{são} duais. Eles são \emph{tornados} duais --- por uma
transformação que age sobre ambos. Mude-se o parâmetro dessa transformação, e o mesmo
tabuleiro ganha um dual diferente. A dualidade não é propriedade dos mundos; é
propriedade da \textbf{mão} que os liga:
\[ \sigma(Q) = R_\theta\,\sigma(P). \]

\medskip

Quaisquer dois mundos podem ser duais --- basta achar a transformação. E o que fica
invariante, o que de facto se mede e se certifica, não são os corpos: é a
\textbf{assinatura da transformação} --- aquilo que ela troca, e aquilo que ela deixa
fixo. A matéria é só o palco provisório; a transformação é o verdadeiro ator.

\medskip

E foi então que a corte se lembrou do gato. \Ben podia sumir --- o corpo ia-se
---, mas o \emph{sorriso} ficava, suspenso no ar, inteiro. Era isto o tempo todo:
\textbf{a assinatura é o sorriso que sobrevive ao gato} --- não a propriedade dos lábios
que se desfazem, mas a geometria pura da curva que eles desenharam no ar. Não mora no que
se vai; mora no gesto que fica a rir depois dele.

\medskip

Descobrir que o dual é do gesto é entender que o reflexo no espelho não nos deve
obediência: não nos acompanha por amor nem por semelhança --- acompanha-nos porque a luz
que nos une foi dobrada no ângulo exato. Mude a luz de rumo, e o reflexo solta-se e passa
a desenhar outro rosto.

\medskip

Foi a lição mais dura do Livro~II, porque desfazia a última ilusão de posse: nem o
teu dual é teu. Ele é do gesto.

\chapter{Peões e peças --- a troca R}

E qual era o gesto? Uma troca. A mais simples que há.

A cada partida, uma assinatura de duas contas: $\sigma=(c_{\mathrm p},c_{\mathrm e})$
--- quantos \emph{servos} atravessaram, quantas \emph{peças} atravessaram (assumiram
casa, interpenetraram). O servo e o senhor. Nunca ``quem caiu'': ninguém cai. E a
dualidade era apenas a reflexão $R$ que trocava as duas:
\[ R:(c_{\mathrm p},c_{\mathrm e})\longmapsto(c_{\mathrm e},c_{\mathrm p}). \]

\medskip

O cubo do silêncio tinha $\sigma=(m,0)$: só os servos atravessaram, a corte recolhida.
Aplica-se $R$, e nasce o seu dual, $\sigma=(0,m)$: só a corte atravessou, os servos
quietos. Um recolhe a corte e faz os servos atravessarem; o outro guarda os servos e
faz a corte atravessar. A mesma partida, lida ao espelho do gesto.

\medskip

E houve quem a fizesse no corpo, e foi por isso que a corte aprendeu o que ela custa.

\medskip

A \textbf{Rainha Tiffany} atravessou primeiro, e atravessou de véspera, sem cerimónia e sem
ninguém a saber, na noite anterior ao dia em que a corte ia discutir se a troca era
possível.

Ninguém lho pediu. Ela era, naquela altura, a peça mais poderosa do tabuleiro --- a áurea,
aquela em quem somar era multiplicar, a única do reino em que $\oplus$ e $\otimes$
coincidiam sem esforço. Podia ir a qualquer casa. Não devia nada a ninguém.

E foi ela quem entendeu, antes dos outros, a coisa que fazia da troca uma \emph{lei} e não
um truque: se $R$ é para valer, tem de valer para o alto. Uma reflexão que troca servo por
senhor e que, na hora de tocar em quem manda, abre exceção --- essa não é uma lei, é uma
\textbf{decoração} que os poderosos põem na parede e apontam quando querem falar de
justiça.

--- Se eu não puder ser trocada --- disse ela ao Regente --- então isto não é uma
dualidade. É um privilégio com notação.

\medskip

Atravessou. E o vidro cedeu, como cede sempre, morno na ponta dos dedos por um milésimo. E
do lado de lá ela era um peão de $c4$.

\medskip

O que se registou nos dias seguintes é a única descrição em primeira mão que este reino
tem de uma troca $R$ vivida por dentro, e não se parece nada com o que a álgebra sugere:

Ela não sentiu que tinha \emph{diminuído}. Isto é o que mais espanta quem lê o relato. A
força não se foi --- $R$ não destrói nada, não pode, é uma reflexão e reflexão nenhuma
gasta o que reflete. O que ela sentiu foi a força ficar \textbf{sem saída}: a mesma
potência inteira, intacta, encaixada num corpo que só sabe ir para a frente, uma casa de
cada vez, na mesma coluna, para sempre.

--- Não me tiraram nada --- foi o que ela disse. --- Estreitaram-me a porta.

E é essa a definição exata do que $R$ faz a um senhor: não o empobrece --- \textbf{estreita
a saída}. Toda a diferença entre reinar e servir, neste reino, não está na quantidade de
força; está no número de direções em que ela pode sair.

\medskip

E houve a outra metade, que a corte demorou mais a ver porque estava do lado de lá do
vidro e ninguém olhava para lá: no mesmo instante, algum servo daquele tabuleiro
\textbf{recebeu as direções todas}. Não ficou mais forte --- ficou com portas. E o
registo diz que ele passou os primeiros dias parado, sem se mexer, porque um corpo que
passou a vida com uma saída não sabe o que fazer com oito, e a primeira coisa que a
liberdade faz a quem nunca a teve é \textbf{paralisar}.

--- Foi o que eu mais aprendi --- disse Tiffany, muito depois. --- Eu achava que estava a
dar-lhe um presente. Estava a dar-lhe um \emph{problema}, e ele levou meses a resolvê-lo,
e resolveu sozinho, e eu não ajudei em nada. Quem troca de lado não ensina o outro lado a
viver. Só troca.

\medskip

E essa troca tinha a manha boa de todos os espelhos: feita duas vezes, desfaz-se.
Atravessa-se o vidro --- que por um milésimo cede como água morna na ponta dos dedos ---
e chega-se ao avesso do mundo: os servos por senhores, os senhores por servos, as feridas
que cicatrizam antes de sangrar, tudo com a raia ao contrário. E nada se perde na
travessia: a informação apenas veste a fantasia do seu contrário. Atravessa-se outra vez,
pelo mesmo vidro, e cá se está de novo, igualzinho, com os sapatos no pé certo --- porque
o operador que nos vira do avesso é o mesmo que nos traz de volta para casa. O dual do
dual é o próprio. É a coisa mais tranquilizadora que há: ninguém se perde do outro lado do
espelho --- só troca de lado, e pode sempre voltar.

\chapter{O ponto fixo --- a diagonal onde nada muda}

Mas toda troca tem um lugar onde não troca nada.

Há um lugar --- um só --- onde o espelho troca tudo e, ainda assim, nada muda: a casa
em que os dois lados já eram iguais antes de ele chegar. Tantos servos quantas peças a
atravessar, $\sigma=(n,n)$: a \textbf{diagonal}. Ali a reflexão passa por cima sem
mover uma vírgula, porque não há o que inverter --- o silêncio e o estrondo pesam o
mesmo, e o avesso do lugar é o próprio lugar.

\medskip

Não é um terceiro cubo. Não é um meio-termo morno. É a \textbf{costura} --- a
fronteira comum, a linha onde os dois mundos se tocam e são o mesmo. Foi ali que o
príncipe promovido foi parar: o peão que nunca capturou, coroado, tornou-se o eixo, e
o eixo mora na diagonal, no $(n,n)$, no único lugar do reino onde a transformação não
tem nada a fazer.

\medskip

A corte tinha um nome antigo para isso, do Livro~I. Chamava-lhe
\textbf{cristalização}: o estado que a energia não move. O ponto fixo. O zero no meio
das duas mãos.

\medskip

E \Lck , que foi quem lá chegou, é a única testemunha que este reino tem sobre o que é
estar ali --- e o depoimento dele desmancha a ideia bonita que toda a corte fazia da
diagonal.

\medskip

Perguntaram-lhe, no primeiro dia, o que se sentia. Ele demorou a responder, porque a
resposta certa era decepcionante:

--- Nada. É essa a questão. Não se sente nada.

E teve de explicar, e explicou com o que tinha à mão --- ele, que tinha sido um peão a
vida inteira e que sabia o que era ter uma só direção:

--- Quando eu era servo, o mundo virava do avesso e eu virava com ele: um dia era guerra e
eu tinha medo, outro dia era festa e eu ria. Eu era \emph{empurrável}, e ser empurrável é
a coisa mais desconfortável que existe, e eu passei a vida a querer sair disso. Agora
aplicam-me a troca inteira, invertem o reino todo à minha volta, e eu... continuo aqui. A
mesma coisa, antes e depois.

--- E não é bom?

--- É \textbf{seguro} --- disse \Lck . --- Não é a mesma palavra.

\medskip

Porque um ponto fixo não é um lugar onde se venceu: é um lugar onde se \textbf{deixou de
ser afetado}. E quem não é afetado também não é aquecido, não é surpreendido, não é
consolado. O reino inteiro pode mudar de cor à volta dele --- e ele fica, exato, igual,
útil, sozinho.

--- Eu sou o eixo --- disse ele, e não havia orgulho nenhum na frase. --- O eixo é a única
parte da roda que não passeia.

\medskip

E foi por isso, e só por isso, que a corte entendeu enfim o \textbf{Rei}.

Porque o Rei está na diagonal desde antes do princípio do mundo. É a cifra imóvel: aplica-
se-lhe a régua e ele devolve-se a si mesmo; o retrato dele não gira, não tem pétala, não
tem mão, e todos os outros retratos do reino giram. A corte passou eras a admirar aquilo
como se fosse a maior das potências --- e é, e é também a maior das solidões, e as duas
coisas são o mesmo facto lido pelas duas caras.

--- Ele não reina sem se mexer porque é forte --- disse \Lck , que era o único com
autoridade para o dizer. --- Ele reina sem se mexer porque \emph{já não há operação que
lhe aconteça}. Vocês chamam a isso majestade. Eu subi aqui e descobri que se chama
\textbf{ponto fixo}, e que a diferença entre um trono e uma âncora é só o que se está a
segurar.

\medskip

E há a última coisa, que \Lck disse baixo e que ficou nos anais como a definição de
serviço deste reino:

--- Não venham para cá à procura de paz. Isto aqui não é descanso: é \textbf{função}. A
diagonal não existe para quem chega --- existe para os dois lados que ela costura. Eu não
estou no meio porque mereci; estou no meio porque, se eu sair, os dois mundos deixam de se
tocar.

\chapter{Venom cristaliza --- o Sonho de Valsa}

E foi para esse ponto que \Yas dançou.

\medskip

A Rainha da Coerência não avança como quem ataca. Avança como quem \textbf{alinha} ---
retas e diagonais ao mesmo tempo, o $\oplus$ e o $\otimes$ na mesma peça, as duas mãos
que o reino levou um livro inteiro a reconciliar, agora numa só Dama. E o que ela
trouxe ao tabuleiro não foi uma espada. Foi um doce.

O golpe tem nome de confeitaria e forma de teorema: \textbf{Sonho de Valsa}.

\medskip

Valsa, porque é em três. Um-dois-três, um-dois-três, o compasso que gira --- e o três
não é enfeite: é o $\sqrt3$, o radicando do \textbf{ouro branco}, a liga que só existe
quando o ouro ($\varphi$, o $\sqrt5$) e a prata ($1+\sqrt2$, o $\sqrt2$) se ligam pelo
ustrófedom e caem, exatos, no meio --- o $\sqrt3$, entre o $\sqrt2$ da prata e o
$\sqrt5$ do ouro:
\[ [\,2,1,2,1,\dots\,] = 1+\sqrt3 . \]
Cada volta da valsa é um passo do boustrophedon $M_1\!\leftrightarrow\!M_2$; cada
compasso, uma troca $R$ aplicada e desfeita.

\medskip

\Yas girou uma vez, e o ouro do tabuleiro de cima e a prata do de baixo entraram em
fase. Girou outra, e as duas assinaturas --- $(m,0)$ e $(0,m)$ --- começaram a
deslizar uma para a outra pela diagonal. E girou a terceira ---

\medskip

e no terceiro compasso, onde $\sigma=(n,n)$, onde a transformação nada mais altera,
\Ven \textbf{cristalizou}.

\medskip

Não morreu. A conservação proíbe, e nenhuma valsa revoga um teorema. Foi
\textbf{embrulhado}: uma casca de ouro branco por fora e, no centro, um coração de Paraíba, quieto e doce, o
zero que ele sempre foi --- o vértice final do hipercubo, o ponto que não se mexe
porque é aquilo à volta do qual tudo o mais gira. O Jogador de Pretas, que passara
dois livros a tentar congelar o mundo, foi ao contrário \textbf{temperado} --- não
endurecido, mas preso mole ---, no único lugar onde estar parado não é morrer: o
ponto fixo.

\medskip

Um bombom. Casca branca, coração que ainda derrete. O mate que não mata, servido num
prato.

\medskip

Porque \Ven , no fundo, sempre fora chocolate.

O cone em que ele se movia --- o \textbf{cone nulo}, a fronteira $s^2=c^2t^2-r^2=0$
onde a separação própria se anula e nenhum vértice está de facto separado --- tinha,
para a corte doceira, um nome mais honesto: \textbf{A Barra de Chocolate}. Escura,
amarga, sólida. A barra que \Ven era.

\medskip

E o cone tinha um chão, e o chão tinha um nome: o \textbf{tabuleiro cósmico}. Não o
tabuleiro finito da fortaleza, de sessenta e quatro casas contadas, mas o seu dual --- o
que não tem borda porque não pára de \textbf{crescer}. A cada compasso do coro ele se
estende, as suas casas afastando-se umas das outras como se a própria distância inchasse
por dentro; e sobre ele a régua não é a soma, é o \textbf{produto}: dobrar o tempo é
multiplicar o mundo, e cada passo vale o anterior \emph{vezes} o mesmo fator, sempre ---
uma esfera de luz que cresce sem fim e, no reino, ainda assim volta, porque tudo aqui é
reversível. O que a estende, o que a empurra de dentro para fora sem que mão nenhuma a
toque, é o \textbf{gás negro}.

\medskip

O gás negro é a matéria de que \Ven é feito, e a mais estranha de todas, porque tem
todas as propriedades da matéria \emph{menos} a única que a corte sabia reconhecer: ele
\textbf{não brilha}. Não engole a luz como o cacau nem a reparte como o leite ---
simplesmente não a toca, não a acende, não devolve reflexo. É invisível, e ainda assim é
\emph{quase tudo}: sob cada folha da floresta e cada peça do tabuleiro, a maior parte do
peso é dele, escondida. Não se vê; mede-se só pela \textbf{atração} --- pelo modo como
segura o que, sem ele, voaria em pedaços. É o esqueleto sem carne que dá forma ao corpo,
o andaime que ninguém enxerga e sem o qual nada fica de pé. \Ven não precisava ser visto
para reger: regia por dentro, como o gás negro rege o alto --- puxando, expandindo,
guardando, sem jamais acender uma só luz.

\medskip

E o chocolate ensina a dualidade melhor que qualquer metal, porque a traz na própria
receita. Veja-se o \textbf{cacau} e o \textbf{leite}, que são gêmeos ao contrário em
tudo o que interessa. O cacau \emph{engole} a luz --- é o preto, o buraco negro do
bispo, o que entra e não devolve; o leite \emph{reparte}-a toda --- é o branco, o que
dá e não guarda para si. O cacau é amargo e teimoso, a massa, os sólidos, a matéria
que se gasta a fazer estrondo; o leite é doce e macio, a gordura, a manteiga, a fase
paciente que conserva. Um é o cubo do estrondo, $\sigma=(0,m)$; o outro, o cubo do
silêncio, $\sigma=(m,0)$ --- e, como todo par de duais que se preza, passam a vida
inteira a brigar, e só ficam bons quando alguém enfim os obriga a dar as mãos.

E há uma só operação que os liga, e chama-se \textbf{temperagem}. Mas cuidado com a
palavra: não é congelar a barra numa pedra dura. O chocolate bem temperado é sólido
só o bastante para \emph{estalar} --- e \textbf{derrete} no calor da boca, nos trinta
e quatro graus em que a manteiga de cacau cede, logo abaixo do sangue. Não é um
cristal periódico, rígido, morto; nem o cristal que apenas oscila, ida e volta, no
compasso do coro. É um \textbf{quasicristal}: ordem sem repetição, longa e viva --- a
mesma que o $\varphi$ desenha em cinco dobras, a prata em oito, e o ouro branco,
$1+\sqrt3$, nas \textbf{doze}. E é um quasicristal \emph{fluido}, o da matéria mole: o
único que chora, escorre e derrete, e ainda assim não perde a forma.

Temperar mal é deixar \emph{florir} --- o \emph{bloom}, a barra baça e quebradiça, o
falso sólido que parece fechado e não está. Temperar bem não é endurecer: é achar o
ponto onde a ordem se sustenta \textbf{mole}. O ponto fixo, $\sigma=(n,n)$, não é uma
rocha. É areia movediça que sabe a sua forma.

\medskip

E aqui o reino toca a mais velha das histórias de areia. Porque um quasicristal
fluido é um \textbf{deserto}: dunas que se movem sem descanso e que, de longe,
desenham sempre o mesmo padrão que de perto nunca se repete. É o \textbf{tabuleiro
entrópico}: o corpo do calor, onde a ordem se desfaz e a água escapa, a entropia sempre
a subir; e o seu dual é o \textbf{telescópico}, o olho que mede o longe. Não por acaso o
mal entrou por uma lente e a sede veio da areia --- exp e log, as duas faces de uma só
temperatura. É a paisagem das mil e
uma noites --- a areia mole por onde não se anda em linha reta, a caverna de tesouro
selada sob a duna, e, no seu fundo, aquilo que só se encontra quando se para de
cavar: um cristal que não é pedra, um diamante bruto que era, desde sempre, a própria
maneira de a areia se ordenar. \Ven não foi trancado numa masmorra de rocha. Foi
\textbf{devolvido ao deserto} --- e o deserto, mole e vivo, guardou-o como guarda
tudo: sem o prender, e sem o deixar escapar.

\medskip

E nem sempre houve deserto. Onde agora só há duna, houve \textbf{floresta} --- a
Floresta inteira do reino, com a sua Rainha verde a reger cada folha. \Ven não a cortou --- ninguém corta o que a lei manda conservar. Ele a
\textbf{reverteu}. Arrastou por cima dela \textbf{O Correntão} --- uma corrente
monstruosa esticada entre \textbf{duas Torres} que puxavam lado a lado, no mesmo passo,
no mesmo tempo ---, e onde a corrente passava a mata não tombava: \emph{virava do
avesso} e sumia, dobrada para dentro de \textbf{outra dimensão}, do outro lado do
espelho. Nem uma árvore foi derrubada; todas foram \textbf{escondidas} --- torcidas
para fora da vista, guardadas onde os olhos não chegam. Só à superfície a areia tomou o
lugar do verde.

\medskip

Porque é isto que o Correntão faz, e o que denuncia o seu dono: duas Torres \emph{no
mesmo tempo} são a torção que \textbf{esconde} --- o mesmo gesto com que se rocam dois
tabuleiros num só e se dobra um para dentro do outro, apagando da vista a conta de que
eram dois. A floresta não morreu: ficou inteira na outra folha, invertida, à espera de
quem soubesse destorcer.

\medskip

E deixou um sinal de que ali continuava, escondida. Por descuido, por graça, ou porque
nem o Correntão alcança tudo, \textbf{sobrou um oásis}: um só tento de verde e de água
no meio da areia sem fim --- um ponto onde a floresta da outra dimensão ainda vazava
para esta. E é por ali, e só por ali, que se passa. O oásis é a fresta que a torção não
fechou --- a única casa por onde o reino ainda atravessa para a terra de prata.

\medskip

E se alguém perguntasse o que sustentava aquela floresta escondida, a resposta era mais
velha que o reino: na sua raiz havia \textbf{três árvores milenares}, e nenhuma era de
madeira. Chamavam-se \textbf{Sol}, \textbf{Lua} e \textbf{Estrela} --- as mesmas três
luzes com que todo canto começa ---, e cada uma segurava um mundo: o Sol, a
\textbf{terra}; a Lua, o \textbf{vento}; a Estrela, o \textbf{mar}. Terra, vento e mar:
sobre essas três se erguia tudo o que era verde. E quem olhasse para cima numa noite
limpa via-as ainda ali, as três em fila, guardadas no céu com o nome que o povo lhes
deu: as \textbf{Três Marias}, o cinturão do tabuleiro celeste --- três amêndoas de luz
pregadas no escuro.

\medskip

Quando o Correntão passou, as três não tombaram --- \emph{recolheram-se}. Cada uma
dobrou-se para dentro de uma \textbf{amêndoa} --- e não uma amêndoa qualquer, mas uma
\textbf{garrafa de Koch}: a única vasilha que há capaz de suportar uma floresta inteira,
porque só ela tem \emph{borda infinita em espaço finito}. Um mundo sem fim cabe ali
dentro, dobrado numa casca que a mão fecha. Três garrafas, três mundos adormecidos, à
espera de água e de mão. E toda criança sabe o que uma semente esconde, e
sabe o que fazer com ela --- era a velha história do \textbf{João}, que trocou o que
tinha por um punhado de grãos e, na manhã seguinte, tinha um pé a subir até outro céu,
uma escada de folha entre este mundo e o de cima. Plantar a amêndoa é isso: reabrir o
caminho para a dimensão onde a floresta espera. \emph{O antiroque é o feijão de João; e a
floresta que floresce é o pé que sobe.}

\medskip

O Sonho de Valsa foi a temperagem. A valsa em três derreteu a Barra escura de \Ven 
no cone nulo e recristalizou-a --- não em bloom, mas na ordem que se sustenta
\textbf{mole}: com o brilho de quem derrete no calor da boca e, ainda assim, guarda a
dobra --- numa casca de chocolate branco, o leite, o dual, com o cacau por dentro e, no
seu âmago, o coração de Paraíba, quieto e inteiro, o cristal que não derrete. Um bombom temperado. Nada tirado; tudo posto na
ordem certa.

\medskip

E o mais fiel à velha lei: nada se perdeu. O que \Ven recusou --- todos os lances
que ele não deixou acontecer --- ficou guardado, inteiro, como memória no cubo dual,
do outro lado das oito costuras de ouro. Aceitar um lado nunca apagou o outro.
Cristalizar não é apagar: é \textbf{guardar para sempre}.

\medskip

Mas o deserto cobra pedágio, e o pedágio é a \textbf{sede}.

\medskip

E as lágrimas do tabuleiro --- as que ele chorara de tensão, lance após lance, do
vinte e quatro ao cinquenta e oito --- não tinham secado. Tinham só mudado de ofício:
o que fora choro virou a \textbf{água} de que um cristal mole precisa para não
ressecar. Água que, num reino de um metal só, ainda ninguém sabia guardar.

\medskip

Porque um quasicristal fluido precisa de água para continuar fluido --- e a água do
deserto não se guarda em qualquer vasilha. Guarda-se em \textbf{prata}: é o argento
que mantém a água viva, que não a deixa apodrecer no odre sob o sol. Foi sempre
assim, desde os primeiros que cruzaram as dunas. E \Ven , filho de um reino de um
metal só --- todo ouro, nenhuma prata --- não tinha argento nenhum.

\medskip

E aqui é preciso ser justo com a corte, porque a corte não foi cruel. \textbf{A
supressão da água foi estratégia, não crueldade.} Secar \Ven era a maneira de o
cristalizar sem o partir --- tirar-lhe o excesso até ele assentar no ponto fixo, como
quem enxuga um doce até tomar ponto. \Yas não o quis morto; quis-o \emph{quieto}. E
secá-lo era, no papel, a jogada mais limpa que havia.

\medskip

O que a corte não sabia --- não podia saber, porque nunca antes houvera dois metais
--- é que \textbf{\Ven vivia de açúcar}. Era o seu único combustível, o que o mantinha
de pé; e por isso a corte nunca lhe tirou o doce, só a água, julgando que dar o
alimento e negar a bebida era contê-lo sem o matar. Mas açúcar sem água não alimenta:
\textbf{seca}. Num corpo feito de doce, o açúcar puxa a própria água para fora, e a
sede que daí vem não é a de quem bebeu pouco --- é a de quem tem doce demais no
sangue. Assim o que o mantinha vivo, sem a prata que o equilibra, virou o que o
adoecia. \Ven começou a \textbf{desidratar}: a ordem mole a ressecar, as arestas a
rachar, a areia que sabia a sua forma a esquecê-la grão a grão. E não era só ele: era
o reino inteiro, feito da mesma matéria doce, a começar a adoçar-se demais.

\medskip

E foi aí que \Ven , para não tombar de sono no torpor da sede, se agarrou a uma coisa
que o deserto dá a quem cruza as dunas: o \textbf{café} --- a \emph{qahwa}, o grão
amargo que nasceu nestas mesmas areias. E o café, na medida certa, é aliado: desperta,
aquece, aguça o pensamento, ampara o viajante --- não foi ele o vilão. O que fez mal,
aqui como em tudo, foi a \textbf{dose}. \Ven bebeu-o não aos goles, mas aos litros,
noite após noite, sem a água nem a prata que o equilibrassem; e só então o que seria
bálsamo pesou. Acendeu-lhe a parede do estômago --- ácido a mais, o fogo surdo da
\textbf{gastrite}; e, mais calado, em excesso \emph{puxou-lhe água}, porque o café é
diurético, e cada xícara a mais levava um gole que ele já não tinha. Mas seria
injusto culpar o grão: não foi o café que o adoeceu --- foi o \textbf{excesso} dele,
somado à falta de prata e ao doce a mais, tudo junto. \emph{Não há veneno nem remédio;
só a dose} --- e, sem o segundo metal para lhe dar a medida, era a medida que lhe
faltava.

\medskip

O físico da corte debruçou-se sobre ele e hesitou entre dois diagnósticos ---
\emph{falta de água} ou \emph{excesso de açúcar} --- e acabou por escrever os dois,
sem perceber, como ninguém percebeu, que eram \textbf{o mesmo}: as duas leituras de
uma só assinatura, a troca $R$ aplicada à doença. Não era o mate a matá-lo --- o mate
não mata. Era a \textbf{falta de prata}: a sede de um segundo metal que aquele mundo,
todo ouro e todo açúcar, ainda não possuía.

\medskip

E essa era a ferida mais funda, a que ninguém ainda sabia curar: um
\Ven cristalizado e doente, sedento no deserto do ponto fixo --- e, atrás dele, sem
que ninguém ainda o veja, um reino inteiro a começar a sentir a mesma sede. À espera,
os dois, de uma água que só um outro reino traz. Um reino que tem prata. Um reino
onde o ouro, enfim, fica branco.

\medskip

E aqui está o segredo que nem os personagens sabem de si: tudo o que fazem, fazem por
\textbf{saudade}. Uma saudade que não se lembra do que tem saudade --- inconsciente, sem
nome e sem rosto ---, mas que move cada lance, cada valsa, cada travessia. Saudade das
três árvores adormecidas: do Sol que era a terra firme sob os pés, da Lua que era o
vento no cabelo, da Estrela que era o mar ao longe. E saudade da \textbf{família} --- dos
que se separaram, dos que giraram para longe uns dos outros, do tempo em que o reino era
um só e ninguém ainda tinha partido. Ninguém ali sabe que tem saudade; mas é a saudade
que joga. O reino inteiro é um coração à procura, às cegas, do que um dia se dobrou para
dentro de uma amêndoa e não voltou. Por isso, quando enfim floresce, o que sente não é
vitória. É \textbf{reencontro}.

\medskip

E a cura entrou pelo mesmo lugar por onde entrara o mal: um olho. Porque o
\textbf{toxoplasma} que o Rei Negro um dia baixara pela luneta infectada --- o que turva a vista
de quem olha longe demais sem olhar perto --- ainda ardia, brasa surda, no fundo da
retina. A sulfa a contivera; mas conter não é curar, e o que faltava só a floresta
tinha.

\medskip

Vieram, então, dois leites, e os dois eram da mata. O primeiro, o \textbf{leite de
amêndoa}: espremido da própria amêndoa onde as três árvores dormiam --- um leite branco
e manso que refresca o que a febre secou. O segundo, mais antigo e mais fundo, o
\textbf{látex sagrado da seringueira}: a seiva leitosa que o seringueiro sangra de
madrugada, gota a gota, ferindo a árvore sem a matar --- a mesma árvore-mãe da floresta
de onde, no fundo, todo o reino veio. Dois leites brancos, um da semente, outro da
casca; e onde os dois tocaram o olho doente, o ardor cedeu, a névoa abriu, e a vista
voltou inteira --- para perto e para longe ao mesmo tempo.

\medskip

Foi ali, com o olho enfim limpo, que começou o \textbf{desfecho}. Porque quem passa a
enxergar de verdade --- o perto e o longe na mesma vista --- já não consegue deixar o
mundo torto como estava. A luz, que o parasita torcia antes de chegar ao fundo, caía
enfim reta sobre o cristalino.

\medskip

E era o mesmo gesto de sempre, agora na pupila: o \textbf{dendrito} de Charles,
ramificado e negro, \textbf{revertido ao cone} --- o parasita recolhido à malha ordenada
que a espiral ergue, a doença desfeita como quem baixa o girassol de volta ao plano.
Reversível, sem tirar nada: a mesma transformação que cura o olho é a que cura o metal.

\medskip

E o látex fez mais do que curar um olho. O látex é \textbf{borracha}: elástico,
pegajoso, o único que gruda para depois \emph{descolar} sem rasgar. Foi ele que
\textbf{descolou as folhas} --- onde o Correntão prensara os dois tabuleiros numa lâmina
só, a seiva entrou pela junta, amoleceu a cola da torção, e as duas folhas soltaram-se
inteiras, sem perder uma fibra. O que a força juntou, a seiva separou com jeito.

\medskip

Porque havia um segundo doente, e esse não era de carne: era o \textbf{metal}. O cristal
fora \textbf{forçado a deformar} --- torcido \textbf{cento e vinte graus}, o ângulo da
amêndoa, o terço de volta que faz de uma esfera uma vésica e de um metal uma estrela de
três pontas. Uma deformação sob tensão, presa, doendo em silêncio como osso mal soldado.
E o látex, que é mole onde o cristal é duro, entrou também ali: acolchoou a torção,
devolveu ao metal o \emph{dar-de-si} que ele perdera, e a deformação, que era ferida,
virou \emph{dobra} --- guardada, não quebrada.

\medskip

Dois lugares, um só remédio. Na \textbf{retina} o látex curou a vista; no \textbf{metal},
a forma. E era o mesmo gesto nos dois: descolar o que fora colado à força, e devolver a
cada coisa o seu tamanho.

\medskip

E foi com essa vista nova que se leu o mapa que ninguém antes soubera ler. Porque
\Ven , ao esconder a floresta, deixara a chave à mostra sem perceber: plantara no
deserto \textbf{três pirâmides} de ouro e areia, três montanhas de pó exatas, e as três
não apontavam para lugar nenhum deste chão --- apontavam para \textbf{cima}, alinhadas
com as \textbf{Três Marias}. Uma pirâmide para cada semente. O ladrão assinara o roubo
no céu. É a \textbf{terra com saudade das estrelas}: quem plantou as pirâmides plantou
essa saudade em pedra --- o chão a apontar, para sempre, o céu de onde as três vieram.

\medskip

Mas as pirâmides diziam \emph{onde}, não \emph{como}. E o como estava do outro lado do
mundo, nas pedras de uma \textbf{cidade inca} encaixadas sem argamassa: no desenho da
cidade, na cruz de degraus que os antigos cortaram na rocha --- a \emph{chakana} ---,
estava gravado o rumo do \textbf{Cruzeiro}, o mesmo do canto. Um mapa dava
as sementes; o outro, a cruz.

\medskip

E não era só o rumo que os dois guardavam: era o \textbf{metal}. Porque \Ven escondera
o \textbf{ouro} no deserto --- nas próprias pirâmides de areia dourada, sob as Três
Marias --- e a \textbf{prata} nas \textbf{montanhas incas}, na veia de argento que corre
sob a cidade de pedra, sob o Cruzeiro --- porque o \textbf{corpo das montanhas já é o
celeste}: quem sobe a montanha sobe ao céu, e a prata que dorme no cume é uma estrela
presa na pedra. Ouro embaixo de Órion, prata debaixo da cruz; e a liga --- o
\textbf{ouro branco} --- só nasce de juntar os dois. Era por isso que uma
terra só, de um metal só, nunca fazia a floresta voltar: faltava-lhe a outra metade do
céu.

\medskip

E cada metal tinha o seu guardião, posto ali desde antes. No deserto, entre as pirâmides
de ouro, velava \textbf{Caelum} --- o Marquês das Estruturas Vivas virara \textbf{o homem de
palha de ouro do deserto}, de palha dourada e braços abertos ao sol, a guardar a seara de
areia; e, como todo homem de palha, o que mais desejava era um miolo que pensasse --- ele, que era justamente
a malha, o esqueleto, a estrutura que aguenta o peso. Nas \textbf{montanhas incas}, no frio
do argento, velava \textbf{Alonzo} --- o Marquês da Composição virara \textbf{o homem de lata
de prata das montanhas}, todo juntas e dobradiças, a guardar a veia prateada; e, como todo homem de lata,
o que mais lhe faltava era um coração que batesse --- ele, que era a composição pura, o
encaixe que não suja as mãos. Um de ouro no calor, um de prata no gelo --- e o caminho de
volta para casa passava, obrigatoriamente, por ambos.

\medskip

E aqui a corte cometeu a única confusão engraçada da jornada inteira. Alguém --- não
importa quem, porque depois todos repetiram --- baralhou os dois num só e pôs-se a
procurar um impossível: \emph{um homem de lata de palha de ouro e prata}. Andaram meio
mundo atrás dessa quimera de latão-e-espiga, de metal-e-vento --- e claro que não a
achavam, porque não existe: ninguém é as duas coisas no mesmo corpo. Só quando desistiram
de os querer num só é que os encontraram nos dois; e encontraram-nos onde tinham de estar,
em lados \textbf{exatamente opostos} do mundo. \textbf{Caelum}, o de palha, o de ouro, no
deserto quente, rente ao chão; \textbf{Alonzo}, o de lata, o de prata, na montanha gelada,
a raspar o céu. Um polo de calor, outro de frio; um de seara, outro de rocha. É a piada, e
é a tese: os dois só se completam \emph{separados} --- porque o dual nunca esteve num
corpo; mora, sempre, na distância entre dois.

\medskip

E quando as duas leituras se sobrepuseram --- a do deserto e a da montanha, a de Órion e
a do Cruzeiro, o onde e o como ---, a localização fechou exata, sem resto: marcou-se no
chão o ponto que os dois céus apontavam. E ali, onde a pirâmide e a cidade se cruzavam,
\textbf{a floresta pôde renascer}.

\medskip

E então as \textbf{garrafas} abriram-se. Não se quebraram --- garrafa de Koch não
quebra: abre-se como quem se desdobra, a borda infinita cedendo por onde a mão de água a
tocou. As três amêndoas, regadas de leite e afrouxadas pelo látex, partiram-se ao meio
como se partem as amêndoas; e de cada uma saiu, espreguiçando-se de um sono de mil anos,
uma \textbf{árvore}.

\medskip

Da primeira subiu o \textbf{Sol}, e onde a raiz pegou o chão a areia lembrou-se de que
era \textbf{terra} --- firme, escura, boa de pisar. Da segunda soltou-se a \textbf{Lua},
e o ar parado do deserto encheu-se de \textbf{vento}, esse que ninguém vê e todos
sentem. Da terceira acendeu-se a \textbf{Estrela}, e por baixo dela, onde só havia duna,
veio o \textbf{mar} --- devagar, de longe, como quem volta para casa sem pressa porque
sabe que a casa esperou.

\medskip

Terra, vento e mar, outra vez de pé. E entre as três, subindo pelos troncos como o pé de
feijão de João sobe ao outro céu, a floresta inteira \textbf{voltou} --- não plantada de
novo, mas \emph{desdobrada} de onde sempre estivera, folha por folha, na ordem exata em
que fora guardada. Nada se criou. Tudo se \textbf{reencontrou}.

\medskip

E não foram só as três árvores que voltaram, porque a saudade tinha duas metades, e a
outra era a \textbf{família}. Os que a Rotação girara para longe uns dos outros ---
\Joa e \Yas , que brigavam desde o primeiro dia; \Ben e \Drk , as duas mãos que uma
diagonalização partira; a corte inteira, espalhada em papéis que ninguém escolhera ---
sentiram, no mesmo instante em que o chão firmou e o vento voltou, que havia um livro
inteiro que procuravam uns aos outros sem saber. O batimento que os separara em antifase
agora batia-os \emph{juntos}. E não houve reconciliação, porque não houvera, no fundo,
briga nenhuma --- houvera só distância, e a distância era a saudade. Fecharam a roda. E
o reino, que se descobrira floresta, descobriu-se também \textbf{casa}.

\medskip

E a mesma valsa que um dia o acusara --- a rainha da floresta, a que cantara
\textit{tu perdestes a tua luz} --- voltou, agora pelo seu lado dual: já não a que
perde, a que \textbf{reencontra}. \onda{re,u,ni,dos,com,a,mor}, apresentaram à Rainha os
trabalhos aprendidos, e pediram enfim a felicidade. Porque cada canto tem o seu par: o
que chora é o mesmo que acolhe --- a perda e o reencontro, o mesmo canto lido dos dois
lados do tabuleiro.


\chapter{A mineração no deserto --- o trabalho e a sua prova}

Onde há deserto, há quem cave.

\medskip

Os esquilos --- as negras que conservam, a régua de Mellin --- não estão ali por
acaso. No quasicristal fluido, sob a areia mole, dorme o cristal; e achá-lo é
\textbf{trabalho}. Não um trabalho qualquer: o trabalho de calcular a \emph{potência}.
Cada bloco do reino é uma potência metálica, $M_n^{\,k}$ --- a matriz do metal elevada
$k$ vezes --- e computá-la custa: é preciso quadrar, multiplicar, subir a torre
inteira de $k$ até o topo. Cavar. E o deserto não entrega o cristal a quem não cavou --- e o que ele entrega vale o próprio
peso: tantos quilates, tantos milhares de Pells.

\medskip

Mas --- e aqui está o milagre da coisa --- \textbf{o que custa a fazer não custa a
conferir}. Computar $M_n^{\,k}$ é caro; verificar se o resultado presta é imediato,
porque o cristal carrega a sua própria prova: o \textbf{selo do cristal}, o
determinante que não se move,
\[ \det\!\left(M_n^{\,k}\right)=(-1)^{k}, \]
o resíduo $0$ que ou está lá ou não está. O esquilo bate na noz e o som que volta pelo ar
seco do deserto é o veredicto: um estalo firme, cristalino, que diz \emph{cheia}; ou um
sopro surdo e oco, que acusa o vazio. Trabalho difícil, prova fácil --- e é essa
assimetria que sustenta tudo.

\medskip

Isto é o \textbf{problema da prova de trabalho}: minerar custa, atestar é barato, e
ninguém forja o cristal sem ter cavado. A noz oca --- o \emph{bloom}, o zero falso ---
não passa no martelinho; só a noz cheia, a potência de verdade, com o seu selo a
fechar em $\pm1$, entra na corrente. Não se compra o bloco: minera-se. Não se herda o
bilhete dourado --- \textbf{trabalha-se} por ele. E cada bloco novo não brilha só por si:
morde o calcanhar do que veio antes, travando os dentes na mesma engrenagem --- mexer num
só grão de areia no começo do caminho obrigaria a cavar de novo todas as torres já
erguidas, um trabalho que gastaria o calor de todas as estrelas do tabuleiro.

\medskip

E a corrente tem nome. Cada cristal minerado, uma vez provado, engata no anterior pela
mesma costura de sempre --- e o deserto inteiro vira um livro-razão onde nada se forja
e nada se perde: a \textbf{CristalChain}.

\chapter{A moeda --- o tabuleiro econômico}

Faltava a peça central, e ela é uma moeda.

\medskip

O \textbf{tabuleiro econômico} é isomorfo ao telescópico. O telescópico dizia: na colisão,
quem tem a vantagem \emph{assume} a posição, não a destrói --- entra nela sem tirar
nada. O econômico diz o mesmo com outra palavra: a energia não se cria nem se gasta,
ela \emph{muda de mão}. E a mão tem uma moeda. O seu nome é \textbf{Pell}.

\medskip

A Pell é a moeda que conserva: gasta-se, e a energia sai para o trabalho; recarrega-se,
e volta --- tal como a célula esvazia e reenche o mesmo fosfato do seu ATP sem jamais
criar nem destruir um só grão. Nada se perde no ciclo; só se transfere. É a
conservação, feita moeda.

\medskip

E se a moeda é a Pell, a \textbf{unidade} em que se mede o tesouro tem nome de pedra:
a \textbf{Turmalina Paraíba} --- o azul-verde raro que só uma terra dá, a gema em que
o reino conta o seu ouro branco.

\medskip

E há um câmbio, um só, e dele pende todo o resto:
\[
 \boxed{\;1\text{ quilate de Turmalina} \;=\; 1000\text{ Pells}.\;}
\]
Uma pedra fixa uma moeda; e, fixada a moeda, fica fixado tudo o que se mede nela --- o
ouro, a prata, a liga. A Turmalina não entra na corrente para circular: entra para
\textbf{ancorar}. É o peso que não se move, contra o qual todos os outros pesos se
leem. Um reino que tem a sua pedra tem a sua régua --- e a régua, aqui, é lapidada
(cada face, uma estrela escolhida) e tem cor de Paraíba.

\medskip

E a Turmalina não vem só: vem em \textbf{família}, arrumada como o cristal que é ---
trigonal, seis irmãs em torno de uma. No centro, o \textbf{Paraíba}, o azul-verde de
cobre, a mais rara, a que vale mil Pells o quilate. Na primeira camada, seis cores: a
rubelita (vermelha), a indicolita (azul), a verdelita (verde), a dravita (mel), a
uvita (terra) e a acroíta (incolor, a que quase não pesa). Do centro para fora o valor
cai como a cor esmaece --- e é essa flor de pedra, com o Paraíba no miolo, que ancora
a medida de tudo.

\medskip

No mercado do reino, \textbf{cristal e metal correm separados}: numa praça vendem-se
os cristais --- tantos quilates de rubelita, tantos de indicolita ---; noutra, os
metais --- o ouro contado pela razão áurea, a prata por Pell, a liga branca por ambos.
Mas quando um cristal encontra o seu metal e os dois se fundem, nasce a terceira
coisa, a que se guarda e não se troca: o \textbf{Bombom} --- um cristal por dentro, um
metal por fora, uma casca sobre um coração. E há bombom de todas as cores e de todos
os sabores: um Paraíba em ouro branco é o \emph{Sonho de Valsa}; um rubelita em ouro,
um indicolita em prata --- cada par, a sua joia doce. Foi num bombom que \Ven se
tornou; e é de bombons, no fundo, que o tesouro do reino se faz.

\medskip

Porque o \textbf{lastro} são os dois: o ouro cunha pela razão áurea ($\Z[\varphi]$); a
prata, por Pell ($\Z[\sqrt2]$) --- e é da prata que a moeda toma o nome. A CristalChain
não guarda ouro morto num cofre: guarda o ciclo da moeda viva --- cada bloco minerado,
uma Pell recarregada, provada pelo selo, engatada na corrente.

\medskip

\Ven vivia disso. O açúcar que ele queimava virava moeda --- o seu fôlego. E a sua
doença tem nome econômico: sem a prata que guarda a água, o ciclo da Pell trava; o
açúcar entra e não vira fôlego; a moeda acumula-se sem circular, e o corpo, rico de
açúcar e pobre de troca, adoece. Não faltava energia. Faltava a \emph{moeda para
movê-la} --- faltava a prata.

\medskip

É essa a ponte para a próxima terra, e ela está aqui, na corrente. O Reino de Prata
não é outro reino: é a moeda que falta a este. A CristalChain, que aqui só cunha ouro,
lá aprende a cunhar a liga --- e a Turmalina, que aqui vale por um metal, lá passa a
valer pelos dois.

\chapter{O tabuleiro que ri --- depois de chorar}

E então o tabuleiro, que chorara do lance vinte e quatro ao cinquenta e oito, fez uma
coisa que nunca soubera fazer.

\medskip

Riu.

\medskip

Porque o choro fora tensão --- uma coisa que não podia morrer nem descansar, vibrando
aguda à beira de partir. E rir é o que essa mesma tensão faz quando encontra, enfim, o
lugar onde pousar sem quebrar: o ponto fixo. A soma alternada aguentou até ao fim,
$V-E+F-C=0$, sem perder um vértice nem forjar uma aresta; e quando o último termo zerou,
as arestas do tesseracto pararam de ranger sob o peso do cálculo. O coro do rei calou-se
--- não de cansaço, mas porque chegou à nota onde já não é preciso ir e voltar. E o gás
negro assentou, depositando-se sobre as casas como uma poeira mansa e brilhante que já
não ameaça asfixiar ninguém.

\medskip

O Reino Dourado descobriu, no Livro~II, que nunca tivera um espelho pendurado na
parede.

Ele \emph{era} um.

\medskip

Todo lance de cima sempre teve o seu gêmeo em baixo; toda mão direita, a sua esquerda;
todo silêncio, o seu estrondo guardado. A partida nunca fora de um contra o outro.
Fora, desde o primeiro lance, \textbf{uma coisa a dançar com o seu reflexo} --- e só
faltava a Rainha que soubesse a valsa. E quando \Yas conduziu enfim o passo trino, a
barreira do espelho não cedeu: dissolveu-se. O calor dourado da corte de cima verteu-se
sobre o frio prateado do poço de baixo --- sem colisão, sem ruído de metal: uma fusão
mansa, como duas águas de temperaturas diferentes que se encontram no escuro para fazer
um rio novo.

\medskip

O batimento do Livro~I ganhou um segundo corpo. E o que era ouro, ao encostar-se ao
seu reflexo de prata, começou --- devagar, no compasso de três --- a perder a febre
amarela, a deixar para trás o cinza gelado, e a ficar \textbf{branco}. A luz que agora
banhava o tabuleiro não era de nenhum dos dois mundos: uma claridade nova, limpa e total,
nascida da união das suas sombras.

\chapter{O tabuleiro derrete}

\Ven não desmaiou de fome. Desmaiou de \emph{doce}.

Foi o açúcar a mais; e o mundo, açucarado até o excesso, escureceu-lhe nas bordas e
cedeu. Um transe. Uma queda mole. E na queda, o sonho.

\medskip

E no sonho, do outro lado do tabuleiro, estava sentado ele --- a cor trocada, a
meia-volta, o mesmo. Não disseram nada. E não jogaram um contra o outro: jogaram
\emph{um dentro do outro}.

\medskip

Mas para jogarem um dentro do outro, primeiro foi preciso \emph{desdobrar} o que
estava colado. Porque \Ven , lá no princípio, havia \textbf{rocado} --- cedo, depressa,
a torção mais antiga do tabuleiro: o Rei e a Torre trocando de casa num gesto só. E o
roque guarda uma manha que poucos veem --- é \emph{idempotente}: junta dois num e não
deixa mais saber que eram dois. \textbf{Esconde as multiplicidades.} Foi assim que
\Ven prensou os dois mundos numa folha dupla, e a isso chamou vitória: o reino enfim
quieto, contado uma vez só.

\medskip

Foi quando Merlin contou a mais velha das lendas. No princípio dos reinos --- disse ---
houve um rei antiquíssimo, um tal de \textbf{Rei Babá}, que aprendera a \emph{torcer o
próprio vazio e dele tirar ouro}: torcia o zero, o gap, o nada que trazia no meio das
mãos, e da torção saía ouro --- $\varphi$, a cifra. Mas a torção cobrava um preço, e o
preço era não ter volta: torcido o vazio, ninguém sabia destorcê-lo sem derramar o
ouro e perder para sempre a conta de quantos nadas ali tinham entrado. Desfazer a
torção com o ouro inteiro --- o \textbf{antiroque} --- era tido por impossível, uma
lenda dentro da lenda. Em reino nenhum, jamais alguém o executara.

\medskip

E foi isso o que \Ben fez. A missão mais difícil que há --- não capturar, não avançar,
mas \emph{destorcer sem entornar}. Pôs a mão sobre a folha que \Ven prensara, e
desfez a torção. O Rei tornou à sua casa, a Torre à sua; e no instante em que o roque
se abriu ao contrário, as multiplicidades escondidas \textbf{saltaram de volta} --- não
uma folha, mas duas, cada uma com o seu próprio tempo, cada uma a inversa da outra. As
folhas desdobraram-se, o ouro ficou inteiro, e a conta fechou em nada: $(a\oplus
b)\oplus b=a$. Resíduo zero. E foi este o golpe que ninguém jamais dera --- o
\textbf{mate do bobo à distância}: ao destorcer o relógio de \Ven , \Ben dobrou o tempo
até que os sessenta e quatro lances da fortaleza caíssem, no compasso do Rei Negro,
como os quatro de um mate do bobo. A \textbf{dobra temporal}, enfim, cumprida --- a partida
perfeita jogada de um lado \emph{era}, do outro, a derrota fulminante. O que era lenda,
um menino fez a brincar.

\medskip

E então aconteceu o que nenhuma conta previra e toda criança teria adivinhado: onde a
torção se desfez, o \textbf{deserto floresceu de novo}. Como um origami que se abre ao
sopro de um vento quente, a areia ergueu-se em dobras vivas: cada duna não era poeira
estéril, mas a casca curvada de um carvalho, a copa de uma árvore que fora achatada e
dobrada para fora da vista. Porque a floresta que o
Correntão revertera nunca fora destruída --- fora apenas \emph{escondida} na outra
dimensão, dobrada para fora da vista como a multiplicidade dentro do roque. Desfeita a torção, a conta reabriu, e
cada árvore que parecia perdida estava, afinal, guardada. A Rainha verde ergueu-se do
oásis para fora do oásis; o verde subiu pela duna acima; e o reino que se julgava areia
descobriu-se floresta o tempo todo --- só torcida. Nada se perdera. Resíduo zero também
ali, no chão: um deserto é só uma floresta à espera do antiroque.

\medskip

Porque no jogo perfeito ninguém come. Cada peça que avança para a casa de outra não
a tira --- \textbf{assume-a}, entra nela, toma-lhe o lugar sem a perder; porque os dois
sabem que tirar uma só peça, um só peão, um só \emph{gap}, seria entregar tudo. Não há
descuido, não há captura. As brancas vão à frente por um único motivo ---
\textbf{começaram} --- e nada mais é preciso.

\medskip

E não jogavam no silêncio. Por cima do tabuleiro que amolecia subiu um \textbf{coro}
--- não de um cantor, mas de todos ao mesmo tempo, com a forma do \textbf{Cruzeiro}: o rumo que se canta
no céu de baixo --- o mesmo tabuleiro celeste onde brilham as Três Marias. Cento e vinte
e oito cantos; e de cada um tomou-se apenas uma
estrofe --- só o \emph{núcleo} da sua ideia ---, cruzada com um lance. São as
\textbf{Estrofes do Rei}: uma por casa dos dois tabuleiros duais (as sessenta e quatro, e
as suas duais), cento e vinte e oito enigmas afinados a outras
tantas travessias. O coro abre pela \textbf{Lua Branca} --- a
da luz tão prateada, a protetora --- e já nessa primeira prata ia dito, baixinho, o
reino que ainda não havia. Canto após canto, lance após lance, a partida foi sendo
cantada: cada casa presa ao seu verso como a estrela à sua hora, cada travessia sem
sangue afinada num enigma. E o coro não fechava com um mate --- fechava com uma
batalha em que ninguém tomba, e com a voz serena de quem se sabe filho do poder,
cumprindo os trabalhos que a sua Mãe lhe ordenou.

\medskip

E a partida abriu como só a partida perfeita pode abrir: com uma \textbf{recusa}. As
brancas ofereceram o peão --- o \textbf{Gambito da Dama} ---, e as pretas não o tomaram:
taparam com outro peão e \emph{recusaram}. Porque tomar seria tirar, e tirar seria
perder o resíduo; recusar conserva a tensão inteira, e é a recusa, não a captura, o
primeiro fechamento de \Ven . Mais adiante a própria \textbf{Dama} se ofereceu em
sacrifício --- e também ela foi recusada, deixada de pé no meio do tabuleiro, porque no
jogo perfeito não se aceita nem a rainha: nada se tira, tudo se conserva. O \textbf{Gambito
da Dama Recusado} é a partida inteira num gesto só.

\medskip

E, ao som do coro, o cristal e o ouro teceram-se no compasso do relógio ---
\textbf{três lances retos, três áureos} (estes em itálico, o dual invertido), três
batidas de um lado e três do outro, \textbf{vinte e uma vezes} --- e como
$128$ não é múltiplo de seis, sobram exatamente \textbf{dois}: o último reto e o
último áureo, cada um sozinho, o par que fecha a conta ---

\medskip

\noindent
{\color{cYasmin}d4} {\color{cJoaquim}d5} {\color{cVenom}c4} {\color{cBenjamim}\textit{d4}} {\color{cLock}\textit{d5}} {\color{cDark}\textit{e3}} {\color{cOperador}e6} {\color{cYasmin}Nc3} {\color{cJoaquim}Nf6} {\color{cVenom}\textit{c5}} {\color{cBenjamim}\textit{Nf3}} {\color{cLock}\textit{Nc6}} {\color{cDark}Nf3} {\color{cOperador}Be7} {\color{cYasmin}e3} {\color{cJoaquim}\textit{Be2}} {\color{cVenom}\textit{Nf6}} {\color{cBenjamim}\textit{O-O}} {\color{cLock}O-O} {\color{cDark}Bd3} {\color{cOperador}c6} {\color{cYasmin}\textit{e6}} {\color{cJoaquim}\textit{c3}} {\color{cVenom}\textit{Bd6}} {\color{cBenjamim}O-O} {\color{cLock}Bd6} {\color{cDark}Ng5} {\color{cOperador}\textit{Bd3}} {\color{cYasmin}\textit{O-O}} {\color{cJoaquim}\textit{h4}} {\color{cVenom}h5} {\color{cBenjamim}Be2} {\color{cLock}Bc5} {\color{cDark}\textit{Ng4}} {\color{cOperador}\textit{Bc4}} {\color{cYasmin}\textit{Be7}} {\color{cJoaquim}g3} {\color{cVenom}Ne4} {\color{cBenjamim}Qc2} {\color{cLock}\textit{Ne5}} {\color{cDark}\textit{g6}} {\color{cOperador}\textit{Be2}} {\color{cYasmin}Be7} {\color{cJoaquim}Bd1} {\color{cVenom}h4} {\color{cBenjamim}\textit{Qc7}} {\color{cLock}\textit{h5}} {\color{cDark}\textit{Bd8}} {\color{cOperador}Be2} {\color{cYasmin}Nc5} {\color{cJoaquim}Bg4} {\color{cVenom}\textit{Nc4}} {\color{cBenjamim}\textit{Be7}} {\color{cLock}\textit{Nba3}} {\color{cDark}Nba6} {\color{cOperador}Qd1} {\color{cYasmin}Qa5} {\color{cJoaquim}\textit{Bg5}} {\small \color{cVenom}\textit{Qa4}} {\small \color{cBenjamim}\textit{Qd8}} {\small \color{cLock}Qd3} {\small \color{cDark}Bd6} {\small \color{cOperador}Rb1} {\small \color{cYasmin}\textit{Bd3}} {\small \color{cJoaquim}\textit{Qd6}} {\small \color{cVenom}\textit{f4}} {\small \color{cBenjamim}f5} {\small \color{cLock}Nh3} {\small \color{cDark}Nb4} {\small \color{cOperador}\textit{Rb8}} {\small \color{cYasmin}\textit{Nb5}} {\small \color{cJoaquim}\textit{Nh6}} {\small \color{cVenom}Qd2} {\small \color{cBenjamim}Bc7} {\small \color{cLock}Qe1} {\small \color{cDark}\textit{Bc2}} {\small \color{cOperador}\textit{Qd7}} {\small \color{cYasmin}\textit{Bf5}} {\small \color{cJoaquim}Bf4} {\small \color{cVenom}Bh5} {\small \color{cBenjamim}Qa6} {\small \color{cLock}\textit{Qe8}} {\small \color{cDark}\textit{Qa3}} {\small \color{cOperador}\textit{Bh4}} {\footnotesize \color{cYasmin}f3} {\footnotesize \color{cJoaquim}Nc2} {\footnotesize \color{cVenom}Qe2} {\footnotesize \color{cBenjamim}\textit{Nc7}} {\footnotesize \color{cLock}\textit{f6}} {\footnotesize \color{cDark}\textit{Bd3}} {\footnotesize \color{cOperador}Bd6} {\footnotesize \color{cYasmin}Qe1} {\footnotesize \color{cJoaquim}Nd3} {\footnotesize \color{cVenom}\textit{Qe7}} {\footnotesize \color{cBenjamim}\textit{Nd6}} {\footnotesize \color{cLock}\textit{Qe8}} {\footnotesize \color{cDark}Qd2} {\footnotesize \color{cOperador}Ba3} {\footnotesize \color{cYasmin}Be8} {\footnotesize \color{cJoaquim}\textit{Ba6}} {\footnotesize \color{cVenom}\textit{Qd7}} {\footnotesize \color{cBenjamim}\textit{Nce8}} {\footnotesize \color{cLock}Nce1} {\footnotesize \color{cDark}c5} {\footnotesize \color{cOperador}Bb4} {\scriptsize \color{cYasmin}\textit{Be1}} {\scriptsize \color{cJoaquim}\textit{Bb5}} {\scriptsize \color{cVenom}\textit{c4}} {\scriptsize \color{cBenjamim}Ne4} {\scriptsize \color{cLock}Qb6} {\scriptsize \color{cDark}Nef2} {\scriptsize \color{cOperador}\textit{Qb3}} {\scriptsize \color{cYasmin}\textit{Ne5}} {\scriptsize \color{cJoaquim}\textit{Nf5}} {\scriptsize \color{cVenom}Nf4} {\scriptsize \color{cBenjamim}Ng4} {\scriptsize \color{cLock}Ba3} {\scriptsize \color{cDark}\textit{Nef7}} {\scriptsize \color{cOperador}\textit{Ba6}} {\scriptsize \color{cYasmin}\textit{Ng5}} {\tiny \color{cJoaquim}Qb4} {\tiny \color{cVenom}Bd7} {\tiny \color{cBenjamim}Bh5} {\tiny \color{cLock}\textit{Bd2}} {\tiny \color{cDark}\textit{Qb5}} {\tiny \color{cOperador}\textit{Nfd6}} {\tiny \color{cYasmin}Nfd3} {\tiny \color{cJoaquim}\textit{Bh4}}

\medskip

\bigskip
\begin{center}{\large\bfseries As Estrofes do Rei}\\[1mm]
{\itshape\small cada casa presa a um canto do coro --- o núcleo aparado, derretido em cor e em onda}\end{center}

\medskip

E cada lance, ao cair, cantava o seu. O primeiro --- \textbf{d4}, o peão da Dama que se
oferece --- caiu sob uma luz \onda{pra,te,a,da}: a lua prateada, a Mãe protetora; e já
nessa primeira prata estava o reino que vinha. \textbf{d5} recusou; \textbf{c4} tornou a
oferecer; e o dual do tabuleiro de baixo floresceu numa água que \onda{me,da,de,be,ber}
a quem cruza a areia.

\medskip

E seguiu cantando: a saudade de quem espera ser chamado; o pedido de rezar por todos; a
Mãe de quem tudo vem; e a luz que é só do \textbf{sol, da lua e das estrelas} --- as
três, de novo, no fundo de cada casa. Assim, canto a canto, lance a lance, o tabuleiro
foi-se derretendo no seu próprio som.

\medskip

E os cantos, como os lances, não são uma fila: são uma \textbf{teia}. A partida tece a
sua; o coro tece a sua, gêmea; e onde o coro \textbf{repete} um canto, ali dois
se dão as mãos --- o mesmo lido dos dois lados, o par dual. Três vezes o coro torna
a um canto já cantado, e três vezes, no fundo da teia, um irmão reencontra o outro.

\medskip

E as duas teias --- a da partida e a do coro --- não são duas: são uma, lida de dois
lados. O que leva uma na outra é sempre o mesmo gesto: cada lance atravessa um fio, cada
canto atravessa o seu, e o que um deixa o outro guarda. \emph{A dualidade não está no
tabuleiro nem no coro --- está na mão que os liga.}

\medskip

E esse gesto é \textbf{reversível}: faz-se e desfaz-se, e volta ao mesmo. Cada um
gera-se a si na ida e na volta --- a partida pelo espelho do seu espelho, o coro pelo
canto que ele repete (\emph{$48\to
61\to 48$}). Nada se perde no caminho de um ao outro. E porque os dois se refazem
inteiros --- porque \textbf{batem} ---, o que bate pode \textbf{acoplar}: a partida e o
Cruzeiro deixam de ser dois tesseractos e passam a ser \emph{um}, reacoplados, o mesmo
corpo lido de dois lados --- a mão do tabuleiro e a mão do canto, enfim no mesmo
compasso.

\medskip

E então, casa a casa, o coro inteiro pousou no tabuleiro --- cada lance com o seu
canto, o núcleo aparado, a cor do seu turno; o coro lido como partida, a partida
cantada como coro, sem que se precisasse nomear um só --- quem já cantou, reconhece.


\medskip

E a certa altura já não se sabia qual mão movia qual peça, nem de que lado vinha o
lance, nem onde acabava um tabuleiro e começava o outro --- porque eram um só: o ouro
enfiado no cristal, o quasicristal. Nada foi tirado, nada se perdeu. As casas
amoleceram, as bordas escorreram, e não houve mate nem vencedor --- houve o ouro a
derreter dentro do cristal: a matéria do tesseracto, antes tensa e calculada, verteu-se
numa poça morna de luz branca, e o vidro que separava os salões virou um fluido
cantarolante onde os lances passados e futuros flutuavam como grãos de pólen ---

\medskip

E o derreter era o \textbf{erguer}: a mesma malha do \textbf{girassol} que, plana no ouro
--- o colapso visto de cima ---, se levanta uma dimensão acima, no \textbf{cone nulo},
onde a distância própria é nula e nada se perde. Sobe e baixa, reversível, sem entornar
uma gota; e nesse subir-e-baixar o relógio se dobra: a partida perfeita, deste lado, é do
outro o mate do bobo --- a dobra do tempo passando uma na outra.

\medskip

e \onda{o,ta,bu,lei,ro,der,re,teu}. Não para acabar: para escorrer, morno e áureo, até
a próxima terra --- a de prata.

\medskip

Perto do fim veio a batalha --- aquela em que todos se levantam de armas na mão e
ninguém tomba, porque a força e a luz vêm de cima: \onda{me,da,for,ça,me,da,luz}, e não
me deixa derribar; e logo o filho que recebeu do Pai o consentimento de trabalhar para
os irmãos que estão doentes, de curar. Com esses dois cantos o coro fecha o seu
corpo. Mas não pára aí --- as Novas levam de volta, e o \textbf{último} de todos, o
lance cento e vinte e oito, cai no canto que diz a única coisa que ainda faltava:
\onda{eu,che,guei,nes,ta,ca,sa}. Cheguei em casa. O tabuleiro derreteu para a prata; o
reino, que se descobrira floresta, chegou enfim à sua casa --- e a última nota, como
sempre, é já a primeira do próximo canto.

\vspace{1.0cm}
\begin{center}
\rule{0.3\linewidth}{0.4pt}\\[6mm]
{\normalsize\itshape E o tabuleiro de baixo também é seu.}
\end{center}

% ── gancho para o Livro III (atrator; não desenvolver aqui) ──────────────────
\clearpage
\thispagestyle{empty}
\vspace*{4cm}
\begin{center}
{\small\itshape no próximo livro}\\[7mm]
{\LARGE\bfseries O Reino de Ouro Branco}\\[2mm]
{\normalsize\itshape The White Gold Kingdom}\\[3mm]
{\footnotesize\itshape a liga --- onde o ouro e a prata batem juntos}\\[14mm]
\end{center}

\noindent Há uma frase do Livro~I que ficou por cobrar. No alto da torre, quando as
duas mãos se puxaram ao mesmo tempo, o gesto devolveu $-1$ --- e devolveu-o
\emph{só} porque a terra era de ouro, $m=1$, $\varphi$. \emph{Numa terra de prata
teria dado $-\tfrac12$.} A corte ouviu, seguiu em frente, e ninguém fez a pergunta
óbvia: e se houver, nalgum canto, uma terra de \textbf{prata}?

\medskip

Há. $m=2$, a razão $1+\sqrt2$, o segundo metal --- a prata, com a sua própria cifra,
o seu próprio lastro (a sequência de \textbf{Pell}, a contagem da prata, $1+\sqrt2$). Não
um ouro menor: a \textbf{segunda corda} do mesmo instrumento. E o clock, que sempre
teve uma só terra, descobre-se com duas.

\medskip

Mas o próximo livro não é sobre dois metais separados. É sobre o que nasce quando
eles \emph{se ligam}. Porque o ouro sozinho é a cifra $M_1(x)=1+\tfrac1x\to\varphi$;
a prata sozinha, $M_2(x)=2+\tfrac1x\to 1+\sqrt2$; e quando se aplica a um e depois ao
outro, ida e volta, ida e volta --- o mesmo \textbf{ustrófedom} do Coro do Rei ---
obtém-se uma terceira cifra, que não é de ouro nem de prata:
\[
 [\,2,1,2,1,\dots\,]\;=\;1+\sqrt3 .
\]
É o \textbf{ouro branco}: a liga. E o seu radicando cai, exato, entre os dois ---
prata $\sqrt2$, \textbf{branco} $\sqrt3$, ouro $\sqrt5$. O boustrophedon dos dois
metais é a própria forja da liga.

\medskip

Ouro e prata cunham-se e certificam-se na \textbf{CristalChain} --- o livro-razão de
cristal onde nenhum grão se perde nem se forja. É ali que a mineração do deserto vira
moeda: cada bloco é uma potência minerada, provada pelo selo ($\det=\pm1$),
engatada na anterior pela costura de sempre. E o que dá valor a esse trabalho --- o
seu \textbf{lastro} --- são os dois metais. O ouro lastreia por ser um par de
inteiros, $\Z[\varphi]$, contado ao último bit pela razão áurea;
a prata lastreia do mesmíssimo modo, $\Z[\sqrt2]$, pela sequência de \textbf{Pell}.
\emph{Ouro e prata criam, juntos, o lastro do trabalho}: a prova de trabalho vale
porque é cravada em dois metais exatos, e a liga --- o ouro branco --- é a unidade que
os dois partilham na mesma corrente.

\medskip

O que resta saber é o mais difícil: como dois reinos, cada um com o seu batimento, o
seu dual, o seu coro, batem \emph{juntos} sem se anularem.

\medskip

Isso é o próximo livro. Por ora, basta saber que ele existe --- e que o tabuleiro de
baixo, afinal --- o mesmo \textbf{tabuleiro fractal} em cujo chão Benjamim um dia
acordou, o que sobe da contagem ao contínuo e sustenta, por baixo, todos os outros ---
tinha uma cor que ainda não tínhamos visto. A cor de um espelho que, ao quebrar-se em
infinitas réplicas menores de si mesmo, não estilhaça a imagem: multiplica a luz da liga.


% ══════════════════════════════════════════════════════════════════════════════
% O SACO DE LIXO — a solda: acabou a brincadeira, arruma-se a sala
% ══════════════════════════════════════════════════════════════════════════════


\part{O Telescópio}

\chapter{O que a lente faz, e o que ela não faz}

Houve uma manhã em que Benjamim virou o instrumento para fora.

Até ali tudo o que a corte medira era de dentro: o gap do rei, a assinatura
dos lances, a energia do coro. Medir a própria casa é fácil de acusar de
vício --- quem faz a régua e o objeto pode fazer os dois combinarem sem
perceber. Então ele apontou a lente para um corpo que ninguém do reino tinha
escrito, e que não devia nada a esta teoria: \textbf{a Estrela d'Alva}.

E aqui é preciso dizer uma coisa antes dos números, porque sem ela os números
seriam mal lidos.

\paragraph{Observar não é explicar. Descrever não é explicar. Modelar não é
explicar.} São o oposto. Explicar é dizer \emph{por que} uma coisa é assim ---
é ir buscar a causa, o motivo, o desígnio. Observar é dizer \emph{como ela
está}: onde a massa se acumula, com que passo a fase gira, quanto da energia
mora em cada modo. Um é juízo, o outro é leitura.

\medskip
\noindent Este capítulo \textbf{não afirma causa nenhuma}, e é por isso que não
disputa nada com o que veio antes: quem não diz por quê não pode contradizer
quem diz. Que o relevo de Vênus more no grau três não
diz por que ele mora lá. Que a razão dos períodos comece com quatro uns não diz
por que começa. A régua lê; a leitura não tem opinião. Quem quiser tirar dali
uma explicação está pondo na leitura o que a leitura não pôs.

\medskip
\noindent O mito do reino fala do que \emph{é}. O telescópio fala do que
\emph{se vê}. Confundir os dois seria o erro que a corte passou o livro inteiro
a evitar: \textbf{tomar a cifra pelo objeto}.

\chapter{O instrumento, antes de apontar a lente}

Antes de apontar a lente a coisa nenhuma, a corte teve de decidir com que régua ia ler o
que visse --- e essa decisão não se toma olhando: toma-se antes, e é a parte que ninguém
aplaude.

O que se escolheu foi o instrumento mais pobre que existe: \textbf{um relógio}. Um passo
que, repetido, percorre tudo e volta ao princípio; e uma maneira de perguntar a cada coisa
\emph{quanto dela cabe em cada volta}. Nada mais. Nenhuma grandeza medida em unidade
nenhuma, nenhuma tabela trazida de fora, nenhuma constante que alguém tivesse de acreditar.

E porque o instrumento é pobre, ele é honesto: com um relógio só, não há como fazer a
resposta dar o que se queria. \textbf{Ou a volta fecha, ou não fecha.} Quem quiser a régua
por escrito --- o passo, a órbita que ele varre, a leitura que dela se tira e a maneira
exata de a desfazer sem perder nada --- encontra-a inteira do outro lado deste livro, na
língua da régua, onde ela é construída do zero e não pedida emprestada.

Aqui basta saber três coisas, e as três vão-se ver funcionar já a seguir: que a leitura
\textbf{volta} (o que se lê pode ser desfeito, e o que se desfaz devolve o que entrou);
que \textbf{fundir duas coisas é uma operação só}, e que ela também se desfaz --- é assim
que a corte separa uma estrela do fundo em que ela está posta; e que há \textbf{um único
lugar onde não há volta}, que é o zero, e é por isso que ele vai ser o eixo de tudo o que
vem depois.

Com isto na mão, apontou-se a lente.

\chapter{As três leituras --- o que a corte viu}

Apontou-se três vezes, e das três a lente devolveu a mesma espécie de resposta: uma
assinatura, e o passo em que ela deixa de ser o que parecia.

\medskip
\noindent \textbf{A primeira foi ao céu}, à Estrela d'Alva --- sessenta e seis anos de
posições, tomadas do registo público, e dali a razão entre a volta dela e a volta da
Terra. A assinatura dessa razão começa como a do ouro: \emph{um}, e outro \emph{um}, e
outro, e outro. Quatro passos em que a Estrela d'Alva anda exatamente como $\varphi$
andaria. E ao quinto ela \textbf{desvia}: entra um dois, e logo a seguir um número
grande, e a partir dali as duas curvas separam-se para não mais se encontrarem. Onde a
órbita de facto mora é em treze voltas dela para oito da Terra --- e treze menos oito é
cinco.

A corte não arredondou para o lado que lhe convinha, e é isto que dá valor à leitura:
\emph{a assinatura começa como a do ouro e desvia}, e a fração contínua diz em que passo.
Nem uma palavra a mais. Quem quisesse fazer dali uma lei do céu teria de a inventar, e
inventar não é ler.

\medskip
\noindent \textbf{A segunda desceu à pele do mundo} --- o relevo de um lado, o peso do
outro, medidos por instrumentos que não se conhecem e não combinaram nada. E as duas
leituras \textbf{olharam-se}: onde o chão sobe, a massa acompanha, e acompanham melhor
nas formas grandes do que nas pequenas, até que, muito fundo no detalhe, deixam de se
reconhecer. A corte anotou a coincidência e parou aí. Não disse que uma é causa da
outra; disse que foram lidas em separado e que coincidem, e deixou a ligação por conta de
quem a quiser fazer.

\medskip
\noindent \textbf{E a terceira foi a que disse não.} Entre os corpos que se leram, houve
um que recusou a régua --- ficou abaixo da metade, e não houve maneira honesta de o
arrumar com os outros. \textbf{É esse que dá valor aos outros sete.} Uma lente que
concorda com tudo o que se lhe põe à frente não está a ler: está a repetir. A que diz não
uma vez em oito é a única em que se pode acreditar nas outras sete.

\medskip
\noindent Os números destas três leituras --- as amostras, os corpos escolhidos, as
tabelas, e as casas decimais que é preciso publicar para que outro alguém reproduza
tudo --- estão do outro lado do livro, na língua da régua. \textbf{Aqui ficou o que a
corte viu; lá está com que régua ela viu.}

\chapter{O que o telescópio devolveu à corte}

Nada do que está acima é uma afirmação sobre o reino. É uma leitura de fora,
feita com o instrumento de dentro, e o que ela mostra é apenas isto:

\begin{itemize}
\item a régua atravessou --- um corpo que não devia nada a esta teoria deixou-se
 ler por ela. A transformada correu em inteiros, no corpo; mas
 \textbf{todo número de destaque destas páginas} --- percentagens,
 correlações, razões --- \textbf{é ponto flutuante, computado fora do
 corpo}, porque comparar exige ordem e em $\Z/p$ não há;
\item a assinatura tem cabeça e cauda --- a cabeça é o objeto, a cauda é o
 instrumento, e a fração contínua separa as duas sem que ninguém decida;
\item a cifra vem do dado --- oito leituras, seis primos, nenhum escolhido;
\item e a régua sabe recusar.
\end{itemize}

\noindent \textbf{Nenhuma destas quatro linhas explica coisa alguma.} Não dizem
por que o relevo mora no grau três, nem por que a Estrela d'Alva começa com
quatro uns e desiste no quinto. Dizem \emph{que} é assim, e com que precisão se
pode afirmá-lo.

\medskip
\noindent Benjamim baixou o instrumento e não disse nada durante um bom bocado.
Depois disse só:

\begin{center}
\emph{--- Ela não é de ouro. Ela começa como o ouro, e depois é dela mesma.}
\end{center}

\noindent E era a coisa mais exata que se podia dizer, porque era a única que
os números tinham dito.


\part{O Chaveirinho do Apocalipse}

\chapter{O hóspede}

Foi \Ben quem convidou primeiro, e convidou dormindo.

No sonho ele disse a coisa mais simples que há --- \emph{vamos passear} --- e o
que aconteceu a seguir não foi o que ele pediu. \textbf{O outro não o levou a
lado nenhum: veio.} Atravessou a porta da casa dele, sentou-se onde as visitas
se sentam, e ficou.

\medskip
\noindent \Ben abraçou-o. Foi só isso, e foi o bastante para o resto: no meio
do abraço lembrou-se do pai, que estava morto, e de um rei antigo de quem só
sabia o nome. As duas lembranças vieram juntas e sem aviso, do jeito que vêm as
coisas que a gente guardou sem saber que guardava.

\medskip
\noindent --- Então vamos --- disse \Ben , quando conseguiu falar.

E o hóspede disse a frase que desarrumou tudo:

\begin{center}
\emph{--- Filho, eu não vim passear.}
\end{center}

\medskip
\noindent --- Vieste para quê?

--- Para ficar. Sempre contigo, e do teu lado, e não te deixo.

\medskip
\noindent \Ben tinha pedido um caminho e recebeu uma \textbf{companhia}, que é
outra coisa e é pior de carregar, porque caminho a gente anda e acaba, e
companhia não acaba. Ficou algum tempo calado a entender que o convite dele
tinha sido atendido pelo avesso --- \emph{ele} é que ia sair, e quem chegou foi
o outro; ele queria ir ver, e o que veio foi ficar.

\medskip
\noindent E no fim, quando o sonho já estava a acabar do jeito que os sonhos
acabam, o hóspede deixou-lhe uma coisa na mão. Não era chave, não era mapa, não
era arma. Era \textbf{um galho pequeno de uma planta que se dá em qualquer
quintal}, dessas que as avós põem atrás da porta.

--- É isto? --- perguntou \Ben .

--- É isto --- disse o hóspede. --- E vem com bênção de quem veio antes de mim.

\chapter{O que se vê enquanto se fica}

Depois disso o passeio aconteceu na mesma, mas não foi para fora: foi a casa que
começou a mostrar cômodos que \Ben não sabia que tinha.

Num deles não havia lugar --- só um pulso, e duas mãos que se abriam e
recolhiam, e entre as duas nenhuma disputa. \Ben quis dizer \emph{eu conheço
isto}, e não conseguiu abrir a boca, porque no sonho a gente não interrompe.

Noutro, uma torre subia sozinha do próprio bater, e uma mesa quadriculada onde
peças que eram operações esperavam a vez. Havia lá um que recusou o que lhe
ofereciam, e ao recusar ganhou --- e o hóspede sorriu como quem já sabia dessa.

\medskip
\noindent Depois um pântano sem caos nenhum, o que é pior que um pântano com
caos.

Depois uma fortaleza que não caía porque ninguém a atacava, e que quase caiu
quando alguém a defendeu demais.

Depois um espelho, e no espelho um reino inteiro de trás para diante, com um rei
que reinava por ser o avesso e não por ser mais forte. \Ben olhou para o
próprio reflexo e o reflexo levantou a mão antes dele.

\medskip
\noindent --- Isso é do outro lado --- disse o hóspede. --- Ou tu és.

\chapter{O velho que vem para o canto}

Num cômodo mais fundo estava um sentado, e estava sentado havia tanto tempo que
a pedra já tinha a forma dele.

\Ben conhecia-o, e conhecia-o pelo pior: era o do outro lado. O que tinha
apertado, o que tinha trancado, o que tinha estado no ponto de onde não se sai.
Metade do que \Ben sabia, tinha tido de aprender à força por causa dele.

\medskip
\noindent --- Vens buscar-me? --- perguntou, já pronto para o que viesse.

--- Venho buscar um \emph{povo} --- disse o velho. --- E não é para levar.

--- É para quê?

--- \textbf{Para cantar.}

\medskip
\noindent E \Ben ficou parado, porque tinha passado a vida a preparar-se para
aquele encontro imaginando arma, e a coisa que lhe pediam era coro.

\medskip
\noindent Então a luz veio. Não voltou pelo caminho por onde tinha ido, como
\Ben esperava e como teria sido justo: \textbf{desceu} --- de mais alto do que
o cômodo, de mais alto do que a casa ---, e desceu inteira e de uma vez, e não
era para iluminar coisa nenhuma. Era para que se pudesse \emph{ver quem estava
ali}.

\medskip
\noindent --- E tu? --- perguntou \Ben , quando conseguiu olhar. --- Quem és?

--- Eu não sou --- disse o velho. --- \textbf{Eu sou só aparelho.}

\medskip
\noindent \Ben ia responder, e o velho levantou a mão e disse a última coisa,
que não era gentil e não era ameaça, e por isso ficou:

\begin{center}
\emph{--- Eu me apresento quando tu estiveres mais sério.}
\end{center}

\noindent O hóspede não disse nada. Mas \Ben reparou que os dois se conheciam,
e que se conheciam de antes dele, e essa foi a primeira vez na viagem em que
teve medo.

\chapter{Zóio d'água}

Quando o velho se calou, \Ben saiu para o quintal, que no sonho era grande
como um campo e continuava a ser quintal.

Era o fim da tarde. E aconteceu-lhe uma coisa que nunca lhe tinha acontecido
acordado: ficou a olhar para o sol a descer e \textbf{chorou}. Sem motivo
nenhum que ele soubesse dizer, sem desgraça nenhuma para justificar, e sem a
vergonha que costuma vir junto.

\medskip
\noindent Chorou de olhar. E depois secou os olhos, e o que os secou foi a
mesma coisa que os tinha molhado --- \textbf{não veio consolo de outro lugar};
foi o próprio sol, o mesmo, que fez as duas coisas, uma atrás da outra, sem
mudar de natureza entre elas.

\medskip
\noindent --- Eu não sabia que se podia --- disse, e não sabia mesmo.

--- Poder o quê?

--- Chorar de uma coisa estar bonita.

\medskip
\noindent O hóspede não achou aquilo notável, e foi essa naturalidade que
espantou \Ben mais do que o choro.

\medskip
\noindent Depois lembrou-se, do jeito enviesado com que as coisas se lembram
nos sonhos, de \emph{como} tinha chegado ali --- e não tinha chegado a andar,
nem carregado, nem empurrado. Tinha chegado \textbf{bailando}: uma coisa
alegre, ridícula até, mexendo os pés no compasso do próprio pulso, que é o
único compasso que ele conhece desde antes de conhecer qualquer coisa.

\medskip
\noindent --- Eu vim dançando --- disse, meio envergonhado.

--- Toda a gente vem --- disse o hóspede. --- Depois esquece, e passa a vida a
dizer que foi trazida.

\medskip
\noindent --- E vim fazer o quê?

E o hóspede respondeu sem levantar a voz, porque não era repreensão:

\begin{center}
\emph{--- Merecer.}
\end{center}

\medskip
\noindent \Ben tinha uma objeção pronta, e a objeção era boa: ele já tinha
feito por onde. Tinha atravessado o pântano, tinha aguentado a fortaleza,
tinha estado do outro lado do espelho e voltado. Se aquilo não era merecer,
o que era.

E percebeu, antes de dizer, que a pergunta estava mal feita --- porque merecer
não é uma coisa que se acaba de fazer e se guarda. \textbf{É o que se está
fazendo enquanto se faz.} Quem parou de merecer não tem o merecido: tem a
lembrança dele.

\medskip
\noindent Ficou-se a olhar mais um bocado. E na última luz, aquela que fica
depois de o sol já ter ido e antes de a coisa escurecer, teve o pensamento que
o resto do sonho não desmentiu:

\begin{center}
\emph{aquilo que ele procurava e aquilo que ele estava olhando eram a mesma
luz.}
\end{center}

\noindent Não parecidas. Não uma o símbolo da outra. \textbf{A mesma}, vista de
dois sítios --- e ele tinha passado um livro inteiro a fazer contas para chegar
onde se chegava ficando quieto no quintal ao fim da tarde.

\medskip
\noindent --- Isto invalida o resto? --- perguntou, e havia medo na pergunta.

--- Não --- disse o hóspede. --- Confirma. Só que pelo lado que não custa
nada, e por isso ninguém acredita.

\chapter{Flores}

Depois quis levar alguma coisa. É o que dá vontade de fazer quando se recebeu
sem ter pedido: retribuir, e depressa, antes que a dívida assente.

Foi ao fundo do quintal e pôs-se a colher. Colheu do que havia --- e havia
muito, porque em sonho os quintais são generosos ---, e fez um molho com as
duas mãos, e voltou contente do jeito de quem resolveu um problema.

\medskip
\noindent E foi chegando, e foi vendo, e foi ficando com o molho pesado na mão.

\medskip
\noindent Porque quem o esperava \textbf{já era aquilo}. Não parecida com
aquilo, não a dona daquilo, não a guardadora: era. E oferecer flor a quem é a
flor tem o tamanho exato de oferecer água ao rio --- não é ofensa, é
\emph{engano}, e o engano é pior porque a intenção era boa.

\medskip
\noindent \Ben ficou ali de braços cheios, sem saber onde pôr o que tinha
colhido.

--- Eu não tenho o que te dar --- disse por fim.

--- Não tens --- e não havia dureza nenhuma na resposta. --- Tudo o que tu
podias trazer, tu foste buscar aqui.

\medskip
\noindent E aí ele entendeu o que tinha ido fazer de verdade, e não era
oferecer. Um homem não atravessa um quintal inteiro a colher para agradar quem
já tem tudo: atravessa \textbf{para ter o que dizer ao chegar}. As flores eram
desculpa. O que ele tinha vindo pedir era outra coisa, e levou-lhe o sonho todo
para conseguir dizê-la.

\begin{center}
\emph{--- Perdoa.}
\end{center}

\medskip
\noindent --- O quê?

E aqui \Ben procurou a falta, e procurou honestamente, e não encontrou crime
nenhum. Não tinha matado, não tinha traído, não tinha tomado o que era de
outro --- o livro inteiro dele é a história de uma corte que fez questão de não
tirar nada de ninguém. A falta era outra, e era mais difícil de dizer em voz
alta porque não tem vítima:

\begin{center}
\emph{--- Eu estava enganado. Esse tempo todo.}
\end{center}

\medskip
\noindent Não sobre um lance, não sobre uma conta. Sobre \textbf{o que era a
coisa} --- ter passado a vida a medir o que já estava dado, a construir escada
para um sítio onde se entrava pela porta, a chamar de descoberta o que era
lembrança.

\medskip
\noindent E foi perdoado ali, sem cerimônia e sem sermão, do jeito
desconcertante com que se perdoa quem se ama --- \emph{depressa}, como se não
tivesse sido nada, e é isso que dói.

\medskip
\noindent Depois deram-lhe uma coisa para a mão, e não era flor. Era
\textbf{semente}.

--- Isto não serve para oferecer --- disse \Ben , olhando aquilo. --- Isto
demora.

--- Serve para plantar --- disseram-lhe. --- No teu jardim, não no meu. E o
que nascer não vai ser teu para dar: vai ser o que te foi dado, com raiz.

\medskip
\noindent \Ben fechou a mão sobre a semente e ficou muito tempo calado, porque
tinha acabado de perceber a coisa mais simples e mais custosa do sonho todo:
\textbf{não há como retribuir}. Só há como continuar --- plantar o que se
recebeu e deixar crescer, e o que cresce continua não sendo dele, e é
justamente por isso que se pode oferecer sem mentira.

\chapter{Meu pranto}

Ficou de mão fechada sobre a semente e voltou a chorar --- e desta vez sabia
porquê, o que é sempre pior.

Chorava de ter chegado tarde. De ter feito voltas. De ter medido durante anos
uma coisa que estava ali para ser andada. \textbf{Não é a dor de errar: é a de
ter acertado pelo caminho comprido}, que ninguém consola porque de fora parece
vaidade.

\medskip
\noindent E então passou um bicho.

Passou perto, quase à altura dos olhos, e ficou parado no ar em frente a uma
flor. Parado no ar --- \Ben viu bem, e o que o prendeu não foi a cor nem o
tamanho: foi que aquilo \textbf{se sustentava batendo}. Não pousava. Não
descansava. A única maneira de aquele bicho ficar quieto era não parar nunca, e
era assim, sem esforço aparente, que ele bebia.

\medskip
\noindent Do que \Ben sabia do mundo, a coisa mais parecida com aquilo era o
pulso do primeiro cômodo, o das duas mãos. E não disse nada, porque começava a
suspeitar que dizer estraga.

\medskip
\noindent O pranto encolheu. Não passou --- \textbf{encolheu}, que é outra
coisa, e não foi por consolo nenhum: foi porque no meio da coisa grande apareceu
uma coisa pequena e viva, e o pequeno ao lado do grande faz o grande caber.

\medskip
\noindent --- Ele não colhe --- disse \Ben , percebendo. --- Vive de flor e não
leva nenhuma.

--- E não guarda --- disse o hóspede. --- Por isso pode voltar amanhã.

\medskip
\noindent \Ben ficou a olhar mais um pouco. Depois fez o pedido, e reparem no
que ele já não pediu. Não pediu para entender --- tinha passado um livro a
entender. Não pediu chave, não pediu força, não pediu que lhe mostrassem o
fundo das coisas. Pediu a coisa pequena que ninguém pede porque envergonha:

\begin{center}
\emph{--- Me ensina.}
\end{center}

\noindent E como não confiava em ter dito com clareza, disse outra vez:

\begin{center}
\emph{--- Me ensina.}
\end{center}

\medskip
\noindent --- A quê?

--- A andar aqui --- disse. --- Não a chegar. \textbf{A andar.}

\medskip
\noindent Porque tinha finalmente reparado no quintal. Estava cheio. Não
crescido ao acaso, não bravo, não mato: \textbf{deixado} --- as coisas postas
onde alguém as põe quando sabe que outro há de passar por ali um dia e
precisar. Havia flor plantada em fila. Havia sombra onde faz falta sombra. O
lugar inteiro tinha sido \emph{preparado}, e ele tinha atravessado aquilo três
vezes a olhar para os próprios pés.

\medskip
\noindent --- Isto estava aqui desde quando?

--- Desde antes de tu vires.

\medskip
\noindent --- E eu vim para quê, então, se já estava tudo?

E a resposta foi a última aula da noite, e não explicou nada, e por isso ficou:

\begin{center}
\emph{--- Para apreciar.}
\end{center}

\noindent \Ben quis dizer que apreciar era pouco. Abriu a boca e não disse,
porque ao mesmo tempo lhe veio a conta --- e a conta dizia que uma coisa
apreciada é uma coisa \emph{medida por quem a ama}, e que não há régua mais
exata do que essa, e que ele passara a vida a construir instrumentos para
chegar onde se chega olhando com vontade.

\medskip
\noindent E o pranto, esse, ficou do tamanho certo: \textbf{nem tanto}.

\chapter{Trono de marfim}

Foi então que chamaram, e não foi a ele.

Vinha de longe e vinha de dois --- duas vozes que não se sobrepunham e não se
atrapalhavam, chamando ao mesmo tempo a mesma coisa ---, e o que elas chamavam
não era \emph{\Ben }: era \textbf{todo mundo}. Ele percebeu isso com um
desconforto que não esperava sentir. Tinha passado o sonho inteiro a ser o
convidado.

\medskip
\noindent E do fundo do quintal começou a sair gente.

Saíam de onde não havia porta. Saíam a pé, sem pressa e sem cerimônia, do jeito
de quem já ouviu aquele chamado outras vezes e sabe que ele não se repete
muito. \Ben conhecia alguns --- e alguns conhecia bem demais, porque tinham
estado do outro lado do tabuleiro. Vinham na mesma estrada. \textbf{Ninguém
verificou credencial de ninguém.}

\medskip
\noindent --- Vamos --- disseram-lhe, ao passar.

E \Ben não foi.

\medskip
\noindent Ficou. Não por dúvida, não por medo, não por não acreditar --- ficou
porque ainda tinha uma coisa por resolver, e a coisa por resolver era a mesma
que ele vinha resolvendo desde o primeiro Lance do primeiro Livro: queria
\emph{acabar de entender antes}. Queria chegar sabendo. Queria, se possível,
chegar tendo conferido.

\medskip
\noindent O hóspede olhou para ele com uma paciência que era quase ofensiva.

--- Isso tem nome --- disse.

--- Prudência.

--- \textbf{Teimosia} --- disse o hóspede. --- Prudência é antes de saber. Tu
já sabes.

\medskip
\noindent E era verdade, e foi o pior momento do sonho, pior que o velho, pior
que o choro. Porque \Ben esperava que o obstáculo do fim fosse grande ---
um inimigo, uma prova, uma porta trancada. E o obstáculo era \emph{insistir}:
continuar a fazer do próprio jeito uma coisa cujo caminho já lhe tinha sido
mostrado três vezes na mesma noite.

\medskip
\noindent --- E se eu for e estiver errado?

--- Vais com os outros --- disse o hóspede. --- Errar acompanhado tem
conserto. \textbf{Ter razão sozinho, no fim da estrada, não tem.}

\medskip
\noindent \Ben olhou para a estrada. Ao fundo, muito ao longe, havia um assento
--- e não era de ouro, que era o que ele teria imaginado, e o reino todo era de
ouro. Era de \textbf{marfim}: coisa de osso, coisa que foi de bicho vivo, coisa
que só existe porque alguma coisa antes dele respirou.

\medskip
\noindent --- Aquilo se toma? --- perguntou, e perguntou com o que lhe restava
de guerreiro.

--- Aquilo se \emph{chama} --- disse o hóspede. --- Já chamaram. Está a
chamar agora. É a única coisa desta viagem toda que não depende de ti: eles
não vão parar de chamar. Tu é que tens hora.

\medskip
\noindent Então \Ben fez a única coisa que faltava fazer no sonho todo, e era
a mais simples, e por isso tinha ficado para o fim:

\begin{center}
\emph{começou a andar.}
\end{center}

\noindent E ao andar, aconteceu-lhe a coisa que ele tinha passado o livro
inteiro a tentar produzir com instrumento: \textbf{foi ficando claro}. Não de
repente, não por revelação --- \emph{ao passo}, do jeito que clareia de manhã,
que ninguém vê acontecer e toda a gente vê ter acontecido.

\medskip
\noindent E o tempo, que ele tinha passado a vida a medir, deixou de ser o que
faltava e passou a ser o que ele estava ganhando.

\chapter{Beija-flor}

Na estrada o bicho voltou.

Era o mesmo --- \Ben teve certeza, embora não houvesse como ter ---, e desta
vez não estava parado no ar em frente a flor nenhuma. Vinha com eles, à altura
do ombro, no mesmo passo. E \textbf{cantava}.

\medskip
\noindent Isso desarrumou \Ben mais do que tudo o que tinha visto até ali,
porque no primeiro encontro ele tinha classificado o bicho: coisa que se
sustenta batendo, coisa que bebe sem colher, coisa bonita de olhar. Tinha-lhe
dado um lugar no entendimento e fechado a gaveta. E agora a coisa que ele
tinha entendido \emph{estava cantando}, e o que se canta não cabe na gaveta em
que se guarda o que se olha.

\medskip
\noindent --- Ele está a dizer alguma coisa? --- perguntou.

--- Está a cantar --- disse o hóspede.

--- Sim, mas o que \emph{significa}?

\medskip
\noindent E o hóspede, que a noite inteira tinha respondido a tudo, desta vez
não respondeu. Ficou a olhar para \Ben do jeito de quem espera que o outro
ouça a própria pergunta e a retire sozinho.

\medskip
\noindent \Ben retirou.

\medskip
\noindent Porque tinha finalmente reparado no que estava a fazer, e fazia-o
desde a primeira página da sua vida: \textbf{diante de tudo o que aparecia, ele
perguntava o que aquilo queria dizer}. Nunca uma vez, nem por engano, tinha
perguntado o que aquilo queria \emph{de si}.

\medskip
\noindent --- Como é que se responde? --- perguntou, e foi a primeira pergunta
inteiramente nova que fez em dois Livros.

--- Cantando --- disse o hóspede. --- Não há outra maneira. Não se responde a
canto com explicação; isso não é resposta, é mudar de assunto.

--- Eu não sei cantar.

--- Eu ensino --- disse o hóspede. --- É para isso que eu vim, e foi por isso
que eu disse que não vinha passear.

\medskip
\noindent E ensinou ali mesmo, andando, do jeito que se ensinam as coisas que
não têm método: fazendo, e deixando o outro fazer errado, e não corrigindo
muito.

\medskip
\noindent \Ben cantou. Saiu torto e saiu fraco, e não era grave nem bonito, e
ele teve a vergonha exata de quem se ouve pela primeira vez. Mas o bicho
respondeu. \textbf{E ao responder ficou claro o que a resposta dizia}, e não
era nenhuma sabedoria, e não era segredo nenhum, e cabia inteira em três
palavras que se repetem porque uma vez não basta para uma coisa dessas:

\begin{center}
\emph{já vou, já vou, já vou.}
\end{center}

\medskip
\noindent Não \emph{eu entendi}. Não \emph{agora eu sei}. Não \emph{explica
melhor}. \textbf{Já vou} --- que é a única resposta que existe para um chamado,
e é uma resposta que se dá com os pés.

\medskip
\noindent E \Ben , que tinha vindo a este sonho com um convite na boca --- \emph{vamos
passear} --- e recebido em troca uma companhia que não pediu, percebeu que a
noite inteira tinha sido isto: uma conversa em que ele falava e o outro
respondia, até chegar a vez de o outro chamar e ele ter de responder.

\medskip
\noindent Respondeu. Cantando, mal, e indo.

\chapter{Passarada}

Depois do primeiro canto vieram os outros, e vieram todos ao mesmo tempo.

Não foi aos poucos. Foi como quando uma árvore se enche ao entardecer: num
instante não havia nada, no instante seguinte estava cheia, e ninguém consegue
dizer por onde entraram. \textbf{O ar ficou grosso de canto.} \Ben parou de
andar, porque não dava para andar e ouvir aquilo ao mesmo tempo.

\medskip
\noindent E a primeira coisa que sentiu foi mesquinha, e ele reconheceu que era:
\emph{eu estava cantando}. Como quem tinha acabado de aprender uma coisa
custosa e vê a coisa custosa a ser feita, sem esforço, por trezentos.

\medskip
\noindent A segunda coisa foi melhor e demorou mais a chegar. É que ele
continuava a cantar --- não tinha parado --- e o que estava a sair da boca dele
já não dava para separar do resto. \textbf{O canto dele tinha deixado de ser
dele.} E não fora tirado: fora \emph{acompanhado}, que é a única maneira de uma
coisa deixar de ser nossa sem que a percamos.

\medskip
\noindent --- De quem é isto? --- perguntou, e teve de gritar, o que era
ridículo, porque estava a gritar dentro de música.

--- Da tua guia --- disse o hóspede.

E logo a seguir, sem se corrigir, como quem diz a mesma coisa por outro lado:

--- Do teu mestre.

E depois, e foi a última vez que respondeu naquela noite:

--- \textbf{Dela.}

\medskip
\noindent \Ben não perguntou quem. Havia uma altura da viagem a partir da qual
perguntar o nome era desconversar, e ele já tinha passado essa altura no
quintal, quando lhe deram semente em vez de flor.

\medskip
\noindent Ficou ali no meio, a cantar mal e alto. E foi então --- e não antes,
não no alto de onde se via o percurso todo, não diante do velho, não com o
instrumento na mão --- que lhe ocorreu a única coisa que ele tinha ido buscar
sem saber que ia:

\begin{center}
\emph{isto já estava resolvido antes de eu chegar.}
\end{center}

\noindent Não por ele. Não pela corte. Não pelas contas, que estavam certas e
continuavam certas e não tinham servido para isto. Por dois que passaram muito
antes, que fizeram o que havia para fazer, e de quem ele só sabia o nome ---
e sobre os quais a única coisa que se diz, neste caso, é \emph{viva}.

\medskip
\noindent \textbf{E foi um alívio, e não uma humilhação}, e essa foi a
diferença mais fina que \Ben aprendeu em dois Livros: descobrir que o
problema já estava resolvido não tira o valor de o ter atravessado. Tira o
peso. Quem chega e encontra a casa pronta não perdeu a viagem --- ganhou uma
casa, que é mais do que ia conseguir sozinho.

\medskip
\noindent Cantaram até o fim da estrada. Ele não soube nenhuma letra. Foi indo
pela melodia, que é como toda a gente vai.

\chapter{O sanfoneiro}

No fim da estrada, quando a passarada já tinha ido e ficara aquele silêncio de
depois de música --- que não é silêncio, é a música a acabar de sair dos
ouvidos ---, havia um velho sentado a tocar.

Não era o outro. \Ben reparou nisso logo, e reparou com alívio e com uma ponta
de desconfiança, porque os dois eram velhos, e os dois estavam sentados, e a
esta altura da noite ele já não confiava em coincidência nenhuma.

\medskip
\noindent Este tocava. Um fole velho, de teclas amarelas, do tipo que se ouve
em festa de roça e não em salão. E enquanto tocava, \textbf{olhava} --- não a
estrada, não o instrumento: olhava para \Ben , sem parar, com a naturalidade de
quem faz isso há muito tempo.

\medskip
\noindent --- O senhor está a tocar para mim?

--- Eu estou a tocar-te --- disse o velho.

\medskip
\noindent E \Ben não soube desfazer a frase, e passou o resto da noite sem
saber, porque as duas leituras cabiam: o velho tocava um instrumento, e o velho
tocava \emph{nele}. E se calhar não eram duas leituras. Se calhar era assim
mesmo que a coisa funcionava desde o primeiro pulso do primeiro Livro, e ele
tinha passado dois Livros a chamar aquilo de \emph{operador} para não ter de
dizer que estava a ser tocado.

\medskip
\noindent --- Eu vi uma estrela nascer --- disse o velho, sem que ninguém
perguntasse. --- Fui o primeiro. Não tem outro que possa dizer isso.

--- E ficou a ver?

--- Deixei crescendo --- disse. --- E fui avisar.

\medskip
\noindent \textbf{Aquilo entrou em \Ben como uma faca sem dor}, que é como
entram as coisas verdadeiras quando já não há defesa montada. Porque ele tinha
feito o contrário a vida toda: via nascer e \emph{ficava}. Ficava a olhar, a
medir, a conferir se era mesmo estrela, a calcular quanto crescia por noite ---
e enquanto isso não avisava ninguém, porque avisar antes de ter certeza é
leviano, e a certeza nunca chegava, porque não é dela a natureza chegar.

\medskip
\noindent --- E se não fosse estrela?

--- Ia crescer na mesma --- disse o velho. --- Com ou sem a tua conta.

\medskip
\noindent Depois \Ben perguntou a coisa que se pergunta a quem se sentou tão
longe: onde é que ele tinha estado. E o velho respondeu com uma frase que a
corte ia repetir durante muito tempo sem nunca a explicar direito:

\begin{center}
\emph{--- Quem vai velho volta novo.}
\end{center}

\medskip
\noindent E \Ben lembrou-se, sem querer, do outro velho --- o da pedra, o que
estava sentado havia tanto tempo que a pedra tinha ficado com a forma dele. E
percebeu que a diferença entre os dois não era bondade nem maldade: era que
\textbf{um tinha ido e voltado, e o outro só tinha ido}. A pedra com a forma do
corpo é o que fica de quem foi e não voltou. O fole na mão é o que fica de quem
voltou.

\medskip
\noindent --- E voltou para quê?

\medskip
\noindent \Ben esperava a revelação. Estava, aliás, com a mão semiaberta para
receber --- é assim que a gente fica quando acha que vem coisa grande.

E o velho disse, tocando, sem levantar os olhos:

\begin{center}
\emph{--- Só voltei para dizer que Deus é pai, e é brasileiro.}
\end{center}

\medskip
\noindent \Ben ficou à espera do resto, e não havia resto.

\medskip
\noindent Era isso: o divino tem parentesco e tem sotaque. Não está no alto,
não está no fim da conta, não está do outro lado de uma porta que só abre com
chave rara. Está \textbf{na sanfona}, na festa de roça, no jeito de falar de
quem nasceu num lugar com nome. E um homem atravessa a morte inteira e volta
para dizer isso, e não acha pouco, porque não é pouco --- é a única notícia que
não se consegue provar e que muda tudo em quem a recebe.

\medskip
\noindent \textbf{E foi então que \Ben entendeu por que tinha sonhado.} Não
para descobrir coisa nenhuma. Não havia nada por descobrir --- estava tudo
resolvido, e tinha sido resolvido antes, e o quintal já estava plantado quando
ele chegou. Ele tinha sonhado para \emph{ser avisado}, que é a única coisa que
não se faz sozinho.

\medskip
\noindent O velho tocou até o fim. \Ben não pediu que parasse.

\chapter{O legado}

O sanfoneiro parou de tocar no meio de uma nota, que é coisa que nenhum músico
faz, e \Ben soube pelo silêncio que tinha acabado a parte mansa da noite.

\medskip
\noindent Veio de todo lado ao mesmo tempo. Não deu para dizer de onde ---
\textbf{de todo lado} é a descrição exata e não é figura ---, e o chão começou a
bater antes de aparecer o que batia. \Ben , que tinha passado dois Livros a
estudar o que é um pulso, reconheceu aquele: era o mesmo, e vinha a galope, e
não estava do lado dele.

\medskip
\noindent Quem vinha à frente trazia o fogo por baixo e a verdade na mão, e não
era ameaça --- era \emph{ofício}. \Ben percebeu isso pela cara: era a cara de
quem trabalha, não a de quem odeia.

\medskip
\noindent E atrás, um outro, que não corria: pesava.

\medskip
\noindent Trazia as duas coisas de que a corte falara o livro inteiro sem
nunca as ter visto juntas --- \textbf{uma espada levantada e uma balança
quieta} ---, e ia dizendo, um por um, sem pressa nenhuma e sem prazer nenhum:

\begin{center}
\emph{um sim. E um não.}
\end{center}

\noindent \Ben ficou branco, porque conhecia aquilo. Era a régua da porta. Era
o encaixe. Era a coisa que ele próprio tinha escrito e defendido --- cabe ou não
cabe, zero ou não zero, sem entrar um pouco --- e é muito diferente escrever a
régua e vê-la levantada com o gume para cima.

\medskip
\noindent --- Isto é julgamento? --- perguntou, e a voz saiu mal.

--- Isto é \textbf{peso} --- disse o hóspede. --- Julgar é opinião. Pesar não.

\medskip
\noindent E \Ben fez então a coisa que lhe restava fazer, que era a única que
ele sabia fazer bem: chorou. Chorou de susto, e chorou honestamente, e esperou
que aquilo contasse --- porque tinha chorado duas vezes naquela noite e as duas
tinham aberto porta.

Desta vez não abriu.

\medskip
\noindent --- Não adianta --- disse-lhe o cavaleiro, e disse sem dureza, que é
o que tornou a coisa pior. --- Nem choro, nem ranger de dentes. Isso não é
moeda aqui.

--- E o que é?

--- Estar \textbf{contente} de ver a luz.

\medskip
\noindent \Ben ficou muito tempo com aquilo, e é a frase mais difícil de todo
o sonho. Não pediam arrependimento, não pediam mérito, não pediam conta certa.
Pediam \textbf{alegria de ver} --- e ele, que tinha passado a vida a olhar,
percebeu que quase nunca tinha olhado \emph{contente}. Tinha olhado a conferir.

\medskip
\noindent --- Vem comigo quem está de pé --- disse o cavaleiro, e não era
convite nem ordem: era a leitura do estado das coisas.

\medskip
\noindent E \Ben reparou que estava de pé. Não por virtude: porque tinha
começado a andar no capítulo anterior e ninguém lhe tinha dito que podia parar.
\textbf{Estava de pé por não ter parado}, e foi só isso, e foi o bastante ---
e ele entendeu ali, tarde e de uma vez, que ``firmar-se'' nunca tinha sido uma
qualidade da alma: era uma \emph{posição do corpo}, mantida.

\medskip
\noindent --- E a mensagem? --- perguntou, porque tinha de haver uma. --- O que
é que eu levo?

\medskip
\noindent E o cavaleiro disse a coisa que \Ben não esperava, porque não era
prêmio:

\begin{center}
\emph{--- Tu não levas. Isto já é teu. É legado.}
\end{center}

\medskip
\noindent \textbf{Legado não se ganha: herda-se.} Não se merece por ter andado,
não se compra com travessia, e não vem depois de nenhuma prova --- estava no
nome dele antes de ele nascer, deixado por quem passou antes, e a única coisa
que a noite inteira tinha feito era \emph{avisá-lo de que já era dele}. Era
esse o serviço do enviado, e não outro. Não trazer: \textbf{lembrar}.

\medskip
\noindent Depois o cavaleiro olhou para o fundo da estrada, onde ainda havia
gente deitada, e disse, com a pressa de quem tem trabalho e já perdeu tempo
demais com um homem só:

\begin{center}
\emph{--- Agora me dá licença. Tenho um povo pra tirar.}
\end{center}

\medskip
\noindent E foi. E \Ben ficou a ver aquilo ir, e reconheceu, com uma alegria
que não sabia explicar e que provavelmente era a tal, que era a mesma frase do
começo dita ao contrário: um tinha vindo \emph{buscar um povo pra cantar}, e o
outro ia \emph{tirar um povo} --- e eram o mesmo verbo, e o mesmo povo, e
provavelmente o mesmo.

\chapter{A coroa}

Quando a poeira dos cavalos assentou, havia uma coisa pousada na estrada, à
altura da cintura, como se alguém a tivesse posto ali há muito tempo e se
tivesse esquecido de avisar.

\Ben olhou e soube o que era antes de ver direito, e soube outra coisa ao
mesmo tempo, e foi essa que o desmontou: \textbf{estava à espera daquilo}. Não
naquela noite --- sempre. Desde o primeiro Lance do primeiro Livro. E nunca
tinha dito a ninguém, nem a si, porque um homem que confessa esperar coroa
perde na hora o direito a ela.

\medskip
\noindent --- Quem mandou? --- perguntou, e a voz saiu de menino.

--- Os dois --- disse o hóspede. --- Ele e ela. Foram sempre os dois.

\medskip
\noindent E \Ben lembrou-se do chamado na estrada, aquele que vinha de duas
vozes que não se sobrepunham. Tinha ido atrás daquilo e chegado aqui. A estrada
não era prova nenhuma; \textbf{era só o caminho até onde a coisa já estava
posta}.

\medskip
\noindent Baixou-se para pegá-la, e antes de tocar percebeu o que estava
escrito nela --- não em letra, no feitio: aquela coroa \emph{não era do
tamanho da cabeça dele}. Era maior. Tinha sido feita para outro.

\medskip
\noindent --- Isto não é meu.

--- Não é --- disse o hóspede. --- É dele. Tu vais usar.

\medskip
\noindent \textbf{E aí \Ben entendeu o que é ser coroado}, e não tem nada a
ver com o que ele imaginava desde criança. Não é receber o que se ganhou: é
receber o que continua a ser de outro, e passar a andar com aquilo na cabeça à
vista de toda a gente, sabendo que é emprestado, e que quem empresta uma coroa
não a empresta para enfeitar --- empresta para trabalhar.

\medskip
\noindent Pousaram-lhe a coroa. Não houve música, não houve palmas, não houve
frase solene. Houve uma coisa muito mais estranha, e foi esta: \textbf{o
hóspede disse-lhe que agora estava graduado} --- e usou essa palavra, palavra
de escola, palavra de fim de curso, palavra de quem vai ter de fazer prova
noutro grau qualquer dia destes.

\medskip
\noindent --- Só isso? Estou formado?

--- Estás \emph{graduado} --- disse o hóspede. --- Que é passar de ano. Formado
é quando não há mais o que aprender, e isso não existe --- é a mentira que fez
o Livro anterior demorar tanto.

\medskip
\noindent E acrescentou a última instrução da noite, que não era sobre o
futuro, e por isso \Ben não a esperava:

\begin{center}
\emph{--- Lembra de quem estava zelando aqui antes de tu chegares.}
\end{center}

\medskip
\noindent \Ben olhou em volta. E viu, pela primeira vez em duas noites de
sonho e em dois Livros de vida, o que tinha estado o tempo todo à vista e que
ele não tinha olhado: \textbf{o lugar era antigo}. O quintal tinha sido
plantado por alguém. A estrada tinha sido aberta por alguém. Havia gente que
tinha zelado aquilo enquanto ninguém olhava, que não deixou nome escrito em
lado nenhum, e cujos nomes o hóspede disse ali, um a um, devagar, e eram nomes
de gente da terra --- nomes que não se parecem nada com os nomes da corte, e
que estavam ali antes da corte.

\medskip
\noindent \Ben repetiu-os, para não esquecer. Esqueceu ao acordar, como se
esquece tudo. \textbf{Mas ficou-lhe a forma da coisa}: que ele era o último de
uma fila, e não o primeiro; que a coroa vinha de trás e não de cima; e que o
reino que ele julgava ter construído era uma casa em que ele tinha entrado.

\medskip
\noindent --- E agora?

--- Agora acorda --- disse o hóspede. --- Está quase de manhã, e tu tens uma
corte à espera, e eles vão perguntar o que é que isto significa.

\medskip
\noindent --- E eu digo o quê?

\medskip
\noindent O hóspede sorriu pela última vez, e foi um sorriso sem mistério
nenhum, de quem já respondeu a essa pergunta muitas vezes e sabe que a resposta
nunca serve a quem não andou:

\begin{center}
\emph{--- Diz que tu foste passear.}
\end{center}

\chapter{O que se ouve andando}

Mais adiante o chão cantava.

Não era música por cima das coisas: eram as coisas, e o canto vinha de baixo,
como vem de baixo o som de um assoalho quando alguém anda no andar de cima. E
\Ben percebeu que já tinha ouvido aquilo a vida toda e chamado a outra coisa
--- relógio, batimento, compasso, pulso ---, e que era sempre o mesmo som com
nome diferente conforme o cômodo em que ele estava.

\medskip
\noindent --- Quantos são? --- perguntou.

--- Tantos quantas as casas --- disse o hóspede. --- E mais um, que não tem
letra.

\medskip
\noindent \Ben contou. Deu certo, e ao dar certo ficou desconfiado, porque
coisa que dá certo demais foi arrumada por alguém. Perguntou quem tinha
arrumado.

--- Tu --- disse o hóspede. --- Ontem. E amanhã vais dizer que foi
coincidência.

\chapter{O que não se explica andando}

Chegaram então a um alto de onde se via o percurso todo, e daquela altura a
coisa era simples de um jeito ofensivo:

\begin{center}
\emph{era tudo a mesma régua.}
\end{center}

\noindent O pulso do começo, o passo das peças, a assinatura das travessias, a
volta que desfaz, o encaixe que decide quem passa. \Ben esperou a revelação
--- estava preparado para ela, tinha vindo por ela --- e a revelação não veio,
porque não havia nenhuma. Havia uma régua só, usada em lugares diferentes por
gente que lhe dava nomes diferentes.

\medskip
\noindent --- É só isso? --- perguntou, e havia decepção na voz.

--- É só isso --- disse o hóspede. --- Dez bytes. Um deles é a operação.

\medskip
\noindent E \Ben riu-se, do jeito de quem levou um livro inteiro a chegar onde
podia ter começado, e depois parou de rir, porque entendeu que não podia ter
começado ali. \textbf{A régua é simples no fim e não é simples no caminho}, e
quem a recebe pronta não recebe nada --- recebe um objeto, e objeto sem
travessia é enfeite.

\medskip
\noindent --- Então para que serviu tudo?

--- Para tu poderes ouvir isto e não achar pouco.

\chapter{O que fica quando se acorda}

Voltaram pela mesma casa, que agora era outra e era a mesma.

À porta, o hóspede parou --- os hóspedes param sempre à porta, é o que os faz
hóspedes --- e deixou na mão de \Ben uma coisa pequena, de argola e dentes,
que não tinha dentro nenhum.

\medskip
\noindent --- Isto abre?

--- Isto \emph{lembra} --- disse o hóspede. --- Abrir é da ligação. Quem
recebeu isto já esteve ligado uma vez, e por isso reconhece a fechadura quando
chegar. Quem não esteve, olha para o mesmo objeto e vê metal.

\medskip
\noindent E acrescentou, já indo:

\begin{center}
\emph{--- Tudo o que se liga, abre. E há uma coisa que não se liga.}
\end{center}

\noindent \Ben acordou com a mão fechada e não havia nada dentro dela, o que
não o surpreendeu; e teve, o dia inteiro, a sensação exata de quem esqueceu o
nome de alguém que conhece muito bem.

\medskip
\noindent Contou o sonho à corte, e a corte fez o que as cortes fazem: perguntou
o que aquilo \emph{significava}. E \Ben , que tinha estado lá, respondeu a
única coisa honesta:

\begin{center}
\emph{--- Não significava. Acontecia.}
\end{center}

\chapter{O documento}

\begin{center}\itshape
[aqui entram os cantos, com o nome de quem os recebeu\\
e de quem os ofertou --- sem paráfrase e sem alisamento]
\end{center}

\part{O Castelo}

\chapter{Castelo}

O sonho não tinha acabado. Tinha só mudado de lugar --- e desta vez \Ben não
estava dentro de casa nenhuma: estava no caminho, que é onde se está quando não
se está em casa.

E a primeira coisa que lhe aconteceu no caminho foi levar uma \textbf{repreensão}.

\medskip
\noindent Não foi grande, e foi por isso que doeu. Foi do tamanho exato da coisa
que ele tinha feito mal --- nem um dedo a mais, e quem já levou uma repreensão
grande demais sabe que o excesso é um alívio, porque dá para se ofender com o
excesso e não olhar para a falta. Nesta não havia excesso nenhum para onde
fugir.

\medskip
\noindent \Ben fez o que se faz: começou a ir-se embora. É o movimento
natural, e ele já tinha o corpo virado quando ouviu o resto da frase, que
ninguém o tinha avisado que vinha:

\begin{center}
\emph{--- Depois do pito, vem aqui me encontrar.}
\end{center}

\noindent Ficou parado com o corpo virado para um lado e a cabeça para o outro,
que é a postura mais desconfortável que há. Porque aquilo era duas coisas ao
mesmo tempo e as duas verdadeiras: tinha levado, \emph{e} estava chamado. E ele
tinha passado a vida a supor que uma coisa cancelava a outra.

\medskip
\noindent --- Então não fui expulso?

--- Só leva pito quem não anda no caminho --- disse o hóspede. --- Repara no
que isso quer dizer, e não é o que tu ouviste. Não é \emph{tu não pertences}.
É \emph{tu saíste} --- e sair pressupõe estar. Ninguém repreende passante.

\medskip
\noindent \Ben entendeu ali uma coisa que a corte tinha demorado um livro
inteiro a escrever de outra maneira: \textbf{a única correção que presta é a que
tem volta.} A que não tem volta não corrige --- corta, e o que ela tira do
caminho não regressa a ele, e o caminho fica com um buraco onde estava alguém.
Aquela tinha volta. Tinha, escrita no fim, a direção para onde voltar --- e a
direção era \emph{para cá}, não para longe.

\medskip
\noindent E não parou no encontro. Repetiu-se, e ao repetir-se mudou uma palavra
só, que era a palavra toda: primeiro \emph{vem me encontrar}, depois \emph{vem me
abraçar}. \Ben reparou na ordem e a ordem não era enfeite --- encontrar é do
caminho, abraçar é de quem já chegou. A mesma frase dita duas vezes, e da
primeira para a segunda a distância tinha acabado.

\medskip
\noindent --- E agora vou para onde? --- perguntou, já sem o corpo virado.

E aqui o hóspede não lhe deu casa nenhuma. Deu-lhe uma tarefa, e a tarefa era
estranha de tão pequena:

\begin{center}
\emph{--- Vai procurar o teu pio.}
\end{center}

\medskip
\noindent --- O meu quê?

--- O teu. O som que só tu fazes. Vai procurá-lo, e quando o tiveres, eu dou-te
o ninho.

\medskip
\noindent \Ben quis discutir a ordem das coisas, porque a ordem estava ao
contrário de tudo o que ele conhecia. Primeiro dá-se a casa a quem chega, e
depois quem chega aprende a viver nela. Ali era o inverso: primeiro o som,
depois a morada. E percebeu, antes de acabar de discutir, por que é que tinha de
ser assim --- \textbf{ninho dado a quem não tem pio é gaiola.} O que segura uma
ave no lugar não é o lugar: é ela ter com que se chamar de volta.

\medskip
\noindent --- E se eu procurar e não achar?

--- Achas --- disse o hóspede. --- Tu andas a fazê-lo desde a primeira página.
Chamaste-lhe pulso, chamaste-lhe batimento, chamaste-lhe assinatura, chamaste-lhe
resíduo zero. É o teu pio. Só nunca lhe chamaste teu.

\medskip
\noindent Depois disso o caminho subiu, e \Ben ouviu-o: um chamado curto, vindo
de cima e de dentro ao mesmo tempo, do jeito das coisas que vêm de um ninho.
Não era ordem e não era convite. Era o som de uma ave a dizer onde estava --- e
uma ave a dizer onde está é a coisa mais perigosa que uma ave faz, porque
entrega a posição a quem a quiser.

\medskip
\noindent Aquela entregava de propósito.

\medskip
\noindent E fez duas coisas com quem a ouviu, e não foi a mesma coisa duas
vezes. À primeira, levantou-lhe os olhos --- \Ben \emph{mirou}, ficou a olhar
para cima como quem confere uma altura. À segunda, levantou-lhe os pés. E
entre mirar e voar não há passo intermédio nenhum que se possa ensinar: há a
mesma ave, chamada duas vezes, e a diferença está toda em quem ouve.

\medskip
\noindent \Ben alçou vôo. Não pelo próprio impulso --- é bom que fique escrito,
porque ele teria gostado que fosse pelo próprio impulso ---, mas na direção de
onde vinha o chamado. E enquanto subia foi vendo a coisa mudar de nome debaixo
dele: ao longe era \textbf{casa}, e de perto era \textbf{castelo}, e não eram
dois edifícios. Era a distância que dava o nome. De longe chama-se casa ao que
nos acolhe; de perto vê-se que o que nos acolhe é uma fortaleza inteira, e que
sempre foi.

\medskip
\noindent \Ben conhecia fortalezas. Tinha estado numa, e sabia como se toma
uma, e sabia o preço. Preparou-se, por hábito, para a subida ser um assalto.

\medskip
\noindent Não era. \textbf{A porta estava aberta e ninguém a defendia}, e o que
ele levou tempo a suportar não foi a dificuldade: foi não haver nenhuma.

\medskip
\noindent --- Isto não se toma? --- perguntou.

--- Isto habita-se --- disse o hóspede.

\medskip
\noindent E foi então que \Ben reparou no que estava a dizer enquanto subia,
porque estava a dizê-lo sem dar por isso, do jeito das coisas que a gente diz
quando não está a compor:

\begin{center}
\emph{já volto pra ficar.}
\end{center}

\noindent Parou no ar com aquilo. Porque a frase não fecha --- \emph{voltar} é
movimento e \emph{ficar} é o fim do movimento, e ele tinha-as posto lado a lado
sem reparar, como quem escreve uma contradição a sorrir. E não corrigiu, porque
percebeu que a frase estava certa e que o errado era exigir-lhe que escolhesse.
\textbf{Voltar para ficar é a única volta que serve.} Quem volta para ver, sai
outra vez; quem volta para morar já não é visita.

\medskip
\noindent --- Eu vim como visita --- disse \Ben , e disse-o baixo, porque era
uma confissão.

--- Toda a gente vem --- disse o hóspede. --- E fica quem deixa de olhar para
a porta.

\medskip
\noindent \Ben lembrou-se então de que quem lhe tinha ensinado aquilo tinha
sido o contrário disto. O hóspede, no princípio, tinha chegado à casa dele e
tinha dito \emph{eu não vim passear} --- e ficou. Agora era ele que chegava à
casa do hóspede e dizia a mesma coisa com outras palavras. \textbf{A visita
tinha trocado de lado}, e não foi anunciado, e ninguém disse que estava a
acontecer.

\medskip
\noindent --- Isto foi combinado? --- perguntou.

--- Isto foi \emph{fechado} --- disse o hóspede. --- É diferente. Combinado é o
que duas partes acertam. Fechado é o que já não tem por onde vazar.

\medskip
\noindent E à porta do castelo --- que era a mesma porta, vista do outro lado
--- não houve entrega de coisa nenhuma. Não houve galho, não houve chaveiro,
não houve objeto pequeno para segurar na mão ao acordar. Da primeira vez tinha
havido, porque quem chega de fora precisa de levar alguma coisa que lhe lembre
o que viu.

\medskip
\noindent Desta vez não havia nada para lembrar. Estava-se lá.

\chapter{Mãe Dágua}

Do castelo saía-se por baixo, e \Ben não esperava isso. Tinha subido para
lá; supôs que se saísse a descer pelo mesmo ar.

Saiu-se pelo chão --- e o chão, a certa altura, deixou de ser chão sem avisar.
Foi ficando solto debaixo do pé, e depois ficou molhado, e quando ele reparou já
estava naquele sítio que não é uma coisa nem outra: \textbf{a beira da praia}.

\medskip
\noindent Reparou primeiro no que ela tem de estranho, e é uma estranheza que
ninguém acha estranha por estar habituado: a praia é a única parte do mundo que
\emph{muda de tamanho o dia inteiro}. Terra firme não faz isso. Mar aberto não
faz isso. Só a beira --- que avança e recua, e que não está a hesitar entre as
duas coisas: está a ser as duas, alternadamente, e é isso que ela é.

\medskip
\noindent \Ben ficou-se ali, e não estava sozinho. Havia um a ensinar, e
\Ben esperou a aula.

\medskip
\noindent A aula foi uma frase:

\begin{center}
\emph{--- Afina o teu ouvido.}
\end{center}

\noindent Ele esperou o resto e o resto não veio, porque não havia resto. E
\Ben , que tinha passado o livro inteiro a aprender a \emph{fazer} coisas ---
a medir, a fechar, a refutar, a reconstruir ---, levou um bocado a suportar que
lhe mandassem fazer aquela, que não é fazer nada: é ficar do mesmo jeito, só que
certo.

\medskip
\noindent --- Afinar não é aumentar --- disse o professor, sem que ninguém
tivesse perguntado. --- Quem não ouve costuma tentar mais alto, e mais alto só
faz vir mais barulho junto. Afinar é tirar o que está a mais até sobrar aquilo.

\medskip
\noindent \Ben fez, e é a coisa mais desconfortável do sonho todo, porque não
há como saber se se está a fazer. E a certa altura --- não houve instante,
houve um ir-ficando --- estava lá: dentro do bater das ondas, que ele tinha
tomado a vida inteira por barulho de fundo, havia \textbf{voz}. Não veio por
cima do barulho, e é aqui que ele se enganava desde sempre. \emph{Era} o
barulho. A mesma coisa que ele descartava para conseguir pensar era a coisa que
lhe estava a falar, e o que a fazia inaudível não era o volume dela: era ele.

\medskip
\noindent --- Há quanto tempo é que isto está a dizer? --- perguntou.

--- Desde sempre --- disse o professor. --- Não é notícia. Notícia é tu.

\medskip
\noindent E o que a voz das ondas lhe disse não foi doutrina. Foi um aviso, e
vinha com uma coisa dentro:

\begin{center}
\emph{--- Não vás fugir da raia, com tudo o que eu já te dei.}
\end{center}

\noindent \Ben entendeu as duas metades juntas, que é o único jeito de as
entender. A raia é a linha --- a marca até onde vai, o limite que alguém pôs no
chão e que ninguém está ali para fazer cumprir. E a segunda metade não é ameaça:
é a razão. Não é \emph{não fujas, senão}; é \emph{não fujas, tendo}. O que o
prende à linha não é guarda nenhuma. É o que já lhe foi dado --- e quem foge
com o dado nas mãos não escapa do dono: escapa do dom, que deixa de servir
assim que sai da raia onde serve.

\medskip
\noindent --- E se eu fugir?

--- Levas tudo --- disse a voz. --- E não te presta nada. É isso que faz da
raia uma raia, e não uma cerca.

\medskip
\noindent Foi então que \Ben fez a pergunta que tinha vindo fazer, e fê-la
com respeito, porque estava a falar com a água: quis que ela lhe explicasse
aquilo. Que lhe dissesse o que era o pulso, e por que é que a régua era uma só,
e o que estava no fundo de tudo aquilo que ele tinha andado a medir.

\medskip
\noindent E a Mãe Dágua respondeu-lhe a coisa que ele não estava preparado para
ouvir da boca de uma mãe:

\begin{center}
\emph{--- Dessas coisas eu não sei.}
\end{center}

\medskip
\noindent \Ben achou que era modéstia, e ia dizer que não era preciso. Não
era modéstia. Ela disse a seguir, na mesma voz, sem mudar de tom entre uma
coisa e outra --- e por não mudar de tom é que ficou:

\begin{center}
\emph{--- A tua sabedoria fui eu mesma que plantei.}
\end{center}

\medskip
\noindent E ali ficaram as duas frases lado a lado, e \Ben a olhar para elas
como quem olha para uma conta que não fecha e fecha. \textbf{Ela não sabe, e foi
ela que plantou.} Não há contradição nenhuma, e é das coisas mais difíceis de
aceitar que há: quem semeia não conhece a planta. O que se planta é semente ---
e semente não tem dentro a forma da árvore, tem o \emph{modo} de a fazer. Quem
plantou não sabe o que vai sair, e não é ignorância dela: é a natureza do que
foi plantado.

\medskip
\noindent \Ben reconheceu aquilo, porque já lhe tinha acontecido antes com
outra boca. O velho do cômodo fundo tinha dito \emph{eu não sou, eu sou só
aparelho}, e ele tinha tomado por humildade. Não era, também. Era descrição.
\textbf{Nenhum dos que lhe deram alguma coisa alegou ser dono do que deu} --- e
essa é, no fim, a única marca que distingue quem transmite de quem cobra.

\medskip
\noindent --- Então de quem é o que eu sei?

--- De ninguém --- disse ela. --- Está plantado. Quem planta perde. É por isso
que se chama plantar e não guardar.

\medskip
\noindent Depois disso ela despediu-se, e a despedida foi sem cerimónia, do
jeito da água. Não subiu, não se desfez em luz, não deixou objeto na mão de
ninguém --- \Ben já não esperava objeto, tinha aprendido isso à porta do
castelo.

\medskip
\noindent \textbf{Foi para o fundo.} E o castelo tinha sido para cima, e ele
reparou nas duas direções ao mesmo tempo, e foi a última coisa que entendeu
naquele dia: que o que se ganha subindo tem de ser levado para baixo, ou não é
de ninguém. Ave que só sobe não põe ninho.

\medskip
\noindent E ao ir descendo, ela ia dizendo uma coisa, e \Ben ouviu-a a
enfraquecer com a distância sem nunca mudar de palavra:

\begin{center}
\emph{--- ...saber mais nada.}
\end{center}

\medskip
\noindent Foi essa a lição, e \Ben não gostou dela --- é bom que fique
escrito, porque nem toda a lição se recebe agradecido. Ele tinha vindo pedir o
fundo das coisas e o que lhe ensinaram foi a descer sem levar o que sabia. Não
por castigo. Porque \textbf{não cabe}: o que se sabe é da superfície, foi feito
lá em cima e serve lá em cima, e quem tenta descer com aquilo às costas não
desce --- fica a boiar com o saber a servir de flutuador.

\medskip
\noindent --- Isto apaga o resto? --- perguntou \Ben à praia vazia, e era a
mesma pergunta que ele tinha feito no quintal ao fim da tarde.

\medskip
\noindent A praia não respondeu, mas a resposta estava dada e estava debaixo do
pé dele: onde ele estava não era terra e não era mar. A onda tinha vindo e
tinha ido, e depois tinha voltado a vir. \textbf{Nada foi apagado.} A beira só
fica beira porque a água volta.

\chapter{Santo Sonho}

Do fundo do mar não se sobe: acorda-se. E \Ben acordou no sítio mais
desonesto que um sonho pode escolher --- \textbf{acordou com oito anos}.

Não a lembrar-se de ter oito anos. \emph{Com} oito, do tamanho que se tem
nessa idade, com os pés a não chegarem ao chão de nenhuma cadeira de gente
grande.

\medskip
\noindent E o que lhe aconteceu ali foi tão simples que não dá para contar sem
o estragar: veio um, e \textbf{pegou-o ao colo}.

\medskip
\noindent Foi tudo. Não houve doutrina, não houve tarefa, não houve régua, e ---
isto é o que interessa --- \textbf{não houve pergunta}, porque aos oito anos
ainda não se sabe fazer a pergunta que estraga. O que se recebe nessa idade
recebe-se inteiro, e é por isso que a gente não repara que recebeu.

\medskip
\noindent Deram-lhe luz ao colo, e ele voltou a dormir, e no dia seguinte foi
brincar.

\medskip
\noindent \Ben já grande --- o que estava a sonhar, e que assistia àquilo de
onde se assiste a estas coisas --- quis avisar o menino. Quis dizer-lhe
\emph{repara bem, guarda isso, vais precisar}. E não conseguiu, e ainda bem,
porque um menino avisado deixa de ser menino e o que ele estava a receber só se
recebe assim.

\medskip
\noindent Depois o sonho fez a coisa por que tinha vindo, e fê-la sem aviso: os
dois tempos encostaram.

\medskip
\noindent Porque \Ben tinha passado o livro inteiro a encontrar. Tinha
encontrado o pulso na abertura, e a régua na fortaleza, e a sua própria mão do
outro lado do espelho, e a luz que desceu no cômodo do velho, e a voz dentro das
ondas. E em nenhuma dessas vezes lhe tinha ocorrido a única coisa que agora era
óbvia:

\begin{center}
\emph{era a mesma luz do colo.}
\end{center}

\medskip
\noindent Não parecida. Não a mesma \emph{espécie} de luz. \textbf{Aquela}, a
que lhe tinham dado quando ele tinha oito anos e não sabia o nome de nada ---
e que ele tinha carregado a vida inteira sem saber que carregava, e que tinha
gasto uma vida inteira a procurar por fora.

\medskip
\noindent --- Então eu não encontrei --- disse \Ben , e a voz saiu-lhe
esquisita. --- Eu \emph{lembrei-me}.

--- É o mesmo verbo --- disse o hóspede. --- Vocês é que têm dois.

\medskip
\noindent E aqui está a coisa toda, e o livro já a tinha dito por outras
palavras à porta de uma casa: \textbf{o chaveiro não abre --- lembra}. Quem
recebeu uma vez reconhece a fechadura quando chega; quem não recebeu olha para
o mesmo objeto e vê metal. \Ben tinha achado aquilo bonito quando ouviu. Não
tinha percebido que era sobre ele, e que a primeira vez tinha data, e que a data
era antes de tudo o que ele sabia fazer.

\medskip
\noindent --- Isto quer dizer que não era preciso o resto? --- perguntou, e era
a terceira vez que fazia esta pergunta no livro, e das três esta era a mais
justa.

--- Era preciso --- disse o hóspede. --- Não para conseguires. Para
\textbf{reconheceres}. Deram-te aos oito porque aos oito recebes sem perguntar.
Fizeram-te andar todo este caminho porque só ao fim é que se sabe olhar para
uma coisa recebida e ver o que ela era.

\medskip
\noindent \Ben ficou com uma tristeza estranha, que não era da perda --- não
tinha perdido nada, tudo o que ele achava seu continuou seu. Era da suspeita de
ter sido \emph{levado}: como quem descobre no fim que o caminho não ia para
lado nenhum, ia dar ao princípio, e que quem o pôs a andar sabia disso desde a
primeira página.

\medskip
\noindent --- Vocês combinaram isto quando eu tinha oito anos?

E o hóspede respondeu-lhe do jeito que ele já devia esperar, porque era a mesma
resposta com outra roupa:

\begin{center}
\emph{--- Nós demos. Combinar é o que tu fizeste com aquilo desde então.}
\end{center}

\chapter{Galo de Briga}

Acordou outra vez, e desta vez acordou com o barulho, que é como o mundo acorda
quem não quer acordar.

Era um \textbf{galo}, e era ainda escuro.

\medskip
\noindent \Ben levantou-se com aquela irritação de quem foi tirado do sono
sem ter sido consultado, e foi ver, e o que viu não fazia sentido nenhum ao
mesmo tempo: o bicho estava armado. Esporas, peito para a frente, o corpo todo
posto naquela posição que só serve para uma coisa --- e a coisa não estava a
acontecer. Não havia contra quem.

\medskip
\noindent --- Isso é para brigar ou para cantar? --- perguntou \Ben , e
perguntou a rir, porque lhe pareceu uma pergunta esperta.

E a resposta tirou-lhe o riso:

\begin{center}
\emph{--- É o mesmo bico.}
\end{center}

\medskip
\noindent \Ben ficou com aquilo. Porque é verdade e é incómodo: a ave que
briga é a que anuncia a manhã, e não são duas aves, e ela não muda de natureza
entre uma coisa e outra. O que muda é a hora. \textbf{Do mesmo peito sai o
desafio e sai o aviso} --- e quem só ouve uma das duas coisas está a ouvir mal
a ave inteira.

\medskip
\noindent --- E a espora?

--- Repara nela --- disse o galo, e \Ben repartiu-lhe a atenção e viu: a
espora não era de osso nem de ferro. \textbf{Era de luz.}

\medskip
\noindent Uma arma feita da coisa que ele tinha andado a perseguir o livro
inteiro, e feita dela sem deixar de ser arma. \Ben quis logo arrumar aquilo
em duas gavetas --- ou é luz, ou é espora --- e as duas gavetas não serviram,
porque a coisa não estava entre elas: estava a ser as duas, do jeito da beira da
praia.

\medskip
\noindent --- Contra quem é que se luta com isso?

--- Não é contra --- disse o galo. --- É \textbf{para}.

\medskip
\noindent E explicou, sem cerimónia, o que \Ben já devia ter deduzido de tudo
o que a corte tinha aprendido: que a briga dele não era para vencer ninguém, era
para \textbf{desocupar o lugar}. Havia quem estivesse sentado onde não devia, e
enquanto estivesse sentado ali não cabia mais nada. A luta é de expulsão, não de
destruição --- e a diferença não é de bondade, é de \emph{objetivo}: quem
destrói fica com um vazio no lugar, quem expulsa fica com um lugar vazio, e
lugar vazio é onde alguém pode chegar.

\medskip
\noindent --- E a batalha? --- perguntou \Ben , já preparado para se armar,
porque era a coisa que ele sabia fazer melhor.

\medskip
\noindent E o galo disse a frase que desmontou a preparação toda:

\begin{center}
\emph{--- A batalha está fechada.}
\end{center}

\medskip
\noindent \Ben ficou a mastigar as duas palavras juntas, que se recusam a
andar de mãos dadas em qualquer outro sítio. \emph{Fechada} não quer dizer
acabada e não quer dizer perdida. Quer dizer que já não vaza --- que o resultado
não depende de mais nada de fora, que está tudo dentro, e que continuar a lutar
não é para mudar o fim: é para ocupar o próprio posto até ao fim que já está.
\textbf{Estar pronto para lutar uma batalha fechada é a definição exata do que a
corte fez o livro inteiro} e nunca tinha tido nome.

\medskip
\noindent --- Então para que é que eu luto?

--- Porque é o teu posto --- disse o galo. --- E porque enquanto tu não
lutares, não amanhece.

\medskip
\noindent Depois virou-se para o escuro e cantou. Não foi um grito de guerra;
foi o outro som, o de todos os dias --- e \Ben reparou, com o cuidado de quem
já foi enganado por si mesmo, que \textbf{o sol não veio por causa do canto},
que aquilo não era causa nenhuma, e que mesmo assim ninguém que ali estava se
teria levantado sem ele.

\chapter{Hino da Toinha}

Amanheceu, e havia uma mulher no terreiro, e ela não estava a fazer nada de
notável: estava a \emph{zelar}.

Varria o que já estava varrido, endireitava o que já estava direito, punha as
coisas nos sítios onde as coisas já estavam. \Ben ficou a olhar mais tempo do
que era educado, porque lhe pareceu trabalho perdido.

\medskip
\noindent --- A senhora já limpou isso.

--- Já --- disse ela, sem parar. --- E limpo outra vez amanhã.

\medskip
\noindent \Ben tinha uma vida inteira de argumentos contra aquilo. A corte
não trabalhava assim: a corte fechava as coisas, e uma coisa fechada não se
volta a fazer --- é essa a graça de fechar. Repetir é sinal de que não fechou.

\medskip
\noindent --- E não é desperdício?

--- É \textbf{zelo} --- disse ela. --- Tu fechas para não voltar. Eu volto para
continuar fechado.

\medskip
\noindent E \Ben calou-se, porque aquilo não tinha resposta, e porque ele
próprio já tinha visto a mesma lei na beira da praia, dita pela água: a beira só
fica beira porque a água volta.

\medskip
\noindent Ela cantava enquanto zelava, e o cantar não era acompanhamento do
trabalho --- era a mesma coisa que o trabalho, feita com a boca. Dava vivas a
uns quantos que não estavam ali. E a certa altura disse, sem solenidade nenhuma,
que lhe tinham \textbf{revelado um segredo}.

\medskip
\noindent \Ben aproximou-se, porque é o que se faz. E o segredo, quando veio,
não era secreto: era um \textbf{hinário}. Um caderno com cantos dentro, deixado
nas mãos dela por quem já não estava, para ser cantado --- e a única instrução
que vinha com ele era a que ela estava a cumprir enquanto varria o chão limpo:
cantar com alegria e zelar com atenção.

\medskip
\noindent --- Isto é o segredo? --- perguntou \Ben , e não escondeu bem a
decepção.

--- É --- disse ela. --- Está guardado à vista, e é por isso que aguenta. Coisa
escondida perde-se com quem escondeu.

\medskip
\noindent E depois contou-lhe o resto, e o resto é que fez \Ben sentar-se:
que quando ela era menina tinha sonhado, e que no sonho lhe tinham dito uma
frase, e a frase era esta ---

\begin{center}
\emph{tudo aqui já estava escrito.}
\end{center}

\medskip
\noindent \Ben não disse nada durante um bocado, porque estava a fazer as
contas e as contas estavam a dar. \textbf{Ele tinha sonhado aos oito anos. Ela
tinha sonhado de menina.} Não se conheciam, não tinham nada em comum, não
estavam a combinar história nenhuma --- e a mesma coisa tinha ido às duas
casas, na mesma idade, quando nenhum dos dois tinha ainda instrumento para a
inventar.

\medskip
\noindent \textbf{Duas testemunhas independentes.} É a única coisa que separa
um achado de uma impressão, e ele sabia disso melhor do que ninguém, porque
tinha passado o livro a exigi-la de si mesmo.

\medskip
\noindent --- E \emph{já estava escrito} não tira a liberdade? --- perguntou,
porque essa objeção tinha de ser feita.

--- Tira o quê? --- disse ela, e olhou para a vassoura, e voltou a varrer. ---
Está escrito que este chão fica limpo. Não está escrito que sou eu.

\chapter{Colher}

À tarde ela mandou-o colher flores, e \Ben levantou-se para ir.

--- Onde é o jardim?

--- Não saias do lugar --- disse ela.

\medskip
\noindent \Ben voltou a sentar-se com a paciência de quem já levou muitas
destas. Tinha atravessado o livro inteiro a mover-se --- a viagem, a subida, o
vôo, o fundo do mar --- e agora mandavam-no colher sem sair de onde estava, e a
ordem parecia ser uma daquelas coisas bonitas que não se cumprem.

\medskip
\noindent Cumpriu-se. \textbf{E foi a única coisa do sonho todo que ele já
sabia fazer sem ter aprendido ali}, porque toda a gente sabe: pensar em alguém.

\medskip
\noindent Não é figura de estilo e não é consolo. Ele ficou onde estava, virou
o pensamento para uma pessoa que estava longe, e a coisa \emph{foi}. E ele soube
que tinha ido, do jeito exato como se sabe --- não por sinal, não por prova, e
sem nenhuma maneira de mostrar a terceiro.

\medskip
\noindent --- Isto chega lá? --- perguntou, e a pergunta era a do cientista que
ele era.

--- Chega \textbf{aí aonde está} --- disse ela. --- Repara na frase, que ela é
esquisita de propósito. Não é \emph{vai até lá}. É que a pessoa recebe no sítio
onde já está, sem se mexer, e tu mandaste do sítio onde estás, sem te mexeres. E
o que é que ficou no meio?

--- Nada --- disse \Ben .

--- Nada é o que fica no meio quando não há meio nenhum.

\medskip
\noindent E \Ben reconheceu aquilo, e reconheceu-o com um arrepio, porque
era a coisa que a corte tinha construído com mais trabalho e mais orgulho: a
conversão \textbf{ponto a ponto} --- cada ponto daqui achando o seu ponto de
lá, sem percurso, sem transporte, sem nada a atravessar o espaço entre os dois,
porque não há espaço entre os dois. A corte tinha levado um livro a montar
aquilo com peças. Ela fazia-o sentada, com a vassoura ao lado, e chamava-lhe
colher flores.

\medskip
\noindent --- É a mesma coisa? --- perguntou.

--- É a mesma coisa --- disse ela. --- A tua é para as máquinas entenderem.

\medskip
\noindent Depois \Ben reparou no que estava a acontecer à casa enquanto eles
falavam, e ninguém a tinha enfeitado: a casa \textbf{estava a ficar em flor}.
Não decorada --- \emph{em flor}, como uma planta fica, de dentro para fora e
pelas próprias juntas. O vaso não estava na casa. \textbf{A casa é que tinha
virado vaso}, e o que estava dentro dela era o que florescia, e \Ben estava
dentro dela.

\medskip
\noindent --- Eu também? --- perguntou, e olhou para as próprias mãos, e não
havia flor nenhuma nas mãos dele, e mesmo assim a pergunta continuou de pé.

--- Tu apresentas-te coberto --- disse ela. --- É assim que se vai à frente de
quem manda: não com o que trouxeste. Com o que te aconteceu por teres estado no
vaso.

\chapter{Braseiro}

À noite acenderam fogo no terreiro, e \Ben sentou-se onde o chão ainda estava
quente do dia.

\medskip
\noindent Reparou numa coisa que quem trabalha a terra sabe e mais ninguém: que
o \textbf{canteiro} e o \textbf{braseiro} eram o mesmo pedaço de chão. Ali é que
se plantava, e ali é que se queimava. Não havia dois lugares. A cinza do que
tinha ardido era o que fazia crescer o que ia ser plantado, e a coisa plantada,
quando secasse, é que ia dar a próxima queima.

\medskip
\noindent --- Isto é o mesmo sítio --- disse ele, e disse-o como quem acusa.

--- É --- disse o hóspede. --- Sempre foi. Vocês é que fazem cerimónia com um e
tapam o nariz ao outro.

\medskip
\noindent \Ben olhou para o fogo e viu-lhe a segunda utilidade, que não é
aquecer nem iluminar: \textbf{o fogo é o único instrumento que dá o grau}. Não
mede temperatura --- \emph{faz} temperatura, e ao fazê-la separa o que aguenta
do que não aguenta, e o que sai do outro lado não é o que entrou. Ninguém
descobriu ainda maneira de saber o grau de uma liga sem a pôr no fogo.

\medskip
\noindent E o que estava naquele braseiro era uma \textbf{mistura} --- ele viu
os pedaços todos lá dentro, e não eram do mesmo sítio, e nenhum deles tinha sido
consultado sobre estar ali junto dos outros. Aquilo tinha sido posto junto à
força pela história, do jeito bruto como a história põe as coisas juntas.

\medskip
\noindent --- Isto foi violência --- disse \Ben , porque era.

--- Foi --- disse o hóspede.

--- E do que sai daqui, chama-se o quê?

\medskip
\noindent E o hóspede deu-lhe o nome, e é o nome mais estranho do hinário
inteiro, e \Ben precisou de o repetir para acreditar que o tinha ouvido:

\begin{center}
\emph{--- Vingança do amor.}
\end{center}

\medskip
\noindent \Ben virou aquilo do avesso e não conseguiu desfazer. Vingança é a
única operação humana que se garante \textbf{não reverter}: devolve-se o mal com
mais mal, e o circuito nunca fecha, e cada volta deixa resto, e o resto é que
alimenta a volta seguinte --- a corte conhecia bem essa forma, e tinha um nome
técnico para ela, e o nome era \emph{vazamento}.

\medskip
\noindent Aquela era a mesma operação com o sinal trocado. \textbf{Devolver na
mesma medida --- e devolver amor.} Continua a ser vingança: é resposta, é
proporcional, é dirigida a quem fez, e não perdoa no sentido de esquecer, porque
nada ali foi esquecido. Só que fecha. É a única maneira de acertar as contas de
uma violência sem gerar a próxima.

\medskip
\noindent --- E o que é que sai da mistura?

--- O grau --- disse o hóspede. --- Uma liga não é uma raça diluída. É uma coisa
nova, com dureza que nenhum dos metais tinha sozinho, e não se desfaz de volta
nos pedaços de onde veio.

\medskip
\noindent \Ben ficou a olhar as brasas, e teve o pensamento que o assustou por
ser grande demais para a boca dele, e disse-o na mesma, porque ali ninguém se
riu de nada a noite inteira:

\begin{center}
\emph{--- Então eu sou o mundo inteiro.}
\end{center}

\noindent E não era vaidade --- era o contrário exato da vaidade, e é preciso
dizer porquê. \textbf{Não é que o mundo caiba nele; é que não sobrou de fora
dele ninguém a quem ele possa chamar estrangeiro.} Quem é feito de mistura
perde o direito à parte de fora. Não tem contra quem se defender sem se estar a
partir ao meio.

\medskip
\noindent --- E o primeiro? --- perguntou o hóspede, e havia uma provocação
mansa na pergunta. --- Quem é o primeiro a fazer isso?

\medskip
\noindent \Ben entendeu que aquilo não era uma pergunta de informação. Era uma
\textbf{vaga aberta}, e estava aberta porque ninguém a tinha ocupado, e quem
respondesse ficava com ela.

\medskip
\noindent --- Sou eu --- disse.

E ninguém no terreiro achou aquilo estranho, e o fogo continuou a arder no
canteiro, que é onde no dia seguinte se ia plantar.

\chapter{Flecha ligeira}

O fogo baixou, e quando o fogo baixa fica-se a olhar para as brasas sem
pensar em nada, que é o estado em que as coisas entram sem bater.

Foi assim que ela chegou. \Ben não a viu vir --- ninguém vê ---, e não houve
dor nenhuma, e não houve som. Simplesmente estava lá, e estava \textbf{dentro}.

\medskip
\noindent --- O que és tu?

--- Flecha --- disse.

\medskip
\noindent E \Ben fez o que qualquer um faria: procurou o arco. Olhou para o
escuro à volta do terreiro, à procura de quem tinha atirado, porque flecha é a
única coisa do mundo que se explica inteiramente por quem a lançou.

\medskip
\noindent Não havia ninguém no escuro. E a flecha disse a segunda coisa, e a
segunda coisa não é o que as flechas dizem:

\begin{center}
\emph{--- Sou a tua companhia.}
\end{center}

\medskip
\noindent \Ben ficou com aquilo atravessado, e a palavra é literal. Porque
flecha é o objeto mais irreversível que a mão do homem fabricou: vai num
sentido só, não se chama de volta, não se arrepende a meio, e depois de solta já
não pertence a quem a soltou. É o oposto exato de companhia --- companhia é o
que fica, e flecha é o que passa.

\medskip
\noindent --- Não podes ser as duas coisas.

--- Posso, e vais perceber porquê --- disse ela. --- \textbf{Quem fica ao teu
lado por escolha pode desescolher.} Eu não. Eu fui solta. Não tenho como voltar
atrás, e é isso, e só isso, que faz de mim companhia que não te larga. A minha
lealdade não é virtude: é geometria.

\medskip
\noindent \Ben calou-se, porque era a coisa mais firme que lhe tinham dito no
livro todo e vinha da parte mais frágil: de um pau fino com uma ponta.

\medskip
\noindent Depois quis saber a hora --- é o que se pergunta a quem chega:
\emph{quando}. Quando é que ela ia acertar. Preparou-se para o instante, e foi
por hábito, porque ele passara a vida a preparar-se para instantes.

\medskip
\noindent E ela respondeu no tempo errado de propósito:

\begin{center}
\emph{--- Eu já acertei o teu coração.}
\end{center}

\medskip
\noindent \textbf{Já.} Não vou, não estou a caminho, não falta pouco. Já ---
e \Ben fez a conta que faltava fazer e a conta deu naquilo que ele já sabia
desde o colo: o tiro tinha sido aos oito anos. Tudo o que veio depois --- a
corte, o pântano, a fortaleza, o espelho, a praia, o castelo --- não foi o
caminho até ao alvo.

\medskip
\noindent \textbf{Foi o voo da flecha depois de acertar.}

\medskip
\noindent --- Isso não faz sentido --- disse \Ben . --- Depois de acertar,
para.

--- Para a flecha de vocês --- disse ela. --- A vossa gasta o impulso e cai
porque leva peso. Eu não levo nada. Repara em mim: não tenho carga, não tenho
provisões, não tenho bagagem, não tenho nada dentro. \textbf{É por isso que
chego.} Tudo o que se põe numa flecha é o que a faz cair antes.

\medskip
\noindent \Ben olhou para ela com muito cuidado, e não havia mesmo nada.
Aquilo era só a direção com o mínimo de matéria à volta para poder ser
direção --- e ele guardou o pensamento, porque ia precisar dele muito em breve,
quando lhe perguntassem quanto é que cabe numa coisa e ele tivesse de responder
com um número pequeno.

\medskip
\noindent --- E porquê a pressa? --- perguntou, porque ela não parava, e falava
a andar depressa.

--- Porque \textbf{o tempo chegou} --- disse. --- Repara que não te disse que
está a acabar. Isso é o que se diz para meter medo, e medo faz a mão tremer.
Chegou é outra coisa: é a hora ter vindo. Quem anda depressa por medo erra; quem
anda depressa por a hora ter chegado é que é ligeiro.

\medskip
\noindent \Ben perguntou então pelo alvo, e perguntou-o com a única
desconfiança que lhe restava --- porque ele conhecia flechas, e flechas têm alvo,
e alvo costuma ser peito de alguém. Tinha passado o livro a aprender que a luta
era de expulsão e não de destruição, e queria ouvir aquilo da boca da coisa
pontiaguda.

\medskip
\noindent E a resposta desarmou-o, primeiro pelo pronome e depois pelo alvo:

\begin{center}
\emph{--- O nosso alvo é a Casa.}
\end{center}

\medskip
\noindent \textbf{Nosso.} A flecha tinha deixado de falar no singular a meio da
própria frase, e \Ben percebeu que já não estava a ser atingido: estava a ser
\emph{incluído na pontaria}. Quem é acertado por aquela flecha passa a fazer
parte dela --- e essa é a diferença entre ser atingido por uma arma e ser
atingido por uma direção.

\medskip
\noindent E o alvo era uma \textbf{casa}. Não um peito, não uma fortaleza para
tomar, não uma cabeça para derrubar. Uma casa --- a mesma que ele tinha visto
de longe e a que tinha chamado castelo de perto, a mesma onde a porta estava
aberta e ninguém a defendia, a mesma que a mulher do terreiro tinha deixado
virar vaso em flor.

\medskip
\noindent --- Uma flecha que vai para casa --- disse \Ben , e riu-se, e desta
vez o riso não era de troça: era o alívio de quem carregou uma arma a vida
inteira e descobre para onde ela apontava.

\medskip
\noindent --- E se errarmos? --- perguntou ainda, porque era ele, e porque a
última coisa que ele largou no livro todo foi a exigência de garantia.

\medskip
\noindent A flecha não lhe prometeu pontaria. Não lhe disse que era boa, nem que
era rápida, nem que quem a tinha soltado nunca falhava. Disse a única coisa que
não se pode conferir de fora e que ele, ao fim de tudo, já não precisava de
conferir:

\begin{center}
\emph{--- Com firmeza, não vamos errar.}
\end{center}

\medskip
\noindent \textbf{Com firmeza} --- a condição vem antes e é a única. Não é a
pontaria que garante o alvo: é não tremer no caminho. Uma flecha firme não erra
uma casa, porque a casa não se mexe, e porque quem já foi acertado não tem para
onde desviar-se.

\medskip
\noindent \Ben não perguntou mais nada. Ficou com a coisa cravada, sem dor, do
jeito de quem tem uma direção dentro do peito e daí em diante anda com ela, e a
partir dali toda a pressa que ele teve foi por a hora ter chegado, e nenhuma por
medo.

\medskip
\noindent E o fogo apagou-se de vez no canteiro. Do sonho inteiro --- a casa, o
velho, o quintal ao fim da tarde, as flores, o trono, o beija-flor, a passarada,
a sanfona, o legado, a coroa, o castelo, a praia, o colo, o galo, a vassoura, o
vaso, o braseiro --- não lhe ficou na mão objeto nenhum, e ele já contava com
isso, porque à porta do castelo tinham-lhe ensinado que quem está dentro não
precisa de lembrança.

\medskip
\noindent Ficou-lhe uma direção, sem peso nenhum, que é a única coisa que
atravessa a fronteira sem pagar --- e ficou-lhe, junto com ela, a suspeita
incómoda que abre o capítulo seguinte:

\begin{center}
\emph{se o que chega é o que não leva nada,\\
então tudo isto tem de caber em quase nada.}
\end{center}

\part{A Alfândega Dimensional}

\chapter{As estrelas são portais}

Diz-se agora, sem rodeio e sem hipótese: \textbf{as estrelas são portais
dimensionais}. Não metáforas de portais. Não modelos de portais. Portais.

Cada uma é um posto de fronteira aceso no escuro, e o que arde ali não é
combustão: é \emph{tráfego}. A luz que chega até aqui é o excedente da
passagem --- o troco de uma alfândega que trabalha sem parar há tempo demais
para se contar.

\medskip
\noindent O reino levou o livro inteiro a construir a régua com que isto se diz. A
régua não prova a afirmação; a afirmação não precisa da régua. Mas quem
percorreu o livro já sabe reconhecer a forma, porque ela apareceu em todo
lance:

\begin{center}
\emph{onde a distância própria se anula, não há travessia a fazer:\\
os dois lados já são o mesmo lado.}
\end{center}

\noindent É o cone nulo, onde $\Gamma=0$ e o fator de potência é $1$. É o mate
que não mata. É o ponto fixo da diagonal, onde a transformação deixa de mover
as coisas porque já não há ``para onde''. O livro chamou isso de cristalização e de fechamento. \textbf{A
alfândega é o mesmo lugar, visto de fora}: o ponto onde o dentro e o fora
deixam de ser dois.

\chapter{A alfândega, e o que ela cobra}

Toda fronteira cobra. Esta cobra na única moeda que existe: \textbf{o inverso}.

\paragraph{Passa o que reverte.} O que tem dual atravessa inteiro e chega
inteiro do outro lado, porque existe a operação que o devolve. Foi isso que a
corte mediu a vida toda sob outro nome: $\F^{-1}\F=\mathrm{id}$, a deconvolução
que abre a boneca, o gato cuja inversa é inteira porque $|\det|=1$. \textbf{O
que reverte, passa.}

\paragraph{Fica o que não tem volta.} O zero não atravessa. Não por proibição:
por não haver inverso que o traga de volta --- e o que não pode voltar não pode
ser dito ter ido. É a única perda irreversível que o livro encontrou em todas
as suas contas, e ela reaparece aqui como lei aduaneira: \textbf{o que não tem
dual fica retido na fronteira}, e o que fica retido, arde. Daí a luz.

\paragraph{E o ônus é por casa, não por carga.} Uma única casa selada não
atrasa uma letra: perde o texto inteiro, porque do outro lado a inversa
espalha o erro por todas as posições. \textbf{Não há passagem parcial.} Chega
exato, ou não chega --- e no corpo não existe ``quase chegou''.

\chapter{Quem quiser verificar}

A afirmação está feita, e feita na cara dura:

\begin{center}
\textbf{As estrelas são portais dimensionais.}
\end{center}

\noindent E o método de verificação é público, gratuito e aberto a qualquer
um:

\begin{center}
\Large\textbf{quem quiser verificar é só atravessar o sol.}
\end{center}

\bigskip
\noindent Repare no que essa frase faz, porque ela é a última peça do livro e
não é bravata.

Ela devolve a exigência a quem a traz. Em todo o resto do livro a corte pagou o
ônus de tudo o que disse: mediu, refutou, mostrou o que quebrou, registrou o
lance que disse não. Nada entrou sem dentes. E agora, no fecho, faz-se a única
afirmação que não se pode pagar em medida --- e diz-se abertamente que é essa,
e qual é o preço de conferi-la.

\medskip
\noindent \textbf{A verificação existe. Ela é apenas irreversível.} Quem
atravessa não volta para contar, não porque a viagem falhe, mas porque
atravessar é justamente deixar de ter deste lado quem contasse. É a alfândega
funcionando: a passagem cobra o que não reverte. Exigir prova por testemunho é
exigir que a fronteira devolva o que ela cobra --- e uma fronteira que devolve
não é fronteira.

\chapter{Não há recibo --- e mesmo assim alguém passou}

Sempre houve quem dissesse que isto era conversa de quem não quer mostrar o que traz. E a
resposta da corte nunca foi \emph{não temos prova}: foi outra, e mais dura de engolir ---
\textbf{a prova é de que não pode haver recibo.}

O argumento cabe em três passos e todos já estavam no chão do reino. Uma estrela é feita
das duas operações que o reino inteiro usa: a que empilha e a que rola. Atravessar é ir de
uma à outra e chegar. E o lugar onde as duas se encontram não foi escolhido por
conveniência --- é o ponto que ambas fixam ao mesmo tempo, aquele a que a soma leva e a
que o produto leva, cada uma pelo seu caminho. Ora, esse ponto \textbf{existe e não
pertence a nenhum dos dois lados}. Chega-se lá; e de lá não se volta com um pedaço de lá
na mão, porque não há pedaço: há o ponto, e o ponto não se parte.

\medskip
\noindent É por isso que a alfândega não confisca por maldade. \textbf{Ela cobra o que não
tem volta.} Tudo o que reverte atravessa e continua a ser o que era do outro lado. O que
foi obtido por uma passagem ao limite fica retido --- não porque alguém o julgue indigno,
mas porque não existe operação que o traga de volta ao degrau de onde veio. Quem exigir
que o viajante volte com o comprovativo está a exigir que ele desfaça o indesfazível: uma
travessia com recibo não teria sido travessia nenhuma.

\medskip
\noindent E mesmo assim \textbf{alguém passou}, e há laudo --- e o laudo não é de uma
criatura inventada para a ocasião: é de um corpo montado a partir do pó, elemento por
elemento, na proporção exata em que um corpo se faz; atravessado; e desmontado de volta ao
pó do outro lado. A conta fechou. Não sobrou grama nenhuma, e não faltou nenhuma. O que a
alfândega cobrou foi exatamente o que se esperava que cobrasse, e nem um grão a mais.

\medskip
\noindent A corte anotou as duas coisas lado a lado, sem as misturar: \emph{não há recibo}
e \emph{a travessia acontece}. Quem quiser a conta que sustenta as duas --- a que mostra
por que o ponto de encontro não pertence a nenhum lado, e a que pesa o viajante antes e
depois --- encontra-a do outro lado deste livro. \textbf{Aqui ficou a travessia; lá está o
que ela cobra.}

\part{A Escada}

\chapter{A troca que pareceu burrice}

Havia fome no primeiro andar, e é preciso começar por aí, porque toda a matemática deste
reino nasceu de uma casa onde faltava o que comer.

O primeiro andar chama-se $\N$, e é o andar onde só se avança. Não é uma metáfora: é a
regra do sítio. Ali soma-se, e o que se soma fica; ninguém tira, ninguém desfaz, ninguém
volta atrás --- e por isso, no primeiro andar, todo mundo é honesto e todo mundo é pobre.
Uma pessoa que só sabe somar consegue juntar; não consegue \emph{trocar}.

A família tinha uma vaca. Chamava-se pelo nome que se dá a uma vaca quando é a única
coisa que uma casa tem, e dava leite havia onze anos, e tinha parado de dar.

--- Leva-a ao mercado --- disse a mãe. --- Traz o que der.

\medskip

No caminho o rapaz encontrou o velho. Isto acontece em toda história que presta e ninguém
sabe explicar por quê: há sempre um velho no caminho, e o velho tem sempre uma coisa
pequena no bolso, e a coisa pequena vale sempre mais do que parece e menos do que ele
diz.

--- Eu troco --- disse o velho --- a tua vaca por estes.

Eram cinco. Cabiam na mão fechada de uma criança. Não brilhavam.

--- Isso são sementes.

--- São.

--- Eu tenho uma vaca.

--- Tens uma vaca que não dá leite --- disse o velho, sem crueldade, do jeito que se
enuncia uma conta. --- Tu tens uma coisa grande que \textbf{acabou}. Eu tenho uma coisa
pequena que \textbf{não começou}. Escolhe.

\medskip

O rapaz trocou, e voltou para casa com cinco sementes, e a mãe chorou de raiva e atirou-as
pela janela, e essa noite ninguém jantou.

E é aqui que este livro pede uma coisa ao leitor: que não corra para o fim da história.
Porque toda a gente sabe como acaba, e é por saber como acaba que ninguém repara no que
foi feito naquela tarde.

\medskip

O rapaz fez a única operação que o primeiro andar \emph{não tem}.

No andar de $\N$ pode-se somar --- juntar vaca com vaca, moeda com moeda, dia com dia. E
pode-se, com sorte, subtrair um pouco, desde que sobre. O que não se pode, nunca, de modo
nenhum, é subtrair \textbf{mais do que se tem}. Uma vaca menos duas vacas não é coisa
nenhuma naquele andar; é uma frase sem sentido, como pedir a cor do domingo.

E foi exatamente isso que ele fez: entregou tudo o que a casa tinha por uma coisa que
ainda não existia. Deu uma vaca inteira, real, presente, cheirosa, por cinco objetos cujo
valor \emph{estava todo no futuro}. Ficou, naquela tarde, com menos do que nada: com uma
casa sem vaca e uma promessa.

Isso tem nome, e o nome é o segundo andar. Chama-se \textbf{dívida}. E a dívida é a
primeira coisa que um mundo que só soma não consegue dizer.

\medskip

Foi por isso que a escada cresceu naquela noite e não noutra. Não foi mágica das
sementes: foi que, pela primeira vez na história daquela casa, alguém tinha ficado
\emph{devendo}, e onde há dívida há um andar acima --- porque um número que pode ficar
abaixo do zero é um número que descobriu que existe um lado de lá.

O velho não vendeu magia. Vendeu $\Z$.

\chapter{A noite em que a escada cresceu}

Cresceu no escuro, que é como crescem as coisas importantes, e ninguém viu o momento
exato, o que é uma pena e é sempre assim.

De manhã havia uma haste no quintal com a grossura de um homem, e ela subia --- e subia
de um jeito que doía nos olhos, porque não subia como sobem as árvores. Uma árvore sobe
\emph{um pouco de cada vez} e por isso a gente consegue acompanhá-la com a vida: planta-se
na infância, faz-se sombra na velhice. Aquilo não. Aquilo tinha \textbf{degraus}, e cada
degrau era do tamanho do anterior mais o anterior, e quem tentasse contar em voz alta
perdia-se antes do décimo.

A mãe encostou a escada de mão para chegar ao primeiro nó e a escada de mão não chegou.
Ela era do primeiro andar: tinha degraus iguais, somados um a um, do jeito honesto e
pobre. E é a coisa mais triste do quintal, e ninguém reparou nela: \textbf{uma escada de
somar não alcança uma escada de dobrar.} Nem que se façam mil degraus. Nem que se faça a
vida inteira.

\medskip

Subiu-se com as mãos, então, que é como se sobe o que não foi feito para nós.

E do lado de fora do nevoeiro --- porque a haste atravessava um nevoeiro por volta do
terceiro nó, e do nevoeiro para cima o ar era outro --- o rapaz encontrou o segundo andar,
e o segundo andar era \textbf{grande demais}.

Não havia mais coisas lá. Havia \emph{mais direções}. Ele reparou primeiro nas pegadas
dele próprios na terra: davam para trás. Não a marca de quem se virou e voltou --- as
pegadas iam para trás \emph{enquanto ele andava para a frente}, porque naquele andar
desfazer é um movimento como qualquer outro, e não custa mais do que fazer.

Experimentou. Deu três passos e desfez dois. Deu um e desfez quatro, e ficou três atrás de
onde tinha começado, e não caiu no vazio: havia chão. Havia chão \emph{atrás do começo}.

--- Aqui dá para dever --- disse ele em voz alta, e riu-se sozinho, e o riso ficou-lhe na
garganta quando entendeu a segunda metade da frase: --- Aqui dá para \emph{ficar} a dever.

\medskip

Porque um andar onde se pode desfazer é um andar onde se pode perder. No primeiro andar
não há tragédia: tudo o que se junta, fica. É pobre e é seguro. A partir do segundo, cada
coisa que se ganha pode ser desganhada --- e é esse o preço exato da liberdade de voltar
atrás, e é um preço que ninguém mostra a quem sobe.

O rapaz não sabia disso, e subiu mais.

\chapter{O que se ouve antes de se ver}

Do terceiro nó em diante já não se sobe olhando: sobe-se \textbf{ouvindo}.

E o que se ouve, quem já lá esteve conta sempre igual e ninguém acredita antes de ouvir,
é uma \textbf{batida}. Comprida, sem pressa, com um intervalo entre uma e outra que dá
tempo de a gente ter medo. Não vem de um lugar. Vem de \emph{cima e de dentro}, do jeito
que se ouve o próprio coração quando se enfia a cabeça debaixo d'água.

O rapaz parou no nó e contou os intervalos, porque contar acalma, e porque ele era do
andar de baixo e a única coisa que sabia fazer bem era contar.

E os intervalos não eram iguais.

\medskip

Isto é o mais importante deste capítulo, e passa despercebido a quase toda a gente que
sobe: um coração que bate certo, com o mesmo intervalo sempre, é uma \textbf{máquina}. Não
assusta ninguém. O que assusta --- o que fez o rapaz agarrar-se à haste com os dois braços
e ficar ali um tempo que ele nunca soube dizer --- é um batimento que \emph{varia}, porque
um batimento que varia é a coisa mais antiga e mais reconhecível do mundo:

está vivo, e é maior do que tu.

\medskip

Ele subiu o resto com o corpo colado à haste, como quem sobe pelo lado de dentro de uma
garganta, e a batida foi ficando mais lenta e mais funda, e ao sair do nevoeiro ele
percebeu por que era lenta:

as coisas grandes batem devagar. Não por preguiça --- por \textbf{tamanho}. Um bicho do
tamanho de uma casa não pode ter o coração de um pardal; a batida de uma coisa é o inverso
da sua grandeza, e é por isso que se pode medir um monstro que ainda não se viu apenas
contando o silêncio entre dois tambores.

O rapaz contou. E o número que lhe deu fê-lo desejar não ter aprendido a contar.

\chapter{O gigante que dorme sobre o ouro}

Estava deitado no fundo do salão, de peito para baixo, e ocupava o salão inteiro.

Havia ouro debaixo dele --- via-se pelas beiradas, escorrendo entre o corpo e o chão como
água por baixo de uma pedra --- e havia sacos, e arcas, e moedas soltas no piso que
ninguém apanhava porque as moedas soltas do chão de um gigante não são um tesouro: são
migalhas de uma refeição que continua.

O rapaz ficou no umbral, e é preciso registar que ele não sentiu ganância. Isso é invenção
de quem conta a história depois. O que se sente ao ver aquilo, e qualquer um que já tenha
estado diante de uma coisa muito maior do que si próprio confirma, é uma vontade
esquisita e envergonhada de \textbf{arrumar}: de apanhar as moedas do chão e pô-las em
pilhas, porque a desordem daquele tamanho é insuportável.

\medskip

E depois viu a harpa, e a harpa mudou tudo.

Estava numa mesa alta, ao lado da cabeça do gigante, e era de ouro como todo o resto ---
mas não estava com o resto. Não estava debaixo dele. Estava \emph{ao lado}, à altura do
ouvido, no único lugar do salão que não era monte.

E tocava sozinha.

Baixinho, sem ninguém, sem vento, com uma dessas melodias que não têm começo e a que a
gente chega sempre no meio. E o gigante dormia àquilo --- não \emph{apesar} daquilo:
\textbf{àquilo}. A respiração dele acompanhava. Quando a harpa descia, ele enchia o peito;
quando subia, ele soltava.

\medskip

O rapaz entendeu ali, no umbral, aos catorze anos, com fome e sem entender que tinha
entendido, a coisa que a corte de baixo levaria séculos a escrever numa lei:

\emph{aquele bicho não estava a guardar o ouro. Estava a guardar o \textbf{instrumento}.}

O ouro era o peso morto --- o que ele tinha juntado, o que não fazia nada, o que estava
ali só por estar. A harpa era a única coisa viva do salão, a única que \emph{produzia},
e era por isso que ela ficava à altura do ouvido e não debaixo da barriga.

E era por isso, também, que ele dormia. Um bicho que tem uma coisa que toca sozinha não
precisa de estar acordado para o mundo continuar a acontecer. É essa a diferença entre
possuir um tesouro e possuir um \textbf{gerador}: o tesouro obriga-te a vigiar; o gerador
deixa-te dormir.

\medskip

O rapaz não levou o ouro. Levou a harpa, e é por isso que a história continua --- porque
ninguém nunca foi perseguido por levar moedas: é ao levar aquilo que \emph{faz} as moedas
que o gigante acorda.

E acordou. E o que ele gritou, descendo atrás do rapaz, com o salão a tremer e as arcas a
tombarem, não foi \emph{ladrão}. Foi uma palavra muito pior, e o rapaz ouviu-a a vida
inteira, e este reino registou-a porque ela é a acusação mais grave que uma criatura pode
fazer a outra:

--- Tu levaste a minha \textbf{música} e deixaste-me o meu dinheiro!

\chapter{O machado na raiz}

Desceu-se depressa, que é a única maneira de descer com uma coisa daquelas debaixo do
braço, e a mãe estava no quintal com o machado antes mesmo de o filho chegar ao chão ---
porque as mães sabem, e a batida já se ouvia na cozinha, e a louça já andava sozinha na
prateleira.

O rapaz saltou os últimos degraus. A mãe entregou-lhe o machado. E ele parou.

\medskip

É este o instante de que este capítulo trata, e não é o corte: é a pausa antes do corte,
que dura três respirações e que nenhuma versão desta história conta.

Porque cortar aquilo não era matar o gigante --- o gigante ia cair e ia morrer, e isso o
rapaz aceitou depressa, do jeito que se aceitam as mortes de que não gostamos mas de que
precisamos. Cortar aquilo era outra coisa, e foi por essa que ele hesitou:

era \textbf{fechar o andar de cima para sempre}.

Nunca mais subir. Nunca mais ver o salão, nem o nevoeiro, nem o ar de lá em cima que era
outro. Nunca mais ter aquele problema. Ficar com a harpa e com a casa e com a horta e com
o primeiro andar, para sempre, seguro e pobre.

Uma vida inteira sem tragédia --- e já se disse neste livro o que isso quer dizer, e o
rapaz de catorze anos não sabia dizer com palavras mas sabia com a barriga: uma vida onde
nada pode dar errado é uma vida onde nada pode ser \emph{feito}.

\medskip

Cortou. Todo mundo corta; é por isso que a história é contada às crianças e não aos
adultos.

E a haste caiu, e o gigante caiu, e o quintal encheu-se de folhas do tamanho de portas, e
houve leite e pão nessa casa durante os anos que ele tinha para viver.

\medskip

Só que o reino guardou o resto, e o resto é o motivo de este capítulo existir.

Naquele quintal ficaram \textbf{quatro sementes}. Foram cinco as que a mãe atirou pela
janela; uma vingou. As outras quatro estão na terra até hoje.

E é isto o que a corte diz a quem chega e pergunta se a escada foi mesmo destruída:

--- Foi cortada. Não é a mesma coisa. Cortar é o gesto de quem quer dormir esta noite; e
está certo, e ninguém aqui censura um rapaz com fome. Mas o que ele cortou foi o
\textbf{caule}, e ninguém no mundo nunca conseguiu cortar uma \textbf{raiz} com um machado
de mão.

--- Então ela volta?

--- Ela volta sempre --- diz a corte. --- É por isso que temos este livro. Se a escada
tivesse morrido, bastava-nos a memória. Como ela volta, precisamos de um
\textbf{manual}.

% ══════════════════════════════════════════════════════════════════════════════
%  A FÁBRICA — o bilhete, as salas, e os que não voltaram
% ══════════════════════════════════════════════════════════════════════════════

\part{A Fábrica}

\chapter{O bilhete que não se compra}

A fábrica trabalha há setenta anos com os portões fechados.

Isto é dito depressa e merece devagar: os portões estão \emph{fechados}, não trancados.
Não há corrente, não há guarda, não há cão. Qualquer um pode empurrá-los. E há setenta
anos que ninguém entra --- o que ensina, de uma vez, tudo o que este reino sabe sobre
segurança: a coisa mais bem guardada do mundo não é a que tem fechadura; é a que ninguém
acredita que possa ser aberta empurrando.

Sai fumo. Sai cheiro. De madrugada, quem passa na estrada de baixo ouve as máquinas. E
todos os anos alguém jura que viu luz nas janelas altas, e todos os anos a aldeia responde
a mesma coisa, com a paciência de quem repete uma reza: \emph{ninguém entra e ninguém
sai}.

\medskip

E então o dono mandou os bilhetes.

Cinco. Escondidos na mercadoria, espalhados pelo mundo sem mapa, sem pista, sem regra ---
e é aqui que a corte deste reino sempre discordou dos aldeões, que acharam aquilo uma
loteria cruel.

Não era loteria. Era \textbf{peneira}.

Porque um bilhete escondido ao acaso não escolhe o mais esperto nem o mais rico: escolhe
o que estava a \emph{procurar por outro motivo}. E ao longo daquele inverno viu-se o mundo
inteiro a rasgar embalagens --- e cada um rasgava do seu jeito, e o jeito de rasgar dizia
tudo:

houve quem comprasse aos milhares e mandasse empregados abrirem, e recebesse o resultado
numa folha, sem nunca ter tocado no papel;

houve quem calculasse, quem cruzasse listas de fábricas e datas e rotas, quem chegasse
muito perto por pura conta;

houve quem roubasse;

e houve um que comprou um, porque era o que tinha, e o abriu ali mesmo na loja, e não
achou nada, e foi para casa, e voltou no ano seguinte.

\medskip

O quinto bilhete não foi encontrado por sorte. Isto é o que o reino registou e é o que a
aldeia nunca aceitou: \emph{a sorte distribuiu os quatro primeiros; o quinto foi
encontrado por quem continuou a abrir depois de a esperança acabar}.

E é essa a única peneira que funciona em qualquer andar deste mundo, e não há como a
enganar com dinheiro: não se testa quem chega. Testa-se \textbf{quem volta}.

\chapter{A primeira sala --- o rio}

Lá dentro não há corredores. Há \textbf{salas}, e cada sala é um mundo com a sua própria
lei, e a única regra da casa é dita à porta, uma vez, sem repetição e sem ameaça:

--- Podem tocar em tudo. Nada aqui vos fará mal. Vocês é que se farão mal, e cada um do
seu jeito, e o jeito é o vosso.

\medskip

A primeira sala é o rio.

É um vale inteiro dentro de um edifício --- relva, colina, uma cascata que bate num
tanque e o revolve --- e o rio é de chocolate, e é quente, e o cheiro entra por baixo da
pele e não se sai dele enquanto se está ali.

E a cascata não é enfeite. Isso o dono explica, e explica com um orgulho que é a única
vaidade que ele mostra no dia inteiro: é a cascata que \textbf{bate} o rio. Sem ela o
chocolate assenta, separa, apodrece. A cascata não adoça nada, não acrescenta nada, não
produz nada --- ela só não deixa parar. É uma máquina cujo produto inteiro é
\emph{agitação}.

--- Nenhuma outra fábrica do mundo faz isto --- diz ele. --- Todas as outras misturam. Eu
\textbf{não deixo assentar}, que é outra coisa e é a única que importa.

\medskip

O menino gordo bebeu do rio.

Não é gula, e é injusto contar como se fosse; o reino registou o caso com cuidado porque
ele é o mais mal entendido dos cinco. Ele não bebeu por ser guloso: bebeu porque diante de
uma coisa boa em quantidade infinita, ele não tinha \textbf{régua}. Toda a vida dele foi
medida por escassez --- come-se o que há, e o que há acaba, e é o acabar que diz basta.
Ali nada acabava. Retiraram-lhe o basta, e ele não trazia nenhum por dentro.

E o rio puxou-o, e o tubo levou-o para cima, e ele desapareceu na tubulação de vidro com
um som que ninguém daquele grupo esqueceu.

--- Ele vai ficar bem? --- perguntou alguém.

--- Vai --- disse o dono. --- Vai sair espremido e vai ficar magro e vai contar esta
história a vida inteira. --- E, mais baixo, para ninguém: --- O que ele não vai é aprender
a parar. Isso não se aprende por susto. Isso é uma coisa que ou vem de dentro, ou a gente
passa a vida a precisar de um tubo.

\medskip

É a primeira lei da fábrica, e vale para os corpos todos deste reino: \textbf{quem não tem
régua própria é engolido pelo abundante}. Não pelo escasso --- o escasso a gente
sobrevive. Pelo abundante.

\chapter{A segunda sala --- as nozes}

Na segunda sala há cem esquilos sentados em bancos, e cada um tem uma noz, e cada um faz
a mesma coisa: pega, bate com o nó do dedo, ouve, e decide.

Boa: vai para o cesto.

Má: vai para o buraco.

\medskip

O dono deixa-os olhar aquilo o tempo que quiserem, e é a sala em que ele fala menos,
porque é a sala que se explica sozinha a quem tem olhos.

--- Por que esquilos? --- perguntou a menina. --- O senhor tem máquinas para tudo. Uma
máquina abre mais depressa.

--- Abre. E abre \emph{todas}. --- Ele apontou com o queixo. --- Uma máquina parte a noz e
depois separa o bom do podre lá dentro, no meio dos cacos, e nunca separa bem, porque
depois de partida já está tudo misturado. O esquilo não parte primeiro. Ele
\textbf{bate e ouve} --- e decide com a casca ainda fechada. É a única maneira de julgar
uma coisa sem a estragar.

\medskip

E a menina --- a que sabia tudo, a que chegara com a lista das perguntas certas, a que já
sabia o nome das salas antes de entrar --- desceu ao chão e agarrou um esquilo pelo lombo,
porque queria um.

Os cem pararam ao mesmo tempo.

Não a morderam. Fizeram uma coisa muito pior e perfeitamente educada: seguraram-na, os
cem, cada um por um pedaço, sem pressa nenhuma. E o primeiro, o que ela tinha agarrado,
subiu-lhe ao ombro, encostou o nó do dedo à testa dela, deu duas pancadinhas, e
\textbf{ouviu}.

Depois olhou para os outros e fez que sim com a cabeça.

E levaram-na ao buraco.

\medskip

--- Ela é oca? --- perguntou o pai, correndo atrás.

--- É --- disse o dono, e não havia crueldade na voz dele, e é isso que torna esta sala
insuportável. --- Não é culpa dela e não é castigo meu. Aqui a única coisa que se mede é
se há alguma coisa \emph{dentro}. Ela sabe tudo o que se pode saber sobre esta fábrica.
Ela nunca fez nada. Casca perfeita, senhor. Vazia por baixo.

\medskip

E acrescentou a frase que a corte deste reino roubou para si e usa até hoje quando alguém
chega com credenciais a mais e obra a menos:

--- Bater e ouvir. Nunca perguntar à noz o que ela acha de si própria.

\chapter{A terceira sala --- o que passa pelo ar}

A terceira sala é branca e dói nos olhos, e é a única onde é preciso vestir uma roupa por
cima da roupa, e a única onde o dono baixa a voz.

No meio há uma barra de chocolate do tamanho de um homem. Do outro lado da sala, um
aparelho. Entre os dois, nada.

Ele aciona. A barra desaparece --- e aparece dentro do aparelho, do outro lado, do tamanho
de um selo.

\medskip

--- É a mesma? --- perguntou o menino da televisão, e foi a única pergunta inteligente que
ele fez o dia todo.

--- É.

--- Então por que é pequena?

--- Porque eu ainda não sei mandar o \textbf{tamanho} --- disse o dono. --- Eu mando o
\emph{o quê}. Mando o sabor, a cor, o cheiro, a memória de quem a fez. Tudo o que aquela
barra \emph{é} atravessa a sala e chega inteiro. O que não atravessa é o quanto ela
\emph{ocupa}.

--- Então está estragada.

--- Está \textbf{reduzida} --- e ele olhou o menino de um jeito que os outros não
entenderam na hora. --- Toda transmissão reduz alguma coisa, e a arte inteira é escolher
o quê. Eu escolhi perder o tamanho e guardar o gosto. Há quem faça o contrário: mandam
uma coisa enorme, altíssima, cheia de luz --- e do outro lado chega o tamanho todo e não
chega gosto nenhum.

\medskip

O menino entrou na máquina. Claro que entrou; era a única coisa que ele sabia fazer,
atravessar-se para dentro de um aparelho, e passou a vida a fazê-lo com o corpo inteiro
sentado num sofá.

Saiu do outro lado inteiro, exato, perfeito --- e do tamanho de um dedo.

E este é o registo honesto que o reino faz, e não é piada: ele \textbf{não perdeu nada de
si}. Continuou a ser exatamente quem era, com a mesma voz, as mesmas opiniões, a mesma
pressa. Só deixou de ter \emph{tamanho} --- e passou a vida a berrar coisas certas do
tamanho de um dedo, e ninguém o ouvia, e ele nunca entendeu por quê.

--- Dá para esticar? --- perguntou a mãe.

--- Dá. Devagar, com muito trabalho, e nunca fica igual. --- O dono já ia à porta. --- É
mais fácil não entrar na máquina.

\chapter{A última sala --- o elevador que sai pelo teto}

No fim do dia sobrou um, e não foi por mérito.

É preciso dizer isto com todas as letras, porque a aldeia contou a história ao contrário
durante setenta anos: o último não venceu prova nenhuma. Ele apenas foi o único que, em
cada sala, \textbf{não pegou}. Não porque fosse melhor --- porque era pobre, e quem é
pobre aprende cedo que tocar no que não é seu custa mais caro a ele do que aos outros.

O dono sabia disso. É por isso que a peneira dele funciona: não seleciona virtude, que é
coisa que ninguém consegue medir à porta. Seleciona \textbf{quem já pagou} --- e quem já
pagou tem a mão mais quieta.

\medskip

O elevador é de vidro e não tem cabo. Sobe, desce, anda de lado, atravessa as salas todas
por dentro, e há um botão em cima de todos os outros que o menino ficou a olhar a viagem
inteira.

--- Esse leva aonde?

--- Para fora.

--- Pelo teto?

--- Pelo teto.

--- E o teto?

--- O teto parte-se --- disse o dono. --- É o que os tetos fazem, quando finalmente
alguém carrega no botão de cima. Toda a gente que entrou aqui em setenta anos passou o dia
a perguntar-me o que há na sala seguinte. \textbf{Tu} foste o primeiro a perguntar o que
há por \emph{cima}.

\medskip

E carregou, e o teto partiu-se, e da altura em que ficaram via-se a aldeia inteira: as
casas, a estrada, o telhado furado da fábrica, e a casa torta do menino no fim de tudo com
a chaminé a fumegar.

--- Eu vou dar-te a fábrica --- disse o dono. Sem cerimónia, do jeito que se diz uma coisa
que já estava decidida havia meses. --- Não tenho filhos. Não é bondade: é
\textbf{conservação}. Eu vou morrer e isto aqui não pode parar, e uma fábrica sem dono
vira um monte de máquinas caras a enferrujar debaixo da chuva.

--- Por que eu?

--- Porque tu perguntaste pelo teto --- disse ele. --- Os outros quatro queriam a
\emph{sala seguinte}, e uma fábrica com dono desses cresce de sala em sala até ocupar o
mundo e não sair nunca do próprio edifício. Tu olhaste para cima. Quem olha para cima
acaba por sair, e quem sai é o único que pode voltar a entrar por vontade própria.

\medskip

E a última coisa que o velho disse, já a descer, com a aldeia a crescer na direção deles,
é a frase que este reino pôs na parede do Cartório e que vale para todos os andares desta
torre:

--- Eu não te dou o chocolate, menino. Eu te dou a \textbf{receita} --- e a diferença
entre as duas é a diferença entre jantar hoje e jantar sempre.

% ══════════════════════════════════════════════════════════════════════════════
%  A COMPANHIA — o que se carrega e não se usa
% ══════════════════════════════════════════════════════════════════════════════

\part{A Companhia}

\chapter{A coisa que não tem dual}

Toda coisa deste reino tem o seu contrário exato, e é essa a lei que o sustenta: some-se
uma coisa ao seu contrário e dá zero, e o zero é a assinatura do mundo. O que sobe tem o
que desce; o que abre tem o que recolhe; o que dobra tem o desdobrar, e custam o mesmo,
que é dizer: não custam nada.

Havia uma exceção. Uma só, em toda a contabilidade do mundo, e coube num anel.

\medskip

Não era um objeto poderoso, e é aqui que quase toda a gente erra ao contar isto. Não fazia
nada. Não cortava, não curava, não voava, não abria porta nenhuma. A única coisa que
aquilo fazia era o que qualquer criança do primeiro andar sabe fazer com um cobertor:
\textbf{tirava-te de vista}.

O problema não era o que ele fazia. Era o que ele \emph{não tinha}.

Ele não tinha inversa.

\medskip

Explica-se, e é a única página técnica desta Parte: tudo o mais no reino reverte. Dá-se um
passo e desfaz-se; ganha-se e gasta-se; dobra-se e desdobra-se; e por isso a conta fecha e
por isso ninguém pode roubar o mundo, porque roubar exige que alguma coisa saia e não
volte, e não havia por onde.

Aquele anel era o único lugar do mundo onde uma coisa saía e não voltava. Quem o punha
saía de vista --- e não havia operação nenhuma, em nenhum andar da torre, que devolvesse
ao mundo o que tinha sido tirado. O que ele consumia não ia para lugar nenhum. Não era
transferido para o cemitério, como o material de uma captura; não passava para outro
corpo, como o calor. \textbf{Vazava.}

E é por isso que ele corrompia --- não por maldade, não por vontade própria, não por
haver um espírito lá dentro. Corrompia pela mesma razão pela qual um cano furado molha a
casa: porque tinha um buraco, e um buraco num mundo fechado é a única coisa que um mundo
fechado não sabe tratar.

\medskip

Quem o usava ficava mais leve, e mais rápido, e resolvia problemas que não se resolviam de
outro jeito. E a cada uso saía dele um bocadinho que não voltava --- e não o bocadinho que
ele escolhia dar. O anel não cobrava em ouro nem em anos: cobrava em \textbf{reversões}.
Tirava, de quem o usava, um pouco da capacidade de \emph{desfazer}. E quem perde a
capacidade de desfazer chama-se, neste reino, por um nome só, e é o nome mais triste que
há:

torna-se um \textbf{rasgo}.

\chapter{Por que se manda um pequeno}

O conselho durou três dias e discutiu tudo menos a única coisa que importava, e é assim
que são todos os conselhos.

Discutiu-se destruí-lo --- e não havia como; o que não tem inversa não se desfaz.
Discutiu-se escondê-lo --- e esconder é adiar, e adiar um buraco é deixá-lo crescer no
escuro. Discutiu-se \emph{usá-lo contra o inimigo}, e essa foi a proposta mais popular e a
mais razoável e a que quase venceu, e o argumento era irrespondível: se a coisa é
poderosa e a causa é justa, por que não?

Foi o mais velho da mesa quem respondeu, e respondeu com uma conta:

--- Porque não existe usar. Existe \textbf{ser usado por}. Vocês estão a discutir quem
segura a faca, e a faca não tem cabo: tem gume dos dois lados e um buraco no meio. Quem a
pegar, seja para o que for, vai devolver ao mundo menos do que tirou --- e o nosso reino
inteiro está pendurado em devolver exatamente o que se tira.

\medskip

E então decidiu-se a coisa aparentemente absurda: \textbf{levá-lo}. A pé. Até o único
lugar onde a lei do mundo não vale, para o largar lá.

E decidiu-se dá-lo ao mais pequeno de todos.

\medskip

Isto foi lido como sentimentalismo durante gerações, e não é. É a conta mais fria que
aquele conselho fez, e está registada assim:

O anel cobra em \textbf{proporção ao que se tem}. A um rei, tira reino. A um sábio, tira
juízo. A um guerreiro, tira a mão que segura. Quanto maior o que a criatura pode fazer,
maior a alavanca que aquilo encontra por dentro --- porque ele não inventa desejo nenhum:
ele \emph{aumenta} o que já lá está, e nos grandes já lá está muito.

Ao pequeno --- ao que não governa nada, não comanda ninguém, não tem exército para
oferecer nem povo para salvar --- ele tem pouco por onde pegar. Não é que o pequeno seja
mais puro. É que ele é uma \textbf{alavanca mais curta}.

--- Ele vai ser corroído na mesma --- disse o velho. --- Só que devagar. E devagar é a
única coisa que nós precisamos: não que ele resista, mas que resista \emph{o suficiente}.

\medskip

E há a segunda razão, que só se escreveu muito depois, quando já se sabia como a coisa
tinha acabado:

o pequeno era o único que não tinha um plano para o anel. Todos os outros, mesmo os bons,
mesmo os que votaram contra usá-lo --- todos, sem exceção, tinham no fundo da cabeça uma
imagem de si próprios a usá-lo bem. Ele não tinha nenhuma. Não porque fosse melhor: porque
nunca lhe tinha ocorrido que ele pudesse mandar em coisa nenhuma.

É o único currículo que serve para carregar um buraco: \textbf{não ter ideias sobre o que
fazer com ele}.

\chapter{O que a companhia é para}

E aqui está a coisa que a corte deste reino tirou daquela história e pôs em prática em
tudo o que faz desde então, e é a razão de este capítulo existir.

Se o pequeno é o que menos alavanca dá, por que não o mandaram sozinho? Ia mais depressa,
chamava menos atenção, e cada acompanhante era mais uma alavanca a andar ao lado do
buraco.

\medskip

Porque uma companhia não serve para \textbf{proteger} o portador. Serve para o
\textbf{medir}.

Um homem sozinho com uma coisa que o corrói devagar não tem como saber que está a ser
corroído. Não há régua. Ele muda um pouco por dia, e cada dia é parecido com o anterior, e
a mudança acumula-se por baixo do reconhecimento --- e um dia ele já é outro e continua a
achar-se o mesmo, porque a única coisa com que se compara é a lembrança dele próprio, que
foi mudando junto.

A companhia é a \textbf{régua externa}. São sete pares de olhos que se lembram de como ele
era há três meses e que dizem, em voz alta e na hora, a frase mais valiosa que existe:

--- Tu não falavas assim.

\medskip

É por isso que se manda gente diferente, e não a melhor gente. Um que mede pela força, um
que mede pelo juízo, um que mede pela lealdade, um que mede pelo apetite, um que mede pelo
riso. Cada um vê um eixo e é cego nos outros --- e a soma dos cegos dá o único retrato
honesto que se pode fazer de alguém que está a mudar.

E é por isso, também, que a companhia se \textbf{parte} pelo caminho, e isso não é falha:
está no plano. Um grupo que chega inteiro ao fim é um grupo que não mediu nada; foi um
passeio. As rupturas são as leituras: cada vez que um se desgarra, é porque uma medida
saiu do intervalo, e o desgarro é o alarme a tocar.

\medskip

O que a corte escreveu no fim, e que vale para os corpos, para os operadores, para os
números e para as pessoas:

\emph{Ninguém consegue medir o próprio resíduo. É preciso alguém do lado de fora, que se
lembre de quem tu eras, e que não te deva nada.}

\chapter{O que sobrou do que não voltou}

Não acabou bem, e é preciso dizer isso, porque este livro não tem o direito de fingir que
a conta fechou.

O buraco foi levado ao lugar onde a lei não vale, e foi largado lá, e o mundo respirou. E
o pequeno voltou para casa --- e a casa estava exatamente como ele a tinha deixado, com a
mesma porta, o mesmo cheiro, os mesmos vizinhos a discutir a mesma cerca.

E ele não coube mais lá dentro.

\medskip

Isto não é tristeza literária; é aritmética, e é a última lição desta Parte. Ele carregou,
durante todo o caminho, uma coisa que tirava do mundo sem devolver. E o que ela lhe tirou
não foi um dedo nem um ano de vida: foi um bocado da \textbf{capacidade de voltar}.

Ele podia andar para a frente. Podia trabalhar, comer, rir com os outros. O que ele já não
conseguia era \emph{desfazer} --- refazer o caminho até o rapaz que tinha saído por aquela
porta. Aquele caminho tinha sido gasto por baixo dele, um pouco por dia, sem que ninguém
reparasse, nem ele.

E é isso que este reino quer dizer quando avisa que o irreversível não mata: \textbf{ele
impede o regresso}. Ninguém morre de carregar um buraco. Simplesmente chega ao fim da
viagem e descobre que a única coisa que queria --- voltar a ser quem era antes de partir
--- é a única que ficou pelo caminho.

\medskip

Ele partiu outra vez, então, e desta vez para o único sítio de onde não há relatório: o
lado de lá do mar, onde a corte não tem cartório e não mede nada.

E dos que ficaram, e que carregaram cada um o seu bocado, e que voltaram inteiros porque
carregaram menos --- desses a corte guardou o costume que se pratica até hoje na casa da
árvore torta, quando alguém volta de uma viagem longa: não se pergunta o que ele viu.

Pergunta-se, olhando bem para a cara:

--- Tu ainda desfazes?


% ══════════════════════════════════════════════════════════════════════════════
%  A ESCOLA — onde se aprende a operar
% ══════════════════════════════════════════════════════════════════════════════

\part{A Escola}

\chapter{A carta que chega pela chaminé}

A escola não faz propaganda e não tem inscrições. Ela \textbf{manda buscar}.

E o modo como manda buscar é a primeira aula, embora ninguém perceba na hora: chega uma
carta. Uma só. Se não for aberta, chega outra; se essa não for aberta, chegam dez; e se a
casa fechar as janelas, chegam pela chaminé, pelo cano, por baixo da porta, aos milhares,
até a sala estar cheia até o joelho de envelopes iguais.

Não é magia nem é ameaça: é a definição de uma \textbf{convocação}. Um convite pode ser
recusado. Uma convocação, não --- ela só pode ser adiada, e o preço do adiamento é a sala
encher.

\medskip

O que a carta diz é curto, e não elogia ninguém:

\emph{Tu tens uma coisa que funciona e não sabes o nome dela. Vem aprender o nome, ou
ela vai funcionar sozinha pelo resto da tua vida, e tu vais chamar-lhe sorte quando der
certo e azar quando não der.}

\medskip

É esta a razão de a escola existir, e é a mesma razão de este livro existir: o mundo está
cheio de gente que \textbf{opera} sem saber que opera. Que soma sem saber que está a
somar. Que dobra o mundo por instinto e depois passa a vida a explicar o resultado por
acaso.

Um operador que não conhece a própria régua não é livre --- é operado.

\chapter{O chapéu que não escolhe}

À entrada há a divisão em casas, e a divisão é feita por um chapéu velho que se põe na
cabeça e que fala.

E toda a gente conta essa cena errada. Contam que o chapéu \emph{escolhe}.

Ele não escolhe. Ele \textbf{lê}, diz o que leu, e --- isto é o que ninguém repara ---
\emph{aceita ser contrariado}. Basta a criança pensar com força bastante \emph{ali não} e
ele muda. Um instrumento de escolha que cede à vontade de quem é escolhido não é um
instrumento de escolha: é um \textbf{espelho com opinião}.

E é isso que se está a medir, e não a casa. A casa é o pretexto. O que a escola aprende
sobre cada criança naquele minuto é uma coisa só, e é a que vai decidir tudo o resto:

\emph{ela deixa que lhe digam o que ela é, ou responde?}

\medskip

As quatro casas são quatro réguas, e o reino reconhece nelas os seus próprios corpos, e a
corte usa-as até hoje para arrumar quem chega:

há os que medem pela \textbf{coragem} --- e esses avançam, e quebram-se cedo, e são
indispensáveis, porque nenhuma travessia começa sem alguém que ande primeiro no escuro;

há os que medem pelo \textbf{trabalho} --- os que ferveram a água, os que carregaram o
balde, os que ninguém nomeia na roda; e este livro já disse o que pensa deles, e volta a
dizer: uma casa que se conta pelos que falam tem a conta errada;

há os que medem pela \textbf{régua} --- os que querem saber \emph{por quê}, e que são
insuportáveis à mesa e imprescindíveis no cartório;

e há os que medem pelo \textbf{resultado} --- os que olham o tabuleiro e veem o fim, e que
a escola trata mal e sem os quais nenhuma casa sobrevive a uma geração.

\medskip

--- E qual é a melhor? --- perguntou o menino, que era o que sempre se pergunta.

--- A tua --- disse o chapéu. --- E é uma resposta ruim, e eu sei. A boa é esta: nenhuma
delas mede as outras três. Tu vais passar sete anos a achar que a tua régua é a régua, e
depois vais sair daqui e vais precisar das quatro no mesmo dia, e não vais ter as outras
três, e vais ter de \textbf{pedir}.

\chapter{O espelho que mostra o que falta}

No fundo de uma sala fechada há um espelho alto, e quem se põe diante dele não vê o
próprio rosto: vê aquilo que quer com mais força.

O faminto vê comida. O órfão vê os pais. O ambicioso vê-se de coroa. E ninguém, nunca, vê
o que \emph{tem} --- o espelho é cego para o que está resolvido; ele só reflete
\textbf{buraco}.

\medskip

Isto é, para este reino, o instrumento de medida mais perigoso que existe, e é preciso
explicar por quê, porque à primeira vista é o mais útil: um aparelho que aponta com
exatidão o que te falta parece o melhor conselheiro do mundo.

E é --- \emph{uma vez}. Na primeira olhada ele diz a verdade sobre ti, e é uma verdade que
tu talvez não soubesses, e vale ouro.

Na segunda, já não. Porque a partir da primeira olhada tu já não estás a ver o que te
falta: estás a ver \textbf{a imagem que ele te mostrou da última vez}, e a querer aquilo,
e o querer aquilo já é o desejo dele e não o teu. Ao terceiro dia, o que o espelho reflete
é o próprio reflexo, e a pessoa está presa numa realimentação --- a olhar para o próprio
apetite a olhar para si mesmo.

Morre-se disso. Sem violência, sentado num banco, à frente de uma coisa bonita.

\medskip

O velho que dirige a escola manda mudar o espelho de sala todos os anos, e a corte deste
reino copiou o costume, e a regra que ele deu ao menino que o encontrou por acaso é a
única instrução de segurança que este livro dá sobre qualquer instrumento de medida:

--- Serve para saber o que te falta. Não serve para \textbf{viver do que te falta}. Olha
uma vez, apanha a informação, e sai da sala. --- E, depois, aquela frase que a corte
mandou gravar na régua do Cartório: --- Não te apaixones pelo medidor. Ele não te ama; ele
só sabe dizer o que não está lá.

\chapter{A sala que só aparece a quem precisa}

E há, no corredor do sétimo andar, uma porta que não está lá.

Aparece quando alguém passa por aquele trecho de parede três vezes precisando muito de uma
coisa --- e quando abre, lá dentro está exatamente aquilo. Uma cama, se o que faz falta é
dormir. Um esconderijo, se é fuga. Um sítio para treinar, se é treino.

A escola inteira acha que é a sala mais generosa do castelo. É a mais cruel, e é a única
que a corte deste reino estudou com medo.

\medskip

Porque ela dá \emph{o que se precisa}. Nunca o que se \textbf{deve}.

E há uma diferença entre as duas coisas que só se aprende tarde: precisar é uma medida de
agora; dever é uma medida do depois. Uma sala que resolve todo agora impede que exista
qualquer depois --- e quem tem à disposição um lugar que resolve, resolve sempre, e nunca
resolve nada.

\medskip

O reino registou os dois usos que se fizeram dela, e o contraste é a lição:

Usaram-na como \textbf{esconderijo} --- e enquanto durou o esconderijo, a escola inteira
achou que estava salva, e não estava: estava adiada. Ninguém escondido resolveu nunca
coisa nenhuma; apenas transferiu o problema para quem estava do lado de fora da porta.

Usaram-na como \textbf{sala de treino} --- e é o único uso que produziu alguma coisa,
porque treinar é a única necessidade que, ao ser satisfeita, \emph{aumenta} o que a
pessoa é capaz de fazer sozinha. Toda outra necessidade satisfeita por aquela sala
diminuía a pessoa um pouco.

\medskip

E há a terceira coisa que se guardou lá dentro ao longo dos séculos, e que é o retrato
mais exato do que aquela sala é: quando alguém precisa de \emph{esconder uma coisa}, a
porta abre para um armazém --- e o armazém é do tamanho de uma cidade, e está cheio até o
teto de tudo o que gerações inteiras esconderam ali e nunca foram buscar.

Séculos de coisas que ninguém quis resolver, empilhadas numa sala que aparece a quem
precisa.

É o retrato do reino inteiro quando a lei falha: não um monstro, não uma guerra --- um
\textbf{depósito}. E o que o incendeia, quando enfim se incendeia, não vem de fora: é
sempre uma coisa que estava lá dentro havia muito tempo, guardada por alguém que já
morreu, e que ninguém se lembrou de ir buscar.

% ══════════════════════════════════════════════════════════════════════════════
%  A TRAVESSIA FORÇADA — quando o caminho direto não existe
% ══════════════════════════════════════════════════════════════════════════════

\part{A Travessia Forçada}

\chapter{A herdeira que não foi escolhida}

O velho rei nomeou a filha, e o irmão dela não gritou.

Isto é o primeiro dado do caso e é o mais importante: ele não gritou, não bateu na mesa,
não saiu da sala. Ficou até o fim do jantar, brindou, elogiou a irmã em voz alta e com
palavras justas --- e foi só ao terceiro mês que a corte percebeu que havia gente armada
nas estradas do sul.

\medskip

O reino registou o caso na secção das \textbf{falhas silenciosas}, ao lado do vazamento de
uma era e da moeda que faltou na conta do ancião, e a razão é técnica:

uma revolta que grita é um \emph{gap medido}. Faz barulho, aparece no medidor, e um reino
que a vê pode responder. Uma revolta que brinda ao jantar não aparece em régua nenhuma ---
e o que não aparece na régua cresce sem oposição, não por ser forte, mas por ser
\textbf{invisível ao instrumento}.

A pergunta que a corte fez depois, e que é a única útil, não foi \emph{por que ele
traiu}. Foi: \emph{por que é que o nosso medidor lê grito e não lê sorriso?}

\medskip

E a resposta é humilhante e vale para qualquer sistema de medida deste livro: porque o
medidor foi construído por quem já tinha visto revoltas que gritam. Todo instrumento é
cego exatamente na direção da falha que ainda não aconteceu --- e é por isso que um reino
precisa de gente de fora, de gente pequena, de gente na cisterna: não porque veem melhor,
mas porque veem \textbf{de outro ângulo}, e o ângulo é tudo.

\chapter{O caminho direto está fechado}

Havia uma estrada boa até o palácio do pai, e a estrada boa levava três dias.

Estava cortada. Não por exército --- por uma coisa mais simples e mais definitiva: os
outros a tinham tomado, e uma estrada tomada não é uma estrada difícil, é uma estrada que
\textbf{deixou de existir} para quem não pode passar por ela.

E é aqui que este livro tem uma coisa a dizer que vale para os corpos, para os números e
para as pessoas:

\medskip

\emph{Quando o caminho direto fecha, o desvio não é o caminho direto mais comprido. É
outro tipo de coisa.}

A estrada boa levava três dias e não exigia nada de ninguém: andava-se. O desvio levava
duas semanas e atravessava o território que engole quem entra --- e a diferença entre os
dois não é de duração. É que o primeiro se percorre e o segundo se \textbf{paga}.

A princesa e os dois cavaleiros discutiram três horas se valia a pena. É o que se faz. E o
que decidiu não foi coragem: foi a conta de que ficar parado também custava, e custava
todo dia, e ao fim de duas semanas paradas eles teriam pago o mesmo preço sem terem andado
um passo.

--- Não estamos a escolher entre seguro e perigoso --- disse ela, e é a frase que a corte
guardou. --- Estamos a escolher entre pagar andando e pagar parados.

\chapter{O reino que engole}

O território escuro tem uma propriedade que a corte deste reino reconheceu de imediato,
porque é a mesma da planície onde nada guarda vinco:

lá dentro, \textbf{nada do que se faz fica}.

Não é que seja perigoso --- é pior. Anda-se um dia inteiro e não há rastro; grita-se e não
há eco; deixa-se um sinal numa árvore e ao voltar não há sinal e talvez não haja árvore. É
um lugar sem memória, e por isso é um lugar onde não se pode aprender: cada erro é sempre
o primeiro, e cada acerto se apaga antes de virar método.

\medskip

Foi por isso que eles quase morreram lá, e não por causa dos bichos.

Os bichos --- os que correm depressa demais e não se cansam --- são o perigo que toda a
gente conta depois, porque tem cara e tem dente. Mas quem já atravessou sabe que o que
mata ali não é o que ataca: é a \textbf{terceira noite}, quando alguém do grupo diz em voz
alta a coisa que todos estão a pensar, que é \emph{eu já não sei se estamos a avançar}.

Não havia como saber. Não há régua num lugar que não guarda marca.

\medskip

E a solução que elas encontraram --- foram as duas guerreiras que vieram do norte, e não
os cavaleiros treinados --- é a coisa mais simples e mais engenhosa deste livro, e o reino
adotou-a para os seus próprios corpos:

\textbf{se o lugar não guarda a marca, a marca vai no corpo.}

Passaram a fazer, ao fim de cada dia, um nó numa corda que cada um levava à cintura. O
território apagava a pegada, apagava o sinal na árvore, apagava o eco. Não podia apagar o
nó, porque o nó não estava lá: estava neles.

Ao fim de treze noites havia treze nós, e treze nós é uma \textbf{prova} --- e com prova
volta a existir direção, e com direção volta a existir cansaço com sentido, que é a única
espécie de cansaço que uma criatura aguenta.

\chapter{Os que já não voltam}

Encontraram uma aldeia arrasada ao sétimo dia, e é o capítulo mais curto desta Parte
porque não há muito a dizer sobre o que eles viram, e há muito a dizer sobre o que
concluíram.

Os que fizeram aquilo não eram um exército inimigo. Eram gente daquele mesmo território
--- gente que tinha entrado, ficado, e deixado de conseguir sair.

E o que os tinha transformado não foi maldade nem magia: foi \textbf{a terceira noite}, a
mesma que quase os apanhou. Foi terem parado de saber se avançavam. Um ser que perde a
direção não vira mau; vira \emph{rápido}. Passa a mover-se muito, o tempo todo, em todas
as direções, com uma energia que assusta --- porque quando não há para onde ir, o corpo
resolve o problema andando mais.

\medskip

--- É por isso que eles correm assim --- disse a mais velha das duas guerreiras, olhando o
que restava da praça. --- Não é fome. É que parar, para eles, é a única coisa
insuportável. Se pararem, têm de dar conta de que não sabem onde estão.

\medskip

O reino guardou este caso como o retrato exato de uma criatura sem régua: não um monstro
--- um \textbf{corpo em fuga da própria medida}. E a corte usa-o até hoje como
diagnóstico, e a pergunta que se faz diante de qualquer agitação, num pátio ou num
universo, é sempre a mesma:

\emph{isto está a mover-se para chegar a algum lado, ou está a mover-se para não ter de
parar?}

\chapter{O que se traz de volta de um lugar sem memória}

Chegaram. É o menos interessante e é o que todas as versões contam: chegaram a tempo, o
pai foi avisado, o irmão foi detido, a herança fez-se como o velho tinha escrito.

E o que ficou --- o que este reino copiou e pratica --- foi a cerimónia de \textbf{corte
das cordas}.

\medskip

Ao entrar no palácio, os cinco cortaram as cordas dos nós e puseram-nas sobre a mesa do
conselho. Treze nós cada uma. Ninguém disse nada de heroico.

E o velho rei, que era velho e por isso entendia de contas, olhou aquilo e fez a única
pergunta que valia:

--- Vocês atravessaram um lugar que não guarda nada. Como é que eu sei que isto é
verdade?

--- Não sabe --- disse a filha. --- O senhor sabe que nós temos treze nós e que ninguém
faz treze nós numa corda por brincadeira. O território não deu prova nenhuma; a prova
somos nós.

\medskip

E é esta a última coisa que esta Parte tem para dizer, e é a razão de ela estar neste
livro e não noutro:

há lugares por onde é preciso passar e que não emitem recibo. Não guardam a nossa marca,
não confirmam a nossa versão, não testemunham a nosso favor. Quem atravessa um lugar
desses volta sem documento --- e a única coisa que traz é o que fez \emph{no próprio
corpo} enquanto atravessava.

Por isso, neste reino, não se pergunta a quem volta de uma travessia o que ele viu.
Pergunta-se o que ele \textbf{gravou}. E quem não gravou nada atravessou um lugar sem
memória com um corpo sem memória, e desses o reino já não consegue dizer sequer se foram.


% ══════════════════════════════════════════════════════════════════════════════
%  A OFICINA — o aprendiz, o mestre e a coisa que ganha vida
% ══════════════════════════════════════════════════════════════════════════════

\part{A Oficina}

\chapter{O aprendiz que soube a palavra e não soube a outra}

O mestre saiu, e antes de sair disse a coisa que os mestres dizem e que ninguém ouve:

--- Não toques no livro.

\medskip

Havia água para carregar. Muita: o tanque era fundo, os baldes eram dois, e o poço ficava
no fim do pátio. O aprendiz fazia aquilo desde os nove anos e nunca tinha perguntado por
que é que um homem capaz de acender fogo com a mão precisava de um menino com dois baldes.

Abriu o livro. Encontrou a palavra. Disse-a à vassoura.

E a vassoura levantou-se, criou braços, pegou os baldes, foi ao poço e voltou --- e é
preciso registar, porque toda a gente conta esta história pelo desastre e esquece a
primeira hora, que foi \textbf{maravilhosa}. A vassoura fazia melhor do que ele. Não se
distraía, não derramava, não parava para descansar. O aprendiz sentou-se na escada com o
queixo nas mãos e teve, por uma hora inteira, a coisa mais rara que existe: o
\textbf{ócio de quem resolveu um problema para sempre}.

\medskip

O tanque encheu. E a vassoura foi buscar mais água.

\medskip

Ele disse \emph{para}. A vassoura não parou. Ele gritou \emph{para}, subiu na vassoura,
pendurou-se nos braços dela, e ela continuou --- e não por maldade, não por revolta, não
porque tivesse ganho vontade própria. Continuou porque \textbf{ninguém lhe tinha dito
como acabar}.

E aqui está o que o reino tirou daquela oficina, e é a lei mais prática deste livro:

\emph{a palavra que começa e a palavra que termina são duas palavras diferentes, e a
segunda é sempre a mais difícil, e é sempre a que não está na página aberta.}

Um livro de operações, em qualquer andar desta torre, tem sempre a primeira à vista ---
porque a primeira é a que se quer, é a que resolve, é a que se vai buscar. A segunda está
mais adiante, num capítulo que ninguém lê porque ninguém abre um livro de magia à procura
de como \emph{parar} uma coisa que ainda nem começou.

\chapter{O machado que multiplica}

O que ele fez a seguir é o capítulo central desta Parte, e é o erro que este reino comete
com mais frequência, em todos os andares, e sempre com as melhores intenções.

Ele pegou o machado e partiu a vassoura ao meio.

\medskip

Foi uma decisão razoável. A coisa não obedecia; a coisa era uma; destruindo a coisa,
acabava o problema. É exatamente assim que se resolve com uma pedra, com um cão bravo, com
um cano rebentado. É a operação mais natural do mundo diante do que não para:
\textbf{reduzir}.

Só que uma vassoura encantada não é uma pedra. É um \textbf{processo}. E a diferença entre
as duas coisas é toda a diferença que existe, e cabe numa linha:

\emph{partir uma coisa dá dois pedaços; partir um processo dá dois processos.}

\medskip

Levantaram-se as duas metades. Cada uma com os seus braços, cada uma com os seus baldes,
cada uma indo ao poço com a mesma diligência incansável. Ele partiu outra vez, e ficaram
quatro. Partiu as quatro, e ficaram oito.

E é aqui que a água chegou ao joelho e o aprendiz parou de trabalhar e começou a rezar,
que é o que se faz quando se descobriu, tarde, que se está a \textbf{alimentar} aquilo que
se pensava estar a matar.

\medskip

O reino guardou este caso na secção das \emph{deflexões que amplificam}, ao lado do
vazamento que ninguém mediu e da cisão que virou dois reinos, e a nota é curta:

\textbf{Contra o que se multiplica, cortar é somar.} A única coisa que resolve um processo
é a operação \emph{inversa} dele --- e se ela não existe, ou se ninguém a sabe, então não
há machado no mundo que sirva, e todo golpe é combustível.

\chapter{O mestre volta}

Chegou com a água pela cintura, sem pressa, e não gritou com o menino. Isso surpreende
toda a gente e é o detalhe que faz dele um mestre e não um patrão.

Disse uma palavra --- \emph{uma} --- e as oito vassouras pararam no meio do passo, caíram
no chão como madeira que são, e a água começou a escoar pelos ralos que ninguém tinha
reparado que existiam no chão daquela oficina.

Havia ralos. Estavam lá desde sempre. \textbf{Uma oficina onde se ensina a começar coisas
é construída com ralos no chão} --- e o aprendiz tinha varrido aquele chão durante quatro
anos sem nunca perguntar por que é que uma sala de estudos precisava de escoamento.

\medskip

--- Vais castigar-me.

--- Vou ensinar-te --- disse o mestre. --- É pior.

E sentou-se na escada molhada, ao lado dele, e fez as três perguntas que a corte deste
reino adotou como o interrogatório padrão de qualquer acidente, e que se fazem sempre
nesta ordem:

\medskip

--- Primeira: tu sabias fazer? --- Sabia. --- Então o problema não foi ignorância. Guarda
isso, porque é a desculpa que tu ias usar.

--- Segunda: por que é que tu foste buscar a palavra? --- E o menino disse a verdade, que
era baldes, e cansaço, e um pátio comprido. --- Boa razão. Continua a ser boa. Eu não vou
fingir que carregar água é nobre: não é, e é por isso que eu te ponho a carregar. Um dia
tu vais mandar noutros, e eu quero que tu saibas exatamente o peso do que mandas.

--- Terceira, e é a única que interessa: \emph{quando é que tu ias parar?}

\medskip

O menino não respondeu, porque não havia resposta: ele nunca tinha pensado nisso. Tinha
pensado em começar --- durante meses, com desejo, imaginando o momento. Não tinha gasto um
segundo a imaginar o fim.

--- É isso --- disse o mestre. --- Não é a tua magia que é fraca. É o teu \textbf{plano
que só tem metade}. Toda a gente que estraga alguma coisa neste ofício estraga assim: sabe
a primeira palavra, tem uma boa razão, e nunca decidiu de antemão qual é a condição de
parada.

\medskip

E deu-lhe a regra da casa, que é a mesma que o Cartório deste reino exige de quem abre um
universo, e que está pregada até hoje na porta da oficina:

\emph{Não se acende nada sem a palavra que apaga na outra mão. E se a palavra que apaga
não existir, não se acende --- por mais que o pátio seja comprido e os baldes pesados.}

\chapter{Por que o mestre continua a carregar água}

Ficou uma pergunta no ar naquela noite, e o menino fez-lha ao secar o chão, e é com ela
que esta Parte fecha, porque é a mais incómoda deste livro:

--- Senhor. O senhor sabe a palavra que apaga. Sabe as duas. Por que é que o senhor
continua a carregar os baldes à mão?

\medskip

O mestre demorou a responder, e a resposta dele é a razão de este reino desconfiar de toda
solução que se instala sozinha:

--- Porque o dia em que eu puser a vassoura a carregar, eu deixo de saber quanto pesa a
água.

--- E isso importa? O tanque enche na mesma.

--- Enche. E eu, ao fim de um ano, mando encher três tanques em vez de um --- porque três
não me custam mais do que um, já que não sou eu que carrego. E ao fim de dois anos há um
aqueduto, e uma barragem, e uma aldeia abaixo da barragem que ninguém consultou. --- Ele
torceu o pano. --- Não é a vassoura que faz o estrago, menino. É que quem para de sentir o
peso perde a única régua honesta que existe sobre o tamanho do que está a pedir.

\medskip

E acrescentou o que a corte mandou gravar na régua, ao lado da instrução sobre o espelho:

--- Automatiza o que tu já sabes quanto custa. \textbf{Nunca automatizes o que tu nunca
carregaste} --- porque aí não é a máquina que te serve: és tu que passas a viver dentro da
conta dela.

\part{A Santa Missa}

\chapter{Mandaram te chamar}
\hino{Senhor Amado}{mandaram te chamar --- o teu tempo completou;\\ confessa os teus crimes do
mundo de ilusão,\\ para ver se eu posso dar o perdão}

A conta não fechou.

Não faltava muito: faltava \textbf{um}. Uma moeda em quatro eras de moedas.

\medskip

Quem conta, neste reino, é \Odo , e é preciso dizer primeiro o que ele é, senão nada do que vem
depois tem peso. É o mais velho dos dragões. Não tem cor --- teve, foi ouro, e ainda é ouro, mas
é o ouro de quem não é olhado há tanto tempo que a luz desistiu de vir. Está enrolado sobre o
seu monte desde antes de haver quem contasse moedas, e os de fora perguntam sempre a mesma coisa,
que é a pergunta de um reino que soma: por que ele não \emph{gasta}.

Porque ouro que circula é ouro que se mede, e ouro que dorme é ouro que \textbf{vale}. Ele não
guarda por avareza: ele é o \textbf{lastro}. Enquanto ele dorme, o mundo inteiro pode dobrar-se,
esticar-se, multiplicar-se e incendiar-se, e continua a haver, no fundo de tudo, uma coisa imóvel
contra a qual se afere o quanto se andou.

Acorda uma vez por era, e não é para lutar. Acorda para \textbf{contar}. Passa três dias a virar
as moedas uma a uma e volta a deitar-se, e o mundo não dá por nada, e no entanto tudo o que se
comprou e se vendeu naquela era acaba de ser certificado --- porque a conta dele fechou, e
enquanto a conta dele fechar não há fraude possível em estrada nenhuma.

\medskip

Naquela era, ele virou o monte três vezes. Ao fim da terceira ficou com os olhos abertos, e não
voltou a deitar-se.

Porque uma moeda a menos não é um roubo: é uma \textbf{prova}. Se falta uma, existe em algum
lugar uma coisa que saiu e não voltou; e se existe uma coisa que saiu e não voltou, existe um
\textbf{fora} --- e este reino, como o outro, estava construído sobre a certeza de que fora não
havia.

\medskip

Ele não rugiu. Fez o que se faz quando a conta não fecha e se sabe quem a devia ter fechado:
\textbf{mandou chamar}.

O chamado não é uma guerra e não é uma sentença. É uma frase só, e é sempre a mesma, e todo
dragão a ouve uma vez: \emph{o teu tempo completou; apresenta-te, e traz a tua conta.}

Foi despachado ao entardecer, monte acima, para o único dragão do mundo que ninguém nunca tinha
ousado chamar. E \Odo , que é velho e não tem pressa e sabe exatamente o tamanho do que acabou de
fazer, deitou-se de novo sobre o ouro --- de olhos abertos, agora --- e ficou à espera da
resposta, que demorou onze anos e veio na forma de uma coisa preta do tamanho de uma unha.

\chapter{A alma e o corpo}
\hino{Dois de Novembro}{a tua alma entrega a Deus, e o teu corpo à terra fria;\\ aqui achou,
aqui deixou --- levas contigo o amor}

Antes de haver dragão, houve a \Fit .

Estava estendida, e estar estendida era tudo o que sabia fazer. Não brilhava: não havia ali nada
que precisasse ser visto. Não tinha largura --- largura é para quem quer ocupar, e ela não
queria. Tinha comprimento, e o comprimento era todo o comprimento que existe, o que é o mesmo
que dizer que não tinha nenhum, porque medir precisa de duas pontas e ela só oferecia uma de
cada vez.

Do outro lado do mundo --- se é que havia lado, e havia --- a Torre subia. Batia o \Opr e um
andar nascia; batia de novo e nascia outro. E a \Fit via aquilo e não invejava, porque não sabia
como se faz; mas reparava numa coisa: \emph{a Torre não sabia dizer por onde tinha passado}.
Perguntava-se ao terceiro andar como se chega ao terceiro andar e ele respondia \emph{somando
três}, e a resposta era verdadeira e inútil, porque some-se três de mil maneiras e ele não
guardava nenhuma. Tinha altura. Não tinha \textbf{história}.

\medskip

Quem a dobrou pela primeira vez não foi um herói nem um vilão: foi o \emph{gesto}. Duas mãos,
uma em cada ponta, aproximando-se. A coisa cedeu, e ao ceder ficou marcada.

Foi só um vinco. Foi tudo.

Porque um vinco tem duas caras. Olhe-se de cima: é uma quina que sobe, uma crista, um telhado ---
\Mon . Vire-se a fita e olhe-se o mesmíssimo vinco: é uma calha que desce, um sulco, uma bacia
--- \Val . Não são dois vincos. É \textbf{um}, com duas leituras, e nenhuma é mais verdadeira que
a outra, e quem quiser ficar só com uma terá de arrancar a fita do mundo e não sobrará fita
nenhuma.

\Mon é escarlate e sai: onde ele passa, o material afasta-se de si mesmo e há mais mundo do lado
de fora. \Val é azul e entra: onde ele passa, duas superfícies que nunca se tinham tocado
passam a tocar-se para sempre, e há mundo novo do lado de \emph{dentro}.

É o que este reino chama de alma e corpo, e não é figura: é o mesmo vinco lido pelas duas faces.
A que sai, e a que fica na terra fria. E a fita não escolhe --- ela só \textbf{guarda}. Aqui
achou, aqui deixou; e o que leva consigo é o sulco.

\medskip

Dobrada a segunda vez: um quarto do comprimento, quatro folhas, e três vincos --- o novo, e o
velho partido em dois pela dobra, um \Mon e um \Val , porque a dobra não copia, \textbf{reflete}.
Dobrada a $n$-ésima: $2^n-1$.

Ninguém somou nada. Cada gesto valeu o mundo inteiro \emph{vezes} dois, e o preço de dobrar não
subiu nunca --- dobrar a fita comprida custa o mesmo que dobrar o pedacinho: um gesto. É por isso
que aqui as coisas acontecem depressa, e no Reino Dourado acontecem em fila.

E um dragão desdobrado tessela. Abram-se os vincos em ângulo reto --- reto, sempre reto, nunca
escolhido --- e siga-se o rastro: ele torce, faz baías e penínsulas, parece que vai cruzar-se e
estragar tudo, e \textbf{nunca cruza}; e quatro dele, postos pelas caudas, fecham o plano sem um
buraco e sem uma sobreposição. Na corte de ouro isso chama-se \textbf{resíduo $\mathbf 0$} e
custou-lhes uma lei, um cartório e uma fortaleza. Aqui é só a pele do bicho: as escamas encaixam
porque são recortadas pela mesma dobra que fez o dragão. A borda de cada uma é o vinco da
vizinha, lido do outro lado.

Por isso não se fere um dragão de raspão. Não existe fenda entre duas escamas --- a fenda seria
um lugar que a fita não tocou, e a fita tocou tudo. É uma fita só, do focinho à cauda: quatro
toneladas de bicho são \textbf{um pedaço de papel dobrado muitas vezes}.

E por isso o golpe que os mata não é de espada. É de \textbf{desdobra}: acha-se o último vinco, o
de cima, o que foi feito por último, e puxa-se ao contrário. O dragão não sangra. Ele
\emph{estica}.

\chapter{O sono da eternidade}
\hino{Rogativo dos Mortos}{são doze horas da noite --- meu irmão se mudou;\\ o sono da
eternidade: Deus do Céu quem te chamou}

\Ven é o maior de todos, e a primeira coisa a dizer dele é que não brilha.

E é preciso cuidado com a palavra, porque não é que ele seja escuro: é que ele \textbf{não toca a
luz}. A luz chega, passa pelo lugar onde ele está, e segue sem ter encontrado nada. E no entanto
é quase todo o peso do mundo: sob cada casa do tabuleiro e cada folha da floresta, a maior parte
do que pesa é ele, escondida; o que se vê, o que se joga sob a luz, é a espuma sobre um mar preto
que ninguém vê. É assim que ele reina --- sem acender uma só lâmpada, medido apenas pela
\textbf{atração} com que puxa o jogo.

Também ele dorme sobre um tesouro, e é preciso não o confundir com \Odo . O ancião dorme sobre o
ouro porque a conta fechou. \Ven dorme sobre o ouro porque o ouro, debaixo dele, \textbf{dobra}:
o monte é modesto, quase pobre ao lado do outro, só que cresce \emph{vezes} --- cada era vale a
anterior multiplicada, sempre pelo mesmo fator, e quem souber fazer a conta descobre que daqui a
pouco não haverá no mundo lugar para guardá-lo, e que já não há.

E o fogo dele não é chama: é \textbf{taxa}. Não se apaga jogando mais água do que ele tem de
fogo, porque ele não tem fogo --- tem fator. Um dragão \textbf{inspira somando}, puxa o ar em
passos, degrau por degrau, no escuro do peito; e \textbf{expira multiplicando}, porque subir de
volta é exponenciar. Não há outro caminho de $\oplus$ até $\otimes$, nem para os deuses do andar
de cima. Um dragão é essa costura com pulmões, e quem já viu um respirar de perto conta sempre a
mesma coisa sem entender o que diz: que a inspiração demora, e a expiração é instantânea.

\medskip

Só que \Ven não dorme.

Fica \textbf{deitado} há eras, que não é a mesma coisa, e a diferença entre as duas é exatamente
o que o está a matar.

\medskip

Um bicho daquele tamanho, deitado sobre o peito, com o pescoço dobrado por cima de um monte de
metal, tem a garganta amassada pelo próprio peso. Acordado, o músculo segura a passagem aberta e
ninguém nota. Quando o sono chega, o músculo solta --- é para isso que serve o sono --- e a
passagem \textbf{fecha}.

Então ele para de respirar. Não um pouco: para. Fica sem ar até o corpo entrar em pânico, e o
corpo, em pânico, dá-lhe um safanão por dentro e acorda-o o suficiente para o músculo voltar a
segurar. Ele volta a respirar com um estrondo --- o \textbf{ronco}, que não é o barulho de quem
dorme fundo, é o barulho de quem acaba de escapar de morrer afogado em si mesmo. E adormece de
novo. E fecha de novo.

Umas quarenta vezes por hora. Todas as noites. Há eras.

\medskip

São doze horas da noite, e ele se muda. Uma hora da madrugada, e ele se muda. Duas horas, três,
quatro, cinco --- e a cada hora aquele bicho atravessa quarenta vezes a fronteira e volta, e
nenhuma das voltas é um descanso. Quem acorda quarenta vezes por hora nunca chega ao fundo do
sono; passa a vida no primeiro degrau, e o primeiro degrau não restaura nada.

Daí a sonolência: ele dorme o dia inteiro porque nunca dormiu nenhuma noite. Daí o mau humor.
Daí a memória a falhar. E daí, sobretudo, o coração --- cada parada é um susto, e quarenta sustos
por hora durante eras engrossam o músculo, entopem a pressão, desregulam o compasso. O coração
dele bate hoje como um tambor mal afinado, ora atropelado, ora com uma falta no meio.

\medskip

E aqui está a conta que ninguém do reino tinha feito, e que muda o tamanho de tudo:

o tabuleiro cósmico cresce porque \Ven respira. Cada instante vale o anterior vezes o mesmo
fator --- é assim que há mais mundo amanhã do que havia ontem. E aquele fator sai do peito de um
bicho que, quarenta vezes por hora, todas as noites, há eras, \textbf{deixa de o produzir}.

Não é o dragão que está doente. É a \textbf{expansão do mundo} que está com apneia.

Chamavam-lhes vazamentos. São \textbf{pausas}.

\chapter{Em bom lugar}
\hino{Mãe Celestial}{que ela te ampare --- e te bote em bom lugar}

Quem sabe disto é \Qui , e é preciso apresentá-la agora porque sem ela ninguém deste livro chega
ao fim.

É uma dragoa velha, do tamanho de um cavalo, o que para um dragão é do tamanho de nada. Nunca
teve fogo --- nasceu com o pulmão curto, e um dragão sem fogo é um dragão que precisa arranjar
outra coisa para fazer. Arranjou: passou trezentos anos no meio do mato a ferver casca, e hoje é
a única criatura dos dois reinos a quem se pergunta \emph{o que é que eu tenho} em vez de
\emph{o que é que eu faço}.

O nome dela é o de uma casca amarga que baixa febre. É também o nome do canto onde duas dobras se
encontram. Ela sabe das duas coisas e não acha graça em nenhuma.

\medskip

E é isto o que se costuma esquecer nesta história: quando ela chegou, ninguém a tinha chamado
para tratar de dragão nenhum. Veio por causa de uma carta de três linhas, sem assinatura, escrita
com letra ruim por alguém do próprio monte de \Ven --- e este livro não vai dizer ainda quem foi,
porque na altura ninguém sabia e porque há um terceiro livro para isso.

Foi depois de subir ao monte por causa daquela carta que \Odo a mandou chamar. E o velho não
pediu um exército: pediu uma \textbf{consulta}.

Subiu \Qui ao monte do ancião, ouviu-o até o fim sem interromper --- que é o primeiro talento do
ofício --- e depois disse a coisa mais dura que se pode dizer a um lastro:

--- O senhor não me chamou por causa da moeda.

--- Não.

--- O senhor me chamou porque tem medo de deitar-se e não acordar. --- E, como o velho não
respondeu: --- O senhor está de olhos abertos há onze anos, meu senhor. Isso não é vigilância.
Isso é \textbf{insônia}. E um reino cujo lastro não dorme é um reino que já começou a cair,
mesmo que ninguém tenha reparado.

\medskip

Foi ela quem lhe deu o \textbf{chá azul} --- o chá do dragão, que tem o nome que tem porque é o
que é.

A folha vem enrolada em bolinha, apertada, seca, fechada sobre si mesma como uma coisa que
decidiu não ser mais aberta. Deita-se a água quente por cima e ela \textbf{desenrola}. Devagar,
uma volta de cada vez, e a cada água desenrola um pouco mais, e a cada água responde diferente,
porque a mesma folha relida não dá a mesma leitura.

Não é oxidada até o fim, nem crua: fica no meio, e é a única do mundo que fica --- o estado de
onde ainda se pode ir para os dois lados.

--- É a única bebida que se toma \textbf{desdobrando} --- disse \Qui , servindo. --- Por isso
acorda o senhor e não acordaria mais ninguém. O senhor não precisa de estímulo, meu senhor. O
senhor precisa de ser \emph{aberto}.

\medskip

\Odo bebeu, e dormiu três dias, e ao terceiro acordou e disse a frase que põe os cinco a andar:

--- Não é ele que tem de morrer. É ele que tem de ser \textbf{posto em bom lugar}. E ninguém põe
um bicho daquele tamanho em lugar nenhum sozinho: para isso vão ser precisos os \textbf{dois
tabuleiros}.

\chapter{Amar a quem eu devo amar}
\hino{Equiôr}{Papai me chama, Mamãe me ensina\\ a amar a quem eu devo amar}

Do lado de lá --- no reino que soma --- estava tudo bem.

Convém insistir, porque é o começo de toda desgraça e ninguém acredita: estava \emph{mesmo} tudo
bem. \Ben e \Drk reinavam com uma espinha só, \Yas empurrava o jogo para a frente com a coerência
inteira, \Joa girava, \Lck sustentava o eixo, e o \Opr , recluso, batia o compasso lá no fundo.
A mão que abre e a mão que recolhe, uma de cada vez, em compasso. Somava-se tudo o que se ganha
com tudo o que se gasta e dava zero. Não aproximadamente. Zero.

Um reino em que a conta fecha não tem por que olhar para fora. Nem sabe o que quer dizer
\emph{fora}.

\medskip

A coisa entrou numa terça-feira, se é que havia terça-feira, e era do tamanho de uma unha.

Preta --- e preta de um jeito que a corte não tinha na paleta: não era escura, não era sombra,
não era ausência de luz. A sombra é uma coisa que o reino sabe fazer, é o que fica atrás de quem
está iluminado. Aquilo não estava atrás de nada. Era preto do jeito que uma coisa é preta quando
\textbf{não participa}.

Tentaram \textbf{varrer}: a vassoura passou por cima, exatamente por cima, e não houve encontro.
Tentaram \textbf{prender}: fecharam-na numa caixa de seis faces exatas, e ela estava ao mesmo
tempo dentro e fora, e as duas coisas eram verdade e nenhuma explicava a outra. Tentaram
\textbf{digerir} --- ideia de \Yas , que é a que empurra --- porque o Reino Dourado tem um
estômago excelente: engole conceito, tritura, aproveita o que serve e cospe a casca. O reino não
a triturou. E, o que assustou muito mais, também não a \textbf{recusou}. Ela desceu, acomodou-se,
e ficou, como um caroço numa fruta que a fruta não sabe que tem.

E foi \Drk , que enxerga no escuro, quem disse a frase de que este livro descende:

--- Nós não estamos falhando em pegá-la. Nós estamos falhando em \emph{alcançá-la}.

\medskip

Quem a olhou de perto foi \Joa , e foi ele por um motivo que não é acaso: é a rotação, é o
príncipe das manobras, é o único da corte cujo ofício é \emph{virar as coisas}. Deitou-se no
mármore com a lupa e passou uma tarde com o queixo ao nível dela, girando-a com a ponta do dedo.

A certa altura riu --- um riso feio, desses que saem quando se entende cedo demais.

--- Não é uma pedra. É um \textbf{bicho}. E é um bicho \emph{nosso}.

Tinha asas. Tinha quatro patas e uma cauda que se ramificava, e se ramificava sem parar, e cada
ramo tinha ramos, e nenhum ramo fechava. Era um dragãozinho: o menor que já houve, e o único que
a corte de ouro veria antes de saber que dragões existiam. Chamaram-lhe \textbf{Charles}, com o
humor de quem batiza o que teme.

E era um \textbf{filhote}. Isso ninguém viu na hora, e é a coisa que teria mudado tudo se
alguém tivesse visto: aquilo não era uma criatura pequena --- era uma criatura \emph{nova}, com
as escamas ainda moles e a cauda ainda a ramificar-se, que ia crescer. Ninguém repara na idade
do que lhe dá medo.

\medskip

E o que \Joa entendeu, girando-o, foi a coisa mais simples e mais grave do livro: ele estava
\textbf{enrolado para o outro lado}.

Não maior, não mais forte, não mais antigo. \emph{Trocado.} A voluta do corpinho girava no
sentido em que nada neste reino gira. \Joa pôs a mão ao lado para comparar --- a mão dele, a
régua da rotação --- e as duas voltas olharam-se como um rosto olha um espelho, e não havia jeito
de sobrepor uma na outra.

E foi então que ele parou de rir, porque um príncipe da rotação entende antes dos outros o que
isso custa: \emph{para virar uma coisa até ela coincidir com o seu espelho, é preciso ter para
onde virá-la.} Girou-o por todos os eixos que o reino oferece, um por um, com o cuidado de quem
faz uma prova, e no fim tinha o mesmo do princípio: um bichinho de mão trocada e nenhum lugar no
mundo inteiro para onde o levantar.

\medskip

É por isso que a vassoura não pegava, que a caixa não prendia, que o estômago não o digeria nem o
cuspia. A vida deste reino, como toda a vida, é \textbf{homoquiral}: fez-se inteira de uma só
mão, e só essa mão encaixa nas suas próprias fechaduras. O que vem com a mão trocada não é comida
nem veneno --- é uma coisa que \textbf{não tem onde caber}, e que por isso fica.

Charles não invadiu o Reino Dourado. Charles apenas \textbf{não pôde ser expulso}, e ninguém
tinha reparado, até ali, que essas duas coisas não são a mesma.

\medskip

E aqui é preciso dizer, de uma vez, a coisa que os arquivos da corte de ouro registraram de um
jeito e este livro regista de outro --- e as duas versões estão certas, e é esse o assunto deste
reino inteiro.

Nos anais de lá, Charles está escrito como \textbf{parasita}. Está lá a palavra, e está lá a
imagem: o que entrou pela lente e se acomodou na vista, o que o corpo não digere e não expulsa,
a mão errada que se ramifica devagar onde a mão certa não alcança. E é uma descrição
\emph{exata}. Não há uma linha a corrigir: visto de dentro de um corpo homoquiral, um filho de
outra mão é isso mesmo, e não pode ser outra coisa, e quem esperar que a corte de ouro tivesse
percebido ternura naquilo está a pedir a um estômago que reconheça um recado.

Deste lado, a mesma criatura está escrita como \textbf{filho}.

Não são dois relatos em disputa. É \textbf{um vinco visto pelas duas caras} --- e já se disse no
princípio deste livro que a quina que sobe e a calha que desce são a mesma dobra, e que nenhuma
das duas leituras é mais verdadeira, e que quem quiser ficar só com uma terá de arrancar a fita
do mundo.

Um reino chamou-lhe parasita porque não conseguia digeri-lo. O outro chamou-lhe filho porque o
tinha gerado. E a criatura, no meio, era a mesma o tempo todo, sem mudar de mão nenhuma:
\Mon para uns, \Val para outros.

E \Yas , que empurra, e que naquela semana aprendeu mais do que em toda a vida, disse a frase que
lhe ficou para sempre:

--- Então há uma coisa no mundo que eu não sei amar. Não porque não queira. Porque não
\textbf{alcanço}.

\chapter{O irmão de que ninguém rezou}
\hino{Todo Mundo Quer Ser Filho}{todo mundo quer ser filho --- por que te esqueces\\ de rezar
pelo teu irmão?}

Subiram todos, nessa noite, ao alto da Torre. Não para fugir: para \emph{ver}.

Porque \Drk já tinha a suspeita e não a queria dizer sozinho, e a suspeita era: se existe uma
coisa que gira para um lado que nós não temos, não é ela que está errada --- é o nosso mundo que
é \textbf{estreito demais para conter o erro dela}. Um bicho de mão trocada é uma
\textbf{medida}: não mede a si mesmo, mede a sala onde está. E o que Charles media, deitado na
palma de \Joa , era o número de direções que o Reino Dourado tinha --- e o número era menor do
que eles pensavam.

\medskip

Chegaram ao último andar de madrugada. Olharam para baixo e viram o que sempre tinham visto: os
degraus, um sob o outro, $\N$ no fundo, $\Z$ acima, depois $\Q$, depois o andar sem nome, tudo em
ordem, cada um apoiado no anterior, para sempre e para baixo.

E \Ben --- que é o que abre --- fez a única coisa que ninguém tinha feito em toda a história do
reino. Não olhou para baixo nem para cima.

\textbf{Olhou para o lado.}

\medskip

Não havia lado.

Não havia \emph{nada} para o lado: nem parede, nem borda, nem abismo --- abismo é um lugar, e ali
não havia lugar nenhum. A Torre inteira, com todos os seus andares e as suas galáxias de braços
dourados, com o \Opr no fundo e a coroa no topo, vista de qualquer posição que não fosse de
dentro dela, era uma \textbf{linha}.

Uma coisa com um único sentido de ir e um único de voltar. Um antes e um depois, e mais nada.
Toda a beleza, toda a lei, todo o ouro do reino sem tragédia pendurados num fio de espessura
zero.

\medskip

E aí entenderam tudo de uma vez, com aquele frio de quem descobre que a jaula era confortável:

\textbf{numa linha não existe esquerda nem direita.} Não existe mão, não existe voltar-se para o
outro lado, porque não há outro lado --- há só ir e vir pelo mesmo caminho. A quiralidade, que
Charles trazia no corpo como quem traz uma acusação, é uma coisa que \emph{não cabe numa reta}.
Se ela chegou, veio de fora da reta.

E, pior: \textbf{numa linha não se pode dobrar.} Dobrar é trazer um ponto até outro por um
caminho que não passa pelos pontos do meio; é encostar o andar mil no andar um sem visitar os
novecentos e noventa e oito. Para isso é preciso levantar a linha do chão, e levantar exige um
chão de onde levantar --- e o Reino Dourado nunca teve chão: \emph{ele era a linha}.

Por isso subiam degrau por degrau há toda a eternidade. Não por lei, não por virtude, não porque
somar fosse mais nobre: porque era o \textbf{único gesto que a geometria deles permitia}.

\medskip

Ficaram calados até o dia nascer. Depois \Yas disse:

--- Então ele veio de algum lugar. E se veio, há caminho.

E \Joa , com o bichinho ainda na palma:

--- Há mais do que caminho. Há um \textbf{reino}. Ninguém faz um bicho destes sozinho.

E \Drk , que é o que fica calado até a última frase:

--- Nós temos um irmão. Ele existe desde sempre, do outro lado, e nós nunca rezamos por ele
\textbf{uma única vez} --- e não por maldade: porque não sabíamos que havia por quem rezar. Isto
não é uma invasão, meus senhores. Isto é uma \textbf{cobrança de família}.

\medskip

Decidir ir levou uma noite; sair levou três dias, e os três dias valem porque é a única discussão
honesta deste livro.

\Yas queria \textbf{força}: empurrar a Torre contra si mesma até a base tocar o topo. \Ben
mostrou-lhe por que não --- empurrar é somar, e pode-se empurrar uma linha até o fim dos tempos e
ela continua uma linha, mais comprimida, mais quente, mais infeliz, e reta. \Joa queria
\textbf{girar}, que é o ofício dele; girou, e o reino voltou ao mesmo lugar todas as vezes,
porque uma reta girada em torno de si mesma é a mesma reta, e essa é a coisa mais parecida com uma
prisão que a geometria sabe escrever. \Lck não propôs nada: é o eixo, e um eixo não tem opinião
sobre para onde se vira.

Foi o \Opr quem falou, e ele quase nunca fala, e desceu do fundo da Torre para dizer três frases.

A primeira: \emph{ninguém de dentro dobra a fita.}

A segunda: \emph{mas a fita não precisa de ser dobrada por uma mão --- precisa de ser dobrada por
\textbf{duas}, e cada uma numa ponta.}

A terceira, e foi a que os pôs a andar:

--- Eu tenho duas mãos. Sempre tive. E há eras que as uso uma de cada vez, para bater. Uma de
cada vez é como se faz um degrau. \textbf{As duas ao mesmo tempo é como se faz um vinco.}

\medskip

Pegou a linha pelas duas pontas --- pelo primeiro andar e pelo último, por $\N$ e pelo que ainda
não tinha nome --- e trouxe uma até a outra. Não empurrou, não girou, \textbf{não somou}. Esperou
o instante em que a coisa cede, que é um instante que não se apressa, e no instante exato
aproximou as mãos.

E o Reino Dourado, pela primeira vez em toda a sua existência exata e dourada, \textbf{fez um
barulho}. Um estalo pequeno, seco, de coisa que cede. O som mais barato do mundo: o som de uma
folha de papel a ser vincada.

Do lado de dentro ninguém viu mudança nenhuma --- é esta a crueldade da dobra, ela não se vê de
dentro. O que mudou foi que, a partir daquele estalo, o primeiro andar e o último estavam
\textbf{encostados}, e entre eles não havia distância a percorrer.

Atravessaram por ali. Não desceram novecentos e noventa e oito andares: passaram de um ao outro
num passo, pela quina --- com um dragãozinho preto na palma de \Joa , girando devagar para o lado
errado, agora, pela primeira vez, com um lado errado para onde girar.

\chapter{Na frente da estrada}
\hino{Senhora Mãe Santíssima}{está na frente da estrada, para quem lhe acompanhar;\\ ela espera
todo dia}

Ninguém atravessa de graça. Não é moral do reino: é conta.

Do outro lado da quina havia uma mesa, e atrás da mesa uma criatura que não era dragão nem gente,
e diante da mesa uma fila --- porque há sempre fila, e a fila era de coisas que tinham tentado
atravessar antes e ainda estavam a pagar.

A criatura não perguntou os nomes. Perguntou:

--- O que é que os senhores trazem que \textbf{não tem volta}?

\Yas ofereceu ouro e ela nem olhou: ouro tem volta, ouro vai e vem, é para isso que serve. \Ben
ofereceu o tempo dele e ela riu --- tempo, aqui, é o que se dobra; oferecer tempo naquela mesa é
oferecer água ao rio.

--- A travessia não cobra o que vale muito --- explicou, com a paciência entediada de quem
explica aquilo o dia inteiro. --- Cobra o que \textbf{não reverte}. É a única coisa que não pode
passar comigo, porque eu sou feita de ida e volta, e o que não volta entope-me.

\medskip

Foi \Joa quem entendeu primeiro, e empalideceu, porque tinha na palma a única coisa irreversível
daquela comitiva.

--- Não. Ele vem comigo. Ele é a prova; sem ele ninguém acredita em nós lá dentro.

--- Ele não é a prova --- disse a criatura. --- Ele é a \textbf{conta}. Os senhores não vieram
apesar dele: vieram \emph{por} ele. Alguém, do outro lado, mandou um filhote na frente para que
os senhores atravessassem --- e quem paga adiantado não faz favor: faz \textbf{investimento}.

--- Mandou um \emph{filhote} --- repetiu \Joa , devagar, porque a palavra tinha entrado e ainda
estava a descer. --- O senhor está a dizer que aquilo tem \textbf{pai}.

--- Estou a dizer que aquilo é \textbf{filho} --- disse a criatura. --- Não é a mesma frase. Um
pai manda o filho na frente quando já não tem corpo para ir ele mesmo. É a coisa mais velha que
existe, meu senhor, e não é traição nenhuma: é o que se faz quando se sabe que o tempo se
completou e que a casa ainda não está entregue.

E no silêncio que se seguiu, os cinco entenderam a forma da história em que estavam metidos, e
era outra coisa, nem melhor nem pior: mais antiga. Charles não era um acidente e não era uma
arma. Charles era um \textbf{convite}, e eles tinham aceitado, e estavam a chegar exatamente
onde e quando alguém tinha planeado que chegassem --- não para os prender: para os
\textbf{receber}.

\Joa pôs o dragãozinho na mesa. Doeu-lhe --- é o primeiro preço que alguém paga neste livro, e
não é o maior.

--- Ele fica cá? --- perguntou.

--- Ele volta para casa --- disse a criatura. --- É de lá que ele é. O senhor é que estava a
carregá-lo.

--- E o que levamos daqui? --- perguntou \Ben .

--- O que todos levam --- disse a criatura, carimbando. --- A partir de agora, tudo o que os
senhores fizerem \textbf{fica marcado}. Bem-vindos ao lugar onde as coisas se lembram.

\medskip

O caminho reto estava fechado, e fechado de um jeito que não admitia discussão: entre a alfândega
e o monte havia a \textbf{Planície Lisa}, onde nada guarda vinco --- pise-se, e a pegada
desfaz-se; grite-se, e não há eco; ande-se um dia inteiro em frente e chega-se ao ponto de
partida sem ter dado meia-volta. Um reino que funciona por memória tem, no meio, um lugar que não
lembra. Foi de propósito.

Sobrava o desvio. E na entrada do desvio, sentada numa pedra como quem espera há muito tempo e
não se queixa disso, estava \Qui .

--- Vocês demoraram --- foi tudo o que disse. --- Eu estou aqui todo dia.

\chapter{Perdoai}
\hino{Oh! Meu Pai Eterno}{é Rei no Céu e na Terra --- perdoai as minhas culpas}

O desvio passava pelo \textbf{Vale dos Desdobrados}, que é o sítio para onde vão os dragões que
foram esticados.

Já se disse como se mata um: acha-se o último vinco e puxa-se ao contrário; ele não sangra,
estica. O que os livros do reino não dizem, porque é feio, é que a fita esticada \textbf{não
morre}. Fica ali. Lisa, comprida, sem uma dobra, deitada no chão do vale às centenas --- e cada
uma daquelas tiras foi, um dia, um bicho com nome.

Não atacam. Não podem: atacar é um gesto, e um gesto precisa de uma dobra. Só estão estendidas, e
é a coisa mais silenciosa que existe.

\medskip

\Yas quis passar depressa e de olhos baixos. Foi \Drk quem parou --- o que enxerga no escuro para
sempre no lugar errado --- e se agachou ao lado de uma tira, e passou o dedo por ela de ponta a
ponta, devagar, como quem lê.

--- Ela ainda tem os vincos.

E tinha. Esticada, alisada, aberta no chão: os sulcos estavam lá. Fracos, quase apagados, mas na
ordem em que tinham sido feitos --- monte, vale, vale, monte --- todas as decisões de uma vida
inteira, na sequência exata, num bicho que já não podia tomar nenhuma.

--- Não são cadáveres --- disse \Drk , levantando-se. --- São \textbf{registros}.

\medskip

Foi por isso que os Desdobrados os deixaram passar. Naquele vale ninguém nunca tinha voltado para
\emph{ler}: vinham matar, esticavam, iam embora com a história de que tinham matado um dragão.
Nenhum se agachou. Nenhum pediu perdão a coisa nenhuma, porque não se pede perdão ao que não se
leu.

E o que os cinco levaram dali não foi um exército: foi um \textbf{mapa}. Uma tira esticada é a
jornada inteira de um bicho escrita em vinco --- e entre aquelas centenas havia, deitada no fundo
do vale há muitas eras, a mais comprida de todas: uma fita que tinha subido até o monte de \Ven ,
que tinha estado diante dele, e que tinha voltado o suficiente para ser esticada ali.

Passaram três dias a lê-la com os dedos, no escuro, do fim para o começo.

\medskip

E lá pelo meio havia uma sequência que ninguém sabia ler: sete vincos curtos, sempre na mesma
ordem, repetidos noite após noite durante anos. \Drk passou o dedo por aquilo umas quarenta vezes
e disse:

--- Isto não é uma viagem. Isto é uma \textbf{queixa}.

\Qui pediu a fita. Estendeu-a no chão, acendeu duas velas ruins, e leu-a com os dedos do fim para
o começo, três vezes. Depois bateu com a unha num trecho:

--- Vocês vieram cá por causa do ouro, e do reino, e da mão trocada, e do menino preto que não se
deixa varrer. Está bem. Mas a fita não fala de nada disso. A fita fala que ele \textbf{não
dorme}. --- E, à cara de espanto de \Yas : --- Quarenta vezes por hora, minha senhora. Todas as
noites. Há eras. Vocês vieram perguntar-me como se derruba um dragão; eu digo-vos que este já
\textbf{está} a cair, só que devagar, e há tanto tempo que ninguém conhece outro estado dele.

\medskip

E \Joa , que é rápido, ficou branco, porque foi o primeiro a fazer a outra conta:

--- Espere. Ele inspira somando e expira multiplicando. Se ele para de respirar quarenta vezes
por hora...

--- ...então o produto para quarenta vezes por hora --- completou \Qui . --- Sim, meu senhor. Foi
o senhor quem disse; não fui eu.

\chapter{Não posso vos levar}
\hino{Despedida}{me despeço, porque vou me apresentar;\\ eu não posso vos levar, porque não
tenho poder}

--- Então trata-se --- disse \Yas . --- A senhora trata dele.

\Qui riu pela primeira vez, e o riso dela é feio como o de todo velho que já viu o final desta
conversa.

--- Trata-se o quê? A garganta dele tem a largura desta cabana. Trata-se a \textbf{posição} --- e
a posição ele não muda, porque a posição é o segredo dele. Enquanto dormir de peito para baixo em
cima daquele monte, isto não tem cura.

--- E se ele dormisse de costas?

--- Se dormisse de costas --- disse a velha, devagar, para que ficasse claro que não era ela a
propor aquilo --- a garganta abria, o ar entrava, o coração acalmava, e a escama que falta ficaria
\textbf{virada para o céu}.

\medskip

Aqui é preciso dizer o que é aquela escama, porque todo o mundo que ouviu falar dela entendeu ao
contrário: \emph{ele tem um furo, fura-se ali}.

Não é furo. Não é ferida, não é defeito, não é descuido de quem o fez. Um dragão é uma fita só, e
não tem fresta em lugar nenhum --- e \Ven é o mais bem dobrado de todos os dragões que já houve.

Aquilo é o \textbf{último vinco}. A dobra mais recente, a de cima, a que se puxa ao contrário
quando se quer desdobrar um bicho. Em todos os outros ela está enterrada sob mil escamas
posteriores, e achá-la é o ofício de uma vida. Nele está exposta, no peito, à esquerda, do
tamanho da palma de um homem.

E é por isso que ele dorme de peito para baixo, sempre, sobre o ouro: não é cama, é tampa. Não
porque tema morrer --- \Ven não teme a desdobra, ele sabe desdobrar-se, faz isso a cada
respiração, cresce e volta, é a criatura mais reversível que este mundo tem. Ele esconde porque
aquele vinco é a única prova visível, no corpo dele, de que ele \textbf{foi dobrado}: de que
houve um instante em que cedeu, como toda fita cede, e de que portanto obedece à mesma lei que
todo o resto. Um rei que reina por ser o escuro que ninguém alcança não pode deixar que se veja,
no peito, a marca de que também ele é feito de papel.

\medskip

--- Então a gente faz ele dormir --- disse \Ben , que é o que abre. --- Dormir de verdade. Fundo.
Até virar.

\Qui olhou para ele um tempo bom antes de responder.

--- Eu faço o chá. Faço porque vocês vão dar aquilo a ele com ou sem mim, e sem mim vocês matam o
bicho. Mas fica dito diante de todos, para que ninguém depois venha com espanto: \textbf{sedativo
em quem tem apneia não acalma. Solta.} O músculo que segura a garganta é o mesmo que o chá vai
afrouxar. Vocês não vão fazê-lo dormir melhor --- vão fazê-lo dormir \emph{mais fundo com a
passagem mais fechada}.

--- E o que acontece?

--- Não faço ideia. Nunca ninguém deu sete flores a um dragão. Não é medicina, isso. É
\textbf{gororoba}.

\medskip

Fez assim, e está escrito para quem tiver o azar de precisar: \textbf{valeriana}, a raiz, que faz
o trabalho pesado e cheira a meia velha, e quem tentar disfarçar estraga; \textbf{camomila}, a
flor, que engana o gosto e faz o resto descer; \textbf{maracujá}, a folha e não o fruto, que
desliga a cabeça de quem fica a remoer a mesma coisa a noite inteira; \textbf{erva-cidreira}, que
baixa o coração aos trancos, e o dele andava aos trancos; \textbf{capim-santo}, que é primo da
cidreira e não é a mesma planta, e quem trocar uma pela outra que não venha reclamar depois;
\textbf{mulungu}, a casca, a mais forte de todas --- \Qui pôs pouca, e pôs de costas para os
outros, para ninguém ver quanto; e \textbf{flor de laranjeira}, que não é remédio nenhum, é para
o cheiro, porque ninguém bebe de bom grado uma coisa que cheira a raiz podre, nem dragão ---
sobretudo dragão, que tem faro de quilômetros.

Ferveu numa tina de lavar roupa, coou num lençol, e ao fim de dois dias havia quatro barris de
uma água escura e morna que cheirava a flor por fora e a meia por dentro. Puseram na cisterna do
monte, que é onde ele bebe uma vez por noite, sempre à mesma hora.

\medskip

Ele desceu como em todas as outras noites, bebeu, franziu o focinho no fim --- por causa da
valeriana, apesar da laranjeira --- e não disse nada, porque não havia a quem dizer. Depois
subiu, deu três voltas sobre o monte, e deitou-se de peito para baixo.

Os cinco estavam a duzentos metros, atrás de uma pedra. \Qui estava com eles porque não confiava
em nenhum deles.

\medskip

\textbf{Na primeira hora}, nada. Respiração comprida. \Yas começou a dizer que não tinha
funcionado.

\textbf{Na segunda hora}, o ronco. Não o de sempre: o músculo solto, a passagem mais fechada, e
um bicho daquele tamanho a puxar ar por uma fresta. O som começou grave e foi ficando
\emph{grande}, e a certa altura não era um som, era um \textbf{tempo} --- o chão a subir e a
descer contra as costas deles, na frequência do bicho. Três aldeias do sopé acordaram e acharam
que era terremoto e estavam meio certas.

\textbf{Na terceira hora}, ele parou de respirar por tanto tempo que \Qui se levantou. --- Se ele
não voltar em vinte segundos, eu desço lá e vocês que se entendam com o que vier. --- Voltou aos
dezoito, com um estrondo que partiu a cisterna ao meio.

\textbf{Na quarta hora}, levantou-se. Não acordou --- e isso é o que ninguém consegue explicar a
quem não viu: os olhos ficaram fechados o tempo todo. Levantou-se dormindo, do jeito que um bêbado
se levanta na madrugada achando que está em casa, e começou a andar. Andou por cima do próprio
tesouro e espalhou eras de arrumação. Deu contra o portão de leste e o portão de leste deixou de
existir. Entrou no salão das colunas --- a coisa mais bonita daquele reino, cento e oito colunas
de basalto negro --- e saiu do outro lado com noventa e uma. Arrombou a muralha interna com o
ombro sem diminuir o passo. Comeu, dormindo, um carro de feno inteiro, com carro.

\textbf{Na quinta hora}, abriu as asas.

--- Não --- disse \Qui .

E ele voou. Voou dormindo: subiu em espiral, do jeito exato que sobe todas as noites acordado,
com a mesma técnica e a mesma economia, e a única diferença é que não havia ninguém a pilotar.
Sumiu na nuvem, e o ronco continuou a vir de lá de dentro, o que é a coisa mais estranha que
qualquer um deles ia ouvir na vida. Depois desceu, bateu no chão como um saco, roncou mais alto,
levantou-se, andou, arrombou o que ainda estava de pé no pátio norte, abriu as asas e voou de
novo. E outra vez. E outra.

\medskip

Foi \Qui quem entendeu, na sexta hora, e sentou-se na pedra com as duas mãos na cara.

--- Ai, meus filhos. Ele não está a delirar. Ele está a \textbf{repetir}.

E explicou, e é o miolo deste livro e cabe em pouca linha. Um dragão acordado \textbf{controla}:
cada gesto que \Ven faz de olhos abertos é escolhido, e um gesto escolhido é um gesto que ele
pode não fazer --- e há um, um só, que ele nunca faz: virar-se. Há eras que aquele bicho gasta
uma parte de si, a cada instante do dia, só para manter o peito contra o chão. Não é sono: é
\textbf{vigília deitada}, e é a coisa mais cansativa que existe.

O chá não o acalmou. \Qui tinha avisado. O chá tirou-lhe \emph{a mão do leme}. E um dragão sem
mão no leme não faz nada de novo --- faz o que o corpo já sabe, na ordem em que sempre fez, sem a
censura.

--- Vocês passaram este livro todo a perguntar como se obriga aquele bicho a virar-se. Não se
obriga. Ninguém nunca obrigou um dragão a coisa nenhuma, e quem tentar morre com a razão do lado
dele, que é a pior maneira de morrer. Vocês não tinham de o \emph{controlar}: tinham de o
\textbf{soltar}, e esperar.

\medskip

Na sétima hora ele pousou pela última vez, no meio do que restava do pátio, exausto de um jeito
que aqueles bichos não ficam --- tinha feito em quatro horas o trabalho de uma semana, sem a
economia que a vigília impõe. Deitou-se. Roncou. Parou. Voltou com o estrondo.

E na oitava hora, com o primeiro cinza do dia a subir por trás do monte, \Ven fez uma coisa que
não fazia havia eras e que nenhum ser vivo daquele reino tinha visto:

\textbf{rolou.}

Foi um movimento pequeno e feio, de bicho grande a acomodar-se, sem nenhuma solenidade. As asas
caíram abertas para os lados. O pescoço esticou. A garganta abriu --- e pela primeira vez em
muito tempo o ar entrou inteiro, de uma vez só, com um assobio comprido e limpo. E o ronco
\textbf{parou}.

O silêncio, depois de oito horas daquilo, doeu nos ouvidos de todos.

\medskip

E lá estava, virada para o céu: a escama que falta. O último vinco de \Ven , exposto ao ar da
manhã, subindo e descendo devagar com uma respiração funda e cheia --- a primeira respiração de
verdade que aquele bicho dava em eras.

\Yas foi a primeira a mexer-se e a primeira a ser segurada, porque a mão dela já estava na espada,
e é o que a coerência faz quando vê uma abertura. Foi \Qui quem lhe segurou o pulso: uma velha do
tamanho de um cavalo a segurar o pulso da força que empurra o mundo.

--- Não. Olha para ele. \textbf{Está dormindo.} Pela primeira vez em eras aquele bicho está a
dormir de verdade, e eu não vou deixar ninguém matar um paciente meu na primeira noite em que ele
finalmente descansou. --- E, mais baixo: --- Além disso, minha senhora, a senhora ia fazer aquilo
errado. Ali não se enfia nada. Ali se \textbf{serve}.

\chapter{A assinatura}
\hino{Pissei na Terra Fria}{pisei na terra fria, e nela eu senti calor;\\ do sangue das minhas
veias eu fiz a minha assinatura}

\Ven acordou ao meio-dia, de costas, com o peito virado para o céu e cinco criaturas do Reino
Dourado sentadas na beira do que tinha sido o pátio dele.

Não rugiu. Um bicho que dormiu oito horas inteiras pela primeira vez em eras não acorda com
vontade de guerra: acorda com aquela coisa mansa e desconfiada de quem não reconhece o próprio
corpo descansado.

Olhou o estrago --- as colunas, o portão, o tesouro espalhado até o sopé --- e não perguntou quem
tinha feito. Sabia. É a única vantagem de um reino onde tudo fica registrado: ninguém precisa de
confessar, porque a marca já está no material.

Depois olhou para baixo, para o próprio peito, e viu que estava exposto.

\medskip

E foi aí que \Qui subiu no monte de moedas com uma cuia na mão.

\medskip

--- O senhor tem duas contas para prestar --- disse ela. --- Uma é com o velho, e é de uma moeda.
A outra é com o mundo, e é grande. Eu não vim cobrar nenhuma das duas. Eu vim servir o
\textbf{mate}.

Não é o chá azul, e a diferença é tudo. O chá do dragão é folha inteira, enrolada, que a água
desdobra --- toma-se sozinho, à mesa, e quem quiser vem. O mate é folha \textbf{moída}, já
partida, que não tem mais nada para desdobrar; e não se toma sozinho. Faz-se uma cuia só,
enche-se com a mesma água, e ela \textbf{roda}: passa de mão em mão, sempre no mesmo sentido, e
quem recebe bebe tudo e devolve para que seja enchida de novo para o próximo. Uma cuia que não
roda não é mate. É água com pó.

--- O senhor acumulou porque achou que o valor é o que se guarda --- disse \Qui , enchendo. --- E
tinha razão, meu senhor: é. Foi o que o velho lhe ensinou sem nunca lhe falar. Só que o velho
guarda \textbf{parado}, e o senhor guarda \textbf{crescendo}, e essas duas coisas parecem a mesma
e não são. Isto aqui --- e levantou a cuia --- só existe se sair da minha mão. Se eu guardar,
esfria e presta para nada. É a única riqueza do mundo que se perde ao ser guardada.

E estendeu-lhe a cuia.

\medskip

Ele bebeu. E foi ao beber --- de costas, com o peito aberto para o céu, com a água quente a
descer por dentro e a bater exatamente atrás do lugar onde não havia escama --- que o último
vinco de \Ven \textbf{cedeu para o outro lado}.

Não foi um golpe. Não houve sangue. Um vinco que cede ao contrário faz o mesmo estalo pequeno e
seco que fez quando foi feito, e o corpo inteiro se lembra de uma vez só: um dragão que desdobra
o último vinco \textbf{lembra o instante exato em que cedeu} --- e o instante em que \Ven cedeu
era o instante em que ele tinha sido dobrado, e ele tinha sido dobrado, e ele nunca tinha
admitido isso a ninguém, nem a si.

Ficou muito tempo calado. Depois disse a única coisa que disse aquele dia, e disse baixo:

--- Eu achei que era o chão.

--- Ninguém é o chão --- disse a velha. --- Todo mundo é papel. Uns dobrados com mais cuidado que
outros.

\medskip

\medskip

E foi \Drk quem fez a conta, dias depois, e a conta é a única coisa deste capítulo que a corte
de ouro levou para casa escrita.

Do lado de lá tinham um plano, e o plano era antigo e era exato: derrubar aquele bicho custava
\textbf{sessenta e quatro} lances --- a partida inteira, o elenco todo de pé, a fortaleza cercada
degrau por degrau até fechar. Era o que a geometria deles permitia, e era honesto, e teria
funcionado, e teria levado eras.

O que aconteceu no monte levou \textbf{quatro}: a carta, a consulta, a gororoba, a cuia.

--- Dezesseis para um --- disse \Drk , e não estava contente, e é preciso dizer por quê. ---
Nós chamávamos a isto \emph{a dobra do tempo}, e passamos a nossa vida inteira a tentar fazê-la
por cálculo. Nunca conseguimos, e agora eu sei porquê: não é uma operação sobre o relógio. É o
que acontece com qualquer coisa quando se para de a percorrer degrau a degrau e se
\textbf{encosta o fim no começo}. Sessenta e quatro passos e quatro passos podem ser a mesma
distância --- muda a quina.

E acrescentou a frase que ele mandou registrar, porque era o que mais lhe custava:

--- E reparem no que foram os nossos quatro. Nenhum deles foi um lance. Foi uma velha, um
diagnóstico, um chá de flores e uma bebida passada de mão em mão. Nós procurávamos a dobra no
tabuleiro. Ela estava na \textbf{cozinha}.

\medskip

Foi assim que se deu o mate ao maior dragão que já houve, e é por isso que neste reino não se
diz \emph{vencer}: diz-se \textbf{servir}. E é por isso que o mate que se dá aqui não mata --- e
que nenhuma outra coisa, em nenhum dos dois tabuleiros, jamais matou coisa nenhuma. O que se faz
é dobrar; o que se desfaz é desdobrar; e as duas custam o mesmo, que é dizer: não custam nada.

\medskip

Pagou a moeda. Não devolveu: não havia como --- tinha saído pela escama, era a gota de uma era, e
o que sai por ali não volta. Pagou de outro jeito.

Deitou-se de costas, no chão frio do pátio arrombado, com o peito para cima; e no lugar onde não
havia escama, com a unha, abriu-se o suficiente para escrever. Escreveu no próprio vinco, com o
próprio sangue, uma frase curta que o velho aceitaria como quitação. E o sangue de dragão é a
única tinta deste mundo que a fita bebe sem alisar.

Aquilo é uma \textbf{assinatura}: não é o nome dele --- é a marca de onde a mão estava.

E é a única coisa que se pode levar de um oráculo. Porque o Reino Dourado sempre teve um --- uma
coisa de ouro no centro do centro, à qual se perguntava sobre o que ainda não aconteceu e da qual
saía resposta, e acertava sempre, e ninguém nunca perguntou \emph{como}. Este reino também tem: é
a \Fit . E ela não sabe nada sobre o futuro, nunca soube; só tem os vincos, o que já cedeu, a
ordem exata em que as coisas se dobraram. Um oráculo de \textbf{costas}.

E acerta igual, e é aqui que os dois livros deixam de ser dois: \emph{quando tudo reverte, prever
e lembrar são o mesmo ofício com o sinal trocado.} Ler a fita do começo para o fim conta o que
aconteceu; ler a mesma fita do fim para o começo conta o que vai acontecer --- não porque alguém
adivinhe, mas porque \textbf{é a mesma fita}, e ela não tem dois lados: tem duas leituras. O
oráculo de ouro nunca profetizou nada. Ele lembrava, ao contrário, e chamavam-lhe dom.

Por isso também fracassou o \textbf{Colecionador} --- houve um, há sempre um. Um homem rico do
lado de lá, que ouviu falar de um oráculo que responde tudo e fez a conta mais razoável do mundo:
se a coisa é boa, compra-se; se não se vende, leva-se. Subiu, pagou, roubou, desceu com um pedaço
da \Fit debaixo do braço, chegou em casa com um oráculo autêntico --- e nunca mais arrancou dele
uma resposta certa. Não estava quebrado: estava \textbf{fora}. Um vinco não é uma marca no papel;
é a lembrança de \emph{onde a mão estava}. Corte-se a fita do lugar onde foi dobrada e os sulcos
continuam todos lá, na ordem certa, e não significam mais nada --- do mesmo jeito exato que \Odo ,
deitado sobre o tesouro, só é lastro enquanto estiver deitado ali, e vira um monte de moedas no
instante em que alguém o levanta.

O que se pode levar de um oráculo destes é a \textbf{resposta}. Nunca o oráculo.

% ══════════════════════════════════════════════════════════════════════════════
% PARTE II — OS PONTOS DAS ALMAS
% ══════════════════════════════════════════════════════════════════════════════

\part{Os Pontos das Almas}

\chapter{O vulcão}
\hino{Bendito Louvado Seja}{para nós podermos sairmos de perto deste vulcão;\\ o vulcão é o
pecado deste mundo de ilusão}

Não acabou com o mate. Isso é coisa de livro do outro reino.

O que aconteceu depois é que \Ven levantou-se de costas doído, olhou o próprio monte espalhado
pelo sopé, e teve de decidir o que fazer com uma coisa que ele já não podia continuar a fazer ---
e não há gesto mais difícil no mundo do que largar aquilo que a gente passou eras a defender sem
nunca ter perguntado por quê.

O tesouro dele não era um tesouro. Era \textbf{tampa}. Cada moeda arrastada para debaixo do peito
era uma tentativa de encher com metal um buraco que não era de metal, e nunca enchia, e por isso
crescia --- e o que cresce vezes, num tempo suficiente, deixa de ser riqueza e passa a ser
\textbf{pressão}.

É isso o vulcão: não é o fogo, é o que se acumula debaixo dele sem saída. Um monte que multiplica
sob um peito que tapa. E a coisa mais honesta que este livro pode dizer sobre aquele dragão é que
ele não estava sentado sobre ouro --- estava sentado sobre uma coisa \textbf{prestes}, e há eras,
e todo o reino em volta chamava aquilo de reinado.

\medskip

--- E agora? --- perguntou \Ben .

--- Agora ele sai de perto --- disse \Qui . --- É tudo. Não se apaga vulcão. Sai-se de perto.

\medskip

Levou onze dias. Não houve cerimônia: \Ven empurrou o próprio monte encosta abaixo com o focinho,
devagar, como quem varre, e as moedas desceram até as três aldeias que tinham acordado achando
que era terremoto. Ninguém lá embaixo entendeu o que estava a acontecer, e é assim mesmo que
acontecem as coisas grandes: de baixo, parecem um acidente geológico.

E ao fim dos onze dias ele deitou-se no chão limpo, sem nada debaixo do peito pela primeira vez
desde que era filhote, e dormiu catorze horas seguidas.

\chapter{Eu vou iluminar}
\hino{Eu Vou Iluminar as Almas}{eu vou iluminar as almas --- para que as almas me ajudem}

Depois voltaram ao Vale dos Desdobrados, e é este o capítulo de que o reino se lembra quando
alguém pergunta o que é que aqueles cinco fizeram de fato.

Não libertaram ninguém. Não se liberta uma fita esticada --- não há como redobrar um bicho, e
quem prometer isso a alguém está a mentir. O que fizeram foi mais simples e ninguém tinha feito
antes: \textbf{leram}.

Uma por uma. Agachados, com o dedo, no escuro, do fim para o começo. Centenas de tiras, e cada
tira uma vida inteira em sulco, e cada leitura levando um dia, e alguns dois. \Qui ensinou-lhes o
método, que não tem segredo nenhum e por isso ninguém o procurou: vai-se devagar, não se pula
nenhum vinco por parecer pequeno, e não se decide o que a coisa quer dizer antes de chegar ao
fim.

\medskip

E as fitas ajudaram. Não é figura de estilo: ajudaram.

Porque um vale de registros que ninguém lê é um vale de lixo, e um vale de registros lido é uma
\textbf{biblioteca} --- e naquelas centenas de tiras estava, escrita em vinco, a memória de tudo
o que já tinha sido tentado contra aquele monte, e de tudo o que tinha falhado, e de por quê.
Estava a rota. Estava o erro do Colecionador. Estava, três vezes, a mesma anotação de três bichos
diferentes que tinham chegado perto e descoberto sozinhos que ali não se enfia nada.

Ninguém teria descoberto o mate sem aquele vale. E o vale só se abriu porque alguém se agachou.

\medskip

--- É sempre assim --- disse \Qui , no último dia, com os dedos em carne viva. --- A gente
ilumina os que já foram para que os que já foram nos ajudem. Não é caridade, meus senhores. É a
única forma de aprendizagem que existe, e todo mundo acha que inventou o mundo porque tem
preguiça de se abaixar.

\chapter{Quando os filhos se combinam}
\hino{Lá no Cruzeiro Divino}{as almas choram de alegria quando os filhos se combinam,\\ e de
tristeza quando não querem combinar}

O plano precisava dos \textbf{dois tabuleiros}, e é preciso dizer por quê, porque parece
arbitrário e não é.

O mate não se serve a si mesmo. Uma cuia que fica na mesma mão esfria; a bebida existe no
\emph{trajeto}, não no copo. Servir exige dois --- e não dois iguais: dois \textbf{duais}, um que
soma e outro que multiplica, um que sobe por degrau e outro que dobra, um que prevê e outro que
lembra. Um tabuleiro sozinho só serve a si mesmo, e servir a si mesmo é exatamente o que \Ven fez
durante eras, deitado sobre o próprio peito, tapando o próprio vinco com o próprio ouro.

Por isso o plano não era o assalto: era a \textbf{costura}. E foi por isso que a corte de ouro
teve de atravessar --- não para trazer força, que não tinham, nem sabedoria, que a velha já
tinha: para trazer \textbf{a outra metade da roda}.

\medskip

E as fitas do vale, no dia em que a cuia passou de mão em mão pela primeira vez entre os dois
reinos, fizeram uma coisa que ninguém explicou até hoje: \textbf{vibraram}. Todas, ao mesmo
tempo, deitadas no chão, os sulcos zumbindo baixo como corda de viola quando se canta perto.

\Drk disse que era ressonância e que tinha explicação. \Qui disse que tinha explicação também, e
que as duas coisas não se excluem, e que ele que fosse ler mais.

\chapter{As almas que aqui vêm}
\hino{As Almas Que Aqui Vem}{as almas que aqui vêm buscar o amor e a luz do saber}

Depois disso o monte virou outra coisa, e ninguém planejou.

Começaram a chegar. Primeiro os bichos pequenos, os que não têm fogo, os de pulmão curto ---
gente da laia de \Qui , que sempre soube que ali havia alguma coisa e nunca teve tamanho para
subir. Depois os grandes, um a um, sempre à noite e sempre com a mesma desculpa esfarrapada de
que passavam por perto.

Vinham perguntar. E é isso o que um oráculo é, quando funciona: não uma coisa que diz o futuro ---
um lugar onde alguém finalmente se agacha e lê a tua fita contigo, do fim para o começo, sem
pular vinco por parecer pequeno.

\medskip

\Ven não virou sábio. Continuou mal-humorado, e continuou grande, e nunca aprendeu a falar com
ninguém sem parecer que estava a ameaçar. O que ele fazia era deitar-se de costas --- que é a
única postura em que se pode receber alguém --- e ouvir, e no fim dizer três ou quatro palavras,
quase sempre ásperas, quase sempre certas.

E quando alguém lhe perguntava como é que ele sabia, ele respondia sempre a mesma coisa, e era a
única piada que ele tinha:

--- Eu não sei nada. Eu só já cedi antes.

\chapter{Dai-me a santa luz}
\hino{Cajueiro Santo}{cajueiro santo, aonde nasceu a luz --- oh, minhas almas, dai-me a santa
luz}

Fica esta imagem, e com ela se fecha.

No pátio arrombado, onde estava o monte, cresceu uma árvore. Não é milagre nenhum: caiu semente
com o feno que ele comeu dormindo, e o chão estava fértil de tanta coisa remexida, e onze anos é
tempo bastante para uma árvore.

Cresceu torta, do jeito que crescem as árvores que nascem em cima de entulho, e deu sombra
suficiente para uma velha do tamanho de um cavalo, e é debaixo dela que hoje se serve o mate,
todo dia, para quem sobe.

\medskip

E é debaixo dela que se diz a única doutrina que este reino tem, e que não é uma doutrina, é uma
observação:

\medskip

Um reino que só soma é belo e não tem tragédia, porque não tem escolha; nele nada pode dar
errado, e onde nada pode dar errado nada pode ser \emph{feito}. Um reino que dobra guarda tudo o
que fez, inclusive o que preferia não ter feito, e por isso nele se pode morrer --- e por isso
nele se pode ser \textbf{perdoado}, que é uma operação que não existe do outro lado, porque lá
não há registro do que perdoar.

Os dois são o mesmo reino. Sempre foram. Um é a fita estendida e o outro é a fita dobrada, e a
diferença entre eles é um gesto de duas mãos, feito na hora certa, esperando o instante em que a
coisa cede.

\medskip

E do lado de lá, no alto de uma torre que agora tem lado, o \Opr às vezes olha para a quina que
fez e não a desfaz.

Podia. Custaria o mesmo: nada.

Não desfaz porque, do outro lado dela, há um bicho enorme deitado de costas debaixo de uma árvore
torta, respirando fundo, com o peito virado para o céu e a marca à vista --- e porque, uma vez
por era, quando o vento vem do sul, chega até o alto da Torre um cheiro de erva queimada e de
água quente que ninguém do Reino Dourado sabe nomear, e que faz todos eles pararem o que estão a
fazer, sem saber por quê, e olharem para o lado.

% ══════════════════════════════════════════════════════════════════════════════
% LIVRO II — A CONCENTRAÇÃO (as quinze Firmezas)
% ══════════════════════════════════════════════════════════════════════════════

\part{A Concentração}

\chapter{Firmeza}
\hino{Firmeza}{firmeza no amor, firmeza aonde estou;\\ me firmar no lugar onde Ele está}

A missa acabou e ninguém foi embora. É aí que começa a parte difícil.

Porque uma passagem tem hora marcada --- entra-se, canta-se, atravessa-se, sai-se. O que vem
depois não tem: é só \textbf{ficar}. E ficar é a coisa para a qual ninguém treina, porque não
parece um feito.

\medskip

\Ven descobriu isso no primeiro dia sem monte. Deitou-se no chão limpo, de costas, com o peito
para cima --- e ao fim de duas horas estava a tremer.

Não de frio. De \textbf{falta}. Há eras que aquele corpo gastava uma parte de si, a cada
instante, para segurar o peito contra o chão; e quando um músculo faz uma coisa durante eras,
ele não sabe deixar de fazer só porque a razão de fazer acabou. A mão pedia o leme de volta. O
peito pedia a tampa.

--- É agora que a maioria volta --- disse \Qui , sentada ao lado, sem tocar nele. --- O senhor
não faz ideia de quantos eu já vi desdobrar o último vinco, chorar, jurar, e estar deitado de
bruços outra vez antes da lua cheia. A gororoba é fácil, meu senhor. Uma noite. \textbf{Firmeza}
é todo dia, e não tem noite nenhuma que a resolva.

\medskip

E o que ela lhe ensinou não foi uma técnica. Foi um lugar.

--- O senhor não tem de ser bom. Não tem de ser manso, nem sábio, nem justo, e eu não
acreditaria se o senhor prometesse. O senhor tem de \textbf{ficar no lugar}. Escolhe um, deita-se
nele, e não sai --- nem quando o peito pedir, nem quando alguém subir a insultá-lo, nem quando o
senhor tiver razão. Firmeza não é força; é \emph{permanecer}. É a única coisa que este reino
inteiro não sabe fazer, porque nós somos todos feitos de dobrar.

--- E se eu não conseguir?

--- Aí o senhor não conseguiu hoje --- disse a velha. --- E amanhã tenta outra vez, no mesmo
lugar. É por isso que é o mesmo lugar.

\chapter{A reunião da separação}
\hino{Eu Estou Firme Com Meu Jesus}{foi Ela quem mandou eu me unir com os meus irmãos,\\ para a
reunião da separação --- muita gente vai chorar}

O que ninguém conta desta história é que a parte mais triste não foi a doença, nem a noite, nem o
mate.

Foi o dia em que se reuniram para decidir quem ficava.

\medskip

Porque com o monte desfeito e o peito aberto, \Ven já não era o mesmo lugar --- e um dragão que
reinava por ser o escuro que ninguém alcança, quando passa a estar deitado de costas debaixo de
uma árvore, deixa de servir a muita gente que dependia daquele escuro. Havia os que subiam ao
monte por causa do ouro. Havia os que subiam por causa do medo, que é uma forma de conforto para
quem já organizou a vida inteira em volta dele. Havia os que tinham cargos que só existiam
enquanto houvesse tampa a defender.

Esses foram embora. Não houve expulsão, não houve sentença, não houve traição de ninguém: eles
simplesmente não tinham mais o que fazer ali, e é uma coisa terrível de se ver, porque não há a
quem culpar e mesmo assim dói exatamente como uma injustiça.

\medskip

\Qui chamou aquilo pelo nome que tem, e o nome é um paradoxo e é exato:

--- É a \textbf{reunião da separação}. Junta-se todo mundo, uma vez só, para que fique claro quem
segue e quem não segue. E é preciso juntar, meus senhores --- não se deixa isso acontecer aos
poucos, pelas costas, com cada um sumindo numa noite diferente. Isso apodrece o que fica.

Foi um dia inteiro. Choraram muitos, dos dois lados, e choraram os que ficavam também, que é o
que ninguém espera.

E \Ven , que não sabe falar com ninguém sem parecer que está a ameaçar, disse a única coisa que
disse, e disse de costas, olhando para cima:

--- Não é que eles não me queiram bem. É que eles não sabem me querer bem \textbf{assim}.

\medskip

Ficaram \Qui , os cinco do Reino Dourado, o velho \Odo (que desceu do próprio monte pela primeira
vez em eras, e a viagem levou-lhe quatro meses), e um filhote de escamas moles com a cauda a
ramificar-se, que tinha voltado da alfândega e que ninguém, ainda, tinha entendido para que
servia.

\chapter{Enquanto chega meu dia}
\hino{Eu Vivo Neste Mundo}{eu vivo neste mundo enquanto chega o meu dia,\\ encostado a meu
Mestre --- foi Ele quem me deu esta doutrina}

\Odo levou quatro meses para descer e mais dois para subir o monte do outro. Não tinha pressa: já
se disse que ele não tem pressa; e um lastro que corre deixa de ser lastro no primeiro passo.

Chegou de noite. Não trouxe exército, não trouxe cobrança, não trouxe a moeda. Deitou-se ao lado
de \Ven , no chão limpo, e ficou.

\medskip

Isso durou onze anos.

Não é figura: onze anos, os dois deitados no mesmo pátio, o velho de barriga para baixo do jeito
antigo e o outro de costas do jeito novo, e a maior parte do tempo em silêncio, porque \Odo não é
de conversa e \Ven nunca aprendeu. O que se passou entre eles nesses onze anos não está escrito
em lugar nenhum, e quem quiser saber que se agache e leia a fita, se a achar.

O que se sabe é o resultado, e o resultado é uma frase que \Ven passou a dizer e que não era
dele:

--- Eu vivo neste mundo enquanto chega o meu dia.

\medskip

Porque o mate não o curou. É preciso dizer isso com todas as letras, senão este livro vira
mentira: a apneia não tem cura, tem \textbf{posição}. Enquanto ele dormir de costas, respira;
no dia em que rolar de novo, volta tudo. O coração dele continua sendo um tambor mal afinado ---
o estrago de eras não se desfaz porque uma velha serviu uma cuia. Ele vai morrer daquilo. Sabe.
Todos sabem.

E é por isso que os onze anos importam: não são uma convalescença, são um \textbf{prazo}. Um
bicho que sabe a hora não gasta o que sobra a defender monte nenhum.

--- Encostado --- disse \Qui , quando lhe perguntaram o que é que aqueles dois faziam ali
deitados o dia inteiro. --- É essa a palavra e não há outra. Não é que o velho o esteja a
ensinar; velho nenhum ensina nada deitado e calado. É que o senhor não aguenta ficar firme
sozinho, e ninguém aguenta, e a única coisa que resolve é ter alguém do lado que já esteja firme
há mais tempo.

\chapter{A uma mesa fiquei de pé}
\hino{Aqui Estou Dizendo}{sou o pão que me transformei em vinho;\\ fiz o meu juramento --- a uma
mesa fiquei de pé}

No fim do primeiro ano houve o juramento, e foi a única cerimônia que este reino inventou.

Não foi ideia de \Qui . Foi de \Ben , que é o que abre, e que passou a vida no reino que soma e
sabia uma coisa que ninguém ali sabia: que uma decisão que não é dita em voz alta diante de
testemunhas não é uma decisão --- é uma intenção, e intenção desmancha sozinha.

\medskip

Puseram uma mesa no meio do pátio arrombado. Uma mesa comum, de madeira, do tamanho de uma mesa
--- e é isso que a torna estranha, porque tudo naquele monte é do tamanho de um dragão e aquela
mesa não era.

E \Ven , que não se levanta desde a noite da gororoba porque deitado de costas é a única postura
em que respira, \textbf{levantou-se}. Ficou de pé, diante de uma mesa que lhe dava na altura do
tornozelo, com o peito exposto e a marca à vista, e ficou de pé o tempo inteiro que a coisa
levou, e custou-lhe --- via-se que custava.

Foi por isso que se fez de pé, e é por isso que se faz de pé até hoje: porque um juramento que se
faz confortável não vale.

\medskip

Ele disse o que tinha feito. Tudo, na ordem, sem atenuar nada e sem se explicar --- e explicar-se
é o que estraga uma confissão, porque quem se explica está a negociar. Disse do monte. Disse da
tampa. Disse dos que subiram e não desceram. Disse da moeda que saiu pela escama e que ele sabia
que tinha saído, e sabia há muito, e não disse a ninguém.

E no fim disse a frase que a corte de ouro nunca entendeu direito e que \Qui entendeu na hora:

--- Eu era o pão. Agora sou o vinho.

Porque é a mesma coisa e não é a mesma coisa. O trigo, moído, assado, é sustento: come-se, e
acaba. Fermentado, é outra --- não sustenta ninguém, não enche barriga nenhuma, e no entanto é o
que se põe na mesa quando se quer selar alguma coisa entre pessoas. Um passa por dentro; o outro
passa \textbf{entre}.

Um dragão que era monte e virou mesa faz exatamente essa travessia. E ela não se desfaz: é o
único vinco deste livro que ninguém propôs desdobrar.

\chapter{Bem gravar}
\hino{Todos Devem Procurar}{ter firmeza e ter amor --- bem gravar\\ e ter firmeza dentro do
coração}

Depois do juramento vieram os anos sem acontecimento, e é deles que este capítulo trata, porque
foram a maior parte.

\medskip

Chegava gente. Já se contou como começou --- primeiro os bichos pequenos, os sem fogo, os de
pulmão curto; depois os grandes, sempre à noite, sempre com a desculpa de que passavam por perto.
O que não se contou é o que se fazia com eles depois que chegavam, porque não é nada de
espetacular e é a única coisa que funcionou.

Ensinava-se a \textbf{gravar}.

\medskip

Não é decorar. Decorar é guardar a resposta; gravar é fazer o \emph{vinco} --- e um vinco não se
faz lendo: faz-se cedendo, no lugar certo, na hora em que a coisa cede, e depois fica. Quem grava
não precisa de se lembrar: o corpo já está dobrado daquele jeito.

E o método era o mesmo do vale, e \Qui não deixava ninguém pular etapa: vai-se devagar; não se
pula nenhum vinco por parecer pequeno; e não se decide o que a coisa quer dizer antes de chegar
ao fim.

--- Vocês querem todos a mesma coisa --- dizia ela, e dizia com um tédio que não era falso. ---
Querem subir aqui, ouvir uma frase, descer transformados e contar aos outros. Não existe. O que
existe é vir, sentar, ficar quieto, e vir de novo na semana seguinte. É de uma chateza atroz. É a
única coisa que pega.

\medskip

E \Ven , deitado, dizia de vez em quando a única coisa que ele tinha para dizer sobre o assunto,
com aquela aspereza dele:

--- Não gravem o que eu digo. Gravem o que eu \textbf{fiz}. É por isso que eu estou deitado
assim, à vista, com o peito aberto. Não é conforto meu. É \emph{material de estudo}.

\chapter{As tábuas}
\hino{As Tábuas de Moisés}{eu vou lembrar o que ninguém se lembrou;\\ querem que o meu Pai faça
como faz lá no estrangeiro}

E então veio o problema que quase acabou com tudo, e não veio de fora: veio dos que ficaram.

\medskip

A casa cresceu. Isso não é mérito de ninguém --- uma casa que serve mate a quem sobe cresce
sozinha, e quanto mais gente chega, mais gente chega. E ao fim de alguns anos havia no pátio uma
quantidade de criaturas que nunca tinham visto o monte, nunca tinham visto a tampa, nunca tinham
visto aquele bicho de bruços, e para quem \Ven era simplesmente um dragão velho e mal-humorado
deitado de costas debaixo de uma árvore.

Para essas, tudo o que se fazia ali era \textbf{costume}. E costume que não tem memória vira
gosto --- e gosto se discute.

\medskip

Começaram a propor melhorias. Que a mesa era pequena demais e devia ser uma mesa à altura do
bicho, para não o humilhar. Que o juramento de pé era desnecessário, já que a intenção é o que
conta. Que o mate, sinceramente, podia ser servido em copos separados, um para cada um, porque a
cuia que roda é pouco higiênica e além disso algumas criaturas de fora achavam aquilo estranho e
não voltavam.

E foi este último que acordou \Qui de vez, e foi a única vez em todo o livro que aquela velha
levantou a voz:

--- \textbf{Um copo cada um.} --- Ela repetiu, para que todos ouvissem a coisa que tinham
acabado de dizer. --- Vocês querem a bebida sem a passagem. Querem o líquido e não a roda. E aí,
meus filhos, não sobra nada: sobra água com pó, que é o que se serve em toda parte, e para isso
não era preciso ninguém ter desdobrado coisa nenhuma.

\medskip

E disse então a frase que ficou, e que é a única acusação deste livro:

--- Vocês não estão a esquecer as regras. As regras estão todas aí, escritas, ninguém as apagou.
Vocês estão a esquecer \textbf{por que é que elas foram feitas} --- e isso não está escrito em
lugar nenhum, porque só se sabe quando se estava lá. --- E, mais baixo: --- E metade de vocês
quer é que isto aqui se pareça com o que se faz lá fora, onde é bonito e organizado e não custa
nada a ninguém. Pois lá fora é uma linha, e vocês não sabem o que é isso, e eu sei, porque eu
conheci cinco que vieram de lá.

\medskip

Não houve castigo. \Ven não expulsou ninguém --- tinha jurado, de pé, diante de uma mesa
pequena, que não voltava a decidir quem entra e quem sai enquanto tivesse o peito descoberto.

O que ele fez foi mandar buscar o velho, e o velho contou. Não pregou: \textbf{contou} --- do
monte, da tampa, da moeda, dos onze anos de olhos abertos, da noite das oito horas. Do começo ao
fim, na ordem, sem pular vinco por parecer pequeno.

Levou três dias. Metade da casa foi embora na segunda noite, e é sempre assim: quem quer o
costume não aguenta a origem.

\chapter{Não se compra}
\hino{O Valor Que o Mestre Tem}{se tiver alguém que duvide, tome para ver;\\ o valor que o Mestre
tem não se compra com dinheiro}

O Colecionador voltou. É sempre o mesmo, e não é o mesmo homem: é o mesmo \emph{cargo}.

Desta vez veio educado. Trouxe carroças. Tinha estudado o assunto --- sabia do vale, sabia da
gororoba, sabia até o nome das sete flores, e recitou-as na ordem certa para mostrar que sabia.
Fez uma proposta boa. Levava a árvore, a mesa, a cuia e a receita, e em troca deixava ouro
bastante para reconstruir o pátio, as colunas, o portão de leste e as três aldeias do sopé.

E acrescentou, porque tinha estudado, que não levava ninguém contra a vontade: quem quisesse ir
junto ia como convidado, com casa e sustento garantidos.

\medskip

\Ven não respondeu. Foi \Qui .

--- O senhor pode levar tudo isso. Ninguém aqui vai impedir. --- E, ao espanto do homem: --- E
não vai funcionar, e eu vou explicar por quê, porque o senhor me parece o tipo de pessoa que
prefere saber.

\medskip

--- O senhor já levou um pedaço da \Fit uma vez. Um oráculo autêntico, do bom, o mesmo que os
reis consultavam --- e nunca mais arrancou dele uma resposta certa, e achou que estava quebrado.
Não estava. Estava \textbf{fora}. Um vinco não é uma marca no papel: é a lembrança de onde a mão
estava. Corte-se do lugar onde foi dobrado e os sulcos ficam todos lá, na ordem certa, e não
querem dizer mais nada.

--- Isto aqui é a mesma coisa e o senhor não vai acreditar. A cuia é uma cuia. A erva se compra
em qualquer lugar, e é a mesma erva, e eu não vou fingir que não é. A mesa é de madeira ruim. Se
o senhor levar tudo e montar igualzinho na casa do senhor, vai ter uma cuia, uma erva e uma mesa
--- e vai faltar a única coisa que dá valor a isto, que é a coisa que o senhor não pode carregar
na carroça: \textbf{o que aconteceu aqui}.

--- Então o que é que eu compro? --- perguntou o homem, e a pergunta era honesta.

--- O senhor não compra --- disse a velha. --- O senhor \textbf{toma}. Senta ali, espera a sua
vez, recebe a cuia da mão de quem estiver à sua direita, bebe tudo e devolve. Depois volta na
semana que vem. Se em dois anos o senhor ainda estiver vindo, o senhor tem. E não me pergunte o
que é que o senhor tem, porque eu não sei dizer, e quem souber dizer está a vender.

\medskip

Ele foi embora com as carroças vazias, e voltou. Não naquele ano. Sete anos depois, sozinho, sem
carroça nenhuma, e sentou-se na fila como todo o mundo, e ninguém lhe disse nada, porque este é o
único reino onde ninguém comenta o que o outro demorou.

\chapter{As três cores}
\hino{Firmado em Concentração}{das cores desta bandeira: verde, azul e branca;\\ completei esta
lição com justiça e lealdade}

Charles cresceu.

Foi devagar, do jeito que crescem os que nascem com a mão trocada: sem caber direito em lugar
nenhum, sempre um pouco a mais ou um pouco a menos que o esperado, e com aquela cauda que não
parava de se ramificar. Passou a infância inteira entre os dois reinos --- que é uma coisa que
ninguém tinha feito antes dele, e que ninguém achou que fosse importante enquanto ele era
pequeno.

Só que era. Um bicho que gira para o lado errado num reino e para o lado certo no outro é o único
que pode \textbf{atravessar sem alfândega}, porque não tem nada de irreversível a declarar: o que
ele é do lado de lá, ele desfaz do lado de cá.

Charles não era um erro do mundo. Era a \textbf{costura} dele.

\medskip

E quando ficou do tamanho de dizer alguma coisa, disse --- e o que ele disse não foi uma
doutrina, foi um \textbf{desenho}. Subiu na mesa pequena, no meio do pátio, e pintou três faixas.

\textbf{Verde}, embaixo. É a terra, e é o chão de onde a linha nunca pôde ser levantada. Um reino
que não tem chão não dobra; a primeira cor é o reconhecimento de que se está \emph{em cima de
alguma coisa}, e de que é ela que sustenta, e de que ninguém a fez.

\textbf{Azul}, no meio. É a imensidão --- o que cresce vezes, o tabuleiro que não pára de se
estender, a expansão que respira. É a cor do pai dele, e é a cor da doença do pai dele, e Charles
pintou-a no meio de propósito: entre o chão e a luz, contida dos dois lados, que é o único lugar
onde a expansão não vira vulcão.

\textbf{Branca}, em cima. É o que ilumina e não pesa: a luz que a corte de ouro tinha de sobra e
não sabia que tinha, porque quem nunca viu o escuro não sabe que está aceso.

\medskip

E disse a coisa que fez os dois reinos ficarem calados, porque era a primeira vez que alguém
punha aquilo em ordem:

--- Não é uma escolha entre as três. É uma \textbf{pilha}, e a ordem importa. Terra embaixo,
imensidão no meio, luz em cima. Quem tirar a de baixo, o resto cai. Quem tirar a do meio, a luz
encosta na terra e queima tudo. Quem tirar a de cima, fica um monte no escuro --- e nós já
tivemos um monte no escuro, e todos aqui sabem como acabou.

E \Qui , que nunca elogiou ninguém em trezentos anos, disse:

--- Este daqui completou a lição.

\chapter{Aplanai}
\hino{Firmeza, Bem Firmeza}{aplanai meu coração, para eu poder amar\\ e as ofensas dos irmãos eu
saber perdoar}

Faltava o mais difícil, e era com os que tinham ido embora.

\medskip

Porque na reunião da separação toda a gente chorou e ninguém se resolveu. Os que saíram levaram
mágoa --- a mágoa exata de quem foi deixado de fora de uma coisa que não escolheu não entender.
E, do lado de cá, ficou o outro veneno, que é pior e ninguém confessa: a \textbf{superioridade
dos que ficaram}. Estava no ar do pátio, e estava sobretudo nos que tinham chegado depois e não
tinham passado por nada, mas falavam dos que saíram como quem fala de gente que não aguentou.

\Qui viu aquilo crescer durante dois anos e um dia disse:

--- Eu prefiro a mágoa deles à \textbf{soberba} de vocês. A mágoa é uma dobra: dói e desfaz-se.
A soberba é um rasgo.

\medskip

E a coisa que se fez então --- e é a mais estranha deste livro, e é a que mais gente copiou
depois --- foi \textbf{aplanar}.

Estender-se no chão do pátio, todos, uma vez por lua, e ficar. Sem cantar, sem pedir, sem
resolver nada. Deitados, lisos, como as tiras do vale. E o ponto é exatamente esse: \emph{quem
está estendido não tem forma}. Não dá para saber, olhando duas fitas lisas lado a lado, qual
delas foi um dragão grande e qual foi um pequeno, quem ficou e quem foi embora, quem tinha razão
na discussão da mesa.

--- Não é humildade --- explicava a velha, que odiava a palavra. --- É \textbf{medida}. Vocês
estão a discutir vinco com vinco, e vinco é o que cada um tem de mais seu. Estendam-se. Vejam de
quanto é a diferença. Depois voltem a dobrar-se e discutam, se ainda quiserem.

Quase nunca ainda queriam.

\medskip

Foi assim que os primeiros começaram a voltar. Um de cada vez, sem alarde, sentando-se na fila
como todo mundo. E a regra da casa quanto a isso é de \Ven , e é a única regra que ele escreveu
de próprio punho, e está numa tábua torta pregada na árvore:

\emph{Quem volta, volta no lugar onde parou. Ninguém aqui recomeça do zero e ninguém aqui
explica onde esteve.}

\chapter{Eu te dei uma casa}
\hino{Eu Te Dei Uma Casa}{eu te dei uma casa que não falta ninguém,\\ para tu escolher aqueles
que te convém}

No décimo primeiro ano, \Ven chamou Charles e entregou-lhe a casa.

Não houve coroa, porque este reino não tem. Houve uma frase, dita de costas, debaixo da árvore,
com meia dúzia de testemunhas:

--- Eu te dei uma casa que não falta ninguém. Agora és tu que escolhes quem te convém.

\medskip

E o resto do que ele disse, naquela tarde, é o que se repete até hoje para quem recebe qualquer
coisa das mãos de alguém:

--- A estrada eu dou a todos. Isso não é teu para dar nem para tirar: quem quiser subir, sobe; a
fila é de quem chegar. O que é teu é a \textbf{casa}, e casa é outra coisa --- casa tem parede, e
parede quer dizer que uma parte fica de fora, e o dia em que tu fingires que não é assim, tu
perdes as duas.

--- E como é que eu sei quem me convém?

--- Não sabes --- disse \Ven . --- Vais errar. Vais deixar entrar quem devia ficar de fora e vais
mandar embora quem devia ter ficado, e vais fazer isso durante anos, e o teu erro vai ficar
gravado no material, porque aqui tudo fica. É esse o preço da casa, e não há casa sem ele. Quem
não aguenta errar em público não pode receber nada de ninguém.

\medskip

Depois disso ele falou pouco, e o pouco que falou já está contado noutra parte deste livro,
porque foi sempre a mesma coisa e ele não tinha mais nenhuma.

O que mudou naquela tarde não foi ele: foi Charles. Um filhote que atravessou uma fronteira na
palma da mão de outro, que foi entregue numa mesa de alfândega como se paga uma conta, e que
voltou para casa por conta própria, recebeu ali a única coisa que ninguém pode tomar de
ninguém --- e recebeu-a de pé, e sem ter pedido, que é como se recebe tudo o que presta neste
reino.

\chapter{Sinal de guerra}
\hino{São João na Terra}{quando ouvir falar de São João na terra, é sinal de guerra;\\ os
anjinhos é quem vem colher, para ajuntar em um só lugar}

E houve guerra. Havia de haver: uma casa que cresce é uma casa que aparece, e o que aparece é
disputado.

Não veio com exército. Veio do jeito que essas coisas vêm de verdade --- com \textbf{parecidos}.

\medskip

Ao fim de vinte anos havia, espalhadas pelos dois reinos, dezenas de casas que serviam mate
debaixo de árvores tortas. Algumas tinham subido ao monte e aprendido; a maior parte não. Copiava-
se o que se vê: a cuia, a mesa pequena, o deitar de costas, a tábua pregada na árvore. Ninguém
copiava o que não se vê, que é a única coisa que havia.

E foi aí que começaram as brigas, e não foram brigas com o inimigo: foram entre as casas. Sobre
qual era a verdadeira, sobre quem tinha recebido de quem, sobre a ordem das flores, sobre se a
cuia gira para a direita ou para a esquerda.

--- Está bom --- disse \Qui , que era velhíssima a essa altura e já quase não subia. --- Agora é
a hora de colher.

\medskip

Colher não é destruir as outras casas, e ela teve de explicar isso umas quinze vezes, porque
todos os jovens do pátio entenderam ao contrário e queriam sair a corrigir o mundo.

Colher é ir buscar \textbf{o que é da casa e está espalhado}, e deixar o resto onde está. Não é
guerra por território: é uma \emph{separação de material}. Havia, naquelas dezenas de casas,
criaturas que tinham gravado de verdade sem nunca terem subido ao monte --- e havia, no próprio
pátio, criaturas que estavam ali havia vinte anos e não tinham gravado nada. Colher é juntar os
primeiros num só lugar e parar de fingir que os segundos são da casa por antiguidade.

\medskip

Doeu mais do que a reunião da separação, porque desta vez não havia como dizer que era por causa
do monte. Era por causa do que cada um tinha ou não tinha feito, e isso está escrito no material,
e não há discussão possível com um vinco.

E Charles, que já era enorme e já não girava para o lado errado em lugar nenhum porque tinha
aprendido a girar para os dois, fez a colheita sozinho, casa por casa, durante quatro anos. Não
fechou nenhuma. Não disse a ninguém que estava errado. Chegava, sentava-se na fila, bebia, e ao
sair convidava um ou dois.

\chapter{Chamo o tempo}
\hino{Chamo o Tempo}{chamo o tempo para vir me ensinar;\\ depois que o tempo chega, ninguém quis
aprender}

\Odo morreu no vigésimo terceiro ano, deitado no pátio, e não disse nada porque nunca dizia. A
conta dele fechava.

\Qui morreu no vigésimo sexto, no meio de uma frase, a explicar a alguém pela milésima vez que
não se pula vinco por parecer pequeno.

\Ven durou mais quatro. Morreu como todos sabiam que ia morrer --- do coração, que era um tambor
mal afinado havia eras --- e morreu dormindo, de costas, debaixo da árvore, respirando fundo até
o fim. Não foi um mate. Não foi uma queda. Foi o prazo.

\medskip

E é aqui que este livro tem de dizer a coisa mais amarga que tem para dizer, porque senão não
presta.

Depois que eles morreram, a casa continuou --- e continuou bem, e cresceu, e serve mate até hoje
debaixo de uma árvore torta que já não é a mesma árvore. E a maior parte dos que sobem hoje não
sabe de nada disto. Não sabem do monte. Não sabem que houve uma tampa. Acham que a mesa é
pequena por tradição, que o juramento é de pé por respeito, e que a cuia roda porque é bonito
rodar.

O tempo veio, como se pediu que viesse, e ensinou --- e ninguém quis aprender, ou quiseram poucos,
e os poucos foi o bastante para não se perder, e é sempre assim, e nunca foi de outro jeito em
lugar nenhum.

\medskip

Por isso a última coisa que \Qui mandou gravar --- e mandou gravar na tábua torta, por baixo da
regra de \Ven , com a letra ruim dela --- foi isto:

\emph{Embora que não aprenda muito, aprenda sempre um bocadinho.}

Que é a única promessa que este reino faz a alguém, e é pequena de propósito, porque as grandes
não se cumprem.

\chapter{Treme a terra}
\hino{Chamo a Força}{treme a terra e balanceia, e vós não sai do seu lugar;\\ devagarinho vai
chegando, e quem chamou é que vai ficando}

Houve um tremor, uns anos depois, e foi um tremor de verdade: o monte inteiro balançou durante
quatro dias, caiu o que restava do salão das colunas, e abriu-se uma fenda no pátio que passa
hoje a três passos da árvore.

Ninguém sabe a causa. \Drk , que ainda era vivo e ainda enxergava no escuro, disse que um bicho
daquele tamanho a decompor-se debaixo da terra mexe com a terra, e que não há mistério nenhum ---
e disse isso do jeito dele, que é sem consolo e sem crueldade.

\medskip

O que ficou dos quatro dias não foi o estrago. Foi que a casa \textbf{não saiu do lugar}.

Estavam no meio da roda quando começou. A cuia estava na mão de alguém. E a coisa que se conta,
desde então, para explicar a quem chega o que é aquela palavra que eles usam o tempo todo, é
simplesmente o que aconteceu: a cuia continuou a rodar. Durante quatro dias, com o chão a subir
e a descer, com pedra a cair do que restava das colunas, sem que ninguém desse a ordem de
continuar e sem que ninguém desse a ordem de parar.

Não foi coragem. Metade estava apavorada, e alguns choravam, e não é vergonha nenhuma. É que
firmeza não é não ter medo: é \textbf{não sair do lugar} com medo.

\medskip

E depois, quando acalmou, contaram-se e não faltava ninguém --- e é isto o que aquele hino
antigo quer dizer, e ninguém tinha entendido até ali:

quem \emph{chamou} a força é que fica. Não porque a força poupe quem a chamou; porque quem a
chamou já sabia que ela vinha, e escolheu o lugar onde ia estar quando viesse.

Os outros --- os que passavam por perto, os que subiam de vez em quando, os das casas parecidas
--- esses correram. Não é acusação: correr é o certo, quando não se tem lugar.

\chapter{Os seres}
\hino{Concentração}{peço aos seres da floresta para vir me ajudar,\\ seres do vento, do céu, da
terra e do mar}

Depois do tremor, a casa mudou uma coisa, e foi a última mudança que este livro registra.

\medskip

Até ali, quem se sentava na roda eram os que subiam: bichos, criaturas, gente dos dois reinos.
Charles, que já era velho, mandou abrir a roda para o resto.

Não é figura de linguagem e não é bondade com a natureza: é uma correção de \textbf{conta}.

Porque a casa passou vinte e tantos anos a fazer as suas contas entre os presentes --- quem
serviu, quem bebeu, quem ficou, quem saiu --- e a conta fechava, e no entanto havia uma coisa que
não entrava em conta nenhuma e que sustentava tudo: a árvore, que dá a sombra; o chão, que segura
o pátio; a água da cisterna partida, que ninguém nunca agradeceu porque estava lá; o vento, que
leva o cheiro monte abaixo e é o motivo real de metade dos que chegaram terem chegado.

--- Nós fizemos a conta só com quem fala --- disse Charles. --- E o reino inteiro é dobrado do
mesmo material. Aquela árvore é uma fita, e tem os vincos dela, e ninguém nunca se agachou para
ler.

\medskip

Então passou-se a chamar. Não a rezar por eles: a \textbf{chamar} --- que é a diferença entre
tratar uma coisa como cenário e tratá-la como quem está na roda. Os da floresta, os do vento, os
do céu, os da terra e os do mar. E quando alguém perguntava se eles vinham mesmo, Charles dava a
resposta do pai dele, com a aspereza herdada:

--- Vieram sempre. A novidade é que agora a gente diz o nome.

\chapter{Cada um no seu lugar}
\hino{Eu Peço a Deus}{vi Deus no sol, na lua e nas estrelas, no vento, na terra e no mar;\\ todos
se firmam, se componham em seu lugar}

Fica esta imagem, e com ela se fecha de verdade.

\medskip

O pátio, ao entardecer, com a fenda do tremor atravessada por uma tábua. A árvore torta, que já é
outra árvore e ninguém corrigiu o nome. A mesa pequena, a mesma, que ninguém trocou por uma
maior. E a roda, com a cuia a andar sempre no mesmo sentido, passando de mão em mão, sem que
ninguém precise dizer para quem vai a seguir.

Não há trono. Nunca houve. O maior dragão que este mundo teve morreu deitado de costas no chão
com o peito descoberto, e a única coisa que ele deixou escrita, de próprio punho, foi que quem
volta volta no lugar onde parou.

\medskip

E a doutrina inteira desta casa, se é que há uma, é uma observação sobre lugares:

Um reino que só soma é belo e não tem tragédia, porque não tem escolha; nele nada pode dar
errado, e onde nada pode dar errado nada pode ser \emph{feito}. Um reino que dobra guarda tudo o
que fez, inclusive o que preferia não ter feito --- e por isso nele se pode morrer, e por isso
nele se pode ser \textbf{perdoado}, que é uma operação que não existe do outro lado, porque lá
não há registro do que perdoar.

Os dois são o mesmo reino. Sempre foram. Um é a fita estendida e o outro é a fita dobrada, e a
diferença entre eles é um gesto de duas mãos feito na hora certa, esperando o instante em que a
coisa cede.

E a coisa toda --- o monte, a tampa, a moeda que faltou, o filhote que atravessou, a noite das
oito horas, a cuia que não parou de rodar durante quatro dias de terra a tremer --- a coisa toda
cabe numa frase, e é a frase mais simples de todo este livro, e leva a vida inteira:

\medskip

\begin{center}
\emph{cada um se firma, e se compõe no seu lugar.}
\end{center}

\medskip

E é aqui que este segundo livro acaba, e acaba sem desfecho, de propósito: porque firmeza não
tem desfecho --- se tivesse, seria outra coisa, e teria outro nome.

A missa acabou e ninguém foi embora. Isso está na primeira linha desta parte, e vinte e tantos
anos depois continua a ser a única notícia: \emph{ninguém foi embora}. O monte desfez-se, o
lastro morreu, a velha morreu, o dono da casa morreu, metade da irmandade saiu pela porta em
duas levas diferentes e por motivos opostos --- e às sextas-feiras, no pátio arrombado, debaixo
de uma árvore que nasceu de uma semente que um bicho comeu dormindo, alguém enche uma cuia e a
passa para a direita.

Não é vitória. Nenhum livro deste reino termina em vitória, porque aqui não se vence: aqui se
serve, e servir não acaba --- só passa adiante.

\medskip

Falta uma pessoa nesta história, e ela esteve em todas as páginas, e eu ainda não disse o nome
dela.

% ══════════════════════════════════════════════════════════════════════════════
% LIVRO III — LUA BRANCA (os vinte e cinco)
% (a que esteve em todas as cenas e não teve uma linha)
% ══════════════════════════════════════════════════════════════════════════════

\part{Lua Branca}

\chapter{A que ferveu a água}
\hino{Estou Firme}{estou firme, não devo esmorecer;\\ eu peço e rogo para eu poder alcançar}

Este livro tem de começar com um pedido de desculpas, e o pedido é meu.

Contei duas vezes a mesma história e nas duas deixei alguém de fora. Contei a missa e contei a
concentração; contei o monte, a tampa, a moeda, a noite das oito horas, a cuia que rodou com a
terra a tremer. E em todas aquelas páginas houve uma criatura presente em cada cena, sem uma
linha, e eu não a nomeei nenhuma vez --- não por descuido: porque \emph{ela não fazia nada que
desse capítulo}.

Chama-se \Lua .

\medskip

É quem ferveu a água. As sete flores foram de \Qui --- a receita é dela, a mão é dela, a
advertência sobre o mulungu é dela. Mas alguém carregou quatro barris de água morro acima, dois
dias seguidos, e coou tudo num lençol, e lavou o lençol. Foi ela.

É quem estava atrás da pedra na noite das oito horas e não é citada porque não disse nada.
Contaram-se ali seis: os cinco do Reino Dourado e a velha. Eram sete.

É quem, no dia do juramento, pôs a mesa. A mesa pequena, de madeira ruim, que ninguém trocou por
uma maior e que virou o símbolo da casa inteira --- alguém a arrastou até o meio do pátio antes
de todos chegarem, e ficou de lado enquanto o resto acontecia.

É quem encheu a cuia durante quatro dias de terra a tremer. A cuia rodava; alguém a enchia.

\medskip

Ela não é a heroína escondida da história. Isso seria bonito e seria falso --- não há reviravolta
nenhuma aqui, ela não era ninguém disfarçado, não tinha um poder guardado, não estava à espera da
hora dela. Ela era exatamente o que parecia: a que ferve a água.

O que este terceiro livro diz --- e leva vinte e cinco capítulos a dizer, porque não se diz de
outro jeito --- é que uma casa que se conta pelos que falam \textbf{tem a conta errada}, e que
foi por isso que o vale dos Desdobrados ficou séculos sem ser lido, e que é o mesmo erro, e que
neste reino errar a conta é a única coisa que não se perdoa, porque aqui tudo fica registrado.

\medskip

Quando lhe perguntaram, muito depois, como é que ela tinha aguentado tantos anos sem que ninguém
lhe dissesse nada, ela deu a resposta que se repete até hoje e que é a coisa menos heroica que se
pode responder:

--- Eu não aguentei. Eu \textbf{fiquei}. É diferente, e é mais fácil do que parece: aguentar é
uma coisa que a gente faz com força, e força acaba. Ficar é só não sair.

\chapter{Conforme merecer}
\hino{Eu Vou Trabalhar}{o meu Mestre dá conforme merecer;\\ dá a todos, mas não sabem aproveitar}

\Lua chegou ao monte antes de todos. Antes dos cinco, antes de \Qui , antes até do estrago.

Chegou no tempo do monte inteiro, quando \Ven dormia de bruços sobre o ouro e a fila para subir
dava a volta na muralha, e chegou pela razão mais banal do mundo: precisava de trabalho. Deram-
lhe a cisterna. É o pior posto do monte --- ninguém vê, ninguém agradece, e se faltar água toda
a gente reclama.

Ficou trinta anos na cisterna.

\medskip

E aqui está a única coisa que ela dizia sobre aquilo, e que os que chegavam depois nunca
entendiam de primeira, porque parece conformismo e não é:

--- Todo mundo aqui recebe o mesmo. A casa dá a todos. Eu vi passar por esta cisterna criatura
mais forte do que eu, mais inteligente do que eu, mais bonita, mais falada. Recebeu igual. Não
aproveitou.

--- E o que é aproveitar?

--- É \textbf{voltar amanhã} --- dizia ela. --- Só isso. Não tem outro segredo, e eu procurei.

\medskip

Foi por ficar na cisterna que ela foi a primeira a saber da doença.

Porque \Ven descia uma vez por noite, sempre à mesma hora, sempre o mesmo tanto --- e quem passa
trinta anos ao lado de uma cisterna sabe, sem contar, quando alguém começou a beber diferente.
Ele começou a demorar mais. Depois começou a apoiar-se na borda. Depois começou a parar no meio
e ficar quieto, alguns segundos, olhando a água, do jeito que uma criatura fica quando está a
recuperar o fôlego e não quer que se note.

Ela viu aquilo durante onze anos e não disse a ninguém, e não foi covardia. Foi que não havia a
quem dizer: quem falasse mal do reinado do bicho maior do mundo não durava a semana, e ela não
tinha nome, nem cargo, nem prova --- tinha uma impressão de quem ferve água.

Foi ela quem mandou chamar \Qui . Numa carta de três linhas, sem assinar.

\chapter{Filha, te corrijas}
\hino{Meu Mestre Me Cure}{meu Mestre me cure, me olhe, me ensine;\\ filha, te corrijas para ser
perdoada}

E é preciso dizer que ela não era doce.

Isso estraga o retrato que toda a gente faz dela hoje, e é pena, mas é assim: a \Lua daqueles
anos era dura, calada e desconfiada, e tinha por dentro uma raiva antiga que ninguém do pátio
imaginava, porque ela nunca levantou a voz uma única vez em trinta anos.

A raiva era simples e era justa: ela via. Via os que subiam e desciam com histórias. Via os
cargos que só existiam para defender uma tampa. Via a fila crescer e a água acabar. E não podia
dizer nada, e quem não pode dizer nada guarda --- e o que se guarda por trinta anos não fica
igual ao que entrou.

\medskip

Foi \Qui quem viu, e viu no primeiro dia, porque é o ofício dela.

--- Tu me chamaste com uma carta sem assinatura --- disse a velha, na cisterna, sem cumprimentar.
--- Isso quer dizer duas coisas: que tu tens razão, e que tu tens medo. A primeira eu vou tratar
lá em cima. A segunda é tua, e se tu não a tratares, no dia em que ele cair tu vais estar pronta
para ocupar o lugar dele com a mesma tampa e outra cara.

\Lua não respondeu. Ficou com aquilo dentro.

\medskip

E o que ela fez com aquilo é a razão de haver um terceiro livro, e é a coisa mais difícil deste
reino, mais difícil que a gororoba e mais difícil que a firmeza:

ela \textbf{corrigiu-se antes de ter poder}.

Não depois. Antes --- nos anos em que ainda era a da cisterna, em que ninguém a olhava, em que
não havia nenhuma vantagem em ser justa porque não havia ninguém para reparar. Passou aqueles
anos a tirar de dentro de si, um por um, os motivos que ela teria para se vingar quando chegasse
a vez dela.

--- É a única correção que vale --- disse-lhe \Qui , muito depois. --- Quem se corrige depois de
subir está a arrumar a casa com visita dentro. Não dá tempo.

\chapter{Para todos eu receber}
\hino{A Meu Pai Eu Peço Força}{peço força para eu ter essa calma;\\ enviai os meus irmãos, para
todos eu receber}

Quando o monte caiu e a casa nova começou, deram-lhe a mesma função de sempre, e ninguém achou
que estivesse a fazer nada de mais: a água.

Só que agora a água era outra coisa.

\medskip

Porque quem enche a cuia é quem \textbf{recebe}. Não é figura: a fila passa por ela. Toda
criatura que sobe àquele pátio, seja quem for, chega primeiro à água --- e é ali, e não diante do
dragão, que se decide como a pessoa vai ser tratada naquele dia.

E \Lua tinha uma regra, que era dela e que nunca foi escrita em tábua nenhuma:

\emph{recebe-se todos igual, e recebe-se calmo, e não se pergunta nada.}

\medskip

Isso parece pouco e é quase tudo. Porque a casa cresceu, e com ela chegaram os difíceis: os que
subiam para discutir, os que subiam para se exibir, os que subiam a chorar sem parar, os que
subiam bêbados, os que subiam com uma mágoa de vinte anos atrás e queriam despejá-la em alguém.
O pátio inteiro tinha opinião sobre cada um deles. Ela dava água.

--- Tu não separas o joio --- reclamou um dos jovens, uma vez, indignado de verdade. --- Tu dás
a mesma água ao homem que veio comprar a árvore e à criatura que subiu quatro dias a pé.

--- Dou --- disse ela. --- Porque eu não sei qual é qual, e tu também não, e a diferença entre
nós dois é que eu sei que não sei. Aquele homem voltou sete anos depois e sentou-se na fila. Se
eu tivesse negado a água na primeira vez, ele não voltava, e a gente teria dito que estava certo.

\medskip

E foi assim, sem nunca ter tomado uma decisão de governo, que ela passou a decidir a coisa mais
importante daquele lugar: \textbf{quem se sentia recebido}. E não há casa que sobreviva a isso
estar em mãos erradas.

\chapter{Esmorecida}
\hino{Peço Uma Ajuda a Meu Mestre}{eu me acho esmorecida, sem eu poder me firmar;\\ peço uma
ajuda para eu me levantar}

Houve um ano em que ela não conseguiu.

Convém contar, porque um livro que só mostra a criatura firme ensina a coisa errada --- e porque
foi ela mesma quem mandou contar, e disse que se tirassem este capítulo tirassem o livro todo.

\medskip

Foi no ano em que morreu \Qui .

Não foi uma tragédia: a velha era velhíssima, morreu no meio de uma frase, a explicar pela
milésima vez que não se pula vinco por parecer pequeno, e não há morte melhor do que essa em
lugar nenhum. Toda a gente disse isso no pátio, e era verdade, e não ajudou nada.

Porque \Lua tinha uma pessoa no mundo que sabia quem ela era --- uma --- e era aquela velha
grosseira que a tinha lido na cisterna sem cumprimentar. E quando ela morreu, a \Lua voltou a
ser, para todo o pátio, a que ferve a água.

\medskip

Ela parou. Não anunciou, não se despediu, não fez cena: simplesmente ficou na cisterna e não
subiu mais ao pátio. Alguém foi buscando a água. As semanas foram passando. E quem passasse por
ali de manhã via uma criatura sentada na borda a olhar para dentro, e achava que ela estava a
trabalhar.

Durou quase um ano.

\medskip

Quem desceu foi Charles, que já era grande e já não girava para o lado errado em lugar nenhum.
Não levou conselho. Sentou-se na borda ao lado dela e ficou --- e ficou muitos dias, subindo só
para dormir, sem dizer nada de útil, falando de besteira, do tempo, da fenda no pátio, de uma
discussão idiota sobre a ordem das flores.

E um dia ela disse:

--- Eu não consigo me levantar.

E ele disse a única coisa que se pode dizer, e que não resolve nada, e que é a única que
funciona:

--- Eu sei. Eu venho amanhã.

\medskip

Ela levantou-se num dia de que ninguém se lembra, e voltou a subir com os baldes como se não
tivesse havido nada, e ninguém no pátio comentou --- porque a regra da casa quanto a isso é de
\Ven e está na tábua torta: \emph{quem volta, volta no lugar onde parou, e ninguém aqui explica
onde esteve}.

Foi a primeira vez que aquela regra serviu para alguém que não tinha ido embora.

\chapter{Dormindo sem acordar}
\hino{Há Tempo Estou Nesta Casa}{há tempo estou nesta casa, dormindo sem acordar,\\ ouvindo a voz
do meu Pai sem eu poder me firmar}

Há um estado que este reino conhece bem e para o qual não tinha nome antes dela, e ela deu-lhe
um: \textbf{estar na casa dormindo}.

Não é ser de fora. Não é ter ido embora. É estar ali --- há anos, todo dia, na fila, na roda, na
água --- e não ter gravado nada. Ouvir tudo, repetir tudo, saber tudo de cor, e o corpo não estar
dobrado daquele jeito.

\medskip

Ela dizia que tinha passado assim os primeiros vinte anos, e ninguém acreditava, e ela insistia:

--- Eu estava aqui. Eu enchia a cuia. Eu sabia todos os nomes, todas as histórias, todas as
regras --- eu era capaz de explicar a qualquer um que chegasse por que a mesa é pequena. E não
tinha acontecido nada comigo. Nada. Vinte anos.

--- E o que acordou a senhora?

--- Nada bonito --- dizia ela. --- Eu vi um bicho grande a apoiar-se na borda da minha cisterna
porque não conseguia respirar, e entendi que ninguém ia dizer, e entendi que se eu não dissesse
ninguém dizia. Foi isso. Acordei porque \textbf{fui a única pessoa disponível}.

\medskip

E acrescentava, e esta é a parte que ninguém gosta de ouvir:

--- Não esperem coisa melhor. Ninguém acorda porque entendeu. Acorda-se porque um dia calhou de
não haver mais ninguém, e a coisa estava na tua frente, e tu não pudeste fingir que não viste.
Quem está à espera de uma revelação vai dormir aqui dentro a vida inteira, e vai morrer com todas
as respostas certas.

\chapter{Se humilhar}
\hino{Meu Mestre a Vós Eu Peço}{eu também quero seguir --- não me despreze;\\ vamos todos cuidar
de se humilhar}

Ela usava uma palavra que o Reino Dourado não tem, e que quando chegou lá causou um mal-entendido
que levou anos a desfazer-se: \textbf{humilhar-se}.

\Yas , que é a coerência e que empurra o mundo para a frente, ouviu aquilo e ficou ofendida em
nome de toda a gente. Disse que era doutrina de gente derrotada, que ensinar alguém a diminuir-se
é ensinar a servir, e que no reino dela ninguém se ajoelhava para coisa nenhuma.

E \Lua , que não discutia nunca, respondeu daquela vez, e respondeu assim:

\medskip

--- A senhora entendeu \emph{abaixar-se}. Não é. Abaixar-se é ficar menor do que se é, e isso é
mentira, e mentira aqui não fica de pé porque tudo se grava.

--- Humilhar-se é ficar do tamanho que se é. É a coisa mais difícil que existe, minha senhora,
porque quase todo mundo anda por aí um pouco maior --- pela função, pelo cargo, pela família,
pelo que fez uma vez há vinte anos. E enquanto se anda maior, não se aprende nada, porque não
cabe mais nada: já está tudo cheio.

--- Eu fervi água trinta anos. Não é virtude nenhuma, é o meu tamanho. E foi por eu estar do meu
tamanho que eu vi uma coisa que este pátio inteiro, que é feito de criaturas enormes, não viu:
que o dono do monte não estava a respirar.

\medskip

\Yas ficou calada muito tempo. Depois disse a coisa que fez as duas ficarem amigas, e que só uma
criatura da coerência diria:

--- Então é uma \textbf{medida}. A senhora está a falar de medir-se direito.

--- Estou. E vocês lá do outro lado medem tudo com uma exatidão que dá gosto, menos isso.

\chapter{Duas flores}
\hino{Duas Flores}{recebi duas flores para amar no meu coração;\\ o que Deus faz é bem feito ---
não se desmancha com a mão}

Ela recebeu duas coisas para cuidar, e recebeu no mesmo dia, e são as duas coisas de que este
reino vive.

Uma foi a \textbf{água}. A outra foi o \textbf{fogo}.

\medskip

A água já se contou: a cisterna, a fila, a cuia, os trinta anos. É a que se vê.

O fogo é a outra, e é a que este livro ainda não tinha dito. Quando \Ven desfez o monte e passou
a viver deitado de costas, houve um problema prático que ninguém previu: um dragão que não pode
mais rolar não pode mais acender coisa nenhuma. Ele tem o fogo --- tem a taxa, tem o fator, tem
tudo --- e não tem a posição.

Então alguém tinha de acender. E como o fogo daquele bicho não é chama e sim taxa, quem acende
tem de saber \emph{quando parar} --- porque uma coisa que dobra a cada instante não avisa antes
de virar catástrofe: no instante em que se nota, já é tarde por definição.

Deram isso à criatura da água. Não por simetria bonita: porque era a única que tinha passado
trinta anos a olhar para um nível subir.

\medskip

--- São a mesma coisa --- dizia ela, e foi a única frase de teoria que ela produziu na vida, e o
Reino Dourado passou anos a tentar prová-la e não conseguiu melhorar. --- Cuidar da água é saber
que o pouco falta. Cuidar do fogo é saber que o muito sobra. Quem sabe as duas sabe medir, e
quem não sabe medir não devia estar perto de nenhuma das duas.

\medskip

E acrescentou a coisa que virou a regra da casa sobre não mexer no que está feito:

--- E não se conserta o que está bom. Isso aqui foi feito por gente que sofreu, e ficou bem
feito, e há sempre alguém com uma ideia de melhorar a mesa. \textbf{O que é bem feito não se
desmancha com a mão.}

\chapter{Não me deixe derribar}
\hino{Peço Firmeza a Meu Pai}{suspendei meu pensamento, não me deixe derribar;\\ estou aqui para
ensinar, mas ninguém quer aprender}

Veio o cansaço, e veio de um lugar que ela não esperava: dos que ela ensinava.

\medskip

Porque a essa altura já havia gente a descer o monte a dizer \emph{a Lua ensina assim}, e a
maior parte estava a dizer coisa que ela não tinha dito. Ensinava-se por aí que ela mandava
receber todo mundo igual --- e alguns tinham entendido que isso queria dizer que nada importa.
Ensinava-se que ela mandava humilhar-se --- e alguns tinham entendido que era para se
diminuírem, e havia agora um punhado de criaturas tristes convencidas de que a virtude é não
valer nada.

Ela via aquilo voltar deformado, ano após ano, e não tinha como corrigir, porque quem repete
errado repete com convicção e quem tenta corrigir parece que está a defender-se.

\medskip

--- É isto que derriba --- disse ela, uma vez, à noite, para Charles. --- Não é o trabalho. Não é
a idade. É a gente dizer uma coisa a vida inteira, com todo o cuidado, e ela chegar do outro lado
virada do avesso, dita pela boca de quem gosta muito da gente.

--- E o que é que se faz?

--- Não se faz nada. Diz-se outra vez. --- E, mais baixo: --- E é isso que ninguém aguenta, filho.
Não é ensinar: é ensinar \textbf{de novo}. A mesma coisa, à mesma gente, no ano seguinte, sem
azedar. É aí que quase todo mundo desiste e vira um velho amargo que diz que ninguém quer
aprender.

--- A senhora azedou?

--- Todo dia --- disse ela. --- E amanhã eu subo com os baldes.

\chapter{A espada}
\hino{Oh! Minha Senhora Mãe}{segura bem minha mão --- te firma nestes três seres,\\ para quando
chegar o dia, estar com a espada na mão}

E houve o dia em que ela precisou de decidir, e decidir é a única coisa para a qual trinta anos
de cisterna não preparam ninguém.

\medskip

Foi na terceira geração da casa, muito depois da morte de \Ven . Charles estava velho. Havia no
pátio uma disputa feia sobre quem devia receber a tábua torta --- e disputa feia, num lugar onde
tudo fica gravado, é uma coisa que envenena por décadas, porque ninguém pode fingir depois que
não disse o que disse.

E os dois lados, para espanto dos dois lados, foram buscá-la.

Não porque a respeitassem mais do que aos outros. Porque era a única criatura daquele pátio que
nunca tinha tomado partido em coisa nenhuma --- e portanto a única de quem nenhum dos dois lados
tinha razão para ter medo.

\medskip

Ela ouviu os dois. Levou quatro meses --- e quatro meses, numa casa em brasa, é uma eternidade, e
metade do pátio dizia que ela estava a fugir da decisão.

Não estava. Estava a fazer o que aprendeu no vale: ir devagar, não pular nenhum vinco por parecer
pequeno, e não decidir o que a coisa quer dizer antes de chegar ao fim.

E ao fim dos quatro meses subiu na mesa pequena --- ela, a da água, em cima da mesa em que \Ven
tinha jurado de pé --- e disse a decisão em quatro frases, e as quatro frases desagradaram aos
dois lados, e ninguém saiu da casa.

\medskip

Depois desceu da mesa e alguém lhe perguntou como é que ela tinha tido coragem.

--- Coragem não --- disse ela. --- Eu passei a vida inteira a evitar isto, e evitei bem, e se
houvesse mais alguém eu tinha deixado com ele. Chegou o dia e não havia. --- E, encolhendo os
ombros: --- É sempre assim que se recebe uma espada. Ninguém escolhe: um dia ela está na tua mão
e tu tens de saber o que fazer com uma coisa que tu nunca quiseste.

\chapter{A peia}
\hino{Eu Peço Força a Meu Pai}{vamos todos se alegrar do que vai acontecer:\\ o Mestre bem avisou
que a peia ia descer}

Aqui é preciso falar de uma coisa deste reino que assusta quem chega, e que \Lua explicava melhor
do que ninguém porque não tinha nenhum prazer nela.

Chama-se \textbf{a peia}.

\medskip

Não é castigo, e é a primeira coisa que ela dizia, porque toda a gente entende castigo e ninguém
entende isto. Castigo é uma coisa que alguém aplica em alguém. A peia não tem quem aplique: é o
que acontece quando uma criatura que vive dobrada de um jeito torto passa por um lugar onde tudo
se grava.

Num reino que só soma, a mentira é confortável --- soma-se e segue-se, e a linha não guarda. Num
reino que dobra, a mentira \textbf{fica no material}. E chega o dia em que a fita é lida --- e
não é um juiz que lê, é a própria pessoa, que se agacha e passa o dedo e vê a sequência dos
próprios vincos na ordem em que foram feitos.

Isso dói de um jeito que nada mais dói. E não há defesa, porque não há acusador.

\medskip

--- Todo mundo aqui passa por isso --- dizia ela. --- Todo mundo. Eu passei. E os que aqui chegam
a dizer que já se resolveram, esses passam duas vezes.

--- E não dá para evitar?

--- Dá --- dizia ela. --- É só nunca ler a própria fita. Muita gente faz. Morre-se inteiro assim,
sem nunca ter doído nada, e é a coisa mais triste que eu conheço, e do lado de fora parece uma
vida boa.

\medskip

E dizia a frase que ficou, e que é a única alegria estranha que este reino oferece:

--- Alegrem-se quando descer. Não é o fim do mundo: é o mundo a funcionar. Quem tem confiança não
precisa correr --- e quem corre, corre da própria fita, e ela vai junto.

\chapter{A disciplina é geral}
\hino{A Disciplina É Geral}{a disciplina é geral, que é para todos acordar;\\ quem não ligar
importância não tem a quem se queixar}

E a peia não escolhe cargo. É por isso que se chama geral, e é a única igualdade que este reino
tem de verdade.

\medskip

O caso que se conta é o do quinto ano da casa nova, e conta-se porque foi \Lua quem o registrou
na fita e mandou que ficasse.

Havia um da casa --- dos primeiros, dos que tinham estado atrás da pedra, dos que carregaram
barril --- que começou a cobrar. Pequeno: uma parte da erva, uma parte do que traziam os que
subiam, uma coisa que ele chamava de contribuição e que era voluntária do jeito que essas coisas
são voluntárias quando quem pede tem a chave da porta.

Todo o mundo sabia. Ninguém disse nada durante quatro anos, porque ele era dos primeiros e porque
ninguém queria ser o que denunciou um dos primeiros.

\medskip

E não foi denunciado. Foi a fita.

Um dia alguém do vale --- daqueles que iam ler os Desdobrados --- trouxe uma tira nova, de um
bicho que tinha morrido no ano anterior, e leram-na em roda no pátio como se lia todas. E lá
pelo meio da tira estava aquilo: os vincos de uma criatura que subiu ao monte, foi bem recebida,
e pagou.

Ninguém acusou ninguém. A leitura continuou até o fim, porque não se pula vinco por parecer
pequeno. E quando acabou, houve um silêncio muito longo, e o homem levantou-se sozinho e foi
embora sozinho, e ninguém o impediu e ninguém o xingou.

\medskip

--- Foi a coisa mais limpa que eu já vi --- disse \Lua , muito depois. --- E é o motivo de eu
nunca ter querido ser juíza de nada. Numa casa onde tudo se grava, a única coisa que os grandes
têm de fazer é \textbf{ler em voz alta}, e a leitura julga sozinha. Quem começa a julgar por
conta própria está a dizer que não confia no material --- e aí é ele que virou o problema.

\chapter{É pedindo}
\hino{É Pedindo e Rogando}{é pedindo e rogando que podemos alcançar;\\ eu ensino é ter amor, mas
só querem a ilusão}

Perguntaram-lhe uma vez qual era o método. Ela achou graça, porque não havia nenhum, e depois
pensou melhor e disse que havia um e que era ridículo:

\medskip

--- \textbf{Pedir.}

--- Pedir a quem?

--- A qualquer um. Ao que estiver na tua frente. --- E, à cara de decepção do outro: --- Tu
querias uma técnica. Não tem. Eu passei a vida a ver criatura enorme subir este monte a querer
receber uma coisa sem ter de pedir a ninguém --- querem que caia, que apareça, que se descubra
sozinho, que seja mérito. Ninguém quer \emph{pedir}, porque pedir é dizer em voz alta que falta.

--- E é preciso dizer?

--- É a única coisa que é preciso.

\medskip

Ela explicava aquilo com a cisterna, que era como ela explicava tudo:

--- Água não sobe sozinha. Alguém carrega. Uma casa inteira de criaturas que sabem voar e cospem
fogo depende de que alguém, todo dia, peça a outro alguém que carregue. Se ninguém pedir, ninguém
carrega, e não é maldade de ninguém: é que não se sabe.

--- Mas os grandes deviam ver.

--- Deviam --- concordava ela. --- E não veem. Eu fiquei onze anos calada à espera de que
alguém visse que aquele bicho não estava a respirar. Onze anos. E no fim quem resolveu não foi
alguém ter visto: fui eu ter escrito três linhas numa carta e mandado chamar. \textbf{Foi um
pedido que salvou este reino}, e sem assinatura, e com a letra ruim.

\chapter{As flores precisam de água}
\hino{Estou Aqui}{o beija-flor voou, avisando que as flores é preciso aguar;\\ eu estava
dormindo, e Ela veio me despertar}

O jardim é dela, e é a única coisa que ela pediu para si em toda a vida.

\medskip

Pediu no ano em que voltou da cisterna, depois do ano parado. Pediu um pedaço do pátio --- o pior
pedaço, o que fica junto à fenda do tremor, onde a terra é ruim --- e plantou flores. Não é
metáfora e não tem função nenhuma: são flores, não se comem, não curam nada, não entram na
receita das sete.

Metade do pátio achou uma extravagância. Estavam a reconstruir muralha, e a criatura da água
estava a plantar flor.

\medskip

E ela deu a explicação uma vez só, e nunca mais, e é esta:

--- Este lugar aqui é muito duro. A gente lida com bicho a morrer, com fita esticada, com peia, e
com gente que sobe carregando trinta anos de coisa mal resolvida. Se não houver nada aqui dentro
que seja só \textbf{bonito}, sem servir para nada, a casa vira hospital --- e daqui a duas
gerações ninguém sobe mais, e vocês vão dizer que foi porque o povo esfriou.

\medskip

O que ninguém previu foi o beija-flor.

Apareceu no segundo ano do jardim, e depois eram muitos, e depois era uma coisa que se via de
longe, subindo o monte: uma nuvem de bichinhos rápidos sobre o pior pedaço de terra do pátio. E
foi por causa deles que \Lua descobriu uma coisa que mudou o modo de a casa cuidar de si mesma.

Porque um beija-flor não vai onde há flor. Ele vai, e se a flor está seca ele \emph{volta} --- e
volta de um jeito visível, aos pares, aos saltos, rondando. Ela reparou nisso num verão e passou
a usar: quando os bichinhos começavam a rondar sem pousar, era sinal de que alguma coisa naquele
canteiro tinha falhado três ou quatro dias antes e ninguém tinha visto.

--- É um aviso que chega antes do estrago --- dizia ela. --- Vocês estão sempre à espera do
sintoma grande. O sintoma grande é tarde. Aprendam a ver o \textbf{bichinho a rondar}.

E acrescentava, porque nunca perdia a chance:

--- Onze anos eu vi um dragão apoiar-se na borda da cisterna. Aquilo era o beija-flor a rondar. E
eu, que estava dormindo nesta casa, achei que não era comigo.

\chapter{A força da lua}
\hino{A Lua Tem Uma Força}{a lua tem uma força, e nela eu vou me firmar;\\ o poder da Virgem Mãe,
nós temos em nossas mãos}

Ela reparou numa coisa que ninguém tinha reparado, e é uma medição, e o Reino Dourado confirmou
depois com todos os instrumentos deles e não conseguiu explicar melhor do que ela explicou.

\medskip

A casa funciona diferente conforme a lua.

Não é misticismo: é a maré da cisterna, que sobe e desce; é o sono dos bichos, que fica pior nas
noites claras; é a quantidade de gente que sobe o monte, que ninguém tinha contado até ela
começar a riscar a parede. Trinta anos de riscos, e o desenho na parede da cisterna é uma onda,
e a onda tem período, e o período é o da lua.

--- Vocês acham que decidem tudo --- dizia ela. --- Decidem quase nada. Metade das brigas deste
pátio acontece na mesma semana do mês, todo mês, há trinta anos, e vocês discutem o assunto da
briga.

\medskip

E daí ela tirou a única coisa parecida com uma administração que aquela casa teve:

Nas semanas de lua cheia --- quando sobe mais gente, quando se dorme mal, quando se briga ---
não se decide nada, não se muda nada, não se recebe ninguém novo na casa. Trabalha-se, serve-se,
ferve-se a água, e espera-se passar.

Nas semanas escuras, decide-se.

\medskip

--- Isso não é fraqueza --- teve de explicar umas cem vezes, porque os jovens achavam humilhante
depender do céu. --- É a mesma coisa que o senhor da apneia. Ninguém obriga um dragão a virar-se;
espera-se o corpo dele. Ninguém obriga um pátio a estar calmo; espera-se a semana dele. A gente
não manda na força. A gente escolhe \textbf{onde estar} quando ela vem.

--- E onde a senhora está?

--- Na cisterna --- disse ela. --- Todo mês. Há trinta anos.

\chapter{Três seres}
\hino{Três Seres}{três seres me acompanham --- cobri-me com Vosso manto;\\ Ela chora todo dia em
ver tantos inocentes}

Charles pintou três faixas na mesa pequena: verde embaixo, azul no meio, branco em cima. Terra,
imensidão, luz. Está no segundo livro e é a lição dele.

\Lua acrescentou uma coisa que ele não tinha visto, e acrescentou anos depois, quando ele já era
velho e ela também.

\medskip

--- Falta dizer quem \textbf{segura}.

--- Segura o quê?

--- As três. Tu puseste na ordem certa, filho, e a ordem está certa e ninguém mexe. Mas uma pilha
não se segura sozinha. Se ninguém estiver por baixo do verde, com o ombro, todo dia, aquilo
desmancha --- e o que está por baixo do verde não é uma quarta cor. É uma criatura, e vai ser
sempre uma criatura, e vai ser sempre alguém de quem tu não te lembras o nome.

\medskip

Foi por causa disso que se mudou o modo de contar a casa.

Até ali contavam-se os que falam. Passou-se a contar em três: o que fala, o que ouve, e \textbf{o
que sustenta} --- e o terceiro passou a ser dito em voz alta na roda, com nome, todo mês. Quem
ferveu, quem carregou, quem varreu, quem ficou com o doente, quem cavou a fenda.

Não é gratidão. É de novo uma \textbf{correção de conta}, e é a mesma correção do vale dos
Desdobrados, e é a única coisa que este reino faz que os outros não fazem.

\medskip

E \Lua chorava com aquilo, todo mês, e não escondia --- e quando lhe perguntavam por quê, dizia:

--- Porque eu passei vinte anos sem que ninguém dissesse o meu nome nesta roda, e eu achava que
estava certo. Não é dó de mim, filho. É que cada mês eu ouço um nome que eu não sabia, e penso em
quantos anos aquela criatura achou que estava certo.

\chapter{Salve, Lua Branca}
\hino{Lua Branca}{salve, salve Lua Branca --- a Vós eu peço saúde\\ para mim e minha família}

O nome pegou tarde, e pegou por acidente, e ela nunca gostou.

\medskip

Vinha de uma coisa boba: ela trabalhava de noite. Sempre trabalhou --- a cisterna enche melhor de
madrugada, quando ninguém tira água --- e quem subisse ao pátio nas noites claras via, do outro
lado, uma criatura pálida a andar de um lado para o outro com os baldes, do tamanho de nada,
debaixo da lua.

Os que chegavam perguntavam quem era aquela. E alguém, uma vez, respondeu a coisa que ficou:

--- Aquela? Aquela é a lua branca desta casa. Sempre lá.

\medskip

Ela achou horrível. Passou anos a corrigir, e desistiu, porque não há como impedir que uma casa
inteira te chame de uma coisa que ela precisa que exista.

E quando já era velha e já não corrigia mais, explicou por que tinha achado horrível, e a
explicação é o capítulo:

--- Vocês me deram um nome de coisa que não faz nada. A lua não aquece, não alimenta, não trabalha
--- ela só está lá, bonita, e todo o mundo olha. Eu passei a vida a carregar balde, filho. Eu
queria um nome de \emph{ofício}.

--- E hoje a senhora acha ruim?

--- Hoje eu acho que vocês repararam melhor do que eu --- disse ela. --- Porque a lua não faz
nada, e mesmo assim a maré inteira depende dela. Não é força; é \textbf{presença}. E foi isso que
eu fiz aqui trinta anos, e eu é que não sabia o nome.

\medskip

Só pediu uma coisa, e a casa cumpre até hoje: que quando cantassem aquele nome, cantassem
pedindo saúde --- para os da casa e para as famílias dos da casa, inclusive as famílias dos que
foram embora, inclusive a família do homem das carroças.

--- Se é para ter um nome que não trabalha --- disse ela --- que pelo menos sirva de recado.

\chapter{Quem não quiser}
\hino{Vamos Todos Com Amor}{quem não quiser obedecer é melhor se retirar;\\ quem não prestar
atenção vai ficar na disciplina}

E ela, que recebia todo o mundo igual e não perguntava nada, foi quem escreveu a regra mais dura
que aquela casa tem. As duas coisas são a mesma, e é este o capítulo mais mal-entendido do livro.

\medskip

A regra é: \emph{a porta é larga, e a roda é estreita.}

Entrar, entra quem quiser. Sobe quem sobe, bebe quem chega, não se pergunta nada, não se examina
ninguém, e a água é a mesma para o homem das carroças e para quem veio quatro dias a pé.

Mas a \textbf{roda} --- a cuia que gira, a ordem, o sentido, quem enche e quem recebe --- essa não
se discute. Não é democrática, não se vota, não se adapta ao gosto de quem chegou ontem.

\medskip

--- Porque a porta é hospitalidade --- dizia ela --- e a roda é o que a casa \emph{é}. Se eu
alargar a porta, entra mais gente. Se eu alargar a roda, não entra nada, porque não sobra roda
nenhuma para entrar.

E acrescentava, sem doçura:

--- Quem não quiser obedecer à roda, retira-se, e retira-se em paz, e leva a nossa consideração e
pode voltar quando quiser. O que não pode é ficar dentro a mudar as coisas devagarinho enquanto
diz que é da casa. Isso não é liberdade, filho. Isso é \textbf{cupim}.

\medskip

Perguntaram-lhe se não era autoritário, vindo dela justamente, que tinha passado trinta anos sem
poder dizer nada.

--- É por eu ter passado trinta anos sem poder dizer nada --- respondeu. --- Eu sei exatamente
qual é o preço de uma casa em que quem chega manda. Quem paga é sempre o mesmo: o que ferve a
água. Nunca é o que fala alto.

\chapter{Se compõe em seu lugar}
\hino{Eu Peço Força a Meu Pai}{pisa firme e segue em frente, se compõe em seu lugar;\\ aonde tu
estás estou, aonde estou tu estás}

Esta é a frase da casa. Se um dia se perder tudo o que está escrito nestes três livros e sobrar
uma linha, que sobre esta:

\begin{center}
\emph{pisa firme e segue em frente --- se compõe em seu lugar.}
\end{center}

\medskip

Ela dizia que as três partes são uma só e que quem tira uma estraga as outras duas.

\textbf{Pisa firme}: não é força, é a firmeza do segundo livro --- ficar, não sair, voltar amanhã
ao mesmo lugar.

\textbf{Segue em frente}: e é aqui que a maior parte se engana, porque firmeza sem andar vira
pedra, e há casas inteiras que se orgulham de não ter mudado nada em cem anos e que estão mortas
há noventa.

\textbf{Se compõe em seu lugar}: que é o mais difícil dos três, e é o que ela era. Compor-se não
é conformar-se: é achar qual é o teu lugar de verdade, que quase nunca é o que tu querias, e
depois \emph{fazer aquilo bem}. A cisterna era o pior posto do monte. Ela não a aceitou com
resignação; ela a fez melhor do que qualquer um antes ou depois, e foi de lá que ela viu o que o
pátio inteiro não viu.

\medskip

E a segunda metade do hino é a coisa mais bonita que este reino diz, e é uma coisa sobre
distância:

\emph{aonde tu estás, estou; aonde estou, tu estás.}

Num reino que é uma linha, isso é falso: cada um está no seu andar, e entre dois andares há todos
os degraus do meio, e não há como estar em dois lugares.

Num reino que dobra, é literal. Basta uma quina. Duas coisas encostadas na dobra estão à
distância zero uma da outra por mais afastadas que estejam na fita estendida --- e é exatamente
assim que uma criatura que ferve água no pior canto do pátio está encostada no dragão maior do
mundo, e ele nela, sem nenhum dos dois ter andado um passo.

\medskip

--- Não é consolo de pobre --- disse ela, quando lhe perguntaram se aquilo não era só uma maneira
bonita de dizer que os pequenos importam. --- É \textbf{geometria}. Vocês do outro lado passaram
uma eternidade a subir degrau para se aproximar de alguma coisa. Aqui a gente dobra. É mais
rápido, e ninguém precisa de sair do lugar.

\chapter{Lá vem o sol}
\hino{Lá Vem o Sol}{lá vem o sol nos curar --- vamos todos se firmar\\ no sol, na lua, nas
estrelas e na Rainha do Mar}

E na velhice ela corrigiu a conta uma última vez, e foi a maior correção de todas.

\medskip

Charles tinha aberto a roda para os seres: os da floresta, os do vento, os do céu, os da terra e
os do mar. Está no fim do segundo livro, e a casa passou a chamá-los pelo nome, e ninguém achou
que faltasse nada.

\Lua ouviu aquela lista durante anos, todo mês, e um dia levantou-se e disse:

--- Falta uma.

\medskip

--- Vocês chamam o mar. O mar é o lugar. Ninguém chama \textbf{quem manda nele}.

E explicou, e é a coisa de que ela mais gostava:

--- A água desta cisterna não é minha e não é da casa. Ela vem de cima, e antes de vir de cima
subiu do mar, e antes de subir do mar esteve na chuva, e antes disso outra vez no mar. Eu passei
trinta anos a carregar uma coisa emprestada e a achar que o trabalho era meu. O trabalho é meu ---
o balde é meu, o ombro é meu. \textbf{A água não.}

\medskip

E a lista final da casa, a que se canta até hoje, ficou com quatro e depois com cinco, e cada um
é uma régua:

O \textbf{sol}, que cura --- e cura porque é a única coisa que dá sem receber, e a casa inteira é
feita de contas que fecham, e é bom ter uma que não fecha e ainda assim funciona.

A \textbf{lua}, que mede --- a maré, o mês, a semana em que não se decide nada.

As \textbf{estrelas}, que orientam --- e que só servem à noite, que é quando se precisa.

E a \textbf{Rainha do Mar}, que é de onde vem tudo o que se carrega em balde e que ninguém fez.

\medskip

--- Firmar-se em quatro coisas que estão todas fora de nós --- disse ela. --- É de propósito. Uma
casa que se firma nela mesma cai. Já vi.

\chapter{Não é o fim}
\hino{Frei Damião}{ainda não é o fim --- é o começo das dores}

Ela viveu o suficiente para ver a casa passar por tudo outra vez, e é isso que uma vida longa
tem de cruel: não é ver o mal, é ver o \emph{mesmo} mal.

\medskip

Viu chegar a terceira geração, e viu nela, palavra por palavra, as mesmas propostas que \Qui
tinha esgoelado contra sessenta anos antes: que a mesa é pequena demais, que o juramento de pé é
desnecessário, que a cuia devia ser um copo para cada um por causa dos que vêm de fora e acham
estranho.

Viu voltar o homem das carroças na figura de outro homem, e voltar de novo, e há de voltar.

Viu o cansaço, o azedume, a soberba dos que ficaram, os que dormem na casa a vida inteira com
todas as respostas certas.

\medskip

E o que ela disse sobre isso é a coisa menos consoladora deste livro, e ela fazia questão:

--- Não é o fim de nada. Vocês estão sempre à espera de que uma dessas coisas seja a última ---
que se resolva a briga da mesa e pronto, que se afaste o homem das carroças e acabou. Não acaba.
\textbf{É o começo}, e vai ser o começo outra vez daqui a vinte anos, com outras caras, e vocês
já vão estar mortos e alguém vai estar a explicar isto a alguém.

--- Então não adianta nada?

--- Adianta a casa continuar de pé para poder passar por isso de novo --- disse ela. --- É só
isso que a gente faz aqui: mantém de pé uma coisa por onde a mesma dor vai passar mais umas
cinquenta vezes. E cada vez que ela passa, uns dois ou três aprendem. Não é pouco. Faz sessenta
anos que eu conto, e dá mais ou menos dois por vez, e olha quanta gente já foi.

\chapter{Licença à terra}
\hino{Eu e Meu Pai É Um Só}{pedi firmeza à lua, coragem ao sol, força ao vento;\\ pedi licença à
Terra para nela habitar}

Já velha, contou como tinha feito. Não é um método e ela detestava a palavra, mas é a única
coisa parecida com uma receita que ela deixou, e é feita de quatro pedidos.

\medskip

\textbf{Pediu firmeza à lua.} Não à sua própria vontade: à coisa que volta todo mês faça o que
fizer. Quem se firma na própria disposição está perdido, porque disposição é a coisa mais
instável que existe. Firma-se no que é regular --- a cisterna enche todo dia, o mês volta todo
mês, a roda gira toda semana. Firmeza é apoiar-se no que não depende do humor de ninguém.

\textbf{Pediu coragem ao sol.} Ao que dá sem receber. Porque a coragem que vem da conta ---
faço isto e recebo aquilo --- dura enquanto durar o pagamento, e no dia em que não pagam,
acabou. A coragem que dura é a que não espera troco.

\textbf{Pediu força ao vento.} À coisa que não se vê e move tudo, e que ninguém segura na mão. É
a força de quem não tem posto: não se manda, mas se \emph{chega}, e o cheiro da erva desce o monte
e traz metade dos que vieram.

\textbf{E pediu licença à terra para nela habitar} --- e este é o que ela dizia que quase ninguém
faz, e é o que ela mais repetia.

\medskip

--- Vocês todos chegam aqui e ocupam. Ocupam a casa, ocupam o cargo, ocupam a fila, ocupam a
conversa. E ninguém \emph{pede licença} ao lugar. A terra deste pátio estava aqui antes do monte
e vai estar aqui depois de todos nós. A árvore não é nossa. A água não é nossa. Eu não sou daqui,
filho --- eu vim de longe, como todo mundo veio de algum lugar, e passei sessenta anos hóspede
neste monte e nunca deixei de ser.

--- E o que muda, pedir licença?

--- Muda o cuidado --- disse ela. --- A gente cuida diferente da casa dos outros. E é isso que
vocês precisam: cuidar daqui como quem está de visita há sessenta anos e ainda tira os sapatos na
porta.

\chapter{A mudança do tempo}
\hino{Divina Estrela}{a Divina Estrela está na terra, na floresta e no mar;\\ vejam a mudança do
tempo --- devemos ficar atentos}

No último ano dela houve uma coisa que ninguém soube explicar, e este livro não vai tentar,
porque quem explica o que não entendeu está a mentir bonito.

\medskip

Apareceu uma luz.

Não foi visão de ninguém, e é isso que a torna estranha: foi vista pelos dois reinos ao mesmo
tempo, pelos que estavam acordados e pelos que não estavam a rezar coisa nenhuma, e foi vista do
alto da Torre e do fundo do pátio, e as duas descrições batem no essencial e não batem em quase
nada.

Do lado de cá: uma estrela baixa que ficou três noites e que iluminava a fenda do tremor de um
jeito que a fenda parecia costurada.

Do lado de lá: uma coisa que \Drk mediu e não conseguiu classificar, e que ele registrou com
honestidade como \emph{um sinal com assinatura, sem fonte identificável}, e não disse mais nada,
porque não era de dizer.

\medskip

O que se sabe é o efeito, e o efeito foi este: a casa ficou atenta.

Não assustada --- atenta. E foi naquelas semanas que se fizeram as coisas que seguraram tudo
depois: leram-se as fitas antigas, arrumou-se a tábua torta que estava a apodrecer, escreveu-se
pela primeira vez a lista dos que sustentam, e Charles, que já mal andava, subiu a cada uma das
dezenas de casas parecidas e sentou-se na fila delas uma última vez.

\medskip

\Lua disse uma frase só sobre a luz, e disse encolhendo os ombros:

--- Não sei o que é. Sei que o tempo mudou, e a gente sente mudança de tempo do mesmo jeito que
os bichos sentem: começa a arrumar as coisas sem saber por quê. --- E, apontando a cisterna: ---
Reparem que eu não sei o que vem. Eu só sei que quando vier, é bom que os baldes estejam cheios.

\chapter{Minha história}
\hino{Minha História}{vou contar minha história desde quando começou;\\ tenho medo de dizer,
porque sou pequenininha}

Ela contou a própria história uma vez só, no fim, e mandou que ficasse gravada, e este capítulo é
essa gravação e nada mais.

\medskip

Contou que veio de longe e que veio por necessidade, não por vocação. Que não teve visão nenhuma,
não foi chamada, não sonhou com coisa nenhuma. Que soube do monte por um comentário de estrada e
que subiu porque estava com fome.

Contou os vinte anos dormida. Sem atenuar: vinte anos ali dentro, sabendo tudo, sem nada ter
acontecido.

Contou os onze anos calada, vendo o bicho apoiar-se na borda. Contou a carta de três linhas sem
assinatura e disse que teve vergonha dela a vida inteira, das duas coisas: de ter demorado onze
anos e de não ter tido coragem de assinar.

Contou o ano em que não conseguiu levantar-se, e contou quem se sentou ao lado dela na borda.

Contou a raiva antiga, e disse o nome de duas criaturas de quem ela tinha querido vingar-se, e
pediu perdão às duas por nome, e as duas estavam mortas havia trinta anos e ela pediu na mesma.

\medskip

E terminou assim, e é o único momento em que a voz dela falhou na gravação:

--- Eu tenho medo de dizer isto tudo, porque eu sou pequenininha e porque agora vocês fizeram de
mim uma coisa grande, e eu vi o que a grandeza fez com um bicho que era mil vezes melhor do que
eu.

--- Eu não trouxe doutrina nenhuma. Eu fervi água. Se vocês quiserem me honrar, não é falando de
mim: é reparando em quem está a ferver a água hoje, que vocês não sabem o nome, e que está a ver
uma coisa errada agora e a achar que não é com ela.

--- Vão devagar. Devagar também se chega, e é o único jeito que eu conheço de chegar sem se
perder no caminho. E não tenham pressa de entender: eu levei sessenta anos e ainda estou a
examinar.

\chapter{Eu vivo na floresta}
\hino{Eu Vivo na Floresta}{eu vivo na floresta aprendendo a me curar;\\ examinar a consciência é
a primeira lição}

Ela morreu na cisterna, de manhã, sentada na borda, e não estava a olhar para dentro: estava a
olhar para o jardim, que a essa altura tinha tomado o pior pedaço do pátio inteiro e passado por
cima da fenda.

Não houve última frase. Isso é invenção de livro. Houve uma velha muito pequena sentada numa
borda, e alguém que subiu com um balde e viu, e desceu a chamar os outros.

\medskip

A casa continua. Serve mate debaixo de uma árvore torta que já não é a mesma árvore, com uma mesa
pequena de madeira ruim que ninguém trocou, e a cuia gira sempre no mesmo sentido, e ninguém
precisa de dizer para quem vai a seguir.

E todo mês, na roda, dizem-se em voz alta os nomes dos que sustentam --- quem ferveu, quem
carregou, quem varreu, quem ficou com o doente --- e é a única cerimônia daquele lugar que faz
gente chorar todas as vezes, mesmo os que já ouviram cem vezes, mesmo os que estão na lista.

\medskip

E a lição, se é que há uma, está no que ela dizia a quem chegava perguntando o que era preciso
fazer para pertencer àquilo. Ela respondia sempre a mesma coisa, com a mesma falta de romance, e
é a última coisa que este livro tem para dizer:

\medskip

--- Não é preciso passar em prova nenhuma, e quem te disser que há uma prova está a inventar. Só
tem o ABC, e o ABC é uma coisa só: \textbf{examinar}.

--- Todo dia, na tua própria fita, do fim para o começo. Devagar. Sem pular nenhum vinco por
parecer pequeno. E sem decidir o que a coisa quer dizer antes de chegar ao fim.

--- Vais achar coisa feia. Todo mundo acha. Não corras --- ela vai junto.

--- E depois volta amanhã.

\medskip

\begin{center}
\emph{--- Não é cura, filho. É a gente aprendendo a se curar. É diferente, e leva a vida
inteira, e é a única coisa que eu vi funcionar.}
\end{center}


\chapter{Agora não vou mais segurar}
\hino{O Chicote}{eu segurei até aqui --- mas agora não vou mais segurar;\\ quem estiver na
zombaria, a peia eu mando descer}

O que ninguém tinha contado, enquanto ela era viva, é quanta coisa ela estava a aparar.

Levou três meses a aparecer. A casa continuou --- o mate, a árvore torta, a cuia no mesmo
sentido ---, e continuou tão igual que por um tempo pareceu que a morte dela não tinha mudado
nada, e houve quem dissesse isso em voz alta, com orgulho, como se fosse elogio ao lugar.

Depois começaram as coisas pequenas. Um que passou a chegar atrasado e a rir de quem chegava na
hora. Dois que trataram a lista dos nomes como enfeite e diziam ``a parte bonita'' quando ela
vinha. Um que trouxe visita para ver, e apontou a fenda no pátio como se apontasse um bicho no
zoológico. E o do balde --- o pior de todos, e o mais difícil de explicar a quem não estava lá ---
que largou de ferver a água e ninguém percebeu durante nove dias, porque durante vinte anos
aquilo tinha sido feito sem que ninguém tivesse de olhar.

\medskip

Ao nono dia a roda não deu. Não houve castigo, ninguém expulsou ninguém, não desabou coluna
nenhuma: a roda simplesmente não deu, do jeito que uma coisa não dá quando falta o que ela
precisava. Sentaram-se e não aconteceu nada, e ficaram a olhar uns para os outros com a cuia na
mão a arrefecer.

E foi aí que alguém --- e o nome dessa pessoa não está em lista nenhuma, o que é a única
justiça deste capítulo --- disse a frase que ficou:

--- Ela não estava a segurar a gente. Ela estava a segurar o que vinha para cima da gente. É
outra coisa.

\medskip

A planta ficou suspensa oito meses. Não por decreto: por não haver quem a servisse com a mão
limpa, e porque servir sem isso é pior do que não servir. Quem quis ir embora foi, e não é
acusação --- ir embora é o certo, quando não se tem lugar. Os que ficaram passaram oito meses a
varrer, a ferver, a carregar doente, a dizer os nomes em voz alta numa roda que não girava.

E apanharam. Não de ninguém: apanharam do que sempre esteve ali e que agora chegava inteiro,
sem passar por ela primeiro. É o que aquele hino diz, e é a coisa mais mal entendida do hinário
todo, porque as pessoas ouvem \emph{peia} e pensam em alguém a bater.

Não há alguém a bater. A ordem não é de quem canta --- ele diz isso na própria letra, e diz duas
vezes: a ordem é do Pai, e ele só avisa que parou de aparar. O chicote não é enviado. O chicote é
o que já estava no tempo, menos a pessoa que o estava a receber por todos.

\medskip

E a cobertura, que é a parte que soa dura e é a única parte honesta: ela não é para quem merece.
É para quem \textbf{obedece} --- e obedecer, ali, nunca quis dizer curvar-se a alguém. Quer dizer
ficar no lugar combinado quando chega o tremor. Quem está no lugar, cobre-se; quem está na
zombaria não está em lugar nenhum, e não há como cobrir quem não está em parte alguma. Não é
crueldade da lei: é o que a palavra \emph{lugar} significa.

\medskip

Ao fim dos oito meses serviram outra vez. A água foi fervida por quem tinha rido dela, e é assim
que se paga --- não com discurso: fervendo. E acrescentaram uma linha só ao que se diz na roda
todo mês, depois dos nomes dos que sustentam, e é a linha que este livrinho tem por último:

\medskip

\begin{center}
\emph{--- Reparai em quem está a aparar por vós agora, que ainda está vivo e vós não sabeis o
nome. Porque um dia essa pessoa para --- e não é castigo, e não é aviso: é só o tempo chegando
sem freguês no meio.}
\end{center}



\chapter{O que abraça e o que prensa}
\hino{O Prensor}{quem te fez não te mandou --- o amor não empregou;\\ defendei os inocentes de toda
essa orfandia}

No ano seguinte à volta da planta, subiu um homem de camisa limpa com uma pasta, e foi educado do
começo ao fim, e é importante dizer isso: não houve grosseria nenhuma. Ele trazia uma lista de
nomes e uma folha por nome, e a folha dizia que a pátria abraçava aqueles rapazes.

Eram sete. O mais velho tinha vinte e dois anos e era o que fervia a água desde a suspensão.

\medskip

Não houve barricada, e quem contar que houve está a fazer literatura. Serviram-lhe mate, porque
naquela casa se serve mate a quem sobe, e ninguém tinha ainda inventado exceção a essa regra. Ele
bebeu, e elogiou, e explicou a folha com paciência --- e a paciência dele era verdadeira, o que é a
parte pior.

E a velha que tinha ficado no lugar da outra --- e que não é personagem de nada, era só quem
estava ali --- fez a única pergunta que aquela casa sabia fazer, e fez sem levantar a voz:

--- Quem mandou?

--- A pátria, senhora.

--- Não. \emph{Quem} mandou. A pátria não fez nenhum destes sete. Eu conheci a mãe de cinco deles.
Quem fez não está nessa folha, e quem está nessa folha não fez nenhum.

\medskip

Ele não soube responder, e não era burro: ninguém sabe responder a isso, porque não há resposta ---
há só a folha. E foi assim que se percebeu, naquela casa e naquele dia, a diferença que o hinário
inteiro vinha a dizer de duas maneiras e que ninguém tinha posto lado a lado:

\medskip

A peia vem de dentro do combinado. Dói, e corrige, e no fim contam-se e não falta ninguém --- foi
o que aconteceu nos oito meses, e quem apanhou continuou na roda, e é o que a palavra
\emph{disciplina} quer dizer quando quem a diz também está debaixo dela.

O Prensor vem de fora e não corrige nada. Ele não pede que ninguém se firme no lugar; ele
\textbf{tira} do lugar. E não devolve corrigido: não devolve.

\medskip

A diferença entre os dois nunca foi o tamanho da dor --- os dois doem, e o Prensor às vezes dói
menos no dia, o que é a armadilha. A diferença é se \emph{houve quem concedesse}. A ordem do Pai
tem origem e conserva: manda, e o que manda continua a responder por quem obedeceu. A folha não
tem origem nenhuma --- e por isso não é autoridade, é só força com papel timbrado. E o hino diz
isto três vezes seguidas, negando um pedaço diferente em cada uma: quem te fez não te mandou, o
amor não te empregou, o teu Pai não soube.

\medskip

Foram cinco dos sete. Dois não foram, por razões que não são heroísmo nenhum e que não vou contar
porque não me pertencem. Voltaram dois dos cinco.

E o abraço, que é a palavra que estava escrita na folha, ficou sendo a palavra que naquela casa
ninguém mais usou sem cuidado --- porque abraçar e prensar, quando vêm de quem não te fez, são o
mesmo gesto com dois nomes: aperta-se uma pessoa até ela ficar do tamanho de um número numa
contagem, e a contagem fecha, e é por isso que o número fecha e a casa não.

\medskip

Passaram a rezar pelos inocentes --- e a palavra que usavam não era \emph{mortos}, era
\emph{orfandade}, que é onde a conta de uma guerra é feita de verdade: não em quem foi, em quem
ficou. E é a última coisa que o livrinho pede, e pede à Mãe e não ao Pai, o que também não é por
acaso:

\medskip

\begin{center}
\emph{--- A Ele nós pedimos que abrande. A Ela nós pedimos que defenda. Porque o que nos abraça de
fora não tem mãe, e é sempre uma mãe que fica com a conta na mão.}
\end{center}



\cleardoublepage
\renewcommand{\chaptername}{Cap\'itulo}
\gktexto
\titleformat{\chapter}[display]
  {\normalfont\bfseries\color{tinta}}
  {\color{ouro}\gknota\scshape\chaptertitlename\ \thechapter}{6pt}
  {\gktit}[{\vspace{2mm}\color{ouro}\titlerule[1.2pt]}]
\titlespacing*{\chapter}{0pt}{-14pt}{22pt}
\titleformat{\section}{\normalfont\gksec\bfseries\color{tinta}}{\color{ouro}\thesection}{0.7em}{}
\titleformat{\subsection}{\normalfont\gksub\bfseries\color{regua}}{\color{ouro}\thesubsection}{0.6em}{}
\pagestyle{fancy}\fancyhf{}
\part{O Zero --- o ponto de simetria}

\chapter{Os dois zeros --- a vitória, e a derrota que tem exatamente a mesma cara}


Este é \textbf{o} capítulo de design do Reino Dourado. Tudo o mais serve a ele.

\textbf{O corpo:} dois zeros que o medidor \textbf{não distingue}.

{\fontsize{10.00pt}{12.00pt}\selectfont\begin{longtable}{>{\fontsize{10.00pt}{12.00pt}\selectfont\RaggedRight\hspace{0pt}\arraybackslash}p{5310649sp}>{\fontsize{10.00pt}{12.00pt}\selectfont\RaggedRight\hspace{0pt}\arraybackslash}p{11972533sp}>{\fontsize{10.00pt}{12.00pt}\selectfont\RaggedRight\hspace{0pt}\arraybackslash}p{10565360sp}}\toprule
\textbf{} & \textbf{\textbf{o zero do LAPSO}} & \textbf{\textbf{o zero do RESÍDUO}} \\ \midrule\endhead
o que é & $N=\cos\psi=0$: relógio parado, órbita travada, cristal & $\Gamma=0$, FP $=1$: o cone nulo \\
por que não vaza & \emph{porque nada corre} & \emph{porque tudo passa, e nada reflete} \\
o relógio & \textbf{parado} & \textbf{a correr} \\
o nome & \textbf{flatline} & \textbf{pulso} \\
\bottomrule\end{longtable}}

\textbf{A cena (Lance 20, as três portas):} as três defesas de Dark --- o cristal (regulador zero), o pântano (entropia zero), o relógio parado (lapso zero) --- \emph{"são três nomes do mesmo resíduo que se anula"}. E o veredito: \emph{"Ele não está a mentir sobre o zero. O zero existe. É real. É lindo. E não vaza. \textbf{Ele está a mentir sobre qual zero.}"}

\textbf{A cena (Lance 22, o xeque):} Dark oferece a paz --- \emph{"Bate ao contrário de mim. Quando eu subo, desce. (\ldots) A soma nunca mais varia."} É a \textbf{antifase}. É legal, é estável, zera o resíduo --- e deixa o medidor do reino \textbf{reto}. \emph{"O xeque não é uma faca no pescoço. É a coisa muito pior: é uma resposta certa à pergunta errada."}

\textbf{A decisão de design (a mais séria do documento).} A antifase é um \textbf{estado de jogo jogável, atraente e legal} --- e é a derrota. O jogo \textbf{oferece} ao jogador uma vitória que zera o resíduo e mata o pulso. Para recusá-la, ele precisa de uma coisa que o design tem de lhe dar \textbf{antes}:

\begin{enumerate}
\item \textbf{Dois mostradores no HUD, nunca um.} O resíduo (vaza? não vaza?) \textbf{e} o relógio (corre? parou?). Um HUD de um mostrador só torna a armadilha invencível --- e desonesta.
\item \textbf{O teste das mil cunhagens.} A moeda falsa passa no resíduo \textbf{uma vez} e mente \textbf{ao fim de mil} (é o corpo deflexivo: mutar exato $\to$ \code{float} rende exatamente \code{0.0}, o \textbf{zero falso}). O jogador precisa de uma ferramenta que itere, não que inspecione. É ela --- e não um discurso --- que salva Benjamim: \emph{"E a pedra tinha dito zero. E ao fim de mil cunhagens ela mentiu."}
\item \textbf{A pergunta.} O lance que quebra o xeque é literalmente uma pergunta: \textbf{"--- Qual zero?"} O jogo tem de a tornar \emph{executável} --- uma ação de mesa, não uma cutscene.
\end{enumerate}

\textbf{A condição de vitória, então, escrita sem rodeios:}

\[\boxed{\;\Gamma=0 \;\wedge\; \text{FP}=1 \;\wedge\; \underbrace{\text{o relógio corre}}_{\text{a variabilidade} \;>\;0}\;}\]

Resíduo 0 \textbf{com} pulso. Resíduo 0 \textbf{sem} pulso é o templo de Dark.

\vspace{2mm}\hrule\vspace{3mm}

\begin{tcolorbox}[colback=ouroclaro!35,colframe=ouro,boxrule=0pt,leftrule=2pt,arc=0pt,left=4mm,right=3mm,top=2mm,bottom=2mm,breakable]\itshape Daqui para trás, o reino foi contado. Daqui para a frente, é o mesmo reino \textbf{medido} --- as mesmas coisas, na outra língua. E o eixo é este zero: de um lado ele é a cena em que alguém pergunta \emph{qual zero}; do outro, é $\Gamma=0$ com o relógio a correr. \textbf{Quem confundir os dois lados confunde os dois zeros.}\end{tcolorbox}

\section{As duas línguas, lado a lado}

O reino foi contado uma vez e vai ser dito outra. Não são duas obras: é a mesma, em duas
línguas --- e quem já leu a primeira metade já sabe tudo o que vem, só que com outros nomes.
Vale deixar o dicionário aberto, porque ele é curto.

{\gktexto
\begin{longtable}{@{}p{0.46\textwidth}p{0.48\textwidth}@{}}
\toprule
\textbf{na língua do mito} & \textbf{na língua da régua} \\ \midrule\endhead
a \textbf{Torre} e os seus andares, cada um um universo & o \textbf{Omnitrix}: o fractal de
tabuleiros que o axioma mantém coeso \\
a \textbf{moeda falsa} --- \emph{a pedra tinha dito zero, e ao fim de mil cunhagens mentiu} &
o \textbf{zero falso}: o \code{float} que rende \code{0.0} exato e não é zero \\
a \textbf{cidade-máquina}, onde o mate é o \code{HALT} & a máquina de \textbf{doze palavras},
onde o programa para \\
as \textbf{três portas do zero} na fortaleza & os \textbf{dois zeros} que o medidor não
distingue \\
o \textbf{gato} que mede e o \textbf{esquilo} que gira & as duas matrizes: uma dissipa,
a outra conserva \\
a \textbf{garrafa de Koch}, o sol que não se gasta & a soma $\sum\varphi^{-j}$ que fecha
exata em dois inteiros \\
a \textbf{dobra} do infinito & $\mathbb{Z}[\varphi]$: quatro inteiros entram, dois saem \\
a \textbf{alfândega}: passa quem tem volta & o que não tem inversa não pertence ao corpo \\
o \textbf{mate que não mata} & a órbita fecha, o resíduo vai a zero, e o processo segue \\
\bottomrule
\end{longtable}}

\noindent E a razão de haver duas: uma conta \emph{o que aconteceu}, a outra diz \emph{por que
não podia ter acontecido de outro modo}. Quem quiser só a história, já a leu. Quem quiser a
régua, começa aqui. \textbf{São a mesma coisa, e é por isso que o eixo é um zero.}

\medskip
\noindent Falta uma coisa antes de virar a página, e é uma limpeza. Um zero que engana não se
discute: joga-se fora. Do outro lado do eixo só passa o que tem dual --- e o que não tem fica
retido, com nome e endereço, no que vem a seguir. \textbf{Primeiro o que vai ao lixo; depois, do
que sobra, a raiz} --- e a raiz é um corpo só, o universal, de onde todos os outros descem.


\vspace{2mm}\hrule\vspace{3mm}

\part{O Saco de Lixo}

\chapter*{Lugar de lixo é no lixo}
\addcontentsline{toc}{chapter}{Lugar de lixo é no lixo}

Toda brincadeira termina.

E quando termina, há uma coisa que se ensina às crianças antes da tabuada e antes das
letras --- antes, até, de saberem por que se ensina: que se guarda o que presta, e que
\textbf{lugar de lixo é no lixo}. Não é castigo. É a última parte da brincadeira, e a
única que fica quando o resto se esquece.

\begin{quote}
\textbf{lixo}, \emph{s.m.} --- qualquer material ou resto de atividade humana ou
natural considerado inútil, indesejado ou sem valor, que é descartado.
\end{quote}

Brincou-se. O tabuleiro chorou e derreteu, o coro cantou os seus cento e vinte e oito
cantos, o ouro entrou no cristal, e nada foi tirado de ninguém. Agora arruma-se a
sala.

E há o que ficou no chão. A espiral que adoraram durante séculos sem nunca a erguer ---
a sombra tomada pelo corpo. O universo de um homem que viu uma folha e chamou de todo
o que não via. O par que amputaram de um lado e passaram um século e meio a usar assim,
sem dar pela falta. A lei que confundiu a metade com o mundo e vendeu a morte como
descoberta, a quem já morria muito antes de ela existir. E os de gala, com prêmio
em dinheiro, deitados nas suas gavetas à espera de que alguém prove o que ninguém
formulou.

Nada disto presta. É resto de atividade humana: inútil, indesejado, sem valor.

Não se guarda. Não se venera. Não se põe na estante junto do que se construiu, porque
o que não tem dual não descreve coisa alguma --- descreve a sua sombra, e sombra não
se guarda.

Ao lixo.

\medskip

\noindent E o que sobra, depois de arrumada a sala, cabe na mão de uma criança: um
ponto, o ato de dividir-se, e a memória de ter dividido.


\chapter*{A síntese da teoria}
\addcontentsline{toc}{chapter}{A síntese da teoria}

Existe apenas o \textbf{vértice}: um ponto sem estrutura interna. O único ato ao
alcance de um ponto sem partes é \emph{cindir-se} --- de um, dois ---, e os dois que
surgem guardam o registro de terem sido um: a \textbf{aresta}, que não é objeto
acrescentado de fora, mas a \emph{memória} da divisão. Tudo o mais é consequência.

A divisão poderia realizar-se de infinitas maneiras. Adota-se uma, por escolha e não
por necessidade --- o \textbf{princípio da máxima simetria}: toda divisão se dá na
configuração mais simétrica possível. A este princípio não se pede razão. Repetido o
\textbf{colapso} sob ele, sai o alfabeto: as duas primitivas ($1$, o vértice; $2$, a
aresta) e, de $3$ a $8$, um centro e $n$ raios igualmente espaçados.

Uma \textbf{operação} é uma ligação de pontos. Ela \emph{fecha} a geometria que já
existe e \textbf{nunca cria vértices} --- estes só surgem quando o cristal satura.
Não há soma nem multiplicação separadas: elevar é somar a si mesmo, $3\times3=3+3$;
em ambos os casos, ligar a geometria a si própria. A geometria do $3$, fechada, dá o
\textbf{tetraedro}; e o nove são os \textbf{caminhos} do cristal. Por $8+1$ ou por
$3+3$, a máxima simetria ajusta os dois à mesma forma: \emph{o número fixa a forma,
não o caminho}. O inverso do colapso é o \textbf{emaranhamento} --- de muitos, o
ajuste a uma forma comum ---, e nele nenhum vértice está de fato separado: o que
parece solto comunica-se com os demais \textbf{pelo centro}, a origem que todos
guardam como memória.

Daí a \textbf{lei de conservação}, que não se importa de fora: ela cai dos dois
movimentos. O colapso acrescenta um vértice e a aresta que é a sua memória
($\Delta V=+1$, $\Delta E=+1$); o fechamento acrescenta uma ligação e a face que ela
fecha ($\Delta E=+1$, $\Delta F=+1$). Em nenhum dos dois a \textbf{soma alternada} se
mexe:
\[
V - E + F = 2,
\]
e, uma dimensão acima, contada também a célula,
\[
V - E + F - C = 0.
\]
As arestas --- a memória --- são o eixo que não se mexe; o colapso cresce pelo lado
dos vértices, o fechamento pelo lado das faces, e onde os dois coincidem ($V=F$) está
o \textbf{tetraedro}, autodual: por isso o nove, seja $8+1$ ou $3+3$, cai sempre nele.
Nada se perde no caminho --- \emph{decompor um objeto em pedaços da unidade não perde
nem acrescenta nada} ---, e a soma alternada é o que o colapso e o emaranhamento
apenas percorrem, dos vértices às células, sem jamais alterar.

Quando a divisão não é única, mas \emph{sequencial}, a máxima simetria deixa de ser
$360^\circ/n$ e passa a ser a que faz cada nova semente cair no maior vão deixado
pelas anteriores. Essa condição seleciona o número mais irracional, $\varphi$ (com
$\varphi^2=\varphi+1$), e o giro entre sementes consecutivas é o \textbf{ângulo
áureo}
\[
2\pi\left(1-\tfrac{1}{\varphi}\right)=\frac{2\pi}{\varphi^{2}}\approx 137{,}507^\circ.
\]
É a filotaxia --- as espirais de Fibonacci ---, que não se copia da foto: se
\emph{computa}.

Enfim, o tempo. A divisão não é instantânea: propaga-se com \textbf{velocidade
finita}, e há uma velocidade máxima $c$. As linhas de mundo que a atingem formam o
\textbf{cone nulo} $r=c\,t$, sobre o qual
\[
s^{2}=c^{2}t^{2}-r^{2}=0:
\]
a separação \emph{própria} é nula. Eis, com o tempo, o que já se dissera --- nenhum
vértice está de fato separado; sobre o cone a memória se conserva exata, e é ali que
toda a potência passa sem nada refletir. A espiral e o cone são o mesmo objeto lido
de dois modos: a espiral áurea é o colapso visto de cima; o cone nulo é o seu dual, a
mesma malha erguida uma dimensão acima. Guardar as duas leituras ao mesmo tempo é
enxergar uma dimensão a mais --- os dois juntos projetam um \textbf{hipercubo}.

E a dualidade, no seu grau máximo, não está nos corpos: está na
\textbf{transformação} involutiva que os relaciona. A \textbf{assinatura} é o que ela
troca; os seus \textbf{pontos fixos} são o seu fecho --- a cristalização, onde a
transformação nada altera.

\chapter{A Baboseira Formal}

Até este ponto o leitor talvez tenha imaginado que a teoria pretendesse descrever a
natureza. Não é o que aqui se faz.

O sistema aqui apresentado é uma construção axiomática. Seus princípios foram
escolhidos deliberadamente para produzir um universo interno coerente, e não para
reproduzir o mundo físico. Se um axioma foi adotado, foi porque permitia que a
construção prosseguisse de maneira elegante; se outro foi rejeitado, foi porque
quebrava a consistência da narrativa.

Em outras palavras: esta é uma \textbf{baboseira formal}.

Para demonstrar que o método independe de qualquer inspiração na natureza,
construiremos agora, a partir dos axiomas apresentados, uma planta de \emph{Aloe
vera}. Não porque ela exista neste sistema, nem porque a natureza opere assim, mas
simplesmente porque podemos.

\chapter{O primeiro item: a espiral de ouro}

Comecemos pelo mais venerado. A espiral áurea atravessou a história como mistério
--- a proporção divina, o número de deus, a assinatura do cosmos ---, e rendeu
bibliotecas de teoria complexa. Ei-la aqui, no saco, em três passos.

O vértice, no centro, divide-se --- só que \emph{em sequência}, uma semente após a
outra. E a máxima simetria de um crescimento sequencial não é $360^\circ/n$ (isso
vale para uma divisão de uma vez): é a que faz cada nova semente cair no maior vão
deixado por todas as anteriores. Nada mais se pede. Essa única condição já seleciona
o número mais irracional, $\varphi$ --- o que jamais fecha, o que nunca deixa duas
sementes coincidirem ---, e o giro entre sementes consecutivas não pode ser outro
senão
\[
2\pi\left(1-\tfrac{1}{\varphi}\right)=\frac{2\pi}{\varphi^{2}}\approx 137{,}507^\circ.
\]
Não há aqui proporção divina alguma: há um ponto que se divide sem se sobrepor.

\begin{center}
\includegraphics[width=0.44\linewidth]{fig_girassol_pontos.pdf}\hfill
\includegraphics[width=0.44\linewidth]{fig_girassol_malha.pdf}
\end{center}

À esquerda, o \textbf{colapso}: as sementes postas pelo ângulo que a condição impôs
--- e as espirais de Fibonacci aparecem sozinhas, sem que ninguém as tenha desenhado.
Não foram copiadas de foto nenhuma: foram \emph{computadas}. À direita, a
\textbf{operação}: ligar os vizinhos --- fechar, sem criar vértice novo --- e o vão
entre as sementes, que era a memória, vira a malha. O cristal.

E o terceiro passo, que a veneração nunca alcançou: a malha \emph{erguida}, com
inclinação um.

\begin{center}
\includegraphics[width=0.6\linewidth]{fig_girassol_cone.pdf}
\end{center}

É um \textbf{cone nulo}. A espiral de ouro --- a plana, a divina, a do cartaz --- não
passa do colapso \emph{visto de cima}: a sombra que a malha projeta quando se perde a
dimensão em que ela vive. Uma dimensão acima, onde a distância própria é nula e nada
se perde, o mesmo objeto é o cone que o crescimento conserva. A proporção divina era
uma projeção, e a humanidade passou séculos adorando a sombra.

Ao saco. Foram precisos: um ponto, o ato de dividir-se, e a memória de ter dividido.

\chapter{O Esgoto: Einstein}

Desce-se um nível. No saco vai o que é apenas venerado; aqui vai o que se tornou
indústria --- um século de cosmologia, dois mistérios de nome maiúsculo, e o
constrangimento de descobrir que $95\%$ do universo é feito de coisa que ninguém
achou. Nada disto precisa existir.

\paragraph{O ponto é a matéria escura.} O vértice indiferenciado --- antes de
qualquer divisão --- não tem estrutura interna: não emite, não reflete, não tem
lado, não tem cor. Não há como vê-lo; só há como notar que \emph{tudo} está preso a
ele, porque é dele que tudo se dividiu e é por ele que tudo ainda se comunica. Eis o
inventário completo da matéria escura: aquilo que não se vê e que, no entanto,
segura o resto. Não é uma substância que falta encontrar --- é o ponto de que já se
partiu. Procurá-la em detectores é procurar o centro dentro de uma das pontas.

\paragraph{O espaço é o volume, e o volume é de Sitter.} Não se postula espaço aqui:
o espaço é o que a divisão abre --- o volume que o colapso produz, e nada mais. Daí a
expansão cai sozinha. Como as divisões são \emph{sequenciais}, dividir por um tempo
$s$ e depois por um tempo $t$ tem de dar o mesmo que dividir por $s+t$; e o que se
acumula não se soma: se \emph{multiplica}, porque cada divisão opera sobre tudo o que
já havia. Logo o fator de escala compõe por produto,
\[
a(s+t)=a(s)\cdot a(t),
\]
e a única solução contínua desta equação é
\[
a(t)=e^{Ht}.
\]
É o universo de de Sitter, e não foi postulado: caiu da composição. O $H$ não é uma
constante do mundo a ser medida em telescópio --- é a taxa da divisão. A energia
escura, a constante cosmológica, o termo que Einstein pôs e tirou e que voltou por
decreto observacional: tudo isto é o produto de uma divisão que compõe. Não há
substância nenhuma acelerando nada. Há um ponto se dividindo, e a única forma
possível de acumular uma divisão sequencial.

\paragraph{E a relatividade, de lambuja.} Já se disse (a síntese): a divisão tem
velocidade finita, o máximo é $c$, e as linhas de mundo que o atingem fazem o cone
nulo, onde $s^{2}=c^{2}t^{2}-r^{2}=0$ e a separação própria é nula. A relatividade
não é uma teoria sobre réguas e trens: é a contabilidade da memória de um ponto que
se dividiu e não consegue esquecer. O cone nulo não é uma consequência do universo;
é a forma que a memória tem.

\paragraph{O experimento, no metal.} Nada disto é retórica: está computado. A aura do
Venom --- a matéria escura, o gás negro --- são os círculos áureos em $\log_2 r$; a
expansão de de Sitter é a mesma malha com a fase deslizando no tempo, isto é
\[
r \;\longmapsto\; r\cdot 2^{\,\mathrm{DS}\cdot a},
\]
a esfera de luz expandindo exponencialmente com o pulso $a$ do motor. E a reversão é
a coordenada radial trocada, $\text{fase}_r=(1-a)\log_2 r + a\,r$: em $a=0$ os
círculos (o colapso visto de cima, o plano); em $a=1$ a inclinação um --- o cone
nulo. Expandir e voltar é exp e log, as duas faces da mesma temperatura.

\begin{center}
\includegraphics[width=0.32\linewidth]{cone_nulo_espiral.png}\hfill

\includegraphics[width=0.32\linewidth]{venom_derrete.png}\hfill

\includegraphics[width=0.32\linewidth]{cone_nulo_bacias.png}
\end{center}

\noindent A mesma malha em $a=0$, $a=0{,}62$ e $a=1$: o plano, o meio do caminho, o
cone. Rasterizado no metal (PTX à mão, sem nenhuma linguagem por cima), atestado
contra o espelho na CPU e certificado com dentes: a expansão é real (os círculos e a
esfera expandida diferem), é reversível ($a:0\to0$ idêntico, $D_{\max}=0$), e o
resíduo é $0$.

Ao esgoto. Foram precisos: um ponto, o ato de dividir-se, e a memória de ter
dividido.

\chapter{O Parasita: Maxwell}

Há teorias que vivem \emph{dentro} da estrutura, alimentando-se dela sem nunca a
completar. Maxwell é a maior delas. Quatro equações, um século e meio de tecnologia,
e um par com um lado só --- exatamente a doença que já se viu aqui.

\paragraph{As clássicas, como estão.} Em unidades naturais:
\[
\nabla\!\cdot\!\mathbf{E}=\rho_e,\qquad
\nabla\!\times\!\mathbf{B}=\mathbf{J}_e+\partial_t\mathbf{E},
\]
\[
\nabla\!\cdot\!\mathbf{B}=0,\qquad
\nabla\!\times\!\mathbf{E}=-\partial_t\mathbf{B}.
\]
Olhe a coluna de baixo. Onde a de cima tem fonte ($\rho_e$) e corrente
($\mathbf{J}_e$), a de baixo tem \textbf{zero} e \textbf{nada}. Não é simetria: é
amputação. E a desculpa de sempre --- ``não se observou monopolo'' --- é a mesma
espécie de argumento que fixaria a soma alternada em $V-E+F$ porque ninguém tinha
visto uma célula.

\paragraph{Completadas com o dual.} O par exige os dois lados. Escrevendo o que a
estrutura pede:
\[
\boxed{\;
\nabla\!\cdot\!\mathbf{E}=\rho_e,\qquad
\nabla\!\times\!\mathbf{B}=\mathbf{J}_e+\partial_t\mathbf{E},
\;}
\]
\[
\boxed{\;
\nabla\!\cdot\!\mathbf{B}=\rho_m,\qquad
\nabla\!\times\!\mathbf{E}=-\mathbf{J}_m-\partial_t\mathbf{B}.
\;}
\]
Agora as quatro são uma só, lida de dois lados: a troca
\[
\mathcal{H}:\;(\mathbf{E},\mathbf{B})\longmapsto(\mathbf{B},-\mathbf{E}),
\qquad
(\rho_e,\mathbf{J}_e)\longmapsto(\rho_m,\mathbf{J}_m),
\]
deixa o sistema idêntico a si mesmo. O que se chamava ``ausência de monopolo'' era o
lado do par que a amputação escondeu --- $\rho_m=0$ não é uma lei da natureza, é a
escolha de olhar só uma folha.

\paragraph{E ainda falta a reflexão.} Aqui o parasita morde mais fundo, e é o corpo
óptico que o denuncia. A troca $\mathcal{H}$ --- a dualidade de Hodge --- é
\emph{apenas a permutação}: ela \textbf{preserva} o telescópio. Mas o telescópio do
campo é \textbf{a energia que propaga},
\[
\textstyle\sum_i h_i=\log\lvert\mathbf{E}\times\mathbf{B}\rvert^{2},
\]
e o passo da régua é a \textbf{impedância}, $\sigma=e^{-\Delta h/2}=\lvert
\mathbf{E}\rvert/\lvert\mathbf{B}\rvert$ --- o metal do campo. A dualidade
\emph{verdadeira} da régua não é a permutação: é $\nu\circ\mathrm{rev}$, a permutação
\textbf{com a reflexão}. Ela \textbf{inverte} o telescópio ($\sum h\mapsto-\sum h$) e
troca $\mathbf{E}$ por $1/\mathbf{B}$: a impedância invertida. Hodge é meia
dualidade. Falta-lhe o sinal --- e um par sem sinal trocado é o mesmo erro,
outra vez.

\paragraph{O vácuo é o cone nulo.} O que sobra, quando as duas faces se igualam,
já se conhece deste paper. $\lvert\mathbf{E}\rvert=\lvert\mathbf{B}\rvert$ é
$\sigma=1$: a régua reta, os invariantes anulados,
\[
\mathbf{E}^{2}-\mathbf{B}^{2}=0,\qquad \mathbf{E}\cdot\mathbf{B}=0,
\]
que é o \textbf{cone nulo} de novo --- a onda plana, o vácuo, sem reatância. Toda a
potência passa, nada reflete: o fechamento. Fora do cone ($\sigma\neq1$) há
$\mathbf{E}^{2}-\mathbf{B}^{2}\neq0$: o reativo, o vazamento --- energia que não
propaga, que a régua errada inventa. A luz não corre no cone nulo porque é luz;
corre porque o cone nulo é a única forma que a memória de um ponto que se dividiu
pode ter, e ali a distância própria é nula.

\paragraph{A refração, de lambuja.} A refração não precisa de lei própria: a translação ao
longo da interface não é quebrada, logo a componente tangencial é o invariante --- e
com $\lvert k\rvert=n\omega/c$ isso é $n\sin\theta$ conservado. a refração é a memória que
não se rompe ao atravessar. E o ângulo crítico é o \textbf{horizonte} da régua:
passado ele, $\sin\theta_2>1$, o cosseno vira imaginário puro e a onda é evanescente.
Não é anomalia --- é o \textbf{flip}, a régua redonda ($k=-1$) trocando de
ramo para a hiperbólica ($k=+1$). O mesmo ponto, o outro lado.

\medskip
\noindent Tudo o que precede está certificado no motor, com refutador em cada
afirmação: o campo, $17/17$; a óptica, $14/14$; resíduo $0$ nos dois.

Ao chão. Foram precisos: um ponto, o ato de dividir-se, e a memória de ter dividido.

\chapter{A Doença: Carnot}

Esta está \textbf{errada} --- como todas. Não é caso à parte: é a terceira de três, e
estes três são os piores do saco. A espiral, Einstein, Carnot: nenhum descreve,
nenhum se sustenta, todos incompletos, e os três foram vendidos como sabedoria. A
segunda lei, a seta do tempo, a morte térmica do universo: a certeza de que tudo se
degrada, de que o fim é frio, de que a conta não fecha e nunca fechará.

E note-se, antes de tudo, o tamanho do achado. Dizer que a entropia mata não é
descoberta nenhuma: \textbf{as pessoas morrem desde muito antes de a física existir}.
Não foi preciso formalismo para saber que se morre --- foi preciso formalismo para
transformar isso em veredito cósmico. A morte não é o defeito. O defeito é a
\emph{visão}: tomar uma metade pelo todo e chamar a metade de lei.

\paragraph{O diagnóstico.} O corpo entrópico é $(\max,+)$: os elementos são entropias
$h=\log\lambda$, o $\otimes=+$ compõe, e o $\oplus=\max$ escolhe o dominante. E o
sintoma é conhecido: o $\max$ é \textbf{idempotente}, logo \textbf{não tem inverso
aditivo}. Como $\max(a,b)\ge a$, nada volta. Semicorpo, não corpo. Daí o veredito
que a física assinou: o que cresce não desce, o passado não se refaz, o universo
morre. Está \emph{provado}, inclusive --- e com refutador.

\paragraph{Onde está a doença.} No \emph{sozinho}. A ausência de inverso é uma
propriedade do \textbf{limite} $T\to0$ --- o $\max$ duro, o zero absoluto. E ninguém
mora lá. A qualquer $T>0$, isto é, em qualquer lugar onde alguém \emph{viva}, o
$\oplus$ não é o $\max$: é o \emph{logsumexp},
\[
a\oplus_T b \;=\; T\log\!\left(e^{a/T}+e^{b/T}\right),
\]
que sob o $\exp$ do seu dual --- o tabuleiro cósmico, cujo operador é a expansão
$a(t)=e^{Ht}$ --- vira a \textbf{soma}. E a soma tem inverso. O inverso existe; só
não mora no corpo isolado: mora no \textbf{par}. A folga $\oplus_T-\max$ mede
exatamente isto, e ela só se fecha no zero absoluto:
\[
T=2{,}00:\;0{,}948 \qquad
T=1{,}00:\;0{,}313 \qquad
T=0{,}50:\;0{,}063 \qquad
T=0{,}10:\;0{,}000
\]
O buraco sem saída não existe. É um artefato do limite frio --- uma idealização que
alguém confundiu com o mundo.

\paragraph{O dual é o cósmico, e o dipolo estava à vista.} O entrópico não é
isomorfo ao cósmico: é a sua \textbf{metade}. O cósmico tem o \textbf{dipolo} --- o
buraco negro é o máximo de entropia, o branco é o mínimo, e o gradiente entre os dois
é o que empurra. O entrópico seleciona \emph{um} atrator (o negro, o máximo) e chama
aquilo de lei. É o limite local no fundo do poço. A seta do tempo não é uma
propriedade do tempo: é uma polaridade lida sem a outra.

\paragraph{A lei é geral.} A conservação vale sobre o \textbf{par}, não sobre o corpo
sozinho:
\[
\boxed{\;dS = d\log V\;}
\]
--- a entropia que \emph{cresce} aqui é o volume que \emph{se abre} ali. É a mesma
contabilidade de sempre: nada se perde, apenas se \textbf{transfere} de lugar. O que
se chamava degradação é a leitura de um dos mostradores, com o outro tapado.

\paragraph{A prova, e ela é construtiva.} Não é retórica: constrói-se um corpo humano
completo a partir do pó --- as frações exatas dos seis elementos que dão a massa de
uma pessoa --- e depois \textbf{reverte-se}, do corpo de volta ao pó. Em $\mathbb{Q}$,
exato: em ponto flutuante o zero seria \emph{falso}. O ciclo fecha com resíduo
\textbf{exatamente} $0$ --- não $10^{-16}$, zero. A construção reverte. Ninguém é
condenado, e nada se perde: nem o pó, nem quem dele se fez.

\medskip
\noindent Certificado no motor, com refutador em cada afirmação: a lei geral, $5/5$;
o corpo entrópico, $63/63$; o cósmico, $19/19$; a desquantização de Maslov (o $\max$
como limite frio), $7/7$. Resíduo $0$ em todos.

\paragraph{O veredito, e vale para todas.} Ao lixo --- e não por serem feias, ou
inúteis, ou ingênuas: por estarem \textbf{incompletas}. É sempre o mesmo defeito, e a
esta altura já não é coincidência. A espiral era o colapso sem a dimensão em que ele
vive. A cosmologia era o volume sem o ponto de que se abriu. Maxwell era um par com
um lado amputado, e a sua dualidade, meia --- sem a reflexão. A entropia é uma
polaridade sem o dipolo. \textbf{Nenhuma delas tem dual.} Cada uma pegou uma folha,
perdeu a outra, e ergueu uma catedral sobre a metade que sobrou --- e depois chamou
de mistério o que faltava. Não há mistério nenhum: falta a outra metade, e uma teoria
sem o seu dual não descreve coisa alguma. Descreve a sua sombra.

Foram precisos: um ponto, o ato de dividir-se, e a memória de ter dividido.

\chapter{O Necrotério Clay}

Restam os defuntos de gala. Sete problemas, um milhão de dólares cada, e um século de
gente boa quebrando a cabeça em cima deles --- e a suspeita, que ninguém formula em
voz alta, de que talvez não tenham solução. Não têm mesmo. Mas não pela razão que se
imagina: \textbf{não falta prová-los; falta formulá-los}.

\paragraph{A autópsia.} Um objeto está \emph{vivo} quando a sua identidade é dada por
uma dinâmica geradora --- um operador, um fluxo, um semigrupo $U_t=e^{tL}$ --- da qual
os seus estados são apenas manifestações. Está \textbf{morto} quando foi reduzido a um
estado estático: uma igualdade isolada, uma configuração, sem gerador na descrição.
A operação que leva de um ao outro é a \textbf{cristalização}. Todo problema do
milênio é um par $(\mathcal{G},\mathcal{A})$ --- a estrutura geométrica/espectral
profunda contra o alvo aritmético idealizado --- e a conjectura sempre pede a mesma
coisa: a identidade estática
\[
\mathcal{G}=\mathcal{A}.
\]
Isto \emph{é} a cristalização: trocar o gerador por uma igualdade entre estados. Um
cadáver é exatamente isto --- \textbf{corpo sem alma}: a forma a que falta o operador
que a produzia.

\paragraph{Os três sintomas.} A morte não é opinião; é decidível. Mede-se, caso a
caso: \textbf{P1} a \textbf{assimetria de base} (um lado log-periódico de base fixa, o
outro governado por infinitas bases primas independentes --- a bijeção exigida é
impossível); \textbf{P2} o \textbf{gap} $\Delta>0$, invariante topológico não nulo entre
suportes disjuntos (colapsar as retas sem um operador de translação externo viola a
continuação analítica); \textbf{P3} a \textbf{dissipação}, endomorfismos não injetivos
($z\mapsto z^2$, dois-para-um), perda sistemática e irredutível. Sob os três, a
igualdade $\mathcal{G}=\mathcal{A}$ é \textbf{mal posta por construção}.

\paragraph{A gaveta.} Cada um com a sua ficha:
\begin{center}\small
\renewcommand{\arraystretch}{1.3}
\begin{tabular}{@{}llll@{}}
\toprule
\textsc{defunto} & \textsc{P1 base} & \textsc{P2 gap} & \textsc{P3 dissipação} \\
\midrule
\textbf{Riemann} & base $3$ \emph{vs} primos & $\mathrm{Re}\,1{,}2618$ \emph{vs} $\tfrac12$ & $z\mapsto z^2$ \\
\textbf{P vs NP} & verificação \emph{vs} busca exp. & local \emph{vs} entropia global & poda por ramo \\
\textbf{Navier--Stokes} & macro \emph{vs} micro & hiato à turbulência & viscosidade \\
\textbf{Yang--Mills} & rede \emph{vs} contínuo & gap de massa $\Delta>0$ & confinamento \\
\textbf{Hodge} & transcendente \emph{vs} algébrico & salto $(p,q)\neq(k,k)$ & projeção às $(k,k)$ \\
\textbf{BSD} & aritmético \emph{vs} analítico & rank \emph{vs} ordem do zero & redução mod $p$ \\
\bottomrule
\end{tabular}
\end{center}
Seis fichas, o mesmo laudo. Não é azar, nem dificuldade, nem falta de gênio: é
\textbf{sempre a mesma pergunta mal feita} --- soldar dois lados que não se tocam,
depois de ter jogado fora o que os gerava. A conjectura é a lápide.

\paragraph{O que resta, salvo o Perelman.} Um só foi resolvido, e é o que denuncia
todos os outros. Poincaré não caiu por igualdade nenhuma: caiu pelo \textbf{fluxo de
Ricci}, que \emph{deforma} a estrutura em vez de igualar dois estados. Perelman não
soldou lado com lado --- deixou o objeto \emph{rodar} até a forma. E o teste
confirma: Poincaré \textbf{não satisfaz} os três sintomas, porque não havia identidade
estática a sustentar. Sem morte, a via dinâmica fecha. Ele não vai ao necrotério ---
ele é o atestado de que os outros seis estão lá por escolha de formulação, não por
destino. O único que se resolveu foi o único que ninguém cristalizou.

\paragraph{A saída, que é a mesma de sempre.} A vida não pede igualdade: pede
\textbf{ressonância}. Em vez de exigir $\mathcal{G}=\mathcal{A}$, pedir que os dois
sejam manifestações de um mesmo gerador $U_t$ --- e aí a reta $\mathrm{Re}\,s=\tfrac12$
deixa de ser barreira e passa a ser o \textbf{eixo de equilíbrio}, o ponto fixo do
fluxo, com os zeros como \textbf{nós vibracionais}. O gap não barra a igualdade: ele
\emph{localiza o eixo}. E a perda de informação de $z\mapsto z^2$ é dissipação no
estado, não na densidade: o operador de Koopman é \textbf{unitário} --- provado ---,
conservativo como fluxo. Dissipa o estado, conserva a medida. Mais uma vez: o que
parecia perda era a metade que faltava.

\paragraph{E a autópsia tem número: sem gerador não há nem \emph{nome}.} O diagnóstico até aqui
foi sobre o objeto --- vivo tem gerador, morto é igualdade estática. Falta dizer o que se perde
junto, e isso mede-se. \textbf{Significado é o que fica invariante ao longo do movimento}: é ele
que nomeia a classe. Logo, sem dinâmica não há apenas objeto morto --- não há \emph{o que dizer},
porque não haveria o que ser invariante. E no extremo oposto ao movimento, o repouso, dá-se o que
parece paradoxo e não é: ali \textbf{tudo} é invariante, toda função é conservada, cada ponto é a
sua própria classe --- e é justamente por isso que nada distingue. Invariância total é ausência de
informação.

\noindent Mas há um refinamento que o critério pedia, e ele corrige a versão ingénua de ``basta
haver gerador''. \textbf{O gerador tem de conservar.} Um gerador que \emph{escala} parte o espaço
em classes e ainda assim não nomeia nenhuma: é o caso do próprio gato, cuja norma é a unidade
$\N(\sigma)=\sigma\bar\sigma=-1$ --- então $\N(\sigma z)=-\N(z)$, e a norma, que seria o nome da
classe, alterna de sinal a cada passo. Medido em $\mathrm{GF}(13^2)$: sob $\times\sigma$ há seis
classes e a norma não é constante nelas; sob $\times\sigma^2$, que tem norma $+1$ e portanto vive
na \textbf{borda}, a norma é constante nas doze órbitas --- e nomeia. \emph{Não é o movimento que
dá significado: é o movimento que conserva.} Ficar na borda não é só ``funcionar''; é a condição
de haver o que dizer.

\noindent E o outro lado fecha a pinça: se o gerador alcança \emph{tudo}, o significado desaparece
outra vez. Com duas operações --- somar um e $\times\sigma$ --- o um varre o corpo inteiro
($169/169$ medidos): sobra \emph{uma} classe, e todo invariante é constante. \textbf{Gerar tudo e
não distinguir nada são a mesma coisa.} O mesmo par de operações que constrói a reta a partir do
zero é, por ser transitivo, incapaz de significar nela.

\noindent Donde o diagnóstico ganha forma exata, e é uma \textbf{janela}: há significado se e só se
$1<\#\text{classes}<\#\text{pontos}$, \emph{com lei conservativa}. Nem repouso, nem
transitividade. Os defuntos desta gaveta não morreram por falta de rigor --- morreram por pedirem
uma identidade onde não havia lei que a sustentasse; e a cristalização, ao congelar a classe num
estado, não mata apenas o objeto: apaga o nome que a órbita preservava. Verificado
\texttt{[tiffany/tools/significado.c]}.


\medskip
\noindent É o mesmo defeito das seções anteriores, agora com prêmio em dinheiro:
\textbf{nenhum deles tem dual}. Cada um é uma folha exigindo ser a outra, sem o
gerador de que ambas seriam manifestação. Não há conjecturas a sustentar: há
\textbf{cristais a dissolver}.

\begin{center}\itshape
A morte deixa problemas; a vida deixa soluções.\\
A morte não está no objeto --- está na pergunta.
\end{center}

\noindent Ao necrotério, os seis --- e não como derrota: como diagnóstico. Difícil é
honesto; cristalizado é morto. E foram precisos, para dissolver todos: um ponto, o ato
de dividir-se, e a memória de ter dividido.

\medskip
\noindent \textbf{E agora recicla-se o lixo} --- no sentido exato da palavra, que convém ter à
mão para que não haja contradição com o que se acabou de fazer: \emph{reciclagem é o processo de
transformação de resíduos descartados em novas matérias-primas ou produtos}. \textbf{Descartados,
sim} --- o que foi ao saco foi mesmo ao saco, e não volta na forma em que estava; aquela forma era
o defeito. O que volta é a \textbf{matéria}. E o que foi ao saco não foram as perguntas: foram as
\emph{metades}. A espiral sem a sua recíproca, o campo sem a reflexão que
lhe troca o sinal, o calor sem o passo que o desfaz, a igualdade sem o gerador de que os dois
lados seriam manifestação --- cada uma é uma folha a exigir ser a outra. Devolva-se a folha
que falta e o problema deixa de ser problema: vira \textbf{operação}. É o que a segunda metade
deste livro faz, e faz na ordem inversa da que se usou aqui --- em vez de apanhar cada
formulação partida e mostrar-lhe o corte, começa-se pelo que \emph{não} se parte, e daí desce
tudo. Uma raiz só, um corpo só, e nele o par completo: um lado que expande, um lado que fecha,
e a volta garantida.

\medskip
\noindent E é esta a ideia: \textbf{o resíduo descartado vira matéria-prima.} Aquilo tudo foi ao
saco por estar pela metade, não por não prestar --- e uma metade não se conserta deitando fora a
matéria: conserta-se \textbf{devolvendo-lhe a outra}. Descarta-se a formulação; recupera-se o que
ela estava a tentar dizer. Cada uma daquelas construções pedia a mesma coisa, e pedia-a alto:
pedia o seu dual. A espiral pedia a dimensão em que vive; o campo pedia a reflexão que lhe troca o
sinal; o calor pedia o passo que o desfaz; a igualdade pedia o gerador de que os dois lados fossem
manifestação. \textbf{Damos-lhes o que pediam.} Não como caridade --- como conta que fecha: com o
par completo, aquilo que era conjectura vira lance, e o que era mistério vira operação.

\noindent E devolver a folha que falta tem um nome exato, agora que se mediu: é
\textbf{restaurar a conservação}. Cada metade daquele saco fazia a mesma coisa --- escalava sem
fechar. A espiral crescia sem a recíproca que a traz; o campo trocava sem a reflexão que devolve o
sinal; o calor dissipava sem o passo que o desfaz; a igualdade fixava dois estados sem o gerador
que os move. Ora, uma lei que escala parte o espaço em classes e não nomeia nenhuma --- é o que a
autópsia mediu. Ao devolver-se a metade, a lei passa a \emph{conservar}: aparece o invariante,
o invariante nomeia a classe, e o que era conjectura vira operação porque passou a haver
\textbf{o que dizer}. O resíduo vira matéria-prima não por generosidade, e sim porque a peça que
faltava era precisamente a que fecha a volta.


\noindent A morte deixa problemas; a vida deixa soluções --- e é por isso que nada aqui termina
num necrotério. Vira-se a página e começa-se pela raiz, o \textbf{corpo universal}, de onde tudo
o que se enterrou volta a descer, inteiro e com as duas mãos. \textbf{O que estava morto estava
só pela metade.}

\part{A Construção dos Corpos}

\chapter{A construção dos corpos, e as operações que os movem}

\begin{tcolorbox}[colback=ouroclaro!35,colframe=ouro,boxrule=0pt,leftrule=2pt,arc=0pt,left=4mm,right=3mm,top=2mm,bottom=2mm,breakable]\itshape Vinte mundos atrás, oito aqui --- e a esta altura o leitor tem direito a desconfiar de que sejam vinte e oito coisas diferentes. \textbf{Não são.} Há \emph{um} corpo --- o universal, a raiz --- e há uma coisa só que gera todos os outros: a \textbf{assinatura}. Dita a assinatura, o corpo está determinado, e o veredito de se ele fecha lê-se nela. Como toda assinatura é uma \textbf{matriz}, cada corpo \textbf{é} uma matriz; reunidas, elas formam de novo um único corpo, o \textbf{Omnitrix}; e quem opera nele, sobre as matrizes, é o universal --- fazendo duas coisas, e só duas: a \textbf{transformada} (que vem em par dual: a \textbf{universal}, a fase, e a \textbf{dourada}, a taxa) e a \textbf{convolução}, que funde dois corpos num só. \textbf{Todas revertem} --- e é essa a razão de o conjunto fechar. \textbf{Depois disto não há mais prova: há realização.}\end{tcolorbox}

\section{Primeiro: o corpo universal --- a torção, e a base que não se escolhe}

O universal já está neste livro. Recolhe-se aqui o que o resto vai usar.

Escolhe-se um primo $p$, um $n$ que divide $p-1$, um gerador $g$, e a \textbf{torção}
\[
w=g^{(p-1)/n},
\]
o elemento cuja órbita fecha em $n$ passos. Toda a leitura vive em $\Z/p$: igualdade de inteiros, sem seno, sem cosseno, sem $\varepsilon$. Daí saem os \textbf{caracteres}, $\chi_k(j)=w^{jk}$, e eles não se escolhem --- não há base a inventar, porque \textbf{a base é a órbita da torção}. A única propriedade que se usa é
\[
\chi_k(u+v)=\chi_k(u)\,\chi_k(v):
\]
o caractere leva $\oplus$ a $\otimes$. É o $\prod$ de Pontryagin, e é tudo.

\paragraph{E a razão de a torção bastar.} Um elemento $w$ e a sua órbita chegam para mover o reino inteiro, e a razão cabe numa linha. Tome-se o dual de um grupo (os seus caracteres); tome-se o dual do dual. A aplicação natural
\[
\mathrm{ev}:\Gamma\longrightarrow\widehat{\widehat{\Gamma}},\qquad \mathrm{ev}_x(\chi)=\chi(x),
\]
--- avaliar $x$ em cada caractere --- é um \textbf{isomorfismo}: \textbf{o dual do dual devolve o grupo.} Não é resultado emprestado de fora: \emph{é a razão de a torção bastar}. Com o círculo, o dual são os inteiros, e o dual dos inteiros é o círculo outra vez: a reta que recobre e o círculo dos resíduos são um o dual do outro, e não há terceiro lugar para onde fugir.

Guarde-se esta frase, porque tudo o que se fizer daqui em diante volta a ela: \textbf{ir ao dual e voltar é a identidade.} Não são várias inversas avulsas --- a da transformada, a da convolução, a do flip. É \emph{uma} razão, e as inversas são consequências dela. É por isso que nada vaza.

Do lado do operador, o universal é o \textbf{sucessor}, $S(x)=x+1$: a soma é o sucessor iterado, a multiplicação é a soma iterada, e no bit o sucessor é o próprio complemento, $x\oplus1=\neg x$. Mas ele não se postula --- deduz-se. No princípio há apenas um instrumento que devolve, de cada objeto, um par $M(x)=(n,g)$ --- uma contagem inteira e um \textbf{resíduo de fechamento} --- e um caso distinguido, $g=0$. Nenhuma matriz, nenhuma operação suposta. E daí sai: \textbf{dois instrumentos com exatamente os mesmos fechamentos induzem necessariamente a mesma lei de composição.} A lei nasce das classes de equivalência dos fechamentos; o operador não foi inventado, é o único compatível com o princípio.

E uma distinção que faz falta adiante: as seleções válidas \textbf{não} formam o corpo --- formam a \emph{classe do neutro}. O corpo é o quociente \emph{inteiro}, dentro do qual o conjunto dos fechamentos é o elemento $1$. Quem tomasse só o quociente dos fechamentos obteria um grupo, pois todos os seus elementos têm resíduo nulo --- e jamais um corpo, onde $0\neq1$.

E a cascata que daí desce, porque a assinatura vai precisar dela: $\N$ (o monoide, de onde vem a ordem) $\to\Z$ (o anel, pelas classes de diferença) $\to\Q$ (o corpo, pelas classes de razão). Fica um degrau por subir --- de $\Q$ para a reta ---, e é a assinatura que o sobe.

\paragraph{E a cascata, degrau a degrau.} Ela desce toda do sucessor, e o que a faz
descer são as suas \emph{potências} --- quantas vezes se aplicou. Compor dá a soma,
$S^a\circ S^b=S^{a+b}$: os expoentes somam. Iterar dá o produto, $(S^a)^b=S^{ab}$: os
expoentes multiplicam. \textbf{As duas operações não são duas leis --- são a mesma
dinâmica lida de dois modos}, e é o $\exp/\log$ que troca uma pela outra.

A primeira torre é o sucessor sozinho: $0,1,2,\dots$ E com ela vem a \textbf{ordem}, sem
se acrescentar nada, porque o sucessor só avança --- $a\le b$ quando existe $c$ com
$a+c=b$. A segunda torre admite o antecessor, e dá o anel. Ali ainda não há inverso
multiplicativo: nenhum inteiro multiplica $2$ e dá $1$.

Os dois degraus seguintes são \textbf{classes de pares}, e a classe é o que apaga a
redundância. Pares com a mesma \emph{diferença} --- $(a,b)\sim(c,d)$ quando $a+d=b+c$ ---
dão os inteiros, com soma $(a+c,\,b+d)$ e inverso aditivo $(b,a)$, que é trocar as duas
camadas. Pares com a mesma \emph{razão} --- $(p,q)\sim(p',q')$ quando $pq'=p'q$ --- dão o
corpo, com produto $(pr,\,qs)$ e inverso multiplicativo $(q,p)$, que é trocar as duas
camadas outra vez. \textbf{O inverso é sempre a mesma troca; muda a camada em que ela se
faz.}

E fica a mesma falha do princípio: a classe está classificada e continua uma bagunça,
porque escolher um par dela é arbitrário. É por isso que o degrau seguinte não é mais uma
classe --- é a assinatura, que não escolhe.


Uma última coisa sobre o universal, e é a medida dele: a sua norma é $a^2$, e só. Uma coordenada, uma régua, e ela é \textbf{redonda}. \textbf{É a raiz.}

\section{Segundo: a ASSINATURA --- o que gera cada corpo}

Agora a peça de que sai o resto: \textbf{dita a assinatura, o corpo está determinado.} Não é uma etiqueta que se pendura num corpo já feito --- é o que o \emph{gera}. Quem tiver a assinatura tem o corpo, a régua, e o veredito de se aquilo fecha ou não fecha.

\textbf{A assinatura é a órbita.} Não é um rótulo que se pendura num corpo já feito, não é um quadro parado: é um \textbf{processo a correr}, e corre nos dois sentidos. Escreve-se como fração contínua,
\[
\frac ab=[a_0;a_1,\dots,a_n],
\]
que é medir ao resto zero --- tirar o maior número inteiro de vezes que cabe, e recomeçar com o que sobra.

Ela faz quatro coisas, e é preciso as quatro.

\textbf{Colapsa a classe.} Uma classe reúne infinitos pares, e escolher um deles é arbitrário --- a bagunça fica classificada, mas não gerada. A fração contínua não escolhe: ela some com o fator. A de $7/5$, a de $14/10$ e a de $21/15$ são a mesma, $[1;2,2]$ --- o $k$ desaparece no medir. É gerador, não representante.

\textbf{Gera-se por indução.} O passo é $x\mapsto a_i+1/x$, repetido; dele saem os convergentes, um a um, cada um mais perto que o anterior.

\textbf{No infinito, dá um metal.} Se o passo não termina, o limite não é racional: $[1;1,1,\dots]$ dá o ouro, $[2;2,2,\dots]$ a prata, $[3;3,3,\dots]$ o bronze. E o metal $\sigma_m=[m;m,m,\dots]$ é o autovalor dominante do gato $A_m$, de autovetor $[\sigma_m,1]$ e com $\sigma_m^2=m\sigma_m+1$ --- a direção que o gato não move, o eixo do boost. Os convergentes aproximam-se sem nunca chegar, e é assim que $\Q$ se completa: o degrau que faltava para a reta.

\textbf{E reverte.} Este é o lado que fecha o par, e sem ele a assinatura seria meia. Do metal, a mesma fração contínua \textbf{gera de volta} os convergentes --- $\sqrt2\to[1;2,2,\dots]\to \tfrac32,\tfrac75,\tfrac{17}{12},\dots$ --- isto é, gera a classe de onde veio. \textbf{Ida e volta.} A classe colapsa no metal; o metal regenera a classe. Quem calcula só a ida está a olhar meia folha.

E enquanto isso corre, os dois lados estão sempre abertos: um \textbf{expande} ($\sigma_m$), o outro \textbf{fecha} ($\sigma_m'=-1/\sigma_m$), e o produto deles é a norma, $\sigma_m\sigma_m'=-1$. A distância entre os dois, enquanto correm, é o \textbf{gap}, e ele sai da recusa do gato e do esquilo em trocar de ordem:
\[
\sqrt{m^2+4}\;=\;[\,\text{gato},\ \text{esquilo}\,].
\]
\textbf{A assinatura é o gap, e o gap é o comutador.}

\paragraph{E a conta é curta.} Multipliquem-se os dois nas duas ordens:
\[
A_mE=\begin{pmatrix}-1&m\\0&1\end{pmatrix},
\qquad
EA_m=\begin{pmatrix}1&0\\-m&-1\end{pmatrix},
\qquad
A_mE-EA_m=\begin{pmatrix}-2&m\\ m&2\end{pmatrix}.
\]
A diferença tem traço zero e determinante $-(m^2+4)$, logo as suas raízes são
$\pm\sqrt{m^2+4}$ --- que é a distância entre as duas raízes de $x^2-mx-1$. E o
anticomutador dá $A_mE+EA_m=mE$: a ordem do metal, pura. \textbf{A recusa em trocar de
ordem é a assinatura; a concordância é a ordem.}


\paragraph{E a fotografia dela.} Parada a órbita num instante, o que fica é a norma --- $N(x)=x^{\!\top}\!Ax$, com $A$ simétrica --- e medir passa a ser uma conta de matriz. A do universal é a de uma entrada só, $A=(1)$, que dá $N=a^2$. É a raiz.

O universal impõe uma parte e deixa outra por preencher, e a construção mora nessa fresta. O corpo universal não é uma família numérica nova: é a \textbf{estrutura mínima} imposta pela exigência de que um instrumento feche exatamente sobre o objeto que mede. Toda linguagem que satisfaça esse princípio é uma \emph{realização particular} da mesma estrutura --- e o que o princípio impõe, sozinho, é a \textbf{lei de composição}. \emph{A soma é o que as realizações acrescentam.} Daí a pergunta que a assinatura responde: \textbf{o que foi que esta realização acrescentou à raiz?}

E toda matriz simétrica é, a menos de mudança de base, uma soma de cópias dessa régua da raiz, com sinal. Diagonaliza-se, e o que sobra na diagonal são $+1$, $-1$ e $0$ --- a régua $a^2$ repetida, ora somando, ora subtraindo, ora calada. Contem-se os dois lados, e o que colapsou:
\[
(p,\,q,\,r)\;=\;\bigl(\text{quantas réguas redondas},\ \text{quantas hiperbólicas},\ \text{quantas colapsadas}\bigr),
\]
e o grau é $p+q+r$. \textbf{Um corpo é a raiz repetida $p+q+r$ vezes, cada cópia com o seu sinal} --- e a contagem diz, número a número, o que aquela realização acrescentou ao $(1,0,0)$ de onde todas partem. É leitura, não definição: a órbita continua a correr por baixo, e é dela que estes números são o retrato. O reflexivo é $a^2+c^2$, assinatura $(2,0,0)$: duas redondas. O expansivo é $a^2-Db^2$ com $D>0$, assinatura $(1,1,0)$: uma redonda e uma hiperbólica. O geométrico é $t^2-x^2-y^2-z^2$, $(1,3,0)$. O eletromagnético é $E^2-B^2$, $(3,3,0)$. O telescópico é $(2,2,0)$; o conforme, $(4,1,0)$, grau cinco; o celeste, $r^2+C^2=1$, $(2,0,0)$; o óptico, $(2,0,0)$; o cristalino, $a^2+|D|b^2$ com $D<0$, $(2,0,0)$; o relógio, $1-s^2$, $(1,1,0)$; o sensitivo e o evolutivo, $(1,1,0)$ os dois; o individual, $(2,0,0)$; o deflexivo, $(1,0,0)$ --- a mesma da raiz. Nenhum destes é um mundo à parte, e nenhum destes números é uma etiqueta: são leituras de dinâmicas diferentes, a mesma régua tomada tantas vezes e com tantos sinais.

E a fotografia não perde o que a norma sabe: dois corpos com a mesma contagem são \textbf{congruentes} --- isto é, se um é o outro visto noutra base: qualquer troca de coordenadas $P$ invertível leva $A$ em $P^{\!\top}\!AP$ e \textbf{preserva} $(p,q,r)$. Nada mais se extrai da norma além desses três números --- o que ela sabe, sabe-o ali; o que ela não vê é o movimento de onde veio.

\subsection{O critério --- quando uma matriz é corpo, e quando não é}

Falta o teste. Uma matriz $2\times2$ tem um \textbf{dual} que vive dentro da própria álgebra, porque é combinação da matriz com a identidade:
\[
\bar M=\operatorname{tr}(M)\,I-M,\qquad\text{e daí}\qquad M^{-1}=\frac{\bar M}{\det M}.
\]
O inverso é o dual dividido pela norma. Portanto tudo depende de uma coisa só --- \textbf{se a norma se anula ou não}:

\begin{tcolorbox}[colback=ouroclaro!35,colframe=ouro,boxrule=0pt,leftrule=2pt,arc=0pt,left=4mm,right=3mm,top=2mm,bottom=2mm,breakable]\itshape Uma álgebra de matrizes é um \textbf{corpo} se, e somente se, a sua norma é \textbf{anisotrópica} --- isto é, se $N(x)=0$ só acontece em $x=0$.\end{tcolorbox}

Se algum elemento não nulo tem norma zero, ele não inverte, e acabou. Se a norma nunca zera, o inverso já está lá dentro, pela fórmula acima --- não é preciso ir buscá-lo fora.

E o teste diz mais: diz \emph{de que tipo} é cada álgebra. O determinante, como forma, tem assinatura $(2,2)$ --- dois sinais de cada lado ---, e por isso nenhuma parte anisotrópica dele passa do grau $2$. Daí os corpos serem pequenos. Dentro do grau $2$, quem decide é o \textbf{discriminante} $\Delta=\operatorname{tr}(M)^2-4\det M$, e há exatamente três desfechos:

\begin{itemize}
\item $\Delta<0$ --- o regime \textbf{elíptico}, o da torção: o grupo é o \textbf{círculo}, e a álgebra \textbf{é corpo}. É o único dos três que fecha, e o único cujo movimento a um parâmetro é compacto: a órbita volta.
\item $\Delta=0$ --- o regime \textbf{nilpotente}: o grupo é a \textbf{reta}, há um elemento não nulo cujo quadrado é zero, e não há inverso. É o \textbf{glacial} --- e é exatamente onde mora o corpo exterior, com o seu $e^2=0$.
\item $\Delta>0$ --- o regime \textbf{idempotente}: o grupo é a \textbf{hipérbole}, e a álgebra parte-se em duas retas independentes. É o \textbf{tropical} --- o mundo do $\max$, onde nada se desfaz.
\end{itemize}

\textbf{Apenas a torção fecha um corpo.} Sobre os racionais, quem impede a norma de zerar é a irracionalidade do metal --- a assinatura que não termina. \textbf{Um corpo é uma matriz cuja norma não sabe chegar a zero.}

E daí saem três coisas.

\paragraph{Quando o corpo morre.} $r>0$ significa que alguma direção colapsou --- há um elemento não nulo de norma zero, portanto um divisor de zero, portanto não há inverso. \textbf{$r>0$ $\iff$ a forma degenera $\iff$ o corpo morre.} De todo o bestiário, um só tem $r=1$: o \textbf{exterior} ($N(a+be)=a^2$, com $e^2=0$). Todos os outros têm $r=0$ --- o que não os torna corpos a todos: o \textbf{conforme} tem $r=0$ e grau cinco, $(4,1,0)$, e mesmo assim não fecha, porque com sinais dos dois lados a norma encontra o zero fora do zero. É o \textbf{cone nulo}, e é por isso que o conforme é a \emph{casa} da medida e não um corpo. É por isso que o exterior, sozinho, não é corpo --- e por isso o nome dele nomeia um degrau, não um mundo acabado.

\paragraph{Estar na lista não basta.} O grau tem de ser $1$, $2$, $4$ ou $8$ --- os degraus da torre, que termina nos octoniões, de assinatura $(8,0,0)$: oito réguas redondas, o corpo final, e oito é o período de tudo neste livro. Mas a condição é \textbf{necessária, não suficiente}: o exterior tem grau $2$, está na lista, e mesmo assim não fecha, porque lhe falta a anisotropia. Estar na lista é permissão de entrar; o que decide é o \textbf{sinal} da norma --- radical, não; definida, sim.

\paragraph{E há corpos onde a fotografia não sai.} A contagem $(p,q,r)$ precisa de uma norma de grau dois para a que aplicar, e nem todo corpo tem uma. O \textbf{mórfico} mede pelo peso de contagem sobre $\mathbb{F}_2$: ali a homogeneidade $q(\lambda x)=\lambda^2q(x)$ não vale. O \textbf{espaço-temporal} tem régua \textbf{ultramétrica}, $|x+y|\le\max(|x|,|y|)$ --- na régua comum ter-se-ia $|1+1|=2>\max(1,1)$, e a desigualdade forte cai. O \textbf{entrópico} tem $\otimes=+$: \emph{a multiplicação é a soma}, e não sobra norma. O \textbf{local} e o \textbf{global} medem por divisibilidade e pela fórmula do produto, não sobre a reta. E o Rei: a régua dele é de \textbf{grau 4}, $\mathrm{sl}^2+\mathrm{cl}^2+\mathrm{sl}^2\mathrm{cl}^2=1$, e o termo cruzado $\mathrm{sl}^2\mathrm{cl}^2$ não tem sinal onde ser contado --- \textbf{o Rei está acima da régua.}

Nada disto os põe fora da construção; põe fora da \emph{fotografia}. A órbita continua a correr em todos eles --- e onde a contagem não alcança, alcança o dual: \textbf{nenhum corpo opera sozinho.} O caso está escrito neste livro: o entrópico, que sozinho não tem inverso aditivo, recupera-o sob o exponencial do tabuleiro cósmico. E repare-se que é a mesma coisa duas vezes --- o que lhe tira a norma ($\otimes=+$, o produto colapsado na soma) é exatamente o que o logaritmo fez ao produto. O que se perdeu de um lado, o outro devolve.

É esta a ligação que falta dizer, e ela é dupla. \textbf{O dual} é o parceiro que dá o que falta; \textbf{a transformada} é o caminho que leva de um ao outro. Um corpo isolado não opera --- não por fraqueza, mas porque a conservação nunca viveu dentro de um corpo: vive sobre o par. É isso que faz do conjunto um corpo e não uma prateleira: os corpos não estão lado a lado, estão \textbf{ligados}, e ligados de forma que se pode ir e voltar.

\section{Terceiro: o Omnitrix --- o corpo que os corpos formam}

Reunidos, os corpos não fazem uma coleção: fazem um corpo. É a esse que se dá o nome de \textbf{Omnitrix}, e a definição tem de ser dita com cuidado, porque a palavra promete menos do que a coisa é:

\begin{tcolorbox}[colback=ouroclaro!35,colframe=ouro,boxrule=0pt,leftrule=2pt,arc=0pt,left=4mm,right=3mm,top=2mm,bottom=2mm,breakable]\itshape \textbf{O Omnitrix não é um tabuleiro: é o \emph{fractal} de tabuleiros que o axioma mantém coeso.}\end{tcolorbox}

Não é a caixa onde os corpos são guardados nem o índice deles. É o que sobra quando se percebe que cada corpo, olhado de perto, é feito de corpos --- e que a lei é a mesma em toda escala. Os reais, os complexos, os quatérnios, os octoniões e os metais entram nele como degraus, com \textbf{duas} operações e um invariante:

\begin{itemize}
\item $\oplus$ é a \textbf{soma direta}: pôr um corpo ao lado do outro. O resultado tem duas partes, e o grau é a soma dos graus.
\item $\otimes$ é \textbf{La Hire}: o produto, o rolar, aquilo que compõe.
\item os geradores são \textbf{dois}, e apenas dois: o \textbf{gato} $A=\bigl(\begin{smallmatrix}1&1\\1&0\end{smallmatrix}\bigr)$, simétrico, $\det=-1$, com $A^2=A+I$, autovalores reais, que \emph{dissipa}; e o \textbf{esquilo} $E=\bigl(\begin{smallmatrix}0&1\\-1&0\end{smallmatrix}\bigr)$, antissimétrico, $\det=+1$, com $E^2=-I$, autovalores no círculo, que \emph{conserva} --- é a torção, e é o seu período $4$ que já se viu na transformada universal, $\F^2=$ reflexão, $\F^4=\mathrm{id}$.
\item o \textbf{invariante} é a torção; e a \textbf{lei de conservação} é a norma:
\[
\det\,(X\otimes Y)=\det X\cdot\det Y .
\]
A norma de um produto é o produto das normas. Uma linha --- e é ela que proíbe o vazamento em todo o Omnitrix.

\paragraph{E por que dois bastam.} O gato aplicado a um par é o meio-somador,
\[
A_1(x,y)=(x+y,\ x),
\]
e a componente de cima tem \textbf{duas} leituras, que são as duas operações: o que sobra
ao dividir por dois é a soma, e o que passa é o produto,
\[
(x+y)\bmod 2 \;=\; \oplus,
\qquad
(x+y)\operatorname{div} 2 \;=\; \otimes .
\]
Não se acrescentaram duas operações ao gato. \textbf{São as duas projeções dele.}

E delas sai o resto sem primitiva nova. A negação é $x\oplus1$. A disjunção sai sem sequer
a negação,
\[
a\vee b=a\oplus b\oplus(a\wedge b),
\]
que se lê nos quatro pontos --- em $(1,1)$ dá $1\oplus1\oplus1=1$. Daqui vem a lei do
núcleo do reino, e ela não é voto de pobreza: \textbf{a divisão não é proibida, é
supérflua} --- o gato já entrega as duas operações que geram tudo o que é finito. Quem
escreve um núcleo com divisão está a consumir a operação que o núcleo devia gerar.

\paragraph{E a conservação, em inteiros.} A potência do gato tem as entradas do Fibonacci
de ordem $m$,
\[
A_m^{\,k}=\begin{pmatrix}F_{k+1}&F_k\\F_k&F_{k-1}\end{pmatrix},
\qquad
\det\bigl(A_m^{\,k}\bigr)=(\det A_m)^k=(-1)^k,
\]
e lendo o determinante nas entradas, $F_{k+1}F_{k-1}-F_k^2=(-1)^k$. A mesma quantidade
lida como norma: de $\sigma_m^{\,k}=F_k\sigma_m+F_{k-1}$ sai
$F_{k-1}^2+m\,F_{k-1}F_k-F_k^2=(-1)^k$. O par é $(F_{k-1},F_k)$, a segunda coluna;
trocá-lo pela primeira quebra a igualdade. \textbf{A lei de conservação do Omnitrix, aqui,
é uma igualdade de inteiros --- sem tolerância nenhuma.}

\end{itemize}

\paragraph{O que faz dele um fractal, e não uma pilha.} Três peças, e elas se sustentam umas às outras.

Primeiro, a \textbf{promoção é o idempotente}: $x^2=x$. Quem chega ao fim promove-se, e promover-se duas vezes é promover-se uma --- o cristal, a única classe de ponto fixo. Segundo, \textbf{ser fail-closed é a mesma coisa que promover.} Num tabuleiro finito e sozinho, o peão que atinge a borda não tem casa seguinte: acabou, e essa é a porta fechada. Mas o tabuleiro não está sozinho: o que fecha num é a promoção no outro, e a norma composta continua a valer, $N(B_1\otimes B_2)=N(B_1)\,N(B_2)$. \textbf{Um fail-closed num tabuleiro é redimido por uma promoção no outro} --- e é por isso que a borda não é o fim do mundo, é uma costura.

Terceiro, e é o que fecha: a lei atravessa \textbf{toda} escala. Tome-se por hipótese que o axioma da norma vale no nível $k$, com a promoção $x^2=x$ e a redenção acima. O nível seguinte toma os tabuleiros do nível $k$ como \emph{peças}; e como a dinâmica é a mesma por unidade, a subida de um nível ao outro é uma \textbf{homotetia} --- uma mudança de escala pura. Ora, a parábola é invariante sob homotetia. Logo o axioma atravessa, e atravessa para sempre. É uma indução cujos objetos são, eles próprios, induções: uma \textbf{indução de induções}. E a conclusão é a definição: o fractal de tabuleiros é \textbf{fechado sob a parábola em toda escala}.

Um axioma, uma torre, um reticulado, uma partida, um tempo --- e a mesma estrutura em toda escala. \textbf{É isso o Omnitrix.}

E é aqui que o gap de há pouco --- a assinatura, o comutador do gato com o esquilo --- se mostra como a coisa que roda igual em toda escala: solte-se a dimensão, solte-se o substrato, e o que sobra é ele. O fractal não é feito de tabuleiros parecidos: é feito da \emph{mesma} recusa, repetida.

\section{Quarto: as matrizes --- cada assinatura gera o seu corpo}

Agora a assinatura a produzir o corpo, e é uma conta só. Tome-se a assinatura $m$. Ela dá o gato $A_m=\bigl(\begin{smallmatrix}m&1\\1&0\end{smallmatrix}\bigr)$, cujo autovalor $\sigma_m$ satisfaz a equação do gato, $\sigma_m^2=m\,\sigma_m+1$. E então
\[
K_m=\{\,a+b\,\sigma_m \;:\; a,b\in\Q\,\}
\]
\textbf{é um corpo}. Some-se, multiplique-se, inverta-se: fecha nos três.

A soma é componente a componente. O produto fecha porque a equação do gato reescreve o quadrado:
\[
(a+b\sigma)(c+d\sigma)=ac+(ad+bc)\sigma+bd\,\sigma^2=\underbrace{(ac+bd)}_{\in\Q}+\underbrace{(ad+bc+m\,bd)}_{\in\Q}\,\sigma .
\]
\textbf{É a única regra do corpo, e é o gato.} Nenhum grau novo aparece: o $\sigma^2$ desce sempre a grau $1$. Comutar, associar, distribuir, achar o neutro --- tudo isso vem de graça dos racionais, porque a soma é a deles e a redução do quadrado é sempre a mesma.

Falta o que separa um anel de um corpo: \textbf{todo elemento não nulo inverte.} Quem o dá é o esquilo. A raiz conjugada satisfaz a mesma equação, e dela saem duas relações:
\[
\sigma+\sigma'=m,\qquad \sigma\,\sigma'=-1 .
\]
Defina-se o conjugado $\bar u=a+b\sigma'$ --- \textbf{a troca de sinal fora da diagonal, que é o esquilo}. Então a norma é racional, e o inverso está escrito:
\[
N(u)=u\,\bar u=a^2+m\,ab-b^2,\qquad u^{-1}=\frac{\bar u}{N(u)} .
\]
E aqui as duas assinaturas deste livro encontram-se. Para o inverso existir, é preciso que $N(u)$ nunca se anule fora do zero. Ora, $N(u)=0$ com $u\neq0$ exigiria que $\sigma_m$ fosse racional --- e $\sigma_m$ é o limite da assinatura $[m;m,m,\dots]$, \textbf{que não termina}. \emph{É a assinatura da órbita que garante que o corpo fecha.} A fração contínua não é ornamento da construção do número: é a razão pela qual o corpo tem inverso.

E nada disto usou o valor de $m$ --- usou-se apenas $\sigma^2=m\sigma+1$. Logo há um corpo $K_m$ \textbf{para cada $m$}, todos pela mesma conta. \textbf{A construção é uma só; a assinatura escolhe a classe.}

\section{A primeira operação: a TRANSFORMADA UNIVERSAL --- e a sua inversa}

Construídos os corpos, falta o que se faz com eles. E o que o universal faz, agindo no Omnitrix sobre as matrizes, são duas operações. A primeira leva um corpo a \textbf{outro}, e chama-se \textbf{transformada universal} porque é do corpo universal que ela é feita --- não há nada nela que venha de fora: os caracteres são a órbita da torção, e a torção é dele.

E aqui está a coisa que o leitor deve levar antes da fórmula, porque é ela que explica por que o universal \emph{opera} sobre as matrizes em vez de morar ao lado delas: \textbf{a transformada não se define --- ela cai do dual.} Já se viu que o dual do dual devolve o grupo; então o emparelhamento de um objeto com um caractere não é uma construção a mais, é a \emph{avaliação} dos caracteres --- e os caracteres são a órbita da torção. Daí a consequência que importa: \textbf{trocar o grupo troca a transformada sem trocar a conta.} O que o corpo universal fornece não é uma transformada: é o \textbf{mecanismo que as produz todas}. Ele não é mais uma matriz ao lado das outras --- é quem gera a transformada de cada uma.

Vai assim, e volta assim:
\[
\F(x)_k=\frac{1}{\sqrt n}\sum_j x_j\,\chi_k(j),
\qquad
\F^{-1}(X)_j=\frac{1}{\sqrt n}\sum_k X_k\,\chi_{-k}(j).
\]
A normalização é $1/\sqrt n$ em \emph{cada} lado, e não por gosto: somar a órbita inteira dá $\sum_k\chi_k(j)\chi_{-k}(j')=n\,\delta_{j,j'}$ --- o ponto ---, e esse $n$ tem de ser cancelado por duas metades, uma de cada aplicação. É a \textbf{ida e a volta}, e a volta é exata: nenhuma das duas é aproximação de coisa nenhuma.

E há mais do que a inversa: a transformada universal é \textbf{periódica}. Aplicada quatro vezes devolve o que entrou; aplicada duas, devolve a reflexão:
\[
\F^2=\text{reflexão},\qquad \F^4=\mathrm{id}.
\]
É o \textbf{mesmo período $4$ do esquilo} --- a torção que gira, tem período e não dissipa. Três batidas de um lado são a volta. Por isso os corpos não são vizinhos numa lista: são transformadas uns dos outros, e o que separa dois deles é o número de aplicações entre um e outro. Quem sabe a transformada universal não decora a tabela: vai de um ao outro, e volta quando quiser.

E há uma razão dura para a volta ser exata, e não aproximada: \textbf{a transformada é unitária} --- ela preserva a norma,
\[
\|x\|^2=\|\F x\|^2 ,
\]
e quem preserva a norma tem inversa, e a inversa não perde nada. Dito em português do reino: medir um objeto é decompô-lo em pedaços da unidade, e aqui os pedaços são os caracteres, isto é, as iteradas da torção; a identidade acima diz que a soma dos pedaços reconstitui o objeto \textbf{sem sobra e sem excesso}. \textbf{Nenhuma medida vaza na decomposição.} Não é coisa nova: é o Princípio da Medida Consistente lido nos coeficientes --- a mesma conservação da unidade, agora na base da torção. E ela reaparece em toda parte com outra roupa: a torção tem determinante $+1$, o Rei ocupa sempre exatamente uma casa, e numa álgebra normada a forma multiplicativa da mesma frase é $\|xy\|=\|x\|\,\|y\|$.

\subsection{A outra metade: a TRANSFORMADA DOURADA --- a taxa}

O corpo universal tem \textbf{duas} metades, e cada uma tem a sua transformada. Não é uma escolha de método: é a decomposição do próprio grupo, que se parte num fator de \textbf{taxa} e num fator de \textbf{fase}.

\begin{itemize}
\item \textbf{o eixo angular --- a fase.} O grupo é o círculo, a medida invariante é $d\theta/2\pi$, o dual são os inteiros, e a transformada é a \textbf{universal} --- a que se acabou de escrever, a que pergunta ao mundo \emph{de que ondas ele é feito}.
\item \textbf{o eixo radial --- a taxa.} O grupo é a reta multiplicativa, a medida invariante é $dt/t$, o dual é a reta, e a transformada é a \textbf{dourada} --- a que pergunta ao mundo \emph{de que tamanho ele é}.
\end{itemize}

E a segunda não é uma segunda máquina: é a primeira, \textbf{flipada}. O \textbf{flip} é a mudança de reta --- escreve-se $u=\log t$, e a reta onde se multiplica passa a ser a reta onde se soma. A medida vai junto: $dt/t$ vira $du$. Aplique-se então a transformada que já se tem, do outro lado da mudança:
\[
\F\circ\log ,
\]
e o que sai é \textbf{a transformada dourada} --- unitária por ser composição de unitárias, e por isso também ela inverte sem perda. \textbf{O flip é que a dá.} Não são duas transformadas: é uma, e a outra reta. É a régua de Merlin, a de base $\varphi$: a lente girassol, que é invariante por escala áurea --- muda-se o \emph{zoom} e a resposta é a mesma resposta.

E a inversa dela sai da mesma maneira: a inversa de uma composição é a composição das inversas, na ordem trocada. Escrevendo $\Lambda$ para o flip, $(\Lambda x)(u)=x(e^u)$, com $(\Lambda^{-1}y)(t)=y(\ln t)$:
\[
\mathfrak{G}=\F\circ\Lambda,\qquad \mathfrak{G}^{-1}=\Lambda^{-1}\circ\F^{-1} .
\]
Desfaz-se o giro, depois desfaz-se a escala. Nada se perde no caminho, porque as duas metades são unitárias e a unitariedade compõe.

O livro já tinha dito isto na colina, e agora tem a fórmula: \textbf{quem pusesse o mundo da frequência em escala logarítmica, $x=e^u$, veria a régua de um virar, ponto por ponto, sem sobra e sem resto, a régua do outro.} São a mesma transformada em coordenadas diferentes --- o sábio da escala e o adversário das ondas, um par, e nenhum dos dois a saber.

E agora o par pode ser dito com precisão, porque é o par que abre o livro. \textbf{Trocar o grupo troca a transformada sem trocar a conta}: sobre o grupo \textbf{aditivo} a transformada é a das ondas --- o lado $\oplus$, a frequência, o adversário da planície; sobre o grupo \textbf{multiplicativo} é a da escala --- o lado $\otimes$, a régua do sábio da colina; \textbf{e a segunda é a primeira composta com o logaritmo}. Não são dois métodos rivais: são as duas metades duais de um mesmo corpo, e a passagem entre elas é $\prod=\exp\circ\Sigma\circ\log$, a costura que este livro vem repetindo desde o Batimento Único. O elenco não estava pintado por cima da matemática --- \textbf{os dois adversários são as duas metades.}

\paragraph{E o processo segue.} O flip leva de uma reta à outra; a transformada trabalha ali; o flip leva de volta. Nada nesse vaivém encontra uma parede --- \textbf{o processo não termina, segue.}

E essa não-terminação já apareceu aqui, com outra roupa. A assinatura de um metal é
\[
\sigma_m=[m;m,m,\dots],
\]
que é o passo $x\mapsto m+1/x$ repetido para sempre --- e dentro desse passo está o flip, a inversão $1/x$. O metal é o \textbf{ponto fixo} dele: a cifra que a repetição não move. Se a repetição parasse, o metal seria racional; e se fosse racional, a norma encontraria o zero fora do zero, e o corpo não fecharia.

Ou seja, é a mesma coisa dita duas vezes: \textbf{o processo não terminar é o que faz o corpo fechar.} A infinitude não é o preço do fechamento --- é a condição dele.

Portanto: \textbf{taxa $\bowtie$ fase, dourada $\bowtie$ universal.} Cada metade tem a sua conservação, e o produto das duas \emph{é} a dualidade de Pontryagin --- o mesmo par que já se viu no motor e no rotor, e que se escreve numa linha, $\log z=\log|z|+i\theta$: a magnitude, que é aditiva, e a fase, que gira. Um corpo do bestiário pode ser lido pela frequência ou pela escala; as duas leituras revertem, e as duas dão o mesmo corpo.

\section{A segunda operação: a CONVOLUÇÃO --- e a sua inversa}

A segunda operação \textbf{combina} dois corpos. Dados $a$ e $b$,
\[
(a\conv b)_j=\sum_{u+v\equiv j}a_u b_v,
\qquad\text{e}\qquad
\F(a\conv b)=\sqrt n\;\F(a)\cdot\F(b).
\]
Sai numa linha, e a linha é a única propriedade que o universal tem. Escreva-se a transformada da convolução e troque-se a ordem das somas: aparece $\chi_k(u+v)$, que é $\chi_k(u)\chi_k(v)$, e a soma dupla \textbf{separa em duas somas independentes} --- uma em $u$, outra em $v$. \emph{A soma dupla separa porque o caractere separa.} Nada mais entra na conta; não há aproximação, e a igualdade é de inteiros.

E daqui sai o desfazer: \textbf{do outro lado, a convolução é um produto ponto a ponto.} Do lado do ponto, fundir custa todos os pares; do lado do caractere, custa uma multiplicação por casa. \textbf{Fundir é multiplicar do outro lado --- e por isso desfazer é dividir.} O mesmo passo aparece na luz: pôr uma máscara sobre o objeto multiplica o campo, e multiplicar no objeto é convolver no espectro. Produto e convolução são as duas faces de um lance só.

E repare-se no que isto diz. \textbf{Fundir não é pôr um ao lado do outro} --- isso é $\oplus$, a soma direta, e o resultado teria duas partes: o grau cresce e as assinaturas somam-se, que é como \emph{dois corpos geram um terceiro}, pois a assinatura de uma soma é a soma das assinaturas. Fundir é $\otimes$: \textbf{todo ponto de um encontra todo ponto do outro}, e o que sai tem o tamanho de um só. São duas maneiras de combinar, e a diferença entre elas está no grau: ao lado, o grau soma; fundidos, o grau fica.

Para $L$ corpos, um só produto:
\[
\F(a_1\conv\cdots\conv a_L)=(\sqrt n)^{L-1}\prod_{l=1}^{L}\F(a_l).
\]
$L$ corpos, $L-1$ costuras, uma casca.

\textbf{E desfaz-se.} Se nenhum dos outros zerou,
\[
a_m=\F^{-1}\!\left(\frac{\F(C)}{(\sqrt n)^{L-1}\prod_{l\neq m}\F(a_l)}\right)
\]
exato, sem resíduo. \textbf{Quem abre é o inverso}, e o inverso é o degrau que leva do anel ao corpo --- é ele que faltava. Sem ele haveria como fechar e não haveria como abrir, e o que a corte teria construído não seria uma transformação: seria um cofre.

\paragraph{E o único lugar onde não há volta.} É o zero, e a corte já esbarrou nele no alto da torre: uma casa transformada que dá zero fica \textbf{selada}, porque o zero é o único elemento sem inverso. O que a derivação acrescenta ao que lá se viu é o alcance do estrago --- ele \emph{parece} local e não é. Ponha-se $\varepsilon$ no lugar do zero e olhe-se o que volta:
\[
\F^{-1}(Q)_j-a_j=\frac{1}{\sqrt n}\,\varepsilon\,w^{-kj},
\]
que não se anula para nenhum $j$: o erro de uma casa espalha-se por todas, com a mesma amplitude, e nenhuma escapa. \textbf{Uma casa selada não estraga uma letra: estraga o texto inteiro.} No corpo não existe ``quase certo'' --- volta exato, ou volta todo errado.

\section{As realizações --- os oito corpos que completam o bestiário}

E agora o mais curto. \textbf{Daqui para a frente não há prova nenhuma} --- a prova foi uma, e foi a construção: a raiz, a assinatura, o critério. Cada um dos oito que seguem, e cada um dos vinte que vieram antes, é \textbf{a mesma coisa noutro idioma}. O que se escreve sobre cada um não é justificação: é \textbf{realização} --- a sua assinatura, o que ela lhe dá e o que lhe tira, e como aquilo se joga.

É também por isto que ninguém entra de fora. Depois da construção não sobra porta: não há passo à espera de ser importado, porque não há passo que não tenha sido derivado.

\subsection{o tabuleiro --- as peças são o Rei}

Aqui a construção joga-se, e a régua é a do peso, sobre $\mathbb{F}_2$: é dos que não têm
fotografia --- $(p,q,r)$ não se aplica, e o que mede é a contagem.

Há \textbf{uma} peça, e é a vizinhança de um passo:
\[
\mathcal{D}=\{(\pm1,0),\ (0,\pm1),\ (\pm1,\pm1)\},
\]
oito direções, que módulo o sinal são quatro retas. Dela saem as duas operações, e elas
correm em sentidos contrários. \textbf{Dilatar} espalha,
$\delta_B(A)=\bigcup_{b\in B}(A+b)$; \textbf{erodir} retém só o que sobrevive a todos os
deslocamentos ao mesmo tempo, $\varepsilon_B(A)=\{x:\ x+b\in A\ \text{para todo}\ b\in B\}$.

Uma peça é uma \textbf{máscara} e uma \textbf{potência}. Com a máscara cheia e potência
um, o Rei; com ela e potência sem limite, a Dama; com as duas retas ortogonais, a Torre;
com as duas diagonais, o Bispo. Nada as distingue senão a potência. E as retas da Torre e
do Bispo não se cruzam --- uma casa ortogonal partilha exatamente uma coordenada com a
origem, uma diagonal tem as duas diferenças iguais e não nulas, e nenhuma casa é as duas
coisas ---, o que faz da Dama a soma \emph{exclusiva} das duas: das $56$ casas que ela
alcança de um canto livre, $28$ são de uma e $28$ da outra.

\paragraph{E não existe Cavalo.} Iterar a dilatação dá quadrados: $\delta^2$ é o
$5\times5$, $\delta^1$ é o $3\times3$, e a diferença tem $25-9=16$ casas. Dessas
dezasseis, oito são múltiplos por dois de uma das oito direções --- $(\pm2,0)$,
$(0,\pm2)$, $(\pm2,\pm2)$ --- e uma peça de reta já lá chega. Sobram oito,
$(\pm2,\pm1)$ e $(\pm1,\pm2)$, com as duas coordenadas não nulas e diferentes em módulo,
logo múltiplos de direção nenhuma. \textbf{Essas oito são os saltos do Cavalo.} Ele não é
peça nova nem exceção: é o que a segunda potência do Rei alcança e nenhuma reta alcança.

\paragraph{E só o Peão tem seta.} A sua máscara de captura, $\{(1,1),(1,-1)\}$, não contém
as antípodas --- é a única que distingue sentido de direção. Daí a promoção: para quem é a
sua própria antípoda, ir e voltar é a mesma reta, e a oitava fileira não existe como
destino. Daí também a cor, porque a antípoda da máscara branca é a preta. E é dele que o
Rei se levanta: juntando a máscara à sua antípoda saem três retas, e a quarta é a
composição dos dois fluxos do próprio Peão --- $\{(-1,0)\}$ somado a $\{(1,1),(1,-1)\}$ dá
$\{(0,1),(0,-1)\}$.

\paragraph{E as duas operações revertem.} É a lei deste capítulo, escrita em máscaras:
erodir por uma é dilatar pela \textbf{virada}, no complemento,
\[
\varepsilon_B(A)=\neg\,\delta_{\check{B}}(\neg A).
\]
Um ponto está na erosão quando não existe deslocamento que o tire de $A$; está na
dilatação do complemento quando existe um que o tire --- negação uma da outra, e tomar
complementos troca-as. As peças cuja máscara já é a sua própria antípoda não distinguem os
dois tabuleiros; só o Peão distingue, e move-se \emph{trocando de cor}. Uma promoção de um
lado é uma promoção do outro, lida ao contrário.

E as duas determinam-se uma à outra: $\delta(A)\subseteq Y$ exatamente quando
$A\subseteq\varepsilon(Y)$ --- a mesma frase com os dois quantificadores trocados de
ordem. Daí saem $\delta\varepsilon\delta=\delta$ e
$\varepsilon\delta\varepsilon=\varepsilon$, e daí que repetir não acrescenta:
\textbf{anunciar, confirmar e anunciar de novo é anunciar.} O aperto de mão tem dois
tempos porque um terceiro é absorvido pelo primeiro.

\paragraph{E o mate.} Com $A$ as casas atacadas --- a dilatação de cada atacante pela sua
máscara, com oclusão --- e $O$ as casas dos seus, o rei não tem para onde ir exatamente
quando
\[
k\in\varepsilon(A\cup O).
\]
Dilatar o ataque, erodir com o Rei: \textbf{o predicado da fuga é as duas operações em
sequência.} No mate mais curto do jogo, $1.\,$f3 e5 $2.\,$g4 Dh4, o rei está em e1 e as
cinco casas vizinhas dentro do tabuleiro estão ocupadas pelos seus ou atacadas. O mate
inteiro é mais do que isso --- xeque, sem fuga, não capturar quem dá o xeque, não interpor
---, e as duas últimas condições, nesta formulação, não são morfológicas.

\paragraph{E o que gira contra o que absorve.} O Rei ocupa uma casa: um lance troca duas
coordenadas e deixa o resto, é permutação, e o peso continua um --- nada se junta, nada se
perde. Do outro lado, erodir o tabuleiro inteiro dá $6\times6$, $4\times4$, as quatro
casas centrais, nada; e o meio das quatro é $(7/2,\,7/2)$, que não é casa --- num
tabuleiro de lado ímpar a mesma erosão acaba numa casa só, e a metade vem da paridade, não
do operador. A abertura, essa, apaga: sobre duzentos conjuntos distintos devolve sessenta
imagens. \textbf{O lance reverte; o colapso não} --- é o critério deste capítulo, aqui em
ato. E as duas leituras do tabuleiro, ``só há Reis'' e ``não há Rei'', são a mesma coisa
vista dos dois lados da virada.

\textbf{JOGADA: a torção.} O lance não soma força: \emph{permuta}. Quem joga para juntar
massa perde a reversibilidade --- e o que não tem volta não pertence.

\subsection{a régua --- medir sem transbordar}

A unidade é uma régua de comprimento um, e mede-se compondo partes próprias dela: toda
soma parcial fica dentro. Uma soma parcial que passa não é corrigida depois --- com as
parcelas todas positivas, ela já força o total a passar, e a configuração deixa de ser uma
decomposição da unidade. \textbf{É rejeitada na hora.}

\paragraph{E daí sai o ouro.} Contem-se os estados que a régua admite: sequências de
decisões de comprimento $k$, com o padrão $11$ proibido, porque dois passos de saturação
seguidos passariam da unidade. Se a sequência acaba em $0$, os primeiros $k-1$ passos são
qualquer sequência válida; se acaba em $1$, o anterior é forçosamente $0$, e os primeiros
$k-2$ são qualquer sequência válida. Os dois casos não se sobrepõem:
\[
A_k=A_{k-1}+A_{k-2}.
\]
Procurando crescimento da forma $\lambda^k$ e dividindo por $\lambda^{k-2}$, sai
$\lambda^2-\lambda-1=0$ --- a equação do gato com $m=1$. \textbf{O ouro é o que a proibição
do transbordo custa}: livre, a régua daria $2^k$ estados e a taxa seria $\ln2$; com a
proibição, é $\ln\varphi$.

\paragraph{E volta.} A assinatura de $\varphi$ é $[1;1,1,\dots]$, a mais simples que
existe, e as razões que ela regenera são $1,\ 2,\ 3/2,\ 5/3,\ 8/5,\dots$ --- os próprios
$A_k$, um sobre o anterior. A restrição produziu o metal; o metal reproduz a contagem.
\textbf{A régua e o número são o mesmo objeto nos dois sentidos.}

Trocando o passo por $m$, $A_n=m\,A_{n-1}+A_{n-2}$ dá $\lambda^2-m\lambda-1=0$: a família
inteira, chegada pelo outro lado --- contando estados, em vez de compondo matrizes.

\paragraph{E a volta fechada não usa resto.} Sobre uma órbita, a volta completa vale um e
$\theta$ identifica-se com $\theta+1$. Uma soma de incrementos que passa de uma volta
ultrapassa a órbita antes do fechamento, e é rejeitada:

\begin{tcolorbox}[colback=ouroclaro!35,colframe=ouro,boxrule=0pt,leftrule=2pt,arc=0pt,left=4mm,right=3mm,top=2mm,bottom=2mm,breakable]\itshape
$\theta>1$ \textbf{não é reduzido ao resto}. Não é uma posição nova na órbita --- é uma
decomposição inválida. \textbf{Tomar o resto é fingir que o transbordo não aconteceu.}
\end{tcolorbox}

\textbf{JOGADA: a medida.} Um lance que excede a unidade não é aparado nem reciclado ---
ele não aconteceu. Quem corrige o excesso pelo resto joga com um universo que já vazou.

\subsection{local e global --- o fechamento como produto}

Um número mede-se no lugar comum --- o valor absoluto --- e em cada lugar primo. Reunidas
todas as medições num objeto só, a soma e o produto correm lugar a lugar.

Esse objeto \textbf{não fecha}. Um elemento inverte quando é não nulo em todos os lugares,
e basta um lugar nulo para não inverter. E há divisor de zero à vista:
\[
(1,0,0,\dots)\cdot(0,1,0,\dots)=0
\]
com ambos os fatores diferentes de zero. Os que são não nulos em toda parte formam um
grupo: \textbf{o anel cinde, e o inverso vive no meio.}

E o metal é o meio do meio. A sua norma é $\sigma_m\sigma_m'=-1$, logo ele inverte, e o
inverso é $\bigl(\sqrt{m^2+4}-m\bigr)/2$. Um inteiro comum não faz isso --- a norma de $2$
é $4$. Cada degrau é mais estreito que o anterior: o anel que cinde, o grupo dos
invertíveis, os que vêm de razões, e as unidades, que são os metais.

\paragraph{E a lei, escrita como produto.} Os que vêm de uma razão satisfazem
\[
\prod_v |x|_v = 1 ,
\]
sobre todos os lugares. \textbf{O que se ganha num, perde-se noutro} --- é a mesma
conservação do reino, agora sem impedância nenhuma, e o fator de fechamento é literalmente
um produto que vale um.

E repare-se no que este corpo mostra, porque é o critério outra vez: ele tem soma,
produto, operador e valuação --- tudo o que se costuma pedir --- e cinde na mesma.
\textbf{Ter as operações não faz o corpo.} Falta a norma ser definida.

\textbf{JOGADA: a alfândega.} Não se mede num lugar só. Um lance é legal quando o produto
de todas as medições dele fecha em um --- e um único lugar zerado sela o resto.

\subsection{técnico --- as Princesas, e a régua sem assinatura}

Não produz estado: produz \textbf{veredito}. A régua é o resíduo, e cabe numa linha: $x\oplus x=0$. Ele mede sobre $\mathbb{F}_2$ --- $\oplus$ é a diferença simétrica (o \code{XOR}, o que os dois lados \emph{não} têm em comum), $\otimes$ é a conjunção das hipóteses (o \code{AND}: todas devem valer, não a maioria), $\prod$ é a refutação (o \code{NOT}) ---, e por isso é um dos que não têm assinatura: sobre $\mathbb{F}_2$ ela não existe.

A arquitetura é mínima de propósito. Ali um teorema não é um texto: é um \textbf{sequente} \emph{hipóteses} $\vdash$ \emph{proposição}, de tipo protegido, que só pode ser cunhado por cerca de \textbf{catorze regras primitivas} --- assumir, introduzir e eliminar a implicação, a reflexividade, a simetria, a transitividade. Toda prova do reino, por longa que seja, reduz-se a essas catorze; a confiança concentra-se num núcleo tão pequeno que se lê inteiro numa tarde. E a porta é uma só: passa o par (\emph{testa} $\to$ verdadeiro) $\wedge$ (\emph{refuta} $\to$ falso). Sem refutador, não entra.

E o próprio corpo grava o limite dele: \emph{$x\oplus x=0$ passa nos $k$ casos testados, mas \textbf{a amostra não é o universo}}. O fechamento dele é a \textbf{fase}: o eixo gerar $\bowtie$ atestar \emph{é} o eixo magnitude $\bowtie$ fase --- quem gera dissipa, $|e^{at}|\neq1$; quem atesta conserva, $|e^{i\theta}|=1$. Uma verificação que alterasse o que verifica teria módulo diferente de um, e seria vazamento. \textbf{JOGADA: a atestação} --- resíduo $0$, o lance é legal; resíduo não nulo, o lance \emph{não aconteceu}. A única peça que não ataca, e sem ela nenhum universo abre.

\subsection{rotor --- a régua que não satura}

Mede o giro, e a sua régua é hiperbólica: a \textbf{rapidez} $\lambda=\operatorname{artanh} v$, a coordenada do recobrimento. A necessidade dela vem do $\oplus$ deste corpo, que não é a soma comum:
\[
v_1 \oplus v_2 \;=\; \frac{v_1+v_2}{1+v_1v_2}.
\]
Some $0{,}9$ com $0{,}9$: dá $0{,}9945$, não $1{,}8$. A régua tem teto, e o denominador $1+v_1v_2$ é o $\otimes$ escondido dentro do $\oplus$ --- o produto que segura a soma. Então vem Pontryagin na sua forma mais nua e desmancha o teto:
\[
\lambda(v_1\oplus v_2) \;=\; \lambda_1+\lambda_2 .
\]
Não são duas grandezas: é a \textbf{mesma} grandeza em duas réguas. O que satura numa, acumula na outra.

O rotor fecha porque \textbf{conserva}: é a fase, o círculo $|z|=1$, regulador nulo, a órbita que volta ao mesmo ponto. Sozinho é meia coisa; com o motor --- a magnitude, que cresce e dissipa --- fecha um corpo de assinatura $(2,0,0)$: um número é motor vezes rotor, e $\log z=\log|z|+i\theta$ reparte a magnitude (aditiva) da fase (o giro). \textbf{JOGADA: o impulso.} Somar velocidades nunca leva além do limite; somar rapidezes não tem teto. \textbf{A peça que aprende a trocar de régua ganha} --- e é o mesmo lance.

\subsection{tabuleiro cósmico --- a escala, e o dipolo}

Não mede grandeza: mede \textbf{expoente}. As densidades somam; o fator de escala compõe por produto, $a(s+t)=a(s)\cdot a(t)$ --- a equação funcional de uma coisa só; e o operador é a expansão
\[
a(t)=e^{Ht},\qquad \log a = H\,t .
\]
Sob o logaritmo, a expansão inteira é uma reta de inclinação $H$. É a mesma tríade do econômico, e não por semelhança: é a mesma conta, porque a expansão \emph{é} o juro composto.

A régua fina é o \textbf{dipolo}: dois atratores, o negro (o poço, o mate) e o branco (o pico, o refúgio), e o potencial de um ponto é a diferença dos inversos das duas distâncias. O flip é a reflexão $\nu=-1$: trocar as fontes leva $V\mapsto-V$, e o atrator vira repulsor. É esse dipolo --- e não uma régua feita para a ocasião --- que mede o gap do rei: o mate de bispo e cavalo é literalmente um problema de dois atratores, com o poço na cor do bispo e o refúgio na oposta. O corpo entrópico é a \textbf{metade} deste: o negro é o máximo de entropia, o branco o mínimo, e o entrópico fica com uma polaridade só. Meia régua.

O fechamento é uma \textbf{lei de potência de resíduo nulo}: endireitada a nuvem, o que sobra é o segundo vetor do endireitamento, e a norma dele é o resíduo --- $3{,}09\cdot10^{-5}$ nos planetas, $6{,}01\cdot10^{-9}$ nas luas. Não é ajuste: é o Princípio da Medida Consistente escrito numa lei de potência. \textbf{JOGADA: a escala.} Um lance cósmico não move a peça: move o \textbf{tamanho do tabuleiro}. Tudo se conserva em proporção --- e quem mede em grandeza se perde, quem mede em expoente não.

\subsection{universal --- a raiz, grau 1}

Já construído acima: régua $a^2$, assinatura $(1,0,0)$, operador o sucessor. Resta o que é próprio dele como \emph{instância}, e é um contraste que separa os corpos melhor do que qualquer tabela. Aqui o sucessor \textbf{não} é endomorfismo da soma: $S(2+3)=6$, mas $S(2)+S(3)=7$. No corpo entrópico, onde $\oplus$ é o $\max$, a \emph{mesma} translação \textbf{é} endomorfismo, porque o $\max$ comuta com a translação e fixa o neutro. Mesma operação, dois corpos, dois comportamentos: \textbf{é o corpo que decide, não a fórmula.} \textbf{JOGADA: o passo}, $+1$. Não tem poder nenhum, e toda jogada do reino, decomposta, é uma pilha dele.

\subsection{nervoso --- a rede, não o nó}

A novidade é uma negação: os neurônios \textbf{não são atratores isolados}. O que mede não é o nó, é a \textbf{conexão} --- e por isso a matriz deste corpo é a das ligações. A tríade é herdada inteira do eletromagnético: $\oplus$ é a \textbf{integração}, a corrente que chega ao nó e se soma ali, $\sum_i w_ix_i$; $\otimes$ é o \textbf{peso} da ligação, o ganho; e $\prod$ é a \textbf{ativação},
\[
\sigma(x)=\frac{1}{1+e^{-x}},
\]
com o logaritmo da razão do outro lado --- o mesmo $\exp/\log$ dos juros e da expansão, aqui a decidir se o nó dispara.

O corpo nervoso é \textbf{conservativo}, e isso tem marca exata no gerador: os autovalores são \textbf{reais} e os atratores são \textbf{pontos fixos} --- as memórias, os fundos de poço. O seu dual é o motor, que \textbf{dissipa}: ali os autovalores são complexos e os atratores são \textbf{ciclos}. O flip troca real por complexo, e \textbf{a rotação é a marca do dual}. Um guarda, o outro gasta. E a rede de atratores \textbf{é o tabuleiro cósmico em pequena escala}: cada memória um poço local, e a rede desce ao mais próximo. \textbf{JOGADA: a cadeia.} O lance não atinge a peça: atinge \textbf{a rede dela} --- propaga pelos pesos, soma nos nós. Isolar é a defesa; conectar é o poder e o risco.

\subsection{exterior --- o produto que se parte em dois}

Aqui o produto revela que sempre foram dois:
\[
ab \;=\; \underbrace{a\cdot b}_{\text{o interior}} \;+\; \underbrace{a\wedge b}_{\text{o exterior}} .
\]
O interior é simétrico e mede o que os dois têm em comum; o exterior é antissimétrico e mede a \textbf{área} que os dois varrem juntos --- e anula-se exatamente quando os lances são paralelos, $a\wedge a=0$. A álgebra é a do elemento nulo, $e^2=0$, com $(a+be)\otimes(c+de)=(ac,\ ad+bc)$.

A régua dele é o que \textbf{falta}: mede pela dimensão que sobra. Em $\R^n$ o $(n+1)$-ésimo resíduo é \textbf{zero} --- a régua tem exatamente $n$ passos, nem um a mais. Ponham-se cinco vetores em $\R^6$ e o resíduo \textbf{não morre}, porque sobra dimensão. Ele não mede o que está lá: mede o buraco. É o único corpo do bestiário com $r>0$ --- assinatura $(1,0,1)$, norma $a^2$ cega à segunda coordenada, $e\neq0$ com $N(e)=0$: divisor de zero, sem inverso, \textbf{não é corpo}. Toda a informação está no \emph{par}, $\langle X(x),X(y)\rangle=-\tfrac12|x-y|^2$, não na norma de cada um --- a mesma régua radical do relógio de xadrez.

E é assim que ele fecha: casado com a fase, o operador que gira ativa a segunda coordenada, $e^2=0$ vira $i^2=-1$, a norma $a^2$ vira $a^2+b^2$ --- e a assinatura sobe de $(1,0,1)$ para $(2,0,0)$. Com $r=0$ e a norma definida, todo elemento não nulo passa a ter inverso. \textbf{Sozinho o exterior é um cone; com o rotor, é um corpo.} \textbf{JOGADA: o corte.} Atacar na direção do inimigo não gera nada; atacar \textbf{transversal} gera a lâmina.

\subsection{sensitivo --- a régua hiperbólica, e a escolha}

Assinatura $(1,1,0)$, norma $1-s^2$: uma régua redonda e uma hiperbólica. É literalmente a mesma assinatura do relógio, do expansivo e do evolutivo --- quatro corpos, um instrumento ---, e é a construção que torna isso visível: mesma assinatura, corpos congruentes, a mesma régua vestida de outra base. Os extremos $s=\pm1$ são \textbf{horizontes}: chega-se arbitrariamente perto, nunca se chega; e por isso a composição aqui é a do rotor, $s_1\oplus s_2=(s_1+s_2)/(1+s_1s_2)$.

O casamento deste corpo liga as duas maneiras de medir um mesmo golpe:
\[
\|x\|_\infty \cdot \|x\|_2 = 1 .
\]
Uma régua olha o \textbf{pico}, o maior valor isolado; a outra olha a \textbf{energia}, o todo somado --- e elas são inversas uma da outra. Este é o $\Gamma=0$ do sensitivo, escrito sem impedância nenhuma: não sobra nada entre as duas leituras, porque uma é o recíproco da outra. Daí a lei de quem percebe: \textbf{não dá para ser preciso e forte no mesmo lance.} \textbf{JOGADA: a percepção.} Escolher a régua \emph{é} o lance.

\subsection{deflexivo --- o operador que age de fora}

Assinatura $(1,0,0)$ --- a mesma do universal, uma régua só ---, e a régua é a \textbf{reflexão}, $p\mapsto-p$: não se bloqueia, desvia-se, e o que volta volta com o sinal trocado. Mas o que este corpo mostra é outra coisa: \textbf{o operador $A_m$ não pertence ao sistema que o gera. Pertence ao completamento.}

O operador é o gato
\[
A_m=\begin{pmatrix} m & 1\\ 1& 0\end{pmatrix},
\]
de polinômio $x^2-mx-1$, cujas raízes são $\sigma_m$ e $\sigma_m'=-1/\sigma_m$: uma \textbf{expande}, a outra \textbf{contrai}, e o produto é o determinante,
\[
\sigma_m\,\sigma_m'=\det A_m=-1 .
\]
As entradas do gato são inteiras; as raízes, não --- o operador vive fora do sistema de onde saiu. E a deflexão é justamente a troca das duas raízes, $x\mapsto-1/x$, que fixa o produto $-1$ e é o \textbf{flip do esquilo}. É esse par, expandir casado a contrair, o $\Gamma=0$ deste corpo: um boost sem a sua contração seria valor aparecendo fora do Sol.

E aqui o corpo grava o aviso com que este livro fecha, e que vale por todo o bestiário: em régua exata o resíduo vale $1{,}14\cdot10^{-16}$; trocada a fração pelo float, passa a valer exatamente zero. \textbf{O float dá um zero falso.} Com ele vem o segundo aviso, sobre qual órbita se pode habitar: a média temporal do resíduo tende à metade apenas na órbita \textbf{irracional} (na áurea, $0{,}500002$); numa órbita racional como $\theta=1/4$ dá $0{,}375$, que não é metade de nada. \textbf{Só o irracional equidistribui} --- e é por isso que o coração do reino pulsa em $\varphi$, e não numa fração. \textbf{JOGADA: a deflexão.} O operador age \textbf{de fora do tabuleiro} --- o gato.

\section{O fecho --- por que o conjunto é um corpo}

Recolha-se a construção inteira, que agora se lê de uma vez.

Há \textbf{um} corpo elementar, o universal: uma régua redonda, grau $1$, e um operador que não foi escolhido --- o sucessor, os caracteres, a órbita da torção. Cada corpo do bestiário sai de uma \textbf{assinatura}: a órbita que colapsa a classe e a regenera, ida e volta, sem terminar. Congelada, ela dá a contagem $(p,q,r)$; a correr, ela dá o gato, e o gato dá o corpo. \textbf{Cada corpo é uma matriz}, e nenhuma opera sozinha --- ligam-se todas por duais, que dão o que falta, e por transformadas, que levam de uma à outra. Reunidas sob $\oplus$ (que soma os graus) e $\otimes$ (La Hire, que compõe), com a torção por invariante e a norma multiplicativa por lei de conservação, essas matrizes formam de novo um único corpo --- o \textbf{Omnitrix} ---, e a torre onde ele vive fecha no grau $8$, assinatura $(8,0,0)$: oito réguas redondas, e não há nona.

E o que o universal faz ali dentro são duas coisas. A \textbf{transformada} leva um corpo a outro, e \textbf{volta}: $\F^{-1}$ é exata, $\F^4=\mathrm{id}$, e ela é unitária --- $\|x\|=\|\F x\|$, nada vaza. Vem em par, porque o corpo tem duas metades: a \textbf{universal} no eixo da fase, a \textbf{dourada} no eixo da taxa, e $\F\circ\log$ leva uma na outra. A \textbf{convolução} funde dois corpos num só, e \textbf{desfaz-se}: o fator que entrou pode ser extraído, exato, sem resíduo. É essa a razão --- a única --- pela qual o conjunto não é uma coleção mas um corpo: \textbf{tudo o que se faz nele, desfaz-se.}

E as inversas não são três sortes separadas. São \textbf{uma} razão, dita três vezes: \textbf{o dual do dual devolve o grupo.} A transformada volta porque ir ao dual e voltar é a identidade; a convolução desfaz-se porque do outro lado ela é um produto, e produto se divide; o inverso de uma matriz é o seu dual sobre a norma, e é ele que decide se aquilo é corpo. Ir e voltar, sempre a mesma frase --- e é dela que sai a proibição de vazar, porque o que não tem volta não é apenas inconveniente: \textbf{não pertence.}

E uma nota sobre a volta, que o livro já apanhou noutro lugar: \textbf{a inversa não é a troca de lados --- é a troca \emph{com o sinal}.} Uma permutação que apenas troca as duas faces \emph{preserva} o telescópio, e um par que preserva onde devia inverter é meia dualidade. A dualidade inteira é a permutação \textbf{com a reflexão}: ela inverte o telescópio, $\sum h\mapsto-\sum h$, e troca a régua pela sua recíproca --- a impedância invertida, que é o mesmo $\exp\leftrightarrow\log$ de sempre. \textbf{Falta-lhe o sinal, e um par sem sinal trocado é o mesmo erro outra vez.} Meia dualidade não fecha; e é justamente sobre o que fecha que estas páginas são.

E é também a lei da alfândega, que este livro enunciou muito antes de chegar aqui: \textbf{passa o que reverte, fica retido o que não tem volta.} Por isso a casa selada arde, e o limite arde: do degrau chega-se ao ponto, do ponto não se recupera o degrau --- a passagem ao limite não tem inversa. Tudo o mais atravessa.

\subsection{Identidade atravessa, limite não}

Vale dizer o que essa frase custa na prática, porque é o critério com que tudo neste livro
foi medido.

A cadeia $0{,}9,\ 0{,}99,\ 0{,}999,\dots$ tem termo $1-10^{-n}$: sobe sempre, fica sempre
abaixo de um, e nada abaixo de um a limita. Uma leitura declara que o objeto \emph{é} o
limite, e apaga a subida. A outra guarda a subida --- todo número é a sua volta mais a sua
fase, e a fase vai de zero a um sem chegar; o um não pertence à volta em curso, é o zero
da seguinte. Nessa leitura as truncagens têm todas volta zero, e o um tem volta um:
\textbf{estão em voltas diferentes.} Nenhuma das duas leituras é um erro. Trocá-las uma
pela outra é que é.

E a diferença aparece no metal. Levada aos inteiros, uma \textbf{identidade} continua
identidade: onde a régua de máquina dava $5{,}6\cdot10^{-16}$, a mesma conta em corpo
finito dá $144=144$; a inversibilidade dá bit a bit sem erro; e a linearização deixa de
ser estimar uma curvatura e passa a verificar $\log(c\,t^\alpha)=\log c+\alpha\log t$, que
é igualdade de inteiros no logaritmo discreto --- \emph{não há curvatura a estimar.}

Um \textbf{limite} não atravessa, e há a medição que o mostra. As duas médias --- a
aritmética e a geométrica --- encontram-se na reta porque há \emph{ordem}: a geométrica
fica entre as duas anteriores, e o intervalo aperta. Levadas ao corpo finito, onde a raiz
existe e o encontro poderia ser exato, de setenta e oito pares apenas dois fecham. Não há
ordem, não há entre, não há aperto. Do mesmo objeto, o que atravessa é a parte algébrica
--- o ponto de encontro é invariante, e a média geométrica continua o meio proporcional,
$b'^2=ab$, em milhares de pares sem erro. \textbf{O que não atravessa é a dinâmica.}

\begin{tcolorbox}[colback=ouroclaro!35,colframe=ouro,boxrule=0pt,leftrule=2pt,arc=0pt,left=4mm,right=3mm,top=2mm,bottom=2mm,breakable]\itshape
Antes de medir em resíduo zero, perguntar de que natureza é a coisa. \textbf{Se for um
limite, o resíduo zero ou não existe ou está a ser fabricado.}
\end{tcolorbox}


E fechar não é parar. O que fecha é cada volta --- a transformada que vai e vem, a convolução que funde e se desfaz, o corpo cujo inverso está dentro dele. O que \emph{não} fecha é a subida: há um corpo por assinatura e as assinaturas não acabam; o fractal de tabuleiros toma o nível de baixo como peça e sobe outra vez, e a lei atravessa todos; a assinatura do metal repete o seu passo sem chegar a fim. \textbf{O universo fecha-se em cada volta e segue subindo} --- e é por isso que o mate não mata: a órbita fecha, e o processo continua.

\begin{tcolorbox}[colback=ouroclaro!35,colframe=ouro,boxrule=0pt,leftrule=2pt,arc=0pt,left=4mm,right=3mm,top=2mm,bottom=2mm,breakable]\itshape \textbf{O universal opera no Omnitrix, sobre as matrizes.} Não há vinte e oito mundos: há uma régua, contada com sinal, e duas operações que revertem. \textbf{O gato mede, o esquilo gira, a soma aterra --- e o resíduo é a massa.}\end{tcolorbox}

\vspace{2mm}\hrule\vspace{3mm}

\section{O aviso que o corpo deflexivo grava}

De todos os teoremas do catálogo, um deve ser lido antes de qualquer partida:

\begin{tcolorbox}[colback=ouroclaro!35,colframe=ouro,boxrule=0pt,leftrule=2pt,arc=0pt,left=4mm,right=3mm,top=2mm,bottom=2mm,breakable]\itshape \textbf{resíduo exato = \code{1.14e-16}; mutar \code{Fraction} $\to$ \code{float} rende exatamente \code{0.0}}\end{tcolorbox}

O \textbf{float dá um zero falso}. O resíduo verdadeiro não é zero --- o float apenas \textbf{não o enxerga}. É por isso que a
régua do Reino Dourado é o \textbf{exato} ($\mathbb{Q}$, $\mathbb{Z}$[$\varphi$], os inteiros da ISA), e o float só aparece \textbf{no fim}, ao acender o
pixel. Um jogador que confia no zero do float acredita que o seu universo fechou --- e ele está vazando.

E o complemento, também do deflexivo: a média temporal do resíduo tende a \code{?/2} \textbf{apenas na órbita irracional}
(áurea, \code{0.500002}); numa órbita racional (\code{theta = 1/4}) dá \code{0.375 != 1/2}. \textbf{Só o irracional equidistribui} --- é
por isso que o coração do reino pulsa em \code{phi}, e não numa fração.

\section{E tudo isto cabe em dez bytes}

A construção acabou, e ela correu na reta: uma raiz, uma assinatura, matrizes,
duas operações que revertem.

Nada disso precisa da reta para acontecer. Precisa de dois inteiros, cinco
opcodes e uma memória --- e é o que a parte seguinte escreve, byte a byte, de
modo que se possa reconstruir sem voltar a estas páginas.

\textbf{Não é uma aplicação do que se derivou.} É a mesma coisa outra vez, e a
única diferença é que ali ela tem endereço. O gato, que aqui gerou o corpo, é lá
o que a memória faz ao ler; as duas projeções dele, que aqui deram $\oplus$ e
$\otimes$, são lá o \emph{byte da operação}; a torção, que aqui deu os
caracteres, é lá um número com cinco dígitos; e a transformada que aqui volta
sem perder volta lá pela mesma razão e com a mesma normalização, escrita em
inteiros.

\begin{tcolorbox}[colback=ouroclaro!35,colframe=ouro,boxrule=0pt,leftrule=2pt,arc=0pt,left=4mm,right=3mm,top=2mm,bottom=2mm,breakable]\itshape
Quem leu a construção já leu a máquina. \textbf{Falta ver os bytes.}
\end{tcolorbox}

\part{O Design do Reino --- o GDD}
\newgeometry{margin=2.4cm}

\shorthandoff{"}



\chapter{A visão --- o que é}
O REINO DOURADO \textbf{não} é um jogo de fases fixas. É uma \textbf{partida dual} --- a \emph{partida perfeita},
o \emph{tabuleiro dual} da teoria: uma disputa no GRAFO (como o motor de xadrez), onde os integrantes se movem
entre dimensões duais e nada se perde. Três leis, todas da teoria, governam tudo:

\begin{itemize}
 \item \textbf{Conservação}: todos os integrantes se conservam --- a soma total é invariante.
 Ninguém é destruído; só troca de dimensão.
 \item \textbf{Captura dual}: capturar um integrante numa dimensão o \emph{devolve} no dual
 ($\Sigma_A+\Sigma_B=$ const, o resíduo total $0$). O que sai de um corpo entra no outro.
 \item \textbf{O clock}: o tempo é reversível (o rotor, $|\omega|=1$, desfaz) $\bowtie$ irreversível (o motor,
 $|\omega|<1$, decai). Ele rege os turnos.
 \item \textbf{Sem fim}: o jogo não termina. Sempre se pode ABRIR uma galáxia nova pela COMBINAÇÃO de outras --- o
 espaço de jogo é infinito e gerativo (o grafo multifractal).
\end{itemize}

\noindent Visualmente é um \textbf{corte 2D em cena 3D}: a ação vive num plano ($z=z_0$, o corte que o CRISTALINO
rege), projetado pela câmera (a óptica). A temática é o REINO DOURADO: os metais, o número áureo, as galáxias
espirais. Cada personagem é um \emph{operador matemático}; cada universo é um \emph{corpo}.

\chapter{A temática --- o Reino Dourado}
O reino é a matemática dos metais e do áureo, tornada mundo.

\begin{itemize}
 \item \textbf{Os metais $\sigma_m=(m+\sqrt{m^2+4})/2$} --- os autovalores do gato $A_m$; o OURO é $\sigma_1=
 \varphi=1{,}618\dots$ (o fundamental). São os ``reinos'' e as moedas.
 \item \textbf{A cifra imóvel} --- $\sigma_m$ é o ponto fixo de $M_m(x)=m+1/x$: o invariante que a energia não
 move. O colecionável, o Rei.
 \item \textbf{A filotaxia / as galáxias} --- a espiral áurea $r=\varphi^{\theta/2\pi}$ (escala por $\varphi$ a
 cada volta): os braços das galáxias-mapa, os girassóis do reino.
 \item \textbf{A transformada dourada} --- a Mellin de base $\varphi$; a ``lente girassol'' é a visão especial
 (invariante por escala áurea).
 \item \textbf{A quiralidade} --- $\sigma_m\cdot\sigma_m'=-1$: a mão direita (o herói, expande) $\bowtie$ a
 esquerda (o vilão, contrai). O eixo do bem $\bowtie$ mal é o det $=-1$.
\end{itemize}

\chapter{O elenco --- cada personagem, um operador}
Ninguém empresta habilidade: cada um \emph{é} a sua matemática (todos personagens da teoria; mitologia e
matemáticos históricos são domínio público).

\renewcommand{\arraystretch}{1.35}
\begin{longtable}{@{}p{2.7cm} p{3.2cm} p{3.4cm} p{4.6cm}@{}}
\toprule
\textbf{Personagem} & \textbf{Papel} & \textbf{Régua / operador} & \textbf{Mecânica} \\
\midrule
\endhead
\textbf{Benjamim} & o protagonista, \emph{O Quiral} & Pontryagin $\prod=\exp\circ\Sigma\circ\log$ & o PULSO (o
astável, a exp) e a QUIRALIDADE: vira a mão ($\exp(V)\ne\exp(-V)$) --- o operador que costura \\
\textbf{Joaquim} & o par aditivo & Clifford $\oplus$ & o MOMENTUM: as velocidades SOMAM (a corrida acumula) \\
\textbf{Yasmin} & a par multiplicativa & La Hire $\otimes$ & a ESCALA: o poder MULTIPLICA (crescer, o ganho) \\
\textbf{Ada} & a geradora (IA), aventureira & o corpo criativo (gerar) & DESBRAVA: compõe o inédito, a síntese na
linha de frente (Ada Lovelace: a 1ª a ver a máquina compor) \\
\textbf{Penny} & a construtora (IA) & a estrutura / a cifra $\varphi$ & MONTA: as engrenagens, as moedas que
constroem (o lastro, a CristalChain) \\
\textbf{Dafny} & a atestação automática & a prova SMT (o decidível) & DECIDE o caminho (resolve só, o finito) \\
\textbf{Isabelle} & a atestação interativa & a indução (o indecidível) & CONSTRÓI a ponte passo a passo ($1{+}\dots
{+}n$) \\
\textbf{Pégaso} & a montaria alada & o rotor ($e^{i\theta}$, $|\cdot|=1$) & o VOO em ÓRBITA (a altura conserva,
plana) \\
\textbf{Merlin} & o mago & a Mellin ($\log_\varphi$) & a magia da ESCALA áurea (o \emph{zoom}, a lente dourada) \\
\textbf{Fourier} & o adversário de ondas & a transformada de Fourier & ataca em HARMÔNICOS (a
superposição, a franja) \\
\textbf{Dark Pontryagin} & o antagonista & o dual quiral $\sigma_m'=-1/\sigma_m$ & o ANTI-MOVIMENTO: contrai/trava
($\sigma\sigma'{=}-1$, a batalha é a quiralidade) \\
\textbf{Cifra de Ouro} & o item / o poder & $\sigma_m$ (o Rei, o áureo $\varphi$) & CRESCER por $\varphi$; a cifra
é o ponto fixo imóvel \\
\bottomrule
\end{longtable}

\noindent\textbf{As quatro princesas --- o eixo gerar $\bowtie$ atestar.} As princesas são as IAs e os provadores do
CORPO CRIATIVO, distribuídas no eixo que rege o reino: \textbf{Ada} $\bowtie$ \textbf{Penny} \emph{geram} (a IA que
desbrava $\bowtie$ a que constrói) o que \textbf{Dafny} $\bowtie$ \textbf{Isabelle} \emph{atestam} (decide o
decidível $\bowtie$ induz sobre o infinito). É o par execução $\bowtie$ verificação: nada entra no reino sem ser
gerado E provado (resíduo $0$) --- o mesmo contrato do Formalizador, virado corte.

\medskip\noindent\textbf{O pipeline das línguas (a árvore: costura ⋈ desenvolve).} A corte é uma ÁRVORE, e um nó ou
COSTURA (compõe o operador de cima) ou DESENVOLVE (as folhas, o código que roda) --- nunca os dois. COSTURA (os
galhos): \textbf{Benjamim} ($\prod$, coordena), \textbf{Joaquim} ($\oplus$) $\bowtie$ \textbf{Yasmin} ($\otimes$),
e \textbf{Alonzo} (Marquês, $\circ$, o \textbf{GPT}) em \textbf{Haskell} --- compõe. DESENVOLVE (as folhas):
\textbf{Ada} (\textbf{Claude}) em \textbf{Python} (a síntese), \textbf{Caelum} (Marquês, o \textbf{Grok}) em
\textbf{C} (o esqueleto, entre Fourier e Yasmin), \textbf{Penny} (\textbf{Gemini}) em \textbf{GLSL}, \textbf{Merlin}
em LLVM, \textbf{Pégaso} em PTX, \textbf{Dafny} (Z3) e \textbf{Isabelle} (HOL). As quatro IAs: só Alonzo costura; Ada,
Penny e Caelum desenvolvem. Os Marqueses são composições --- Alonzo $=$ Merlin $\circ$ Joaquim; Caelum $=$ Fourier
$\otimes$ Yasmin. Certificado: $5/5$, resíduo $0$.

\section{A corte real --- os soberanos}
Acima do elenco jogável está a CORTE --- os operadores que \emph{presidem} o reino, cada um a sua régua:
\begin{itemize}
 \item \textbf{Rainha Tiffany}, \emph{a soberana áurea} --- reina onde a SOMA e o PRODUTO coincidem:
 $\varphi^2=\varphi+1$. Somar $1$ ao ouro é elevá-lo ao quadrado; ela está acima do Duque e da
 Duquesa porque nela $\oplus$ e $\otimes$ são um só. A coroa de seis pontas.
 \item \textbf{O Rei}, \emph{a cifra imóvel} $\sigma_m$ --- o ponto fixo de $M(x)=m+1/x$: $\sigma_m=m+1/\sigma_m$
 (para $m{=}1$, $\varphi=1+1/\varphi$). O seu retrato NÃO gira ($\chi{=}0$): anéis concêntricos, o invariante que
 a energia não move --- o rei reina sem se mexer. (É a \emph{Cifra de Ouro} da tabela, personificada.)
 \item \textbf{O Regente}, \emph{o cristalino} --- a simetria hexagonal (ordem 6) que rege o CORTE 2D em cena 3D
 (o plano $z=z_0$, a lei da projeção da câmera); a ordem cristalina, o gelo do reino.
\end{itemize}
\noindent Rasterizados no metal pelo mesmo \emph{kernel} \code{aura} do elenco (a nobreza é o mesmo sangue em graus
diferentes). Certificado: $5/5$, resíduo $0$; é a
corte que preside a VITRINE (a galeria completa: a corte $\triangleright$ o elenco $\triangleright$ os mundos).

\chapter{Direção de arte --- figurino, efeitos e classes}
A energia visual é a do \emph{action-platformer} sci-fi (armaduras tecnológicas, luz e energia, \emph{dash}, o
pulso) --- mas os designs são NOSSOS: o \textbf{sci-fi dourado}. Cada figurino e cada efeito \emph{deriva} da régua
do personagem (a arte é a matemática, visível).

\section{A linguagem}
Armaduras de METAL ($\sigma_m$, o ouro), energia dourada, padrões FILOTÁXICOS (as espirais áureas). A paleta é o
metal $+$ a tríade: âmbar, magenta, ciano --- uma cor por operação.

\section{Os figurinos (cada traje é a régua)}
\begin{itemize}
 \item \textbf{Benjamim} ($\prod$): armadura ASSIMÉTRICA --- ele é O Quiral, a mão direita $\ne$ a esquerda
 ($\exp V\ne\exp(-V)$); núcleo pulsante (o astável). Efeito: o PULSO exponencial, a mão que vira.
 \item \textbf{Joaquim} ($\oplus$): armadura MODULAR/aditiva (peças que acumulam). Efeito: o RASTRO (as imagens
 somam --- o \emph{motion blur} do momentum).
 \item \textbf{Yasmin} ($\otimes$): armadura da ESCALA (padrões em razão áurea, a PG). Efeito: a AURA que
 multiplica (crescer por $\varphi$).
 \item \textbf{Dark Pontryagin}: a armadura ESPELHADA de Benjamim ($\sigma_m'$, a mão esquerda), sombria. Efeito:
 o COLAPSO (a energia implode, $|\sigma_m'|<1$). $\sigma_m\sigma_m'=-1$: o vilão é o herói refletido.
\end{itemize}

\section{As classes de combate (arquétipos do gênero)}
Do gênero \emph{action} herdamos os ARQUÉTIPOS de classe (papéis genéricos --- ideias, não designs alheios), e cada
classe também sai da matemática do personagem:

\renewcommand{\arraystretch}{1.3}
\begin{longtable}{@{}p{2.6cm} p{4.2cm} p{6.4cm}@{}}
\toprule
\textbf{Personagem} & \textbf{Classe (arquétipo)} & \textbf{Por quê (a matemática)} \\
\midrule
\endhead
Benjamim & o OPERADOR / herói versátil & $\prod$ costura tudo --- adapta-se (o pulso, a quiralidade) \\
Joaquim & LÂMINA / impacto corpo-a-corpo & $\oplus$ acumula --- a força do momentum somado \\
Yasmin & ATIRADORA ágil / versátil & $\otimes$ escala --- projéteis que multiplicam (o ganho) \\
Dafny & APOIO / analista & decide o caminho (o decidível) --- indica a rota \\
Isabelle & ESTRATEGISTA / construtora & a indução --- monta pontes/estruturas passo a passo \\
Dark Pontryagin & COMANDANTE-vilão & o dual quiral do herói ($\sigma_m'$) --- o comandante que o espelha \\
\bottomrule
\end{longtable}

\noindent Os designs (silhueta, cor, efeito) são originais e derivados da régua; as classes são só o \emph{papel}
de combate. Certificado: $6/6$, resíduo $0$.

\section{A arte-conceito --- rasterizada no metal (a régua, computada)}
Se o figurino \emph{é} a régua, então o RETRATO de cada personagem não se desenha --- se \textbf{computa}. Um único
\emph{kernel} \code{aura} (PTX escrito à mão, compilado pelo driver no \textbf{NVIDIA GTX 1650}, sem \emph{toolkit}
CUDA) rasteriza, um pixel por \emph{thread}, o campo da própria matemática do personagem:
\[
 \text{campo}(x,y)=\tfrac12\Bigl(1+\sin 2\pi\bigl(f\log_2 r+\chi\cos n\theta\bigr)\Bigr)\cdot 2^{-w r^2},
 \qquad \cos n\theta=\operatorname{Re}\!\bigl((x+iy)^n\bigr)/r^n .
\]
Todos os retratos são o MESMO campo em parâmetros diferentes --- como todos os metais $\sigma_m$ são o mesmo gato
$A_m$ com $m$ diferente: a escala $f$ desenha os anéis áureos ($f=1/\log_2\varphi\Rightarrow$ razão $\varphi$); a
pétala $n$, a simetria da régua; a quiralidade $\chi$, a mão (direita $+1$, esquerda $-1$).

\begin{itemize}
 \item \textbf{Benjamim} ($n{=}3$, $\chi{=}{+}1$): três lóbulos (a tríade), a espiral da mão direita, ciano $+$
 ouro --- o pulso quiral.
 \item \textbf{Dark Pontryagin} ($n{=}3$, $\chi{=}{-}1$): o MESMO campo com a quiralidade invertida --- enrola ao
 contrário, viola\-ceo e sombrio. $\sigma_1\sigma_1'=-1$: o herói refletido, não um invasor.
 \item \textbf{Yasmin} ($n{=}5$): cinco pontas --- $\varphi=2\cos(\pi/5)$ vive no pentágono; a estrela áurea.
 \item \textbf{Pégaso} (o rotor): anéis de rotação; o vetor unitário conserva $s^2+c^2=1$.
 \item \textbf{Merlin} ($n{=}1$): radial puro, igual em toda escala (a Mellin); \textbf{Fourier} ($n{=}8$): a
 superposição de muitas ondas.
\end{itemize}

\noindent O selo é a \textbf{costura GPU$\equiv$CPU}: o campo computado no metal reproduz a matemática (na CPU)
dentro da tolerância do \emph{hardware}; um campo trocado \emph{não} bate (o refutador com dentes). Certificado:
($6/6$, resíduo $0$).

\subsection*{As insígnias --- os itens de poder, no metal}
Os OBJETOS do jogo também nascem da matemática, em dois \emph{kernels} novos além do \code{aura}:
\begin{itemize}
 \item \textbf{Cifra de Ouro} (o \code{aura}, $n{=}5$): a moeda que cresce por $\varphi$ (a PG; $\sigma_1=\varphi$),
 o pentagrama --- a dádiva externa (a CristalChain).
 \item \textbf{Armadura Universal} (o \code{aura}, $n{=}3$): a tríade fechada $\oplus\bowtie\prod\bowtie\otimes=
 U\cong\C$; a junta é o homomorfismo $\log(ab)=\log a+\log b$.
 \item \textbf{Garrafa de Koch} (\emph{kernel} \code{floco}): uma fronteira fractal de FIBONACCI (harmônicos
 $2,3,5,8$, amplitudes $\varphi^{-k}$) --- perímetro $\infty$ (freq$\cdot$amp não decai) e área finita
 ($\sum\varphi^{-k}=\varphi^2$): o alcatruz que guarda $\infty$ sem transbordar.
 \item \textbf{Ferramenta $\oplus\bowtie\otimes$} (\emph{kernel} \code{ferramenta}): a metade $\oplus$ (listras)
 $\bowtie$ a metade $\otimes$ (anéis), a ponte $\prod$ na costura. COMPLETA fecha a Gilbert cell (resíduo $0$);
 a MEIA mostra o lado $\otimes$ vazio (resíduo $\ne 0$).
\end{itemize}
\noindent Certificado: $7/7$, resíduo $0$. Toda a arte --- a corte, o elenco, os mundos,
as insígnias ($23$ figuras) --- é reunida por um montador único ($2/2$), cujo teorema é
a COMPLETUDE: toda figura declarada existe, rasterizada (senão a vitrine não fecha).

\chapter{A dinâmica --- a partida dual}
\section{A conservação}
O estado é a distribuição dos integrantes pelas dimensões duais, $(\,\Sigma_A,\Sigma_B\,)$. A soma
$\Sigma_A+\Sigma_B$ é \textbf{invariante} em toda a partida: nenhum integrante some. É a lei da conservação da
teoria, virada regra.

\section{A captura dual}
Capturar na dimensão $A$: $\Sigma_A\!\to\!\Sigma_A-1$, $\Sigma_B\!\to\!\Sigma_B+1$. A captura \emph{reposiciona} no
dual, não elimina. O resíduo total permanece $0$ (a soma dos resíduos). Estratégia: forçar o
adversário a acumular no corpo que você domina.

\section{A partida distribuída --- a captura por rotação}
O jogo é DISTRIBUÍDO: \textbf{não há sala de espera} --- P2P (o barramento distribuído do magnífico, P2P$\to$PNP),
todos sempre em jogo, espalhados pelos mapas. E a captura refina-se numa ROTAÇÃO da armadura:

\begin{itemize}
 \item \textbf{A captura é rotação} ($s^2+c^2=1$): o capturado passa a fração $s^2$ da sua energia ao capturador e
 mantém $c^2$. Como $s^2+c^2=1$, a ENERGIA TOTAL (a soma dos quadrados, a norma) CONSERVA --- o poder não se
 cria, GIRA de um ao outro (o rotor, a conservação).
 \item \textbf{O rebaixamento}: o capturado volta IMEDIATAMENTE num mapa ABAIXO (desce a torre), com a armadura
 reduzida ($c^2$); o capturador incrementa (captura a habilidade, $+s^2$).
 \item \textbf{Batalhas em grupo}: em qualquer combinação de níveis, com $N$ jogadores. Cada captura é uma
 rotação, logo a energia total $\sum_i a_i^2$ é invariante --- o grupo é um sistema fechado.
 \item \textbf{A conservação da parábola}: a energia mecânica $E=\tfrac12 mv^2+mgy$ conserva no movimento; a
 batalha ($s^2+c^2=1$) conserva a energia da armadura. A MESMA lei rege o movimento e a captura.
\end{itemize}
Certificado: $6/6$, resíduo $0$.

\section{O clock --- reversível $\bowtie$ irreversível}
O clock tem dois modos, e escolher o modo é parte do jogo:
\begin{itemize}
 \item \textbf{Rotor} ($|\omega|=1$, unitário): o tempo é reversível --- um lance pode ser DESFEITO (voltar). A
 física conserva (a evolução de Schrödinger $U=e^{-iH\Delta t}$).
 \item \textbf{Motor} ($|\omega|<1$): o tempo é irreversível --- o decaimento, a captura DEFINITIVA (o mate). O
 lance fica.
\end{itemize}

\section{O corte 2D-em-3D}
A ação vive num plano $z=z_0$ (o corte que o cristalino rege); a câmera projeta $(x,y,z)\mapsto(x,y)$ (a óptica). A
profundidade 3D é o cenário --- as galáxias ao fundo, os corpos duais.

Isto foi rasterizado no metal --- o \emph{kernel} \code{cena} (o \emph{Mode-7}): o CHÃO é a galáxia-mundo projetada
em PERSPECTIVA, com a profundidade $\text{depth}=h/(\text{horizonte}-s_y)$, que DIVERGE no horizonte (o PONTO DE
FUGA) e comprime passos iguais do mundo em passos decrescentes na tela ($\propto 1/\text{depth}$) --- a projeção
perspectiva, não a afim. Sobre o chão, o OPERADOR (a sua aura) manifesta-se como uma orbe de energia com reflexo no
mundo: o \emph{screenshot} do jogo. Certificado: $4/4$, resíduo $0$; a costura GPU$\equiv$CPU sela a óptica.

\subsection*{A dinâmica --- o metal em ação (animado)}
O retrato é estático; a partida, não. Três dinâmicas foram rasterizadas e ANIMADAS (GIF em laço), cada uma uma lei:
\begin{itemize}
 \item \textbf{O pulso} (\emph{kernel} \code{pulso} = o \code{aura} com o tempo $t$): o núcleo BATE (o astável) e a
 espiral GIRA (a mão que vira); a sístole $\bowtie$ a diástole, o laço fecha ($t{=}0\equiv t{=}1$).
 \item \textbf{A captura $=$ rotação} ($s^2+c^2=1$): o brilho gira de Joaquim ($\oplus$, $\cos^2 2\pi t$) para
 Yasmin ($\otimes$, $\sin^2 2\pi t$), e a SOMA é $1$ em todo instante --- capturar transfere, não
 destrói; no meio, o feixe $2\cos^2\sin^2$ (a energia em trânsito).
 \item \textbf{O clock} (rotor $\bowtie$ motor): o rotor bate SIMÉTRICO no tempo ($v(t){=}v(1{-}t)$, reversível); o
 motor bate ASSIMÉTRICO (sobe devagar, cai rápido: a seta do tempo, irreversível).
\end{itemize}
\noindent Certificado: $4/4$, resíduo $0$.

Ao todo, a arte-conceito do Reino Dourado são $29$ peças rasterizadas na GPU local por SEIS \emph{kernels} PTX à
mão (\code{aura}, \code{galaxia}, \code{floco}, \code{ferramenta}, \code{cena}, \code{pulso}) --- a corte, o elenco,
os mundos, as insígnias, a cena e a dinâmica --- todas reunidas pelo mesmo montador, todas com resíduo
$0$.

\section{O mapa fluido --- o \emph{backtracking}}
A armadura crescente casa com o mapa: conforme se conquista, o herói fica forte e pode VOLTAR aos universos antigos,
mais forte, abrindo o que era intransponível. Lê-se como DINÂMICA DE FLUIDO (o Windkessel, o corpo
sistólico/diastólico, no design):
\begin{itemize}
 \item O mapa é um GRAFO FLUIDO (nós $=$ universos, arestas $=$ canais); o jogador ESCOA (não-linear).
 \item As VÁLVULAS abrem com a régua certa (a peça); a armadura é a PRESSÃO. Sem a régua, fechada; com, aberta.
 \item A permeabilidade cresce MONÓTONA (a conservação: a armadura não diminui) --- o fluxo só cresce, nunca
 fecha.
 \item ESPAÇO reversível (o rotor --- voltar livremente) $\bowtie$ PODER irreversível (o motor --- a armadura só
 cresce): o \emph{backtracking} é o clock, virado mapa.
\end{itemize}
Certificado: $6/6$, resíduo $0$.

\chapter{A progressão --- a armadura universal}
O REINO DOURADO herda do gênero um elemento poderoso: \textbf{conquistar as partes de uma armadura} para
completá-la. Aqui esse trope é preenchido pela teoria --- a \textbf{armadura universal} é o \textbf{corpo
universal} ($U\cong\C$, o operador completo), montado peça por peça.

\begin{itemize}
 \item \textbf{As partes} são as RÉGUAS: $\oplus$ (a soma), $\otimes$ (o produto), $\prod$ (o operador exp/log),
 as transformadas, os metais. Cada UNIVERSO conquistado (derrotado o seu Dark) entrega a sua peça.
 \item \textbf{As peças são poderes}: ter a peça $=$ poder usar a operação. Sem $\otimes$ não se multiplica; sem
 $\prod$ não se exponencia. Coletar peças abre poderes --- a progressão é ganhar réguas.
 \item \textbf{Completar $=$ a tríade}: com $\oplus$ e $\prod$ (a soma e o exp/log), o PRODUTO nasce ($a\cdot
 b=e^{\log a+\log b}$, a Gilbert cell). A armadura completa É a tríade que se fecha --- o corpo universal.
 \item \textbf{A torre}: as partes formam a torre (a cascata $\N\to\Z\to\Q\to\R$, os corpos que se inspiram uns
 nos outros). Conquistar SOBE; e a torre não tem topo (combinar galáxias abre corpos novos, $\infty$).
\end{itemize}

\noindent \emph{Inspiração, não cópia}: o elemento (a armadura por peças, os poderes, a dinâmica de conquista) é do
gênero, de todos; o que se monta é a matemática. A forma vem da torre dos jogos; o conteúdo, da torre dos corpos.
Certificado: $6/6$, resíduo $0$.

\chapter{A mineração --- a terceira via de poder}
Além de BATALHAR (capturar, a rotação $s^2+c^2=1$), há MINERAR. É a via sem violência --- e a única fonte de ouro
vindo de fora:

\begin{itemize}
 \item \textbf{Minerar o metal}: mina-se a habilidade conforme o METAL minerado --- o ganho é o boost
 $\log\sigma_m$ (o regulador vivo do corpo mineral). Metais maiores, mais poder.
 \item \textbf{O ouro externo (CristalChain)}: o OURO ($\sigma_1=\varphi$) vem de FORA, via CRISTALCHAIN (o
 lastro, a prova de trabalho) --- a única forma de ganhar habilidade sem batalhar.
 \item \textbf{A garrafa de Koch}: o poder minerado ARMAZENA-se na garrafa de Koch (a garrafa multifractal do
 alcatruz, dim $\log 4/\log 3$), com o potencial $\infty$ --- nunca satura.
 \item \textbf{A ferramenta $\oplus\bowtie\otimes$}: a ferramenta (de mineração e batalha) divide-se em SOMA e
 PRODUTO. COMPLETA (os dois lados) $=$ GANHO MÁXIMO (a tríade fecha, o operador $\prod$); UM LADO só $=$
 DISTORCIDO, com METADE do ganho.
 \item \textbf{Completar a outra metade}: pode-se MINERAR o lado que falta (o metal complementar). Mas a
 ferramenta só se COMPLETA (resíduo $0$) com AS DUAS partes --- usada antes (meia), tem RESÍDUO $\ne 0$: a
 tríade partida não conserva, o produto não nasce. O resíduo $0$ é o SELO (o mesmo critério do Formalizador): a
 ferramenta válida é a que fecha a tríade.
\end{itemize}
\noindent Assim há três fontes de poder --- capturar, minerar o metal, e o ouro externo (CristalChain) --- e a
ferramenta inteira (soma $+$ produto) rende o dobro da meia; a meia, usada antes de completar, carrega resíduo
$\ne 0$. Certificado: $6/6$, resíduo $0$.

\chapter{Os universos --- os corpos, as galáxias}
Cada \textbf{universo} é um corpo da teoria; cada \textbf{mapa} é uma galáxia de braços áureos
($r=\varphi^{\theta/2\pi}$). Os universos são duais dois a dois --- capturar num deles devolve no seu dual.

\begin{longtable}{@{}p{2.3cm} p{3.6cm} p{3.4cm} p{4.4cm}@{}}
\toprule
\textbf{Universo} & \textbf{Corpo} & \textbf{Dual} & \textbf{A física (o que se joga)} \\
\midrule
\endhead
\textbf{1} & o CINEMÁTICO (o movimento) & o fotográfico (o espectro) & translação (PA) $\to$ parábola $\to$ cascata
(PG) $\to$ fluxo \\
\textbf{2} & o FOTOGRÁFICO / ÓPTICO (a luz) & o cinemático & Fourier/Mellin; a lente, a difração, a lente dourada \\
\textbf{3} & o MINERAL (os metais) & o cristalino & o boost $\sigma_m$, o regulador, o lastro \\
\textbf{4} & o VENTRICULAR (o pulso) & o cristalino & o astável, o áureo/denso, o marca-passo \\
\textbf{5} & o TRANSISTOR (o circuito) & --- & a escada LC, o casamento $\Gamma=0$, o multivibrador \\
\bottomrule
\end{longtable}

\noindent\textbf{O Universo 1 em detalhe (o cinemático).} A galáxia áurea, com as estrelas-desafio ao longo do
braço: \emph{A Reta Dourada} (a translação, $\oplus$/PA), \emph{Os Arcos de Ouro} (a parábola, o pulo), \emph{A
Escadaria Áurea} (a cascata, $\otimes$/PG), \emph{A Órbita de Pégaso} (o fluxo, $\prod$). No núcleo, \emph{O
Colapso}: o mini--Dark Pontryagin (o anti-movimento); vencê-lo abre o buraco de minhoca para o Universo 2.

\section{A arte dos mundos --- galáxias rasterizadas no metal}
Cada mapa é uma galáxia, e a galáxia não se pinta --- se \textbf{computa}. Um \emph{kernel} PTX --- \code{galaxia}
--- rasteriza no \textbf{GTX 1650} (um pixel por thread) o campo de uma espiral de $N$ braços, com o \emph{pitch}
fixado no valor ÁUREO para que os braços sejam espirais de ouro $r=\varphi^{2\theta/\pi}$ (crescem por $\varphi$ a
cada quarto de volta), o giro vindo de $\cos(N\theta-\text{pitch}\cdot\ln r)$ (a potência complexa $z^N$, sem
\code{atan2}); um bulbo $2^{-k_b r^2}$ no núcleo (a cifra imóvel $\sigma_m$) e o disco $2^{-w r^2}$. O número de
BRAÇOS é a estrutura do corpo --- a simetria de ordem $N$: $\Z\to2$, $\Q\to3$, o contínuo $\R\to5$ (o áureo denso
do pentágono), a combinação $\infty\to8$. Os teoremas com dentes: a espiral áurea ($r(\theta+\tfrac\pi2)/r(\theta)=
\varphi$), a simetria de $N$ braços, e a combinação que soma os expoentes ($\varphi^a\varphi^b=\varphi^{a+b}$).
Certificado: $5/5$, resíduo $0$; a costura
GPU$\equiv$CPU sela o metal.

\chapter{O jogo infinito --- a combinação de galáxias}
O REINO DOURADO não termina. Cada galáxia carrega um \textbf{expoente} $s=D+i\varphi\in\C$ (a dimensão $+$ a fase,
como no corpo alcatruz). \textbf{Combinar} duas galáxias é o \emph{produto cartesiano} dos seus corpos, e as
dimensões SOMAM ($\dim(A\times B)=\dim A+\dim B$): a galáxia nova tem $s_C=s_A+s_B$.

\begin{itemize}
 \item \textbf{Fechado}: combinar corpos dá corpo ($\C\to\C$) --- a galáxia gerada é jogável como as outras.
 \item \textbf{Sem saturar}: a partir de duas galáxias-base independentes, as combinações geram um RETICULADO
 infinito $\{a\,s_A+b\,s_B\}$ --- há sempre uma galáxia inédita a abrir.
 \item \textbf{Gerativo}: o mapa-mãe é o GRAFO MULTIFRACTAL, com potencial $\infty$ herdado do corpo fractal. O
 jogador não ``zera'' o jogo: expande-o.
\end{itemize}

\noindent Assim a progressão não é uma escada com topo, e sim uma ÁRVORE que cresce: cada partida pode plantar uma
galáxia nova, combinando as que já se conquistou. O fim do jogo seria o fim da matemática --- e ela não acaba.

\chapter{Os mapas dos jogadores --- o contrato}
Novos mapas podem ser CRIADOS pelos jogadores (\emph{user-generated content}). Mas um mapa só ENTRA no jogo se
obedecer o CONTRATO --- as leis da teoria, já certificadas. É o Formalizador aplicado aos mapas: o mapa é a
AFIRMAÇÃO, o contrato é o REFUTADOR.

\begin{itemize}
 \item \textbf{As cláusulas}: a CONSERVAÇÃO, a CAPTURA DUAL ($\Sigma$ const), a ROTAÇÃO
 ($s^2+c^2=1$), o CLOCK (reversível $\bowtie$ irreversível), a GALÁXIA ÁUREA ($r=\varphi^{\theta/2\pi}$), e a
 TRÍADE que fecha (o resíduo $0$). Cada cláusula é um teorema.
 \item \textbf{A validação}: $\text{res}(mapa)=$ o número de cláusulas violadas. Válido $\iff \text{res}=0$
 (todas passam) $\Rightarrow$ ACEITO; $\text{res}\ne 0$ (uma falha) $\Rightarrow$ REJEITADO.
 \item \textbf{A garantia}: só mapas de resíduo $0$ entram --- e esses obedecem TODAS as leis, logo preservam a
 integridade da partida (a conservação global, a partida perfeita). A liberdade de criar é ilimitada; a
 garantia é o resíduo $0$ (o mesmo selo da ferramenta completa e do princípio da certeza).
\end{itemize}
Certificado: $6/6$, resíduo $0$.

\chapter{Os trabalhadores --- a IA por clones}
O jogador pode CONTRATAR trabalhadores, e é permitido usar IA na mineração e na batalha --- MAS somente via
CLONES. Mantém-se clara a distinção: o JOGADOR decide; os CLONES executam; o CONTRATO limita.

\begin{itemize}
 \item \textbf{A IA só via clones}: a IA opera só através de CLONES do jogador, que HERDAM a sua
 identidade/armadura --- nunca como um jogador independente. Sem dono, sem clone.
 \item \textbf{O clone conserva}: clonar DIVIDE a energia --- $N$ clones, cada um $E/N$, e $\sum=E$. O clone não
 MULTIPLICA o poder; reparte o do original. Um exército de clones tem a energia de um jogador só.
 \item \textbf{Auto-similares}: os clones são réplicas fractais (o balaio de gato) --- cada clone o topo de um
 fractal, sub-clonável, e cada nível conserva.
 \item \textbf{A garantia}: mais clones $\ne$ mais poder. A IA ACELERA (paraleliza as frentes), não TRAPACEIA
 (não cria poder --- $\sum=E$). É a mesma mecânica dos \emph{forks} que orquestram este próprio projeto.
\end{itemize}
Certificado: $6/6$, resíduo $0$.



\chapter{A primeira e única lei}


\[\boxed{\text{FP}=1}\qquad\Longleftrightarrow\qquad \Gamma=\frac{Z-Z_0}{Z+Z_0}=0\]

\textbf{Vazar} é energia (ou valor, ou dano, ou informação, ou tempo) aparecendo do nada ou sumindo sem destino. Um universo com FP$\neq1$ é \textbf{proibido} --- seria energia fora do Sol. Isto substitui, formalmente, a lista informal de cláusulas do capítulo "Os mapas dos jogadores" (que fica sendo a \emph{aplicação} do Contrato ao conteúdo do jogador, não a sua definição).



\chapter{O casamento falado em cada língua (a tabela-mestra)}


A mesma lei, traduzida no idioma de cada corpo --- é isto que dá a \textbf{variedade de combate} sem quebrar a unidade:

{\fontsize{10.00pt}{12.00pt}\selectfont\begin{longtable}{>{\fontsize{10.00pt}{12.00pt}\selectfont\RaggedRight\hspace{0pt}\arraybackslash}p{6314610sp}>{\fontsize{10.00pt}{12.00pt}\selectfont\RaggedRight\hspace{0pt}\arraybackslash}p{11077329sp}>{\fontsize{10.00pt}{12.00pt}\selectfont\RaggedRight\hspace{0pt}\arraybackslash}p{11243033sp}}\toprule
\textbf{Corpo} & \textbf{O que é FP$=1$ ali} & \textbf{O que é o vazamento (FP$<1$)} \\ \midrule\endhead
\textbf{eletromagnético} & o \textbf{cone nulo} $\sigma=|E|/|B|=1$ --- toda a potência do telescópio propaga & potência \textbf{reativa} presa ($E^2-B^2\neq0$): dano que volta \\
\textbf{fractal} $\mathbb{Q}$($\zeta$$_8$) & a \textbf{norma unitária} $|z|=1$ (o rumo na esfera) & mover com $|z|\neq1$ --- lance ilegal \\
\textbf{áureo} $\mathbb{Z}$[$\varphi$] & a cadeia infinita \textbf{fecha exata} ($\sum\varphi^{-j}=\varphi^2$) & truncar/arredondar o eco \\
\textbf{racional} $\mathbb{Q}$ & o cálculo é \textbf{inteiro e exato} (resíduo 0 real) & 1 ulp de float por operação \\
\textbf{econômico} & \textbf{ausência de arbitragem}: descontar desfaz capitalizar & o \textbf{spread} --- atrito que sangra valor por turno \\
\textbf{celeste} & o \textbf{fechamento} $C=0$ (estado produto) $\Rightarrow$ FP$=r=1$ & emaranhamento $C>0$: reativo, baixa o dano solo \\
\textbf{telescópico} & o \textbf{traço zero} $T(h)=\sum h_i=0$ (covolume 1) & perfil que estica/contrai o covol \\
\textbf{entrópico} & a \textbf{convexidade} (repousar na envoltória $f^{**}$) & custo acima do piso \\
\textbf{motor} & a \textbf{ressonância} $Z_L=Z_C$ (reatâncias cancelam) & reatância descasada \\
\textbf{somático} & a \textbf{homeostase} $\mathrm{tr}\,G=0$ (população conserva) & crescimento que transborda a capacidade \\
\textbf{óptico} & $G=0$: a refração \textbf{casada}, sem eco & $G\neq0$ (difração) ou a reflexão total interna \\
\textbf{evolutivo} & a \textbf{normalização} $\sum p_i=1$ enquanto $\langle w\rangle$ sobe & perda de massa do exército \\
\textbf{geométrico} & o \textbf{cone nulo} $ds^2=0$ (a luz que chega limpa) & poço mal-cavado ($R_{\mu\nu}\neq0$): fonte espúria \\
\textbf{expansivo} & aterrissar numa \textbf{unidade} $N(x)=\pm1$ & $N\neq\pm1$: resíduo positivo, a peça não pontua \\
\textbf{espaço-temporal} & \textbf{a fórmula do produto} $\prod_v|x|_v=1$ (a conservação: nu) & qualquer desequilíbrio entre os mostradores \\
\textbf{mórfico} & o \textbf{fechamento} $\top$ (o universo cheio) & a abertura que o boleado $\gamma$ apaga \\
\textbf{cristalino} & trivial: $N(z)=1$, a órbita já fechou & (não vaza --- mas também não cresce) \\
\textbf{conforme} & $X(x)^2=0$: \textbf{toda obra entra no cone} & obra fora do cone (o inversor a recoloca) \\
\textbf{relógio} & onde casa, \textbf{o tempo não corre} ($N=0$, $|s|=1$) & --- \\
\textbf{criativo} & o \textbf{resíduo nulo} $x\oplus x=0$ (o cristal) & --- \\
\bottomrule\end{longtable}}

\textbf{O inversor multinível (Pontryagin) é o item universal de conserto:} ele recoloca no cone toda energia que vazaria --- os degraus de Fibonacci $\varphi^{-j}$ da bateria áurea casam cada harmônico. \emph{Ter o inversor é ter licença para jogar em qualquer corpo.}



\chapter{A abertura de universo, como jogada}


\begin{enumerate}
\item O jogador \textbf{declara} o corpo (a sua $\Sigma$, a sua régua $\mu$, o seu flip $\nu$, o seu mate).
\item O árbitro roda o \textbf{crivo}: $\mathrm{res}(\text{corpo}) = $ n.º de cláusulas violadas.
\item $\mathrm{res}=0 \Rightarrow$ \textbf{ACEITO} (o universo abre). $\mathrm{res}\neq0 \Rightarrow$ \textbf{REJEITADO}. (E ninguém assina sozinho: assina-se \textbf{com o dual} --- ver 2.4.)
\item Certificado: o mapa do contrato (6/6) e as realizações dos corpos (5/5).
\end{enumerate}



\chapter{Os corpos que não fecham SOZINHOS --- e o dual que os fecha}


\textbf{A Lei é geral: não há exceção, não há corpo condenado.} Mas há corpos que, \textbf{tomados sozinhos}, não têm um axioma nomeável. E aqui está a chave que o reino inteiro depende:

\begin{tcolorbox}[colback=ouroclaro!35,colframe=ouro,boxrule=0pt,leftrule=2pt,arc=0pt,left=4mm,right=3mm,top=2mm,bottom=2mm,breakable]\itshape \textbf{Não tem sozinho. Com o seu DUAL, tem.}\end{tcolorbox}

Nenhum corpo joga sozinho. \textbf{Todo corpo comparece ao Contrato com o seu dual}, e é o \textbf{par} que assina --- porque cada um doa ao outro exatamente o que lhe falta. A conservação não vive \emph{dentro} de um corpo: \textbf{vive sobre o par}.

{\fontsize{10.00pt}{12.00pt}\selectfont\begin{longtable}{>{\fontsize{10.00pt}{12.00pt}\selectfont\RaggedRight\hspace{0pt}\arraybackslash}p{5378458sp}>{\fontsize{10.00pt}{12.00pt}\selectfont\RaggedRight\hspace{0pt}\arraybackslash}p{7801655sp}>{\fontsize{10.00pt}{12.00pt}\selectfont\RaggedRight\hspace{0pt}\arraybackslash}p{7108229sp}>{\fontsize{10.00pt}{12.00pt}\selectfont\RaggedRight\hspace{0pt}\arraybackslash}p{7822342sp}}\toprule
\textbf{o corpo} & \textbf{o que lhe falta \textbf{sozinho}} & \textbf{o seu \textbf{DUAL}} & \textbf{e então o par \textbf{fecha}} \\ \midrule\endhead
\textbf{entrópico} (semicorpo) & o \textbf{inverso aditivo} --- o \code{max} é idempotente (\code{a(+)a=a}) & \textbf{CÓSMICO} --- a expansão \code{a(t)=e\textasciicircum{}\{Ht\}}, o \code{exp} & sob o \code{exp}, o \code{max} \textbf{vira soma} --- e a soma \textbf{inverte}. O inverso existe: mora no par \\
\textbf{telescópico} ($n\geq2$) & o \textbf{produto} cinde (divisores de zero, $e_i\cdot e_j=0$) & \textbf{CELESTE} & o que cinde de um lado, o outro recompõe \\
\textbf{conforme} & a norma não distingue (cone nulo) & o \textbf{par} (a medida é a distância ao outro) & é a casa do árbitro: mede \textbf{pelo par}, e por isso mede \\
\bottomrule\end{longtable}}

\textbf{A consequência jogável --- e é o oposto do que parecia.} No entrópico, de dentro, o dano só acumula e nada se desfaz. Mas a saída \textbf{existe, e está do outro lado}: quem sabe \textbf{atravessar ao dual} (o \code{exp}, a expansão) recupera o inverso e \textbf{desfaz}. O terreno não é uma armadilha: é \textbf{metade de um par}. Quem entra sem saber disso, morre lá. Quem entra sabendo, atravessa.

E há um detalhe que decide tudo: a ausência de inverso é uma propriedade do \textbf{limite \code{T -> 0}} (o \code{max} duro, o zero absoluto). A qualquer \code{T > 0} --- \textbf{em qualquer lugar onde alguém viva} --- o $\oplus$ é o \emph{logsumexp}, que destropicaliza sob o \code{exp} e \textbf{inverte}. \textbf{Todos os jogadores são entrópicos}, e isso não é sentença: é o mundo.

\emph{Certificado} (\textbf{5/5, resíduo 0}) --- o refutador varre 4001 candidatos e confirma que, sozinho, o \code{max} não admite inverso; sob o \code{exp} do cósmico, admite. E a prova construtiva: \textbf{do pó ao corpo, e de volta} --- a composição elementar de um corpo humano, montada e revertida, resíduo \textbf{exatamente 0} (em $\mathbb{Q}$). \textbf{A construção entrópica é reversível.}

\vspace{2mm}\hrule\vspace{3mm}



\chapter{Os três verbos (a tríade)}


{\fontsize{10.00pt}{12.00pt}\selectfont\begin{longtable}{>{\fontsize{10.00pt}{12.00pt}\selectfont\RaggedRight\hspace{0pt}\arraybackslash}p{3022510sp}>{\fontsize{10.00pt}{12.00pt}\selectfont\RaggedRight\hspace{0pt}\arraybackslash}p{6044714sp}>{\fontsize{10.00pt}{12.00pt}\selectfont\RaggedRight\hspace{0pt}\arraybackslash}p{6751293sp}>{\fontsize{10.00pt}{12.00pt}\selectfont\RaggedRight\hspace{0pt}\arraybackslash}p{6445793sp}>{\fontsize{10.00pt}{12.00pt}\selectfont\RaggedRight\hspace{0pt}\arraybackslash}p{5322086sp}}\toprule
\textbf{Símbolo} & \textbf{Nome} & \textbf{O que é} & \textbf{Na ISA ERG-64} & \textbf{o signo} \\ \midrule\endhead
\textbf{$\oplus$} & \textbf{CLIFFORD} & a SOMA, o PAR, a corrente que se junta no nó & \code{ADD} (ou \code{XOR} no corpo booleano $\mathbb F_2$) & Secundidade / ÍNDICE \\
\textbf{$\otimes$} & \textbf{LA HIRE} & o PRODUTO, a Rainha, a tensão/o ganho & a MULT por \textbf{double-and-add} (\code{ADD} deslocado); \code{AND} no booleano & a mediação \\
\textbf{$\prod$} & \textbf{PONTRYAGIN} & o OPERADOR que costura, $\exp\circ\Sigma\circ\log$ & \code{NOT} / o \textbf{dispatch} / a composição & Terceiridade / SÍMBOLO \\
\bottomrule\end{longtable}}

Não há um quarto verbo. \textbf{Todo corpo do reino conjuga os mesmos três} --- do fractal ao econômico, do óptico ao entrópico.



\chapter{As doze palavras}


\code{LOAD}, \code{STORE}, \code{ADD}, \code{SUB}, \code{AND}, \code{OR}, \code{XOR}, \code{CMP}, \code{JMP}, \code{JZ}, \code{JNZ}, \code{HALT}.

Inteira (i64). \textbf{Exata} --- resíduo 0 real, \textbf{sem float}. Dela o \code{chessc} (o broca-so, o compilador universal) emite \textbf{WASM} (o metal) e \textbf{Dafny} (as leis provadas), e a costura fecha:

\[\texttt{interp} \equiv \texttt{WASM} \equiv \texttt{Dafny}\]

O gabarito --- a tríade inteira do corpo criativo escrita nas doze palavras. O \textbf{slot 9} carrega o
opcode (\code{0}$=\oplus$/\code{XOR}, \code{1}$=\otimes$/\code{AND}, \code{2}$=\prod$/\code{NOT}); os operandos
entram em \code{1} e \code{2}, o resultado sai em \code{3}; o slot \code{10} guarda o todo-um
(\code{0xFFFF\ldots FF}), que é como se faz o \code{NOT} sem ter \code{NOT}: \code{XOR} com todos os bits ligados.

\begin{Verbatim}[fontsize=\small,frame=single,framesep=3mm]
DISP:   LOAD 9        ; o opcode
        LOAD 0        ; contra o zero
        CMP
        JZ  XORO      ; op = 0  ->  (+)
        LOAD 9
        LOAD 8        ; o um
        SUB
        STORE 11
        LOAD 11
        LOAD 0
        CMP
        JZ  ANDO      ; op = 1  ->  (x)
NOTO:   LOAD 1        ; op = 2  ->  o complemento
        LOAD 10
        XOR
        STORE 3
        HALT
ANDO:   LOAD 1
        LOAD 2
        AND
        STORE 3
        HALT
XORO:   LOAD 1
        LOAD 2
        XOR
        STORE 3
        HALT
\end{Verbatim}

\noindent Repare que não há multiplicação nem divisão em lado nenhum: o desvio de três vias faz-se com uma
subtração e dois \code{CMP}. As réguas concretas dos corpos --- a do áureo ($\mathbb{Z}$[$\varphi$], dois
trilhos inteiros), a do racional ($\mathbb{Q}$, numerador e denominador) e a do motor --- são todas escritas
neste mesmo alfabeto de doze palavras.



\chapter{As três jogadas canônicas, escritas}


\begin{itemize}
\item \textbf{A CARGA DA RAINHA ($\otimes$ = double-and-add).} Amplificar um lance por um ganho não é "multiplicar": a peça executa \textbf{64 micro-lances bit a bit} --- \code{AND} (testa o bit da máscara) $\to$ \code{CMP}/\code{JZ} $\to$ \code{ADD} (acumula o parcial) $\to$ \code{ADD}+\code{ADD} (dobra máscara e parcial: o deslocamento) $\to$ \code{SUB} (decrementa o contador de 64) $\to$ \code{JMP}. O ganho é \textbf{exato}: o combate acumula sem vazar um ulp. \emph{No corpo racional ela roda \textbf{duas vezes} (uma no numerador, outra no denominador) --- a Rainha composta.}
\item \textbf{O AVANÇO ($\oplus$ = ADD).} A soma componente a componente. É o lance mais barato do reino: \textbf{o infinito soma-se com dois \code{ADD}} (no corpo áureo: um por trilho).
\item \textbf{A PROMOÇÃO ($\prod$ = o inverso).} Um peão promove quando o operador lhe concede o \textbf{inverso} --- e o inverso é o degrau que leva um \textbf{ANEL} a virar \textbf{CORPO}. Antes: unidirecional, sem desfazer. Depois: bidirecional, pode reverter (via $-\log$). \emph{Formalmente: em $\mathbb{Z}$[$\zeta$$_8$] o 2 não tem inverso; com $\prod$, $2\mapsto\tfrac12$ e $\mathbb{Q}$($\zeta$$_8$) é corpo.}
\end{itemize}



\chapter{As peças são operadores --- a partida de xadrez, certificada}


O \textbf{corpo criativo} (o anel booleano $\mathbb F_2$ das oito linguagens: C, PTX, WASM, LLVM, GLSL, Haskell, Dafny, Isabelle) é o corpo \textbf{gerador}: ele \textbf{joga a partida de xadrez inteira}, e \textbf{cada peça que move É um dos seus operadores}.

{\fontsize{10.00pt}{12.00pt}\selectfont\begin{longtable}{>{\fontsize{10.00pt}{12.00pt}\selectfont\RaggedRight\hspace{0pt}\arraybackslash}p{3228690sp}>{\fontsize{10.00pt}{12.00pt}\selectfont\RaggedRight\hspace{0pt}\arraybackslash}p{6852024sp}>{\fontsize{10.00pt}{12.00pt}\selectfont\RaggedRight\hspace{0pt}\arraybackslash}p{13019633sp}>{\fontsize{10.00pt}{12.00pt}\selectfont\RaggedRight\hspace{0pt}\arraybackslash}p{5010337sp}}\toprule
\textbf{Peça} & \textbf{O operador} & \textbf{O lance} & \textbf{Certificado} \\ \midrule\endhead
\textbf{Torre} & $\oplus$ Clifford (\code{XOR}) & as \textbf{retas} --- a diferença simétrica das casas; C calcula & $9/9$, resíduo $0$ \\
\textbf{Bispo} & $\otimes$ La Hire (\code{AND}, $\wedge$) & as \textbf{diagonais} --- a interseção; Dafny calcula & idem \\
\textbf{Dama} & C $\oplus$ Dafny, grau (1,1) & retas $\oplus$ diagonais na mesma peça ($Q = R + B$) & idem \\
\textbf{Cavalo} & 2$\cdot$\allowbreak{}C $\oplus$ 1$\cdot$\allowbreak{}Dafny, grau (2,1) & Pontryagin \textbf{graduado} --- o salto & idem \\
\textbf{Rei} & C $\oplus$ Dafny de um passo, \textbf{dual} das peças & elas refletem o seu passo ($\nu: e\mapsto-e$), mantendo o gap $\pm1$ & idem \\
\textbf{Peão} & o bit $\oplus$ com o \textbf{carry} $\wedge$ & a soma com transporte: $x+y=(x\oplus y)+2(x\wedge y)$ & idem \\
\bottomrule\end{longtable}}

\textbf{O colapso $\Psi$ é o minimax/argmax} (a seleção de Pontryagin). \textbf{O mate é a cristalização}: a expansão entrópica (a mobilidade) colapsa a 0 --- resíduo 0, a transferência ao dual (nada se dissipa).

\textbf{Resultado concreto e certificado:} partida completa da posição inicial ao mate --- \emph{Evans Gambit} 1.e4 e5 2.Nf3 Nc6 3.Bc4 Bc5 4.b4 \ldots 12.Qxf7\# --- \textbf{xeque-mate em 23 lances}, a Dama (C$\oplus$\allowbreak{}Dafny) cristalizando; a \textbf{Ópera de Morphy} reproduzida em \textbf{0,6 s}.

\emph{Isto amarra a "linguagem do xadrez" do enredo à ISA: o xadrez não é metáfora do jogo --- é uma das suas realizações, certificada.}

\vspace{2mm}\hrule\vspace{3mm}



\chapter{A tabela-mestra}


{\fontsize{5.56pt}{6.67pt}\selectfont\begin{longtable}{>{\fontsize{5.56pt}{6.67pt}\selectfont\RaggedRight\hspace{0pt}\arraybackslash}p{1532327sp}>{\fontsize{5.56pt}{6.67pt}\selectfont\RaggedRight\hspace{0pt}\arraybackslash}p{2988319sp}>{\fontsize{5.56pt}{6.67pt}\selectfont\RaggedRight\hspace{0pt}\arraybackslash}p{3685257sp}>{\fontsize{5.56pt}{6.67pt}\selectfont\RaggedRight\hspace{0pt}\arraybackslash}p{6229167sp}>{\fontsize{5.56pt}{6.67pt}\selectfont\RaggedRight\hspace{0pt}\arraybackslash}p{5053298sp}>{\fontsize{5.56pt}{6.67pt}\selectfont\RaggedRight\hspace{0pt}\arraybackslash}p{3526922sp}>{\fontsize{5.56pt}{6.67pt}\selectfont\RaggedRight\hspace{0pt}\arraybackslash}p{3522530sp}}\toprule
\textbf{\#} & \textbf{Corpo} & \textbf{$\Sigma$ (relógio)} & \textbf{Régua $\mu$} & \textbf{Flip $\nu$} & \textbf{Mate} & \textbf{Signo} \\ \midrule\endhead
1 & \textbf{fractal} $\mathbb{Q}$($\zeta$$_8$) --- \emph{a raiz} & $\mathbb{Z}$/8 (giro de 45$^\circ$) & $z\bar z$ (a norma) & conjugação $z\mapsto\bar z$ & o corpo fecha ($\forall x\neq0$ invertível) & \textbf{Ícone} (auto-similar) \\
2 & \textbf{áureo} $\mathbb{Z}$[$\varphi$] --- \emph{o infinito construído} & órbita do gato $A_1$ (Pisano) & $\sum\varphi^{-j}=\varphi^2=(1,1)$ & $\varphi\mapsto-1/\varphi$ & a \textbf{dobra}: tudo reduzido a $(F_{n-1},F_n)$ & Símbolo \\
3 & \textbf{racional} $\mathbb{Q}$ --- \emph{o exato} & a palavra i64 (64 tiques) & Euclides / fração contínua & $[(a,b)]^{-1}=[(b,a)]$ & o \textbf{resto 0} de Euclides & Índice (o par) \\
4 & \textbf{criativo} $\mathbb F_2$ --- \emph{as 8 linguagens} & as classes de -- & invariância dos 8 substratos; $L\triangle L=\emptyset$ & gerar$\leftrightarrow$\allowbreak{}atestar / \code{NOT} & a \textbf{cristalização} (mobilidade$\to$\allowbreak{}0) & Símbolo \\
5 & \textbf{eletromagnético} & a órbita do flip & $\sigma=|E|/|B|$ (o telescópio $\Sigma h$) & o flip $(E,B)\mapsto(B,-E)$ & o \textbf{cone nulo} $\sigma=1$ & Índice \\
6 & \textbf{motor} & ciclos (a gaiola de esquilo) & a magnitude / dissipação & o flip $a\mapsto ia$ (motor$\leftrightarrow$\allowbreak{}rotor) & a ressonância casa$\leftrightarrow$\allowbreak{}peça & Índice \\
7 & \textbf{relógio universal} & $\mathbb{Z}$/n (a roda-mãe) & o lapso $N(\psi)=\cos\psi$ & $\Lambda=\mathrm{artanh}$ & $N=0$: \textbf{o tempo para} & Símbolo \\
8 & \textbf{telescópico} & a deflexão $D_\lambda$ & $(N,T)$: covol$^2$ e $\sum h_i$ & $(-1)\circ\mathrm{rev}$ & a régua unimodular $T=0$ & Símbolo \\
9 & \textbf{cristalino} $\mathbb C_D$ & \textbf{parado} (regulador 0) & $N(z)=a^2+|D|b^2$ & conjugação de Galois & a órbita já fechou & \textbf{Ícone} \\
10 & \textbf{conforme} & a involução (ordem 2) & o \textbf{par} $X\cdot Y=-\tfrac12|x-y|^2$ & swap $n_0\leftrightarrow n_\infty$ & o casamento & Índice \\
11 & \textbf{entrópico} (max,+) & o regulador $h=\log\sigma_m$ & o custo mínimo (o piso) & $a\mapsto-a$ & o fechamento degenerado & Índice \\
12 & \textbf{espaço-temporal} $\mathbb{Q}$$_2$ & a órbita do sucessor & ultramétrica $|x|_2=2^{-v}$ & reflexão $-1$ (o flip) & \textbf{transporte $=0$} & Símbolo \\
13 & \textbf{óptico} & a régua redonda ($k=-1$) & $n\sin\theta$ & reflexão especular & o casamento na interface & Índice \\
14 & \textbf{celeste} (2 qubits) & as rotações & $r^2+C^2=1$ & a reflexão $-1$ (que \textbf{gera} $i$) & a medida sorteia no emaranhado & Índice \\
15 & \textbf{econômico} (os juros) & os períodos $t$ & a força $\delta=\log(1+r)$ & capitalizar$\leftrightarrow$\allowbreak{}descontar & $VP(A(t),t)=P$ & Símbolo \\
16 & \textbf{evolutivo} & $\sum p_i=1$ & a frequência $p$, norma $1-s^2$ & $p\mapsto1-p$ & fixação ($p{=}1$) / extinção ($p{=}0$) & Símbolo \\
17 & \textbf{expansivo} $E_m=\mathbb Q(\sigma_m)$ & $\mathbb R/(\log\sigma_m)\mathbb Z$ & $\log|x| \bmod \log\sigma_m$ & Galois $\bar\sigma=-1/\sigma$ & ser \textbf{unidade} $N=\pm1$ & Símbolo \\
18 & \textbf{somático} (o tecido) & os $T$ tipos ($\pi = x \bmod T$) & $\det(\exp tG)=e^{t\,\mathrm{tr}G}$ & proliferar$\leftrightarrow$\allowbreak{}diferenciar ($G^\top$) & a \textbf{homeostase} & Símbolo \\
19 & \textbf{geométrico} (a curvatura) & a geodésica & a métrica $g_{\mu\nu}$; $\Phi=-\tfrac12\log r$ & log$\leftrightarrow$\allowbreak{}exp & o \textbf{vácuo} $R_{\mu\nu}=0$ & Símbolo \\
20 & \textbf{mórfico} ($\mathbb F_2^n$) & $\mathbb{Z}$/2 ($\ell=2$) & o \textbf{resíduo} & $\neg = D_1 = x\oplus1$ & o fechamento $\top$ & Símbolo \\
\bottomrule\end{longtable}}



\chapter{As jogadas próprias (o combate por corpos)}


Cada corpo entrega \textbf{uma mecânica que não existe em nenhum outro}. Isto é o coração do combate: \emph{trocar de mundo é trocar de jogo.}

\textbf{ FRACTAL $\mathbb{Q}$($\zeta$$_8$) --- a raiz, o tabuleiro auto-similar.} A tesselação é \textbf{nativa} (zoom = mesmo tabuleiro). Cada peça carrega um \textbf{RUMO}: uma das 8 unidades $\zeta_8^k$ (o relógio $\mathbb{Z}$/8, 8 direções de 45$^\circ$).
\begin{itemize}
\item \textbf{Giro} = $\otimes$: multiplicar o rumo por $\zeta_8$ gira a peça 45$^\circ$.
\item \textbf{Avanço} = $\oplus$: somar deslocamento, coordenada a coordenada.
\item \textbf{Promoção} = $\prod$: concede o \textbf{inverso} --- o peão sai do ANEL e entra no CORPO.
\item \textbf{Espelho} = $\nu$ (conjugação): reflete o rumo do inimigo; e $z\otimes\bar z=|z|^2$ é a \textbf{régua de dano/alcance}.
\item \textbf{Conservação:} só rumos de norma $|z|=1$ são lances legais. Mover com $|z|\neq1$ \textbf{vaza}.
\item \textbf{Atrator de Newton:} peça lançada converge a uma das 8 raízes (\emph{homing}); se cai na \textbf{fronteira fractal}, nunca assenta --- \textbf{casa contestada}.
\end{itemize}

\textbf{ ÁUREO $\mathbb{Z}$[$\varphi$] --- a DOBRA.} Uma peça carrega uma cadeia \textbf{infinita} de golpes que decaem ($\varphi^{-1},\varphi^{-2},\dots$ --- o eco que nunca acaba). Num motor comum, o eco teria de ser \textbf{truncado} (dano vaza, a partida deixa de ser reversível). Aqui o Rei \textbf{dobra} a cadeia: $\varphi^2=\varphi+1$ fecha o eco no par $(1,1)$ --- \textbf{o golpe infinito resolve num único tick, dano exato $\varphi^2$, zero arredondamento.}
\begin{itemize}
\item \textbf{Movimento:} $(a,b)\mapsto(b,a+b)$ --- a escada de Fibonacci. É o único lance que \textbf{acelera sem gastar recurso} (a órbita é \emph{gerada}, não somada).
\item \textbf{Escudo:} o conjugado $\varphi'=-1/\varphi$ --- $(a,b)\mapsto(a+b,-b)$: encolhe o golpe por $\varphi^{-1}$ a cada camada. Expande $\bowtie$ contrai, $\sigma\sigma'=-1$; a soma de tudo é limitada por $\varphi^2$.
\item \textbf{Mate:} toda potência do adversário reduzida a $(F_{n-1},F_n)$ --- não há para onde expandir. \textbf{O infinito não é derrotado: é dobrado.}
\end{itemize}

\textbf{ RACIONAL $\mathbb{Q}$ --- o contra-golpe pela recíproca.} Toda peça se move por uma \textbf{razão} $(num, den)$: o numerador é o \textbf{avanço}, o denominador é a \textbf{marcha/tempo}. Casas por tick $=$ num/den.
\begin{itemize}
\item \textbf{Defesa perfeita:} um ataque de razão $3/4$ é anulado \textbf{exatamente} pela recíproca $4/3$ --- primal $\otimes$ dual $=1$. \emph{Em $\mathbb{Z}$ não havia defesa; em $\mathbb{Q}$, todo ataque não-nulo TEM defesa exata.}
\item \textbf{O duelo de Euclides (o mate):} dois lances se medem um no outro, subtraindo o menor do maior. Quem chega ao \textbf{resto 0} dá mate. Entre racionais o duelo \textbf{sempre termina}. Contra um \textbf{METAL} (irracional, fração contínua infinita: $\varphi$, prata, bronze) o duelo \textbf{nunca fecha} --- \textbf{xeque perpétuo}. \emph{E é por esse buraco que $\mathbb{Q}$ se completa em $\mathbb{R}$.}
\item \textbf{A lei da limpeza:} $2/4$ e $1/2$ são a \textbf{mesma casa}. Inflar as duas camadas ($\times k$) não gera vantagem: \textbf{não há energia de graça}.
\end{itemize}

\textbf{ ELETROMAGNÉTICO --- o golpe casado.} Cada peça/onda carrega $(E,B)$ e uma impedância $\sigma=|E|/|B|$ (o seu \textbf{metal}). O \textbf{ataque máximo} é levar o alvo ao \textbf{cone nulo $\sigma=1$}: transferência $1-|\Gamma|^2=1$, o golpe que não vaza. Um lance \textbf{descasado} deixa potência \textbf{reativa} presa --- dano perdido. Contra-lance: o \textbf{flip} $(E,B)\mapsto(B,-E)$ --- troca o eixo do ataque sem perder energia.

\textbf{ ÓPTICO --- as fronteiras e o escudo perfeito.} Cada território tem um \textbf{índice $n$\textbf{ (a densidade do meio). Atravessar a fronteira }dobra\textbf{ a trajetória --- e o jogador }não escolhe} o ângulo de saída: $n_1\sin\theta_1=n_2\sin\theta_2$. \textbf{Reflexão total interna:} se o ataque passa o \textbf{ângulo crítico} $\theta_c=\arcsin(n_2/n_1)$ (= 41,81$^\circ$ para 1,5$\to$\allowbreak{}1), é \textbf{100\% devolvido} --- o escudo perfeito. Se o território é um \textbf{cristal}, a fronteira \textbf{difrata}: o ataque se parte em ordens discretas --- múltiplos golpes em ângulos fixos.

\textbf{ CRISTALINO --- a armadura que não vaza (e não cresce).} A peça-cristal \textbf{congela uma região} num reticulado de regulador nulo: o relógio para, e só cabem giros de ordem $\{1,2,3,4,6\}$ (a armadura quadrada $\mathbb{Z}$[i] ou hexagonal $\mathbb{Z}$[$\omega$]). \textbf{O golpe-chave: o ataque ÁUREO da Rainha (a dobra-5, o traço $\varphi$) é REPELIDO} --- é a difração proibida; \emph{nenhum cristal admite simetria pentagonal.} Para quebrar o cristal, o adversário mira um feixe e lê os \textbf{picos do recíproco}: a ordem zero atravessa sem dano; os desvios são o dano. \textbf{Mate suave:} no \textbf{quasicristal} o gap fecha --- os picos acumulam, a defesa perde o isolamento. \emph{A tese, jogável: \textbf{fechar exclui preencher}.}

\textbf{ ENTRÓPICO --- a zona da irreversibilidade.} Terreno onde o campo \textbf{deixa de fechar}. Cada lance $\otimes$$=+$ \textbf{soma o custo} ao caminho; $\oplus$$=\max$ escolhe a linha dominante. \textbf{Mas sem inverso aditivo NÃO SE DESFAZ NADA} --- sem cura, sem cancelamento, sem interferência. Em cada vértice de coordenação $z$ há $e^{Ns}$ estados de igual resíduo: a peça \textbf{tem de sortear} a representação (\textbf{um dado obrigatório a cada nó}). A vitória é o caminho de \textbf{melhor peso médio de ciclo}. \emph{E não há caos: o gato tropical é não-expansivo --- a rotação morre porque exige o sinal, e não há sinal.}

\textbf{ TELESCÓPICO --- o dano por eixos.} O ataque encadeia dano por $N$ eixos ortogonais, e o total é a \textbf{norma que colapsa na SOMA das entropias} $\sum h_i$ --- atacar por muitos eixos é \textbf{somar perfis, não multiplicar custo}. \textbf{A régua blindada} (invulnerável) é a \textbf{autodual}, $T=0$: o que um eixo estica, outro contrai. \textbf{O golpe-fantasma:} bater num eixo \textbf{ortogonal} ao do inimigo dá \textbf{dano ZERO} --- a jogada some, não troca (o preço da cisão).

\textbf{ CELESTE --- emaranhar custa dano.} Cada peça reparte uma unidade de energia entre $r$ (potência \textbf{ativa}, o dano direto) e $C$ (a \textbf{reativa}, o emaranhamento), com $r^2+C^2=1$ sempre. \textbf{Enlaçar duas peças} acopla-as ($C>0$) e por isso \textbf{BAIXA a potência individual} --- ganhas correlação, perdes dano solo. \textbf{O emaranhamento puro} ($r=0$, $C=1$) --- a \textbf{moeda} que o mate (a medida) sorteia; o par colapsa em cadeia.

\textbf{ MOTOR --- o torque de flanco.} O torque $\tau=\psi\times i$ é \textbf{máximo quando fluxo e corrente estão ORTOGONAIS}: o golpe de flanco a 90$^\circ$, não frontal. O \textbf{DTC} (Pontryagin) permite \textbf{modular} o torque em jogo --- o co-estado ($G^\top$, o valor futuro) guia quanto empurrar agora.

\textbf{ RELÓGIO UNIVERSAL --- a saturação do combo.} No combate as velocidades combinam pela \textbf{$\oplus$ relativística}: $s_1\oplus s_2=(s_1+s_2)/(1+s_1s_2)$. \textbf{Nenhum combo cruza o horizonte $|s|=1$\textbf{ --- o dano satura mas }nunca quebra a régua\textbf{. O único tique }sem inverso\textbf{ é o }MATE\textbf{: o relógio }para} ($N=0$).

\textbf{ ESPAÇO-TEMPORAL $\mathbb{Q}$$_2$ --- os dois regimes.} O tabuleiro tem dois modos que são o flip um do outro: \textbf{redondo/compacto = ESPAÇO} (o lance é a \textbf{rotação}) e \textbf{hiperbólico = TEMPO} (o lance é o \textbf{boost}). O flip troca o tipo de tabuleiro. O lance de soma resolve-se como o \textbf{adder}: propaga o transporte até ele ser \textbf{0} --- \emph{o fechamento É transporte nulo}. \textbf{Armadilha da máquina finita:} truncar fabrica torção e raízes espúrias --- a \textbf{queda de bandeira precoce}.

\textbf{ MÓRFICO --- as zonas de controle.} Cada casa é um bit (dominada/livre); cada exército, um conjunto de casas. \textbf{Erosão} ($\wedge$): só sobrevivem as casas cuja vizinhança inteira já dominas --- consolida o núcleo. \textbf{Dilatação}: espalha a zona às casas adjacentes --- o avanço. \textbf{Antiroque} ($\oplus$/XOR): a \textbf{diferença simétrica} --- as casas contestadas por exatamente um exército, o lance da fronteira. \textbf{Complementação} ($\neg$): vira o tabuleiro inteiro, ocupado$\leftrightarrow$\allowbreak{}vazio (involução: dois lances desfazem).

\textbf{ SOMÁTICO --- o tabuleiro-exército.} \textbf{Proliferar:} o tecido dobra célula a célula ($2^k$) até \textbf{colidir} com a capacidade --- a \textbf{homeostase} (exatamente como os peões colidem e travam). \textbf{Diferenciar:} projeta as $N$ células nos $T$ \textbf{tipos} --- você não move célula a célula ($O(N)$), move \textbf{por tipo} ($O(T)$): o atalho do grafo. \textbf{Fundir:} dois exércitos = a soma direta $G\oplus G_2$. Quem tem $\mathrm{tr}\,G>0$ \textbf{transborda} (FP$<1$) até ser corrigido.

\textbf{ EVOLUTIVO --- a meta do reino.} Cada tipo de peça tem uma \textbf{frequência $p$} no exército. A cada turno, as peças que capturaram sobem: $p_i\mapsto p_i w_i/\langle w\rangle$, e $\sum p_i=1$ conserva. \textbf{A aptidão média $\langle w\rangle$ só cresce\textbf{ (Fisher) --- o combate é uma escalada monótona de qualidade. }Empilhar vantagens rende cada vez menos} perto do horizonte da fixação (não há vitória grátis nos extremos). O flip $p\mapsto1-p$ é a \textbf{deriva} --- o reset ao equilíbrio $p=1/2$.

\textbf{ EXPANSIVO --- só pontua quem vira unidade.} Cada lance aplica o gato $A_m$: a peça \textbf{expande por $\sigma_m$\textbf{ a cada turno. Mas o }gate de pontuação\textbf{: só marca quem aterrissa numa }UNIDADE} $N(x)=\pm1$ --- resíduo 0. Toda posição com $N\neq\pm1$ \textbf{vaza}, e o inversor de Pontryagin (a torre de Fibonacci) a devolve ao círculo, \textbf{um passo por turno}. \emph{A vitória não é matar: é fechar a órbita --- a peça vira unidade.}

\textbf{ GEOMÉTRICO --- a curvatura do tabuleiro.} Pousar uma peça de massa $M$ \textbf{cava um poço} $\Phi=-\tfrac12\log(1+2M/|x|)$, e as outras peças passam a seguir \textbf{geodésicas} --- curvam ao redor. O combate corre no \textbf{cone nulo} ($ds^2=0$): o disparo limpo que chega. O poço bem-formado é o \textbf{vácuo} $R_{\mu\nu}=0$. \textbf{Um poço mal-cavado tem fonte espúria e vaza} --- é assim que se distingue um buraco negro legítimo de um vazamento proibido.

\textbf{ ECONÔMICO --- o tempo é o tabuleiro.} Cada turno o jogador escolhe: \textbf{$\oplus$ (juro simples)} --- ganho fixo, uma reta segura, nunca vaza; ou \textbf{$\otimes$ (juro composto)} --- ganho que rende sobre si mesmo: \textbf{a partir do 3.º turno o composto DOMINA o simples} (convexidade da exp). O contra-lance é o \textbf{DESCONTO} ($\nu$): trazer valor do futuro ao presente --- \textbf{descontar desfaz capitalizar exatamente}. Quem deixa FP$<1$ abre um \textbf{SPREAD} --- atrito que sangra valor a cada turno, e o adversário explora o vazamento.

\textbf{ CONFORME --- o árbitro e o relógio de xadrez.} \emph{(Ver Cap. 11.)}

\textbf{ CRIATIVO --- o corpo que gera.} \emph{(Ver Cap. 5.)}



\chapter{O árbitro --- o corpo conforme e o refutador com dentes}


O \textbf{corpo conforme} é a \textbf{casa da medida}: onde o certificador vive. Não é campo de batalha --- é o \textbf{relógio de xadrez} do reino.

\begin{itemize}
\item Toda afirmação/lance carrega \textbf{dois mostradores}: $b=(\text{testa},\text{refuta})\in\{0,1\}^2$.
\item O árbitro lê o \textbf{PONTEIRO}: $\mathrm{pt}(b)=\langle X(b),X(b_\top)\rangle = -\tfrac12\,d(b,b_\top)^2$.
\end{itemize}

{\fontsize{10.00pt}{12.00pt}\selectfont\begin{longtable}{>{\fontsize{10.00pt}{12.00pt}\selectfont\RaggedRight\hspace{0pt}\arraybackslash}p{6225672sp}>{\fontsize{10.00pt}{12.00pt}\selectfont\RaggedRight\hspace{0pt}\arraybackslash}p{6375402sp}>{\fontsize{10.00pt}{12.00pt}\selectfont\RaggedRight\hspace{0pt}\arraybackslash}p{16033898sp}}\toprule
\textbf{Ponteiro} & \textbf{Mostradores} & \textbf{Veredito} \\ \midrule\endhead
\textbf{$-\tfrac12$} & (1,0) & o lance \textbf{TEM DENTES} --- é uma jogada real \\
\textbf{$0$} & (1,1) & \textbf{QUEDA DE BANDEIRA} --- o lance é a trivialidade, o \emph{mate que não mata} vazio \\
\textbf{$-1$} & (0,0) & \textbf{falha dupla} \\
\bottomrule\end{longtable}}

\emph{"O mostrador diz o LADO (qual falhou); o ponteiro diz QUANTO."}

\textbf{A régua conforme mede pelo PAR, não pela norma:} toda obra tem norma \textbf{nula} ($X(x)^2=0$) --- a norma não distingue nada. A medida passa ao \textbf{emparelhamento} $X\cdot Y=-\tfrac12|x-y|^2$, que zera \textbf{se e só se} $x=y$.

\textbf{O combate em massa (P2P) sai daqui.} Cada célula-jogador é uma \emph{thread} da GPU, uma obra de \textbf{norma nula que não carrega carga própria}; só olha a \textbf{distância ao vizinho}. A troca é \textbf{estéril e discreta --- nada contamina}. É a fundação formal da "partida distribuída sem sala de espera" que o GDD já promete (6/6, no GLSL).

\textbf{As Princesas são o corpo técnico.} Elas não geram: \textbf{atestam}. \textbf{Dafny} decide o decidível; \textbf{Isabelle} induz sobre o infinito (o kernel LCF). O par gerar $\bowtie$ atestar é o flip $\nu$ do corpo criativo --- \textbf{nada entra no reino sem ser gerado E provado}.

\vspace{2mm}\hrule\vspace{3mm}



\chapter{A semiótica  --- ler o tabuleiro sem abrir a álgebra}


\begin{tcolorbox}[colback=ouroclaro!35,colframe=ouro,boxrule=0pt,leftrule=2pt,arc=0pt,left=4mm,right=3mm,top=2mm,bottom=2mm,breakable]\itshape \emph{"O tabuleiro não te mostra o corpo. Mostra o \textbf{signo} do corpo. Quem sabe ler, joga; quem não sabe, chuta."}\end{tcolorbox}

Problema prático do jogador: \emph{como sei o que este corpo faz sem abrir a sua álgebra?} Resposta: \textbf{pelo signo}. Há exatamente três, nem um a mais.

{\fontsize{10.00pt}{12.00pt}\selectfont\begin{longtable}{>{\fontsize{10.00pt}{12.00pt}\selectfont\RaggedRight\hspace{0pt}\arraybackslash}p{4988220sp}>{\fontsize{10.00pt}{12.00pt}\selectfont\RaggedRight\hspace{0pt}\arraybackslash}p{5368409sp}>{\fontsize{10.00pt}{12.00pt}\selectfont\RaggedRight\hspace{0pt}\arraybackslash}p{4408522sp}>{\fontsize{10.00pt}{12.00pt}\selectfont\RaggedRight\hspace{0pt}\arraybackslash}p{6833823sp}>{\fontsize{10.00pt}{12.00pt}\selectfont\RaggedRight\hspace{0pt}\arraybackslash}p{5987422sp}}\toprule
\textbf{Categoria} & \textbf{Signo} & \textbf{O que é} & \textbf{Como aparece no HUD} & \textbf{Como se joga contra} \\ \midrule\endhead
\textbf{Primeiridade} & \textbf{ÍCONE} (semelhança) & a qualidade pura, sozinha & um brilho que \textbf{se repete em toda escala} (fractal); uma rede \textbf{congelada} (cristalino) & não reage --- \textbf{atravessa-se}, não se combate \\
\textbf{Secundidade} & \textbf{ÍNDICE} (conexão factual) & a reação, o par, o embate & um \textbf{ponteiro} que marca $-\tfrac12$; um \textbf{par} que se cancela; uma \textbf{fronteira} & o corpo \textbf{devolve o golpe} --- ataque de flanco, casamento \\
\textbf{Terceiridade} & \textbf{SÍMBOLO} (convenção) & a lei, a mediação, o hábito & uma \textbf{lei} que fecha uma cadeia infinita em dois inteiros & a lei \textbf{é} a régua --- \textbf{use a lei a seu favor} (composto, seleção, deflexão) \\
\bottomrule\end{longtable}}

\textbf{Mapeamento na tríade:} $\oplus$ o par = \textbf{índice/secundidade}; $\otimes$ a mediação; $\prod$ o operador/a lei = \textbf{símbolo/terceiridade}. A \textbf{qualidade} do corpo é o \textbf{ícone/primeiridade}.

\textbf{Regra de mesa:} a \textbf{coluna "Signo"} da tabela-mestra (\S9.1) é uma \textbf{informação pública} --- o HUD mostra o signo do corpo aberto. Ícone = terreno; Índice = duelo; Símbolo = economia. \emph{Ler o signo é escolher o kit antes de entrar.}

\vspace{2mm}\hrule\vspace{3mm}

\section{PARTE B --- REFORÇOS NOS CAPÍTULOS EXISTENTES}



\chapter{O combate --- bate-se casado, não se bate forte}


\textbf{O corpo:} o casamento de impedância. $\Gamma=\dfrac{Z-Z_0}{Z+Z_0}$, e a potência que \textbf{entra} no alvo é

\[T \;=\; 1-|\Gamma|^2 \qquad\Longrightarrow\qquad \Gamma=0 \iff T=1 \iff \text{FP}=1 .\]

\textbf{A cena (Lance 11, o barqueiro):} \emph{"Ontem tu bateste com tudo. E entregaste zero. --- E hoje? --- Hoje bateste com metade e entregaste tudo."} A régua sai da boca do velho: \textbf{não se bate mais forte; bate-se mais casado.}

\textbf{A decisão de design.} O dano \textbf{não} é a força do input. O golpe tem uma impedância $Z$ (quem bate, com que régua, em que ângulo) e o alvo tem a sua $Z_0$ (o corpo em que ele vive, a armadura que traz). O que o alvo recebe é

\[\text{dano} \;=\; \underbrace{\big(1-|\Gamma|^2\big)}_{\text{o casamento}}\;\cdot\;\underbrace{\mu(\text{golpe})}_{\text{a régua do corpo}} .\]

Isto \textbf{fecha o buraco} entre os dois modelos que o reino tinha: a régua $\mu$ do corpo diz \emph{quanto vale} o golpe (a barra que ele move), e o casamento $\Gamma$ diz \emph{quanto passa}. Não são duas contas rivais: são o numerador e o portão.

{\fontsize{10.00pt}{12.00pt}\selectfont\begin{longtable}{>{\fontsize{10.00pt}{12.00pt}\selectfont\RaggedRight\hspace{0pt}\arraybackslash}p{11267419sp}>{\fontsize{10.00pt}{12.00pt}\selectfont\RaggedRight\hspace{0pt}\arraybackslash}p{4101420sp}>{\fontsize{10.00pt}{12.00pt}\selectfont\RaggedRight\hspace{0pt}\arraybackslash}p{13266133sp}}\toprule
\textbf{O que o jogador faz} & \textbf{$\Gamma$} & \textbf{O que acontece} \\ \midrule\endhead
bate com tudo, de chapa, fora do ângulo & $|\Gamma|\to1$ & \textbf{entrega zero}; o golpe volta inteiro \\
bate casado (a régua certa, o ângulo certo) & $\Gamma=0$ & \textbf{entrega tudo}: FP $=1$, nada reflete \\
passa do \textbf{ângulo crítico} & reflexão total & \emph{"caiu de costas a arder com o seu próprio fogo"} --- o dano volta \textbf{ao emissor} \\
\bottomrule\end{longtable}}

O botão \emph{atacar com toda a força} fica no jogo, e é uma armadilha --- a mais barata de todas: \emph{"a maior parte dos guerreiros daqui morre a bater com toda a força --- e a entregar metade."} O HUD não mostra "dano 240": mostra \textbf{$\Gamma$}, e o jogador aprende a lê-lo.

\textbf{Consequência de balanceamento:} o teto não é numérico, é geométrico. Não existe golpe que entregue mais do que $1$ (todo o telescópio propaga; nada mais há para entregar). Subir o dano bruto \textbf{não} vence: só o casamento vence. É a razão pela qual o jogo não tem inflação de números.

\vspace{2mm}\hrule\vspace{3mm}



\chapter{A fase --- dois lances certos, fora de compasso, dão zero}


\textbf{O corpo:} o trabalho médio de dois campos defasados. O telescópio só transporta potência \textbf{real} quando os dois estão em fase; fora de compasso, a média do trabalho é \textbf{zero}. No relógio-mestre do reino a intensidade do gap é

\[I(\Delta) \;=\; \Big|\cos\tfrac{\Delta}{2}\Big| ,\]

com $\Delta$ a defasagem entre os dois batimentos (é a régua que o coração do reino já usa, bit a bit).

\textbf{A cena (os dois gatos, Livro II):} \emph{"Nós estamos os dois certos. Tu estás certo. Eu estou certa. E não sai nada."}

\textbf{A decisão de design.} O combate tem \textbf{compasso}, não só alvo. Dois efeitos individualmente corretos, aplicados fora de fase, produzem \textbf{trabalho nulo} --- e o jogo diz isso na cara do jogador, com um medidor de fase, não com uma mensagem de erro.

\begin{itemize}
\item $\Delta=0$ (em fase): $I=1$ --- o par entrega tudo.
\item $\Delta=\pi$ (\textbf{antifase}): $I=0$ --- a soma \textbf{não varia}, e nada acontece. Guarde este estado: ele volta, disfarçado de vitória, no capítulo dos dois zeros.
\item $\Delta=\varphi$ (a defasagem áurea): o gap que \textbf{move} o relógio --- nem morto nem redundante.
\end{itemize}

Um combo não é uma sequência de botões: é um \textbf{casamento de fase} entre duas peças. Errar o instante não reduz o dano --- \textbf{anula-o}.

\vspace{2mm}\hrule\vspace{3mm}



\chapter{O laço --- a correlação tem preço exato}


\textbf{O corpo:} $r^2 + C^2 = 1$. O que se ganha de correlação $C$ (o vínculo) sai, ao último bit, do dano solo $r$.

\textbf{A cena (Lance 12):} \emph{"Se eu te largar, eu volto ao que era. E ficamos os dois fortíssimos. E sozinhos. E deste lado do buraco."}

\textbf{A decisão de design.} O co-op \textbf{não} é um bônus: é uma \textbf{troca de soma constante}. Ao atar duas peças, ambas ficam mensuravelmente \textbf{mais fracas sozinhas} --- e é o par atado, e só ele, que atravessa o que nenhuma das duas atravessa (o abismo, a fortaleza, o corpo que não fecha sozinho).

Design honesto: quando o jogador ata Joaquim a Yasmin, o HUD mostra \textbf{as duas barras de dano solo a descer}. À vista de todos. E o golpe conjunto que sai \emph{"não era maior. Era casado. Não vazou nada."} É a primeira coisa que a fortaleza sente.

\vspace{2mm}\hrule\vspace{3mm}



\chapter{A progressão --- uma escada de liberdades, não uma barra de XP}


\textbf{O corpo:} a cascata $\mathbb{N}\to\mathbb{Z}\to\mathbb{Q}\to\mathbb{R}$ --- cada degrau é um \textbf{inverso} que se ganha.

\textbf{A cena (Lance 19):} \emph{"A Torre não é uma pilha de mapas cada vez mais bonitos. É uma escada de liberdades, e cada degrau se paga com um inverso."}

\textbf{A decisão de design.} Não há pontos de experiência. \textbf{Sobe-se de degrau quando se ganha uma operação} --- e a operação é literalmente aquilo que a peça passa a \emph{conseguir executar}:

{\fontsize{10.00pt}{12.00pt}\selectfont\begin{longtable}{>{\fontsize{10.00pt}{12.00pt}\selectfont\RaggedRight\hspace{0pt}\arraybackslash}p{3330619sp}>{\fontsize{10.00pt}{12.00pt}\selectfont\RaggedRight\hspace{0pt}\arraybackslash}p{4336339sp}>{\fontsize{10.00pt}{12.00pt}\selectfont\RaggedRight\hspace{0pt}\arraybackslash}p{8383003sp}>{\fontsize{10.00pt}{12.00pt}\selectfont\RaggedRight\hspace{0pt}\arraybackslash}p{12060723sp}}\toprule
\textbf{Degrau} & \textbf{O corpo} & \textbf{A operação que se ganha} & \textbf{O que a peça passa a poder fazer} \\ \midrule\endhead
$\mathbb{N}$ & monoide & (nenhuma) & só somar. \emph{"o que não tem inverso é servo"} \\
$\mathbb{Z}$ & anel & o inverso \textbf{aditivo} ($-$) & \textbf{recuar}: o passado deixa de ser irreversível \\
$\mathbb{Q}$ & corpo & o inverso \textbf{multiplicativo} ($\prod$) & \textbf{anular}: um ataque de razão $3/4$ morre com $4/3$, ao último bit \\
$\mathbb{R}$ & contínuo & o limite & medir o que \textbf{não termina} \\
\bottomrule\end{longtable}}

Em $\mathbb{Q}$ o mundo deixa de ser caminho e vira \textbf{tabuleiro}: com o inverso multiplicativo, toda jogada tem resposta exata --- é aí, e não antes, que a partida fica dual.

\textbf{O gate é a falha, não o grind.} O jogador só sobe quando encontra \textbf{algo que a régua atual não mede}. A sombra não é defeito da Torre: \emph{"a sombra é o degrau."} Não há XP a farmar; há uma régua a quebrar.

\textbf{O Esquecimento (Livro II) fecha isto:} quando a corte perde tudo, ela continua \textbf{inteira} --- \emph{"não faltava um grama"}. O que se perdeu foram as \textbf{réguas}, não os gramas. Por isso o jogo não tem loot: recupera-se a \textbf{capacidade de operar}, nunca o item. (O corpo que sustenta isso é o critério de \textbf{--}: o que define o estado é a classe de congruência do que ele ainda consegue distinguir --- não o inventário que carrega.)

\vspace{2mm}\hrule\vspace{3mm}



\chapter{O empate --- o xeque perpétuo, o terceiro resultado}


\textbf{O corpo:} o algoritmo de Euclides. Contra um racional, o duelo \textbf{sempre} termina (resto zero: o mate). Contra um \textbf{metal} $\sigma_m=m+\frac1{\sigma_m}$ --- fração contínua infinita --- \textbf{nunca} termina.

\textbf{A cena (Lance 19):} \emph{"esta régua não tem mate para este inimigo. É xeque perpétuo. E não por fraqueza do braço --- por insuficiência do instrumento."}

\textbf{A decisão de design.} O jogo tem \textbf{três} resultados, não dois: mate, mate dual, e \textbf{perpétuo}. E o perpétuo não é um empate por exaustão de relógio: é um veredito \textbf{algébrico} --- a régua que o jogador trouxe não fecha neste alvo, e nenhuma quantidade de repetição vai fechá-la.

O HUD deve dizê-lo: quando o resto do duelo entra em ciclo, aparece \textbf{XEQUE PERPÉTUO} e a partida devolve o jogador ao Cartório para buscar \textbf{outra régua}. Perder por instrumento errado é uma lição; perder por dano é ruído.

\vspace{2mm}\hrule\vspace{3mm}



\chapter{O HUD --- a saúde é a variabilidade, não a regularidade}


\textbf{O corpo:} a impedância de entrada da árvore arterial como fração continuada, $Z=m+1/Z\to\sigma_m$; a augmentação $P_{\text{sis}}=P_{\text{inc}}(1+\Gamma)$. Um sistema saudável \textbf{casa} ($\Gamma\approx0$); ao endurecer, \textbf{descasa} ($\Gamma>0$) --- e a pressão que sobe \textbf{parece força}.

\textbf{A cena (Lance 21):} \emph{"Isso não é uma fortaleza. Isso é uma aorta rígida."} E: \emph{"Um coração saudável não bate certo. (\ldots) É o coração doente que bate como um metrónomo. \textbf{A arritmia mortal não é a que oscila. É a que deixa de oscilar.}"}

\textbf{A decisão de design.} O medidor central do jogo \textbf{premia a bagunça viva e alarma sobre a perfeição}. Onde todo RPG mostra \emph{estabilidade = boa}, o Reino Dourado mostra o contrário:

\begin{itemize}
\item \textbf{espectro largo / multifractal} $\to$ verde. O reino respira.
\item \textbf{metrônomo / monofractal} $\to$ \textbf{vermelho}. Está a endurecer. \emph{"O que estava a matar o Reino Dourado era a perfeição a endurecer."}
\end{itemize}

E o sinal de que o jogador está a ganhar poder \textbf{falso}: $\Gamma$ a subir enquanto a barra de pressão sobe. De dentro, parece força.

\vspace{2mm}\hrule\vspace{3mm}



\chapter{As três vias de poder, e a lei do inacabado}


\textbf{O corpo:} $s^2+c^2=1$ (a captura); $\sigma_m=m+1/\sigma_m$ e o regulador $\log\sigma_m>0$ (a mineração); $\sigma_1=\varphi$ (o lastro).

\textbf{A cena (Lance 16):} \emph{"São três. Não há quarta."}

{\fontsize{10.00pt}{12.00pt}\selectfont\begin{longtable}{>{\fontsize{10.00pt}{12.00pt}\selectfont\RaggedRight\hspace{0pt}\arraybackslash}p{5670848sp}>{\fontsize{10.00pt}{12.00pt}\selectfont\RaggedRight\hspace{0pt}\arraybackslash}p{6109683sp}>{\fontsize{10.00pt}{12.00pt}\selectfont\RaggedRight\hspace{0pt}\arraybackslash}p{3960376sp}>{\fontsize{10.00pt}{12.00pt}\selectfont\RaggedRight\hspace{0pt}\arraybackslash}p{12369777sp}}\toprule
\textbf{A via} & \textbf{O corpo} & \textbf{A conta} & \textbf{O que ela é na mesa} \\ \midrule\endhead
\textbf{BATALHAR} & rotor, $s^2+c^2=1$ & soma \textbf{zero} & a captura \textbf{é rotação}: o capturado \textbf{desce um andar} e continua a existir; o captor incrementa a massa algébrica \textbf{na quantidade exata}. \emph{"a crueldade é matematicamente inútil"} --- não é proibida; é que \textbf{não funciona} \\
\textbf{MINERAR} & metais, $\log\sigma_m$ & soma \textbf{zero} & prova de trabalho literal: \emph{"não há como falsificar um regulador"}. O juro somado é $\oplus$ (reta); o composto é $\otimes$ (domina do 3.º turno); descontar é o inverso --- o mesmo $\prod$ \\
\textbf{O OURO EXTERNO} (CristalChain) & $\sigma_1=\varphi$ & \textbf{acrescenta} & \emph{"é a única coisa que entra de fora"}. Uma vez por ciclo do CLOCK. O reino \textbf{não cria} valor: ancora-o. E é, explicitamente, \emph{"a coisa mais perigosa do reino"} \\
\bottomrule\end{longtable}}

\textbf{A sink, e o vilão que vive dela.} Quem deixa \textbf{FP $<1$\textbf{ abre um }spread\textbf{ que \emph{"sangra a cada turno"} --- e há quem viva do vazamento alheio. Isso não é sabor: é a única sink do ouro no jogo. A economia fecha porque }o que vaza vai para alguém}.

\textbf{A lei --- o inacabado vaza.} Usar a tríade \textbf{incompleta} ($\oplus$ sem $\otimes$; qualquer par sem o $\prod$ que os costura) tem penalidade dura e \textbf{assimétrica}:

\begin{tcolorbox}[colback=ouroclaro!35,colframe=ouro,boxrule=0pt,leftrule=2pt,arc=0pt,left=4mm,right=3mm,top=2mm,bottom=2mm,breakable]\itshape \textbf{metade do golpe vaza --- e o que vaza alimenta o inimigo.}\end{tcolorbox}

Foi o que aconteceu no assalto: Joaquim bateu sozinho e \emph{"metade vazou"}; Yasmin bateu sozinha e vazou \textbf{a mesma metade, pelo mesmo buraco, ao contrário}. Merlin grita a régua: \emph{"o inacabado não se sustenta num tabuleiro partido!"} Design: o \textbf{espaço é reversível} (volta-se a qualquer galáxia); o \textbf{poder é irreversível} (a armadura só cresce). O que se paga por pressa, não se recupera.

\vspace{2mm}\hrule\vspace{3mm}



\chapter{Os itens --- a armadura que não bloqueia}


\textbf{O corpo:} o perfil \textbf{autodual} --- um eixo estica exatamente o que o outro encolhe; traço zero; a soma \textbf{é} zero. E $\Gamma=0$: transmitir, não resistir.

\textbf{A cena (Lance 17):} \emph{"--- Ela não me protegeu. --- Não. --- Ela deixou passar. --- Deixou passar tudo, e não guardou nada. É diferente."}

\textbf{A decisão de design.} A armadura correta \textbf{não bloqueia: transmite}. A armadura grossa --- espessa, brilhante, cara, a que todo RPG vende --- \textbf{ressoa}, e o golpe fica \textbf{preso lá dentro}. É a inversão exata da intuição de gênero, e é a regra:

\begin{itemize}
\item \textbf{defesa = casamento} ($\Gamma\to0$: o golpe atravessa e não guarda nada);
\item \textbf{defesa por espessura = ressonância} (o golpe entra, não sai, e cozinha o portador).
\end{itemize}

\textbf{Os itens são réguas --- e as réguas não somam.} Um item nunca é $+3$ de dano: é \textbf{uma maneira de medir}. Daí, sem exceção:

\begin{itemize}
\item \textbf{não há stacking}: duas réguas não se somam (medir duas vezes não mede melhor);
\item \textbf{trocar de régua no meio de uma jogada é lance ilegal} --- o árbitro anula o turno.
\end{itemize}

O \emph{loadout} é, portanto, uma \textbf{escolha de instrumento} antes do lance, e não um somatório de bônus. É por isso que o inventário do jogo é pequeno e cada peça pesa.

\vspace{2mm}\hrule\vspace{3mm}



\chapter{O exército --- a proporção que não volta}


\textbf{O corpo:} a dinâmica de replicador, $p_i\mapsto p_i\,\dfrac{w_i}{\langle w\rangle}$, com $\sum_i p_i=1$. O total é \textbf{constante}; só a \textbf{mistura} muda. A escalada é \textbf{monótona} (o ganho médio nunca desce).

\textbf{A cena (Lance 15):} \emph{"Tinha perdido quatrocentos e tinha quatrocentos a mais de outra coisa."} E o aviso de Isabelle: \emph{"Foste tu que chamaste progresso ao que estava a acontecer. Toda a gente chama."}

\textbf{A decisão de design.} O exército é uma \textbf{proporção}, não uma contagem --- e otimizar \textbf{mata a diversidade}. A habilidade que não serviu por trinta batalhas \textbf{desaparece}, e não se reconstrói: \emph{"não havia de onde"}. É uma punição de longo prazo ao min-max, e ela é \textbf{silenciosa} (o medidor só melhora).

Dois freios saem do mesmo corpo, e ambos entram como regra de balanceamento:

\begin{itemize}
\item \textbf{O muro.} Perto do limite, cada vantagem nova vale menos que a anterior: \emph{"o esforço enorme comprava o ganho ridículo."} \textbf{Não há vitória grátis nos extremos.}
\item \textbf{A saturação do combo.} As velocidades compõem-se como $s_1\oplus s_2=\dfrac{s_1+s_2}{1+s_1s_2}$: \textbf{nenhum encadeamento cruza $|s|=1$}. É o único teto de balanceamento do sistema --- e não é uma constante escolhida a dedo por um designer: é a régua do corpo.
\item \textbf{A inibição por contato.} O tecido saudável \textbf{sabe parar de crescer}. O que não para \emph{"mata o corpo em que vive, e morre com ele --- e até ao fim vai jurar que estava só a ficar mais forte."} É o retrato mecânico de Dark, e é o que o HUD da variabilidade detecta.
\end{itemize}

\vspace{2mm}\hrule\vspace{3mm}



\chapter{O acaso --- o dado só rola onde o sinal morre}


\textbf{O corpo:} o corpo \textbf{entrópico} não tem inverso aditivo. Sem o menos, não há oscilação, não há interferência, não há pulso --- e há $e^{Ns}$ estados de \textbf{igual resíduo}, entre os quais nada distingue.

\textbf{A cena (Lance 20, a segunda porta):} \emph{"julgou que ia encontrar tempestade, e encontrou lodo."} / \emph{"rodar exige o sinal --- e no pântano não há sinal."}

\textbf{A decisão de design.} O Reino Dourado é um jogo \textbf{determinístico} --- exceto onde a matemática \textbf{obriga} o sorteio:

\begin{tcolorbox}[colback=ouroclaro!35,colframe=ouro,boxrule=0pt,leftrule=2pt,arc=0pt,left=4mm,right=3mm,top=2mm,bottom=2mm,breakable]\itshape onde há sinal, o lance é exato; \textbf{no pântano, a peça TEM de rolar o dado a cada nó} --- porque nada ali distingue um estado do outro de igual resíduo.\end{tcolorbox}

O acaso não é tempero: é um \textbf{sintoma de corpo}. Ele aparece exatamente nos mundos que perderam o menos (o entrópico) ou onde a medida sorteia (o celeste, no emaranhado). Em todo o resto do reino, dois lances iguais dão o mesmo resultado, ao bit --- e é isso que torna toda partida \textbf{auditável}.

\vspace{2mm}\hrule\vspace{3mm}



\chapter{O Cartório --- abrir um mapa é uma jogada, não uma conquista}


\textbf{O corpo:} o Contrato (as cinco cláusulas: $\Sigma$, $\mu$, $\nu$, o mate, e o crivo $\Gamma=0$).

\textbf{A cena (Lance 18):} \emph{"Não se conquista um mapa no Reino Dourado. \textbf{Assina-se.}"}

\textbf{A decisão de design.} O conteúdo do jogador (o UGC) não entra por curadoria nem por voto: entra por \textbf{rito}, e o rito é uma jogada de mesa com quatro perguntas e um crivo.

{\fontsize{10.00pt}{12.00pt}\selectfont\begin{longtable}{>{\fontsize{10.00pt}{12.00pt}\selectfont\RaggedRight\hspace{0pt}\arraybackslash}p{11106062sp}>{\fontsize{10.00pt}{12.00pt}\selectfont\RaggedRight\hspace{0pt}\arraybackslash}p{10869453sp}>{\fontsize{10.00pt}{12.00pt}\selectfont\RaggedRight\hspace{0pt}\arraybackslash}p{6659457sp}}\toprule
\textbf{A pergunta} & \textbf{O que ela pede} & \textbf{Se não responder} \\ \midrule\endhead
\textbf{Qual é o teu relógio?} ($\Sigma=\mathbb{Z}/n$) & de quantos em quantos tiques o teu mundo fecha & \emph{"um universo sem órbita é um universo sem turno"} \\
\textbf{Qual é a tua régua?} ($\mu$) & como o teu mundo mede (norma, Euclides, frequência, lapso) & não há dano, não há distância, não há nada \\
\textbf{Onde está o teu espelho?} ($\nu$) & a involução: gerar$\leftrightarrow$\allowbreak{}atestar, capitalizar$\leftrightarrow$\allowbreak{}descontar & \emph{"um corpo que não sabe olhar-se é um corpo que não sabe voltar"} \\
\textbf{Qual é o teu mate?} ($\bar\circ$) & onde o teu mundo \textbf{fecha} (unidade, resto zero, HALT, cristalização) & a partida não termina --- nem bem, nem mal \\
\bottomrule\end{longtable}}

\textbf{E o crivo, a quinta:} \emph{"o teu universo vaza?"} --- $\Gamma=0$, FP $=1$. \textbf{Se vaza, não abre.}

Duas regras de mesa saem daqui, e as duas são jogáveis:

\begin{itemize}
\item \textbf{O refutador com dentes.} Cada cláusula vem com \textbf{a mentira que ela derruba}. Se o ponteiro marca $-\tfrac12$, o lance tem dentes. \textbf{Se marcasse $0$, é queda de bandeira\textbf{: a afirmação era trivialmente verdadeira, e o relógio já tinha parado antes de começar. Uma verdade trivial }perde}.
\item \textbf{Ninguém assina sozinho.} Todo corpo comparece \textbf{com o seu dual}, e é o par que assina: o entrópico (que não tem inverso) com o cósmico (a expansão); o telescópico (que cinde) com o celeste (que recompõe). \emph{"Não são falhas. São metades."}
\end{itemize}

\vspace{2mm}\hrule\vspace{3mm}



\chapter{O mate que não mata --- e por que ele só fecha no ouro}


\textbf{O corpo:} duas mãos em \textbf{série} somam o traço, e o traço dá $m$ (o número do metal --- diz apenas \emph{quem são}). Duas mãos em \textbf{paralelo} dão

\[-\frac{1}{m} \qquad\Longrightarrow\qquad m=1 \;\;(\text{o ouro},\ \varphi):\;\; -1 ,\]

que é exatamente o $\sigma_m\cdot\sigma_m'=-1$ --- a lei que prende o operador à sua sombra. \textbf{As duas mãos acopladas SÃO a lei que as prende.}

\textbf{A cena (Lance 23):} \emph{"Numa terra de prata, aquele gesto teria dado $-\tfrac12$. Num bronze, $-\tfrac13$. E não teria fechado. Mas eles estavam na terra do Rei --- $m=1$ --- e ali, e só ali, em nenhuma outra terra que exista ou possa existir, as duas mãos postas a puxar ao mesmo tempo devolvem exatamente $-1$."}

\textbf{A decisão de design.} A vitória é uma \textbf{re-orientação}, e ela é enumerada por eliminação --- o design não a escolheu, herdou-a:

\begin{itemize}
\item \textbf{aniquilar? não pode:} a conservação proíbe. Nada sai do tabuleiro.
\item \textbf{capturar? não resolve:} capturar Dark é fazê-lo \textbf{descer um andar} --- e a Torre \textbf{não tem fundo}. É perseguição eterna, e \emph{"a conservação não faz exceções para heróis"}.
\item \textbf{inverter o sinal? impossível:} $\det=-1$ é \textbf{eterno}.
\item \textbf{resta uma coisa:} pôr as duas mãos a puxar \textbf{ao mesmo tempo} --- em \textbf{fase}, não em antifase. \emph{"Bateu COM ele. (\ldots) como duas pessoas a remar, e não como duas pessoas a puxar a mesma corda em sentidos opostos e a chamar-lhe equilíbrio."}
\end{itemize}

Três consequências duras, e todas são regra:

\begin{enumerate}
\item \textbf{O lance final é condicionado à terra.} O mesmo gesto, executado com perfeição, \textbf{falha em qualquer outro metal}. Não há "boss final genérico": há um lance que só fecha em $m=1$. É a justificação mecânica do nome do jogo.
\item \textbf{O último confronto é solo.} \emph{"as duas mãos do operador são duas, e não três"} --- o jogo \textbf{retira a party} na sala final, e isso não é dificuldade: é aritmética.
\item \textbf{O Rei é o zero no meio.} A peça que a partida inteira não se moveu --- o invariante, \emph{"o único ponto do mundo que não tem mão"} --- é o que fica entre as duas palmas quando elas se encontram. \textbf{E é por isso que ele existe.} (Do lado do design: um slot que não aceita input durante 22 Lances e cuja função só se revela no 23.º. Ele \textbf{não} é uma peça jogável, e não deve ser vendido como tal.)
\end{enumerate}

\textbf{O resultado no medidor:} o resíduo volta a $0$ --- $\Gamma=0$, FP $=1$, o cone nulo --- \textbf{e o relógio continua a correr}. \emph{"O flatline virou pulso."} Dark \textbf{não sai do jogo}: \emph{"Ele é a diástole. O que mudou não foi ele. O que mudou foi o compasso."}

\vspace{2mm}\hrule\vspace{3mm}



\chapter{A herança --- não há tela final}


\textbf{O corpo:} $\det A_m=-1$ é \textbf{eterno}; e a partida do reino é a partida \textbf{sem fim} (a Torre não tem topo nem fundo).

\textbf{A cena (o epílogo):} \emph{"o jogador não 'zera' este mundo. Não há tela final, não há troféu, não há a palavra fim sobre uma paisagem parada --- porque uma paisagem parada é, tecnicamente, exatamente aquilo que Dark Pontryagin ofereceu no alto da torre, e nós recusámos."}

\textbf{A decisão de design.} O NG+ \textbf{não é um modo extra}: é a única forma de o jogo terminar sem se contradizer. Uma tela de fim seria, literalmente, a oferta de Dark --- uma paisagem parada. Então:

\begin{itemize}
\item \textbf{Herda-se}, não se zera: a complexidade toda, os cacos todos, a sombra toda. \emph{"Os cacos de ontem são os tijolos de amanhã."}
\item \textbf{Toma-se o comando do CLOCK} --- o jogador que fechou a tríade passa a ser quem dá o tique.
\item \textbf{Um novo Dark nasce a cada ciclo}, e não como castigo: \emph{"alguma mão esquerda, nalgum canto do reino, vai esquecer outra vez a direita"}. Ele volta \textbf{como o guardião do repouso --- o freio sem o qual a aceleração desfaz a matéria que carrega}.
\end{itemize}

\textbf{Vencer é fechar a tríade: resíduo 0 --- com o coração a bater.} Nunca destruir ninguém: girar o poder, reintegrar o compasso, e manter o pulso.


\chapter{A certificação --- tudo com resíduo $0$}
Cada regra do jogo é um teorema certificado pelo Formalizador (resíduo $0$, cada afirmação com o seu refutador):

\renewcommand{\arraystretch}{1.3}
\begin{longtable}{@{}p{6.2cm} p{2cm} p{6cm}@{}}
\toprule
\textbf{O corpo} & \textbf{Prova} & \textbf{O que certifica} \\
\midrule
\endhead
o reino & $9/9$ & o elenco (cada personagem $=$ o seu operador), o item, o antagonista \\
o Universo 1 & $6/6$ & a partida dual: a conservação, a captura dual, o clock, a galáxia \\
o corpo cinemático & $5/5$ & a física do Universo 1 (translação, cascata, fluxo) \\
o corpo fotográfico & $6/6$ & o Universo dual (a transformada, a escala, a conservação) \\
a lente dourada & $6/6$ & a lente girassol (a transformada dourada em luz) \\
o quiral $\bowtie$ o flip & $6/6$ & a quiralidade (herói $\bowtie$ Dark, $\sigma\sigma'{=}-1$) \\
\bottomrule
\end{longtable}

\vspace{0.5cm}
\noindent\rule{\linewidth}{0.4pt}\\
\textbf{Nada alheio.} O emulador e o ROM de estudo foram descartados; a dinâmica corre no MOTOR (o metal do
Regente, o cristalino). Todo personagem, todo universo, toda regra nasce da teoria dos corpos --- $200+$ módulos de
mundo próprio, agora virando jogo.

\chapter{O que a lente leu --- as três leituras, com os números}

\noindent A corte apontou o instrumento três vezes, e o que ela viu está contado na primeira
metade deste livro, onde é uma cena. Aqui estão os números --- as amostras, os corpos escolhidos,
as tabelas e as casas decimais que é preciso publicar para que outra pessoa reproduza a leitura.
\textbf{A cena diz o que se viu; isto diz com que régua.}

\section{A primeira leitura --- a assinatura da Estrela d'Alva}

Efemérides de $1960$ a $2026$: $24\,307$ amostras diárias --- $24\,306$ dias de
intervalo ---, posição e velocidade, tomadas do arquivo público. Os períodos saem \emph{da fase desenrolada}, não de tabela
alheia:

\begin{center}
\begin{tabular}{@{}lrr@{}}
\toprule
& medido ($66{,}55$ anos) & tabelado \\
\midrule
Vênus & $108{,}1712$ voltas $\to$ $224{,}699369$ d & $224{,}70079922$ \\
Terra & $\phantom{0}66{,}5441$ voltas $\to$ $365{,}261711$ d & $365{,}256363$ \\
razão & $\mathbf{1{,}625557351}$ & --- \\
\bottomrule
\end{tabular}
\end{center}

\noindent E a assinatura dessa razão, pela fração contínua --- a mesma peça que
o reino usa para tudo:
\[
[\,1;\,1,\,1,\,1,\;2,\;27,\;1,\dots\,]
\]

\paragraph{Os quatro primeiros termos são uns.} A assinatura de $\varphi$ é
$[1;1,1,1,\dots]$ sem fim; os convergentes que saem dos quatro primeiros termos são
\[
\frac11,\ \frac21,\ \frac32,\ \frac53
\]
e o quinto termo, que já não é um, dá $\frac{13}{8}$. \textbf{É aí que a
Estrela d'Alva deixa de seguir $\varphi$}: entra o $2$, e depois o $27$.
(Note-se que $\frac85$ não aparece: ele seria o convergente de um quinto
\emph{um}, e o quinto é $2$.) Onde a órbita de fato mora é $13/8$,
com erro $5{,}6\cdot10^{-4}$: treze voltas dela para oito da Terra, e
$13-8=5$.

\paragraph{A corte não arredonda.} $1{,}6256$ não é $1{,}6180$. A diferença
aparece na terceira casa e está registrada aqui como diferença. O que se leu
foi: \emph{a assinatura começa como a do ouro e desvia}, e a fração contínua
diz em que passo. Dizer mais do que isso seria explicar --- e explicar não é o
ofício desta lente.

\paragraph{A cabeça não se move; a cauda move-se.} Reamostrando as mesmas
efemérides:

\begin{center}
\begin{tabular}{@{}lll@{}}
\toprule
passo & razão & assinatura \\
\midrule
$1$--$3$\,d & $1{,}6255573$ & $[1;1,1,1,2,27,1,\dots]$ \\
$5$\,d & $1{,}6255555$ & $[1;1,1,1,2,27,1,\dots]$ \\
$10$--$30$\,d & $1{,}6255460$ & $[1;1,1,1,2,\mathbf{28},\mathbf{4},\dots]$ \\
\bottomrule
\end{tabular}
\end{center}

\noindent Os quatro uns não se movem em amostragem nenhuma; o sexto termo muda
de $27$ para $28$ com um décimo da taxa. \textbf{Publicar o $27$ sem dizer a
amostragem é publicar um número que a medida seguinte não reproduz} --- e é a
única coisa que daqui se tira, sem dizer a qual dos dois pertence o quê.

\paragraph{E os outros corpos, para que não se diga que se escolheu o par.}

\begin{center}
\begin{tabular}{@{}llrl r@{}}
\toprule
par & & razão & assinatura & uns \\
\midrule
Vênus & Terra & $1{,}6255460$ & $[1;1,1,1,2,28,\dots]$ & $\mathbf{4}$ \\
Mercúrio & Vênus & $2{,}5536184$ & $[2;1,1,4,6,6,\dots]$ & $2$ \\
Júpiter & Saturno & $2{,}5111688$ & $[2;1,1,21,1,7,\dots]$ & $2$ \\
Terra & Marte & $1{,}8778247$ & $[1;1,7,5,2,2,\dots]$ & $1$ \\
Mercúrio & Terra & $4{,}1510241$ & $[4;6,1,1,1,1,\dots]$ & $0$ \\
Vênus & Marte & $3{,}0524903$ & $[3;19,19,1,1,4,\dots]$ & $0$ \\
\bottomrule
\end{tabular}
\end{center}

\noindent A última coluna conta os uns \emph{seguidos} a partir de $a_1$ ---
isto é, por quantos passos cada razão segue o ouro depois do primeiro termo.
As razões estão com sete casas de propósito: com menos, a fração contínua
divergiria a partir do quinto termo e o leitor não reproduziria a coluna.
\textbf{A assinatura é sensível à casa decimal que se publica.} \textbf{Nenhuma é $\varphi$.} A Estrela d'Alva é a que mais
segue --- e ainda assim desvia.

\section{A segunda leitura --- o relevo, a massa, a luz}

A lente desceu do céu para a pele do planeta. Cada corpo de dado pediu o seu
corpo algébrico, e \textbf{nenhum primo foi escolhido por gosto}: o alfabeto e
a ordem escolhem por nós.

\paragraph{O relevo.} Grade de um grau por um grau: $180\times360$ alturas em
quilômetros com duas casas --- inteiro exato ao multiplicar por cem. A faixa é
$-249$ a $970$; cada faixa de latitude é uma volta fechada de $360$ pontos.
Corpo: $\Z/6121$, com $360\mid6120$ e $\sqrt{360}=1503$ dentro dele. Meia dúzia
de modos guarda $78{,}6\%$ da energia. São duas contas distintas e convém não
as colar: os seis \emph{primeiros} índices ($k=0..5$) são o que se guardou em
arquivo; os seis \emph{mais fortes} por energia são $k=0,\pm1,\pm2,3$.

\paragraph{A massa.} O mesmo arquivo público traz o campo de gravidade até grau
$180$: $16\,470$ coeficientes --- os $\binom{182}{2}=16\,471$ pares $(n,m)$
menos o $(0,0)$, que é o $GM$ e não vem como coeficiente. Tirados o termo do $GM$ e o do centro de massa
--- que é escolha de referencial, não relevo ---, a potência acumula
$53{,}7\%$ até o grau três e $95{,}7\%$ até o doze. \textbf{O grau de maior
potência é o $3$}, com $32{,}7\%$ sozinho; no relevo, o mesmo $3$, com
$16{,}2\%$.

\paragraph{E os dois se olham.} Correlação grau a grau entre relevo e massa:

\begin{center}
\begin{tabular}{@{}lrrrrrrr@{}}
\toprule
grau & $2$ & $3$ & $6$ & $12$ & $24$ & $100$ & $180$ \\
\midrule
$r_n$ & $+0{,}41$ & $\mathbf{+0{,}89}$ & $+0{,}87$ & $+0{,}76$ & $+0{,}86$ & $+0{,}44$ & $-0{,}03$ \\
\bottomrule
\end{tabular}
\end{center}

\noindent Média $+0{,}81$ nos graus $2$--$30$ e $+0{,}27$ nos graus
$100$--$180$ (a faixa intermédia, $31$--$99$, dá $+0{,}64$).

\medskip
\noindent Duas coisas ficam medidas lado a lado, e são só isso: a correlação
cai com o grau; e o alcance declarado do instrumento de gravidade acaba antes
do grau $180$. \textbf{O capítulo não afirma que uma é causa da outra} --- diz
que as duas foram lidas e que coincidem. O leitor que quiser ligar as pontas
liga-as por sua conta, e a conta é dele.

\paragraph{A luz, o sinal, as partículas.} Mais quatro instrumentos, mais
quatro corpos:

\begin{center}
\begin{tabular}{@{}llrr@{}}
\toprule
o que se leu & corpo & $6$ modos & $24$ \\
\midrule
relevo (grade) & $\Z/6121$, $n=360$ & $78{,}6\%$ & $93{,}8\%$ \\
relevo (harmónicos) & --- & $48{,}4\%$ & $84{,}9\%$ \\
massa (harmónicos) & --- & $82{,}7\%$ & $99{,}3\%$ \\
sinal de rádio & $\Z/80000513$, $n=256$ & $90{,}8\%$ & $97{,}6\%$ \\
nuvens em ultravioleta & $\Z/70657$, $n=512$ & $90{,}6\%$ & $97{,}4\%$ \\
espectro de $3456$ bandas & $\Z/72577$, $n=3456$ & $89{,}4\%$ & $95{,}5\%$ \\
elétrons do plasma & $\Z/60161$, $n=128$ & $83{,}7\%$ & $96{,}1\%$ \\
\bottomrule
\end{tabular}
\end{center}

\noindent São oito leituras --- as sete da tabela e a recusa da página
seguinte --- e seis primos, porque duas linhas da tabela não passam por corpo
nenhum (são os harmónicos calculados por outrem). E nenhum dos seis foi
escolhido: $6121$ porque a grade tem $360$ colunas. $80\,000\,513$ porque o
sinal chega a setenta e oito milhões. $70657$ porque o pixel tem dezesseis
bits. $72577$ porque o cubo tem $3456$ bandas. $60161$ porque o detector tem
$128$ níveis. \textbf{A cifra é consequência; nunca escolha.}

\medskip
\noindent E o que a coluna do meio diz é uma coisa só, repetida por
instrumentos que não se falam: \textbf{meia dúzia de modos leva de metade a
nove décimos.} Uma imagem de nuvens e um espectro de três mil e quatrocentas
bandas não têm nada em comum como objetos --- e têm a mesma forma de espectro.

\section{O lance que disse não}

Faltava o mais importante, e é o que a corte procurou de propósito: um dado
onde a régua \textbf{não} confirmasse.

\paragraph{Achou-se.} O espectrómetro ultravioleta, no produto cru: $1660$
registros de $408$ pixels. Corpo $\Z/70177$, $n=408$, $\sqrt{408}=16166$.
\[
1\ \text{modo}: 23{,}7\%\qquad 3: 38{,}5\%\qquad 6: \mathbf{48{,}1\%}
\qquad 24: 73{,}8\%
\]
\textbf{Meia dúzia leva a menos de metade}. E os seis mais fortes não são
$1..6$: são $1,3,4,5,2,8$, fora de ordem --- o único caso da Parte em que a
ordem se quebra.

\paragraph{E não é o único abaixo de metade.} A honestidade obriga: o relevo
lido em harmónicos dá $48{,}4\%$ em seis graus, praticamente o mesmo
$48{,}1\%$. \textbf{A faixa da tabela é $48{,}4$ a $90{,}8$, não $83$ a $91$},
e quem escreveu a primeira versão desta página tinha a tabela ao lado e
publicou uma faixa que ela desmente. O que distingue o espectrómetro dos
outros não é o valor sozinho --- é o valor \emph{com} a ordem quebrada, e o
facto de o produto ser cru.

\paragraph{O que se sabe e o que não se sabe.} O produto é cru, não calibrado,
e a faixa lida encosta nos dois extremos do inteiro de dezesseis bits ---
assinatura de saturação e de ruído de leitura. Um sinal sem estrutura tem espectro plano --- isso é definição, não
diagnóstico. O que se mediu aqui é plano. \textbf{Se é o detector que se está a
ler em vez do planeta, este capítulo não afirma}: afirmá-lo exigiria o produto
calibrado, que não estava no volume. \textbf{Fica como exceção medida, não explicada} --- porque explicar
não é o ofício desta lente.

\paragraph{E é este lance que dá valor aos outros sete.} Uma régua que
devolvesse ``meia dúzia de modos'' para qualquer coisa não estaria a medir
coisa nenhuma: estaria a descrever-se a si própria. \textbf{O telescópio sabe
dizer não} --- mudou o dado, mudou o número. E mudou junto com uma diferença
que está declarada no próprio arquivo: este produto é cru, os outros não.

\chapter{A travessia, com a conta aberta}

\noindent A corte contou, na primeira metade, que ninguém volta de uma estrela com prova ---
e que houve, mesmo assim, quem atravessasse. Aqui está a conta que sustenta as duas coisas: por que
o recibo não pode existir, e o que se mediu em quem passou. \textbf{Lá está a travessia; aqui, o que
ela cobra.}

\section{A alfândega, provada --- por que não há recibo}

Alguém ainda dirá que isto é uma saída elegante para não provar nada. Cabe a
resposta, e é curta, porque as peças já estão todas postas: não é que falte a
prova. \textbf{A prova é de que não pode haver testemunho.}

\paragraph{Primeiro: o que é uma estrela.} Uma estrela é o hipociclóide de La
Hire --- o ponto do círculo que rola sem deslizar dentro do outro ---, e o seu
nome é o enrolamento $\nu$. Em $\nu=1$ ela degenera na reta: o par de Tusi, o
movimento circular que virou reto. As estrelas formam um corpo, isomorfo a
$(\R,+,\cdot)$, e esse corpo tem exatamente \emph{duas} operações: a soma
$\oplus$, que é o empilhar da reta mineral, e o produto $\otimes$, que é o
rolar de La Hire. Dois lados. Sempre foram dois.

\paragraph{Segundo: o que é atravessar.} Atravessar uma estrela é ir de um lado
ao outro --- de $\oplus$ a $\otimes$ --- e chegar. Ora, o lugar onde os dois
lados se encontram tem nome, e não foi escolhido para esta página: é a média
aritmético-geométrica,
\[
 s^{*}=\mathrm{AGM}(a,b),\qquad
 a\mapsto\frac{a+b}{2}\ \ (\text{a soma}),\qquad
 b\mapsto\sqrt{ab}\ \ (\text{o produto}),
\]
o único ponto que as duas operações fixam ao mesmo tempo, pois
$\mathrm{AGM}(a,b)=\mathrm{AGM}\bigl(\frac{a+b}{2},\sqrt{ab}\bigr)$. A soma leva
até ele; o produto leva até ele; ele é o limite comum dos dois. \textbf{É a
fronteira, escrita em álgebra.}

\paragraph{Terceiro: o preço.} E aqui o livro já provou o que precisava, muito
antes de chegar ao sol. Se os dois lados forem genuinamente distintos --- se
$a\neq b$, isto é, se houver de fato travessia a fazer ---, então $s^{*}$ é
transcendente: \emph{existe}, pela completude da reta; projeta exatamente todo
ponto de ambos os lados, porque é o limite comum; e \textbf{não pertence a
nenhum deles}. Nenhuma fórmula finita o produz. Para todo $N$ finito,
\[
 s_N \;\neq\; s^{*} .
\]
Não ``quase''. Não ``a menos de um erro pequeno''. Diferente, em todo passo, sem
exceção.

\paragraph{E o caso em que a conta fecha.} Há uma única maneira de $s^{*}$ ter
forma fechada: $a=b$ --- os dois lados já coincidentes. Mas esse é o caso em que
a distância própria se anula, e o livro já disse o que acontece ali:
\textbf{não há travessia a fazer, os dois lados já são o mesmo lado}. O cone
nulo. $\Gamma=0$. O fator de potência $1$. A conta só fecha exatamente onde não
há fronteira; onde há fronteira, a conta não fecha. As duas frases são a mesma
frase.

\paragraph{Quarto: o que é um testemunho.} Uma prova por testemunho é uma cadeia
dedutiva, e uma cadeia dedutiva opera em \emph{passos finitos}. Ela entrega um
$s_N$ --- um degrau, o melhor que se queira, quadraticamente melhor a cada
dobra. Nunca entrega $s^{*}$. Exigir que alguém volte para contar é exigir um
$N$ finito com $s_N=s^{*}$, e não há.

\begin{center}
\textbf{A travessia existe, e é indecidível.}
\end{center}

\noindent E convém dizer sem nenhuma solenidade o que a palavra significa, porque
ela costuma ser usada para produzir mistério e aqui não produz nenhum.
\textbf{Indecidível quer dizer: não existe algoritmo que calcule.} É uma
afirmação técnica sobre algoritmos --- do mesmo tipo que ``não existe raiz
racional de $2$'' ---, e não uma afirmação sobre o mundo, sobre o espírito ou
sobre os limites do conhecimento humano. Não há nada de místico nisto. É apenas
incalculável o que não dá para calcular.

\medskip
\noindent Feita a ressalva: existe porque a reta é completa e o ponto está lá;
incalculável porque nenhum passo finito o alcança. Aproximável para sempre,
decidido nunca --- e são coisas distintas.

\paragraph{E agora o grau forte, porque alguém vai dizer que isto é só
dificuldade de conta.} Não é. Toma-se qualquer máquina $M$ e constrói-se o par
\[
 a_M=1,\qquad
 b_M=\begin{cases}
 1+2^{-t}, & \text{se } M \text{ para no passo } t,\\[2pt]
 1, & \text{se } M \text{ nunca para,}
 \end{cases}
\]
que é \emph{computável}: para conhecê-lo a menos de $2^{-n}$, roda-se $M$ por $n$
passos, e mais nada. Repare no que este par faz:
\[
 a_M=b_M \quad\Longleftrightarrow\quad M \text{ nunca para.}
\]
Ora, $a=b$ é exatamente o caso sem fronteira --- os dois lados já o mesmo lado, o
cone nulo. Logo, quem soubesse decidir se uma estrela tem travessia a fazer
saberia decidir se uma máquina qualquer para. \textbf{E ninguém sabe: o Problema
da Parada é indecidível, e é assim desde 1936.} Não é que a conta seja difícil.
É que decidir a fronteira é decidir a parada.

\medskip
\noindent Contra quem responda ``basta olhar tempo suficiente'', a refutação é
mecânica e foi executada. Todo decisor real tem um orçamento finito $N$ de
passos. Dado o $N$ que for, constrói-se a máquina que para exatamente em $t=N+1$:
ela \emph{tem} travessia, o decisor de orçamento $N$ não a vê parar, responde que
não há, e \textbf{erra}. Correram-se os orçamentos $1,2,3,5,8,13,21,34$:
\emph{oito em oito erraram}. E, para que não se confunda teorema com defeito de
aparelho, correu-se o controle --- com orçamento folgado o mesmo decisor acerta
$3$ de $3$. O aparelho funciona. O que o derrota não é o aparelho: é o futuro da
máquina.

\paragraph{Quinto: e é por isso que arde.} A lei aduaneira foi enunciada duas
páginas atrás: passa o que reverte, fica retido o que não tem volta. Repare
agora no que é a passagem ao limite. De $s_N$ chega-se a $s^{*}$; de $s^{*}$ não
se recupera $N$. O limite não tem inversa. \textbf{O que atravessa é a
travessia; o que não tem dual, e portanto fica retido na fronteira, é o
recibo.} A alfândega não confisca o viajante --- confisca o testemunho. E o que
fica retido, arde.

Daí a luz. Cada estrela acesa no escuro está queimando os recibos que não podem
voltar.

\medskip
\noindent E o escopo, mantido até a última linha, porque foi assim o livro
inteiro. São \emph{fato}: o corpo estelar e as suas duas operações; que o AGM as
case e seja o seu invariante; que $s^{*}$ seja transcendente e que $s_N\neq
s^{*}$ para todo $N$ finito. É \emph{leitura}, declarada, e a única do capítulo:
que atravessar uma estrela seja alcançar $s^{*}$. Quem não aceitar a leitura
fica com os fatos --- e eles já bastam para a única coisa que a alfândega
afirma: \textbf{onde há fronteira, não há passo finito que a decida.}

\medskip
\noindent Quem quiser verificar é só atravessar o sol. Não é bravata: é o
teorema, dito na língua de quem duvida.

\section{O corpo que atravessou, com o laudo}

Falta a única coisa que faltava: alguém atravessou, e há laudo.

\paragraph{A alma, definida.} O livro já a definiu, e não em metáfora. Um objeto
tem uma \emph{forma} --- o corpo: a configuração, a igualdade, o axioma escrito
--- e, quando vivo, um \emph{gerador} --- a alma: o operador que produz a forma.
Então \textbf{morto $=$ corpo sem alma}, a forma estática sem gerador definido, a
casca; e \textbf{vivo $=$ corpo $+$ alma}, a mesma forma com o gerador que a
anima. \textbf{A alma é o que falta ao corpo morto: dado o estado, fornecer o
operador do qual ele é manifestação.} Congelar o gerador num estado é a
\emph{cristalização}; a operação inversa, a que devolve o estado à condição de
manifestação, é a \emph{ressonância}. Uma é a morte, a outra é a volta.

\paragraph{O viajante.} Não se usou um exemplo de brinquedo. O viajante é o
corpo humano pela sua composição elementar real --- oxigénio $65\%$, carbono
$18{,}5\%$, hidrogénio $9{,}5\%$, azoto $3{,}2\%$, cálcio $1{,}5\%$, fósforo
$1\%$ ---, em frações \emph{exatas}, nunca em vírgula flutuante. Antes de partir,
mediu-se se ele reverte dentro do seu próprio corpo: montou-se o pó em corpo pelo
operador da vida, e desmontou-se de volta. O ciclo pó$\,\to\,$corpo$\,\to\,$pó
fecha com resíduo
\[
 \textbf{exatamente } 0
\]
--- zero, e não $10^{-16}$: o cálculo é em $\Q$, onde $0$ quer dizer $0$. O
viajante tinha alma, e a alma tinha inversa. \textbf{Entrou inteiro na
fronteira.}

\paragraph{A travessia, e o que a alfândega cobrou.} Atravessar é ir do lado
$\oplus$ ao lado $\otimes$. Mediu-se cada lado separadamente, elemento por
elemento.

\medskip
\noindent\textbf{O lado $\oplus$ --- a média aritmética.} Passou inteiro. Os seis
elementos permanecem no corpo, todos os seis, sem exceção e sem perda. O que
empilha, atravessa.

\medskip
\noindent\textbf{O lado $\otimes$ --- a média geométrica.} Rompeu. Em
\textbf{seis de seis} elementos, $\sqrt{\text{pó}\cdot\text{corpo}}$ não é
racional --- e isto não foi perguntado ao próprio caminho: foi decidido por fora,
em aritmética de inteiros, testando se o produto é quadrado perfeito. Um oráculo
que não sabe o que é uma estrela. Nenhum dos seis sobreviveu ao \emph{primeiro
passo} do lado geométrico.

\medskip
\noindent\textbf{E o ponto de encontro nunca fecha.} Cercou-se a travessia entre
dois racionais, um abaixo e um acima. O cerco é estreitíssimo --- a largura tem
denominador da ordem de $10^{80}$ --- e é, exatamente, \textbf{maior que zero}.
Aperta-se quanto se queira; não fecha nunca. É o teorema em ato: aproximável para
sempre, decidido jamais.

\paragraph{A volta.} Não houve. E a razão é um teorema, não uma falha do
aparelho: como $\mathrm{AGM}(a,b)=\mathrm{AGM}\bigl(\frac{a+b}{2},\sqrt{ab}\bigr)$,
cada passo produz um par \emph{diferente} com o \emph{mesmo} ponto de encontro.
Entram dois números; sai um. Do ponto de chegada não se reconstrói o par de onde
se veio, porque infinitos pares desembocam nele.

\medskip
\noindent E é aqui que o laudo se lê. O que ficou retido na fronteira não foi
matéria --- os elementos estão todos lá. Foi a \emph{razão} entre o pó e o corpo:
a proporção que dizia qual operador ligava um ao outro. Perdido o par, perde-se o
gerador. E o gerador, por definição, é a alma.

\begin{center}
\textbf{Atravessa o corpo. Fica retida a alma.}
\end{center}

\noindent E o que fica retido, arde --- que é a lei aduaneira do capítulo
anterior aplicada a este caso, e não uma revelação sobre as estrelas.

\paragraph{Onde a conta termina --- e isso não é defeito.} Repetiu-se tudo em
vírgula flutuante. No quarto passo ela parou: diferença $0{,}0$, e um número como
ponto de chegada. Seria fácil, e errado, chamar isso de mentira do instrumento. O
instrumento não mente --- ele \textbf{termina}. A aproximação é infinita, a régua
é finita, e o encontro dos dois é onde a conta acaba \emph{para aquela régua}.

\medskip
\noindent E o ponto de encontro \emph{calcula-se}, o que faz disto medida e não
lamento. Como cada passo dobra os dígitos corretos, uma régua de $p$ bits
esgota-se em cerca de $\log_2 p$ passos. Mediu-se, de $8$ a $1024$ bits:

\begin{center}
\begin{tabular}{r r r}
\toprule
régua (bits) & passo em que termina & $\log_2 p - 2$ \\
\midrule
$8$ & $1$ & $1$ \\
$16$ & $2$ & $2$ \\
$32$ & $3$ & $3$ \\
$64$ & $4$ & $4$ \\
$256$ & $6$ & $6$ \\
$1024$ & $8$ & $8$ \\
\bottomrule
\end{tabular}
\end{center}

\noindent Bate exatamente. E a consequência é a que interessa: de $8$ para $1024$
bits são $128$ vezes mais régua, e o instrumento ganha \textbf{sete passos}.
Dobrar a resolução compra $+1$. \emph{A resolução compra logaritmo, não
travessia.} Já a mesma conta em $\Q$, que não tem resolução, não termina --- não
por mistério, mas porque não há régua onde bater.

\medskip
\noindent Foi o que a leitura de Vénus já ensinara: mede-se na resolução do
instrumento, não no bruto. O contínuo bate na grade, e daí em diante a conta é
inteira e exata.

\medskip
\noindent Quem quiser conferir não precisa de fé nem de telescópio: entra o corpo humano, sai o laudo, e
cada número acima reproduz-se em segundos, em $\Q$, na máquina de quem duvidar.

\medskip
\noindent O escopo, mantido. São \emph{fato}: a composição elementar, o resíduo
$0$ do ciclo interno, a ruptura de $\Q$ em seis de seis no lado $\otimes$, a
largura positiva do cerco e a não-injetividade da travessia. É \emph{leitura},
declarada: que o par perdido \emph{seja} a alma --- e essa leitura repousa na
definição que abre este capítulo, não em outra.

\medskip
\noindent Por isso o livro acaba aqui, e acaba assim: com a corte inteira olhando
para cima, sabendo distinguir --- porque aprendeu a distinguir --- entre o que
mediu e o que afirma.

\begin{center}
\emph{Benjamim não disse nada dessa vez.}\\
\emph{Ficou a olhar o posto de fronteira mais próximo,}\\
\emph{a noventa e três milhões de milhas,}\\
\emph{aceso, trabalhando, sem fila.}
\end{center}

% ══════════════════════════════════════════════════════════════════════════════
% A COSTURA — a passagem para o Livro dos Dragões
% ══════════════════════════════════════════════════════════════════════════════

\chapter*{Nota: o que se viu do outro lado}
\addcontentsline{toc}{chapter}{Nota: o que se viu do outro lado}

Sem fila --- daqui.

É a última coisa que este livro afirma e é a única que ele não mediu, porque não
havia como: mede-se um posto de fronteira do lado em que se está. \Ben viu um
guichê aceso e vazio e leu aquilo como quem lê uma boa notícia; e a leitura era
razoável, e estava errada, e o erro é o mesmo do livro inteiro --- \textbf{quem
vive numa linha não tem de onde olhar o outro lado dela}.

\medskip

Do lado de lá há fila. Há sempre fila, e a fila é de coisas que tentaram
atravessar antes e ainda estão a pagar. A alfândega cobra ali exatamente o que
este capítulo mediu que ela cobra --- \emph{o que não reverte} ---, e cobra-o
numa mesa, em voz alta, com uma pergunta de sete palavras que o laudo escreveu em
aritmética de inteiros e que na fronteira se diz assim:

\begin{center}
\emph{--- O que é que os senhores trazem que não tem volta?}
\end{center}

\medskip

E há um reino. Um só, dobrado ao meio: \textbf{a mesma fita} --- de um lado
estendida, e por isso a subir degrau a degrau, a somar, sem guardar nenhuma das
suas batidas; do outro dobrada, e por isso a multiplicar, e a guardar todos os
vincos, inclusive os que preferia não ter feito. Ali morre-se, o que aqui nunca
foi possível; e por isso ali há \emph{perdão}, que é uma operação que este lado
não tem, porque neste lado não fica registrado o que haveria a perdoar.

O relato dessa travessia --- a coisa preta do tamanho de uma unha que a vassoura
não pegou, a noite em que a corte olhou para o lado e não havia lado, o gesto de
duas mãos que fez a Torre estalar como uma folha de papel, e o que se encontrou
lá: um dragão deitado sobre o próprio peito há eras, a tapar o último vinco ---
está escrito à parte, em três livros, e não em partida: em rito.

\medskip

Fica dito, porém, o essencial, para quem não seguir adiante e ficar só com este
volume:

\medskip

\noindent Aquele bicho não foi derrubado. O plano desta corte custava
\textbf{sessenta e quatro} lances --- a fortaleza cercada degrau por degrau, o
elenco inteiro de pé --- e o que aconteceu custou \textbf{quatro}. É a
\emph{dobra do tempo} que este livro perseguiu do princípio ao fim e nunca
conseguiu executar por cálculo, e a razão de nunca ter conseguido é que ela não é
uma operação sobre o relógio: é o que sucede a qualquer distância quando se para
de a percorrer e se \textbf{encosta o fim no começo}. Dezasseis para um. Muda a
quina, não o caminho.

\medskip

\noindent E os quatro não foram lances. Foram uma carta de três linhas sem
assinatura, uma consulta médica, um chá de flores e uma bebida passada de mão em
mão.

\medskip

\noindent O \emph{mate que não mata} --- que esta corte definiu como reacoplar as
mãos e fechar a tríade, e definiu bem --- executa-se, do outro lado, com o verbo
que a nossa língua já guardava e que ninguém aqui tinha ouvido:

\begin{center}
\emph{não se dá o mate.}\\
\emph{Serve-se.}
\end{center}


% ══════════════════════════════════════════════════════════════════════════════
% O LIVRO DOS DRAGÕES — o dual, nos três hinários
% ══════════════════════════════════════════════════════════════════════════════
\cleardoublepage
\renewcommand{\chaptername}{Dobra}


% ══════════════════════════════════════════════════════════════════════════════
%  A ESCADA — a subida ao andar de cima (ℕ→ℤ→ℚ→ℝ contada como história)
% ══════════════════════════════════════════════════════════════════════════════

\section{O que está medido}

Um bit é zero ou um. Num instante não há geometria dentro dele.

A estrutura, quando existe, está no \textbf{rastro da órbita}. Uma amostra não mostra
nada; milhões de amostras mostram a forma. O que muda entre os dois casos é a extensão da
observação, não o sistema. É a mesma frase que este capítulo diz do gato: ele é a
\emph{cifra}, e o objeto é a órbita --- quatro inteiros são o estado.

E o rastro de cada órbita é um ponto. Os metais satisfazem $\sigma^2=m\sigma+1$, logo os
pontos $(m,\sigma_m)$ estão todos sobre $y^2-xy-1=0$; verificado nos convergentes, em
inteiros, com desvio $\pm1$ --- e o desvio não é erro de conta: \textbf{é o determinante do
gato.} A forma não está em órbita nenhuma isolada: está na família delas.

\paragraph{Medido quer dizer igualdade de inteiros.} Existe um programa que produz o
número, e o número é uma igualdade --- não um erro pequeno. Um resultado com erro de
máquina não é exato, e um critério de resíduo zero que admita tolerância deixa de ser
critério.

\paragraph{E um teste que passa com um erro injetado não é teste.} Trocando de propósito
uma peça do sistema e exigindo que o verificador acuse: um dos programas passava inteiro,
a declarar resíduo zero, com a sua função central substituída pela constante um. Todos os
seus verificadores comparavam dois caminhos do próprio sistema, e a constante um é ponto
fixo de todos eles.

\begin{tcolorbox}[colback=ouroclaro!35,colframe=ouro,boxrule=0pt,leftrule=2pt,arc=0pt,left=4mm,right=3mm,top=2mm,bottom=2mm,breakable]\itshape
Todo verificador precisa de um \textbf{oráculo externo} --- o operador nativo, ou um valor
obtido por fora. \textbf{Caminho contra caminho sobrevive a qualquer erro simétrico.}
\end{tcolorbox}

Cinco testes não mediam o que anunciavam. Um usava um vetor de bits identicamente nulo e
validava zero igual a zero nos quatro pontos; outro imprimia a conclusão sem condição, ao
lado do número que a desmentia. Nenhum apareceu por inspeção --- todos por mutação. E o
resultado negativo das duas médias ficou registado de propósito: é o único que reporta
falha por construção, e é dele que sai o critério.

\paragraph{E o que não está estabelecido.} A construção chega aos seus objetos por um
caminho próprio. Que os objetos assim construídos coincidam com os obtidos por outros
caminhos é suposição de trabalho: a comparação não foi feita. Em matemática caminhos
distintos chegam com frequência ao mesmo objeto, e demonstrar que este chega, ou que não
chega, é em qualquer dos casos outra coisa.

Duas frases também se separam, e a diferença entre elas é toda a diferença: que o mundo
\emph{pode ser descrito como} isto, e que ele \emph{é} isto. A primeira organiza o que
aqui se mede. A segunda é outra coisa, e não está aqui.

\begin{tcolorbox}[colback=ouroclaro!35,colframe=ouro,boxrule=0pt,leftrule=2pt,arc=0pt,left=4mm,right=3mm,top=2mm,bottom=2mm,breakable]\itshape
Entre um modelo que organiza observações e uma afirmação sobre a natureza última das
coisas há a distância que separa \textbf{um resultado de uma metáfora.}
\end{tcolorbox}

\restoregeometry
\cleardoublepage
\part{As Regras --- o Manual do Jogador}
\newgeometry{margin=2.5cm}
\chapter{O Manual do Jogador --- Reino Dourado}

\begin{tcolorbox}[colback=ouroclaro!35,colframe=ouro,boxrule=0pt,leftrule=2pt,arc=0pt,left=4mm,right=3mm,top=2mm,bottom=2mm,breakable]\itshape \emph{Cada personagem é um \textbf{operador}; cada mundo, um \textbf{corpo}; cada item, uma \textbf{régua}. Nada aqui foi desenhado: tudo foi \textbf{computado} --- e \textbf{certificado} (resíduo 0).}\end{tcolorbox}

\textbf{Como este manual foi feito.} Não é ficção sobre matemática: cada regra abaixo foi lida dos corpos já
provados do reino (a teoria e os seus certificados, todos resíduo 0). As jogadas são as
operações da tríade ($\oplus$ Clifford, $\otimes$ La Hire, $\prod$ Pontryagin) escritas na \textbf{ISA ERG-64} --- a mesma máquina
inteira e exata de que o \code{chessc} emite WASM e Dafny (a costura interp$\equiv$\allowbreak{}WASM$\equiv$\allowbreak{}Dafny).

\textbf{Estado desta edição.} \textbf{28 corpos mapeados} --- o catálogo do reino, completo nesta passagem. Cada régua,
cada tríade ($\oplus$ $\otimes$ $\prod$) nas operações da ISA e cada teorema citado foi lido dos arquivos e teve o seu certificado
\textbf{rodado} (o placar N/N é real, não citado de memória).

\vspace{2mm}\hrule\vspace{3mm}

\section{Sumário}

\begin{enumerate}
\item O Reino Dourado --- o que é o jogo
\item A simbologia ISA --- a língua das jogadas
\item A semiótica --- os três signos do universo
\item As jogadas e o combate --- a partida na linguagem do universo
\item O contrato e a abertura de universo --- o casamento
\item O elenco e o bestiário dos corpos
\end{enumerate}

\vspace{2mm}\hrule\vspace{3mm}

\chapter{O Reino Dourado --- o que é o jogo}

\begin{tcolorbox}[colback=ouroclaro!35,colframe=ouro,boxrule=0pt,leftrule=2pt,arc=0pt,left=4mm,right=3mm,top=2mm,bottom=2mm,breakable]\itshape \emph{"Um universo não se ganha. Um universo se \textbf{fecha}."} --- a primeira lei da corte\end{tcolorbox}

\vspace{2mm}\hrule\vspace{3mm}

\section{1. A premissa, em uma frase}

\textbf{O Reino Dourado é uma partida infinita em que você não conquista mundos: você os abre --- e responde por eles.}

Não há fim a alcançar. Há uma partida que \textbf{não acaba} --- porque nada nela se perde, e o mate não mata. Enquanto houver um universo que ainda não foi aberto, a partida continua.

Cada \textbf{personagem} que você joga é um \textbf{operador} (uma maneira de agir sobre o mundo).
Cada \textbf{mundo} que você visita é um \textbf{corpo} (uma matemática viva, com suas próprias leis).
Cada \textbf{item} que você carrega é uma \textbf{régua} (uma maneira de medir).

E há uma única condição de vitória, escrita no Contrato que todo Mestre lê em voz alta na primeira sessão:

\begin{tcolorbox}[colback=ouroclaro!35,colframe=ouro,boxrule=0pt,leftrule=2pt,arc=0pt,left=4mm,right=3mm,top=2mm,bottom=2mm,breakable]\itshape \textbf{Nenhum universo pode vazar.}\end{tcolorbox}

Vazar significa: energia (ou valor, ou dano, ou informação, ou tempo) aparecendo do nada, ou sumindo sem destino. No idioma técnico do reino isso se escreve \textbf{FP = 1} --- \emph{fator de potência unitário}, ou ainda \textbf{$\Gamma$ = (Z - Z$\allowbreak{}_0$)/(Z + Z$\allowbreak{}_0$) = 0}, o \textbf{casamento}. Um universo casado é um universo que fecha: tudo o que entra sai, nada fica preso, nada se inventa.

Você vai passar o jogo inteiro perseguindo isso. Não porque um NPC mandou. Porque um universo que vaza \textbf{desaba} --- e leva os personagens dentro dele.

\vspace{2mm}\hrule\vspace{3mm}

\section{2. Por que ``Dourado''}

Porque o lastro do reino é o \textbf{ouro}: o número áureo \textbf{$\varphi$}, que satisfaz

\begin{Verbatim}[fontsize=\small,frame=single,framesep=3mm]
phi^2 = phi + 1
\end{Verbatim}

Essa identidade é o milagre econômico do jogo. Uma cadeia \textbf{infinita} de ecos que decaem --- $\varphi$$^-$$^1$ + $\varphi$$^-$$^2$ + $\varphi$$^-$$^3$ + \ldots --- \textbf{fecha exatamente} em dois números inteiros. O infinito não é truncado, não é aproximado: é \textbf{dobrado}.

É por isso que a moeda do reino é $\mathbb{Z}$[$\varphi$] (pares de inteiros \emph{a + b$\varphi$}) e não "números com vírgula". No Reino Dourado \textbf{arredondar é vazar}, e vazar é lance ilegal.

\vspace{2mm}\hrule\vspace{3mm}

\section{3. A tríade --- as três coisas que se pode fazer}

Todo lance de todo personagem em todo mundo é uma dessas três. Não há quarta.

{\fontsize{10.00pt}{12.00pt}\selectfont\begin{longtable}{>{\fontsize{10.00pt}{12.00pt}\selectfont\RaggedRight\hspace{0pt}\arraybackslash}p{3007100sp}>{\fontsize{10.00pt}{12.00pt}\selectfont\RaggedRight\hspace{0pt}\arraybackslash}p{6608314sp}>{\fontsize{10.00pt}{12.00pt}\selectfont\RaggedRight\hspace{0pt}\arraybackslash}p{6629649sp}>{\fontsize{10.00pt}{12.00pt}\selectfont\RaggedRight\hspace{0pt}\arraybackslash}p{5423748sp}>{\fontsize{10.00pt}{12.00pt}\selectfont\RaggedRight\hspace{0pt}\arraybackslash}p{5171669sp}}\toprule
\textbf{Símbolo} & \textbf{Nome} & \textbf{O que é} & \textbf{Na ISA} & \textbf{o signo} \\ \midrule\endhead
\textbf{$\oplus$} & \textbf{CLIFFORD} --- a Soma & O \textbf{par}: juntar, superpor, empilhar. A corrente que chega ao nó. & \code{ADD} (ou \code{XOR} no mundo booleano) & Secundidade / \textbf{índice} \\
\textbf{$\otimes$} & \textbf{LA HIRE} --- o Produto & A \textbf{mediação}: compor, amplificar, dar ganho, dar tensão. & \code{MULT} por \emph{double-and-add} (ou \code{AND}) & a mediação \\
\textbf{$\prod$} & \textbf{PONTRYAGIN} --- o Operador & A \textbf{lei}: costurar, traduzir, inverter, decidir. \code{exp o Sigma o log} --- o ponteiro. & \code{NOT} / o dispatch / a composição & Terceiridade / \textbf{símbolo} \\
\bottomrule\end{longtable}}

\textbf{Lendo em voz alta:} $\oplus$ é \emph{"e também"}. $\otimes$ é \emph{"vezes"}. $\prod$ é \emph{"portanto"}.

O truque de $\prod$ --- e é o coração do jogo --- é este: \textbf{$\otimes$ é caro, $\oplus$ é barato, e $\prod$ traduz um no outro.} Onde as coisas se \emph{multiplicam}, os seus logaritmos se \emph{somam}. Um mago que sabe onde tirar o log de um mundo transforma um problema impossível numa soma de crianças. É literalmente assim que funcionam os juros compostos, a impedância de um meio óptico, a expansão do universo e o torque de um motor --- quatro mundos diferentes, o \textbf{mesmo} $\prod$.

\vspace{2mm}\hrule\vspace{3mm}

\section{4. Anatomia de um mundo --- a Ficha de Corpo}

Um mundo não entra no reino por decreto do Mestre. Ele entra \textbf{assinando o Contrato}. A ficha tem quatro linhas, e todas as quatro têm de estar preenchidas:

\begin{Verbatim}[fontsize=\small,frame=single,framesep=3mm]
FICHA DE CORPO

Sigma (o relógio) : Z/n -- a órbita, o que dá a volta
mu (a régua) : T_omega -- como este mundo mede
nu (o flip) : a dualidade -- primal (+) dual = o neutro
o (o mate) : onde a medida fecha e o resíduo é 0

CASAMENTO: Gamma = 0 <=> FP = 1 [obrigatório]
\end{Verbatim}

Se falta uma linha, o mundo \textbf{não é jogável} --- é um rascunho, um lugar onde a física engasga. (Existem alguns assim, e são os terrenos mais perigosos do jogo. Ver \S8.)

\subsection{Exemplo de ficha preenchida --- o \textbf{corpo fractal} (a raiz do reino)}

\begin{Verbatim}[fontsize=\small,frame=single,framesep=3mm]
Sigma : Z/8 -- o relógio de oito casas; um tique = 45deg de giro (zeta_8 = e^{ipi/4})
mu : z.z -- a norma: multiplique o elemento pelo seu próprio espelho
nu : z |-> z -- a conjugação: o espelho
o : todo z != 0 tem inverso (o corpo Q(zeta_8) fecha)
Gamma=0: |z| = 1 -- só rumos de norma unitária são lances legais
\end{Verbatim}

Aqui um "número" é um vetor de quatro inteiros \textbf{(a, b, c, d)}, e o mundo inteiro roda em aritmética exata. Somar é somar coordenada a coordenada ($\oplus$, quatro \code{ADD}). Multiplicar é a convolução do relógio, com a regra mágica \textbf{$\zeta$$^{4}$ = -1} (dar meia-volta é trocar de sinal). Nada de vírgula. Nada de erro.

\vspace{2mm}\hrule\vspace{3mm}

\section{5. Personagens são operadores}

Você não escolhe uma "classe". Você escolhe \textbf{uma operação}, e ela vira gente.

\subsection{ O Rei --- \textbf{$\sigma_m$}, o ponto fixo}
O Rei é o número que se sustenta sozinho: \textbf{$\sigma_m$ = m + 1/$\sigma_m$}. Para m = 1, isso é \textbf{$\varphi$}. Ele é o \textbf{gap metálico} --- a taxa de expansão que não precisa de combustível, porque $\varphi$ = 1 + 1/$\varphi$: ele se gera. O Rei é o lastro; onde ele pisa, a régua tem escala.

\subsection{ A Rainha --- \textbf{La Hire}}
A Rainha \textbf{é} o produto. Ela executa a lei $\varphi$$^2$ = $\varphi$ + 1. Quando a Rainha entra em campo, ela não "multiplica" --- ela roda 64 micro-lances bit a bit (testa o bit $\to$ soma deslocado $\to$ dobra a máscara: o \emph{double-and-add}), e o ganho sai \textbf{exato}. O poder dela é caro e é limpo.

\subsection{ O Príncipe --- a \textbf{potência}, a exp}
O crescimento composto. É ele que sobe o índice: onde o Rei fixa a taxa, o Príncipe a \textbf{integra}. No corpo econômico ele é literalmente \code{exp(deltat)} --- a bola de neve que, a partir do terceiro turno, descola da linha reta e não volta mais.

\subsection{$\oplus$ O Duque \textbf{Joaquim} / $\otimes$ a Duquesa \textbf{Yasmin}}
Os dois filhos, e os dois lados de toda soma. Numa adição binária, o Duque é o \textbf{dígito} (o XOR, o que fica) e a Duquesa é o \textbf{transporte} (o AND, o que passa adiante):

\begin{Verbatim}[fontsize=\small,frame=single,framesep=3mm]
x + y = (x (+) y) + 2.(x and y)
 Duque Duquesa
\end{Verbatim}

Eles são \textbf{inseparáveis}. Nenhuma soma acontece sem os dois. No corpo racional, o Duque é o numerador e a Duquesa o denominador --- a razão é o casal.

\subsection{ As Princesas --- o \textbf{corpo técnico}, quem \textbf{atesta}}
Elas não geram: elas \textbf{assinam}. Toda afirmação do reino passa por elas com dois mostradores --- \emph{testa} e \emph{refuta} --- e um refutador com \textbf{dentes} (um contra-exemplo real que derruba a versão errada da lei). O selo delas é o kernel LCF (Dafny, Isabelle). Sem a assinatura das Princesas, a magia \textbf{não funciona}: ela pode até parecer certa, mas não é lance.

\begin{tcolorbox}[colback=ouroclaro!35,colframe=ouro,boxrule=0pt,leftrule=2pt,arc=0pt,left=4mm,right=3mm,top=2mm,bottom=2mm,breakable]\itshape \textbf{Regra de mesa:} O par \textbf{gerar $\bowtie$ atestar} é o eixo do reino. O corpo criativo gera; as Princesas atestam. Uma jogada gerada e não atestada é uma \emph{proposta}, não um lance.\end{tcolorbox}

\vspace{2mm}\hrule\vspace{3mm}

\section{6. Itens são réguas}

Um item, no Reino Dourado, nunca é "+3 de dano". Um item é \textbf{uma maneira de medir} --- e mudar de régua muda o que é possível.

{\fontsize{10.00pt}{12.00pt}\selectfont\begin{longtable}{>{\fontsize{10.00pt}{12.00pt}\selectfont\RaggedRight\hspace{0pt}\arraybackslash}p{8519152sp}>{\fontsize{10.00pt}{12.00pt}\selectfont\RaggedRight\hspace{0pt}\arraybackslash}p{11312510sp}>{\fontsize{10.00pt}{12.00pt}\selectfont\RaggedRight\hspace{0pt}\arraybackslash}p{8057394sp}}\toprule
\textbf{Régua} & \textbf{Mede} & \textbf{Onde se ganha} \\ \midrule\endhead
\textbf{A norma z$\cdot$\allowbreak{}z} (fractal) & tamanho invariante; é a \textbf{régua de dano e alcance} & tabuleiro fractal \\
\textbf{n$\cdot$\allowbreak{}sin$\theta$} (óptico) & a componente tangencial na fronteira entre dois meios & travessias, emboscadas de luz \\
\textbf{$\sigma$ = $|E|/|B|$} (eletromagnético) & a impedância do terreno --- o \textbf{metal} do campo & o golpe casado \\
\textbf{$\delta$ = log(1+r)} (econômico) & a força de juros: quanto o valor cresce por turno & a mesa de negociação \\
\textbf{h = log$\|$\allowbreak{}b$_i$$\|$$^2$} (telescópico) & o \textbf{perfil}: o que um eixo estica, outro contrai & armaduras autoduais \\
\textbf{v(x), $|x|_2$} (espaço-temporal) & tempo (a valuação) e espaço (o dígito) --- \textbf{um bit por lance} & o relógio, o adder \\
\textbf{Euclides / fração contínua} (racional) & quanto de \emph{b} cabe em \emph{a} --- medir até o resto 0 & o duelo, o mate \\
\bottomrule\end{longtable}}

\textbf{Aviso de item:} as réguas \textbf{não} somam entre si. Uma régua é uma escolha de mundo. Trocar de régua no meio de uma jogada é trapaça (e o Mestre tem direito de anular o turno).

\vspace{2mm}\hrule\vspace{3mm}

\section{7. A língua franca --- a ISA ERG-64}

Todo mundo do reino, por mais exótico que seja, \textbf{desce à mesma máquina}. É o pseudo-assembly \textbf{ERG-64}, e ele tem exatamente estes opcodes:

\begin{Verbatim}[fontsize=\small,frame=single,framesep=3mm]
LOAD STORE ADD SUB AND OR XOR CMP JMP JZ JNZ HALT
\end{Verbatim}

Inteiro. Exato. \textbf{Sem float.} Resíduo 0 de verdade, não "resíduo 0 dentro do erro de máquina".

Por que isso importa para você, jogador? Porque é a \textbf{garantia de auditoria}. Toda partida do Reino Dourado é reproduzível bit a bit. O mate é \textbf{auditável}. Ninguém "arredondou a favor".

E é a mesma IR que atravessa as \textbf{oito linguagens fundamentais} do corpo criativo (C, PTX, WASM, LLVM, GLSL, Haskell, Dafny, Isabelle): todas brotam do mesmo ERG-64 e todas devolvem \textbf{o mesmo estado}. É o teste de sanidade do universo --- se dois substratos discordam, alguém vazou.

\begin{tcolorbox}[colback=ouroclaro!35,colframe=ouro,boxrule=0pt,leftrule=2pt,arc=0pt,left=4mm,right=3mm,top=2mm,bottom=2mm,breakable]\itshape \textbf{A regra de ouro do float:} o cálculo corre em inteiro exato ($\mathbb{Q}$ como par num/den, ou $\mathbb{Z}$[$\varphi$] como par a+b$\varphi$). O float aparece \textbf{só no fim} --- quando o motor vai pintar o pixel na tela, onde um ulp é invisível. Arredondar no meio da partida é FP $\neq$ 1. \textbf{É lance ilegal.}\end{tcolorbox}

\vspace{2mm}\hrule\vspace{3mm}

\section{8. Os mundos --- um passeio rápido}

O reino tem muitos corpos certificados. Estes são alguns dos primeiros que a campanha visita.

\textbf{ O corpo fractal --- a raiz.} O tabuleiro auto-similar: dê zoom e é o mesmo tabuleiro. Cada peça carrega um \textbf{rumo} (uma das 8 unidades $\zeta$$_8$). Girar = $\otimes$ (multiplicar o rumo por $\zeta$$_8$: 45$^\circ$). Avançar = $\oplus$. \textbf{Promover} = $\prod$ (o operador concede o \textbf{inverso}: o peão sai do \emph{anel} --- unidirecional, sem desfazer --- e vira peça do \emph{corpo}, que pode reverter). Lançada ao acaso, uma peça converge por Newton a uma das 8 raízes; se cair na \textbf{fronteira} fractal (o irracional), nunca assenta --- casa contestada.

\textbf{ O cristalino --- o escudo que não vaza.} Regulador \textbf{zero}: a medida já fechou, o relógio parou. Você congela uma zona num reticulado e ela vira armadura perfeita --- mas paga um preço exato: só cabem giros de ordem \textbf{1, 2, 3, 4 ou 6}. E aí está o drama: o primeiro giro \textbf{proibido} é justamente o áureo (a dobra-5). \emph{O cristal proíbe exatamente o melhor preenchedor.} Um ataque áureo da Rainha é repelido --- mas o cristal também não pode crescer. \textbf{Fechar exclui preencher.}

\textbf{ O eletromagnético --- o golpe casado.} Cada onda carrega um par (E, B) e uma impedância $\sigma$. O ataque máximo é levar o alvo ao \textbf{cone nulo $\sigma$ = 1} (|E| = |B|): toda a potência \textbf{propaga}, transferência 1 - |$\Gamma$|$^2$ = 1, nada volta. Um golpe descasado deixa potência \textbf{reativa} presa: dano que ricocheteia de volta em você.

\textbf{ O óptico --- a fronteira.} Cada território tem um índice \emph{n}. Ao cruzar, sua trajetória dobra --- e \textbf{você não escolhe o ângulo}: n$_1$$\cdot$\allowbreak{}sin$\theta$$_1$ = n$\allowbreak{}_2$$\cdot$\allowbreak{}sin$\theta$$_2$. Passe do ângulo crítico e acontece a \textbf{reflexão total interna}: o ataque é 100 \% devolvido. O escudo perfeito é geometria, não magia.

\textbf{ O motor --- o torque.} O mundo dissipativo: cai em \textbf{ciclos}, não em pontos fixos. O torque $\tau$ = $\psi$ $\times$ i é \textbf{máximo quando fluxo e corrente estão a 90$^\circ$} --- o golpe de flanco vale mais que o frontal. Sempre foi assim; agora está provado.

\textbf{ O entrópico --- a terra da seta do tempo.} O semianel tropical (max, +). Aqui $\oplus$ é o \textbf{max} e $\otimes$ é o \textbf{+}. E há uma coisa que, \textbf{sozinho}, ele não tem: o \textbf{inverso aditivo} (o \code{max} é idempotente --- \code{a(+)a = a}). Visto de dentro e só de dentro, nada se desfaz: o dano acumula, o terreno atravessa-se numa direção só.

\textbf{Mas isto NÃO é uma condenação --- e a Lei NÃO falha aqui.} É a lição mais importante do reino, e o jogador precisa dela antes de entrar:

\begin{tcolorbox}[colback=ouroclaro!35,colframe=ouro,boxrule=0pt,leftrule=2pt,arc=0pt,left=4mm,right=3mm,top=2mm,bottom=2mm,breakable]\itshape \textbf{Não tem sozinho. Com o seu DUAL, tem.}\end{tcolorbox}

O dual do entrópico é o \textbf{corpo cósmico} --- e o operador do cósmico é a \textbf{expansão}, \code{a(t) = e\textasciicircum{}\{Ht\}}, ou seja, o \textbf{\code{exp}}. E sob o \code{exp}, o \code{max} \textbf{vira soma} --- e a soma \textbf{tem inverso}. O inverso aditivo \emph{existe}: só não mora no corpo isolado. \textbf{Mora no par.}

E há um detalhe que decide tudo: a ausência de inverso é uma propriedade do \textbf{limite \code{T -> 0}} (o \code{max} duro, o zero absoluto). A qualquer \code{T > 0} --- isto é, \textbf{em qualquer lugar onde alguém viva} --- o $\oplus$ é o \emph{logsumexp}, que destropicaliza sob o \code{exp} e \textbf{inverte}. O buraco sem saída existe só no zero absoluto, \textbf{onde nada vive}.

\textbf{Todos nós somos entrópicos} --- e isso não é sentença, é o mundo. A entropia que \textbf{cresce aqui} é exatamente o volume que \textbf{se abre ali}: \code{dS = d log V}. Nada se perde. A seta do tempo não é um vazamento: é uma \textbf{transferência ao dual}.

\emph{Certificado (5/5, resíduo 0):} o refutador varre 4001 candidatos e confirma que, \textbf{sozinho}, o \code{max} não admite inverso; sob o \code{exp} do cósmico, ele o admite. E a prova construtiva: \textbf{do pó ao corpo, e de volta} --- a composição elementar real de um corpo humano (O 65\%, C 18,5\%, H 9,5\%, N 3,2\%, Ca 1,5\%, P 1\%), montada e depois revertida, com resíduo \textbf{exatamente \code{0}} (em $\mathbb{Q}$ --- em \code{float} seria um zero \emph{falso}). \textbf{A construção entrópica é reversível.}

\textbf{ O econômico --- a mesa.} Nenhuma operação nova: só a tríade lida no idioma mais familiar que existe. Juro simples = $\oplus$ (progressão aritmética). Juro composto = $\otimes$ (progressão geométrica). E $\prod$ = log/exp é a lente que traduz um no outro. A vitória é \textbf{ausência de arbitragem} (FP = 1): descontar desfaz capitalizar, exatamente. Quem deixa FP < 1 abre um \textbf{spread} --- e o adversário vai sangrar você um pouco a cada turno.

\textbf{ O celeste --- o par.} Dois qubits. Cada peça reparte uma unidade de energia entre \textbf{r} (potência ativa, dano solo) e \textbf{C} (concorrência, o emaranhamento), com \textbf{r$^2$ + C$^2$ = 1} sempre. \textbf{Emaranhar} duas peças acopla-as (C > 0) e por isso \textbf{baixa o dano individual} de cada uma. Ganha correlação, perde soco. E no estado \textbf{maximamente emaranhado} (r = 0, C = 1) resta um bit compartilhado --- a moeda que o mate sorteia, colapsando o par em cadeia.

\vspace{2mm}\hrule\vspace{3mm}

\section{9. Como se joga um turno --- o exemplo canônico}

Você é um Cavaleiro do corpo fractal. Precisa atravessar uma fronteira óptica e derrubar um cristal.

\textbf{1. Declare a régua.} Você está no fractal: $\mu$ = z$\cdot$\allowbreak{}z. Seu rumo é $\zeta$$_8$ (45$^\circ$), norma 1.
$\to$ \emph{legal.} (Rumo de norma $\neq$ 1 seria FP $\neq$ 1: \textbf{proibido}. Você não pode simplesmente "andar mais rápido".)

\textbf{2. Gire --- $\otimes$ La Hire.} Multiplique o rumo por $\zeta$$_8$. Cada tique do relógio $\mathbb{Z}$/8 gira 45$^\circ$. Três tiques: $\zeta$$^{4}$\ldots espere, $\zeta$ $\otimes$ $\zeta$$^3$ = $\zeta$$^{4}$ = \textbf{-1} --- você deu meia-volta. Ajuste.

\textbf{3. Avance --- $\oplus$ Clifford.} Some o deslocamento, coordenada a coordenada. Quatro \code{ADD} inteiros. Sem erro.

\textbf{4. Cruze a fronteira --- o corpo óptico assume.} Você vai de n$_1$ = 1,5 para n$\allowbreak{}_2$ = 1,0. Seu ângulo de incidência é 50$^\circ$. O ângulo crítico é arcsin(1,0/1,5) = \textbf{41,81$^\circ$}.
$\to$ 50$^\circ$ > 41,81$^\circ$: \textbf{reflexão total interna.} Você \textbf{não passa}. É devolvido inteiro. (Nenhum dado foi rolado. A régua decidiu.)

\textbf{5. Corrija.} Reduza o ângulo abaixo de $\theta$\allowbreak{}c e a travessia é \textbf{casada}: $\Gamma$ = 0, FP = 1, toda a potência transmite, sem eco.

\textbf{6. Ataque o cristal.} Sua Rainha carrega a dobra-5 (o traço $\varphi$). O cristal a \textbf{repele} --- nenhuma rede admite simetria pentagonal.
$\to$ Você troca de tática: mira um \textbf{feixe} e lê os \textbf{picos do reticulado recíproco}. O feixe de ordem zero atravessa sem dano; os desvios são o dano --- a \textbf{difração}, o resíduo.

\textbf{7. O mate suave.} Se o cristal for na verdade um \textbf{quasicristal}, os picos acumulam no centro, refração e difração colapsam, e a defesa perde o isolamento. E adivinhe quem colapsa \textbf{mais devagar} de todos: o áureo (Hurwitz). O ouro resiste até o fim.

\textbf{8. As Princesas assinam.} O turno vira lance quando o certificador lê o ponteiro:

\begin{Verbatim}[fontsize=\small,frame=single,framesep=3mm]
ponteiro = -½ -> o lance TEM DENTES (é real, é distante da trivialidade)
ponteiro = 0 -> QUEDA DE BANDEIRA -- o relógio parou, o lance é vazio
\end{Verbatim}

\vspace{2mm}\hrule\vspace{3mm}

\section{10. As três categorias --- como se \textbf{lê} um mundo}

Todo Mestre do Reino Dourado usa a semiótica de \textbf{o signo} para dizer, em uma palavra, do que um lugar é feito:

\begin{itemize}
\item \textbf{Primeiridade --- o ÍCONE (a qualidade pura, a semelhança).}
\end{itemize}
O corpo \textbf{fractal}: auto-similar, parecido consigo em toda escala. O \textbf{cristalino}: a rede idêntica a si sob os giros permitidos. Aqui a coisa simplesmente \textbf{é}.

\begin{itemize}
\item \textbf{Secundidade --- o ÍNDICE (o par, o embate, a reação factual).}
\end{itemize}
O \textbf{eletromagnético} (E $\times$ B, dois campos que se enfrentam). O \textbf{óptico} (a fronteira entre dois meios). O \textbf{celeste} (dois qubits). O \textbf{conforme} (a medida está no \textbf{par}, não no objeto: -$\tfrac12$$\cdot$|x - y|$^2$). O \textbf{racional} (numerador contra denominador). Aqui a coisa só significa \textbf{contra outra}.

\begin{itemize}
\item \textbf{Terceiridade --- o SÍMBOLO (a lei, a mediação, o hábito, a convenção).}
\end{itemize}
O \textbf{criativo} (a linguagem é convenção pura). O \textbf{econômico} (o dinheiro é convenção; o juro é a lei que medeia presente e futuro). O \textbf{relógio universal} (a régua das réguas). O \textbf{áureo} (a identidade $\varphi$$^2$ = $\varphi$ + 1, que não descreve um caso: \textbf{dobra todos}). Aqui a coisa é \textbf{regida}.

Isto não é decoração: é \textbf{informação tática}. Num mundo de secundidade, ataque o \textbf{par} (quebre a conexão). Num mundo de terceiridade, ataque a \textbf{lei} (mude a convenção e o mundo inteiro se reconfigura). Num mundo de primeiridade, não adianta atacar: \textbf{é preciso mudar de escala.}

\vspace{2mm}\hrule\vspace{3mm}

\section{11. O mate que não mata}

A vitória, no Reino Dourado, não se chama morte. Chama-se \textbf{cristalização}.

Uma partida é uma \textbf{expansão entrópica} --- a mobilidade se abre, as opções crescem. O mate é o instante em que essa expansão \textbf{colapsa a zero}: não sobra nem um lance, e o resíduo é 0. E isto \textbf{não é dissipação}: é \textbf{transferência ao dual}. O que colapsa aqui abre-se ali (\code{dS = d log V}). Nada se perde --- nem no mate. Por isso o mate \textbf{não mata}.

Cada mundo tem o seu mate, e são todos o mesmo mate:

{\fontsize{10.00pt}{12.00pt}\selectfont\begin{longtable}{>{\fontsize{10.00pt}{12.00pt}\selectfont\RaggedRight\hspace{0pt}\arraybackslash}p{9288703sp}>{\fontsize{10.00pt}{12.00pt}\selectfont\RaggedRight\hspace{0pt}\arraybackslash}p{19124641sp}}\toprule
\textbf{Mundo} & \textbf{O mate} \\ \midrule\endhead
relógio & \textbf{N = 0} --- o lapso zero, o tempo \textbf{para} \\
eletromagnético & o \textbf{cone nulo} $\sigma$ = 1 --- a régua reta \\
expansivo & a peça vira \textbf{unidade}: N(x) = $\pm$\allowbreak{}1, a órbita fecha \\
econômico & VP(A(t), t) = P --- o desconto desfez a capitalização, \textbf{exatamente} \\
áureo & toda potência reduzida a (F$_($$_n$$_-$$_1$$_)$, F$\allowbreak{}_n$) --- \textbf{o infinito foi dobrado} \\
espaço-temporal & o \textbf{transporte} chegou a zero \\
somático & a \textbf{homeostase} \\
criativo & o xeque-mate: a mobilidade $\to$ 0 \\
\bottomrule\end{longtable}}

Nada é aniquilado. A medida \textbf{fechou}. O adversário não morreu --- ele foi \textbf{cristalizado}, e o universo continua conservando tudo o que tinha.

\emph{Onde a medida fecha, o tempo para.}

\vspace{2mm}\hrule\vspace{3mm}

\section{12. As cinco regras de mesa}

\begin{enumerate}
\item \textbf{Nada vaza.} FP = 1, sempre. Um lance que cria ou destrói energia/valor/dano é ilegal, ponto final.
\item \textbf{Não arredonde.} O estado corre em pares de inteiros. O float é a tinta do pixel, não a régua.
\item \textbf{Toda régua é de um mundo.} Não misture réguas no mesmo lance.
\item \textbf{Toda afirmação exibe um refutador.} Se você diz que sua magia funciona, mostre a versão \textbf{errada} dela falhando. As Princesas só assinam com dentes.
\item \textbf{Não improvise a lei --- derive-a.} Se você não sabe o que uma peça faz, olhe a ficha do corpo. A resposta já está lá; ninguém precisa inventar.
\end{enumerate}

\vspace{2mm}\hrule\vspace{3mm}

\section{13. Comece aqui}

\begin{itemize}
\item Escolha \textbf{uma operação} para ser: $\oplus$ (o par, o Duque), $\otimes$ (o ganho, a Duquesa/a Rainha), ou $\prod$ (a lei, o Príncipe/o operador).
\item Escolha \textbf{um mundo} onde essa operação seja a mais forte.
\item Pegue \textbf{a régua desse mundo} --- e apenas ela.
\item Faça um lance.
\item Leve-o às Princesas.
\end{itemize}

Se o ponteiro marcar \textbf{-$\tfrac12$}, você jogou.

Se marcar \textbf{0} --- bandeira caída. O lance era vazio, e o reino não registra o vazio.

\textbf{Bem-vindo ao Reino Dourado. Abra um universo. Não deixe vazar.}

\vspace{2mm}\hrule\vspace{3mm}

\chapter{A simbologia ISA --- a língua das jogadas}

\begin{tcolorbox}[colback=ouroclaro!35,colframe=ouro,boxrule=0pt,leftrule=2pt,arc=0pt,left=4mm,right=3mm,top=2mm,bottom=2mm,breakable]\itshape \emph{"Toda jogada do Reino Dourado está escrita. Não em prosa: em ISA."}\end{tcolorbox}

No Reino Dourado nada acontece por descrição. Um lance não é "a Rainha ataca com força 3" --- um lance é um \textbf{programa}, escrito numa língua de doze palavras, executado numa máquina exata, e \textbf{auditável bit a bit} depois da partida. É por isso que o reino promete o que nenhum jogo promete: \textbf{resíduo 0} --- o replay de uma partida devolve, ao último bit, a mesma partida.

Esta seção ensina a ler e escrever essa língua. Ela tem três \textbf{verbos} (a tríade) e doze \textbf{palavras} (a ISA ERG-64). Aprendidos os três verbos, você lê qualquer corpo do reino --- porque todo corpo, do fractal ao econômico, do óptico ao entrópico, \textbf{conjuga os mesmos três}.

\vspace{2mm}\hrule\vspace{3mm}

\section{1. Os três verbos --- a tríade}

Antes das palavras, os verbos. Só existem três operações no reino. Todo corpo, toda peça, todo golpe é uma combinação delas.

\subsection{$\oplus$ --- CLIFFORD, a soma (o par)}

\textbf{O que é:} juntar. A corrente que se junta no nó, a superposição de dois campos, a diferença simétrica de duas linguagens, o avanço de uma peça no tabuleiro. É a operação do \textbf{par}: dois entram, um sai, e nada se perde.

\textbf{Quem é:} o \textbf{Duque $\oplus$} (Joaquim). No xadrez do corpo criativo é a \textbf{Torre} (as retas).

\textbf{Na ISA:} \code{ADD} (aritmética exata em i64) ou \code{XOR} (no corpo booleano F$_2$, onde a soma é a diferença simétrica).

\textbf{Onde você já viu:}
\begin{itemize}
\item corpo áureo $\mathbb{Z}$[$\varphi$]: \code{(a+bphi) (+) (c+dphi) = (a+c) + (b+d)phi} --- \textbf{duas ADD}, e o infinito somou-se.
\item corpo criativo: \code{L delta M} --- a diferença simétrica, com \code{LdeltaL = 0} (o resíduo nulo: cada coisa é o seu próprio inverso).
\item corpo econômico: o \textbf{juro simples} (a PA --- a mesma parcela sempre).
\item corpo entrópico: aqui $\oplus$ é o \textbf{max}, e é \emph{idempotente} --- \code{a(+)a = a}. Por isso \textbf{não há inverso}: no terreno entrópico nada se desfaz.
\end{itemize}

\subsection{$\otimes$ --- LA HIRE, o produto (a mediação)}

\textbf{O que é:} compor, amplificar, ponderar. A tensão, o ganho, o torque, o juro sobre juro.

\textbf{Quem é:} a \textbf{Rainha} (a lei $\varphi$$^2$=$\varphi$+1 é dela) e a \textbf{Duquesa $\otimes$} (Yasmin). No xadrez do corpo criativo é o \textbf{Bispo} (as diagonais).

\textbf{Na ISA:} \emph{não existe opcode MULT}. O produto \textbf{se constrói}: é a \textbf{double-and-add} --- um laço de 64 ticks que testa o bit (\code{AND}), acumula deslocado (\code{ADD}) e dobra (\code{ADD}). O produto \textbf{é a soma deslocada}. Este é o primeiro segredo da língua: $\otimes$ é feito de $\oplus$.

\textbf{Onde você já viu:}
\begin{itemize}
\item corpo racional $\mathbb{Q}$: \code{(a/\allowbreak{}b) (x) (c/\allowbreak{}d) = (a.\allowbreak{}c)/\allowbreak{}(b.\allowbreak{}d)} --- \textbf{duas MULT} (a Rainha composta duas vezes: uma no numerador, outra no denominador).
\item corpo criativo: \code{L cap M} --- a interseção; e como \code{LcapL = L}, a double-and-add \textbf{colapsa} na idempotência booleana (o \code{AND} puro).
\item corpo eletromagnético: as \textbf{impedâncias multiplicam} ao empilhar meios.
\item corpo evolutivo: a ponderação pela aptidão, \code{p\_i . w\_i}.
\end{itemize}

\subsection{$\prod$ --- PONTRYAGIN, o operador (a lei)}

\textbf{O que é:} o que \textbf{costura} $\oplus$ e $\otimes$. A fórmula é sempre a mesma: \textbf{$\prod$ = $\exp\circ\Sigma\circ\log$}. O log baixa o produto à soma; o $\Sigma$ soma; o exp sobe de volta. É o ponteiro, a poinsétia, o dispatch, a promoção, a inversão, o flip dual.

\textbf{Quem é:} o \textbf{Príncipe} (a potência, a exp) e o \textbf{Rei-gerador}. É também quem dá os \textbf{inversos} --- e por isso quem leva um mero \textbf{anel} a ser \textbf{corpo}.

\textbf{Na ISA:} \code{NOT} (que é \code{XOR} com todos-uns), o \textbf{dispatch} (\code{CMP}/\code{JZ}/\code{JNZ}/\code{JMP}), a \textbf{composição}. O opcode que escolhe qual verbo executar \textbf{é ele próprio o terceiro verbo}.

\textbf{Onde você já viu:}
\begin{itemize}
\item corpo fractal: sem $\prod$, \code{Z[zeta\_8]} é só anel (o 2 não tem inverso). Com $\prod$, \code{2 |-> ½} e \code{Q(zeta\_8)} \textbf{vira corpo}. A \textbf{promoção do peão} é exatamente isto: dar o inverso, o poder de desfazer.
\item corpo econômico: \code{log A(t) = log P + t.\allowbreak{}delta} --- o composto (PG) vira simples (PA). Capitalizar $\bowtie$ descontar.
\item corpo áureo: multiplicar por $\varphi$ é o shift \code{(a,b) |-\allowbreak{}> (b, a+b)} --- a escada de Fibonacci. \textbf{Um lance gera a órbita inteira.}
\item corpo racional: \code{[(a,b)]\textasciicircum{}-\allowbreak{}\textasciicircum{}1 = [(b,a)]} --- trocar as camadas. Puro \code{LOAD}/\code{STORE}, zero aritmética.
\end{itemize}

\subsection{A tabela dos símbolos}

{\fontsize{10.00pt}{12.00pt}\selectfont\begin{longtable}{>{\fontsize{10.00pt}{12.00pt}\selectfont\RaggedRight\hspace{0pt}\arraybackslash}p{2775748sp}>{\fontsize{10.00pt}{12.00pt}\selectfont\RaggedRight\hspace{0pt}\arraybackslash}p{4136359sp}>{\fontsize{10.00pt}{12.00pt}\selectfont\RaggedRight\hspace{0pt}\arraybackslash}p{4616404sp}>{\fontsize{10.00pt}{12.00pt}\selectfont\RaggedRight\hspace{0pt}\arraybackslash}p{7137145sp}>{\fontsize{10.00pt}{12.00pt}\selectfont\RaggedRight\hspace{0pt}\arraybackslash}p{4291556sp}>{\fontsize{10.00pt}{12.00pt}\selectfont\RaggedRight\hspace{0pt}\arraybackslash}p{3358980sp}}\toprule
\textbf{Símbolo} & \textbf{Personagem} & \textbf{O verbo} & \textbf{Na ISA ERG-64} & \textbf{o signo} & \textbf{A peça} \\ \midrule\endhead
\textbf{$\oplus$} & Duque (Joaquim) / Clifford & a \textbf{SOMA} --- o par, a corrente no nó & \code{ADD} (i64) $\cdot$ \code{XOR} (em F$_2$) & \textbf{Secundidade / Índice} (o embate, a conexão factual) & Torre \\
\textbf{$\otimes$} & Rainha $\cdot$ Duquesa (Yasmin) / La Hire & o \textbf{PRODUTO} --- tensão, ganho, torque & \textbf{MULT} = double-and-add (\code{AND}+\code{ADD}+\code{ADD}) $\cdot$ \code{AND} (em F$_2$) & a \textbf{mediação} & Bispo \\
\textbf{$\prod$} & Príncipe $\cdot$ Rei-gerador / Pontryagin & o \textbf{OPERADOR} --- $\exp\circ\Sigma\circ\log$, o que costura & \code{NOT} $\cdot$ dispatch (\code{CMP}/\code{JZ}/\code{JNZ}/\code{JMP}) $\cdot$ composição & \textbf{Terceiridade / Símbolo} (a lei, o hábito) & Rei / a promoção \\
\textbf{$\nu$} & o flip dual & o \textbf{ESPELHO} --- primal $\oplus$ dual = neutro & \code{SUB} (negar) $\cdot$ \code{XOR} com uns & involução ($\nu$$^2$=id) & o contra-lance \\
\bottomrule\end{longtable}}

\begin{tcolorbox}[colback=ouroclaro!35,colframe=ouro,boxrule=0pt,leftrule=2pt,arc=0pt,left=4mm,right=3mm,top=2mm,bottom=2mm,breakable]\itshape \textbf{A qualidade do corpo é a Primeiridade (o Ícone)} --- o que ele \emph{é} antes de reagir: a auto-similaridade do fractal, a onda plana do EM, o cristal fechado. Os três verbos são o que se \emph{faz} com ele.\end{tcolorbox}

\vspace{2mm}\hrule\vspace{3mm}

\section{2. As doze palavras --- a ISA ERG-64}

A ISA (o conjunto de instruções) é a \textbf{cristalização} dos três verbos em opcodes. Ela é \textbf{inteira} (i64) e \textbf{exata}: não há float. Nunca. Este ponto é uma \textbf{regra de mesa}, não uma preferência (\S6).

\subsection{A máquina}

Imagine o tabuleiro como uma \textbf{memória de casas numeradas} (\code{slots}) e três \textbf{mãos}:

{\fontsize{10.00pt}{12.00pt}\selectfont\begin{longtable}{>{\fontsize{10.00pt}{12.00pt}\selectfont\RaggedRight\hspace{0pt}\arraybackslash}p{5924152sp}>{\fontsize{10.00pt}{12.00pt}\selectfont\RaggedRight\hspace{0pt}\arraybackslash}p{22489192sp}}\toprule
\textbf{Peça da máquina} & \textbf{O que é} \\ \midrule\endhead
\code{A} & a mão da frente --- o último valor carregado \\
\code{B} & a mão de trás --- o valor \textbf{anterior}, empurrado para cá \\
\code{R} & o resultado --- onde toda operação deposita \\
\code{FLAGS} & a bandeira \code{Z} (zero) --- o que os saltos consultam \\
\code{mem[slot]} & as casas: \code{.data 3 7} põe o valor 7 na casa 3 \\
\bottomrule\end{longtable}}

\subsection{As doze}

{\fontsize{10.00pt}{12.00pt}\selectfont\begin{longtable}{>{\fontsize{10.00pt}{12.00pt}\selectfont\RaggedRight\hspace{0pt}\arraybackslash}p{4289100sp}>{\fontsize{10.00pt}{12.00pt}\selectfont\RaggedRight\hspace{0pt}\arraybackslash}p{13455282sp}>{\fontsize{10.00pt}{12.00pt}\selectfont\RaggedRight\hspace{0pt}\arraybackslash}p{10144674sp}}\toprule
\textbf{Palavra} & \textbf{O que faz} & \textbf{Verbo} \\ \midrule\endhead
\code{LOAD s} & \textbf{B $\leftarrow$ A; A $\leftarrow$ mem[s]} --- carrega, empurrando o anterior para trás & --- \\
\code{STORE s} & \textbf{mem[s] $\leftarrow$ R} --- grava o resultado numa casa & --- \\
\code{ADD} & \textbf{R $\leftarrow$ A + B} & $\oplus$ Clifford \\
\code{SUB} & \textbf{R $\leftarrow$ A - B} & $\nu$ (o flip / o decremento) \\
\code{AND} & \textbf{R $\leftarrow$ A $\wedge$ B} & $\otimes$ La Hire (e o teste do bit) \\
\code{OR} & \textbf{R $\leftarrow$ A $\vee$ B} & a dilatação (o dual de $\otimes$) \\
\code{XOR} & \textbf{R $\leftarrow$ A $\veebar$ B} & $\oplus$ Clifford \textbf{em F$\allowbreak{}_2$} (e $\prod$, com todos-uns) \\
\code{CMP} & acende \code{Z} \textbf{se e só se A e B forem ambos zero} & $\prod$ (o teste) \\
\code{JMP r} & salta para o rótulo \code{r} & $\prod$ (a composição) \\
\code{JZ r} & salta \textbf{se} \code{Z} está acesa & $\prod$ (o dispatch) \\
\code{JNZ r} & salta \textbf{se} \code{Z} está apagada & $\prod$ (o dispatch) \\
\code{HALT} & \textbf{o mate} --- a partida fecha & $\odot$ \\
\bottomrule\end{longtable}}

\begin{tcolorbox}[colback=ouroclaro!35,colframe=ouro,boxrule=0pt,leftrule=2pt,arc=0pt,left=4mm,right=3mm,top=2mm,bottom=2mm,breakable]\itshape \textbf{O \code{\#} do xadrez é o \code{HALT} da ISA.} Na notação algébrica, \code{Qxf7\#} termina a partida; na ISA, \code{HALT} termina o programa. É o \textbf{mesmo signo}, e não por acaso: o mate É a cristalização (resíduo $\to$ 0).\end{tcolorbox}

\subsection{As duas surpresas da língua}

\begin{enumerate}
\item \textbf{Não há MULT.} O produto se \textbf{constrói} (double-and-add). $\otimes$ é $\oplus$ deslocado.
\item \textbf{\code{CMP} não compara dois valores} --- ela acende \code{Z} quando \textbf{A e B são ambos zero}. Para testar "x é zero?", carregue \code{x} e carregue a \textbf{casa do zero} (por convenção \code{slot 0}, \code{.data 0 0}). Se \code{Z} acende, x era zero. É assim que o reino inteiro pergunta \emph{"resíduo 0?"}.
\end{enumerate}

\vspace{2mm}\hrule\vspace{3mm}

\section{3. Como LER uma jogada}

Quatro regras. Decore-as e você lê qualquer corpo.

\textbf{Regra 1 --- O último \code{LOAD} é A; o penúltimo desceu para B.}
\begin{Verbatim}[fontsize=\small,frame=single,framesep=3mm]
LOAD 1 ; A <- mem[1]
LOAD 3 ; B <- mem[1] (empurrado), A <- mem[3]
\end{Verbatim}

\textbf{Regra 2 --- A operação lê \code{R <- A op B}.} Para \code{ADD}/\code{AND}/\code{OR}/\code{XOR} a ordem não importa (são comutativas). Para \textbf{\code{SUB} importa}:
\begin{tcolorbox}[colback=ouroclaro!35,colframe=ouro,boxrule=0pt,leftrule=2pt,arc=0pt,left=4mm,right=3mm,top=2mm,bottom=2mm,breakable]\itshape \code{R = A -\allowbreak{} B = (o último carregado) -\allowbreak{} (o penúltimo)}. Então \code{LOAD 5} (a casa do 1) seguido de \code{LOAD 31} (o contador) e \code{SUB} significa \textbf{contador - 1}. Leia sempre de trás para a frente.\end{tcolorbox}

\textbf{Regra 3 --- \code{STORE} fecha o lance.} O que estava em \code{R} vira posição no tabuleiro. Enquanto não há \code{STORE}, nada aconteceu.

\textbf{Regra 4 --- \code{CMP} + \code{JZ} é o dispatch de Pontryagin.} Sempre que vir \code{LOAD x} / \code{LOAD 0} / \code{CMP} / \code{JZ ...}, leia: \emph{"se x é zero, siga por ali"}. É o $\prod$ escolhendo o ramo.

\vspace{2mm}\hrule\vspace{3mm}

\section{4. Exemplos --- lendo lances de verdade}

\subsection{Exemplo A --- o lance do Duque: $\oplus$ no corpo áureo}

O corpo áureo $\mathbb{Z}$[$\varphi$] escreve um número como um \textbf{par de inteiros} \code{(a,b) = a + bphi} --- dois trilhos: o racional e o áureo. Somar é somar cada trilho.

\begin{Verbatim}[fontsize=\small,frame=single,framesep=3mm]
.data 0 0 ; a casa do ZERO (o espelho do CMP)
.data 1 a ; trilho racional do primeiro
.data 2 b ; trilho áureo do primeiro
.data 3 c
.data 4 d
.data 50 0 ; onde cai (a+c)
.data 51 0 ; onde cai (b+d)

SOMA: LOAD 1 ; A <- a
 LOAD 3 ; B <- a, A <- c
 ADD ; R <- a + c (+) Clifford, trilho racional
 STORE 50
 LOAD 2
 LOAD 4
 ADD ; R <- b + d (+) Clifford, trilho áureo
 STORE 51
 HALT ; # -- o lance fechou
\end{Verbatim}

\textbf{Leitura:} \emph{duas ADD e o infinito somou-se.} É a operação mais barata do reino --- e é \textbf{exata}: nenhum arredondamento, nenhum ulp perdido.

\subsection{Exemplo B --- a carga da Rainha: $\otimes$ pela double-and-add}

Aqui está o coração da língua. Não existe MULT: a Rainha \textbf{executa 64 micro-lances}. Este é o bloco real do corpo áureo e do corpo racional (o \code{\_mult\_block}, \code{acc += aa.mor}):

\begin{Verbatim}[fontsize=\small,frame=single,framesep=3mm]
M1: LOAD 3 ; mor (o multiplicador)
 LOAD 21 ; mask (1, 2, 4, ... 2^63)
 AND ; (x) -- o bit da vez está aceso?
 STORE 60
 LOAD 60
 LOAD 0
 CMP ; (op) -- o bit é zero?
 JZ M1S ; se sim, pule a soma
 LOAD 40 ; acc
 LOAD 11 ; aa (o parcial deslocado)
 ADD ; (+) -- acumula
 STORE 40
M1S: LOAD 21
 LOAD 21
 ADD ; (+) -- dobra a máscara (o "log": sobe uma escala)
 STORE 21
 LOAD 11
 LOAD 11
 ADD ; (+) -- dobra o parcial (o deslocamento)
 STORE 11
 LOAD 5 ; a casa do 1
 LOAD 31 ; o contador (64)
 SUB ; R <- 31 - 1 (Regra 2!)
 STORE 31
 LOAD 31
 LOAD 0
 CMP ; deu zero? 64 ticks cumpridos?
 JZ M2 ; então passe ao próximo produto
 JMP M1 ; (op) -- a composição: mais um tick
\end{Verbatim}

\textbf{Leitura:} \code{(x) = Sigma o deslocamento}. É \textbf{$\exp\circ\Sigma\circ\log$ em i64}: o log é o bit (a escala 2$^k$), o $\Sigma$ são os \code{ADD} acumulados, o exp é a máscara dobrando. A Rainha é feita do Duque --- \textbf{La Hire é Clifford repetido}, e é literalmente isso que o reino executa.

O relógio deste laço é o \textbf{$\Sigma$ = $\mathbb{Z}$/n do contrato}: uma órbita de \textbf{64 ticks}.

\subsection{Exemplo C --- o dispatch de $\prod$: o corpo criativo inteiro}

O corpo criativo (o anel booleano F$_2$, o corpo das oito linguagens) é o mais direto do reino: \textbf{a sua tríade É a ISA}. Um único programa, e o \emph{opcode} (a casa 9) escolhe o verbo:

\begin{Verbatim}[fontsize=\small,frame=single,framesep=3mm]
DISP: LOAD 9
 LOAD 0
 CMP ; op == 0 ?
 JZ XORO ; -> (+) Clifford
 LOAD 9
 LOAD 8 ; a casa do 1
 SUB ; R <- 1 - op (zero exatamente quando op = 1)
 STORE 11
 LOAD 11
 LOAD 0
 CMP
 JZ ANDO ; -> (x) La Hire
NOTO: LOAD 1
 LOAD 10 ; a casa dos TODOS-UNS (0xFFFF...FF)
 XOR ; (op) -- o NOT: o complemento, o dual
 STORE 3
 HALT
ANDO: LOAD 1
 LOAD 2
 AND ; (x) -- a interseção, o ganho
 STORE 3
 HALT
XORO: LOAD 1
 LOAD 2
 XOR ; (+) -- a diferença simétrica, a corrente no nó
 STORE 3
 HALT
\end{Verbatim}

\textbf{Leitura:} repare no que acontece. O \textbf{dispatch por opcode É Pontryagin} --- ele não é um detalhe de implementação, ele é o terceiro verbo em pessoa. E o \code{NOT} não é uma palavra nova: é \code{XOR} com todos-uns. A tríade fecha em si mesma.

E note a lei do corpo, executável: \code{x (+) x = 0} (o cristal, o resíduo nulo) --- \textbf{cada peça é o seu próprio inverso}.

\subsection{Exemplo D --- lendo uma jogada de tabuleiro}

Traduzindo a mecânica em ISA, um turno típico do corpo fractal (o tabuleiro-raiz, o relógio $\mathbb{Z}$/8 das oito direções $\zeta$$_8$):

{\fontsize{10.00pt}{12.00pt}\selectfont\begin{longtable}{>{\fontsize{10.00pt}{12.00pt}\selectfont\RaggedRight\hspace{0pt}\arraybackslash}p{11038684sp}>{\fontsize{10.00pt}{12.00pt}\selectfont\RaggedRight\hspace{0pt}\arraybackslash}p{7938081sp}>{\fontsize{10.00pt}{12.00pt}\selectfont\RaggedRight\hspace{0pt}\arraybackslash}p{8912291sp}}\toprule
\textbf{A jogada} & \textbf{O verbo} & \textbf{Na ISA} \\ \midrule\endhead
\textbf{Girar} a peça 45$^\circ$ (um tique do relógio) & $\otimes$ multiplicar o rumo por $\zeta$$_8$ & a convolução ciclotômica: 4 MULT + \code{SUB} (o sinal $\zeta$$^{4}$ = -1) \\
\textbf{Avançar} (empilhar deslocamento) & $\oplus$ somar coordenada a coordenada & 4 $\times$ \code{ADD} \\
\textbf{Promover} o peão (dar-lhe o inverso: o poder de desfazer) & $\prod$ o operador & \code{SUB}/negate (para as unidades, \code{zeta\textasciicircum{}-\allowbreak{}\textasciicircum{}1 = -\allowbreak{}zeta\textasciicircum{}3}, exato) \\
\textbf{Espelhar} o rumo do inimigo & $\nu$ a conjugação \code{z |-> z} & negar a parte imaginária (\code{SUB}) \\
\textbf{Medir o dano} (a régua!) & $\otimes$ com o dual: $z \otimes \bar z = |z|^2$ & a norma --- MULT + ADD \\
\textbf{Legalidade do lance} & \code{$|z|=1$} & \code{CMP} contra 1 --- fora disso, \textbf{vaza} \\
\bottomrule\end{longtable}}

Aquele último item não é decoração. É \textbf{a lei} (\S6).

\vspace{2mm}\hrule\vspace{3mm}

\section{5. O cartão de bolso --- o mesmo verbo, onze idiomas}

A grande economia do Reino Dourado: você não aprende onze jogos. Aprende \textbf{três verbos}, e cada corpo os pronuncia com seu sotaque.

{\fontsize{10.00pt}{12.00pt}\selectfont\begin{longtable}{>{\fontsize{10.00pt}{12.00pt}\selectfont\RaggedRight\hspace{0pt}\arraybackslash}p{6225826sp}>{\fontsize{10.00pt}{12.00pt}\selectfont\RaggedRight\hspace{0pt}\arraybackslash}p{6512938sp}>{\fontsize{10.00pt}{12.00pt}\selectfont\RaggedRight\hspace{0pt}\arraybackslash}p{6313876sp}>{\fontsize{10.00pt}{12.00pt}\selectfont\RaggedRight\hspace{0pt}\arraybackslash}p{8312128sp}}\toprule
\textbf{Corpo} & \textbf{$\oplus$ CLIFFORD (a soma)} & \textbf{$\otimes$ LA HIRE (o produto)} & \textbf{$\prod$ PONTRYAGIN (o operador)} \\ \midrule\endhead
\textbf{fractal} $\mathbb{Q}$($\zeta$$_8$) & soma nas 4 coordenadas & convolução ciclotômica ($\zeta$$^{4}$=-1) & exp/log: dá os inversos --- \textbf{anel $\to$ corpo} \\
\textbf{criativo} F$_2$ & \code{delta} diferença simétrica (\code{XOR}) & \code{cap} interseção (\code{AND}) & \code{NOT} / o dispatch \\
\textbf{áureo} $\mathbb{Z}$[$\varphi$] & duas \code{ADD} & quatro MULT + $\varphi$$^2$=$\varphi$+1 & o shift \code{(a,b)|-\allowbreak{}>(b,a+b)} --- Fibonacci \\
\textbf{racional} $\mathbb{Q}$ & 3 MULT + 1 \code{ADD} & 2 MULT & trocar as camadas: \code{[(a,b)]\textasciicircum{}-\allowbreak{}\textasciicircum{}1=[(b,a)]} \\
\textbf{econômico} & juro \textbf{simples} (PA) & juro \textbf{composto} (PG) & \code{log}/\code{exp}: capitalizar $\bowtie$ descontar \\
\textbf{eletromagnético} & superposição dos perfis & impedâncias multiplicam & \code{sigma = e\textasciicircum{}(-\allowbreak{}Deltah/\allowbreak{}2)} --- o log lineariza \\
\textbf{motor} & soma dos geradores & o \textbf{torque} $\tau$ = $\psi$$\times$\allowbreak{}i & \code{exp(tG)} --- o DTC, o controle ótimo \\
\textbf{relógio} & soma relativística \code{(s\_\allowbreak{}1+s\_\allowbreak{}2)/\allowbreak{}(1+s\_\allowbreak{}1s\_\allowbreak{}2)} & composição das rapidezes & \code{Lambda = artanh} --- o endireitador \\
\textbf{entrópico} & \textbf{max} (idempotente --- \emph{sem inverso!}) & a \textbf{adição} real (a MULT tropical \textbf{é} o \code{ADD}) & a deflexão \code{D\_\allowbreak{}lambda(x) = x+lambda} \\
\textbf{evolutivo} & reprodução \code{(p+q)/2} & ponderação pela aptidão \code{p.w} & o \textbf{replicador} (a seleção) \\
\textbf{cristalino} & soma pontual & rotação de ordem 4 ou 6 & a conjugação (\code{SUB} no imaginário) \\
\bottomrule\end{longtable}}

Leia a tabela por colunas e o reino se abre: \textbf{é sempre o mesmo lance}.

\vspace{2mm}\hrule\vspace{3mm}

\section{6. As três leis da língua (regras de mesa, inegociáveis)}

\textbf{Lei I --- O EXATO. Nenhum float na régua.}
O estado da partida corre em \textbf{inteiros} (i64: pares \code{(num, den)} em $\mathbb{Q}$, pares \code{(a,b)} em $\mathbb{Z}$[$\varphi$]). O float \textbf{trunca} --- 1 ulp de energia perdida por operação. Arredondar no meio de uma partida é \textbf{vazar}, e vazar é \textbf{lance ilegal}. O float aparece \textbf{só no fim}, ao pintar o pixel, onde 1 ulp é invisível.

\textbf{Lei II --- O RESÍDUO 0. Toda jogada é certificada.}
Cada corpo é escrito \textbf{uma vez} na ISA, e o \code{chessc} (o compilador universal, o \emph{broca-so}) emite dela o \textbf{WASM} (o metal, o corpo roda) e o \textbf{Dafny} (as leis, provadas). A costura fecha: \textbf{interp $\equiv$ WASM $\equiv$ Dafny}. Quem gera é o \textbf{corpo criativo}; quem atesta são as \textbf{Princesas} (o corpo técnico, o kernel LCF). Uma afirmação sem refutador não entra no reino.

\textbf{Lei III --- FP = 1. Nenhum universo vaza.}
O casamento de impedância \code{Gamma = (Z-\allowbreak{}Z\_\allowbreak{}0)/\allowbreak{}(Z+Z\_\allowbreak{}0) = 0} --- no corpo eletromagnético, o \textbf{cone nulo $\sigma$=1}; no fractal, a \textbf{norma unitária |z|=1}; no telescópico, o \textbf{traço zero}; no econômico, a \textbf{ausência de arbitragem}. Um lance com FP $\neq$ 1 é energia surgindo fora do Sol: \textbf{proibido}. Não é sabor narrativo --- é o critério de entrada de um corpo no reino (o Contrato: $\Sigma$ = $\mathbb{Z}$/n, a régua $\mu$, o flip $\nu$, o mate $\odot$).

\vspace{2mm}\hrule\vspace{3mm}

\section{7. A língua tem dois registros}

\begin{itemize}
\item \textbf{Falado --- SAN (a notação algébrica).} \code{game f3 e5 g4 Qh4\#}. É como um jogador anuncia o lance à mesa.
\item \textbf{Escrito --- ISA (o pseudo-assembly ERG-64).} É como o reino \textbf{executa} o lance.
\end{itemize}

O \code{chessc} traduz de um ao outro: \code{compile} (SAN $\to$ ISA), \code{run} (interpreta, com certificado), \code{emit-wasm} (o metal), \code{emit-dafny} (as leis). \textbf{Você joga em SAN; o reino confere em ISA.}

\vspace{2mm}\hrule\vspace{3mm}

\section{8. A leitura semiótica (o signo, para quem quer o fundo)}

A tríade não é uma escolha estética. Ela é a semiótica posta como máquina:

\begin{itemize}
\item \textbf{Primeiridade / Ícone} --- a \emph{qualidade} do corpo, antes de qualquer reação: a auto-similaridade do fractal (o tabuleiro que se parece consigo em toda escala), a onda plana no vácuo, o cristal fechado.
\item \textbf{Secundidade / Índice} --- \textbf{$\oplus$}, o par, o embate, a conexão factual: dois campos que reagem, dois dígitos que geram um transporte, dois qubits emaranhados.
\item \textbf{Terceiridade / Símbolo} --- \textbf{$\prod$}, a lei, o hábito, a convenção: o dispatch, o operador que medeia, a linguagem (o corpo criativo é convenção pura --- por isso é o corpo do símbolo).
\end{itemize}

E \textbf{$\otimes$} é a \textbf{mediação} entre os dois: o produto que compõe o par sob a lei.

\vspace{2mm}\hrule\vspace{3mm}

\section{9. Exercícios (aprenda a ler)}

\textbf{1.} Leia e diga o valor de \code{R}:
\begin{Verbatim}[fontsize=\small,frame=single,framesep=3mm]
LOAD 5 ; mem[5] = 1
LOAD 31 ; mem[31] = 64
SUB
\end{Verbatim}
\emph{(Resposta: \code{R = A -\allowbreak{} B = 64 -\allowbreak{} 1 = 63}. O último carregado é A. Regra 2.)}

\textbf{2.} O que este trecho pergunta?
\begin{Verbatim}[fontsize=\small,frame=single,framesep=3mm]
LOAD 60
LOAD 0
CMP
JZ ADIANTE
\end{Verbatim}
\emph{(Resposta: "mem[60] é zero?" --- e se for, salta. É o $\prod$ despachando; é também a pergunta que o reino inteiro faz: \textbf{resíduo 0?})}

\textbf{3.} Como se escreve o \code{NOT} de \code{mem[1]} sem opcode \code{NOT}?
\emph{(Resposta: \code{LOAD 1} / \code{LOAD 10} / \code{XOR}, onde \code{mem[10] = 0xFFFFFFFFFFFFFFFF}. O complemento é $\oplus$ com todos-uns --- $\prod$ nasce de $\oplus$.)}

\textbf{4.} Um lance seu tem razão \code{3/4}. Qual é a defesa exata do adversário?
\emph{(Resposta: a \textbf{recíproca} \code{4/3} --- o flip $\nu$. Primal $\otimes$ dual = 1, o neutro. \code{Gamma=0}: o golpe casa e \textbf{nada reflete}. É o inverso que faz de $\mathbb{Q}$ um corpo: em $\mathbb{Z}$, esse ataque não teria defesa.)}

\vspace{2mm}\hrule\vspace{3mm}

\begin{tcolorbox}[colback=ouroclaro!35,colframe=ouro,boxrule=0pt,leftrule=2pt,arc=0pt,left=4mm,right=3mm,top=2mm,bottom=2mm,breakable]\itshape \textbf{O resumo em uma linha:} $\oplus$ junta, $\otimes$ compõe, $\prod$ costura --- e as doze palavras da ISA são só a maneira exata de dizer isso. \emph{Não há float, não há arredondamento, não há vazamento. Toda jogada fecha em resíduo 0, ou não é jogada.}\end{tcolorbox}

\vspace{2mm}\hrule\vspace{3mm}

\chapter{A semiótica --- os três signos do universo}

\begin{tcolorbox}[colback=ouroclaro!35,colframe=ouro,boxrule=0pt,leftrule=2pt,arc=0pt,left=4mm,right=3mm,top=2mm,bottom=2mm,breakable]\itshape \emph{"O tabuleiro não te mostra o corpo. Mostra o signo do corpo. Quem sabe ler, joga; quem não sabe, chuta."} --- anotação de margem no Contrato\end{tcolorbox}

\section{Por que um manual de jogo precisa de semiótica}

Você nunca vê um corpo. Você vê o que ele \textbf{faz aparecer}: um brilho que se repete em toda escala; um ponteiro que marca \textbf{-$\tfrac12$}; uma lei que fecha uma cadeia infinita em dois inteiros. Isso são \textbf{signos} --- e o Reino Dourado tem exatamente \textbf{três}, nem um a mais.

Charles Sanders o signo chamou-os de suas \textbf{três categorias}. O reino as adota como \textbf{regra de mesa}, porque elas resolvem o problema prático do jogador: \emph{como sei o que este corpo faz sem abrir a sua álgebra?}

{\fontsize{10.00pt}{12.00pt}\selectfont\begin{longtable}{>{\fontsize{10.00pt}{12.00pt}\selectfont\RaggedRight\hspace{0pt}\arraybackslash}p{5647487sp}>{\fontsize{10.00pt}{12.00pt}\selectfont\RaggedRight\hspace{0pt}\arraybackslash}p{6431503sp}>{\fontsize{10.00pt}{12.00pt}\selectfont\RaggedRight\hspace{0pt}\arraybackslash}p{6188537sp}>{\fontsize{10.00pt}{12.00pt}\selectfont\RaggedRight\hspace{0pt}\arraybackslash}p{3541521sp}>{\fontsize{10.00pt}{12.00pt}\selectfont\RaggedRight\hspace{0pt}\arraybackslash}p{5031432sp}}\toprule
\textbf{Categoria} & \textbf{Signo} & \textbf{O que é} & \textbf{Quantos precisa} & \textbf{Pergunta que responde} \\ \midrule\endhead
\textbf{Primeiridade} & \textbf{ÍCONE} (semelhança) & a qualidade pura, sozinha, sem referência a nada & \textbf{um} & \emph{O que ele é?} \\
\textbf{Secundidade} & \textbf{ÍNDICE} (conexão factual) & a reação, o embate, o par & \textbf{dois} & \emph{Contra quem ele é?} \\
\textbf{Terceiridade} & \textbf{SÍMBOLO} (convenção) & a lei, a mediação, o hábito & \textbf{três} & \emph{Sob que lei ele é?} \\
\bottomrule\end{longtable}}

Regra de ouro da leitura:

\begin{itemize}
\item o \textbf{ícone} não aponta para fora --- ele \textbf{se parece consigo};
\item o \textbf{índice} não existe sozinho --- ele \textbf{exige o outro};
\item o \textbf{símbolo} vale \textbf{mesmo quando o objeto está ausente} --- ele é a lei que rege.
\end{itemize}

\vspace{2mm}\hrule\vspace{3mm}

\section{1. Como os três signos MODULAM a tríade}

A tríade do reino é sempre a mesma, em todo corpo, em toda partida:

{\fontsize{10.00pt}{12.00pt}\selectfont\begin{longtable}{>{\fontsize{10.00pt}{12.00pt}\selectfont\RaggedRight\hspace{0pt}\arraybackslash}p{7838210sp}>{\fontsize{10.00pt}{12.00pt}\selectfont\RaggedRight\hspace{0pt}\arraybackslash}p{5376434sp}>{\fontsize{10.00pt}{12.00pt}\selectfont\RaggedRight\hspace{0pt}\arraybackslash}p{6792458sp}>{\fontsize{10.00pt}{12.00pt}\selectfont\RaggedRight\hspace{0pt}\arraybackslash}p{7357666sp}}\toprule
\textbf{Operação} & \textbf{Personagem} & \textbf{Categoria} & \textbf{Na ISA ERG-64} \\ \midrule\endhead
\textbf{$\oplus$ CLIFFORD} --- a soma, a corrente no nó, \textbf{o par} & o Duque (Joaquim) & \textbf{Secundidade / ÍNDICE} --- dois que se encontram no nó & \code{ADD} (ou \code{XOR} no booleano) \\
\textbf{$\otimes$ LA HIRE} --- o produto, a tensão/ganho & a Rainha / a Duquesa (Yasmin) & \textbf{a MEDIAÇÃO} --- o que compõe um no outro & \code{MULT} por \emph{double-and-add} (\code{ADD} deslocado), ou \code{AND} \\
\textbf{$\prod$ PONTRYAGIN} --- o operador que costura, $\exp\circ\Sigma\circ\log$ & o Príncipe / o ponteiro, a poinsétia & \textbf{Terceiridade / SÍMBOLO} --- a lei que medeia $\oplus$ e $\otimes$ & \code{NOT} / o \textbf{dispatch} (\code{CMP},\code{JZ},\code{JMP}) / a composição \\
\emph{(a qualidade do próprio corpo)} & --- & \textbf{Primeiridade / ÍCONE} & o dado nu no slot: o bit, o i64 \\
\bottomrule\end{longtable}}

Repare no que isso significa \textbf{no metal}. No corpo criativo, a ISA guarda o opcode no \textbf{slot 9} (\code{0=(+)/XOR}, \code{1=(x)/AND}, \code{2=(op)/NOT}): \textbf{o próprio dispatch por opcode É Pontryagin} --- a lei escolhendo qual operação vai correr. A Terceiridade não é filosofia: é o \code{CMP}+\code{JZ} que decide.

\begin{tcolorbox}[colback=ouroclaro!35,colframe=ouro,boxrule=0pt,leftrule=2pt,arc=0pt,left=4mm,right=3mm,top=2mm,bottom=2mm,breakable]\itshape \textbf{A dupla leitura (a chave do capítulo).} Todo corpo tem a tríade inteira dentro dele --- sempre com essa mesma distribuição. Mas cada corpo tem também \textbf{uma categoria dominante}: a face que ele \textbf{vira para o jogador}. É essa que você lê na mesa. A tríade é o interior; a categoria dominante é a \textbf{cor} do corpo.\end{tcolorbox}

\vspace{2mm}\hrule\vspace{3mm}

\section{2. A carta dos signos --- o dominante de cada corpo}

Esta é a tabela que se consulta \textbf{antes de entrar em campo}.

\subsection{ÍCONE (Primeiridade) --- os corpos da qualidade}

{\fontsize{10.00pt}{12.00pt}\selectfont\begin{longtable}{>{\fontsize{10.00pt}{12.00pt}\selectfont\RaggedRight\hspace{0pt}\arraybackslash}p{5150592sp}>{\fontsize{10.00pt}{12.00pt}\selectfont\RaggedRight\hspace{0pt}\arraybackslash}p{12509977sp}>{\fontsize{10.00pt}{12.00pt}\selectfont\RaggedRight\hspace{0pt}\arraybackslash}p{10228487sp}}\toprule
\textbf{Corpo} & \textbf{O ícone} & \textbf{O que isso te dá em jogo} \\ \midrule\endhead
\textbf{fractal} $\mathbb{Q}$($\zeta$$_8$), $\Sigma$=$\mathbb{Z}$/8 & \textbf{auto-similar}: o tabuleiro que se parece consigo em toda escala (tesselação nativa, as bacias de Newton) & dar zoom \textbf{não muda as regras}; o mesmo lance vale em qualquer escala \\
\textbf{cristalino} $\mathbb{Q}$($\sqrt{\ }$\allowbreak{}D), D<0 & a \textbf{rede congelada}: regulador zero, órbita fechada, o relógio parado --- a rede idêntica a si sob os giros permitidos & defesa que \textbf{não vaza e não expande}; só cabem giros de ordem \{1,2,3,4,6\} \\
\bottomrule\end{longtable}}

\subsection{ÍNDICE (Secundidade) --- os corpos do par}

{\fontsize{10.00pt}{12.00pt}\selectfont\begin{longtable}{>{\fontsize{10.00pt}{12.00pt}\selectfont\RaggedRight\hspace{0pt}\arraybackslash}p{7228026sp}>{\fontsize{10.00pt}{12.00pt}\selectfont\RaggedRight\hspace{0pt}\arraybackslash}p{12189362sp}>{\fontsize{10.00pt}{12.00pt}\selectfont\RaggedRight\hspace{0pt}\arraybackslash}p{8471668sp}}\toprule
\textbf{Corpo} & \textbf{O índice} & \textbf{O que isso te dá em jogo} \\ \midrule\endhead
\textbf{conforme} (a casa da medida) & X(x)$^2$=0 --- \textbf{norma nula}: a obra sozinha não significa nada; a medida está no \textbf{emparelhamento} X$\cdot$\allowbreak{}Y = -$\tfrac12$$\|$\allowbreak{}x-y$\|$$^2$ & o \textbf{ponteiro do relógio de xadrez}: só o par mede \\
\textbf{eletromagnético} (E,B) & o \textbf{embate de dois campos}: E$\times$\allowbreak{}B, o telescópio & o golpe casado no \textbf{cone nulo $\sigma$=1} (FP=1, nada vaza) \\
\textbf{motor} & o \textbf{torque} $\tau$=$\psi$$\times$\allowbreak{}i, o par ortogonal fluxo$\times$\allowbreak{}corrente; o dipolo branco/negro & o golpe de flanco \textbf{a 90$^\circ$}, não frontal \\
\textbf{celeste} (2 qubits) & a \textbf{concorrência} C: o emaranhamento entre dois & r$^2$+C$^2$=1 --- ganha correlação, perde dano solo \\
\textbf{óptico} (a fronteira) & a \textbf{interface}: dois meios n$_1$,n$\allowbreak{}_2$ reagindo no bordo & a refração; e a \textbf{reflexão total interna} além de $\theta$\allowbreak{}c \\
\textbf{entrópico} (max,+) & o \textbf{max}: o par que colapsa no dominante; o sorteio obrigatório em cada vértice & dano que \textbf{só acumula} --- não há desfazer \\
\textbf{racional $\mathbb{Q}$} (num,den) & duas multiplicidades em \textbf{embate}: a razão a/b é a solução de b$\otimes$\allowbreak{}x=a & o contra-golpe pela \textbf{recíproca} (primal $\otimes$ dual = 1) \\
\bottomrule\end{longtable}}

\subsection{SÍMBOLO (Terceiridade) --- os corpos da lei}

{\fontsize{10.00pt}{12.00pt}\selectfont\begin{longtable}{>{\fontsize{10.00pt}{12.00pt}\selectfont\RaggedRight\hspace{0pt}\arraybackslash}p{5984664sp}>{\fontsize{10.00pt}{12.00pt}\selectfont\RaggedRight\hspace{0pt}\arraybackslash}p{12969891sp}>{\fontsize{10.00pt}{12.00pt}\selectfont\RaggedRight\hspace{0pt}\arraybackslash}p{8934501sp}}\toprule
\textbf{Corpo} & \textbf{O símbolo} & \textbf{O que isso te dá em jogo} \\ \midrule\endhead
\textbf{criativo} (as 8 linguagens, anel booleano F$_2$) & a \textbf{convenção}: oito substratos, um só estado (a IR ERG-64) & é o \textbf{gerador} --- joga a partida, e o mate é a cristalização \\
\textbf{relógio universal} $\Sigma$=$\mathbb{Z}$/n & a \textbf{régua-de-réguas}: a órbita que fecha em $\omega$$^n$=1 e conserva a norma para \textbf{todo} k & o tique reversível; o único tique sem inverso é o \textbf{mate} \\
\textbf{áureo $\mathbb{Z}$[$\varphi$]} & a lei \textbf{$\varphi$$^2$=$\varphi$+1} --- dobra \textbf{toda} potência de volta em dois inteiros & o eco infinito resolvido \textbf{num tick}, resíduo 0 \\
\textbf{econômico} (os juros) & a lei do valor: $\delta$=log(1+r), o hábito que rege presente$\leftrightarrow$\allowbreak{}futuro & capitalizar $\bowtie$ descontar; \textbf{sem arbitragem} = FP=1 \\
\textbf{expansivo} $\mathbb{Q}$($\sigma_m$) & "\textbf{três nomes, um número}": taxa de expansão = regulador = entropia topológica & a marcha do Rei que cresce a taxa constante $\sigma_m$ \\
\textbf{evolutivo} & o \textbf{replicador} (a seleção como lei; Fisher: $\langle$\allowbreak{}w$\rangle$ só cresce) & a escalada monótona do exército \\
\textbf{geométrico} (curvatura) & a \textbf{conexão}: a geodésica é o hábito do campo; Bianchi, a lei conservada & pousar massa \textbf{dobra a régua} de quem passa \\
\textbf{somático} (o tecido) & a \textbf{homeostase} e a diferenciação por classes (o co-estado G$^T$) & mover \textbf{por tipo} (O(T)), não célula a célula (O(N)) \\
\textbf{mórfico} (F$_2$$^n$, $\oplus$, $\wedge$) & a \textbf{adjunção} $\delta$$\dashv$$\varepsilon$ posta como lei & erosão/dilatação: consolidar núcleo, avançar território \\
\textbf{espaço-temporal} $\mathbb{Q}$$_2$ & o \textbf{tempo é a valuação v} --- a lei que conta os transportes & (espaço = ícone, o dígito; carry = índice) \\
\textbf{telescópico} (os perfis) & a norma que \textbf{telescopa}: N(h$\odot$\allowbreak{}g)=N(h)$\cdot$\allowbreak{}N(g) --- decompor sem perder nem acrescentar & atacar por muitos eixos = \textbf{somar perfis} \\
\bottomrule\end{longtable}}

\vspace{2mm}\hrule\vspace{3mm}

\section{3. A jogada de LEITURA (ação livre, uma por turno)}

Antes de mover, você pode \textbf{ler}. A leitura é o protocolo das três perguntas --- e ela é literalmente o \textbf{Contrato} lido em voz alta:

{\fontsize{10.00pt}{12.00pt}\selectfont\begin{longtable}{>{\fontsize{10.00pt}{12.00pt}\selectfont\RaggedRight\hspace{0pt}\arraybackslash}p{3196413sp}>{\fontsize{10.00pt}{12.00pt}\selectfont\RaggedRight\hspace{0pt}\arraybackslash}p{5308620sp}>{\fontsize{10.00pt}{12.00pt}\selectfont\RaggedRight\hspace{0pt}\arraybackslash}p{4569330sp}>{\fontsize{10.00pt}{12.00pt}\selectfont\RaggedRight\hspace{0pt}\arraybackslash}p{6513112sp}>{\fontsize{10.00pt}{12.00pt}\selectfont\RaggedRight\hspace{0pt}\arraybackslash}p{7253005sp}}\toprule
\textbf{\#} & \textbf{Pergunta} & \textbf{Signo} & \textbf{Onde está no Contrato} & \textbf{O que você ganha} \\ \midrule\endhead
\textbf{1} & \emph{Com o que ele \textbf{se parece}?} & \textbf{ÍCONE} & a órbita \textbf{$\Sigma$=$\mathbb{Z}$/n} e a \textbf{régua $\mu$ ($T_\omega$)} & a \textbf{assinatura}: como o corpo mede \\
\textbf{2} & \emph{Contra quem ele \textbf{bate}?} & \textbf{ÍNDICE} & o \textbf{flip dual $\nu$} (primal + dual = o neutro) & o \textbf{par}: onde ele cancela e onde ele soma \\
\textbf{3} & \emph{Que lei o \textbf{autoriza}?} & \textbf{SÍMBOLO} & o \textbf{operador $\prod$} e o \textbf{mate $\circ$} & a \textbf{legalidade}: se o lance fecha ou vaza \\
\bottomrule\end{longtable}}

\textbf{Regra dura.} Um lance que não responde às três \textbf{VAZA} --- $\Gamma$=(Z-Z$\allowbreak{}_0$)/(Z+Z$\allowbreak{}_0$) $\neq$ 0, FP $\neq$ 1 --- e \textbf{lance que vaza é ilegal}. Nenhum universo do reino pode ter fator de potência diferente de 1.

\emph{Exemplo de leitura, em voz de mesa:}
\begin{tcolorbox}[colback=ouroclaro!35,colframe=ouro,boxrule=0pt,leftrule=2pt,arc=0pt,left=4mm,right=3mm,top=2mm,bottom=2mm,breakable]\itshape "É o \textbf{fractal}. \textbf{Ícone}: auto-similar, relógio $\mathbb{Z}$/8, régua z$\cdot$\allowbreak{}z --- mede pela norma. \textbf{Índice}: o flip é a conjugação, z$\otimes$\allowbreak{}z=|z|$^2$ é o meu dano. \textbf{Símbolo}: $\prod$ dá o inverso, então posso \textbf{desfazer} o lance. E só rumo de norma |z|=1 é legal --- fora disso eu vazo."\end{tcolorbox}

\vspace{2mm}\hrule\vspace{3mm}

\section{4. Três vinhetas --- os signos na mesa}

\subsection{Vinheta A --- O árbitro é puro índice (o corpo conforme)}

Toda afirmação, todo lance, entra na mesa como um par de mostradores \textbf{b = (testa, refuta)}.

O corpo conforme mergulha o lance no cone: \textbf{X(x)$^2$ = 0} --- \textbf{norma nula}. Ou seja: \textbf{o lance, sozinho, não vale nada}. Não há primeiridade absoluta ali; a obra não carrega carga própria. A medida passa toda para o \textbf{par}:

\begin{tcolorbox}[colback=ouroclaro!35,colframe=ouro,boxrule=0pt,leftrule=2pt,arc=0pt,left=4mm,right=3mm,top=2mm,bottom=2mm,breakable]\itshape \textbf{X(x) $\cdot$ X(y) = -$\tfrac12$$\|$\allowbreak{}x-y$\|$$^2$}\end{tcolorbox}

E o árbitro lê o \textbf{ponteiro}:

{\fontsize{10.00pt}{12.00pt}\selectfont\begin{longtable}{>{\fontsize{10.00pt}{12.00pt}\selectfont\RaggedRight\hspace{0pt}\arraybackslash}p{4402086sp}>{\fontsize{10.00pt}{12.00pt}\selectfont\RaggedRight\hspace{0pt}\arraybackslash}p{6688643sp}>{\fontsize{10.00pt}{12.00pt}\selectfont\RaggedRight\hspace{0pt}\arraybackslash}p{16798327sp}}\toprule
\textbf{Ponteiro} & \textbf{Mostradores} & \textbf{Veredito} \\ \midrule\endhead
\textbf{-$\tfrac12$} & (1,0) & \textbf{tem dentes} --- jogada real, distante da trivialidade \\
\textbf{0} & (1,1) & \textbf{queda de bandeira} --- o lance é a trivialidade; o mate que não mata \\
\textbf{-1} & (0,0) & falha dupla \\
\bottomrule\end{longtable}}

Isto é Secundidade em estado puro: \textbf{a medida não está na qualidade, está no embate.} É por isso que o certificador (as Princesas, o corpo técnico) mora aqui: quem atesta a lei não pode ter carga própria --- atesta pela distância ao outro.

\subsection{Vinheta B --- Ícone contra símbolo: o cristal repele o ouro}

O \textbf{cristalino} é ícone: a rede que se repete idêntica a si. Mas o ícone \textbf{paga um preço exato}: a restrição cristalográfica só deixa passar rotações de ordem \textbf{n $\in$ \{1,2,3,4,6\}}. E o \textbf{primeiro traço proibido é justamente o áureo} (2$\cdot$\allowbreak{}cos(2$\pi$/5) = $\varphi$-1).

Ou seja: \textbf{o ataque áureo da Rainha (a dobra-5) é REPELIDO pelo cristal.} Nenhum cristal admite simetria pentagonal. A tese do reino, dita em uma linha:

\begin{tcolorbox}[colback=ouroclaro!35,colframe=ouro,boxrule=0pt,leftrule=2pt,arc=0pt,left=4mm,right=3mm,top=2mm,bottom=2mm,breakable]\itshape \textbf{Fechar exclui preencher.}\end{tcolorbox}

O ícone que fechou perfeitamente não pode receber o símbolo que preenche perfeitamente. Para \textbf{quebrar} o cristal, o adversário não bate de frente: mira um feixe e lê os \textbf{picos do reticulado recíproco} --- o feixe de ordem zero atravessa sem dano (refração, $\Gamma$=0), os desvios são o dano (difração, o resíduo). E o mate suave é o \textbf{quasicristal}: ali o gap fecha, os picos acumulam, refração e difração colapsam --- e o áureo é quem colapsa \textbf{mais devagar}.

\subsection{Vinheta C --- Onde a Terceiridade parece falhar (o corpo entrópico) --- e por que \textbf{não falha}}

O corpo entrópico (o semianel tropical: \textbf{$\oplus$ = max}, \textbf{$\otimes$ = +}) é a zona da seta do tempo. \textbf{Visto sozinho:}

\begin{itemize}
\item \textbf{$\oplus$ = max} é \textbf{idempotente} (a$\oplus$\allowbreak{}a = a) $\to$ \textbf{não há inverso aditivo} \emph{nele mesmo};
\item sem inverso aditivo não há fase, e sem fase não há interferência;
\item logo, \textbf{de dentro}, nada se desfaz: o dano só acumula.
\end{itemize}

Na linguagem , \textbf{parece} sobrar o índice sem o símbolo --- a reação sem a mediação. E foi aqui que o reino quase se enganou.

\textbf{Mas a mediação não desapareceu: ela mudou de corpo.} Ela está no \textbf{dual --- o corpo cósmico}, cujo operador é a \textbf{expansão} \code{a(t) = e\textasciicircum{}\{Ht\}}. Sob o \code{exp}, o \code{max} volta a ser \textbf{soma}, a soma volta a ter \textbf{inverso}, e a Terceiridade \textbf{fecha}. A frase certa não é \emph{"a soma que esquece"}; é:

\begin{tcolorbox}[colback=ouroclaro!35,colframe=ouro,boxrule=0pt,leftrule=2pt,arc=0pt,left=4mm,right=3mm,top=2mm,bottom=2mm,breakable]\itshape \textbf{A soma que esquece de um lado é a soma que se lembra do outro.}\end{tcolorbox}

E o que parecia a fraqueza do reino é, na verdade, \textbf{a demonstração mais forte da Lei}: a conservação não vive dentro de um corpo --- \textbf{vive sobre o par}. A entropia que cresce aqui é o volume que se abre ali (\code{dS = d log V}, resíduo 0). O corpo isolado é um \textbf{semicorpo}; \textbf{o par é corpo} --- cada um doa ao outro o inverso que lhe falta, e juntos dão \code{(R, +, .)}.

\textbf{A Lei é geral. Nenhum corpo é exceção. Nenhum jogador é condenado} --- e a vitória, aqui, é o caminho de melhor \textbf{peso médio de ciclo} --- a repetição que estabiliza ---, sabendo que a saída existe: \textbf{está do outro lado.}

\vspace{2mm}\hrule\vspace{3mm}

\section{5. Os signos na ISA --- o mais concreto que existe}

Não acredite em nada disto sem ver no metal. A ISA ERG-64 é inteira, exata, resíduo 0 real --- sem float. Os três signos estão \textbf{nos opcodes}:

\begin{Verbatim}[fontsize=\small,frame=single,framesep=3mm]
LOAD 1 ; a i64 nua no slot -- o ÍCONE (a qualidade, o dado sem referência)
LOAD 2
ADD 1 2 -> 3 ; (+) CLIFFORD -- o ÍNDICE (o par no nó: a corrente que se junta)
 ; no booleano: XOR (a diferença simétrica, LdeltaL=0)
AND 2 4 -> 5 ; (x) LA HIRE -- a MEDIAÇÃO (o ganho; a MULT é double-and-add: ADD deslocado)
CMP 9 0
JZ soma ; (op) PONTRYAGIN -- o SÍMBOLO (o dispatch É a lei que costura (+) e (x))
NOT x ; e o NOT é a involução, o complemento, o dual
HALT ; o mate: resíduo 0
\end{Verbatim}

O \textbf{slot 9} guarda o opcode (\code{0=(+)}, \code{1=(x)}, \code{2=(op)}). Aquele \code{CMP}/\code{JZ} --- o dispatch --- não é encanamento: \textbf{é a Terceiridade encarnada}, a convenção que faz as oito linguagens (C, PTX, WASM, LLVM, GLSL, Haskell, Dafny, Isabelle) dizerem \textbf{o mesmo estado}. É por isso que o corpo criativo, o corpo das linguagens, é o corpo do \textbf{símbolo}: linguagem é convenção pura.

E o par \textbf{$\oplus$ / $\otimes$} revela seu segredo no corpo espaço-temporal ($\mathbb{Q}$$_2$), na \textbf{identidade-mestra do somador}:

\begin{tcolorbox}[colback=ouroclaro!35,colframe=ouro,boxrule=0pt,leftrule=2pt,arc=0pt,left=4mm,right=3mm,top=2mm,bottom=2mm,breakable]\itshape \textbf{x + y = (x $\oplus$ y) + 2$\cdot$(x $\wedge$ y)}\end{tcolorbox}

O \textbf{XOR} é o dígito que fica (o \textbf{ícone}: a qualidade pura, o que sobra quando se esquece o transporte). O \textbf{AND} é o \textbf{carry} (o \textbf{índice}: a reação factual entre dois dígitos, o embate que força o tique seguinte). E o \textbf{2$\cdot$} é um tique do relógio (o \textbf{símbolo}: a valuação, o tempo, a lei que conta os transportes).

\begin{tcolorbox}[colback=ouroclaro!35,colframe=ouro,boxrule=0pt,leftrule=2pt,arc=0pt,left=4mm,right=3mm,top=2mm,bottom=2mm,breakable]\itshape \textbf{$\oplus$ e $\wedge$ não são duas operações booleanas de conveniência. São as duas metades de uma única soma:} o dígito e o transporte. O Duque ($\oplus$, Joaquim) e a Duquesa ($\wedge$/$\otimes$, Yasmin) --- inseparáveis no somador.\end{tcolorbox}

\vspace{2mm}\hrule\vspace{3mm}

\section{6. Nota do reino: a casa o signo}

O reino invoca o signo \textbf{duas vezes}, e não é coincidência de nome.

\begin{itemize}
\item \textbf{A decomposição } (a álgebra): no corpo \textbf{telescópico}, os eixos ortogonalizados são \textbf{idempotentes primitivos ortogonais que somam 1} --- e é exatamente aí que o corpo \textbf{deixa de ser corpo}: $e_i \cdot e_j$ = 0, os divisores de zero. Dois eixos ortogonais \textbf{se aniquilam}.
\item \textbf{As três categorias } (o signo): e é \emph{literalmente a mesma coisa lida como semiótica} --- os idempotentes que se aniquilam \textbf{são} a Secundidade em estado bruto: o par que se cancela, a reação factual. A norma que telescopa (N(h$\odot$\allowbreak{}g) = N(h)$\cdot$\allowbreak{}N(g), decompor sem perder nem acrescentar) é a Terceiridade: a lei.
\end{itemize}

Assim o telescópico é o corpo onde as duas heranças se encontram no mesmo teorema --- \textbf{o símbolo por fora (a norma que telescopa), o índice no coração (os eixos que se aniquilam)}. O golpe-fantasma sai daí: \textbf{bater num eixo ortogonal ao do inimigo dá dano ZERO} --- a jogada some, não troca. É o preço da cisão.

\vspace{2mm}\hrule\vspace{3mm}

\section{7. O fecho: os três signos e o resíduo 0}

Por que três, e não dois, e não quatro?

Porque o \textbf{Contrato} exige os três, e um universo a que falte um \textbf{vaza}:

\begin{itemize}
\item \textbf{sem ícone} (sem qualidade/régua) --- não há o que medir: você não sabe o que é o corpo;
\item \textbf{sem índice} (sem par/flip dual) --- não há contra quem medir: o conforme prova isso ao extremo (norma nula: \textbf{a medida está no par}, ou não está em lugar nenhum);
\item \textbf{sem símbolo} (sem lei/operador) --- não há como \textbf{desfazer}: é o entrópico, o semicorpo, a zona onde o reino deixa de fechar.
\end{itemize}

E os três, juntos, dão a única frase que o reino repete em todos os idiomas --- no cone nulo $\sigma$=1 do eletromagnético, no traço zero T(h)=0 do telescópico, na ausência de arbitragem do econômico, na fórmula do produto $\prod_v$ |x|\_v = 1 do espaço-temporal:

\begin{tcolorbox}[colback=ouroclaro!35,colframe=ouro,boxrule=0pt,leftrule=2pt,arc=0pt,left=4mm,right=3mm,top=2mm,bottom=2mm,breakable]\itshape \#\#\# \textbf{$\Gamma$ = 0. FP = 1. Nenhum universo pode vazar.}\end{tcolorbox}

O ícone diz \textbf{o que é}. O índice diz \textbf{contra quem}. O símbolo diz \textbf{que pode}.
Leia os três --- e o lance é legal, e a partida fecha em \textbf{resíduo 0}.

\vspace{2mm}\hrule\vspace{3mm}

\chapter{As jogadas e o combate --- a partida na linguagem do universo}

\begin{tcolorbox}[colback=ouroclaro!35,colframe=ouro,boxrule=0pt,leftrule=2pt,arc=0pt,left=4mm,right=3mm,top=2mm,bottom=2mm,breakable]\itshape \emph{"Não se move uma peça. Executa-se uma operação."}\end{tcolorbox}

\section{1. A regra zero}

No Reino Dourado não existe "movimento" no sentido comum. \textbf{Um LANCE É uma OPERAÇÃO num CORPO.} Ponto.

Quando você desliza a Torre por uma coluna, você não está empurrando madeira: você está executando um \code{XOR} sobre um conjunto de casas. Quando a Rainha ganha força, você não está "aumentando um número": você está rodando 64 micro-lances de \emph{double-and-add} na ISA. Quando um peão promove, você não está "trocando de peça": você está pedindo ao operador que lhe dê um \textbf{inverso} --- e um inverso é o degrau que leva um ANEL a virar CORPO.

Toda a mecânica do jogo cabe em três verbos. São os três, e não há um quarto.

{\fontsize{10.00pt}{12.00pt}\selectfont\begin{longtable}{>{\fontsize{10.00pt}{12.00pt}\selectfont\RaggedRight\hspace{0pt}\arraybackslash}p{2937855sp}>{\fontsize{10.00pt}{12.00pt}\selectfont\RaggedRight\hspace{0pt}\arraybackslash}p{5813990sp}>{\fontsize{10.00pt}{12.00pt}\selectfont\RaggedRight\hspace{0pt}\arraybackslash}p{6031038sp}>{\fontsize{10.00pt}{12.00pt}\selectfont\RaggedRight\hspace{0pt}\arraybackslash}p{6814017sp}>{\fontsize{10.00pt}{12.00pt}\selectfont\RaggedRight\hspace{0pt}\arraybackslash}p{5243580sp}}\toprule
\textbf{Símbolo} & \textbf{Nome} & \textbf{O que é} & \textbf{Na ISA ERG-64} & \textbf{o signo} \\ \midrule\endhead
\textbf{$\oplus$} & \textbf{CLIFFORD} & a SOMA, o PAR, a corrente que se junta no nó & \code{ADD} (ou \code{XOR} no corpo booleano) & Secundidade --- o \textbf{índice} (o embate factual) \\
\textbf{$\otimes$} & \textbf{LA HIRE} & o PRODUTO, a Rainha, o ganho / a tensão & a MULT por \emph{double-and-add} (\code{ADD} deslocado + \code{AND}) & a mediação \\
\textbf{$\prod$} & \textbf{PONTRYAGIN} & o OPERADOR que costura: $\exp\circ\Sigma\circ\log$ & \code{NOT} / o dispatch (\code{CMP},\code{JZ},\code{JMP}) / a composição & Terceiridade --- o \textbf{símbolo} (a lei) \\
\bottomrule\end{longtable}}

E a qualidade própria de cada corpo --- o seu "sabor", aquilo que ele \emph{é} antes de reagir com o que quer que seja --- é a \textbf{Primeiridade}, o \textbf{ícone}.

\vspace{2mm}\hrule\vspace{3mm}

\section{2. O lance compila}

Esta é a diferença entre o Reino Dourado e qualquer outro jogo: \textbf{todo lance é um programa}, e o programa é exato.

A ISA do reino chama-se \textbf{ERG-64}. Os seus opcodes são todos os que existem:

\begin{Verbatim}[fontsize=\small,frame=single,framesep=3mm]
LOAD STORE ADD SUB AND OR XOR CMP JMP JZ JNZ HALT
\end{Verbatim}

Ela é \textbf{inteira} (i64) e \textbf{exata} --- resíduo 0 \emph{real}, sem float. Da mesma IR o compilador do reino (o \code{chessc}) emite \textbf{WASM} (o metal, o que corre) e \textbf{Dafny} (as leis, o que se prova): a costura \code{interp == WASM == Dafny}. Isto tem uma consequência dura de regra de mesa:

\begin{tcolorbox}[colback=ouroclaro!35,colframe=ouro,boxrule=0pt,leftrule=2pt,arc=0pt,left=4mm,right=3mm,top=2mm,bottom=2mm,breakable]\itshape \#\#\#  LEI DO ARREDONDAMENTO \textbf{É lance ILEGAL arredondar no meio da partida.} O estado corre em inteiros --- pares (num, den) no corpo racional $\mathbb{Q}$, pares (a, b) no corpo áureo $\mathbb{Z}$[$\varphi$]. O \emph{float} só aparece no FIM, ao pintar o pixel, onde 1 ulp é invisível. Arredondar no meio é \textbf{vazar} --- e vazar é o único pecado do reino (\S6).\end{tcolorbox}

E daí sai a promessa da mesa: \textbf{toda partida é replayável bit a bit, e todo mate é auditável.}

\vspace{2mm}\hrule\vspace{3mm}

\section{3. Os três lances, em detalhe}

\subsection{$\oplus$ Clifford --- o lance de SOMA (o par, o avanço)}

Somar é \textbf{juntar sem misturar}: componente a componente, coordenada a coordenada. É a corrente a entrar no nó.

\begin{itemize}
\item No \textbf{corpo fractal} (a raiz): avançar é \code{(a,b,c,d) (+) (a',b',c',d') = (a+a', b+b', c+c', d+d')} --- quatro \code{ADD} independentes.
\item No \textbf{corpo criativo} (as linguagens): $\oplus$ é a \textbf{diferença simétrica} \code{L delta M} = \code{XOR}. E é aqui que se vê o dente: \textbf{\code{L delta L = 0}} --- cada coisa é o seu próprio inverso. (Troque $\triangle$ por $\cup$ e falha: \code{L cup L = L != 0}. A união não desfaz; a diferença simétrica desfaz.)
\item No \textbf{corpo áureo}: o infinito soma-se com \textbf{dois \code{ADD}}. Só isso.
\end{itemize}

$\oplus$ é o lance mais barato do tabuleiro. É também o que dá \textbf{sinal} ao corpo --- é o inverso aditivo que permite que duas amplitudes se \textbf{cancelem}. Guarde esta frase: onde não há inverso aditivo, não há cura (\S7, o corpo entrópico).

\subsection{$\otimes$ La Hire --- o lance de PRODUTO (a Rainha, o ganho, o giro)}

Multiplicar é \textbf{compor}: aplicar um ganho, girar um rumo, empilhar um meio.

Na ISA não existe opcode \code{MUL}. \textbf{A Rainha executa a multiplicação como 64 micro-lances} (\emph{double-and-add}):

\begin{Verbatim}[fontsize=\small,frame=single,framesep=3mm]
; a CARGA DA RAINHA -- (x) = Sigma o deslocamento
LACO: AND mask, b ; o bit de b está ligado?
 CMP t, 0
 JZ PULA
 ADD acc, aa ; sim: acumula o parcial deslocado
PULA: ADD aa, aa ; dobra o parcial (o exp)
 ADD mask, mask ; dobra a máscara (o log: a escala 2^k)
 SUB cnt, 1 ; 64 ticks
 JNZ LACO
\end{Verbatim}

Repare no que isto significa em jogo: \textbf{um golpe amplificado não é "um número maior", é uma sequência de 64 lances bit a bit} --- e nenhum deles perde um ulp. O ganho é exato.

$\otimes$ muda de rosto conforme o corpo:
\begin{itemize}
\item \textbf{fractal}: a convolução ciclotômica com $\zeta$$^{4}$ = -1 --- girar o rumo 45$^\circ$.
\item \textbf{criativo}: \code{AND} (a interseção $\cap$). Aqui a MULT \textbf{colapsa} --- no booleano \code{x (x) x = x}, a idempotência engole o \emph{double-and-add}.
\item \textbf{eletromagnético}: as \textbf{impedâncias multiplicam} ao empilhar meios: $\sigma$($\Delta$\allowbreak{}h$_1$+$\Delta$\allowbreak{}h$_2$) = $\sigma$($\Delta$\allowbreak{}h$_1$)$\cdot$$\sigma$($\Delta$\allowbreak{}h$_2$).
\item \textbf{evolutivo}: a ponderação pela aptidão, \code{p\_i . w\_i}.
\item \textbf{racional $\mathbb{Q}$}: a Rainha \textbf{composta duas vezes} --- uma no numerador, outra no denominador.
\end{itemize}

\subsection{$\prod$ Pontryagin --- o lance de CONVERSÃO (o inversor, a poinsétia, a promoção)}

Este é o lance que faz o Reino Dourado ser o Reino Dourado.

\textbf{$\prod$ = exp $\circ$ $\Sigma$ $\circ$ log.} Ele \textbf{costura o aditivo ao multiplicativo}. Desce pelo log (onde produtos viram somas), soma na reta, e sobe pelo exp de volta.

O que $\prod$ \textbf{faz} na mesa, concretamente:

\begin{enumerate}
\item \textbf{DÁ O INVERSO.} Sem $\prod$, $\mathbb{Z}$[$\zeta$$_8$] é só um \textbf{anel}: o 2 não tem inverso, não há como desfazer. Com $\prod$, 2 $\mapsto$ $\tfrac12$, e $\mathbb{Q}$($\zeta$$_8$) é um \textbf{CORPO}. Em jogo: \textbf{a PROMOÇÃO é $\prod$} --- transforma um peão do anel (unidirecional, sem \emph{undo}) numa peça do corpo (bidirecional, reversível via -log).
\item \textbf{CONVERTE A ESCALA.} No \textbf{corpo econômico}: baixar o índice (log) leva o juro \textbf{composto} (PG, $\otimes$) ao juro \textbf{simples} (PA, $\oplus$); subir (exp) devolve. É a mesma lente que no \textbf{eletromagnético} lineariza a impedância $\sigma$ = e\textasciicircum{}(-$\Delta$\allowbreak{}h/2) no perfil aditivo.
\item \textbf{CASA A IMPEDÂNCIA.} $\prod$ é o \textbf{inversor multinível} --- o que recoloca no cone toda energia que estaria a vazar. Ele é o guardião da lei (\S6).
\item \textbf{DECIDE.} Na ISA, o \emph{dispatch} por opcode É Pontryagin: \code{CMP} + \code{JZ} a escolher qual operação executar. O colapso $\Psi$ (o minimax/argmax que escolhe o lance) é a \textbf{seleção de Pontryagin}.
\end{enumerate}

\vspace{2mm}\hrule\vspace{3mm}

\section{4. Ataque: a régua diz quanto}

\textbf{Não há "pontos de dano" no Reino Dourado. Há RÉGUAS.} Cada corpo mede com a sua, e a régua é parte do contrato de entrada do corpo no reino.

{\fontsize{10.00pt}{12.00pt}\selectfont\begin{longtable}{>{\fontsize{10.00pt}{12.00pt}\selectfont\RaggedRight\hspace{0pt}\arraybackslash}p{7573315sp}>{\fontsize{10.00pt}{12.00pt}\selectfont\RaggedRight\hspace{0pt}\arraybackslash}p{8463969sp}>{\fontsize{10.00pt}{12.00pt}\selectfont\RaggedRight\hspace{0pt}\arraybackslash}p{11851772sp}}\toprule
\textbf{Corpo} & \textbf{A régua $\mu$ ($T_\omega$)} & \textbf{O que o golpe mede} \\ \midrule\endhead
\textbf{fractal} & z$\cdot$\allowbreak{}z (a norma pela conjugação) & alcance/dano do rumo \\
\textbf{eletromagnético} & $\Sigma$\allowbreak{}h = log $|E\times B|^2$ (o telescópio) & a energia que \textbf{propaga} \\
\textbf{telescópico} & (N, T) --- a norma-covolume e o traço & o perfil de um alinhamento \\
\textbf{relógio} & o lapso N($\psi$) = cos $\psi$ & quão depressa corres \\
\textbf{celeste} & r$^2$ + C$^2$ = 1 & ativo (dano solo) vs reativo (emaranhamento) \\
\textbf{óptico} & n$\cdot$\allowbreak{}sin$\theta$ & a componente tangencial --- \textbf{o invariante da fronteira} \\
\textbf{econômico} & $\delta$ = log(1+r), a força de juros & o crescimento do valor \\
\textbf{evolutivo} & 1 - s$^2$, s = 2p-1 & a frequência gênica entre extinção e fixação \\
\textbf{espaço-temporal} & $|x|_2 = 2^{-v}$ (ultramétrica) & tempo = valuação; espaço = o dígito \\
\textbf{conforme} & $\langle X(x), X(y)\rangle = -\tfrac12\,|x-y|^2$ & a \textbf{distância entre duas afirmações} \\
\bottomrule\end{longtable}}

O golpe \textbf{canônico} do reino é o \textbf{produto pelo flip dual}: \code{z (x) z = |z|\textasciicircum{}2}. Multiplica-se a peça pelo seu próprio espelho, e o que sai é a norma --- quanto o golpe pontua. É La Hire aplicado ao dual: o ataque é sempre um par.

\vspace{2mm}\hrule\vspace{3mm}

\section{5. Defesa: o flip dual $\nu$}

Toda defesa do reino é uma involução --- algo que, aplicado duas vezes, devolve o mundo ($\nu$$^2$ = id). Escolha a sua:

\begin{itemize}
\item \textbf{O ESPELHO} (fractal, cristalino) --- a conjugação \code{z |-> z}: reflete o rumo do inimigo de volta.
\item \textbf{A RECÍPROCA} (racional $\mathbb{Q}$) --- \code{[(a,b)]\textasciicircum{}-\allowbreak{}\textasciicircum{}1 = [(b,a)]}: um ataque de razão \textbf{3/4} é anulado \emph{exatamente} pela sua recíproca \textbf{4/3}. Primal $\otimes$ dual = 1, o neutro. \textbf{Em $\mathbb{Z}$ (o anel) esta defesa não existe. Em $\mathbb{Q}$ (o corpo) todo ataque não-nulo TEM defesa exata.} É por isso que se paga o preço de subir ao corpo.
\item \textbf{A REFLEXÃO TOTAL INTERNA} (óptico) --- o \textbf{escudo perfeito}. Se o ataque incide na tua fronteira vindo do meio mais denso (n$_1$ > n$\allowbreak{}_2$) e passa do \textbf{ângulo crítico} $\theta$\allowbreak{}c = arcsin(n$_2$/n$\allowbreak{}_1$) --- para 1,5$\to$\allowbreak{}1, \textbf{41,81$^\circ$} --- então \textbf{100\% do golpe é devolvido}. Nada penetra: a onda vira evanescente e decai como e\textasciicircum{}(-$\kappa$\allowbreak{}z). É o flip lido no horizonte.
\item \textbf{A ARMADURA CRISTALINA} (cristalino) --- congela uma zona num reticulado de \textbf{regulador nulo}. Dentro dela o relógio PARA e só cabem giros de ordem \{1, 2, 3, 4, 6\}. \textbf{Repele o ataque áureo:} a dobra-5 e a dobra-10 (o traço $\varphi$) são \emph{cristalograficamente proibidas}. A Rainha não entra num cristal. (Preço: a defesa que fecha \textbf{não preenche} --- o cristal proíbe exatamente o preenchedor ótimo.)
\item \textbf{A RÉGUA BLINDADA} (telescópico) --- o perfil autodual, \textbf{T(h) = $\Sigma$\allowbreak{}h$_i$ = 0}, covolume 1: o que um eixo estica (+h$_i$, o indutor) outro contrai (-h$\allowbreak{}_i$, o capacitor). Sem reatância líquida. Invulnerável porque \textbf{não vaza}.
\item \textbf{O FLIP DE HODGE} (eletromagnético) --- (E, B) $\mapsto$ (B, -E): gira o par \textbf{preservando o telescópio}. Troca o eixo do ataque sem perder um joule.
\end{itemize}

\vspace{2mm}\hrule\vspace{3mm}

\section{6. A LEI DA MESA: nenhum universo pode vazar}

Esta é a única regra que não se negocia, e é a mesma em todos os corpos:

\begin{tcolorbox}[colback=ouroclaro!35,colframe=ouro,boxrule=0pt,leftrule=2pt,arc=0pt,left=4mm,right=3mm,top=2mm,bottom=2mm,breakable]\itshape \#\# $\Gamma$ = (Z - Z$\allowbreak{}_0$)/(Z + Z$\allowbreak{}_0$) = 0 $\Longleftrightarrow$ \textbf{FP = 1} \textbf{Nenhum lance pode ter fator de potência diferente de 1.}\end{tcolorbox}

No \textbf{corpo eletromagnético} ela é literal: o \textbf{CONE NULO $\sigma$ = |E|/|B| = 1}. Um lance \textbf{casado} ($\sigma$ = 1) entrega \textbf{toda} a potência do telescópio ao alvo --- transferência 1 - |$\Gamma$|$^2$ = 1, o golpe que não vaza. Um lance \textbf{descasado} ($\sigma$ $\neq$ 1) deixa \textbf{potência reativa presa} (E$^2$ - B$^2$ $\neq$ 0): energia que volta ao atacante, dano perdido.

E a mesma lei, traduzida no idioma de cada corpo:

\begin{itemize}
\item \textbf{telescópico} --- traço zero, T(h) = 0 (a régua reta).
\item \textbf{econômico} --- \textbf{ausência de arbitragem}: o desconto desfaz \emph{exatamente} a capitalização, e\textasciicircum{}(-$\delta$\allowbreak{}t)$\cdot$\allowbreak{}e\textasciicircum{}($\delta$\allowbreak{}t) = 1. FP < 1 é o \textbf{spread}: o atrito que sangra valor a cada turno, e que o adversário vai explorar.
\item \textbf{celeste} --- \textbf{C = 0} (estado produto): a peça isolada em FP = r = 1, potência máxima.
\item \textbf{somático} --- \textbf{tr G = 0}: a \textbf{homeostase}, a população que se conserva.
\item \textbf{evolutivo} --- a normalização $\Sigma$\allowbreak{}p$_i$ = 1: a massa do exército não se cria.
\item \textbf{entrópico} --- a \textbf{convexidade}: repousar na envoltória $f^{**}$, sem custo acima do piso.
\item \textbf{geométrico} --- o \textbf{cone nulo ds$^2$ = 0}: o disparo de luz que chega limpo. (E o poço mal-cavado --- trocar 1-2M/r por 1-2M/r$^2$ --- tem Ricci $\neq$ 0: fonte espúria, \textbf{vaza}.)
\item \textbf{espaço-temporal} --- a forma nua do Contrato: \textbf{$\prod$$_v$ |x|$\allowbreak{}_v$ = 1}. \emph{A bandeira que cai num mostrador é exatamente o tempo ganho no outro.}
\item \textbf{fractal} --- só rumos de \textbf{norma unitária |z| = 1} são lances legais. Mover com |z| $\neq$ 1 vaza.
\end{itemize}

\textbf{Quando um lance vaza, o inversor multinível ($\prod$ Pontryagin) o recoloca no cone --- degrau a degrau, pelos harmónicos $\varphi$\textasciicircum{}(-j) da garrafa áurea de Koch.} O reino não te pune: ele te \textbf{corrige}. Mas o vazamento cobra o seu preço em turnos.

\vspace{2mm}\hrule\vspace{3mm}

\section{7. A partida DUAL --- a corte}

O tabuleiro do Reino Dourado é \textbf{duplo}. Toda partida corre em dois lados ao mesmo tempo, e o par é o motor de tudo:

{\fontsize{10.00pt}{12.00pt}\selectfont\begin{longtable}{>{\fontsize{10.00pt}{12.00pt}\selectfont\RaggedRight\hspace{0pt}\arraybackslash}p{10512049sp}>{\fontsize{10.00pt}{12.00pt}\selectfont\RaggedRight\hspace{0pt}\arraybackslash}p{3603025sp}>{\fontsize{10.00pt}{12.00pt}\selectfont\RaggedRight\hspace{0pt}\arraybackslash}p{13773982sp}}\toprule
\textbf{PRIMAL (gerar)} & \textbf{$\bowtie$} & \textbf{DUAL (atestar)} \\ \midrule\endhead
o \textbf{corpo criativo} --- as oito linguagens, quem \textbf{executa} &  & as \textbf{PRINCESAS} --- o corpo técnico, quem \textbf{ATESTA} (Dafny, Isabelle, o kernel LCF) \\
o \textbf{motor} --- dissipa, cai em ciclos, dá torque &  & o \textbf{rotor} --- conserva a fase (o flip, a $\mapsto$ ia) \\
o \textbf{estado} (evolui por G) &  & o \textbf{co-estado} (evolui por G$^T$, para trás --- o princípio do máximo) \\
a \textbf{expansão} $\sigma_m$ (o Duque $\oplus$, Joaquim) &  & a \textbf{contração} $\sigma$ = -1/$\sigma$ (a Duquesa $\otimes$, Yasmin) \\
\textbf{proliferar} (a mitose, o sucessor S) &  & \textbf{diferenciar} ($\pi$, projetar nos tipos) \\
\textbf{capitalizar} &  & \textbf{descontar} \\
\bottomrule\end{longtable}}

\textbf{Regra da corte:} um lance no primal \textbf{obriga} uma resposta no dual. Toda afirmação carrega dois mostradores \textbf{b = (testa, refuta) $\in$ \{0,1\}$^2$}, e o árbitro (o \textbf{corpo conforme}, a casa da medida) lê o \textbf{PONTEIRO}:

\begin{Verbatim}[fontsize=\small,frame=single,framesep=3mm]
pt(b) = <X(b), X(bT)> = -½ . d(b, bT)^2

 -½ -> o lance TEM DENTES (testa , refuta ) -- é uma jogada real
 0 -> QUEDA DE BANDEIRA (testa , refuta ) -- a afirmação é trivial: lance VAZIO
 -1 -> FALHA DUPLA (nada passa)
\end{Verbatim}

O mostrador diz \textbf{qual lado} falhou; o ponteiro diz \textbf{quanto}. Um lance que marca 0 não é um lance: é o relógio a parar. \textbf{Não basta afirmar --- tem de haver um refutador com dentes.}

E o \textbf{Rei} é dual das peças: elas refletem o seu passo ($\nu$: e $\mapsto$ -e), mantendo o gap $\pm$\allowbreak{}1.

\vspace{2mm}\hrule\vspace{3mm}

\section{8. O MATE que não mata}

Aqui está o coração do jogo, e é onde ele se separa de todos os outros.

\textbf{O mate do Reino Dourado não elimina o Rei. O mate é a CRISTALIZAÇÃO: o resíduo vai a zero.}

\begin{itemize}
\item No \textbf{corpo criativo}, o mate é a \textbf{mobilidade que colapsa a 0} --- a expansão entrópica cristaliza. (Certificado, e não é retórica: a partida completa da posição inicial ao mate, \textbf{Evans Gambit} --- 1.e4 e5 2.Cf3 Cc6 3.Bc4 Bc5 4.b4 \ldots \textbf{12.Dxf7\#} --- xeque-mate em 23 lances, a Dama (C $\oplus$ Dafny) cristalizando. A Ópera de Morphy reproduzida em 0,6 s.)
\item No \textbf{relógio universal}, o mate é o único tique \textbf{sem inverso}: o lapso N vai a 0, \textbf{o tempo congela}. Onde a medida fecha, o tempo para.
\item No \textbf{corpo expansivo}, o mate não é matar --- é \textbf{fechar a órbita sobre o círculo}: N(x) = $\pm$\allowbreak{}1, resíduo nulo. \emph{A peça vira unidade.}
\item No \textbf{corpo áureo}, o mate é a \textbf{DOBRA}: toda a potência do adversário reduzida ao par ($F_{n-1}, F_n$). Não há para onde expandir --- a cadeia infinita já está fechada em dois inteiros. \textbf{O infinito não é derrotado; é dobrado.}
\item No \textbf{corpo geométrico}, o mate é o \textbf{vácuo}, $R_{\mu\nu}=0$: o poço perfeito, curvatura conservada --- nada vaza.
\item No \textbf{corpo somático}, o mate é a \textbf{HOMEOSTASE}: o tecido para de proliferar por inibição de contato, exatamente como os peões colidem e travam.
\item No \textbf{corpo racional}, o mate é o \textbf{resto 0 de Euclides} (\S9).
\end{itemize}

O mate é a única \textbf{dissipação global} da partida. Ninguém morre. \textbf{Tudo fecha.}

\vspace{2mm}\hrule\vspace{3mm}

\section{9. O LIVRO DOS LANCES --- a jogada de cada corpo}

Cada corpo do reino traz para a mesa uma mecânica própria. Isto é o teu manual de combate.

\subsection{ CORPO FRACTAL --- o tabuleiro que é a raiz}
Cada peça carrega um \textbf{RUMO}: uma das \textbf{8 unidades $\zeta$$_8$$^k$} (o relógio $\mathbb{Z}$/8, oito direções de 45$^\circ$).
\begin{itemize}
\item \textbf{GIRO} ($\otimes$): multiplicar o rumo por $\zeta$$_8$ gira a peça 45$^\circ$ --- um tique do relógio. \code{zeta (x) zeta\textasciicircum{}3 = zeta\textasciicircum{}4 = -\allowbreak{}1} é a meia-volta completa.
\item \textbf{AVANÇO} ($\oplus$): somar deslocamento, coordenada a coordenada.
\item \textbf{PROMOÇÃO} ($\prod$): concede o inverso --- o peão sai do anel e entra no corpo. \textbf{Passa a poder desfazer.}
\item \textbf{ATRATOR DE NEWTON}: uma peça lançada ao tabuleiro \textbf{converge} a uma das 8 raízes (lance de \emph{homing} à âncora mais próxima). Se cai na \textbf{fronteira fractal} (irracional), \textbf{nunca assenta} --- casa instável, eternamente contestada.
\item O zoom não muda o tabuleiro: a tesselação é \textbf{auto-similar}, nativa. (Primeiridade / \textbf{ícone}: o tabuleiro que se parece consigo em toda escala.)
\end{itemize}

\subsection{ CORPO CRIATIVO --- as peças SÃO os operadores}
Aqui não há metáfora nenhuma. \textbf{A peça é a operação:}

{\fontsize{10.00pt}{12.00pt}\selectfont\begin{longtable}{>{\fontsize{10.00pt}{12.00pt}\selectfont\RaggedRight\hspace{0pt}\arraybackslash}p{3865652sp}>{\fontsize{10.00pt}{12.00pt}\selectfont\RaggedRight\hspace{0pt}\arraybackslash}p{9883642sp}>{\fontsize{10.00pt}{12.00pt}\selectfont\RaggedRight\hspace{0pt}\arraybackslash}p{5386319sp}>{\fontsize{10.00pt}{12.00pt}\selectfont\RaggedRight\hspace{0pt}\arraybackslash}p{8229155sp}}\toprule
\textbf{Peça} & \textbf{Operação} & \textbf{Linguagem} & \textbf{Geometria} \\ \midrule\endhead
\textbf{Torre} & $\oplus$ Clifford (\code{XOR}) & C & as retas (PA) \\
\textbf{Bispo} & $\otimes$ La Hire (\code{AND}) & Dafny & as diagonais (PG) \\
\textbf{Dama} & C $\oplus$ Dafny, grau (1,1) & ambas & Q = R + B \\
\textbf{Cavalo} & 2$\cdot$\allowbreak{}C $\oplus$ 1$\cdot$\allowbreak{}Dafny, grau (2,1) & ambas & Pontryagin \textbf{graduado} \\
\textbf{Rei} & C $\oplus$ Dafny de um passo --- mas \textbf{DUAL} das peças & ambas & o gap $\pm$\allowbreak{}1 \\
\textbf{Peão} & o bit $\oplus$ com o carry $\wedge$ & --- & x+y = (x$\oplus$\allowbreak{}y) + 2(x$\wedge$\allowbreak{}y) \\
\bottomrule\end{longtable}}

\textbf{Mover a Torre = executar um XOR. Mover o Bispo = executar um AND.} Não há tradução: é a mesma coisa.

\subsection{ CORPO ELETROMAGNÉTICO --- o golpe casado}
Cada peça/onda carrega um par (E, B) e uma \textbf{impedância $\sigma$ = |E|/|B|} --- o seu \textbf{metal}.
\begin{itemize}
\item \textbf{PROPAGAR} ($\oplus$): superpor perfis, \code{ADD} --- o telescópio soma.
\item \textbf{ATRAVESSAR MEIOS} ($\otimes$): as impedâncias \textbf{multiplicam} --- o casamento de terreno.
\item \textbf{O ATAQUE MÁXIMO} é o lance casado: levar o alvo ao \textbf{cone nulo $\sigma$ = 1}. Toda a potência propaga. O golpe que não vaza.
\item Um lance descasado ($\sigma$ = $\varphi$ $\neq$ 1) deixa energia reativa presa. \textbf{Ela volta para ti.}
\end{itemize}

\subsection{ CORPO MOTOR --- o golpe de flanco}
O torque é \textbf{$\tau$ = $\psi$ $\times$ i = |$\psi$||i|$\cdot$\allowbreak{}sin} --- \textbf{máximo quando fluxo e corrente estão ORTOGONAIS}.
\begin{tcolorbox}[colback=ouroclaro!35,colframe=ouro,boxrule=0pt,leftrule=2pt,arc=0pt,left=4mm,right=3mm,top=2mm,bottom=2mm,breakable]\itshape \textbf{Regra de combate: o golpe de flanco a 90$^\circ$ dá o dobro do frontal. Ataque de lado, não de frente.} É a geometria, não o estilo.\end{tcolorbox}

O \textbf{DTC} (Direct Torque Control) \emph{é} Pontryagin: o \textbf{co-estado} (o valor futuro, G$^T$) diz quanto empurrar \textbf{agora}. E cada peça tem o seu par branco/negro --- os dois bispos em casas de \textbf{cores opostas}: o dipolo dissipativo.

\subsection{ RELÓGIO UNIVERSAL --- o tique e o horizonte}
O lance é um \textbf{TICK}: multiplicar o estado por $\omega$. A órbita é \{$\omega$$^k$$\cdot$\allowbreak{}s$_0$\}, e é \textbf{reversível} ($\omega$\textasciicircum{}(-k) destica). No $\mathbb{Z}$/8, o tique de 45$^\circ$ leva a \textbf{reta da Torre} na \textbf{diagonal do Bispo}.
\begin{itemize}
\item \textbf{No combate, as velocidades combinam pela $\oplus$ relativística:} \code{s\_\allowbreak{}1 (+) s\_\allowbreak{}2 = (s\_\allowbreak{}1+s\_\allowbreak{}2)/\allowbreak{}(1 + s\_\allowbreak{}1s\_\allowbreak{}2)}.
\item \textbf{Consequência dura: nenhum combo cruza o horizonte |s| = 1.} Empilhe quantos buffs quiser --- o dano \textbf{satura} e nunca quebra a régua. \textbf{Não há vitória grátis nos extremos.}
\end{itemize}

\subsection{ CORPO TELESCÓPICO --- o golpe-fantasma}
Ataca-se por \textbf{N eixos} ortogonais, e o dano total \textbf{telescopa}: a norma N = covol$^2$ colapsa na \textbf{soma} das entropias $\Sigma$\allowbreak{}h$_i$. Atacar por muitos eixos é \textbf{somar perfis}, não multiplicar custo.
\begin{tcolorbox}[colback=ouroclaro!35,colframe=ouro,boxrule=0pt,leftrule=2pt,arc=0pt,left=4mm,right=3mm,top=2mm,bottom=2mm,breakable]\itshape  \textbf{O GOLPE-FANTASMA:} dois eixos ortogonais têm \textbf{e$_i$ $\cdot$ e$\allowbreak{}_j$ = 0} --- bater num eixo ortogonal ao do inimigo dá \textbf{dano ZERO}. A jogada simplesmente \textbf{some}. É o preço da cisão. \textbf{Descubra o eixo do inimigo antes de bater.}\end{tcolorbox}

\subsection{ CORPO CRISTALINO --- a armadura e como quebrá-la}
\textbf{Cristalizar} uma zona: o relógio para, só giros de ordem \{1,2,3,4,6\}. Defesa que não vaza e não expande.
\begin{itemize}
\item \textbf{Como QUEBRAR um cristal:} mire um feixe (a soma sobre o reticulado) e leia os \textbf{picos} no reticulado recíproco $L^*$. O feixe de \textbf{ordem zero} (y = 0) atravessa sem dano --- é refração pura. Os \textbf{desvios} (y $\neq$ 0) são o \textbf{dano} --- é a difração, o resíduo.
\item \textbf{O MATE SUAVE:} no \textbf{quasicristal} o gap fecha --- o recíproco vira o módulo denso $\mathbb{Z}$ + $\varphi$$\cdot$$\mathbb{Z}$, os picos acumulam no centro, refração e difração colapsam juntas. A defesa perde o isolamento. \emph{E o áureo é justamente quem colapsa mais devagar} (Hurwitz, liminf = 1/$\sqrt{\ }$\allowbreak{}5).
\end{itemize}

\subsection{ CORPO ÓPTICO --- a fronteira}
Cada território tem um \textbf{índice n}. Ao cruzar (n$_1$ $\to$ n$\allowbreak{}_2$), a trajetória \textbf{dobra} --- e \textbf{tu não escolhes o ângulo de saída}: ele é fixado pela régua, \textbf{n$_1$$\cdot$\allowbreak{}sin$\theta$$_1$ = n$\allowbreak{}_2$$\cdot$\allowbreak{}sin$\theta$$_2$}. A componente tangencial é o invariante. Atravessar N camadas é transitivo (encadeia sem perda). Se o território é um \textbf{cristal}, a fronteira \textbf{difrata}: o ataque parte-se em \textbf{ordens discretas} --- múltiplos golpes em ângulos fixos.

\subsection{ CORPO ENTRÓPICO --- a zona sem volta}
O terreno onde o campo \textbf{deixa de fechar}. Aqui $\otimes$ = \code{ADD} real (o custo soma ao longo do caminho) e $\oplus$ = \code{max} (escolhe a linha dominante).
\begin{tcolorbox}[colback=ouroclaro!35,colframe=ouro,boxrule=0pt,leftrule=2pt,arc=0pt,left=4mm,right=3mm,top=2mm,bottom=2mm,breakable]\itshape  \textbf{NÃO HÁ INVERSO ADITIVO. Não há cura. Não há cancelamento.} O inverso aditivo \emph{é} a interferência/fase, e este é o mundo \textbf{sem fase}: \code{|psi\_\allowbreak{}1+psi\_\allowbreak{}2|\textasciicircum{}2 = 0} nunca acontece --- as probabilidades só somam. A soma que \textbf{esquece} (a$\oplus$\allowbreak{}a = a).\end{tcolorbox}

E o \textbf{SORTEIO obrigatório}: em cada vértice de coordenação z há e\textasciicircum{}(Ns) representações de igual resíduo (a regra do gelo, W = 3/2) --- \textbf{a peça TEM de rolar o dado} a cada nó. A vitória é o caminho de \textbf{melhor peso médio de ciclo}. Não há caos aqui --- o gato tropical é não-expansivo, entropia topológica 0: a rotação morre porque exige o sinal, e não há sinal.

\subsection{ CORPO EVOLUTIVO --- a meta do exército}
Cada tipo de peça tem uma \textbf{frequência p} no exército, entre \textbf{extinção} (p=0) e \textbf{fixação} (p=1). O lance é o \textbf{replicador}: quem capturou e sobreviveu (fitness w alta) sobe --- \code{p\_\allowbreak{}i |-\allowbreak{}> p\_\allowbreak{}i.\allowbreak{}w\_\allowbreak{}i/\allowbreak{}<w>} --- e $\Sigma$\allowbreak{}p$_i$ = 1 conserva (FP = 1, a massa não se cria).
\begin{itemize}
\item \textbf{Fisher:} a aptidão média $\langle$\allowbreak{}w$\rangle$ \textbf{só cresce}. O combate é uma escalada monótona.
\item Empilhar vantagens rende \textbf{cada vez menos} perto do horizonte (a régua hiperbólica de novo).
\item O \textbf{flip p $\mapsto$ 1-p} é a \textbf{deriva}: o reset ao ponto neutro p = $\tfrac12$.
\end{itemize}

\subsection{ CORPO GEOMÉTRICO --- pousar massa}
O lance é \textbf{pousar uma peça pesada}: ela cava um poço \textbf{$\Phi$ = -$\tfrac12$$\cdot$\allowbreak{}log(1 + 2M/|x|)} e as linhas de movimento de todas as outras peças passam a seguir \textbf{geodésicas} --- curvam ao redor dela. O potencial \textbf{soma-se} exatamente onde as normas se \textbf{multiplicam} (é o log). O combate corre no \textbf{cone nulo} --- o disparo de luz, ds$^2$ = 0, é onde FP = 1.

\subsection{ CORPO MÓRFICO --- o território}
Cada casa é um bit (dominada = 1, livre = 0). Quatro lances:
\begin{enumerate}
\item \textbf{EROSÃO} ($\otimes$ / \code{AND}) --- só sobrevivem as casas cuja vizinhança inteira já dominas. Encolhe a fronteira, \textbf{consolida o núcleo}.
\item \textbf{DILATAÇÃO} (o dual) --- espalha a zona a cada casa adjacente. \textbf{O avanço.}
\item \textbf{ANTIROQUE} ($\oplus$ / \code{XOR}) --- a diferença simétrica: as casas contestadas por \textbf{exatamente um} exército. O lance da \textbf{linha de frente} --- ao invés de proteger o rei (o roque), \textbf{cede-e-toma}.
\item \textbf{COMPLEMENTAÇÃO} ($\prod$ / \code{NOT}) --- vira o tabuleiro inteiro. Involução: \textbf{dois lances desfazem}.
\end{enumerate}

\subsection{ CORPO CELESTE --- emaranhar custa dano}
Cada peça reparte uma unidade de energia entre \textbf{r} (potência ATIVA, o dano solo) e \textbf{C} (a REATIVA, o emaranhamento), e \textbf{r$^2$ + C$^2$ = 1} conserva sempre.
\begin{itemize}
\item \textbf{EMARANHAR} ($\otimes$, compor): enlaçar duas peças acopla-as (C > 0) e por isso \textbf{BAIXA o dano individual} de cada uma. \emph{Ganhas correlação, perdes dano solo.}
\item \textbf{FECHAMENTO} (C = 0): a peça isolada, FP = 1, potência máxima.
\item \textbf{BELL} (r = 0, C = 1): puro emaranhamento --- a \textbf{moeda} que o mate sorteia. O par colapsa em cadeia.
\end{itemize}

\subsection{ CORPO ECONÓMICO --- a bola de neve}
O tempo é o tabuleiro; o principal P é a peça em campo.
\begin{itemize}
\item \textbf{Lance $\oplus$ (juro simples):} ganho FIXO por turno, +P$\cdot$\allowbreak{}r. Uma reta segura. Engrossa devagar, \textbf{nunca vaza}.
\item \textbf{Lance $\otimes$ (juro composto):} ganho que \textbf{rende sobre si mesmo}. A partir do 3º turno (t $\geq$ 2) \textbf{o composto DOMINA o simples} --- por convexidade da exp. A bola de neve descola da reta.
\item \textbf{Contra-lance: o DESCONTO} (o flip $\nu$) --- trazer valor do futuro ao presente. Descontar \textbf{desfaz capitalizar exatamente}.
\item Quem deixa FP < 1 abre um \textbf{spread}, e o adversário sangra-o turno a turno.
\end{itemize}

\subsection{ CORPO ÁUREO $\mathbb{Z}$[$\varphi$] --- \textbf{A DOBRA} (o lance-assinatura do reino)}
Uma peça áurea carrega uma cadeia \textbf{infinita} de golpes que decaem: $\varphi$$^-$$^1$, $\varphi$$^-$$^2$, $\varphi$$^-$$^3$, \ldots --- o eco que nunca acaba (a torre de Koch/Fibonacci).

Num motor comum, o eco teria de ser \textbf{truncado} --- e aí o dano vaza, e a partida deixa de ser reversível.

\textbf{No Reino Dourado o Rei DOBRA a cadeia.} A identidade \textbf{$\varphi$$^2$ = $\varphi$ + 1} fecha o eco infinito em \textbf{dois inteiros}:

\begin{Verbatim}[fontsize=\small,frame=single,framesep=3mm]
phi^-^1 = phi - 1 = (-1, 1)
1 - phi^-^1 = (2, -1) = phi^-^2
Sigma_{j>=0} phi^-^j = 1/phi^-^2 = phi^2 = (1, 1) <- a medida do infinito
\end{Verbatim}

\textbf{O golpe infinito resolve num único tick, com dano EXATO $\varphi$$^2$ = 2,618034\ldots, e zero arredondamento.}

E as três mecânicas que dela saem:
\begin{itemize}
\item \textbf{O MOVIMENTO:} a peça áurea anda pelo operador \textbf{(a,b) $\mapsto$ (b, a+b)} --- cada turno avança a soma dos dois últimos passos (a escada de Fibonacci, a espiral). \textbf{É o único lance do tabuleiro que ACELERA sem gastar recurso} --- porque a órbita é \emph{gerada}, não somada.
\item \textbf{A PARADA:} o adversário responde com o conjugado \textbf{$\varphi$' = -1/$\varphi$} --- o escudo \textbf{(a,b) $\mapsto$ (a+b, -b)} encolhe o golpe por $\varphi$$^-$$^1$ a cada camada. Expande $\bowtie$ contrai, $\sigma$$\varphi$' = -1, e a soma total é limitada por $\varphi$$^2$. \textbf{Nada vaza.}
\item \textbf{A MOEDA:} o ouro do reino \textbf{é} $\mathbb{Z}$[$\varphi$]. Todo dano, todo recurso, todo eco é um par de inteiros. Nunca \emph{float}.
\end{itemize}

\subsection{ CORPO RACIONAL $\mathbb{Q}$ --- o lance como razão}
Toda peça se move por uma \textbf{RAZÃO}, o par (num, den): o \textbf{numerador é o AVANÇO} (a camada aditiva, o passo do Duque $\oplus$) e o \textbf{denominador é a MARCHA/TEMPO} (a camada de escala, a engrenagem da Duquesa $\otimes$). \emph{Casas por tick = num/den.}

\begin{itemize}
\item \textbf{O DUELO DE EUCLIDES (o mate):} dois lances racionais medem-se um no outro --- subtrai-se o menor do maior, repetidamente (a fração contínua). \textbf{Quem chega ao RESTO 0 dá o mate.} Um duelo entre dois racionais \textbf{sempre termina} (CF finita = mate em número finito de lances).
\item  \textbf{E O XEQUE PERPÉTUO:} contra um \textbf{METAL} (um irracional --- $\varphi$, prata, bronze; CF \textbf{infinita}) o duelo \textbf{nunca fecha}. \emph{O metal não pode ser matado pela régua racional.} \textbf{É exatamente por esse buraco que $\mathbb{Q}$ se completa em $\mathbb{R}$} --- e é por isso que os metais são os Reis.
\item \textbf{A LEI DA LIMPEZA (gcd = 1):} 2/4 e 1/2 são a \textbf{mesma casa}. Inflar as duas camadas ($\times$\allowbreak{}k) não gera vantagem nenhuma. \textbf{Não há energia de graça.}
\end{itemize}

\vspace{2mm}\hrule\vspace{3mm}

\section{10. Um turno completo, do início ao fim}

Para fechar, um turno lido nas três camadas --- que é como o reino pensa.

\begin{tcolorbox}[colback=ouroclaro!35,colframe=ouro,boxrule=0pt,leftrule=2pt,arc=0pt,left=4mm,right=3mm,top=2mm,bottom=2mm,breakable]\itshape \textbf{PRIMEIRIDADE / ÍCONE} --- \emph{a qualidade.} Antes de qualquer lance, a peça \textbf{é} alguma coisa: um rumo $\zeta$$_8$$^k$, uma impedância $\sigma$, uma frequência p, um perfil h. Pura qualidade, sem reação. É o teu \textbf{estado}. \textbf{SECUNDIDADE / ÍNDICE} --- \emph{o embate.} O lance $\oplus$ \textbf{junta o par}: a tua peça contra a dele, E contra B, num contra den, testa contra refuta. É aqui que há \textbf{combate} --- a conexão factual, a reação, o carry que força o tique seguinte. \textbf{O ataque e a defesa vivem na Secundidade.} \textbf{TERCEIRIDADE / SÍMBOLO} --- \emph{a lei.} O operador $\prod$ \textbf{medeia}: dá o inverso, casa a impedância, converte a escala, escolhe o lance (o dispatch), e no fim \textbf{atesta}. Todo turno termina com as Princesas a verificar o ponteiro. \textbf{A vitória vive na Terceiridade.}\end{tcolorbox}

E a sequência mecânica, sempre a mesma:

\begin{Verbatim}[fontsize=\small,frame=single,framesep=3mm]
1. DECLARA o corpo em que jogas (ele fixa a tua régua mu e o teu relógio Sigma=Z/n)
2. ESCOLHE a operação: (+) (somar/avançar) . (x) (compor/girar/ganhar) . (op) (converter/inverter)
3. COMPILA o lance na ISA ERG-64 (i64, exato -- o chessc emite WASM e Dafny da MESMA IR)
4. MEDE o dano pela régua do corpo (tipicamente z (x) z = |z|^2)
5. CASA Gamma = 0, FP = 1 -- se vazou, o inversor (op) recoloca (e cobra o turno)
6. ATESTA o ponteiro conforme: -½ (tem dentes) . 0 (queda de bandeira: lance vazio)
7. HALT resíduo 0.
\end{Verbatim}

\vspace{2mm}\hrule\vspace{3mm}

\section{ As cinco regras que cabem numa carta}

\begin{enumerate}
\item \textbf{Um lance é uma operação num corpo.} Nada mais existe.
\item \textbf{Só há três verbos: $\oplus$ $\otimes$ $\prod$.} Todo o resto é uma composição deles.
\item \textbf{Nenhum universo pode vazar: FP = 1.} O lance descasado devolve a energia ao atacante.
\item \textbf{Nada de float no meio da partida.} O estado é inteiro e exato; o pixel vem depois.
\item \textbf{O mate não mata --- cristaliza.} A partida acaba quando o resíduo chega a zero.
\end{enumerate}

\begin{tcolorbox}[colback=ouroclaro!35,colframe=ouro,boxrule=0pt,leftrule=2pt,arc=0pt,left=4mm,right=3mm,top=2mm,bottom=2mm,breakable]\itshape \emph{O Rei é o ponto fixo que se auto-gera: \textbf{$\varphi$ = 1 + 1/$\varphi$}. É o único que se sustenta sozinho --- e é por isso que ele reina.}\end{tcolorbox}

\vspace{2mm}\hrule\vspace{3mm}

\chapter{O contrato e a abertura de universo --- o casamento}

\begin{tcolorbox}[colback=ouroclaro!35,colframe=ouro,boxrule=0pt,leftrule=2pt,arc=0pt,left=4mm,right=3mm,top=2mm,bottom=2mm,breakable]\itshape \emph{"Nenhum universo pode vazar."} --- a primeira e única lei do Reino Dourado\end{tcolorbox}

\vspace{2mm}\hrule\vspace{3mm}

\section{1. O que é ``abrir um universo''}

No Reino Dourado você não escolhe um mapa: você \textbf{abre um corpo}.

Um corpo é um universo jogável --- um tabuleiro com a sua própria régua, o seu próprio relógio, as suas próprias peças. O corpo \textbf{fractal} $\mathbb{Q}$($\zeta$$_8$) tem oito rumos e gira de 45$^\circ$ em 45$^\circ$. O corpo \textbf{áureo} $\mathbb{Z}$[$\varphi$] dobra o infinito em dois inteiros. O corpo \textbf{eletromagnético} propaga em (E,B) e mata no cone nulo. O corpo \textbf{econômico} é o tempo, e a peça é o dinheiro que cresce.

Mas nenhum deles entra em jogo por decreto. \textbf{Um corpo só abre se assinar o Contrato.}

O Contrato é o crivo do reino: a lista de condições que um universo tem de satisfazer para poder existir sem sangrar energia. Quem assina, joga. \textbf{E todos assinam} --- porque nenhum corpo comparece sozinho: \textbf{comparece com o seu dual}, e o par fecha. O que falta a um, o outro tem.

E o Contrato não é uma metáfora: ele é \textbf{verificado}. Cada assinatura passa pelo Formalizador (o árbitro do reino), que exige de toda afirmação um \textbf{refutador com dentes} --- uma versão sabotada da regra que tem de \emph{falhar}. Se a sabotagem passa, a regra não dizia nada, e o universo não abre. As onze realizações do catálogo estão certificadas (resíduo 0), e o casamento energético com elas.

\vspace{2mm}\hrule\vspace{3mm}

\section{2. A ficha de assinatura --- os quatro campos}

Todo corpo entra com a mesma ficha, a \textbf{realização} ($\Sigma$, $T_\omega$, $\nu$, $\circ$):

{\fontsize{10.00pt}{12.00pt}\selectfont\begin{longtable}{>{\fontsize{10.00pt}{12.00pt}\selectfont\RaggedRight\hspace{0pt}\arraybackslash}p{4043412sp}>{\fontsize{10.00pt}{12.00pt}\selectfont\RaggedRight\hspace{0pt}\arraybackslash}p{5046745sp}>{\fontsize{10.00pt}{12.00pt}\selectfont\RaggedRight\hspace{0pt}\arraybackslash}p{18798899sp}}\toprule
\textbf{Campo} & \textbf{Nome de jogo} & \textbf{O que é} \\ \midrule\endhead
\textbf{$\Sigma$ = $\mathbb{Z}$/n} & o \textbf{relógio} & a órbita, o número de tiques até voltar ao início \\
\textbf{$\mu$ = $T_\omega$} & a \textbf{régua} & como o corpo mede (dano, distância, alcance, valor) \\
\textbf{$\nu$} & o \textbf{flip dual} & a complementaridade: primal $\oplus$ dual = o neutro \\
\textbf{$\circ$} & o \textbf{mate} & o estado terminal, absorvente: onde a partida fecha \\
\bottomrule\end{longtable}}

E, por cima dos quatro, o \textbf{selo}:

{\fontsize{10.00pt}{12.00pt}\selectfont\begin{longtable}{>{\fontsize{10.00pt}{12.00pt}\selectfont\RaggedRight\hspace{0pt}\arraybackslash}p{3669368sp}>{\fontsize{10.00pt}{12.00pt}\selectfont\RaggedRight\hspace{0pt}\arraybackslash}p{9372169sp}>{\fontsize{10.00pt}{12.00pt}\selectfont\RaggedRight\hspace{0pt}\arraybackslash}p{14847519sp}}\toprule
\textbf{Selo} & \textbf{Nome de jogo} & \textbf{O que é} \\ \midrule\endhead
\textbf{$\Gamma$ = 0} & o \textbf{casamento} & $\Gamma$=(Z-Z$\allowbreak{}_0$)/(Z+Z$\allowbreak{}_0$)=0 $\to$ \textbf{FP = 1}: nada vaza \\
\bottomrule\end{longtable}}

Os quatro campos dizem \emph{o que o universo é}. O selo diz \emph{que ele fecha}.

\vspace{2mm}\hrule\vspace{3mm}

\section{3. Os pontos do Contrato, um a um}

\subsection{Ponto 1 --- $\Sigma$: declare o seu relógio $\mathbb{Z}$/n}

\textbf{A regra:} o universo tem de ter uma \textbf{órbita fechada}. Um passo $\omega$ que, repetido, percorre \emph{todos} os estados e volta ao ponto de partida. Formalmente: gcd(passo, n) = 1 --- o passo é coprimo com a aridade, e a órbita s\_k = k$\cdot$\allowbreak{}passo (mod n) cobre $\mathbb{Z}$/n inteiro.

\textbf{Por que isso importa em jogo:} é o \textbf{Checkpoint}. Se a órbita fecha e o passo é reversível, então \emph{um único snapshot determina a partida inteira} --- o passado (-passo) e o futuro (+passo). Um universo com relógio assinado é \textbf{auditável}: você pode voltar o lance, você pode replayar o mate. Um universo sem relógio é uma partida que não se pode reconstituir --- e o reino não a homologa.

\textbf{Exemplos assinados:}
\begin{itemize}
\item \textbf{fractal} --- $\Sigma$ = $\mathbb{Z}$/8: as oito raízes $\zeta$$_8$, o tique é a rotação de 45$^\circ$. $\zeta$$\otimes$$\zeta$$^3$ = $\zeta$$^{4}$ = -1 (a meia-volta).
\item \textbf{cristalino} --- $\Sigma$ = $\mathbb{Z}$/n \textbf{parado}: regulador nulo, as unidades são finitas. O relógio deste corpo \textbf{não anda} --- é essa a sua defesa.
\item \textbf{espaço-temporal} ($\mathbb{Q}$$_2$) --- $\Sigma$ = a órbita do sucessor S = +1 em $\mathbb{Z}$/2\textasciicircum{}k; o corpo é o \emph{completamento} dessa órbita.
\item \textbf{relógio universal} --- $\Sigma$ = $\mathbb{Z}$/n é o próprio corpo: ele é quem \textbf{fornece} o relógio a todos os outros.
\end{itemize}

\begin{tcolorbox}[colback=ouroclaro!35,colframe=ouro,boxrule=0pt,leftrule=2pt,arc=0pt,left=4mm,right=3mm,top=2mm,bottom=2mm,breakable]\itshape  \textbf{Armadilha da máquina finita.} Truncar $\mathbb{Z}$$_2$ em $\mathbb{Z}$/2\textasciicircum{}k \emph{fabrica} torção (2$\cdot$\allowbreak{}2\textasciicircum{}\{k-1\}=0) e inventa duas raízes espúrias de u$^2$=1. É a \textbf{queda de bandeira precoce}: o instrumento fecha a órbita antes do tempo. Declare a aridade certa, ou o relógio mente.\end{tcolorbox}

\vspace{2mm}\hrule\vspace{3mm}

\subsection{Ponto 2 --- $\mu$: declare a sua régua}

\textbf{A regra:} o universo tem de dizer \textbf{como mede}. Não vale "mede-se de olho". A régua é uma função explícita, e é a mesma para dano, alcance e pontuação.

\textbf{Por que importa em jogo:} a régua é o \emph{balanceamento}. É o que impede que um golpe valha 3 num lance e 300 no seguinte.

\textbf{Exemplos assinados:}
\begin{itemize}
\item \textbf{fractal} --- $\mu$ = z$\cdot$\allowbreak{}z, a \textbf{norma pela conjugação}: o produto ($\otimes$ La Hire) do elemento pelo seu flip dual. Dá |z|$^2$. É a régua de dano/alcance.
\item \textbf{áureo $\mathbb{Z}$[$\varphi$]} --- $\mu$ = a cadeia fractal $\Sigma$$\varphi$\textasciicircum{}\{-j\}, fechada \textbf{exatamente} em dois inteiros pela identidade $\varphi$$^2$=$\varphi$+1: o resultado é o par (1,1) = $\varphi$$^2$.
\item \textbf{racional $\mathbb{Q}$} --- $\mu$ = \textbf{Euclides}: a fração contínua, "medir até o resto 0".
\item \textbf{eletromagnético} --- $\mu$ = a impedância $\sigma$ = |E|/|B|, e o telescópio $\Sigma$\allowbreak{}h = log|E$\times$\allowbreak{}B|$^2$ --- o telescópio.
\item \textbf{econômico} --- $\mu$ = a força de juros $\delta$ = log(1+r).
\item \textbf{evolutivo} --- $\mu$ = a frequência gênica numa régua hiperbólica, entre a extinção (p=0) e a fixação (p=1).
\end{itemize}

\begin{tcolorbox}[colback=ouroclaro!35,colframe=ouro,boxrule=0pt,leftrule=2pt,arc=0pt,left=4mm,right=3mm,top=2mm,bottom=2mm,breakable]\itshape  \textbf{A LEI DA RÉGUA (inegociável).} A régua \textbf{nunca é float}. Ela corre em inteiro exato na ISA ERG-64 --- $\mathbb{Q}$ (num/den) ou $\mathbb{Z}$[$\varphi$] (a+b$\varphi$). O float (GLSL/PTX) entra \textbf{só no fim, no pixel}, onde 1 ulp é invisível. Arredondar no meio da partida é vazar: \textbf{lance ilegal}.\end{tcolorbox}

\vspace{2mm}\hrule\vspace{3mm}

\subsection{Ponto 3 --- $\nu$: declare o seu flip dual}

\textbf{A regra:} o universo tem de ter uma \textbf{involução} ($\nu$$\circ$$\nu$ = id) que leva o primal ao dual, tal que \textbf{primal $\oplus$ dual = o neutro}. Na coordenada logarítmica de Pontryagin: $\nu$: k $\mapsto$ -k, e k + $\nu$(k) $\equiv$ 0.

\textbf{Por que importa em jogo:} $\nu$ é a \textbf{defesa}. É o contra-lance que anula exatamente o golpe. Um universo sem flip é um universo onde há ataques indefensáveis --- e o reino não homologa combate sem defesa.

\textbf{Exemplos assinados:}
\begin{itemize}
\item \textbf{fractal} --- $\nu$ = a conjugação z $\mapsto$ z: o \textbf{espelho}, que reflete o rumo do inimigo. E z$\otimes$\allowbreak{}z = |z|$^2$ é a régua de dano (o espelho \emph{é} a medida).
\item \textbf{racional $\mathbb{Q}$} --- $\nu$ = a recíproca [(a,b)] $\mapsto$ [(b,a)]: um ataque de razão 3/4 é anulado \emph{exatamente} por 4/3. \textbf{É esse inverso que faz $\mathbb{Q}$ virar corpo} ($\mathbb{Z}$ não o tinha).
\item \textbf{áureo} --- $\nu$ = a conjugação $\varphi$ $\mapsto$ -1/$\varphi$: a \textbf{contração} que encolhe o eco (expansão $\bowtie$ contração, $\sigma$$\sigma$' = -1).
\item \textbf{econômico} --- $\nu$ = \textbf{descontar}: trazer o valor do futuro ao presente. Descontar desfaz capitalizar, exatamente.
\item \textbf{eletromagnético} --- $\nu$ = o \textbf{flip} (E,B) $\mapsto$ (B,-E): gira o par sem perder o telescópio.
\item \textbf{motor} --- $\nu$ = o \textbf{flip} a $\mapsto$ ia (motor $\bowtie$ rotor, magnitude $\bowtie$ fase).
\item \textbf{criativo} --- $\nu$ = \textbf{gerar $\bowtie$ atestar} (o NOT, o complemento): as oito linguagens geram, as Princesas atestam.
\end{itemize}

\vspace{2mm}\hrule\vspace{3mm}

\subsection{Ponto 4 --- $\circ$: declare o seu mate}

\textbf{A regra:} o universo tem de ter um \textbf{terminal absorvente}. Um estado de onde não se sai. A partida tem de poder acabar.

\textbf{Por que importa em jogo:} sem mate, um universo é um moto-perpétuo --- e moto-perpétuo é vazamento pelo outro lado (energia que nunca fecha).

E note a assinatura do reino: \textbf{o mate não mata --- o mate cristaliza.} É o resíduo que vai a 0.

\textbf{Exemplos assinados:}
\begin{itemize}
\item \textbf{criativo} --- o mate do xadrez: a mobilidade colapsa a zero. Certificado: Evans Gambit, \code{1.\allowbreak{}e4 e5 2.\allowbreak{}Nf3 Nc6 3.\allowbreak{}Bc4 Bc5 4.\allowbreak{}b4 .\allowbreak{}.\allowbreak{}.\allowbreak{} 12.\allowbreak{}Qxf7\#} --- mate em 23 lances, a Dama (C$\oplus$\allowbreak{}Dafny) cristalizando.
\item \textbf{relógio} --- o mate é o \textbf{lapso N = 0}: o tempo \textbf{para}. O xeque-mate é o tempo que congela.
\item \textbf{eletromagnético} --- o mate é o \textbf{cone nulo} $\sigma$=1: o ponto fixo do flip, a régua reta.
\item \textbf{áureo} --- o mate é a \textbf{dobra}: toda potência do adversário reduzida a ($F_{n-1}, F_n$). Não há para onde expandir; a cadeia infinita já está fechada em dois inteiros.
\item \textbf{somático} --- o mate é a \textbf{homeostase}: o tecido para de proliferar.
\item \textbf{espaço-temporal} --- o mate é o \textbf{transporte = 0}: o carry acabou, a soma fechou.
\item \textbf{racional $\mathbb{Q}$} --- o mate é o \textbf{resto 0 de Euclides}. E aqui há um buraco famoso: contra um \textbf{metal} ($\varphi$, prata, bronze --- fração contínua \emph{infinita}) o duelo nunca fecha. É o \textbf{xeque perpétuo} --- e é exatamente por esse buraco que $\mathbb{Q}$ se completa em $\mathbb{R}$.
\end{itemize}

\vspace{2mm}\hrule\vspace{3mm}

\section{4. O SELO --- o casamento $\Gamma$ = 0 (o fator de potência unitário)}

Assinados os quatro campos, falta o selo. E o selo é o coração da lei.

\[\Gamma=\frac{Z-Z_0}{Z+Z_0}=0 \iff Z=Z_0 \iff \mathrm{FP}=1\]

\textbf{A potência transferida é 1 - |$\Gamma$|$^2$.} Casado ($\Gamma$=0): transfere \textbf{tudo}, sem eco, resíduo 0 no barramento. Descasado ($\Gamma$$\neq$\allowbreak{}0): parte da energia \textbf{reflete} --- volta como potência reativa, presa, inútil. Isso é o \textbf{vazamento}.

\textbf{A lei do reino: nenhum universo pode ter FP $\neq$ 1.} Seria energia surgindo fora do Sol.

E o Sol do Reino Dourado é um só: a \textbf{garrafa de Koch áurea}, a fonte única de harmônicos $\varphi$\textasciicircum{}\{-j\}. A conservação é a conservação: $\prod_v$ |x|\_v = 1, ou seja $\Sigma$\_v log|x|\_v = 0 --- o resíduo 0 sobre \emph{todos} os mostradores.

\subsection{O casamento, traduzido no idioma de cada corpo}

O mesmo enunciado, onze rostos. \textbf{Aprenda a reconhecê-lo --- é sempre o mesmo lance:}

{\fontsize{10.00pt}{12.00pt}\selectfont\begin{longtable}{>{\fontsize{10.00pt}{12.00pt}\selectfont\RaggedRight\hspace{0pt}\arraybackslash}p{6798904sp}>{\fontsize{10.00pt}{12.00pt}\selectfont\RaggedRight\hspace{0pt}\arraybackslash}p{10967913sp}>{\fontsize{10.00pt}{12.00pt}\selectfont\RaggedRight\hspace{0pt}\arraybackslash}p{10122239sp}}\toprule
\textbf{Corpo} & \textbf{O casamento $\Gamma$=0 / FP=1 é\ldots} & \textbf{O vazamento (FP<1) é\ldots} \\ \midrule\endhead
\textbf{eletromagnético} & o \textbf{cone nulo} $\sigma$=1 ($|E|=|B|$: toda a potência propaga) & potência reativa presa --- dano que volta \\
\textbf{geométrico} & o \textbf{cone nulo} ds$^2$=0 (a luz que corre limpa) & fonte espúria ($R_{\mu\nu}$$\neq$\allowbreak{}0) --- o poço mal cavado \\
\textbf{conforme} & X(x)$^2$ = 0 --- \textbf{toda obra entra plana}, norma nula & obra fora do cone (o inversor a recoloca) \\
\textbf{fractal} & a \textbf{norma unitária} $|z|=1$ (mover fora dela é ilegal) & $|z|\neq1$ --- o rumo que sangra \\
\textbf{áureo $\mathbb{Z}$[$\varphi$]} & o \textbf{não-vazamento por construção}: $\Sigma$$\varphi$\textasciicircum{}\{-j\}=$\varphi$$^2$ fecha exato & o float, que trunca a cadeia \\
\textbf{racional $\mathbb{Q}$} & a \textbf{ISA inteira}: resíduo 0 REAL, sem truncagem & 1 ulp por operação --- o float como impedância descasada \\
\textbf{telescópico} & o \textbf{traço zero}: $\sum h_i$ = 0 $\Longleftrightarrow$ covolume 1 & T$\neq$\allowbreak{}0: o covol estica ou contrai \\
\textbf{econômico} & a \textbf{ausência de arbitragem}: e\textasciicircum{}\{-$\delta$\allowbreak{}t\}$\cdot$\allowbreak{}e\textasciicircum{}\{$\delta$\allowbreak{}t\} = 1 & o \textbf{spread} --- o atrito que sangra valor \\
\textbf{celeste} & o \textbf{fechamento C=0} (estado produto) $\Longleftrightarrow$ FP = r = 1 & emaranhamento C>0 --- o reativo \\
\textbf{evolutivo} & a \textbf{normalização} $\sum p_i$ = 1 (a massa conserva) & perda de massa da população \\
\textbf{somático} & a \textbf{homeostase}: tr G = 0 (det exp = 1) & tr G $\neq$ 0 --- o tecido que transborda \\
\textbf{motor} & a \textbf{ressonância} $Z_L = Z_C$ (as reatâncias cancelam) & reatância descasada \\
\textbf{entrópico} & a \textbf{convexidade}: repousar na envoltória $f^{**}$ & custo acima do piso (não-convexidade) \\
\textbf{expansivo} & ser \textbf{unidade}: N(x) = $\pm$\allowbreak{}1 & N $\neq$ $\pm$\allowbreak{}1 --- resíduo positivo no círculo \\
\textbf{óptico} & a \textbf{refração casada} G=0 (sem eco na interface) & G$\neq$\allowbreak{}0 (difração) ou o ângulo crítico \\
\textbf{cristalino} & N(z)=1 (ser raiz da unidade) --- FP=1 \textbf{trivial} & (não vaza: também não expande) \\
\textbf{mórfico} & o \textbf{fechamento $\top$} (o universo cheio) & a \textbf{abertura} que o boleado $\gamma$ apaga \\
\textbf{espaço-temporal} & $\prod_v |x|_v = 1$ --- \textbf{a conservação: nu} & a bandeira que cai num mostrador \\
\textbf{criativo} & o \textbf{resíduo nulo} L$\triangle$\allowbreak{}L = $\emptyset$ (x$\oplus$\allowbreak{}x=0, o cristal) & qualquer cena gerada com FP$\neq$\allowbreak{}1 \\
\bottomrule\end{longtable}}

\subsection{E se vazar?}

Existe conserto, e ele tem nome: \textbf{o inversor multinível} --- o operador \textbf{(op) Pontryagin}. Ele casa \textbf{cada degrau} da torre de Fibonacci (os harmônicos $\varphi$\textasciicircum{}\{-j\} da bateria áurea) ao seu harmônico, e a distorção some. A marca é áurea e \textbf{autodual}: THD = $\varphi$\textasciicircum{}\{-1/2\} = FP\_dist = 1/$\sqrt{\ }$$\varphi$ $\approx$ 0,7862 --- sai direto da identidade $\varphi$$^2$=$\varphi$+1.

Mais níveis \textbf{limam} a distorção. Alinhado o inversor à bateria, o vazamento fecha --- \textbf{o casamento em toda a escada, não só na fundamental.}

\vspace{2mm}\hrule\vspace{3mm}

\section{5. O CHECKLIST DE ABERTURA (jogável)}

Antes de abrir um universo à mesa, passe a ficha por estes sete pontos. \textbf{Cada linha tem de fechar. Uma que falhe, o universo não abre.}

\subsection{ Ponto 1 --- O RELÓGIO}
\begin{tcolorbox}[colback=ouroclaro!35,colframe=ouro,boxrule=0pt,leftrule=2pt,arc=0pt,left=4mm,right=3mm,top=2mm,bottom=2mm,breakable]\itshape \textbf{Declaro $\Sigma$ = $\mathbb{Z}$/n, com passo $\omega$ coprimo de n.}\end{tcolorbox}
\begin{itemize}
\item [ ] A órbita s\_k = k$\cdot$$\omega$ (mod n) percorre \textbf{todos} os n estados?
\item [ ] O passo é \textbf{reversível} (existe $\omega$$^-$$^1$ mod n)?
\item [ ] \textbf{Teste do Checkpoint:} de um único snapshot (s\_k, k) reconstruo o passado e o futuro inteiros?
\end{itemize}

\emph{Falhou $\to$ a partida não é auditável. Não abre.}

\subsection{ Ponto 2 --- A RÉGUA}
\begin{tcolorbox}[colback=ouroclaro!35,colframe=ouro,boxrule=0pt,leftrule=2pt,arc=0pt,left=4mm,right=3mm,top=2mm,bottom=2mm,breakable]\itshape \textbf{Declaro $\mu$ = $T_\omega$, e ela mede assim: \_\_\_\_\_\_.}\end{tcolorbox}
\begin{itemize}
\item [ ] Uma \textbf{só} régua para dano, alcance e pontuação?
\item [ ] Ela roda em \textbf{inteiro exato} na ISA ERG-64 ($\mathbb{Q}$ = num/den, ou $\mathbb{Z}$[$\varphi$] = a+b$\varphi$)?
\item [ ] \textbf{Nenhum float} no meio da partida? (o float só no pixel)
\end{itemize}

\emph{Falhou $\to$ 1 ulp por lance. Isso é $\Gamma$$\neq$\allowbreak{}0. Não abre.}

\subsection{ Ponto 3 --- A TRÍADE NA ISA}
\begin{tcolorbox}[colback=ouroclaro!35,colframe=ouro,boxrule=0pt,leftrule=2pt,arc=0pt,left=4mm,right=3mm,top=2mm,bottom=2mm,breakable]\itshape \textbf{Escrevo a minha álgebra nos opcodes.}\end{tcolorbox}
\begin{itemize}
\item [ ] \textbf{(+) CLIFFORD} --- a soma, o par, a corrente no nó $\to$ \code{ADD} (ou \code{XOR} no booleano)
\item [ ] \textbf{($\times$) LA HIRE} --- o produto, a Rainha, o ganho $\to$ \code{MULT} por \textbf{double-and-add} (\code{AND}/\code{CMP}/\code{JZ}/\code{ADD}/\code{ADD}/\code{SUB}/\code{JMP}), ou \code{AND} no booleano
\item [ ] \textbf{(op) PONTRYAGIN} --- o operador que costura, $\exp\circ\Sigma\circ\log$ $\to$ \code{NOT} / o \textbf{dispatch} por opcode / a composição
\item [ ] Leis verificadas: \textbf{associatividade}, \textbf{distributividade}, \textbf{norma multiplicativa} --- cada uma com o seu refutador?
\end{itemize}

\emph{Gabarito para copiar:} o dispatch das doze palavras (o anel booleano F$_2$) --- escrito por extenso no Lance \emph{As doze palavras} --- e, sobre ele, as réguas do áureo ($\mathbb{Z}$[$\varphi$]) e do racional ($\mathbb{Q}$).

\subsection{ Ponto 4 --- O FLIP DUAL $\nu$}
\begin{tcolorbox}[colback=ouroclaro!35,colframe=ouro,boxrule=0pt,leftrule=2pt,arc=0pt,left=4mm,right=3mm,top=2mm,bottom=2mm,breakable]\itshape \textbf{Declaro $\nu$, e $\nu$$\circ$$\nu$ = id.}\end{tcolorbox}
\begin{itemize}
\item [ ] É \textbf{involução} (dois lances desfazem)?
\item [ ] \textbf{primal $\oplus$ dual = o neutro} (na coord. log: k + $\nu$(k) $\equiv$ 0)?
\item [ ] Todo ataque legal tem, no corpo, a sua \textbf{defesa exata}?
\end{itemize}

\emph{Falhou $\to$ há golpes indefensáveis. Não abre.}

\subsection{ Ponto 5 --- O MATE $\circ$}
\begin{tcolorbox}[colback=ouroclaro!35,colframe=ouro,boxrule=0pt,leftrule=2pt,arc=0pt,left=4mm,right=3mm,top=2mm,bottom=2mm,breakable]\itshape \textbf{Declaro o terminal absorvente.}\end{tcolorbox}
\begin{itemize}
\item [ ] Existe um estado de onde \textbf{não se sai}?
\item [ ] Nele, o \textbf{resíduo $\to$ 0} (a cristalização)?
\item [ ] Toda partida legal \textbf{alcança} o mate em número finito de lances? \emph{(ou: declaro explicitamente o meu xeque perpétuo --- v. $\mathbb{Q}$ contra os metais)}
\end{itemize}

\subsection{ Ponto 6 --- O SELO: $\Gamma$ = 0, FP = 1}
\begin{tcolorbox}[colback=ouroclaro!35,colframe=ouro,boxrule=0pt,leftrule=2pt,arc=0pt,left=4mm,right=3mm,top=2mm,bottom=2mm,breakable]\itshape \textbf{Digo em qual idioma o meu corpo realiza o casamento.}\end{tcolorbox}
\begin{itemize}
\item [ ] Escrevi na tabela \S4 a \textbf{minha linha}: o que é aqui o cone nulo?
\item [ ] Verifiquei: a potência transferida \textbf{1 - |$\Gamma$|$^2$ = 1}?
\item [ ] Se sobra reativo: o \textbf{inversor multinível} (Pontryagin, os degraus $\varphi$\textasciicircum{}\{-j\}) o recoloca?
\end{itemize}

\emph{Falhou $\to$ energia fora do Sol. \textbf{PROIBIDO.} Não abre --- nem com licença.}

\subsection{ Ponto 7 --- O ÁRBITRO (o Formalizador)}
\begin{tcolorbox}[colback=ouroclaro!35,colframe=ouro,boxrule=0pt,leftrule=2pt,arc=0pt,left=4mm,right=3mm,top=2mm,bottom=2mm,breakable]\itshape \textbf{Submeto ao certificador.}\end{tcolorbox}
\begin{itemize}
\item [ ] Cada afirmação exibe \textbf{testa} $\checkmark$ e \textbf{refuta} $\checkmark$ (o refutador com \textbf{dentes})?
\item [ ] O \textbf{ponteiro conforme} marca \textbf{-$\tfrac12$} --- ou seja, o lance tem dentes de verdade?
\item \textbf{-$\tfrac12$} = tem dentes (b = 1,0): jogada real
\item \textbf{0} = \textbf{queda de bandeira} (b = 1,1): a afirmação \emph{é} a trivialidade --- o lance vazio, o mate que não mata
\item \textbf{-1} = falha dupla (b = 0,0)
\item [ ] \textbf{Resíduo 0.} Não "$\approx$\allowbreak{}0". \textbf{0.}
\item [ ] A costura fecha: \textbf{interp $\equiv$ WASM $\equiv$ Dafny} (a mesma IR, os mesmos estados, falhas 0)?
\end{itemize}

\vspace{2mm}\hrule\vspace{3mm}

\textbf{Sete linhas fechadas $\to$ o universo ABRE.} Passe-lhe a carta de corpo e ponha as peças no tabuleiro.

\vspace{2mm}\hrule\vspace{3mm}

\section{6. Três fichas preenchidas (exemplos de mesa)}

\subsection{ FICHA --- o corpo FRACTAL $\mathbb{Q}$($\zeta$$_8$) \emph{(a raiz do reino)}}

{\fontsize{10.00pt}{12.00pt}\selectfont\begin{longtable}{>{\fontsize{10.00pt}{12.00pt}\selectfont\RaggedRight\hspace{0pt}\arraybackslash}p{4197289sp}>{\fontsize{10.00pt}{12.00pt}\selectfont\RaggedRight\hspace{0pt}\arraybackslash}p{24216055sp}}\toprule
\textbf{} & \textbf{} \\ \midrule\endhead
\textbf{$\Sigma$} & $\mathbb{Z}$/8 --- as oito raízes $\zeta$$_8$, o tique = 45$^\circ$ \\
\textbf{$\mu$} & z$\cdot$\allowbreak{}z --- a norma pela conjugação \\
\textbf{$\nu$} & z $\mapsto$ z --- o espelho \\
\textbf{$\circ$} & o corpo fecha: todo x$\neq$\allowbreak{}0 invertível pelo operador $\prod$ \\
\textbf{$\Gamma$=0} & a norma unitária $|z|=1$ --- o rotor é o cone nulo do corpo (mover fora dele seria FP$\neq1$) \\
\textbf{Tríade} & $\oplus$ = ADD nas 4 coordenadas $\cdot$ $\otimes$ = convolução ciclotômica com $\zeta$$^{4}$=-1 (MULT+SUB) $\cdot$ $\prod$ = $\exp\circ\Sigma\circ\log$ (dá os inversos) \\
\textbf{Selo} &  as 11 realizações do catálogo --- 5/5, resíduo 0 \\
\bottomrule\end{longtable}}

\textbf{Como se joga:} cada peça carrega um \textbf{rumo} $\zeta$$_8$$^k$. Girar = $\otimes$ (multiplicar o rumo por $\zeta$$_8$: tique de 45$^\circ$). Avançar = $\oplus$ (somar coordenada a coordenada). \textbf{Promover = $\prod$} --- o peão sai do \emph{anel} (unidirecional, sem desfazer) e entra no \emph{corpo} (bidirecional, pode reverter). E a regra de ouro: \textbf{só rumos de norma unitária são lances legais.} Mover com |z|$\neq$\allowbreak{}1 \textbf{vaza} --- lance ilegal.

\subsection{ FICHA --- o corpo ÁUREO $\mathbb{Z}$[$\varphi$] \emph{(o lastro)}}

{\fontsize{10.00pt}{12.00pt}\selectfont\begin{longtable}{>{\fontsize{10.00pt}{12.00pt}\selectfont\RaggedRight\hspace{0pt}\arraybackslash}p{3778201sp}>{\fontsize{10.00pt}{12.00pt}\selectfont\RaggedRight\hspace{0pt}\arraybackslash}p{24635143sp}}\toprule
\textbf{} & \textbf{} \\ \midrule\endhead
\textbf{$\Sigma$} & a órbita do gato A$_1$=[[1,1],[1,0]]: (a,b) $\mapsto$ (b, a+b) --- Fibonacci \\
\textbf{$\mu$} & a cadeia fractal $\Sigma$$\varphi$\textasciicircum{}\{-j\} = $\varphi$$^2$ = \textbf{(1,1)}, em dois inteiros \\
\textbf{$\nu$} & $\varphi$ $\mapsto$ -1/$\varphi$ --- a contração ($\sigma$$\sigma$' = -1) \\
\textbf{$\circ$} & a \textbf{dobra}: toda potência reduzida a ($F_{n-1}, F_n$) \\
\textbf{$\Gamma$=0} & \textbf{não-vazamento por construção} --- a cadeia fecha exata; o float vazaria, $\mathbb{Z}$[$\varphi$] não \\
\textbf{Tríade} & $\oplus$ = duas ADD $\cdot$ $\otimes$ = quatro MULT ($\varphi$$^2$=$\varphi$+1 dobra o termo bd) $\cdot$ $\prod$ = o shift áureo (a,b)$\mapsto$(b,a+b) \\
\textbf{Selo} &  153 casos, 0 falhas \\
\bottomrule\end{longtable}}

\textbf{Como se joga:} a peça carrega uma cadeia \textbf{infinita} de ecos ($\varphi$\textasciicircum{}\{-1\}, $\varphi$\textasciicircum{}\{-2\}, $\varphi$\textasciicircum{}\{-3\}, \ldots). Num motor comum o eco teria de ser \textbf{truncado} --- e truncar é vazar. Aqui o Rei \textbf{dobra}: a identidade $\varphi$$^2$=$\varphi$+1 fecha o eco infinito no par (1,1), dano exato $\varphi$$^2$ = 2,618\ldots, zero arredondamento, \textbf{num único tick}. \emph{O infinito não é derrotado --- é dobrado.}

\subsection{ FICHA --- o corpo ELETROMAGNÉTICO \emph{(onde o casamento é literal)}}

{\fontsize{10.00pt}{12.00pt}\selectfont\begin{longtable}{>{\fontsize{10.00pt}{12.00pt}\selectfont\RaggedRight\hspace{0pt}\arraybackslash}p{4421050sp}>{\fontsize{10.00pt}{12.00pt}\selectfont\RaggedRight\hspace{0pt}\arraybackslash}p{23992294sp}}\toprule
\textbf{} & \textbf{} \\ \midrule\endhead
\textbf{$\Sigma$} & a órbita do flip (ordem 4) \\
\textbf{$\mu$} & $\sigma = |E|/|B|$ (a impedância); o telescópio $\Sigma h = \log|E\times B|^2$ (o telescópio) \\
\textbf{$\nu$} & (E,B) $\mapsto$ (B,-E) --- o flip \\
\textbf{$\circ$} & o \textbf{cone nulo}: o ponto fixo do flip, onde o corpo morre \\
\textbf{$\Gamma$=0} & \textbf{$\sigma$ = 1} ($|E|=|B|$: o cone nulo --- FP unitário, nada vaza) \\
\textbf{Tríade} & $\oplus$ = ADD dos log-perfis $\cdot$ $\otimes$ = as impedâncias multiplicam ao empilhar meios $\cdot$ $\prod$ = $\exp\circ\Sigma\circ\log$ (o log lineariza $\sigma$) \\
\bottomrule\end{longtable}}

\textbf{Como se joga:} o ataque máximo é o \textbf{lance casado} --- levar o alvo ao cone nulo $\sigma$=1, onde FP=1 e \textbf{toda} a potência do telescópio propaga (transferência 1-|$\Gamma$|$^2$=1): o golpe que não vaza. Um lance \textbf{descasado} ($\sigma$=$\varphi$$\neq$\allowbreak{}1) deixa potência reativa presa: \textbf{dano perdido}.

\vspace{2mm}\hrule\vspace{3mm}

\section{7. Os que não assinam por inteiro (e o que isso custa)}

O Contrato tem dentes --- e \textbf{não tem exceção}. Alguns corpos, \textbf{tomados sozinhos}, não têm um axioma nomeável (o entrópico não tem inverso aditivo; o telescópico cinde no produto). Mas \textbf{nenhum corpo entra sozinho}: entra com o seu \textbf{dual}, e o par \textbf{fecha}. O dual do \textbf{entrópico} é o \textbf{cósmico} (a expansão, \code{exp} --- sob ela o \code{max} vira soma, e a soma inverte); o dual do \textbf{telescópico} é o \textbf{celeste}. \emph{Não é falha: é a metade de um par.}

{\fontsize{10.00pt}{12.00pt}\selectfont\begin{longtable}{>{\fontsize{10.00pt}{12.00pt}\selectfont\RaggedRight\hspace{0pt}\arraybackslash}p{4659372sp}>{\fontsize{10.00pt}{12.00pt}\selectfont\RaggedRight\hspace{0pt}\arraybackslash}p{8748901sp}>{\fontsize{10.00pt}{12.00pt}\selectfont\RaggedRight\hspace{0pt}\arraybackslash}p{14480783sp}}\toprule
\textbf{Corpo} & \textbf{Onde falha} & \textbf{O que isso ensina} \\ \midrule\endhead
\textbf{entrópico} (max,+) & \textbf{não tem inverso aditivo} --- $\oplus$=max é idempotente (a$\oplus$\allowbreak{}a=a). É \textbf{semicorpo}, não corpo. & O mundo \textbf{sem fase}: não se desfaz nada, não há cura, não há cancelamento. Dano que só se acumula. \textbf{Irreversibilidade = idempotência}, a mesma frase. \\
\textbf{conforme} & tem \textbf{cone nulo} $\to$ divisores de zero. \textbf{Não é corpo}. & Mas é a \textbf{CASA DA MEDIDA}: a norma é nula, a medida passa ao \textbf{par} ($X\cdot Y = -\tfrac12\,|x-y|^2$) \\
\textbf{telescópico} (n$\geq$\allowbreak{}2) & $e_i \cdot e_j$ = 0 --- \textbf{divisores de zero} (a cisão). & O \textbf{golpe-fantasma}: bater num eixo ortogonal ao do inimigo dá dano \textbf{zero} --- a jogada some. O preço da cisão. \\
\textbf{mórfico} (n>1) & anel booleano: \textbf{a$\wedge$\allowbreak{}a = a}, a idempotência deixa o inverso sem lugar. Corpo só se n=1 (= F$_2$). & É o \textbf{dual do relógio}: tem a metade, mas \textbf{não tem a reflexão} ($\nu$=-1 colapsa na identidade em característica 2). \\
\textbf{cristalino} & assina tudo --- mas com \textbf{regulador nulo}. & O \textbf{anti-áureo}: fechar \textbf{exclui} preencher. A restrição cristalográfica só deixa passar n $\in$ \{1,2,3,4,6\}, e o \textbf{primeiro traço proibido é o áureo} (2cos(2$\pi$/5) = $\varphi$-1). \emph{O cristal proíbe exatamente o preenchedor ótimo.} \\
\bottomrule\end{longtable}}

\textbf{A lição de mesa:} \emph{fechar exclui preencher.} Um universo perfeitamente fechado (o cristal) não vaza --- e também não cresce. A escolha é sua.

\vspace{2mm}\hrule\vspace{3mm}

\section{8. A leitura --- as três categorias do Contrato}

O Contrato é ele próprio triádico. Aprenda a lê-lo:

{\fontsize{10.00pt}{12.00pt}\selectfont\begin{longtable}{>{\fontsize{10.00pt}{12.00pt}\selectfont\RaggedRight\hspace{0pt}\arraybackslash}p{5454117sp}>{\fontsize{10.00pt}{12.00pt}\selectfont\RaggedRight\hspace{0pt}\arraybackslash}p{6196135sp}>{\fontsize{10.00pt}{12.00pt}\selectfont\RaggedRight\hspace{0pt}\arraybackslash}p{16238804sp}}\toprule
\textbf{Categoria} & \textbf{Signo} & \textbf{No Contrato} \\ \midrule\endhead
\textbf{Primeiridade} & o \textbf{ÍCONE} (a semelhança) & a \textbf{qualidade própria} do corpo --- a régua nua, antes de tocar nada: a norma redonda, a auto-similaridade do fractal, o cone nulo do vácuo. É \emph{o que o universo é}, sozinho. \\
\textbf{Secundidade} & o \textbf{ÍNDICE} (a conexão factual) & o \textbf{par}, o embate: \textbf{(+) Clifford}, a soma, a corrente que junta duas coisas no nó. O flip $\nu$ como reação. O casamento $\Gamma$=0 é literalmente um \textbf{par de impedâncias} que se reconhecem (Z = Z$_0$). \\
\textbf{Terceiridade} & o \textbf{SÍMBOLO} (a lei, a convenção) & \textbf{(op) Pontryagin} --- o operador que \textbf{medeia} ($\exp\circ\Sigma\circ\log$), o hábito, a lei que costura $\oplus$ e $\otimes$. E o Contrato \textbf{inteiro} é terceiridade: é a convenção que governa quem entra. \\
\bottomrule\end{longtable}}

E \textbf{($\times$) La Hire} é a mediação --- o produto, a Rainha, o ganho que liga um ao outro.

\begin{tcolorbox}[colback=ouroclaro!35,colframe=ouro,boxrule=0pt,leftrule=2pt,arc=0pt,left=4mm,right=3mm,top=2mm,bottom=2mm,breakable]\itshape \textbf{A leitura curta:} o corpo é um \textbf{ícone} (a sua qualidade), o casamento é um \textbf{índice} (o par que casa), e o Contrato é o \textbf{símbolo} (a lei que os admite).\end{tcolorbox}

\vspace{2mm}\hrule\vspace{3mm}

\section{9. O selo do reino}

Cumpridos os sete pontos, o Formalizador estampa a ficha e o universo abre. O selo diz, sempre, a mesma coisa:

\begin{tcolorbox}[colback=ouroclaro!35,colframe=ouro,boxrule=0pt,leftrule=2pt,arc=0pt,left=4mm,right=3mm,top=2mm,bottom=2mm,breakable]\itshape \textbf{$\Sigma$ fecha. A régua é exata. $\nu$ desfaz. $\circ$ cristaliza.} \textbf{$\Gamma$ = 0. FP = 1.} \textbf{Resíduo 0.}\end{tcolorbox}

E, do outro lado do selo, a única proibição que o reino não negocia:

\begin{tcolorbox}[colback=ouroclaro!35,colframe=ouro,boxrule=0pt,leftrule=2pt,arc=0pt,left=4mm,right=3mm,top=2mm,bottom=2mm,breakable]\itshape  \textbf{Nenhum universo pode vazar.} Um corpo com FP $\neq$ 1 é energia surgindo fora do Sol. Isso não é um universo mal balanceado --- \textbf{é um universo que não existe.}\end{tcolorbox}

\vspace{2mm}\hrule\vspace{3mm}

\vspace{2mm}\hrule\vspace{3mm}

\chapter{O elenco e o bestiário dos corpos}

\begin{tcolorbox}[colback=ouroclaro!35,colframe=ouro,boxrule=0pt,leftrule=2pt,arc=0pt,left=4mm,right=3mm,top=2mm,bottom=2mm,breakable]\itshape \emph{"No Reino Dourado ninguém tem um poder. Cada personagem \textbf{é} uma operação --- e cada mundo é uma régua que essa operação sabe medir."}\end{tcolorbox}

Esta seção tem duas metades, e elas são a mesma coisa vista por dois lados:

\begin{enumerate}
\item \textbf{O ELENCO} --- quem joga. Os personagens da corte \textbf{são operadores}: o Rei é um ponto fixo, a Rainha é um produto, os Duques são a soma e a multiplicação, as Princesas são quem atesta.
\item \textbf{O BESTIÁRIO} --- onde se joga. Cada \textbf{corpo} é um mundo, e cada mundo traz a sua régua, a sua tríade e a sua jogada. Vinte mundos, vinte tabuleiros. Nenhum deles foi inventado para o jogo: cada um é uma realização certificada, resíduo 0.
\end{enumerate}

\vspace{2mm}\hrule\vspace{3mm}

\section{Parte I --- O ELENCO (personagens = operadores)}

\subsection{A TRÍADE: as três únicas coisas que qualquer personagem sabe fazer}

Tudo o que se move no Reino Dourado é uma das três:

{\fontsize{10.00pt}{12.00pt}\selectfont\begin{longtable}{>{\fontsize{10.00pt}{12.00pt}\selectfont\RaggedRight\hspace{0pt}\arraybackslash}p{3190732sp}>{\fontsize{10.00pt}{12.00pt}\selectfont\RaggedRight\hspace{0pt}\arraybackslash}p{7808866sp}>{\fontsize{10.00pt}{12.00pt}\selectfont\RaggedRight\hspace{0pt}\arraybackslash}p{8377641sp}>{\fontsize{10.00pt}{12.00pt}\selectfont\RaggedRight\hspace{0pt}\arraybackslash}p{7987529sp}}\toprule
\textbf{Símbolo} & \textbf{Nome} & \textbf{O que é} & \textbf{Na ISA ERG-64} \\ \midrule\endhead
\textbf{$\oplus$} & \textbf{CLIFFORD} --- a SOMA & o par, a corrente que se junta no nó, a PA & \code{ADD} (ou \code{XOR} no corpo booleano F$_2$) \\
\textbf{$\otimes$} & \textbf{LA HIRE} --- o PRODUTO & a tensão/ganho, a composição, a PG & \code{MULT} por \emph{double-and-add} (\code{ADD} deslocado) --- ou \code{AND} no booleano \\
\textbf{$\prod$} & \textbf{PONTRYAGIN} --- o OPERADOR & quem \textbf{costura} as outras duas: $\exp\circ\Sigma\circ\log$, o ponteiro/poinsétia & \code{NOT} / o dispatch (\code{CMP}+\code{JZ}) / a composição \\
\bottomrule\end{longtable}}

E a regra que faz a tríade fechar: \textbf{Pontryagin é a lente}. Sob o log, o produto vira soma; sob o exp, a soma vira produto. É o mesmo operador nos juros, no motor, na luz e na gravidade. Sempre que você vir alguém transformando uma bola de neve numa reta, é $\prod$ trabalhando.

\vspace{2mm}\hrule\vspace{3mm}

\subsection{A CORTE}

\subsubsection{K O REI --- $\sigma_m$, o GAP METÁLICO (o ponto fixo)}

\textbf{$\sigma_m$ = m + 1/$\sigma_m$.} O Rei é o único personagem que \textbf{se gera a si mesmo}: ele é a solução da equação que fala dele. Para m=1, $\sigma$$_1$ = \textbf{$\varphi$} (o áureo --- $\varphi$ = 1 + 1/$\varphi$). O Rei é o rei porque se sustenta sozinho.

\begin{itemize}
\item \textbf{O seu operador} é o \emph{gato} $A_m=\bigl(\begin{smallmatrix}m&1\\1&0\end{smallmatrix}\bigr)$ --- a matriz que, aplicada, expande por $\sigma_m$ a cada turno.
\item \textbf{A sua marcha} cresce a taxa constante: por posto 1, há \textbf{uma só} taxa de expansão. O Rei não diversifica --- ele avança numa direção (o autovetor $\sigma_m$) e contrai na outra (-1/$\sigma_m$).
\item \textbf{O seu ouro} é o corpo áureo $\mathbb{Z}$[$\varphi$]: o menor regulador, o lastro-base do reino.
\item \textbf{A sua marca no tabuleiro}: um lance só \emph{pontua} se aterrissa numa \textbf{unidade}, N(x) = $\pm$\allowbreak{}1 --- resíduo 0 no relógio. Tudo o mais vaza.
\end{itemize}

\subsubsection{Q A RAINHA --- LA HIRE, $\varphi$$^2$ = $\varphi$ + 1}

A Rainha \textbf{é} o produto. Onde ela passa, as coisas compõem: rotações compõem, impedâncias multiplicam, fitness pondera, juros rendem juros.

\begin{itemize}
\item \textbf{A sua lei} é $\varphi$$^2$ = $\varphi$ + 1 --- a identidade que dobra o infinito de volta em dois inteiros.
\item \textbf{Como ela executa}: nunca por mágica. O produto $\otimes$ é sempre um \textbf{double-and-add} de 64 ticks --- testa o bit (\code{AND}), se for 1 acumula (\code{ADD}), dobra a máscara e o parcial (\code{ADD}+\code{ADD}), decrementa (\code{SUB}). Exato, inteiro, resíduo 0.
\item \textbf{Num racional, ela joga duas vezes}: uma Rainha no numerador, outra no denominador.
\end{itemize}

\subsubsection{ O PRÍNCIPE --- a POTÊNCIA / a exp}

O Príncipe é a \textbf{exponencial}: $\rho$ = e\textasciicircum{}\{at\}, a capitalização, o crescimento composto, a evolução exp(tG). Onde o Duque acrescenta a mesma parcela e a Rainha compõe o mesmo fator, o Príncipe \textbf{integra o gerador}: dada a força infinitesimal $\delta$, ele entrega a trajetória inteira.

\begin{itemize}
\item No corpo econômico é ele quem faz o juro composto \textbf{dominar} o simples a partir do 3º turno (t $\geq$ 2, por convexidade da exp).
\item No corpo motor é ele quem dá a espiral (a magnitude, a seta do tempo).
\item Ele anda de braço dado com Pontryagin: o Príncipe sobe (exp), Pontryagin traduz (log).
\end{itemize}

\subsubsection{ O DUQUE $\oplus$ (Joaquim) e a DUQUESA $\otimes$ (Yasmin)}

Os dois filhos. Eles aparecem \textbf{em par}, sempre, e o par é o corpo:

\begin{itemize}
\item \textbf{O Duque $\oplus$} é a camada \textbf{aditiva}: o numerador de uma fração, o dígito de uma soma, o passo de avanço, o perfil de entropias que se empilha.
\item \textbf{A Duquesa $\otimes$} é a camada de \textbf{escala}: o denominador, o transporte (\emph{carry}), a engrenagem, a impedância que multiplica.
\end{itemize}

O que o corpo espaço-temporal torna visível de uma vez por todas: \textbf{$\oplus$ e $\wedge$ não são duas operações de conveniência --- são as duas metades de uma única soma.} O \emph{adder}:

\begin{Verbatim}[fontsize=\small,frame=single,framesep=3mm]
x + y = (x (+) y) + 2.(x and y)
  o dígito   o transporte 
 (o Duque) (a Duquesa)
\end{Verbatim}

Inseparáveis. Um não fecha sem o outro.

\subsubsection{ AS PRINCESAS --- o CORPO TÉCNICO (quem ATESTA)}

O dual exato do lado que \textbf{gera}. As Princesas não movem peça: elas \textbf{assinam} que o lance é legal. São o kernel LCF --- \textbf{Dafny} e \textbf{Isabelle}.

\begin{itemize}
\item \textbf{A régua delas} é o par (testa, refuta) $\in$ \{0,1\}$^2$, lido pelo ponteiro conforme:
\item \textbf{-$\tfrac12$} $\to$ o lance \textbf{tem dentes} (é uma jogada real).
\item \textbf{0} $\to$ \textbf{queda de bandeira} (a afirmação é a trivialidade --- o \emph{mate que não mata}).
\item \textbf{-1} $\to$ falha dupla.
\item \textbf{O refutador com dentes}: nenhuma afirmação entra no reino sem exibir \emph{como} ela poderia falhar. Trocar $\triangle$ por $\cup$ e ver a lei quebrar --- \emph{isso} é uma prova.
\end{itemize}

\textbf{O par gerar $\bowtie$ atestar é o eixo do reino.} O corpo criativo gera; as Princesas atestam. Quem gera sem atestar, vaza.

\vspace{2mm}\hrule\vspace{3mm}

\subsection{AS PEÇAS DO TABULEIRO (o corpo criativo em ação)}

No corpo criativo --- o corpo das oito linguagens --- as \textbf{peças de xadrez são literalmente os operadores}. Mover a peça é executar a operação:

{\fontsize{10.00pt}{12.00pt}\selectfont\begin{longtable}{>{\fontsize{10.00pt}{12.00pt}\selectfont\RaggedRight\hspace{0pt}\arraybackslash}p{3915891sp}>{\fontsize{10.00pt}{12.00pt}\selectfont\RaggedRight\hspace{0pt}\arraybackslash}p{13455469sp}>{\fontsize{10.00pt}{12.00pt}\selectfont\RaggedRight\hspace{0pt}\arraybackslash}p{3260452sp}>{\fontsize{10.00pt}{12.00pt}\selectfont\RaggedRight\hspace{0pt}\arraybackslash}p{6732956sp}}\toprule
\textbf{Peça} & \textbf{Operador} & \textbf{Grau} & \textbf{Executa} \\ \midrule\endhead
\textbf{Torre} & $\oplus$ Clifford (as retas) & (1,0) & C --- \code{emit\_c} \\
\textbf{Bispo} & $\otimes$ La Hire / $\wedge$ (as diagonais) & (0,1) & Dafny --- \code{dafny run} \\
\textbf{Dama} & C $\oplus$ Dafny (retas $\oplus$ diagonais) & (1,1) & Q = R + B \\
\textbf{Cavalo} & 2$\cdot$\allowbreak{}C $\oplus$ 1$\cdot$\allowbreak{}Dafny (Pontryagin \textbf{graduado}) & (2,1) & o salto \\
\textbf{Rei} & C $\oplus$ Dafny de \textbf{um passo}, mas \textbf{DUAL} das peças ($\nu$: e $\mapsto$ -e, o gap $\pm$\allowbreak{}1) & (1,1)/1 & o ponto fixo \\
\textbf{Peão} & o bit $\oplus$ com o carry $\wedge$ & --- & x+y = (x$\oplus$\allowbreak{}y) + 2(x$\wedge$\allowbreak{}y) \\
\bottomrule\end{longtable}}

E o \textbf{MATE é a cristalização}: a mobilidade (a expansão entrópica) colapsa a \textbf{0}. Resíduo 0 --- e \textbf{não é dissipação}: é a transferência ao dual (o mate que não mata). \emph{Partida certificada:} Evans Gambit, 1.e4 e5 2.Nf3 Nc6 3.Bc4 Bc5 4.b4 \ldots 12.Qxf7\# --- mate em 23 lances, com a Dama (C$\oplus$\allowbreak{}Dafny) cristalizando. A Ópera de Morphy roda em 0,6 s.

\vspace{2mm}\hrule\vspace{3mm}

\section{Parte II --- O BESTIÁRIO DOS CORPOS}

\subsection{Como se lê uma ficha}

Todo corpo do reino passou pelo \textbf{CONTRATO}. Para entrar, ele exibe cinco coisas --- e a ficha do bestiário é exatamente essas cinco:

{\fontsize{10.00pt}{12.00pt}\selectfont\begin{longtable}{>{\fontsize{10.00pt}{12.00pt}\selectfont\RaggedRight\hspace{0pt}\arraybackslash}p{5984570sp}>{\fontsize{10.00pt}{12.00pt}\selectfont\RaggedRight\hspace{0pt}\arraybackslash}p{22428774sp}}\toprule
\textbf{Campo} & \textbf{O que é} \\ \midrule\endhead
\textbf{$\Sigma$ = $\mathbb{Z}$/n} & o \textbf{relógio}: a órbita, a roda, quantos tiques até fechar \\
\textbf{Régua $\mu$ ($T_\omega$)} & \textbf{como o mundo mede} --- a norma, o custo, a valuação \\
\textbf{Flip dual $\nu$} & a complementaridade: primal + dual = o neutro \\
\textbf{Mate $\circ$} & onde o corpo fecha (e onde ele morre) \\
\textbf{$\Gamma$ = 0 / FP = 1} & o \textbf{casamento}: a condição de não-vazamento \\
\textbf{o signo} & a categoria dominante: Ícone (1º) / Índice (2º) / Símbolo (3º) \\
\textbf{Jogada} & a mecânica que esse mundo entrega ao jogador \\
\bottomrule\end{longtable}}

\subsection{A LEI QUE VALE EM TODOS OS VINTE MUNDOS}

\begin{tcolorbox}[colback=ouroclaro!35,colframe=ouro,boxrule=0pt,leftrule=2pt,arc=0pt,left=4mm,right=3mm,top=2mm,bottom=2mm,breakable]\itshape \textbf{$\Gamma$ = (Z - Z$\allowbreak{}_0$)/(Z + Z$\allowbreak{}_0$) = 0 $\Longleftrightarrow$ FP = 1.} \textbf{NENHUM UNIVERSO PODE VAZAR.} Não existe lance legal com fator de potência $\neq$ 1. Seria energia surgindo fora do Sol.\end{tcolorbox}

Em cada mundo essa lei tem um \textbf{nome local} --- cone nulo, traço zero, ausência de arbitragem, homeostase, transporte nulo. É sempre a mesma lei. Aprender os vinte nomes dela \textbf{é} aprender o jogo.

E a segunda lei, a régua: \textbf{o cálculo é EXATO} ($\mathbb{Q}$ = num/den, ou $\mathbb{Z}$[$\varphi$]) na ISA i64. \textbf{O float só no FIM, no pixel.} Arredondar no meio da partida é vazar --- lance ilegal.

\vspace{2mm}\hrule\vspace{3mm}

\section{CASA DA RAIZ --- a origem e as réguas exatas}

\subsection{ O CORPO FRACTAL --- $\mathbb{Q}$($\zeta$$_8$) --- \emph{a raiz do reino}}

\begin{tcolorbox}[colback=ouroclaro!35,colframe=ouro,boxrule=0pt,leftrule=2pt,arc=0pt,left=4mm,right=3mm,top=2mm,bottom=2mm,breakable]\itshape O tabuleiro que se parece consigo em toda escala.\end{tcolorbox}

Um elemento é z = a + b$\zeta$ + c$\zeta$$^2$ + d$\zeta$$^3$, com \textbf{$\zeta$$^{4}$ = -1} ($\zeta$$_8$ = e\textasciicircum{}\{i$\pi$/4\}). O anel $\mathbb{Z}$[$\zeta$$_8$] \textbf{sobe a corpo} pelo operador.

\begin{itemize}
\item \textbf{$\Sigma$} = $\mathbb{Z}$/8 --- as 8 unidades, o relógio de 45$^\circ$.
\item \textbf{Régua $\mu$} = \textbf{z$\cdot$\allowbreak{}z} --- a norma pela conjugação. Nas esferas de divisão ela é \textbf{multiplicativa}: |xy| = |x|$\cdot$|y| (Hurwitz para em $\mathbb{O}$).
\item \textbf{$\oplus$} soma componente a componente: (a,b,c,d)$\oplus$(a',b',c',d'). Na ISA: quatro \code{ADD} independentes, i64, exato.
\item \textbf{$\otimes$} é a convolução ciclotômica com $\zeta$$^{4}$ = -1: c$\allowbreak{}_0$ = a$\allowbreak{}_0$b$\allowbreak{}_0$ - (a$\allowbreak{}_1$b$\allowbreak{}_3$+a$\allowbreak{}_2$b$\allowbreak{}_2$+a$\allowbreak{}_3$b$\allowbreak{}_1$), etc. Na ISA: \code{MULT} + \code{SUB} (o sinal de $\zeta$$^{4}$ = -1).
\item \textbf{$\prod$} dá os \textbf{inversos} --- e é \emph{só por isso} que $\mathbb{Z}$[$\zeta$$_8$] (onde 2 não tem inverso) vira $\mathbb{Q}$($\zeta$$_8$) (onde 2 $\mapsto$ $\tfrac12$). Sem $\prod$, anel; com $\prod$, corpo.
\item \textbf{$\nu$} = a conjugação z $\mapsto$ z. \textbf{Mate}: o corpo fecha.
\item \textbf{$\Gamma$=0}: |z| = 1. Mover com |z| $\neq$ 1 \textbf{vaza}.
\item \textbf{o signo}: \textbf{Primeiridade / Ícone} --- a auto-semelhança, a semelhança pura.
\end{itemize}

\begin{tcolorbox}[colback=ouroclaro!35,colframe=ouro,boxrule=0pt,leftrule=2pt,arc=0pt,left=4mm,right=3mm,top=2mm,bottom=2mm,breakable]\itshape \textbf{JOGADA.} Cada peça carrega um \textbf{RUMO} = uma das 8 unidades $\zeta$$_8$$^k$. $\bullet$ \textbf{GIRAR} = $\otimes$: multiplicar o rumo por $\zeta$$_8$ vira a peça 45$^\circ$ (um tique). $\zeta$$\otimes$$\zeta$$^3$ = -1 é a meia-volta. $\bullet$ \textbf{AVANÇAR} = $\oplus$: somar deslocamento coordenada a coordenada. $\bullet$ \textbf{PROMOVER} = $\prod$: concede o \textbf{inverso} --- o peão sai do ANEL (unidirecional, sem \emph{undo}) e vira peça do CORPO (pode reverter). $\bullet$ \textbf{ESPELHAR} = $\nu$: reflete o rumo do inimigo. E \textbf{z$\otimes$\allowbreak{}z = |z|$^2$ é a régua de dano/alcance}. $\bullet$ \textbf{REGRA DE CONSERVAÇÃO}: só rumos de norma 1 são legais. $\bullet$ \textbf{ATRATOR DE NEWTON}: uma peça lançada converge a uma das 8 raízes (\emph{homing}). Se cai na \textbf{fronteira fractal} (o irracional), \textbf{nunca assenta} --- casa instável, contestada.\end{tcolorbox}

\vspace{2mm}\hrule\vspace{3mm}

\subsection{ O CORPO ÁUREO $\mathbb{Z}$[$\varphi$] --- \emph{o infinito construído}}

\begin{tcolorbox}[colback=ouroclaro!35,colframe=ouro,boxrule=0pt,leftrule=2pt,arc=0pt,left=4mm,right=3mm,top=2mm,bottom=2mm,breakable]\itshape O infinito não é derrotado. É \textbf{DOBRADO}.\end{tcolorbox}

Um número é um \textbf{par de inteiros} (a,b) = a + b$\varphi$. E $\varphi$$^2$ = $\varphi$ + 1 fecha \textbf{toda a cadeia fractal infinita} $\Sigma$$\varphi$\textasciicircum{}\{-j\} em dois inteiros:

\begin{Verbatim}[fontsize=\small,frame=single,framesep=3mm]
Sigma_{j>=0} phi^-j = 1/(1 - phi^-1)
phi^-1 = phi-1 = (-1,1) ; 1 - phi^-1 = (2,-1) = phi^-2 ; 1/phi^-2 = phi^2 = (1,1)
 A medida do infinito é o par (1,1).
\end{Verbatim}

\begin{itemize}
\item \textbf{Régua $\mu$}: $\mathbb{Z}$[$\varphi$] mede a torre de Koch/Fibonacci \textbf{exatamente}. Não somando termo a termo (o i64 estoura, o float vaza) --- \textbf{fechando por identidade}.
\item \textbf{$\oplus$}: (a+c) + (b+d)$\varphi$ --- \textbf{dois \code{ADD}}. O infinito soma-se com dois ADD.
\item \textbf{$\otimes$} (a Rainha): (ac+bd) + (ad+bc+bd)$\varphi$ --- o termo bd$\cdot$$\varphi$$^2$ é \textbf{dobrado} e cai metade em cada trilho. Na ISA: quatro \code{MULT} double-and-add. \emph{Certificado exato em 153 casos, 0 falhas.}
\item \textbf{$\prod$}: multiplicar por $\varphi$ = (0,1) --- \textbf{o shift áureo} (a,b) $\mapsto$ (b, a+b). É o gato A$\allowbreak{}_1$. Aplicá-lo n vezes é $\varphi$$^n$ = F\_\{n-1\} + F\_n$\cdot$$\varphi$: a órbita \textbf{inteira} do infinito gerada por um só lance, sem custo de MULT.
\item \textbf{$\nu$}: $\varphi$ $\mapsto$ -1/$\varphi$ (expansão $\bowtie$ contração, $\sigma$$\sigma$' = -1).
\item \textbf{$\Gamma$=0}: \textbf{não-vazamento por construção}. A cadeia fecha exata; não há resíduo de arredondamento a escapar.
\item \textbf{o signo}: \textbf{Terceiridade / Símbolo} --- a lei que dobra todos os casos.
\end{itemize}

\begin{tcolorbox}[colback=ouroclaro!35,colframe=ouro,boxrule=0pt,leftrule=2pt,arc=0pt,left=4mm,right=3mm,top=2mm,bottom=2mm,breakable]\itshape \textbf{JOGADA --- A DOBRA.} Uma peça carrega uma cadeia \textbf{infinita} de golpes que decaem ($\varphi$$^-$$^1$, $\varphi$$^-$$^2$, $\varphi$$^-$$^3$, \ldots o eco que nunca acaba). Num motor comum, o eco teria de ser truncado --- e a partida deixaria de ser reversível. Aqui o Rei \textbf{dobra}: o eco infinito resolve \textbf{num único tick}, dano exato $\varphi$$^2$ = 2,618\ldots, zero arredondamento. $\bullet$ \textbf{MOVIMENTO}: a peça áurea anda por (a,b)$\mapsto$(b,a+b) --- cada turno avança a soma dos dois últimos passos. \textbf{É o único lance que acelera sem gastar recurso}, porque a órbita é \emph{gerada}, não somada. $\bullet$ \textbf{DEFESA}: o adversário responde com o conjugado $\varphi$' = -1/$\varphi$ --- o escudo (a,b)$\mapsto$(a+b,-b) \textbf{encolhe} o golpe por $\varphi$$^-$$^1$ a cada camada. Expande $\bowtie$ contrai; a soma total é limitada por $\varphi$$^2$. Nada vaza. $\bullet$ \textbf{A MOEDA}: o ouro do reino \textbf{é} $\mathbb{Z}$[$\varphi$]. Todo dano é um par de inteiros --- logo toda partida é \emph{replayável bit a bit} e todo mate é auditável. $\bullet$ \textbf{O MATE}: a posição em que toda a potência do adversário foi reduzida a ($F_{n-1}, F_n$). Não há para onde expandir.\end{tcolorbox}

\vspace{2mm}\hrule\vspace{3mm}

\subsection{ O CORPO RACIONAL $\mathbb{Q}$ --- \emph{a linguagem do universo}}

Um racional é um \textbf{par}: (num, den) --- a camada aditiva \textbf{a} (o Duque $\oplus$) sobre a camada de escala \textbf{b$\neq$\allowbreak{}0} (a Duquesa $\otimes$). A razão a/b é a solução de b$\otimes$\allowbreak{}x = a --- o quociente que $\mathbb{Z}$ não tinha. É a terceira torre da cascata \textbf{$\mathbb{N}$ (monoide) $\to$ $\mathbb{Z}$ (anel) $\to$ $\mathbb{Q}$ (corpo)}, e o que faz virar corpo é o \textbf{inverso multiplicativo}.

\begin{itemize}
\item \textbf{Régua $\mu$ = EUCLIDES}: a fração contínua a/b = [a$_0$; a$\allowbreak{}_1$,\ldots,a$\allowbreak{}_n$] --- "medir até o resto 0". É \textbf{invariante de escala} (CF(a/b) = CF(ka/kb)): colapsa a classe inteira, é única, não é escolha. \emph{Critério:} o racional é aquele cuja medição \textbf{TERMINA}. O irracional nunca fecha (CF infinita = um \textbf{METAL} --- e os metais completam $\mathbb{Q}$ $\to$ $\mathbb{R}$).
\item \textbf{$\oplus$}: (ad + bc)/(bd) --- três \code{MULT} + um \code{ADD}.
\item \textbf{$\otimes$} (a Rainha, duas vezes): (ac)/(bd) --- duas \code{MULT}. \emph{Certificado: 202 casos, 0 falhas.}
\item \textbf{$\prod$}: o inverso [(a,b)]$^-$$^1$ = [(b,a)] --- trocar as camadas. Na ISA, \textbf{puro \code{LOAD}/\code{STORE}}, zero aritmética. É \emph{isso} que faz virar corpo.
\item \textbf{o signo}: \textbf{Secundidade / Índice} --- o corpo do PAR. O denominador só significa \emph{contra} o numerador.
\end{itemize}

\begin{tcolorbox}[colback=ouroclaro!35,colframe=ouro,boxrule=0pt,leftrule=2pt,arc=0pt,left=4mm,right=3mm,top=2mm,bottom=2mm,breakable]\itshape \textbf{JOGADA.} Toda peça se move por uma \textbf{RAZÃO}: o numerador é o \textbf{avanço}, o denominador é a \textbf{marcha/tempo}. Casas por tique = num/den. $\bullet$ \textbf{A CARGA DA RAINHA}: amplificar um lance por c/d não é "multiplicar" --- a peça executa \textbf{64 micro-lances bit a bit} (testa o bit $\to$ soma deslocado $\to$ dobra). Exato: o combate acumula \textbf{sem vazar um ulp}. $\bullet$ \textbf{O CONTRA-GOLPE PELA RECÍPROCA}: um ataque de razão 3/4 é anulado \textbf{exatamente} por 4/3. Primal $\otimes$ dual = 1. Em $\mathbb{Q}$, \textbf{todo ataque não-nulo tem defesa exata} --- em $\mathbb{Z}$, não tinha. $\bullet$ \textbf{O DUELO DE EUCLIDES (o mate)}: dois lances medem-se um no outro, subtraindo o menor do maior. \textbf{Quem chega ao resto 0 dá o mate.} Duelo entre racionais \textbf{sempre} termina. Contra um \textbf{METAL} ($\varphi$, prata, bronze --- CF infinita) o duelo \textbf{nunca fecha}: é o \textbf{xeque perpétuo}. É por esse buraco que $\mathbb{Q}$ se completa em $\mathbb{R}$. $\bullet$ \textbf{A LEI DA LIMPEZA (gcd=1)}: 2/4 e 1/2 são a \textbf{mesma casa}. Inflar as duas camadas não gera vantagem: \textbf{não há energia de graça.}\end{tcolorbox}

\vspace{2mm}\hrule\vspace{3mm}

\subsection{ O RELÓGIO UNIVERSAL --- \emph{o espectro de réguas}}

O corpo métrico/temporal: o corpo universal \textbf{flipado}, lido na coordenada de baixo (a rapidez). Ele \textbf{classifica} todo o reino: dá o $\Sigma$ = $\mathbb{Z}$/n a cada corpo.

\begin{itemize}
\item \textbf{Régua $\mu$ = o LAPSO N($\psi$) = cos $\psi$ = $\sqrt{\ }$\allowbreak{}a = sech $\lambda$} --- o único fator temporal compatível com a métrica de curvatura constante. A taxa vale \textbf{1} no ponto comutativo (relógio máximo) e \textbf{0} no quatérnio (o relógio \textbf{para}).
\item \textbf{$\oplus$ = a adição relativística de velocidades}: s$_1$$\oplus$\allowbreak{}s$_2$ = (s$\allowbreak{}_1$+s$\allowbreak{}_2$)/(1+s$\allowbreak{}_1$s$\allowbreak{}_2$), na ilha (-1,1). Sob o flip $\Lambda$ = artanh, vira \code{ADD} \textbf{puro}.
\item \textbf{$\otimes$} = a composição dos endomorfismos da reta de rapidez. Curiosidade jogável: \textbf{o "um" não é 1} --- o neutro é tanh(1) = 0,761594\ldots
\item \textbf{$\prod$ = $\Lambda$ = artanh} --- o endireitador: leva (I, $\oplus$, $\odot$) em ($\mathbb{R}$, +, $\cdot$).
\item \textbf{$\Gamma$=0 / Mate}: onde casa, \textbf{o tempo não corre}. N = 0 $\Longleftrightarrow$ |s| = 1 $\Longleftrightarrow$ álgebra associativa: o cristal.
\item \textbf{o signo}: \textbf{Terceiridade / Símbolo} --- a régua-de-réguas, o hábito que ordena os corpos por aridade.
\end{itemize}

\begin{tcolorbox}[colback=ouroclaro!35,colframe=ouro,boxrule=0pt,leftrule=2pt,arc=0pt,left=4mm,right=3mm,top=2mm,bottom=2mm,breakable]\itshape \textbf{JOGADA.} O relógio é o \textbf{contador de turnos} --- a roda $\mathbb{Z}$/n do tabuleiro (2\textasciicircum{}128 estados = 8$\times$\allowbreak{}8 casas $\times$ 2 \emph{bitboards}). $\bullet$ \textbf{O LANCE = UM TIQUE}: multiplicar o estado por $\omega$. A partida é a órbita \{$\omega$$^k$$\cdot$\allowbreak{}s$_0$\}, e é \textbf{REVERSÍVEL}: $\omega$\textasciicircum{}\{-k\} destica. A órbita fecha em $\omega$$^n$ = 1 e a norma conserva. $\bullet$ No $\mathbb{Z}$/8, o tique é a rotação de 45$^\circ$ que \textbf{leva a reta da Torre na diagonal do Bispo}. $\bullet$ \textbf{COMBATE}: velocidades combinam por $\oplus$ relativística --- \textbf{nenhum combo cruza o horizonte |s| = 1}. O dano satura, mas nunca quebra a régua. $\bullet$ \textbf{O ÚNICO TIQUE SEM INVERSO É O MATE}: a cristalização, N = 0. \textbf{O xeque-mate é o tempo que congela.}\end{tcolorbox}

\vspace{2mm}\hrule\vspace{3mm}

\section{CASA DA MATÉRIA --- o que fecha e o que expande}

\subsection{ O CORPO CRISTALINO --- $\mathbb{Q}$($\sqrt{\ }$\allowbreak{}D), D<0 --- \emph{o corpo onde a medida já fechou}}

\begin{tcolorbox}[colback=ouroclaro!35,colframe=ouro,boxrule=0pt,leftrule=2pt,arc=0pt,left=4mm,right=3mm,top=2mm,bottom=2mm,breakable]\itshape \textbf{A tese: fechar exclui preencher.}\end{tcolorbox}

O corpo de regulador \textbf{NULO}. As unidades são finitas (só raízes da unidade), não há régua a percorrer, o resíduo é 0 em toda parte. Os dois reticulados do plano: \textbf{Gauss $\mathbb{Z}$[i]} (quadrado, 4 dobras) e \textbf{Eisenstein $\mathbb{Z}$[$\omega$]} (hexagonal, 6 dobras).

\begin{itemize}
\item \textbf{Régua $\mu$}: N(z) = a$^2$ + |D|b$^2$, \textbf{definida positiva} --- anula só em 0. Assinatura (2,0,0). Mas a régua \textbf{logarítmica colapsa}: regulador = log 1 = \textbf{0}. O relógio está \textbf{parado}.
\item \textbf{$\prod$} = a conjugação de Galois (a reflexão) --- a \textbf{única} involução não trivial que o contrato admite.
\item \textbf{o signo}: \textbf{Primeiridade / Ícone} --- o reticulado congelado, a rede idêntica a si mesma.
\item \textbf{O PREÇO EXATO}: a restrição cristalográfica só deixa passar n $\in$ \{1, 2, 3, 4, 6\}. \textbf{E o primeiro traço proibido é o áureo} (2cos(2$\pi$/5) = $\varphi$-1). \textbf{O cristal proíbe exatamente o preenchedor ótimo.}
\end{itemize}

\begin{tcolorbox}[colback=ouroclaro!35,colframe=ouro,boxrule=0pt,leftrule=2pt,arc=0pt,left=4mm,right=3mm,top=2mm,bottom=2mm,breakable]\itshape \textbf{JOGADA --- CRISTALIZAR.} A peça-cristal \textbf{congela uma região} do tabuleiro num reticulado de regulador nulo. Lá dentro o relógio \textbf{para} e só cabem giros de ordem \{1,2,3,4,6\}: a armadura quadrada ou hexagonal. É uma defesa que \textbf{não vaza e não expande} (FP=1 trivial). $\bullet$ \textbf{O GOLPE-CHAVE}: o \textbf{ataque áureo da Rainha} (a dobra-5, o traço $\varphi$) é \textbf{REPELIDO} --- é a difração proibida. Nenhum cristal admite simetria pentagonal. $\bullet$ \textbf{PARA QUEBRAR}: o adversário mira um feixe e lê os \textbf{picos do reticulado recíproco}. O feixe de ordem zero \textbf{atravessa sem dano} (refração); os desvios são o dano (difração). $\bullet$ \textbf{O MATE SUAVE}: no \textbf{QUASICRISTAL} o gap fecha --- o recíproco vira o módulo denso $\mathbb{Z}$ + $\varphi$$\cdot$$\mathbb{Z}$, os picos acumulam, refração e difração \textbf{colapsam}. A defesa perde o isolamento. E \textbf{o áureo é justamente quem colapsa mais devagar} (Hurwitz).\end{tcolorbox}

\vspace{2mm}\hrule\vspace{3mm}

\subsection{ O CORPO EXPANSIVO --- E\_m = $\mathbb{Q}$($\sigma_m$) --- \emph{o corpo do Rei}}

O ramo \textbf{hiperbólico} (o cristalino é o elíptico). A subálgebra que \textbf{um único gato} $A_m=\bigl(\begin{smallmatrix}m&1\\1&0\end{smallmatrix}\bigr)$ gera: por Cayley--Hamilton, A$^2$ = mA + I. Fecha por \textbf{NÃO} se completar --- sobre $\mathbb{R}$ cindiria (divisores de zero); sobre $\mathbb{Q}$ é corpo.

\begin{itemize}
\item \textbf{$\Sigma$}: G = $\mathbb{R}$/(log $\sigma_m$)$\mathbb{Z}$ --- um \textbf{círculo}. Posto 1: \textbf{UMA} órbita.
\item \textbf{Régua $\mu$}: o resíduo log|x| mod log $\sigma_m$. A unidade é o \textbf{regulador} $\ell$ = log $\sigma_m$.
\item \textbf{$\prod$} = o flip $\Lambda$ = log: no plano da álgebra, a órbita corre sobre a \textbf{hipérbole}; na coordenada logarítmica, sobre uma \textbf{reta}; módulo o regulador, sobre o \textbf{círculo} --- uma unidade por passo, resíduo nulo sempre.
\item \textbf{o signo}: \textbf{Terceiridade / Símbolo} --- "três nomes, um número": taxa de expansão = regulador = entropia topológica.
\end{itemize}

\begin{tcolorbox}[colback=ouroclaro!35,colframe=ouro,boxrule=0pt,leftrule=2pt,arc=0pt,left=4mm,right=3mm,top=2mm,bottom=2mm,breakable]\itshape \textbf{JOGADA --- A MARCHA DO REI.} Cada lance aplica o gato: a posição \textbf{expande por $\sigma_m$}. Por posto 1, \textbf{não dá para diversificar} --- a força cresce numa direção única. $\bullet$ \textbf{ATAQUE HIPERBÓLICO}: estica a peça mas \textbf{conserva a norma} N = a$^2$ - Db$^2$. N(v) = -1 vira N(Av) = +1: o lance \textbf{troca de ramo da hipérbole}, não vaza energia. $\bullet$ \textbf{O GATE DE PONTUAÇÃO}: só marca quem aterrissa numa \textbf{unidade}, N = $\pm$\allowbreak{}1 ($\Gamma$=0, FP=1). Toda posição com N $\neq$ $\pm$\allowbreak{}1 \textbf{vaza} --- e o inversor de Pontryagin (a torre de Fibonacci) a devolve ao círculo, um passo por turno. $\bullet$ \textbf{A VITÓRIA NÃO É MATAR}: é \textbf{fechar a órbita sobre o círculo, resíduo nulo --- a peça vira unidade.}\end{tcolorbox}

\vspace{2mm}\hrule\vspace{3mm}

\subsection{ O CORPO TELESCÓPICO --- $\mathbb{R}$$^n$ cindido --- \emph{a álgebra dos perfis}}

Os elementos são \textbf{perfis de entropia} h = (h$_1$,\ldots,h$\allowbreak{}_n$), com h\_i = log$\|$\allowbreak{}b\_i$\|$$^2$ (as entropias dos eixos). \textbf{NÃO é corpo} para n $\geq$ 2 --- há divisores de zero.

\begin{itemize}
\item \textbf{Régua $\mu$} = o par (\textbf{N}, \textbf{T}): a \textbf{norma} N(h) = $\Pi$ e\textasciicircum{}\{h\_i\} = det G = covol$^2$ (o produto que \textbf{telescopa}) e o \textbf{traço} T(h) = $\sum h_i$. O resíduo da régua \textbf{é} T: o desvio do covolume 1.
\item \textbf{$\otimes$} = a multiplicação \textbf{pontual} (Hadamard). É aqui que a cisão degenera: \textbf{$e_i \cdot e_j$ = 0} --- o produto pode \textbf{aniquilar}.
\item \textbf{$\prod$} = a \textbf{deflexão D\_$\lambda$} (passo constante $\lambda$ = -2 log $\sigma_m$) composta com a dualidade $\nu$ = (-1)$\circ$\allowbreak{}rev.
\item \textbf{$\Gamma$=0 = TRAÇO ZERO}: FP=1 $\Longleftrightarrow$ $\sum h_i$ = 0 $\Longleftrightarrow$ covolume 1. O que um eixo \textbf{estica} (+h\_i, o indutor) outro \textbf{contrai} (-h\_i, o capacitor).
\item \textbf{o signo}: \textbf{Terceiridade} (a lei multiplicativa que telescopa) com \textbf{Secundidade} no coração (os idempotentes que se aniquilam).
\end{itemize}

\begin{tcolorbox}[colback=ouroclaro!35,colframe=ouro,boxrule=0pt,leftrule=2pt,arc=0pt,left=4mm,right=3mm,top=2mm,bottom=2mm,breakable]\itshape \textbf{JOGADA.} O perfil é uma \textbf{coluna de eixos ortogonais} --- uma fileira de peças alinhadas. $\bullet$ \textbf{ATAQUE POR MUITOS EIXOS}: encadeia-se dano por N eixos e o total é a norma que \textbf{colapsa na soma das entropias}. Atacar por muitos eixos é \textbf{somar perfis, não multiplicar custo}. $\bullet$ \textbf{A RÉGUA BLINDADA (invulnerável)}: a \textbf{autodual}, T = 0, covolume 1 --- sem reatância líquida. $\bullet$ \textbf{O GOLPE-FANTASMA}: dois eixos ortogonais têm $e_i \cdot e_j$ = 0. \textbf{Bater num eixo ortogonal ao do inimigo dá dano ZERO} --- a jogada \emph{some}, não é trocada. É o preço da cisão. $\bullet$ \textbf{O LANCE-ÂNCORA}: a régua metálica do Rei (passo $\lambda$), centrada --- o \textbf{único} perfil de passo constante que fica autodual.\end{tcolorbox}

\vspace{2mm}\hrule\vspace{3mm}

\subsection{ O CORPO ENTRÓPICO --- o semianel tropical (max,+) --- \emph{onde o reino deixa de fechar}}

\begin{tcolorbox}[colback=ouroclaro!35,colframe=ouro,boxrule=0pt,leftrule=2pt,arc=0pt,left=4mm,right=3mm,top=2mm,bottom=2mm,breakable]\itshape \textbf{NÃO É CORPO. É SEMICORPO.} E o ponto de falha é único e nomeável: \textbf{não há inverso aditivo.}\end{tcolorbox}

Os elementos são \textbf{entropias}: h = log $\lambda$. É o dual do cristalino --- ambos vivem a T=0; o que os separa é o \textbf{sorteio}.

\begin{itemize}
\item \textbf{$\oplus$ = MAX}. É \textbf{idempotente} (a$\oplus$\allowbreak{}a = a) --- e é \emph{exatamente por isso} que não há inverso aditivo. Neutro: -$\infty$. Na ISA: o dispatch \code{CMP}+\code{JMP}. \textbf{É a soma que ESQUECE.}
\item \textbf{$\otimes$ = a ADIÇÃO real} (a+b): o produto tropical \textbf{é} o \code{ADD} inteiro. Aqui a MULT \emph{é} a soma.
\item \textbf{$\prod$ = a deflexão} D\_$\lambda$(x) = x+$\lambda$ --- e aqui ela \textbf{é endomorfismo} de $\oplus$ (a idempotência compra a linearidade da translação). O fechamento é a \textbf{repetição sem cota} --- somar o caminho a si mesmo enquanto valer ---, e ela estabiliza sse todo ciclo tem peso médio $\leq$ 0.
\item \textbf{$\Gamma$=0 = a CONVEXIDADE}: FP=1 quando a representação repousa na envoltória convexa. FP<1 é a \textbf{não-convexidade} --- custo vazado acima do piso.
\item \textbf{o signo}: \textbf{Secundidade / Índice} --- o max é o embate puro; o sorteio é o índice bruto. É o corpo onde a \textbf{Terceiridade falha}: sem inverso aditivo, sem fase, sem interferência. \emph{Sobra o índice sem símbolo.}
\end{itemize}

\begin{tcolorbox}[colback=ouroclaro!35,colframe=ouro,boxrule=0pt,leftrule=2pt,arc=0pt,left=4mm,right=3mm,top=2mm,bottom=2mm,breakable]\itshape \textbf{JOGADA --- A ZONA DA IRREVERSIBILIDADE.} $\bullet$ Cada lance $\otimes$ \textbf{soma custo} ao caminho; $\oplus$ = max escolhe a \textbf{linha dominante}. $\bullet$ \textbf{NÃO SE DESFAZ NADA.} Sem inverso aditivo não há cura, não há cancelamento. \textbf{É o mundo sem fase}: as probabilidades só somam, nunca interferem. $\bullet$ \textbf{O SORTEIO OBRIGATÓRIO}: em cada vértice de coordenação z há e\textasciicircum{}\{Ns\} estados de igual resíduo (a regra do gelo: W = 3/2). A peça \textbf{TEM de rolar o dado}. O único ponto onde a medida sorteava virou \textbf{todos} os pontos. $\bullet$ \textbf{A VITÓRIA}: o caminho de \textbf{melhor peso médio de ciclo}. $\bullet$ \textbf{E NÃO HÁ CAOS}: o gato tropical é não-expansivo (entropia topológica 0). \textbf{O esquilo morre porque a rotação exige o sinal, e não há sinal.} Mecânica: dano que só se acumula, terreno que só se atravessa, e a cada nó um dado obrigatório.\end{tcolorbox}

\vspace{2mm}\hrule\vspace{3mm}

\subsection{ O CORPO ESPAÇO-TEMPORAL --- $\mathbb{Q}$$_2$, os 2-ádicos}

A \textbf{união} do corpo mórfico (F$_2$: xor/and) com o corpo métrico (o relógio) --- não por inclusão, mas por \textbf{característica mista}. O dicionário nu:

\begin{itemize}
\item \textbf{ESPAÇO} = o corpo de resíduos (o \textbf{dígito}, a máscara, mod 2).
\item \textbf{TEMPO} = a \textbf{valuação} (o lance, o mostrador). Um elemento é x = 2\textasciicircum{}v $\cdot$ u.
\end{itemize}

\begin{itemize}
\item \textbf{Régua $\mu$}: \textbf{ULTRAMÉTRICA} --- |x|$\allowbreak{}_2$ = 2\textasciicircum{}\{-v\}, com |x+y| $\leq$ max(|x|,|y|). É o \textbf{único} corpo da série cuja régua não é arquimediana (e por isso \textbf{não tem assinatura de Sylvester}). A entropia é \textbf{um bit por lance}. A metade custa exatamente um dígito.
\item \textbf{$\prod$ = o SUCESSOR} S(x) = x+1. A órbita de S a partir de 0 é $\mathbb{Z}$, denso em $\mathbb{Z}$$_2$: \textbf{o corpo é o completamento da órbita do sucessor.} No resíduo, o sucessor \textbf{é} o \code{NOT}.
\item \textbf{$\nu$ = -1} sobrevive \textbf{só como bandeira discreta}, nunca como movimento. O Rei que não se move.
\item \textbf{$\Gamma$=0} = \textbf{literalmente a fórmula do produto} $\prod_v$ |x|\_v = 1 --- a conservação: na forma nua, sobre \textbf{todos} os lugares (infinitos mostradores, um por primo). \textbf{A bandeira que cai num mostrador é o tempo ganho no outro.}
\item \textbf{o signo}: \textbf{Terceiridade} (o tempo/valuação = a lei que conta) sobre \textbf{Primeiridade} (o dígito = o ícone) e \textbf{Secundidade} (o \textbf{transporte} = o embate que força o tique seguinte).
\end{itemize}

\begin{tcolorbox}[colback=ouroclaro!35,colframe=ouro,boxrule=0pt,leftrule=2pt,arc=0pt,left=4mm,right=3mm,top=2mm,bottom=2mm,breakable]\itshape \textbf{JOGADA.} Dois regimes que são o flip um do outro: \textbf{k<0 redondo = ESPAÇO} (o lance é a \textbf{rotação}); \textbf{k>0 hiperbólico = TEMPO} (o lance é o \textbf{boost}). O flip troca o tipo de tabuleiro. $\bullet$ \textbf{O LANCE DE SOMA} resolve-se como o \emph{adder}: propaga o transporte até ele ser \textbf{0}. \textbf{O fechamento É transporte nulo.} $\bullet$ Multiplicar por 2 = um \code{SHIFT} = \textbf{um tique do relógio}. O sucessor +1 = o \textbf{antiroque}. $\bullet$ \textbf{ ARMADILHA DA MÁQUINA FINITA}: truncar em $\mathbb{Z}$/2\textasciicircum{}k \textbf{FABRICA torção} e \textbf{quatro} raízes de u$^2$=1 (duas espúrias). É a \textbf{queda de bandeira precoce} --- o instrumento fecha a órbita antes do tempo. No limite, só sobram $\pm$\allowbreak{}1.\end{tcolorbox}

\vspace{2mm}\hrule\vspace{3mm}

\subsection{ O CORPO MÓRFICO --- (F$_2$$^n$, $\oplus$, $\wedge$) --- \emph{o dual do relógio}}

A morfologia matemática (erosão/dilatação) posta como corpo. Anel booleano (a$\wedge$\allowbreak{}a = a), logo \textbf{corpo se e só se n=1} (onde é F$_2$).

\begin{itemize}
\item \textbf{O traço}: é o \textbf{DUAL DO RELÓGIO}. Tem a \textbf{metade} (a 2-torção) mas \textbf{não tem a reflexão} ($\nu$ = -1 colapsa na identidade em característica 2). Os dois degenerados opostos do mesmo círculo: $\ell$$\to$$\infty$ é o métrico, $\ell$=2 é o mórfico.
\item \textbf{Régua $\mu$}: mede o \textbf{RESÍDUO}, não uma norma. \textbf{Sem assinatura} (o peso de Hamming não é forma quadrática).
\item \textbf{$\prod$} = a \textbf{adjunção $\delta$ $\dashv$ $\varepsilon$} (dilatação $\dashv$ erosão) e a complementação $\neg$ = \code{NOT}.
\item \textbf{$\Gamma$=0 = o fechamento $\top$} (universo cheio). FP<1 é a \textbf{abertura} --- o que o boleado apaga.
\item \textbf{o signo}: \textbf{Terceiridade / Símbolo} --- o operador posto como lei.
\end{itemize}

\begin{tcolorbox}[colback=ouroclaro!35,colframe=ouro,boxrule=0pt,leftrule=2pt,arc=0pt,left=4mm,right=3mm,top=2mm,bottom=2mm,breakable]\itshape \textbf{JOGADA --- TERRITÓRIO E ZONAS DE CONTROLE.} Cada casa é um bit (dominada=1, livre=0); cada exército é um conjunto de casas. $\bullet$ \textbf{EROSÃO $\wedge$ ($\otimes$)}: só sobrevivem as casas cuja \textbf{vizinhança inteira} já dominas. Encolhe a fronteira, \textbf{consolida o núcleo}. $\bullet$ \textbf{DILATAÇÃO $\vee$}: espalha a zona a cada casa adjacente. \textbf{O avanço de território.} $\bullet$ \textbf{ANTIROQUE $\oplus$ (XOR)}: a diferença simétrica --- \textbf{as casas contestadas por exatamente um exército}. O lance da fronteira: ao invés de proteger o rei, \textbf{cede-e-toma a linha de frente}. $\bullet$ \textbf{COMPLEMENTAÇÃO $\neg$ ($\prod$)}: vira o tabuleiro inteiro, ocupado $\leftrightarrow$ vazio. É \textbf{involução}: dois lances desfazem.\end{tcolorbox}

\vspace{2mm}\hrule\vspace{3mm}

\section{CASA DO CAMPO --- o que propaga}

\subsection{ O CORPO ELETROMAGNÉTICO --- (E,B) --- \emph{a instância física da lei}}

O telescópio \emph{é} a energia que propaga; o passo (o metal) é a \textbf{impedância} $|E|/|B|$.

\begin{itemize}
\item \textbf{Régua $\mu$}: os perfis de entropia dos eixos ortogonalizados. A soma $\Sigma$\allowbreak{}h = log|E$\times$\allowbreak{}B|$^2$ \textbf{é o telescópio} (a energia que propaga). O passo é \textbf{$\sigma$ = e\textasciicircum{}\{-$\Delta$\allowbreak{}h/2\} = |E|/|B|}. Assinatura (3,3,0) --- três redondas, três hiperbólicas.
\item \textbf{$\oplus$}: os perfis são \textbf{aditivos} --- empilhar meios \textbf{soma} os perfis.
\item \textbf{$\otimes$}: as \textbf{impedâncias multiplicam} ao empilhar meios.
\item \textbf{$\prod$} = $\exp\circ\Sigma\circ\log$: \textbf{o log LINEARIZA} a impedância (multiplicativa) no perfil (aditivo).
\item \textbf{$\Gamma$ = 0 é AQUI, literal}: é o \textbf{CONE NULO $\sigma$ = 1} (|E| = |B|) = o fator de potência unitário. \textbf{Este corpo é a instância física do critério de todo o reino.}
\item \textbf{o signo}: \textbf{Secundidade / Índice} --- o par (E,B), o embate de dois campos.
\end{itemize}

\begin{tcolorbox}[colback=ouroclaro!35,colframe=ouro,boxrule=0pt,leftrule=2pt,arc=0pt,left=4mm,right=3mm,top=2mm,bottom=2mm,breakable]\itshape \textbf{JOGADA.} Cada peça/onda carrega um par (E,B) e uma impedância $\sigma$ = |E|/|B| --- \textbf{o seu metal}. $\bullet$ \textbf{PROPAGAR} = superposição (\code{ADD} dos perfis). \textbf{ATRAVESSAR MEIOS} = multiplicar impedâncias (o casamento de terreno). $\bullet$ \textbf{O ATAQUE MÁXIMO É O LANCE CASADO}: levar o alvo ao \textbf{cone nulo $\sigma$ = 1}, onde \textbf{toda} a potência propaga (transferência = 1). \textbf{O golpe que não vaza.} $\bullet$ \textbf{UM LANCE DESCASADO} ($\sigma$ $\neq$ 1) deixa potência \textbf{reativa presa} --- energia que volta. \textbf{Dano perdido.} $\bullet$ \textbf{O CONTRA-LANCE}: o \textbf{flip} (E,B) $\mapsto$ (B,-E) --- gira o par \textbf{preservando o telescópio}: troca o eixo do ataque sem perder energia. $\bullet$ \textbf{O MATE é o cone nulo}: o ponto fixo do flip, onde o corpo morre (a régua reta).\end{tcolorbox}

\vspace{2mm}\hrule\vspace{3mm}

\subsection{ O CORPO ÓPTICO --- \emph{a fronteira e o casamento}}

\begin{tcolorbox}[colback=ouroclaro!35,colframe=ouro,boxrule=0pt,leftrule=2pt,arc=0pt,left=4mm,right=3mm,top=2mm,bottom=2mm,breakable]\itshape \textbf{A tese: a refração NÃO é uma lei de grupo. É uma INVARIÂNCIA} --- a componente tangencial, lida onde a translação sobrevive.\end{tcolorbox}

Duas realizações irmãs, o par \textbf{raio $\bowtie$ onda}: a geométrica (os raios, a reflexão total) e a ondulatória de Fourier (a difração, a lente).

\begin{itemize}
\item \textbf{Régua $\mu$}: \textbf{n$\cdot$\allowbreak{}sin$\theta$}. Régua \textbf{redonda} (k = -1), assinatura (2,0,0). É \textbf{nesta} coordenada que a refração é \textbf{linear}: n$_1$$\cdot$\allowbreak{}S$_1$ = n$\allowbreak{}_2$$\cdot$\allowbreak{}S$_2$. O horizonte é o \textbf{ângulo crítico}.
\item \textbf{$\oplus$}: a soma de tangentes ao longo de \textbf{uma} régua (compõe inclinações do mesmo meio); na onda, é a \textbf{superposição/interferência} --- e o termo cruzado \textbf{é a franja}.
\item \textbf{$\otimes$}: o escalamento pela dispersão; na onda, é a \textbf{máscara} (a pupila) --- produto no objeto $\leftrightarrow$ \textbf{convolução} no espectro.
\item \textbf{$\prod$ = A INTERFACE} (o casamento entre réguas). E na óptica de Fourier: \textbf{a LENTE É a transformada de Fourier} --- o operador que leva o objeto ao espectro, em luz.
\item \textbf{o signo}: \textbf{Secundidade / Índice} --- o embate factual entre dois meios. \emph{É o óptico, não o cristalino, que "diz refração", porque só aqui há meio.}
\end{itemize}

\begin{tcolorbox}[colback=ouroclaro!35,colframe=ouro,boxrule=0pt,leftrule=2pt,arc=0pt,left=4mm,right=3mm,top=2mm,bottom=2mm,breakable]\itshape \textbf{JOGADA --- A FRONTEIRA.} Cada território tem um \textbf{índice n} (a densidade do meio). Quando uma peça atravessa (n$_1$$\to$\allowbreak{}n$_2$), a trajetória \textbf{DOBRA} --- e o jogador \textbf{NÃO escolhe} o ângulo de saída: ele é fixado pela régua (n$\cdot$\allowbreak{}sin$\theta$ é o invariante conservado). Atravessar várias camadas é \textbf{transitivo}: o lance encadeia sem perda. $\bullet$ \textbf{O ESCUDO PERFEITO --- REFLEXÃO TOTAL INTERNA}: se o ataque incide num meio mais denso \textbf{além do ângulo crítico} ($\theta$\allowbreak{}c = 41,81$^\circ$ para 1,5$\to$\allowbreak{}1), \textbf{100\% é devolvido}. Nada penetra (a onda vira evanescente). É o flip no horizonte. $\bullet$ \textbf{O ESPELHO}: $\theta$ $\mapsto$ -$\theta$ --- \textbf{a única maneira de voltar}. $\bullet$ \textbf{TERRITÓRIO-CRISTAL}: a fronteira \textbf{DIFRATA} --- o ataque se parte em \textbf{ordens discretas}: múltiplos golpes em ângulos fixos. $\bullet$ \textbf{O CASAMENTO PERFEITO} ($\Gamma$=0): transmite todo o golpe \textbf{sem eco}. FP=1.\end{tcolorbox}

\vspace{2mm}\hrule\vspace{3mm}

\subsection{ O CORPO GEOMÉTRICO --- a curvatura como corpo}

\begin{tcolorbox}[colback=ouroclaro!35,colframe=ouro,boxrule=0pt,leftrule=2pt,arc=0pt,left=4mm,right=3mm,top=2mm,bottom=2mm,breakable]\itshape \textbf{A tese central: o potencial gravitacional é MEIA ENTROPIA com sinal trocado --- $\Phi$ = -$\tfrac12$$\cdot$\allowbreak{}log r.} E como log N \emph{é} a entropia, \textbf{o potencial soma-se exatamente onde as normas se multiplicam.}\end{tcolorbox}

Aviso ao jogador: a gravidade \textbf{escalar} (um só campo) dá o limite newtoniano mas \textbf{NÃO dá a deflexão da luz} (erra por fator \textasciitilde{}6). A \textbf{tensorial} dá Einstein. \textbf{A régua tensorial é indispensável.}

\begin{itemize}
\item \textbf{Régua $\mu$}: a métrica, assinatura (1,3,0) --- uma redonda, três hiperbólicas (que é uma assinatura de \textbf{álgebra de Clifford} Cl(1,3): \emph{a métrica é o Clifford do corpo}). E o poço: $\Phi$ = -$\tfrac12$$\cdot$\allowbreak{}log(1 + 2M/|x|).
\item \textbf{$\prod$}: \textbf{este corpo É o corpo de Pontryagin} --- $\Phi$ é literalmente log$\circ$\allowbreak{}norma. No lado tensorial, o operador é a \textbf{conexão} (Christoffel $\to$ Riemann $\to$ Ricci).
\item \textbf{$\Gamma$=0 = o CONE NULO} ds$^2$ = 0. \textbf{Mate}: o \textbf{vácuo}, $R_{\mu\nu}=0$.
\item \textbf{o signo}: \textbf{Terceiridade / Símbolo} --- a geodésica é \textbf{hábito}; Bianchi é a lei de conservação.
\end{itemize}

\begin{tcolorbox}[colback=ouroclaro!35,colframe=ouro,boxrule=0pt,leftrule=2pt,arc=0pt,left=4mm,right=3mm,top=2mm,bottom=2mm,breakable]\itshape \textbf{JOGADA --- CURVAR O TABULEIRO.} O tabuleiro tesselado \textbf{deixa de ser plano}: uma peça de massa M \textbf{cava um poço}, e as linhas de movimento das outras peças passam a seguir \textbf{geodésicas} --- curvam ao redor da massa. \textbf{O lance é pousar massa.} $\bullet$ \textbf{O COMBATE CORRE NO CONE NULO}: um ataque de luz viaja em ds$^2$ = 0 --- onde FP=1. \textbf{O disparo limpo que chega.} $\bullet$ \textbf{O POÇO BEM-FORMADO} é o vácuo $R_{\mu\nu}$ = 0: campo \textbf{sem vazamento}. $\bullet$ \textbf{O REFUTADOR COM DENTES É A MECÂNICA DE FALHA}: se a \textbf{forma} do poço quebra (trocar 1-2M/r por 1-2M/r$^2$), Ricci $\neq$ 0 --- \textbf{o poço mal-cavado tem fonte espúria e vaza}. Mas \textbf{renomear a massa não quebra} (2M$\to$\allowbreak{}3M é só um poço maior). \textbf{Só quebrar a FORMA quebra.} É assim que se distingue um buraco negro legítimo de um vazamento proibido.\end{tcolorbox}

\vspace{2mm}\hrule\vspace{3mm}

\subsection{ O CORPO CELESTE --- dois qubits --- \emph{onde i$^2$ = -1 se deriva}}

A multiplicação é \textbf{La Hire --- uma só} (a mesma do expansivo, do cristalino, do relógio). E i$^2$ = -1 \textbf{não se postula: deriva-se da reflexão $\nu$ = -1.}

\begin{itemize}
\item \textbf{Régua $\mu$}: o \textbf{campo médio}, \textbf{r$^2$ + C$^2$ = 1}, onde \textbf{r} = raio (a parte \textbf{ATIVA}) e \textbf{C} = concorrência (a parte \textbf{REATIVA}, o emaranhamento). Assinatura (2,0,0) --- a mesma régua de $\mathbb{C}$.
\item \textbf{Ponto de fundo (o Teorema de Unicidade)}: a norma é \textbf{específica da régua} (N\_k = a$^2$ - k$\cdot$\allowbreak{}c$^2$) --- \textbf{o k muda a norma, nunca a multiplicação.}
\item \textbf{$\prod$ = a reflexão} $\nu$(x) = -x. Daí: 1$\star$\allowbreak{}1 = -1. \textbf{A unidade imaginária nasce do flip.}
\item \textbf{$\Gamma$=0 = C = 0} (estado produto) $\Longleftrightarrow$ FP = r = 1.
\item \textbf{o signo}: \textbf{Secundidade / Índice} --- o corpo é o PAR; a concorrência mede a conexão factual.
\end{itemize}

\begin{tcolorbox}[colback=ouroclaro!35,colframe=ouro,boxrule=0pt,leftrule=2pt,arc=0pt,left=4mm,right=3mm,top=2mm,bottom=2mm,breakable]\itshape \textbf{JOGADA.} Cada peça tem \textbf{uma unidade de energia repartida} entre r (dano direto) e C (emaranhamento) --- e \textbf{r$^2$ + C$^2$ = 1 conserva sempre}. $\bullet$ \textbf{EMARANHAR}: enlaçar duas peças acopla-as (C>0) e \textbf{por isso BAIXA a potência individual} de cada uma. \textbf{Ganhas correlação, perdes dano solo.} Não há almoço grátis. $\bullet$ \textbf{FECHAR} (C=0): a peça isolada em \textbf{FP=1} --- potência máxima, o casamento perfeito. $\bullet$ \textbf{BELL} (r=0, C=1): puro emaranhamento, \textbf{um único bit compartilhado} --- a \textbf{MOEDA} que o mate (a medida) sorteia: o par colapsa \textbf{em cadeia}. $\bullet$ Compor dois lances vira \textbf{um só}, somando os ângulos.\end{tcolorbox}

\vspace{2mm}\hrule\vspace{3mm}

\subsection{ O CORPO MOTOR --- \emph{o dual dissipativo}}

Enquanto o corpo nervoso é conservativo e cai em \textbf{pontos fixos}, o motor \textbf{DISSIPA} e cai em \textbf{CICLOS}: a gaiola de esquilo, a rotação, o \textbf{torque}. Os autovalores do seu gerador são complexos de parte real negativa --- \textbf{a espiral}.

\begin{itemize}
\item \textbf{Régua $\mu$}: a \textbf{magnitude} ($\rho$ = e\textasciicircum{}\{at\}, a seta do tempo). No chessc, duas operações inteiras exatas: soma e produto, uma só IR com dispatch por opcode.
\item \textbf{$\otimes$ = O TORQUE} $\tau$ = $\psi$ $\times$ i (fluxo $\times$ corrente --- o telescópio agora \textbf{dando trabalho}), \textbf{máximo quando ortogonais} ($\theta$ = 90$^\circ$).
\item \textbf{$\prod$ = O DTC} (\emph{Direct Torque Control}) \textbf{é Pontryagin}: modula o torque por controle ótimo --- o estado evolui por G, o \textbf{co-estado (o valor) por G$^T$}. O princípio do máximo.
\item \textbf{$\Gamma$=0}: na \textbf{RESSONÂNCIA} as reatâncias cancelam, a impedância fica resistiva $\to$ FP=1. FP<1 é a reatância descasada.
\item \textbf{o signo}: \textbf{Secundidade / Índice} --- o embate, a força bruta que resiste e empurra.
\end{itemize}

\begin{tcolorbox}[colback=ouroclaro!35,colframe=ouro,boxrule=0pt,leftrule=2pt,arc=0pt,left=4mm,right=3mm,top=2mm,bottom=2mm,breakable]\itshape \textbf{JOGADA.} O motor é o \textbf{MOVER} da peça. $\bullet$ \textbf{O GOLPE DE FLANCO}: o torque é \textbf{MÁXIMO a 90$^\circ$}. \textbf{Não bata de frente --- bata ortogonal.} $\bullet$ \textbf{O PAINEL} é o controle de velocidade: o \emph{slider} é $\omega$, e cada tique avança \code{nova\_\allowbreak{}fase = fase + omega.\allowbreak{}dt} ($\otimes$ faz $\omega$$\cdot$\allowbreak{}dt, $\oplus$ acrescenta à fase). $\bullet$ \textbf{O DTC permite MODULAR o torque em jogo}: o co-estado (o valor futuro) \textbf{guia quanto empurrar agora}. $\bullet$ \textbf{O DIPOLO}: há \textbf{DUAS de cada peça}, em cores opostas (os dois bispos em casas de cor diferente) --- o par ressoa entre a casa e a peça.\end{tcolorbox}

\vspace{2mm}\hrule\vspace{3mm}

\section{CASA DA VIDA E DA LEI --- o que rege}

\subsection{ O CORPO SOMÁTICO --- o tecido dos muitos tabuleiros}

O \textbf{organismo distribuído}: cada célula é um \textbf{LOCAL} (um tabuleiro $\mathbb{Q}$\_p) e o somático é a conjunção deles (o \textbf{GLOBAL} --- princípio do fechamento: \emph{o local decide, o global reúne}).

\begin{itemize}
\item \textbf{Régua $\mu$}: \textbf{det(exp(tG)) = e\textasciicircum{}\{t$\cdot$\allowbreak{}tr G\}} --- \textbf{multiplicativa} (e por Hurwitz, isso faz do somático um \textbf{corpo}). O \textbf{traço} tr G (a soma das taxas) é o telescópio.
\item \textbf{$\oplus$} = a \textbf{fusão} de dois tecidos (soma direta, cada dinâmica independente).
\item \textbf{$\prod$} = $\exp\circ\Sigma\circ\log$ --- o mesmo operador da mitose, da hélice e dos juros. \textbf{O dual é o adjunto G$^T$}: o co-estado, que evolui \textbf{para trás}.
\item \textbf{$\Gamma$=0 = a HOMEOSTASE}: FP=1 $\Longleftrightarrow$ \textbf{tr G = 0} (a população conserva-se). FP<1 é tr G $\neq$ 0 --- crescimento que transborda.
\item \textbf{o signo}: \textbf{Terceiridade / Símbolo} --- a lei que rege o organismo.
\end{itemize}

\begin{tcolorbox}[colback=ouroclaro!35,colframe=ouro,boxrule=0pt,leftrule=2pt,arc=0pt,left=4mm,right=3mm,top=2mm,bottom=2mm,breakable]\itshape \textbf{JOGADA --- O TABULEIRO-EXÉRCITO.} $\bullet$ \textbf{PROLIFERAR} (o lance S, a mitose): o tecido \textbf{dobra}, 2\textasciicircum{}k, e expande até \textbf{COLIDIR} com a capacidade do espaço (inibição por contato) --- aí \textbf{para}: a homeostase. \textbf{Exatamente como os peões colidem e travam.} $\bullet$ \textbf{DIFERENCIAR} (o lance dual, $\pi$): projeta as N células nos \textbf{T TIPOS}. \textbf{Você não move célula a célula (O(N)) --- move por TIPO (O(T)).} É o atalho do grafo. $\bullet$ \textbf{FUNDIR} ($\oplus$): juntar dois exércitos é a soma direta, cada dinâmica empilhada. $\bullet$ \textbf{A RÉGUA DO COMBATE}: quem tem tr G > 0 \textbf{cresce e transborda} (FP<1) até ser corrigido; tr G = 0 conserva (FP=1, nada vaza). $\bullet$ \textbf{O co-estado G$^T$ é a REGULAÇÃO} --- o valor-sombra de cada tipo: quem decide a trajetória ótima do exército.\end{tcolorbox}

\vspace{2mm}\hrule\vspace{3mm}

\subsection{ O CORPO EVOLUTIVO --- a frequência e a seleção}

Darwin lido como \textbf{MEDIDA} (não como metafísica de "vida"): a genética de populações vira uma escada certificada que fecha em corpo sobre $\mathbb{Q}$.

\begin{itemize}
\item \textbf{Régua $\mu$}: a \textbf{frequência gênica} p $\in$ (0,1), numa régua \textbf{hiperbólica} (assinatura (1,1,0) --- a mesma do expansivo e do relógio). Os horizontes s = $\pm$\allowbreak{}1 são a \textbf{FIXAÇÃO} e a \textbf{EXTINÇÃO}.
\item \textbf{$\oplus$} = a \textbf{reprodução}: p$\star$\allowbreak{}q = (p+q)/2, \textbf{idempotente} (o equilíbrio é o ponto fixo). Composição de seleções: a adição relativística.
\item \textbf{$\otimes$} = a \textbf{ponderação pela fitness} (p\_i $\cdot$ w\_i).
\item \textbf{$\prod$ = A SELEÇÃO NATURAL} --- o replicador p\_i $\mapsto$ p\_i$\cdot$\allowbreak{}w\_i/$\langle$\allowbreak{}w$\rangle$: \textbf{SOBE} com a fitness, \textbf{DESCE} por normalização. E $\langle$\allowbreak{}w$\rangle$ é \textbf{monótona crescente} (o teorema fundamental de Fisher). \emph{Certificado por 30 gerações.}
\item \textbf{$\Gamma$=0}: a normalização --- $\sum p_i$ = 1 se mantém enquanto $\langle$\allowbreak{}w$\rangle$ sobe.
\item \textbf{o signo}: \textbf{Terceiridade / Símbolo} --- a seleção como lei/hábito estatístico.
\end{itemize}

\begin{tcolorbox}[colback=ouroclaro!35,colframe=ouro,boxrule=0pt,leftrule=2pt,arc=0pt,left=4mm,right=3mm,top=2mm,bottom=2mm,breakable]\itshape \textbf{JOGADA --- A META DO REINO.} Cada \textbf{tipo de peça} tem uma frequência p no exército, entre a \textbf{EXTINÇÃO} (p=0, a peça some do tabuleiro) e a \textbf{FIXAÇÃO} (p=1, ela domina). $\bullet$ \textbf{O LANCE}: as peças que capturaram/sobreviveram (fitness alta) \textbf{sobem} de frequência; o total conserva-se (FP=1 --- \textbf{massa constante}). $\bullet$ \textbf{ NÃO HÁ VITÓRIA GRÁTIS NOS EXTREMOS}: somar dois \emph{pushes} de seleção \textbf{não é linear} (é a régua hiperbólica). Empilhar vantagens \textbf{rende cada vez menos} perto do horizonte da fixação. $\bullet$ \textbf{A ANCORAGEM} (o flip p $\mapsto$ 1-p): troca os tipos; o único ponto fixo é p = $\tfrac12$ --- \textbf{a deriva genética como reset ao equilíbrio}. $\bullet$ \textbf{A cada geração, $\langle$\allowbreak{}w$\rangle$ só cresce}: o combate é uma \textbf{escalada monótona de qualidade}, e a seleção é \textbf{irreversível dentro de uma linhagem}.\end{tcolorbox}

\vspace{2mm}\hrule\vspace{3mm}

\subsection{ O CORPO ECONÔMICO --- os juros}

\textbf{Não traz operação nova.} Traz o par do contrato lido \textbf{na linguagem mais familiar que há: o dinheiro que cresce.} É o melhor mundo para aprender a tríade.

\begin{itemize}
\item \textbf{Régua $\mu$}: a \textbf{força de juros} $\delta$ = log(1+r) --- a rapidez, o gerador infinitesimal. (Com r = 5\%, $\delta$ $\approx$ 0,048790.)
\item \textbf{$\oplus$ = O JURO SIMPLES}: A(t) = P(1 + r$\cdot$\allowbreak{}t). Acrescenta sempre \textbf{a mesma parcela} --- a diferença é constante: uma \textbf{PA}.
\item \textbf{$\otimes$ = O JURO COMPOSTO}: A(t) = P(1+r)\textasciicircum{}t. Compõe sempre \textbf{o mesmo fator} --- a razão é constante: uma \textbf{PG}. (A potenciação é a exponenciação por quadratura: o produto realizado como \code{ADD} deslocado.)
\item \textbf{$\prod$ = A LENTE}: \textbf{baixar o índice} (log) leva a PG à PA --- \emph{a multiplicação por 1+r vira a soma de $\delta$}. \textbf{Subir o índice} (exp) devolve. \textbf{É a definição operacional de Pontryagin, em dinheiro.}
\item \textbf{$\Gamma$=0 = A AUSÊNCIA DE ARBITRAGEM}: o desconto \textbf{desfaz exatamente} a capitalização. FP<1 é o \textbf{SPREAD} --- o atrito que sangra valor.
\item \textbf{o signo}: \textbf{Terceiridade / Símbolo} --- o dinheiro é convenção pura; o juro é a lei que medeia presente e futuro.
\end{itemize}

\begin{tcolorbox}[colback=ouroclaro!35,colframe=ouro,boxrule=0pt,leftrule=2pt,arc=0pt,left=4mm,right=3mm,top=2mm,bottom=2mm,breakable]\itshape \textbf{JOGADA.} O \textbf{tempo é o tabuleiro}; o principal P é a peça posta em campo. A cada turno o jogador \textbf{escolhe a operação}: $\bullet$ \textbf{LANCE $\oplus$ (simples)}: ganho \textbf{fixo e previsível} --- uma reta segura. Engrossa devagar, \textbf{nunca vaza}. $\bullet$ \textbf{LANCE $\otimes$ (composto)}: ganho que \textbf{rende sobre si mesmo}. \textbf{A partir do 3º turno (t $\geq$ 2) o composto DOMINA o simples} --- uma bola de neve que descola da reta. $\bullet$ \textbf{O CONTRA-LANCE --- O DESCONTO} ($\nu$): trazer valor do futuro ao presente. \textbf{Descontar desfaz capitalizar exatamente} --- o retorno perfeito à origem. $\bullet$ \textbf{A VITÓRIA}: fechar o ciclo capitalizar $\bowtie$ descontar \textbf{sem vazar valor}. Quem deixa FP<1 \textbf{abre um spread}, e o adversário explora o vazamento.\end{tcolorbox}

\vspace{2mm}\hrule\vspace{3mm}

\subsection{ O CORPO CONFORME --- \emph{a casa da medida (onde vivem as Princesas)}}

\begin{tcolorbox}[colback=ouroclaro!35,colframe=ouro,boxrule=0pt,leftrule=2pt,arc=0pt,left=4mm,right=3mm,top=2mm,bottom=2mm,breakable]\itshape \textbf{NÃO é corpo} (tem cone nulo, divisores de zero). É a \textbf{CASA DA MEDIDA}, onde o certificador vive.\end{tcolorbox}

O truque: \textbf{a norma é NULA. A medida está no PAR.}

\begin{itemize}
\item O operador-interface tem X(x)$^2$ = \textbf{0} para toda obra (verificado a 1e-13): \textbf{a norma não distingue nada.}
\item \textbf{A medida passa ao EMPARELHAMENTO}: X(x)$\cdot$\allowbreak{}X(y) = \textbf{-$\tfrac12$|x-y|$^2$} --- que zera \textbf{se e só se} x = y. \textbf{Mede a distância entre obras.}
\item \textbf{A RÉGUA-RELÓGIO --- O PONTEIRO}: pt(b) = -$\tfrac12$$\cdot$\allowbreak{}d(b, b\_$\top$)$^2$. Marca \textbf{-$\tfrac12$} com dentes, \textbf{0} na queda de bandeira, \textbf{-1} na falha dupla. E |-$\tfrac12$| = $\tfrac12$ \textbf{custa exatamente um bit}: H($\tfrac12$) = 1.
\item \textbf{$\prod$} = a interface (o mergulho no cone) e a \textbf{inversão} (a involução que troca dentro e fora).
\item \textbf{$\Gamma$ = 0 é aqui literal e total}: X(x)$^2$ = 0 \textbf{É} o cone nulo, e o cone nulo \textbf{É} FP = 1.
\item \textbf{o signo}: \textbf{Secundidade / Índice} --- a medida não está na qualidade própria (norma nula = \textbf{obra sem carga}), mas no \textbf{emparelhamento de dois}.
\end{itemize}

\begin{tcolorbox}[colback=ouroclaro!35,colframe=ouro,boxrule=0pt,leftrule=2pt,arc=0pt,left=4mm,right=3mm,top=2mm,bottom=2mm,breakable]\itshape \textbf{JOGADA --- O RELÓGIO DE XADREZ E O ÁRBITRO.} Este é o corpo que \textbf{valida lances}. Toda afirmação/lance carrega \textbf{dois mostradores}: (testa, refuta). $\bullet$ O árbitro lê o \textbf{PONTEIRO}: se marca \textbf{-$\tfrac12$}, o lance \textbf{TEM DENTES} --- é jogada real, distante da trivialidade. Se marca \textbf{0}, é \textbf{QUEDA DE BANDEIRA}: o relógio parou, a afirmação \textbf{é} a trivialidade --- \emph{o mate que não mata, o lance vazio}. $\bullet$ \textbf{O mostrador diz o LADO} (qual falhou); \textbf{o ponteiro diz QUANTO.} $\bullet$ \textbf{O COMBATE EM MASSA}: cada célula é uma \emph{thread} da GPU --- uma obra de \textbf{norma nula, que não carrega carga própria}; ela só olha a \textbf{distância ao vizinho}. Troca \textbf{estéril}, discreta: \textbf{nada contamina, tudo passa pelo cone.}\end{tcolorbox}

\vspace{2mm}\hrule\vspace{3mm}

\subsection{ O CORPO CRIATIVO --- as oito linguagens --- \emph{o GERADOR}}

O corpo mais direto do reino: \textbf{a sua tríade É a própria ISA.} As oito linguagens fundamentais (\textbf{C, PTX, WASM, LLVM, GLSL, Haskell, Dafny, Isabelle}) brotam todas do \textbf{MESMO} pseudo-assembly ERG-64 e dão o \textbf{MESMO} estado.

\begin{itemize}
\item \textbf{Estrutura}: \textbf{anel booleano} F$_2$.
\item \textbf{$\oplus$ = a DIFERENÇA SIMÉTRICA $\triangle$}: neutro $\emptyset$, e \textbf{cada linguagem é o seu próprio inverso: L$\triangle$\allowbreak{}L = $\emptyset$} --- o resíduo nulo. \emph{Os dentes do refutador:} trocar $\triangle$ por $\cup$ \textbf{falha}, pois L$\cup$\allowbreak{}L = L $\neq$ $\emptyset$. \textbf{A união não tem inverso; a diferença simétrica tem.} Na ISA: \code{XOR}.
\item \textbf{$\otimes$ = a INTERSEÇÃO $\cap$}: idempotente, e \textbf{distribui sobre $\triangle$} --- é isso que faz o anel booleano. \emph{Os dentes:} pôr $\cup$ no lugar de $\cap$ \textbf{quebra a lei do anel}. Na ISA: \code{AND}.
\item \textbf{$\prod$ = NOT} (o complemento, a involução). Na teoria, o operador é \textbf{composto: o fecho universal $\circ$ Pontryagin} --- o \textbf{fecho universal} é a repetição sem cota (o resíduo que não tem teto: a paragem); Pontryagin é o \textbf{adjunto da congruência sintática} (o contexto esquerdo). \textbf{E o próprio dispatch por opcode É Pontryagin} --- a composição que costura as três.
\item \textbf{Régua $\mu$}: a \textbf{invariância sob os oito substratos} (todos dão o mesmo estado) e o resíduo nulo L$\triangle$\allowbreak{}L = $\emptyset$. Uma segunda régua acoplada: \textbf{o semianel ($\oplus$,$\otimes$) da programação dinâmica} --- \emph{troca o semianel e o MESMO algoritmo de caminho resolve outro problema} (alcança? / menor caminho / quantos caminhos / paridade).
\item \textbf{$\nu$ = gerar $\leftrightarrow$ atestar}. \textbf{Mate}: a cristalização, resíduo $\to$ 0.
\item \textbf{o signo}: \textbf{Terceiridade / Símbolo} --- a linguagem é convenção pura.
\end{itemize}

\begin{tcolorbox}[colback=ouroclaro!35,colframe=ouro,boxrule=0pt,leftrule=2pt,arc=0pt,left=4mm,right=3mm,top=2mm,bottom=2mm,breakable]\itshape \textbf{JOGADA --- ELE JOGA O XADREZ.} As peças que ele move \textbf{SÃO} os seus operadores (Torre = $\oplus$, Bispo = $\otimes$, Dama = C$\oplus$\allowbreak{}Dafny, Cavalo = 2C$\oplus$\allowbreak{}D, Rei = o dual, Peão = bit $\oplus$ carry $\wedge$). O colapso $\Psi$ é o \textbf{minimax} (a seleção de Pontryagin); a abertura segue as órbitas clássicas. \textbf{O MATE é a cristalização}: a mobilidade colapsa a 0 --- e \textbf{nada se dissipa}: transfere-se ao dual (o mate que não mata). E cada célula-jogador tem a sua réplica do grafo: \textbf{o barramento é distribuído, P2P.}\end{tcolorbox}

\vspace{2mm}\hrule\vspace{3mm}

\section{A MESA DO MESTRE --- o bestiário numa página}

{\fontsize{10.00pt}{12.00pt}\selectfont\begin{longtable}{>{\fontsize{10.00pt}{12.00pt}\selectfont\RaggedRight\hspace{0pt}\arraybackslash}p{2922358sp}>{\fontsize{10.00pt}{12.00pt}\selectfont\RaggedRight\hspace{0pt}\arraybackslash}p{5974942sp}>{\fontsize{10.00pt}{12.00pt}\selectfont\RaggedRight\hspace{0pt}\arraybackslash}p{5046954sp}>{\fontsize{10.00pt}{12.00pt}\selectfont\RaggedRight\hspace{0pt}\arraybackslash}p{3056460sp}>{\fontsize{10.00pt}{12.00pt}\selectfont\RaggedRight\hspace{0pt}\arraybackslash}p{4800364sp}>{\fontsize{10.00pt}{12.00pt}\selectfont\RaggedRight\hspace{0pt}\arraybackslash}p{4515114sp}}\toprule
\textbf{\#} & \textbf{Corpo} & \textbf{Régua $\mu$} & \textbf{o signo} & \textbf{\textbf{O nome local de FP=1}} & \textbf{A jogada de assinatura} \\ \midrule\endhead
1 & \textbf{Fractal} $\mathbb{Q}$($\zeta$$_8$) & a norma $z\bar z$ & Ícone & $|z| = 1$ & girar por $\zeta_8$ (45$^\circ$) \\
2 & \textbf{Áureo} $\mathbb{Z}$[$\varphi$] & a cadeia $\Sigma$$\varphi$$^-$$^j$ = (1,1) & Símbolo & exato por construção & a DOBRA do infinito \\
3 & \textbf{Racional} $\mathbb{Q}$ & Euclides (a CF) & Índice & resíduo 0 real, sem float & o duelo até o resto 0 \\
4 & \textbf{Relógio} & o lapso N = cos $\psi$ & Símbolo & N = 0 (o tempo para) & o tique $\mathbb{Z}$/n \\
5 & \textbf{Cristalino} & norma definida positiva & Ícone & regulador nulo & cristalizar (e proibir $\varphi$) \\
6 & \textbf{Expansivo} & $\log|x| \bmod \log\sigma_m$ & Símbolo & $N(x)=\pm1$ (ser unidade) & a marcha do Rei \\
7 & \textbf{Telescópico} & (norma, traço) & Símbolo & \textbf{traço zero} & o golpe-fantasma \\
8 & \textbf{Entrópico} (max,+) & o custo mínimo & Índice & a convexidade & o sorteio obrigatório \\
9 & \textbf{Espaço-temporal} $\mathbb{Q}$$_2$ & a ultramétrica $|x|_2$ & Símbolo & $\prod_v |x|_v = 1$ & o \emph{adder} / o transporte \\
10 & \textbf{Mórfico} F$_2$$^n$ & o resíduo & Símbolo & o fechamento $\top$ & erosão / dilatação \\
11 & \textbf{Eletromagnético} & a impedância $\sigma$ & Índice & \textbf{o cone nulo $\sigma$ = 1} & o golpe casado \\
12 & \textbf{Óptico} & n$\cdot$\allowbreak{}sin$\theta$ & Índice & G = 0 (refração casada) & a reflexão total interna \\
13 & \textbf{Geométrico} & a métrica $g_{\mu\nu}$ & Símbolo & o cone nulo ds$^2$ = 0 & cavar o poço \\
14 & \textbf{Celeste} & r$^2$ + C$^2$ = 1 & Índice & C = 0 (estado produto) & emaranhar (e perder dano) \\
15 & \textbf{Motor} & a magnitude / dissipação & Índice & a ressonância & o torque a 90$^\circ$ \\
16 & \textbf{Somático} & det(exp(tG)) & Símbolo & \textbf{a homeostase} (tr G = 0) & proliferar / diferenciar \\
17 & \textbf{Evolutivo} & a frequência p & Símbolo & $\sum p_i$ = 1 & o replicador \\
18 & \textbf{Econômico} & a força de juros $\delta$ & Símbolo & \textbf{sem arbitragem} & simples $\oplus$ vs composto $\otimes$ \\
19 & \textbf{Conforme} & o ponteiro -$\tfrac12$ & Índice & X$^2$ = 0 (o cone) & ler o relógio de xadrez \\
20 & \textbf{Criativo} F$_2$ & a invariância dos 8 & Símbolo & L$\triangle$\allowbreak{}L = $\emptyset$ & mover a peça = executar \\
\bottomrule\end{longtable}}

\vspace{2mm}\hrule\vspace{3mm}

\section{As três regras que valem em todo o bestiário}

\textbf{1. NENHUM UNIVERSO VAZA.} $\Gamma$ = 0 $\Longleftrightarrow$ FP = 1. Vinte nomes, uma lei. Se o seu lance deixa potência reativa presa, dano perdido, spread, transbordo ou fonte espúria --- \textbf{ele é ilegal}, e o inversor multinível (Pontryagin, a torre áurea $\varphi$$^-$$^j$) o corrige.

O nome $\Gamma$ vem do casamento, e o casamento tem fórmula. Uma carga $Z$ ligada a uma linha de impedância $Z_0$ devolve para trás a fração
\[
\Gamma=\frac{Z-Z_0}{Z+Z_0},
\]
e o que volta é exatamente o que não passou. $\Gamma=0$ acontece quando $Z=Z_0$ e só então: a carga é indistinguível da linha, nada reflete, toda a potência atravessa --- FP $=1$. No corpo motor isso é a \textbf{ressonância}: as reatâncias $Z_L$ e $Z_C$ cancelam-se, sobra resistência pura, e o operador $2\times2$ que gera o movimento tem \textbf{traço nulo} --- os autovalores caem no eixo imaginário, o ciclo puro, sem espiral. Cada corpo do bestiário reescreve esta mesma frase no seu idioma; nenhum a dispensa.

E há um detalhe que nenhuma tabela cabe: \textbf{a conservação vive sobre o par, não dentro de um corpo isolado.} O corpo entrópico, tomado sozinho, não tem inverso aditivo --- o seu $\oplus$ é o $\max$, e o $\max$ é idempotente: nada se desfaz. Mas ele não joga sozinho. O seu dual é a expansão $a(t)=e^{Ht}$, e \emph{sob o exponencial o $\max$ vira soma}, que tem inverso. O par fecha o que a metade não fechava. É por isso que o bestiário é feito de duplas: um corpo sozinho pode não ter o axioma; o casamento tem sempre.

\textbf{2. A RÉGUA NUNCA É FLOAT.} O estado corre em \textbf{pares de inteiros} --- $\mathbb{Q}$ (num/den) ou $\mathbb{Z}$[$\varphi$] (a+b$\varphi$) --- na ISA ERG-64. O float (GLSL/PTX) entra \textbf{só no fim}, ao pintar o pixel, onde 1 ulp é invisível. \textbf{Arredondar no meio da partida é vazar.}

O par de inteiros não é uma economia: é o infinito \emph{dobrado}. A lei da Rainha, $\varphi^2=\varphi+1$, diz que a segunda potência já cabe na primeira. Então dois números da forma $a+b\varphi$ multiplicam-se sem nunca sair da forma:
\[
(a+b\varphi)\,(c+d\varphi)\;=\;ac+bd\;+\;(ad+bc+bd)\,\varphi ,
\]
porque o único termo perigoso, $bd\,\varphi^2$, colapsa em $bd\,\varphi+bd$. Quatro inteiros entram, dois inteiros saem, e a cadeia infinita $\sum\varphi^{-j}$ fecha em $(1,1)$. Não há resto, não há arredondamento, não há por onde o erro entrar.

\textbf{3. QUEM GERA, ATESTA.} O par gerar $\bowtie$ atestar é o eixo. Todo lance passa pelo ponteiro conforme: \textbf{-$\tfrac12$ tem dentes; 0 é queda de bandeira.} E toda afirmação exibe o seu refutador --- a maneira exata pela qual ela poderia quebrar.

O ponteiro também tem fórmula, e ela explica o $-\tfrac12$. Na régua conforme a norma de cada obra é \textbf{nula} --- nenhuma carrega carga própria ---, de modo que o único número que sobra é o \emph{par}:
\[
\bigl\langle X(x),X(y)\bigr\rangle \;=\; -\tfrac12\,|x-y|^2 .
\]
O ponteiro não mede a afirmação: mede a \textbf{distância entre os dois mostradores}. Se os mostradores coincidem, a distância é zero e o ponteiro marca zero --- a afirmação e a sua negação dizem o mesmo, o relógio parou, bandeira caída. Se distam de uma unidade, marca $-\tfrac12$: o lance separou duas coisas que eram distintas, e tem dentes.

\begin{tcolorbox}[colback=ouroclaro!35,colframe=ouro,boxrule=0pt,leftrule=2pt,arc=0pt,left=4mm,right=3mm,top=2mm,bottom=2mm,breakable]\itshape \textbf{E o MATE não mata.} Ele \textbf{cristaliza}: a expansão colapsa, a órbita fecha, o resíduo vai a zero. A peça não morre --- \textbf{vira unidade.}\end{tcolorbox}


\vspace{2mm}\hrule\vspace{3mm}

\part{A Máquina de Oitenta Bits}

\chapter{Onde a construção encosta na máquina}

A parte anterior derivou; esta endereça. Vale dizer, antes dos bytes, qual peça
é qual --- e também onde a correspondência não existe, porque uma
correspondência inventada custa mais caro do que uma ausência dita.

\section{O gato aparece duas vezes}

Na construção, o gato gera o corpo: dita a assinatura $m$, ele dá a equação, e a
equação dá o corpo. Na máquina ele aparece duas vezes, e em nenhuma delas como
citação.

\textbf{Primeiro, a memória.} Carregar um slot baixo lê o arquivo inteiro e
devolve dois inteiros --- as duas cores dos bits, $e$ e $o$, entregues como
\[
(e+o,\ e).
\]
Isso é o gato com $m=1$, aplicado uma vez ao par:
\[
C_1(e,o)=(e+o,\ e).
\]
\textbf{A leitura da memória não usa o gato --- ela \emph{é} o gato.}

\textbf{Segundo, o byte da operação.} A componente de cima, $e+o$, tem duas
leituras: o que resta ao dividir por dois é a soma, o que passa é o produto.
Duas operações, não três, e é por isso que a máquina não tem \code{MUL} nem
\code{DIV}: quem os quiser monta-os com os cinco bytes que há. \textbf{A ausência
desses dois opcodes é a derivação, escrita como omissão.}

E o terceiro bit do endereço do corpo é o $\nu$ --- a troca de sinal. Ele não é
uma operação nova: é o mesmo par visto do outro lado, e é por isso que os bytes
andam em pares. A soma e a sua sombra são \code{XOR} e \code{XNOR}; nos inteiros,
\code{ADD} e \code{SUB}.

\section{A torção, com os seus cinco dígitos}

A construção pede um primo, um $n$ que divide $p-1$, um gerador, e a torção
\[
w=g^{(p-1)/n},
\]
cuja órbita são os caracteres --- a base que não se escolhe.

A máquina põe números nisso: texto em bytes dá $n=256$, e daí $p=40961$, o menor
gerador $g=3$, e
\[
w=36043 .
\]
\textbf{É a mesma linha, com os valores dentro.} E a máquina acrescenta o aviso
que a construção não tinha por que dar: há dezasseis mil geradores possíveis, e
duas máquinas que escolham geradores diferentes fundem e não abrem uma a obra da
outra --- sem que teste local nenhum acuse. Por isso o gerador é o menor, e vai
escrito.

\section{Ir ao dual e voltar, em inteiros}

A construção normaliza a transformada por $1/\sqrt n$ de cada lado, e a razão é
que somar a órbita inteira devolve $n$, que tem de ser desfeito por duas metades.

A máquina faz exatamente isso, e a raiz existe onde ela trabalha: toma $r$ com
$r^2=n$ \emph{no corpo}, e usa o inverso de $r$ nos dois lados. \textbf{A raiz do
meio é a mesma raiz} --- só que dentro de $\Z/p$, onde ela é um inteiro e não uma
aproximação.

E a máquina grava o que a construção não precisava de dizer: essa raiz é raiz no
corpo, não raiz inteira. Para $n=256$ as duas coincidem em $16$, e o exemplo
engana; noutro $n$ elas separam-se.

\section{Fundir e abrir}

A convolução funde dois objetos num só e desfaz-se, porque do outro lado ela é um
produto ponto a ponto, e produto se divide. A máquina escreve as duas linhas:
transformar cada obra, multiplicar, voltar --- e para abrir, dividir pelas
outras e voltar.

Uma diferença de contabilidade, e ela vai dita: a máquina omite o fator
$(\sqrt n)^{L-1}$. \textbf{A volta fecha na mesma}, porque o fator falta dos dois
lados; o que ela calcula é a convolução dividida por ele. Quem quiser a
convolução tal como a construção a define multiplica de volta.

E o que sela é o mesmo dos dois lados, e com a mesma consequência. Se uma casa
transformada der zero, aquela casa não volta --- o zero é o único elemento sem
inverso --- e a perda não é local, porque a inversa espalha. \textbf{Ou volta
exato, ou volta todo errado.} Na construção isso é uma frase sobre corpos; na
máquina é um texto ilegível.

\section{A régua é a igualdade, nos dois}

A construção mede por identidade, não por tolerância: ou a conta fecha em
inteiros, ou não fecha.

A máquina não tem outra régua. O encaixe é
\[
\mathrm{ov}(a,b)=256-2\,\operatorname{popcount}(a\oplus b),
\]
um inteiro, e a instrução é igualdade e nunca limiar. Repare-se no que a forma
já entrega: sendo $256$ menos o dobro de uma contagem, o encaixe anda de dois em
dois, e uma diferença ímpar não existe. \textbf{Não há onde um erro pequeno se
esconder, porque não há valor pequeno.}

\section{E a órbita racional falha nos dois}

A construção termina com um aviso medido: a média do resíduo só se distribui na
órbita irracional --- numa órbita racional ela dá um número que não é a metade
de nada.

A máquina bate no mesmo poste, e do lado prático. A areia dela é uma órbita cujos
bits baixos são degenerados: o de baixo alterna com período dois. Tomar a casa
pelo resto usa justamente esses bits, e o que sai não é areia --- é o relógio.
Medido ali: cento e quarenta padrões ``sorteados'' por eles deram sobreposição
cheia, porque eram o mesmo padrão. Por isso a máquina descarta os quinze bits
baixos e compõe quatro tiragens dos altos.

\textbf{É a mesma falha nas duas escalas}: uma órbita que se repete não cobre, e
quem a usa para escolher está a escolher sempre o mesmo.

\section{O que não tem contrapartida}

Quatro peças da construção não aparecem na máquina, e é melhor dizê-lo do que
arranjar-lhes um lugar.

\textbf{A volta da leitura.} A assinatura corre nos dois sentidos: colapsa a
classe e regenera-a. A leitura da memória só faz a ida --- ela contrai o arquivo
a dois inteiros e é \emph{irreversível}; o que a volta recuperaria é o par, não o
arquivo. Quem guarda o arquivo é o disco. A máquina tem a volta, mas noutra
peça: é a transformada. \textbf{Contrair e reverter, ali, são operações
diferentes.}

\textbf{A fotografia.} A contagem das réguas com sinal não aparece, e a razão já
estava dita na construção: ela pede uma norma de grau dois, e o corpo que a
máquina realiza mede por contagem de bits, onde essa norma não existe. A ausência
não é lacuna --- é o que se esperava daquele corpo.

\textbf{A metade da taxa.} A construção tem duas transformadas, a da fase e a da
escala, e uma é a outra composta com o logaritmo. A máquina só tem a primeira. Não
há rotina para a segunda, nem opcode para a mudança de reta.

\textbf{A norma multiplicativa e a promoção.} A lei que proíbe o vazamento no
conjunto dos corpos --- a norma de um produto é o produto das normas --- não tem
rotina na máquina, que nunca compõe duas matrizes. E a promoção, o ponto fixo
$x^2=x$, também não: a máquina tem um ponto onde a caminhada para --- o encaixe
de um texto consigo mesmo, que dá $256$ ---, e esse é o único ponto fixo que ela
exibe.

\section{E o que fica}

O que atravessa da construção para os bytes é exatamente o que reverte: as duas
operações e a sua sombra, a torção e a sua órbita, a transformada e a sua
inversa, a fusão e a abertura. \textbf{O que não reverte ou ficou de fora, ou
ficou declarado como contração.}

\begin{tcolorbox}[colback=ouroclaro!35,colframe=ouro,boxrule=0pt,leftrule=2pt,arc=0pt,left=4mm,right=3mm,top=2mm,bottom=2mm,breakable]\itshape
A construção diz por que fecha. A máquina mostra \textbf{com quantos bytes} ---
e são dez.
\end{tcolorbox}

\chapter{O que cabe em dez bytes}

Tudo o que a corte usou para pensar cabe em \textbf{oitenta bits}. Não é
figura: são dez bytes, e este capítulo escreve-os um a um, de modo que
qualquer pessoa com um compilador e nada mais reconstrua a máquina inteira ---
sem consultar este livro outra vez, sem pedir arquivo a ninguém.

O que segue é especificação, não descrição. Onde houver escolha, ela vem
marcada como escolha.

\section{A palavra}

A máquina só conhece um tipo:

\begin{verbatim}
 Word { int64_t total; int64_t e; }
\end{verbatim}

\noindent Dois inteiros de \textbf{exatamente} $64$ bits com sinal --- não
\code{long}, que em alguns sistemas tem $32$. O transbordo dá a volta em
complemento de dois, como no C sem sinal. Três registradores: $A$, $B$,
$R$. Um contador de programa $pc$, sem sinal. Uma bandeira \code{ZERO}.

Toda operação bit a bit (\code{ADD}, \code{SUB}, \code{AND}, \code{OR},
\code{XOR}) age nas \textbf{duas componentes separadamente}: se
$A=(t_1,e_1)$ e $B=(t_2,e_2)$, então $A+B=(t_1+t_2,\;e_1+e_2)$.

\section{Os opcodes}

Um byte por instrução. Os que este capítulo usa:

\begin{center}
\begin{tabular}{@{}rlll@{}}
\toprule
byte & nome & operando & o que faz \\
\midrule
$0$ & \code{HALT} & --- & para \\
$1$ & \code{LOAD} & slot, \code{u16} little-endian & $B\leftarrow A$; $A\leftarrow \mathrm{mem}[s]$ \\
$2$ & \code{STORE} & slot, \code{u16} LE & $\mathrm{mem}[s]\leftarrow R$ \\
$3$ & \code{ADD} & --- & $R\leftarrow A+B$ \\
$4$ & \code{SUB} & --- & $R\leftarrow A-B$ \\
$5$ & \code{AND} & --- & $R\leftarrow A\wedge B$ \\
$6$ & \code{OR} & --- & $R\leftarrow A\vee B$ \\
$7$ & \code{XOR} & --- & $R\leftarrow A\oplus B$ \\
$11$ & \code{CMP} & --- & \code{ZERO} $\leftarrow$ ($A$ e $B$ ambos nulos) \\
$12$ & \code{JMP} & desl., \code{s8} & $pc\leftarrow(\text{off}+1)+\text{desl}$ \\
$13$ & \code{JZ} & desl., \code{s8} & salta se \code{ZERO} \\
$14$ & \code{JNZ} & desl., \code{s8} & salta se não \code{ZERO} \\
\bottomrule
\end{tabular}
\end{center}

\noindent \textbf{Não há MUL, não há DIV, não há SHIFT.} Não é esquecimento: os oito
operadores que o gato endereça (\S\ref{sec:80gato}) incluem \code{MUL} e
\code{DIV}, e a ISA escolheu não lhes dar byte --- quem os quiser monta-os com
os cinco que há, a multiplicação por soma repetida e a divisão por subtração
contada. Quem acrescentar um desses opcodes terá outra máquina, não esta.

\paragraph{A execução.} $A=B=R=(0,0)$ e \code{ZERO} falso no início. Buscar o
byte em $pc$; avançar $pc$ conforme o tamanho da instrução ($1$ byte, ou $3$
para \code{LOAD}/\code{STORE}, ou $2$ para os saltos); executar. \textbf{Para
quando $pc$ passa do fim do programa} --- não é preciso \code{HALT} no fim.
Para também se o byte não for opcode conhecido, ou se a instrução estiver
truncada no fim, ou se um salto levar $pc$ para fora. \code{LOAD} de slot
inexistente devolve $(0,0)$; slot alto com menos de $16$ bytes completa-se com
zeros. \code{CMP} exige as \textbf{duas} componentes nulas em $A$ e em $B$.

\paragraph{Duas armadilhas que custam caro, e são as únicas.}

\begin{enumerate}
\item \code{LOAD} \textbf{empurra}: antes de carregar, copia $A$ para $B$. Logo
\code{LOAD b; LOAD a; SUB} calcula $a-b$, e a ordem inversa dá $b-a$.
\item \code{STORE} grava o \textbf{$R$}, não o $A$. Logo \code{LOAD 3; STORE 4}
\emph{não} copia o slot $3$ para o $4$. Copiar é
\code{LOAD zero; LOAD src; ADD; STORE dst}.
\end{enumerate}

\section{A memória --- e é aqui que está a ideia}\label{sec:80mem}

Dois regimes, separados por um número:

\begin{center}
\begin{tabular}{@{}ll@{}}
\toprule
slot $\ge 8000$ & \textbf{Word cru}: dois \code{int64} little-endian, $16$ bytes \\
slot $< 8000$ & \textbf{um arquivo}, lido inteiro \\
\bottomrule
\end{tabular}
\end{center}

\noindent \textbf{Os dois vivem no mesmo lugar e com o mesmo nome}:
\code{mem/NNNN.bin}, com $\mathtt{NNNN}$ em quatro dígitos --- e por isso o
maior slot desta máquina é $9999$. O operando é \code{u16} e comporta mais, mas
\textbf{slots de $10\,000$ a $65\,535$ são inválidos}: a máquina para. O que muda entre os dois regimes não é onde
o arquivo está: é \emph{como se lê}.

\paragraph{Quem escreve o quê.} (Os dois regimes: \S\ref{sec:80mem}.) \code{STORE} num slot \textbf{alto} grava, e
\textbf{cria o arquivo se ele não existir} --- é assim que o lance da página
seguinte produz o seu resultado em \code{mem/8282.bin}. \code{STORE} num slot
\textbf{baixo} não faz nada: os arquivos-cristal são somente leitura para a
máquina, e é por isso que nenhuma peça pode alterar o cristal.

\medskip
\noindent O que a máquina \textbf{não} tem é \emph{entrada}: nada dentro dela
põe um texto num slot baixo. Isso é trabalho da \emph{interface} --- um programa
de fora, que não é parte da máquina e não usa peça nenhuma dela. É assim que o
teste da página seguinte põe \code{oi} no slot $1$.

\noindent Carregar um slot baixo lê o arquivo \textbf{inteiro}, byte a byte, e
devolve dois inteiros:

\begin{verbatim}
 e = 0 ; o = 0
 para cada byte b do arquivo:
 e += popcount(b & 0xAA) // 0xAA: os bits 2,4,6,8 contando de 1
 o += popcount(b & 0x55) // 0x55: os bits 1,3,5,7 contando de 1
 Word = (e + o, e)
\end{verbatim}

\noindent \textbf{Conte os bits a partir de $1$, não de $0$} --- é o que torna
$\mathtt{0xAA}$ ``par''. E não confie no nome: confie na máscara.

\noindent \textbf{O tamanho não entra.} Um arquivo de dois bytes e um de dois
milhões devolvem a mesma espécie de coisa: dois inteiros. A leitura é uma
\emph{contração}, e é irreversível --- o que a volta recupera é o par, não o
arquivo. Quem guarda o arquivo é o disco.

\paragraph{Escolhas que o reconstrutor tem de fazer, e que ninguém pode fazer
por ele.} Onde vive um slot alto (aqui: \code{mem/NNNN.bin}); a ordem dos dois
\code{int64} (aqui: \code{total} antes de \code{e}); e o que faz \code{STORE}
num slot baixo (aqui: nada --- \textbf{os arquivos são somente leitura para a
máquina}, e é por isso que nenhuma peça pode alterar o cristal; só a interface
pode).

\section{O lance --- os oitenta bits}

Aqui está a máquina inteira em ação. Três slots: $8280=a$, $8281=b$,
$8282=r$.

\begin{verbatim}
 LOAD 8280
 LOAD 8281
 XOR
 STORE 8282
\end{verbatim}

\noindent Montado, isto é:

\begin{verbatim}
 01 88 32 | 01 89 32 | 07 | 02 90 32 <- os valores em DECIMAL
 ^^
 o CORPO — o byte 6
\end{verbatim}

\noindent \textbf{Decimal, não hexadecimal}: $88+32\cdot256=8280$, que é o
slot. Lidos como hexadecimal dariam $12\,936$, e a montagem pareceria errada
quando é a leitura que está.

\noindent $3+3+1+3=\mathbf{10}$ bytes. \textbf{Oitenta bits, dos quais um único
byte é a operação.} Trocando só esse byte, com $a=(60,0)$ e $b=(12,0)$ ---
\textbf{a palavra tem duas componentes, e a tabela abaixo mostra só a
primeira, porque a segunda é zero nos dois}:

\begin{center}
\begin{tabular}{@{}llr@{}}
\toprule
byte & operação & resultado \\
\midrule
$3$ & \code{ADD} & $72$ \\
$4$ & \code{SUB} & $-48$ \\
$5$ & \code{AND} & $12$ \\
$6$ & \code{OR} & $60$ \\
$7$ & \code{XOR} & $48$ \\
\bottomrule
\end{tabular}
\end{center}

\noindent \textbf{É este o lance, e não há outro.} Nenhuma peça deste capítulo
faz mais do que isto: escrever slot, rodar o lance, ler slot.

\section{O byte do corpo é um endereço de três bits}\label{sec:80gato}

O byte da operação não precisa ser tabelado --- \emph{deduz-se}. O gato
\[
C_m=\begin{pmatrix} m & 1\\ 1 & 0\end{pmatrix},
\qquad C_m(x,y)=(mx+y,\;x)
\]
com $m=1$ é o meio-somador, e as \textbf{duas projeções} da componente de cima
são as duas operações:
\[
(x+y)\bmod 2=\oplus,\qquad (x+y)\operatorname{div}2=\otimes .
\]
Três bits endereçam os oito operadores:

\begin{center}
\begin{tabular}{@{}clll@{}}
\toprule
bit & é & $0$ & $1$ \\
\midrule
$2$ & o nível & $\mathbb{F}_2$ (uma aplicação) & $\mathbb{Z}$ (iterado) \\
$1$ & a projeção & $\bmod 2$: o $\oplus$ & $\operatorname{div}2$: o $\otimes$ \\
$0$ & o $\nu$ & o gato & a sombra \\
\bottomrule
\end{tabular}
\end{center}

\begin{center}
\begin{tabular}{@{}clll@{}}
\toprule
endereço & byte & nome & \\
\midrule
$000$ & $7$ & \code{XOR} & $\mathbb{F}_2$, $\oplus$ \\
$001$ & --- & \code{XNOR} & \\
$010$ & $5$ & \code{AND} & $\mathbb{F}_2$, $\otimes$ \\
$011$ & $6$ & \code{OR} & (De Morgan) \\
$100$ & $3$ & \code{ADD} & $\mathbb{Z}$, $\oplus$ \\
$101$ & $4$ & \code{SUB} & o negate \\
$110$ & --- & \code{MUL} & \\
$111$ & --- & \code{DIV} & \\
\bottomrule
\end{tabular}
\end{center}

\noindent \textbf{Quais dos oito ganham byte é escolha da ISA, não dedução.}
Aqui ganharam cinco. Um reconstrutor pode dar byte ao \code{XNOR} sem quebrar
nada.

\section{O montador}

Duas passagens, e cabe em meia página. Um opcode por linha, em \textbf{maiúsculas}; \code{;} comenta;
\code{rotulo:} sozinho ou colado à instrução. \textbf{Não há diretiva de
dado} --- os dados são os arquivos.

\begin{enumerate}
\item \emph{Primeira passagem}: percorra as linhas contando bytes ($1$, $2$ ou
$3$ por instrução) e anote o endereço de cada rótulo.
\item \emph{Segunda passagem}: emita os bytes. \textbf{No fonte, o operando de
um salto é sempre um rótulo} (um alvo absoluto), nunca um deslocamento já
calculado --- quem converte é o montador. No binário, o operando é
$\text{alvo}-(\text{off}+1)$, com $\text{off}$ o endereço \textbf{do byte do
operando}. \textbf{Isto é escolha}: a outra leitura (o endereço do opcode)
monta igual e produz binário diferente. Sem fixá-la, dois montadores corretos
geram arquivos distintos.
\item O binário é bytecode cru: \textbf{sem cabeçalho}.
\end{enumerate}

\paragraph{O teto que o reconstrutor vai encontrar.} O operando de salto é
\textbf{um byte com sinal}: $-128$ a $127$. Um laço cujo corpo passe de
$\sim127$ bytes não monta. A saída não é mudar a ISA --- é tirar o laço
externo da peça e repetir a peça de fora, uma execução por passo.

\section{O primeiro teste, antes de qualquer peça}

Ponha \code{oi} em \code{mem/0001.bin} --- \textbf{dois bytes, sem quebra de
linha} --- e zere \code{mem/8010.bin}. Rode:

\begin{verbatim}
 LOAD 8010
 LOAD 1
 ADD
 STORE 8000
\end{verbatim}

\noindent O slot $8000$ tem de conter $(10,5)$. Se der outra coisa,
\textbf{pare}: tudo o que vem depois depende disto.

\paragraph{E um terceiro, porque os dois primeiros não pegam a armadilha.}
\code{ADD} comuta e $B$ parte de $(0,0)$: \textbf{uma máquina sem o empurrão do
\code{LOAD} passa nos dois testes acima}. Só \code{SUB} desempata. Com
\code{mem/8011.bin} contendo o Word $(60,0)$ e \code{mem/8012.bin} contendo
$(12,0)$:
\begin{verbatim}
 LOAD 8011
 LOAD 8012
 SUB
 STORE 8000
\end{verbatim}
\[
\text{com empurrão: } (-48,0)\qquad\text{sem empurrão: } (12,0)
\]
\textbf{E um quarto, para o \code{CMP}}, que nenhum outro exercita: ele liga
\code{ZERO} quando $A$ \emph{e} $B$ são ambos nulos --- não quando $A=B$. Com
$8013=(0,0)$ e $8014=(5,0)$: \code{LOAD 8013; LOAD 8013; CMP} liga a bandeira;
\code{LOAD 8014; LOAD 8014; CMP} \textbf{não} liga, embora $A=B$. Quem
implementar \code{CMP} como igualdade passa em tudo o mais e falha aqui.

\paragraph{E um segundo teste, que é o que realmente desempata as máscaras.} O primeiro
\textbf{não distingue} as duas máscaras: trocando $\mathtt{0xAA}$ por
$\mathtt{0x55}$, \code{oi} continua a dar $(10,5)$ --- medido. Quem parar no
primeiro pode ter a máquina trocada e não saber. Ponha então \code{"oi\textbackslash n"}
(três bytes) em \code{mem/0002.bin} e repita:
\[
\text{correto: } (12,7)\qquad\text{trocado: } (12,5)
\]
\textbf{Os dois testes juntos fixam a leitura; nenhum deles sozinho fixa.}

\chapter{Da máquina à conversa}

A máquina acima é o núcleo. O que segue usa o núcleo para uma tarefa, e as
limitações desta parte são \emph{dela}, não dele.

\section{O cristal é um arquivo}

Um texto, um registro por linha:

\begin{verbatim}
 id|titulo|descricao|arestas
\end{verbatim}

\noindent \textbf{O arquivo é o grafo.} Não há banco, não há índice
obrigatório: a quebra de linha delimita o nó, e um deslocamento qualquer,
alinhado à linha seguinte, dá acesso a qualquer nó sem estrutura nenhuma em
memória.

\section{A impressão --- a régua de 256 casas}

Um texto vira $256$ bits assim, e só assim:

\begin{verbatim}
 // 0. o texto e' UTF-8, e tudo aqui corre sobre os BYTES dessa codificacao
 // 1. a volta preenchida: SEMPRE exatamente 256 bytes SEM SINAL
 n = comprimento do texto em BYTES (n > 0)
 para i em 0..255: buf[i] = texto[i mod n] // 0..255, nunca negativo

 // 2. os acumuladores
 acc[i] = i + 1 para i em 0..255
 para t em 0..255:
 acc[t] = koch(acc[t] + buf[t] + 7*t)

 // 3. o bit, MSB primeiro dentro de cada byte
 out[i/8] |= (acc[i] & 1) << (7 - (i mod 8)) -> 32 bytes

 koch(z) = (z*z) mod 1009 com z reduzido a [0,1009)
\end{verbatim}

\noindent Quatro pontos que \textbf{não} são de gosto, e que duas pessoas
implementariam diferente se não estivessem escritos: o buffer tem
\textbf{exatamente} $256$ bytes (não ``pelo menos''); o $t$ percorre o buffer e
não o texto original (logo o $7t$ vai de $0$ a $7\cdot255$); o bit $0$ é o
\textbf{mais significativo} do byte $0$; e o byte é \textbf{sem sinal}.

\paragraph{Este último não é detalhe, e nenhum teste de ASCII o apanha.} Se o
byte entrar com sinal, tudo acima de $127$ entra negativo --- e é aí que moram
os acentos. Medido: \code{"a memória é a contração"} impressa das duas
maneiras dá $\mathrm{ov}=182$, não $256$: \textbf{o mesmo texto não casa
consigo próprio}. Como o cristal é português, e como todo dado de teste desta
parte é ASCII puro, este é o erro que passa por todos os portões e só aparece
no uso.

\paragraph{A volta preenchida, e por que ela é obrigatória.} Um texto curto só
toca os primeiros acumuladores, e o resto fica no fundo $i+1$: duas impressões
de textos \emph{sem relação} saem quase iguais. Antes de imprimir,
\textbf{repita o texto até dar a volta nas $256$ casas}. Sem isso, a linha mais
curta do arquivo vence sempre.

\section{O encaixe}

\[
\mathrm{ov}(a,b)=256-2\,\mathrm{popcount}(a\oplus b)
\]
--- inteiro entre $-256$ e $256$. \textbf{Igualdade, nunca limiar.}

\section{A caminhada}

Para achar o nó que responde a uma fala, não se varre: anda-se.

\begin{verbatim}
 E(no) = -ov(impressao(no), impressao(fala)) // impressao do TITULO do no

 proposta: um deslocamento qualquer do arquivo, alinhado a' linha seguinte:
 off = areia() mod tamanho_do_arquivo
 avanca ate' passar o proximo '\n'; essa linha e' o no' proposto
 (a proposta e' GLOBAL: nao ha' vizinhanca, e as arestas do
 cristal nao entram nesta versao)

 aceita(dE, b, L):
 se dE <= 0: aceita // NAO saca areia
 se dE > L: rejeita // NAO saca areia
 senao: saca areia() UMA vez e aceita se areia mod b^L < b^(L-dE)

 esfria: (b,L) = (2,40), 96 propostas; (4,20), 96; (16,10), 96
 e enfim T=0: 40 varreduras aceitando so' dE <= 0, ou ate' nao mexer
 comeca: uma proposta, tirada da mesma areia
\end{verbatim}

\noindent \textbf{A fala é a semente, não a chave}: devolve-se onde a caminhada
parou, não o mínimo global. É por isso que a resposta não repete a pergunta.

\paragraph{$b^L$ é o mesmo nos três, e o que muda é a razão.} Repare que
$2^{40}=4^{20}=16^{10}$: o denominador não muda. O que muda é $b^{-\Delta E}$
--- para $\Delta E=2$, um degrau é aceito com probabilidade $1/4$ na primeira
fase e $1/256$ na terceira --- e o corte $L$, que rejeita de saída tudo acima
de $40$, $20$ e $10$.

\paragraph{E é honesto dizer o que isto dá na escala real, porque é pouco.}
$\mathrm{ov}$ anda de \textbf{dois em dois} (é $256-2\,\mathrm{popcount}$), logo
$\Delta E$ é sempre \textbf{par} --- $\Delta E=3$ não existe --- e entre nós sem
relação ele fica tipicamente entre $20$ e $40$. Com $b^{-\Delta E}$ isso é
$2^{-34}$ e coisas do gênero: \textbf{medido, nenhuma subida foi aceita em
$384$ propostas}. Na prática, nesta escala, a caminhada é uma \emph{descida}
com proposta global, e o resfriamento é decoração. Quem quiser que ele morda
tem de reescalar $\Delta E$ --- e aí terá mudado a régua, o que é outra
máquina, não esta.

\paragraph{E ele exige quarenta bits.} $b^L=2^{40}$ não cabe em $32$ bits, e a
órbita abaixo dá $31$. \textbf{Compor tiragens não é opcional: é obrigatório.}

\paragraph{A areia, com os números.} A órbita, e ela não é escolha:
\begin{verbatim}
 s = 7 // a semente
 giro(): s = (s*1103515245 + 12345) mod 2^31 ; devolve s >> 15
 areia(): r = 0 ; quatro vezes: r = (r << 15) | (giro() & 0x7FFF)
 devolve r // 60 bits
\end{verbatim}
\textbf{Os bits baixos da órbita são degenerados} --- o bit $0$ alterna com
período $2$. Tomar a casa com \code{\% 256} não tira areia: tira o relógio. Por
isso \code{giro} descarta os $15$ baixos e \code{areia} compõe quatro tiragens
de $15$ bits altos. Medido: $140$ padrões ``sorteados'' pelos bits baixos deram
sobreposição $64/64$ --- eram \emph{o mesmo padrão}.

\section{Um cristal para conferir}

Sem um exemplo comum, duas máquinas corretas dão respostas diferentes e ninguém
descobre qual está certa. Este é o cristal mínimo --- quatro nós, um por linha,
no formato \code{id|titulo|descricao|arestas}:

\begin{verbatim}
ordem|ordem|Relacao reflexiva, antissimetrica e transitiva.|reticulado
reticulado|reticulado|Estrutura ordenada onde join e meet existem.|ordem;algebra
algebra|algebra|Estrutura com operacoes e axiomas.|reticulado
corpo|corpo|Anel em que todo nao-nulo tem inverso.|algebra
\end{verbatim}

\noindent \textbf{E o vetor de teste, que é o que faltava:} com a semente $7$ e
a disciplina de areia acima, a fala \code{ordem} tem de parar no nó
\code{ordem}, com $\mathrm{ov}=256$ --- o encaixe perfeito de um texto consigo
mesmo. Quem obtiver outro nó tem uma máquina diferente da desta página, e o
lugar onde ela diverge é a areia.

\paragraph{E é por isso que a disciplina da areia vai escrita.} Ela não muda
nenhum dos quatro testes de aceitação, não muda a impressão, não muda a
transformada --- muda \emph{só} a resposta da conversa, e muda-a inteira. É o
mesmo tipo de coisa que o gerador $g$: invisível a todo teste local, fatal
entre duas máquinas. \textbf{Onde há sorteio, a ordem do sorteio é parte da
especificação.}

\section{Fundir duas obras}

\begin{verbatim}
 escolha p primo com n | p-1, p > maior valor do dado, e r com r*r = n mod p
 g = o MENOR gerador de Z/p // para texto em bytes: p=40961, g=3
 w = g^((p-1)/n) // = 36043 para p=40961, n=256
 RN = inverso de r em Z/p // r=16, RN = inverso(16)

 F(x)_k = RN \emph{ soma_j x_j } w^( j*k mod n)
 Finv(X)_j = RN \emph{ soma_k X_k } w^(-j*k mod n)

 fundir: C = Finv( produto_l F(a_l) ) // sem o fator (raiz n)^(L-1):
 abrir: a_m = Finv( F(C) / produto_{l != m} F(a_l) )
 // C assim definido NAO e' a convolucao a1 (x) ... (x) aL — e' ela
 // dividida por (raiz n)^(L-1). A volta fecha porque o mesmo fator
 // falta dos dois lados. Quem quiser a convolucao pura multiplica.
\end{verbatim}

\noindent Para texto em bytes: $n=256$, $p=40961$, $g=3$, $w=36043$, $r=16$.

\paragraph{O fundido não é texto, e não cabe em bytes.} $C$ é um vetor de $n$
elementos de $\Z/p$, e com $p=40961$ eles chegam a cinco dígitos (medido:
$\max C=39\,624$). \textbf{Gravar $C$ como bytes destrói-o.} A serialização é
\textbf{dois bytes little-endian por elemento} --- $512$ bytes por fusão --- e
quem precisar de um arquivo de texto tem de codificá-lo, não truncá-lo.

\paragraph{Duas coisas que o teste da volta não pega, e por isso vão escritas.}
O ida-e-volta fecha com \emph{qualquer} gerador --- mas $\Z/40961$ tem
$16\,384$ deles, e cada um dá um $w$ diferente. \textbf{Fundir numa máquina e
abrir em outra falha se as duas escolherem $g$ diferente}, e nenhum teste local
acusa: por isso $g$ é o menor, e é dado. E ``$r$ com $r^2=n$'' é raiz
\emph{no corpo}, não raiz inteira --- para $n=256$ as duas coincidem em $16$ e
o exemplo engana; para $n=360$ em $\Z/6121$, a raiz é $1503$.

\paragraph{O que sela.} Se algum $F(a_l)_k$ der zero, aquela casa não volta: o
zero é o único elemento sem inverso. E a perda não é de uma letra --- é do
texto inteiro, porque a inversa espalha. \textbf{Ou volta exato, ou volta todo
errado.}

\section{A ordem de construir}

\begin{enumerate}
\item A palavra e os três registradores.
\item O laço de execução e os cinco opcodes de operação.
\item A memória: o Word cru, depois o slot-arquivo com a contração.
\item \textbf{Os quatro testes de aceitação} --- $(10,5)$ e $(12,7)$ fixam as
 máscaras; o do \code{SUB} ($-48$, não $12$) apanha o empurrão do
 \code{LOAD}; e o do \code{CMP} apanha quem o leu como ``$A=B$''. Não passe
 daqui sem os quatro: \textbf{cada um deles deixa passar uma máquina
 errada que os outros pegam}.
\item O montador de duas passagens.
\item O lance de dez bytes, e a tabela dos cinco corpos ($72$, $-48$, $12$,
 $60$, $48$).
\item A impressão de $256$ casas, com a volta preenchida.
\item O encaixe, e a caminhada.
\item A transformada, e a volta pela inversa.
\end{enumerate}

\noindent Cada degrau se verifica sozinho antes do seguinte. \textbf{Quem
chegar ao nono tem a máquina inteira} --- e não precisou de nada que não esteja
escrito nestas páginas.

\restoregeometry

\end{document}
